\PassOptionsToPackage{hyperfootnotes=false}{hyperref}
% Set-up of the Spanish translation of EGA III, first part:
% document class and packages, statement environments, headings and contents,
% sheaf operators, and the commands for page marks and cross-references.
% Written for these editions and dedicated to the public domain under CC0 1.0.
\documentclass[oneside]{amsart}

% Commutative diagrams.
\usepackage[all]{xy}
\usepackage{tikz-cd}
% T1 text encoding.
\usepackage[T1]{fontenc}
% String handling and argument parsing for the reference commands below;
% colour.
\usepackage{xstring,xparse,xcolor}
% Coloured links: cross-references in deep red, citations in blue; contents
% entries linked by their page numbers.
\definecolor{egalinkred}{rgb}{0.76,0.13,0.28}
\usepackage[
  linktocpage=true,
  colorlinks=true,
  hyperindex,
  hypertexnames=false,
  citecolor=blue,
  linkcolor=egalinkred
]{hyperref}
% A4 paper; margins of 1.25in at the sides and 0.75in at head and foot.
\usepackage[a4paper,hmargin=1.25in,vmargin=0.75in]{geometry}
% Palatino for text and mathematics, AMS symbols and Ralph Smith's script
% alphabet; Palatino is set with its customary slightly opened leading.
\usepackage{amssymb,mathrsfs,mathpazo}\linespread{1.05}
% Headings are not left at the foot of a page; margin notes; lists, long
% tables and multiple columns; the Aurical hand (for sheaf operators);
% microtypography.
\usepackage[nobottomtitles]{titlesec}
\usepackage{marginnote}
\usepackage[shortlabels]{enumitem}
\usepackage{longtable,multicol,aurical,microtype}

% Statements. EGA numbers its statements in the text, so the environments
% are unnumbered and take the number as optional argument. For instance
% \begin{proposition}[2.1.3] starts with the bold heading "Proposition
% (2.1.3). ---" (in the language of the edition) and sets the statement in
% italics; definitions, examples, remarks and notation are set in roman;
% \begin{env}[2.1.4] starts a numbered paragraph with a bold "(2.1.4).".
\newtheoremstyle{egaitem}{}{}{\rmfamily}{}{\bfseries}{.}{ }{\thmnote{(#3)}}
\newtheoremstyle{egastatement}{}{}{\itshape}{}{\bfseries}{. --- }{ }{\thmname{#1}\thmnote{ (#3)}}
\newtheoremstyle{egadescription}{}{}{\rmfamily}{}{\bfseries}{. --- }{ }{\thmname{#1}\thmnote{ (#3)}}
\makeatletter
\newcommand*{\ega@statement}[3]{\theoremstyle{#1}\newtheorem*{#2}{#3}}
\ega@statement{egaitem}{env}{---}
\ega@statement{egastatement}{theorem}{Teorema}
\ega@statement{egastatement}{proposition}{Proposición}
\ega@statement{egastatement}{lemma}{Lema}
\ega@statement{egastatement}{corollary}{Corolario}
\ega@statement{egadescription}{definition}{Definición}
\ega@statement{egadescription}{example}{Ejemplo}
\ega@statement{egadescription}{examples}{Ejemplos}
\ega@statement{egadescription}{remark}{Observación}
\ega@statement{egadescription}{remarks}{Observaciones}
\ega@statement{egadescription}{scholium}{Escolio}
\ega@statement{egadescription}{notation}{Notación}
\makeatother

\renewcommand{\proofname}{Demostración}
\renewcommand{\contentsname}{Índice general}
\renewcommand{\refname}{Referencias}

% Contents and headings. Subsection entries of the contents are indented by
% 2.5pc, with 2.2pc for the number. Section headings are amsart's centred small
% capitals, in bold at \Large size. Heading numbers end with a full stop and a
% tie, and a section number is preceded by the sign \S.
\makeatletter
\renewcommand*{\l@subsection}{\@tocline{2}{0pt}{2.5pc}{2.2pc}{}}
\renewcommand*{\section}{%
  \@startsection{section}{1}{\z@}{.7\linespacing\@plus\linespacing}%
  {.5\linespacing}{\normalfont\bfseries\Large\scshape\centering}}
\@namedef{ega@numbersign@section}{\textsection}
\newcommand*{\ega@numbersign}[1]{%
  \ifcsname ega@numbersign@#1\endcsname\csname ega@numbersign@#1\endcsname\fi}
\renewcommand*{\@seccntformat}[1]{\ega@numbersign{#1}\csname the#1\endcsname.~}
\makeatother

% Sheaves are written in the script alphabet of mathrsfs.
\renewcommand*{\mathcal}{\mathscr}
% Notation of the EGA editions: the abbreviations the transcription uses for alphabets, symbols, operators and
% arrows. Written for these editions and dedicated to the public domain under CC0 1.0.
\makeatletter
% a command that may or may not exist already is (re)defined in one step
\newcommand{\ega@define}[2]{\providecommand{#1}{}\renewcommand{#1}{#2}}
\newcommand{\ega@operator}[2]{\providecommand{#1}{}\renewcommand{#1}{\operatorname{#2}}}

% alphabets: sheaves, bold letters, categories
\ega@define{\sh}{\mathcal}
\ega@define{\bb}{\mathbf}
\ega@define{\cat}{\mathtt}
\ega@define{\C}{\cat{C}}
\ega@define{\Set}{\cat{Set}}
\ega@define{\CHom}{\cat{Hom}}
\ega@define{\op}{^{\cat{op}}}
\ega@define{\GL}{\bb{GL}}
\ega@define{\bbGamma}{\boldsymbol{\Gamma}}

% letters and symbols
\renewcommand{\leq}{\leqslant}
\renewcommand{\geq}{\geqslant}
\ega@define{\setmin}{-}
\ega@define{\emp}{\varnothing}
\ega@define{\vphi}{\phi}
\ega@define{\kres}{k}
\ega@define{\rad}{\mathfrak{r}}
\ega@define{\nilrad}{\mathfrak{N}}
\ega@define{\HH}{\mathrm{H}}
\ega@define{\CHH}{\check{\mathrm{H}}}
\ega@define{\RR}{\mathrm{R}}
\ega@define{\LL}{\mathrm{L}}
\ega@define{\red}{\mathrm{red}}
\providecommand{\dual}{}\renewcommand{\dual}[1]{{#1}^{\vee}}
% the tilde of the sheaf associated with a module, written after it: M\supertilde
\ega@define{\supertilde}{{\,\widetilde{\,}\,}}

% operators
\ega@operator{\rg}{rg}
\ega@operator{\gr}{gr}
\ega@operator{\Gr}{Gr}
\ega@operator{\Hom}{Hom}
\ega@operator{\Proj}{Proj}
\ega@operator{\Tor}{Tor}
\ega@operator{\Ext}{Ext}
\ega@operator{\Supp}{Supp}
\ega@operator{\Ker}{Ker}
\ega@operator{\Im}{Im}
\ega@operator{\Coker}{Coker}
\ega@operator{\Spec}{Spec}
\ega@operator{\Spm}{Spm}
\ega@operator{\Spf}{Spf}
\ega@operator{\grad}{grad}
\ega@operator{\dimc}{dimc}
\ega@operator{\codim}{codim}
\ega@operator{\id}{id}
\ega@operator{\Der}{Der}
\ega@operator{\Diff}{Diff}
\ega@operator{\Hyp}{Hyp}
\ega@operator{\Tr}{Tr}

% arrows: long in displays, short in text, scripts and subscripts
\newcommand{\ega@arrow}[2]{\mathchoice{#1}{#2}{#2}{#2}}
\let\ega@shortmapsto\mapsto
\renewcommand{\to}{\ega@arrow{\longrightarrow}{\rightarrow}}
\providecommand{\from}{}\renewcommand{\from}{\ega@arrow{\longleftarrow}{\leftarrow}}
\renewcommand{\mapsto}{\ega@arrow{\longmapsto}{\ega@shortmapsto}}
\ega@define{\isoto}{\simeq}
\ega@define{\simto}{\xrightarrow{\sim}}

% wide accents
\newcommand{\ega@wide}[2]{\renewcommand{#1}[1]{#2{##1}}}
\ega@wide{\hat}{\widehat}
\ega@wide{\tilde}{\widetilde}
\makeatother


% Sheaf operators (\shHom, \shExt, ...): a script capital followed by the
% rest of the name in slanted Aurical, moved 2.2pt towards the capital. The
% middle argument of \ega@sheafop is the space put after the capital.
\makeatletter
\newcommand*{\ega@sheafop}[3]{\sh{#1}#2\textup{\kern-2.2pt{\Fontauri\slshape #3}}}
\newcommand*{\shHom}{\ega@sheafop{H}{}{om}}
\newcommand*{\shProj}{\ega@sheafop{P}{}{roj}}
\newcommand*{\shExt}{\ega@sheafop{E}{}{xt}}
\newcommand*{\shTor}{\ega@sheafop{T}{}{or}}
\newcommand*{\shGr}{\ega@sheafop{G}{}{r}}
\newcommand*{\shDer}{\ega@sheafop{D}{\,}{er}}
\newcommand*{\shDiff}{\ega@sheafop{D}{\,}{if{}f}\,}
\newcommand*{\shHomcont}{\ega@sheafop{H}{}{om.\,cont}}
\newcommand*{\shAut}{\ega@sheafop{A}{}{ut}}
\makeatother

% Doubtful passages of the translation: \unsure{...} adds a red "(?)" after its
% argument, and \completelyunsure stands for a passage left open, "(???)".
\newcommand*{\unsure}[1]{#1 {\color{red}(?)}}
\newcommand*{\completelyunsure}{{\color{red}(???)}}

% Pages of the printed original, noted in the margin as "volume | page":
% \oldpage[III]{81}. \sectionbreak is a centred row of asterisks.
\newcommand{\oldpage}[2][]{{\marginnote{\normalfont\textbf{#1}~|~#2}}\ignorespaces}
\newcommand*{\sectionbreak}{\begin{center}***\end{center}}

% Give every source label its own zero-size PDF destination.  The inherited
% unnumbered theorem environments otherwise leave many labels pointing to the
% preceding subsection or display anchor rather than to their exact statement.
\let\EGAOriginalLabel\label
\renewcommand{\label}[1]{\phantomsection\EGAOriginalLabel{#1}}

% Cross-references. Labels have the form <chapter>.<number> and only the number
% is printed: \sref{0.4.1.2} gives "(4.1.2)". An optional first argument adds
% a volume in bold, \sref[I]{I.2.3.1} giving "(I, 2.3.1)"; an optional last
% argument adds a further indication, \sref{0.4.1.2}[ii] giving "(4.1.2, ii)".
% \eref{eq:0.4.1.2} refers in the same way to a numbered display. Corrections
% taken from published errata are marked \erratum[volume]{corrected text}, with
% "Err" and the volume in the margin.
\makeatletter
\newcommand{\EGAUnavailableHyperref}[2]{#2}
\newcommand*{\ega@labelnumber}[1]{\StrBehind{#1}{.}}
\newcommand*{\ega@reftext}[3]{%
  \normalfont{(\IfValueT{#1}{\textbf{#1},~}\ega@labelnumber{#2}\IfValueT{#3}{,~{#3}})}}
\NewDocumentCommand{\sref}{o m o}{\@ifundefined{r@#2}%
    {\EGAUnavailableHyperref{#2}{\ega@reftext{#1}{#2}{#3}}}%
    {\hyperref[#2]{\ega@reftext{#1}{#2}{#3}}}}
\newcommand{\eref}[1]{\hyperref[#1]{\normalfont{(\ega@labelnumber{#1})}}}
\newcommand*{\ega@errmark}[1]{\textbf{Err}\textsubscript{#1}}
\newcommand{\erratum}[2][]{{#2}\marginpar{\footnotesize\ega@errmark{#1}}}
\makeatother

\ifdefined\pdfinfoomitdate\pdfinfoomitdate=1\fi
\ifdefined\pdftrailerid\pdftrailerid{}\fi
\ifdefined\pdfsuppressptexinfo\pdfsuppressptexinfo=15\fi
\setlength{\emergencystretch}{6em}

\makeatletter
\let\EGAReaderOriginalHyperref\hyperref
\renewcommand{\hyperref}[2][]{%
  \@ifundefined{r@#1}{#2}{\EGAReaderOriginalHyperref[#1]{#2}}}
\makeatother

\hypersetup{
  pdftitle={Elementos de geometría algebraica III, primera parte — texto vigente},
  pdfauthor={Alexandre Grothendieck y Jean Dieudonné; traducción española},
  pdfsubject={Estudio cohomológico de los haces coherentes}
}

\begin{document}
% BEGIN RIGHTS AND AI TRANSLATION NOTICE
\begingroup
\small

% BEGIN INHERITED HUMAN SOURCE CREDIT
\noindent\textbf{Fuentes y contribuciones}

La obra original es de Alexandre Grothendieck, con la colaboración de Jean Dieudonné: Éléments de géométrie algébrique III, première partie, Publications Mathématiques de l’IHÉS 11 (1961).

Las fuentes inglesas utilizadas en la preparación de esta traducción española incorporan la traducción humana de \href{https://github.com/ryankeleti/ega}{ryankeleti/ega}, versión \href{https://github.com/ryankeleti/ega/tree/3d7afe28c99d1887ccdb85acf2ca1afd93fb1971}{3d7afe28}, en los §§ 8–10 y parte del § 11 del capítulo 0, y los §§ 1 y 3 del capítulo III. El proyecto nombra a Mateo Carmona, Tim Holzschuh, Tim Hosgood, Ryan Keleti, Thiago Solovera \& Nery y Wayne Yang como colaboradores; sus \href{https://github.com/ryankeleti/ega/blob/3d7afe28c99d1887ccdb85acf2ca1afd93fb1971/README.md}{créditos por sección} y su \href{https://github.com/ryankeleti/ega/commits/3d7afe28c99d1887ccdb85acf2ca1afd93fb1971}{historial de contribuciones} precisan la autoría.

La traducción inglesa y la estructura TeX heredadas conservan esa procedencia. La traducción al español, las adiciones posteriores y las correcciones y composición producidas por IA se distinguen de esas contribuciones humanas. No se atribuye a los colaboradores de la traducción inglesa la traducción española ni el material posterior que no figura en su versión. La atribución no implica autorización ni aval.

\par\medskip
% END INHERITED HUMAN SOURCE CREDIT
\noindent\textbf{Derechos, disponibilidad y autoría}

Los derechos sobre las obras originales y sobre las contribuciones protegidas a estas traducciones y transcripciones derivadas corresponden a sus respectivos titulares. En la medida en que el editor tenga derechos sobre la traducción, la transcripción, las correcciones y la composición producidas por IA, los dedica al dominio público bajo CC0 1.0 Universal (https://creativecommons.org/publicdomain/zero/1.0/deed.es); esto no afecta a los derechos sobre las obras originales ni concede permiso para usarlas. Las menciones de derechos de terceros y las autorizaciones ya concedidas siguen siendo aplicables.

Estas ediciones se ofrecen gratuitamente, con fines no comerciales y exclusivamente educativos, bajo la presunción de que los respectivos titulares de derechos no se opondrían a que estas traducciones y transcripciones estuvieran disponibles gratuitamente junto a los originales franceses en NUMDAM (www.numdam.org/). Esta presunción no significa que se haya obtenido autorización; el acceso gratuito a los originales no autoriza por sí solo su traducción, adaptación o republicación.

Si usted es titular de derechos sobre algún material presentado aquí y se opone a su disponibilidad, puede abrir una incidencia en GitHub (github.com/KokunoYumeto/ega-fga-es/issues). El responsable de la publicación retirará el material afectado a petición del titular de los derechos.

La traducción al español está producida por IA. Las correcciones, la edición y la composición también están producidas por IA. El responsable de la publicación no tradujo el texto. Se conservan la autoría y las atribuciones de las obras originales y de las contribuciones heredadas; no se atribuyen a la IA. Los modelos documentados se identifican en cada edición.
\par
\endgroup
\clearpage
% END RIGHTS AND AI TRANSLATION NOTICE


\title[EGA III-1 --- Estudio cohomológico --- Primera parte]{Elementos de geometría algebraica\\
III. Estudio cohomológico de los haces coherentes\\
Primera parte}
\author{Alexandre Grothendieck y Jean Dieudonné}
\date{}
\maketitle

\begin{center}
\large Traducción española — edición actual
\end{center}

\begin{quote}\small
Traducción, correcciones y composición producidas por IA: OpenAI Codex, GPT-5.6 Sol y GPT-6 Astra, esfuerzo Ultra.
\end{quote}

\phantomsection
\label{section:ega0}
\pdfbookmark[0]{Capítulo 0 (continuación)}{chapter.0.continued}
\addcontentsline{toc}{part}{Capítulo 0 (continuación)}
\section*{CAPÍTULO 0 (continuación)}
\setcounter{section}{7}
\setcounter{subsection}{0}
\setcounter{equation}{0}

\section{Funtores representables}
\label{section:0.8}

\subsection{Funtores representables}
\label{subsection:0.8.1}

\oldpage[0\textsubscript{III}]{5}
\begin{env}[8.1.1]
\label{0.8.1.1}
Denotamos por $\Set$ la categoría de conjuntos.
Sea $\C$ una categoría; para dos objetos $X$, $Y$ de $\C$, ponemos $h_X(Y)=\Hom(Y,X)$; para cada morfismo $u:Y\to Y'$ de $\C$, denotamos por $h_X(u)$ la aplicación $v\mapsto vu$ de $\Hom(Y',X)$ en $\Hom(Y,X)$.
Es inmediato que, con estas definiciones, $h_X:\C\to\Set$ es un \emph{funtor contravariante}, es decir, un objeto de la categoría $\CHom(\C\op,\Set)$ de funtores covariantes de la categoría $\C\op$ (la dual de la categoría $\C$) en la categoría $\Set$ (T, 1.7, (d) y \cite{III-29}).
\end{env}

\begin{env}[8.1.2]
\label{0.8.1.2}
Sea ahora $w:X\to X'$ un morfismo de $\C$; para cada $Y\in\C$ y cada $v\in\Hom(Y,X)=h_X(Y)$,
tenemos $wv\in\Hom(Y,X')=h_{X'}(Y)$; denotamos por $h_w(Y)$ la aplicación $v\mapsto wv$ de
$h_X(Y)$ en $h_{X'}(Y)$. Es inmediato que, para cada morfismo $u:Y\to Y'$ de $\C$, el
diagrama
\[
  \xymatrix{
    h_X(Y')\ar[r]^{h_X(u)}\ar[d]_{h_w(Y')} &
    h_X(Y)\ar[d]^{h_w(Y)}\\
    h_{X'}(Y')\ar[r]_{h_{X'}(u)} &
    h_{X'}(Y)
  }
\]
es conmutativo; en otros términos, $h_w$ es una \emph{transformación natural (o morfismo funtorial)
$h_X\to h_{X'}$} (T, 1.2), y también un morfismo de la categoría $\CHom(\C\op,\Set)$ (T, 1.7, (d)). Las
definiciones de $h_X$ y de $h_w$ constituyen, por tanto, la definición de un \emph{funtor covariante
canónico}
\[
  h:\C\to\CHom(\C\op,\Set),\quad X\mapsto h_X.
  \tag{8.1.2.1}
\]
\end{env}

\begin{env}[8.1.3]
\label{0.8.1.3}
Sea $X$ un objeto de $\C$, y sea $F$ un funtor contravariante de $\C$ en $\Set$
(un objeto de $\CHom(\C\op,\Set)$). Sea $g:h_X\to F$ una \emph{transformación natural}: para todo $Y\in\C$,
\oldpage[0\textsubscript{III}]{6}
$g(Y)$ es, pues, una aplicación $h_X(Y)\to F(Y)$ tal que, para cada morfismo $u:Y\to Y'$ de $\C$,
el diagrama
\[
  \xymatrix{
    h_X(Y')\ar[r]^{h_X(u)}\ar[d]_{g(Y')} &
    h_X(Y)\ar[d]^{g(Y)}\\
    F(Y')\ar[r]_{F(u)} &
    F(Y)
  }
  \tag{8.1.3.1}\label{0.8.1.3.1}
\]
es conmutativo. En particular, tenemos una aplicación $g(X):h_X(X)=\Hom(X,X)\to F(X)$ y, por consiguiente, un elemento
\[
  \alpha(g)=(g(X))(1_X)\in F(X)
  \tag{8.1.3.2}
\]
y, en consecuencia, una aplicación canónica
\[
  \alpha:\Hom(h_X,F)\to F(X).
  \tag{8.1.3.3}
\]

Recíprocamente, consideremos un elemento $\xi\in F(X)$; para cada morfismo $v:Y\to X$ de $\C$, $F(v)$ es una
aplicación $F(X)\to F(Y)$; consideremos la aplicación
\[
  v\mapsto(F(v))(\xi)
  \tag{8.1.3.4}
\]
de $h_X(Y)$ en $F(Y)$; si denotamos esta aplicación por $(\beta(\xi))(Y)$,
\[
  \beta(\xi):h_X\to F
  \tag{8.1.3.5}
\]
es una \emph{transformación natural}, pues para cada morfismo $u:Y\to Y'$ de $\C$ tenemos
$(F(vu))(\xi)=(F(u)\circ F(v))(\xi)$, lo cual hace conmutativo (\hyperref[0.8.1.3.1]{8.1.3.1}) para $g=\beta(\xi)$.
Hemos definido así una aplicación canónica
\[
  \beta:F(X)\to\Hom(h_X,F).
  \tag{8.1.3.6}
\]
\end{env}

\begin{proposition}[8.1.4]
\label{0.8.1.4}
Las aplicaciones $\alpha$ y $\beta$ son biyecciones inversas entre sí.
\end{proposition}

\begin{proof}
Calculemos $\alpha(\beta(\xi))$ para $\xi\in F(X)$; para cada $Y\in\C$, $(\beta(\xi))(Y)$ es la aplicación
$g_1(Y):v\mapsto(F(v))(\xi)$ de $h_X(Y)$ en $F(Y)$. Tenemos, pues,
\[
  \alpha(\beta(\xi))=(g_1(X))(1_X)=(F(1_X))(\xi)=1_{F(X)}(\xi)=\xi.
\]
Calculemos ahora $\beta(\alpha(g))$ para $g\in\Hom(h_X,F)$; para cada $Y\in\C$, $(\beta(\alpha(g)))(Y)$
es la aplicación $v\mapsto(F(v))((g(X))(1_X))$; por la conmutatividad de (\hyperref[0.8.1.3.1]{8.1.3.1}), esta aplicación
no es otra que $v\mapsto(g(Y))((h_X(v))(1_X))=(g(Y))(v)$, por definición de $h_X(v)$; en otros términos,
es igual a $g(Y)$, lo que termina la demostración.
\end{proof}

\begin{env}[8.1.5]
\label{0.8.1.5}
Recordemos que una \emph{subcategoría} $\C'$ de una categoría $\C$ se define por la condición de que sus objetos
sean objetos de $\C$, y de que, si $X'$, $Y'$ son dos objetos de $\C'$, el conjunto
$\Hom_{\C'}(X',Y')$ de morfismos $X'\to Y'$ \emph{en $\C'$} sea un subconjunto del conjunto $\Hom_\C(X',Y')$ de
morfismos $X'\to Y'$ \emph{en $\C$}, siendo la aplicación canónica de ``composición de morfismos''
\[
  \Hom_{\C'}(X',Y')\times\Hom_{\C'}(Y',Z')\to\Hom_{\C'}(X',Z')
\]
\oldpage[0\textsubscript{III}]{7}
la restricción de la aplicación canónica
\[
  \Hom_\C(X',Y')\times\Hom_\C(Y',Z')\to\Hom_\C(X',Z').
\]

Decimos que $\C'$ es una subcategoría \emph{plena} de $\C$ si $\Hom_{\C'}(X',Y')=\Hom_\C(X',Y')$ para todo
par de objetos de $\C'$. La subcategoría $\C''$ de $\C$ formada por los objetos de $\C$ isomorfos a
objetos de $\C'$ es entonces también una subcategoría plena de $\C$, \emph{equivalente} (T, 1.2) a $\C'$, como se
comprueba fácilmente.

Un funtor covariante $F:\C_1\to\C_2$ se llama \emph{plenamente fiel} si, para todo par de objetos
$X_1$, $Y_1$ de $\C_1$, la aplicación $u\mapsto F(u)$ de $\Hom(X_1,Y_1)$ en $\Hom(F(X_1),F(Y_1))$ es
\emph{biyectiva}; esto implica que la subcategoría $F(\C_1)$ de $\C_2$ es \emph{plena}. Además,
si dos objetos $X_1$, $X_1'$ tienen la misma imagen $X_2$, existe un único isomorfismo
$u:X_1\to X_1'$ tal que $F(u)=1_{X_2}$. Para cada objeto $X_2$ de $F(\C_1)$, sea $G(X_2)$ uno de
los objetos $X_1$ de $\C_1$ tales que $F(X_1)=X_2$ ($G$ se define mediante el axioma de elección); para
cada morfismo $v:X_2\to Y_2$ de $F(\C_1)$, $G(v)$ será el único morfismo $u:G(X_2)\to G(Y_2)$ tal
que $F(u)=v$; $G$ es entonces un \emph{funtor} de $F(\C_1)$ en $\C_1$; $FG$ es el funtor identidad de
$F(\C_1)$, y lo anterior muestra que existe un isomorfismo de funtores $\vphi:1_{\C_1}\to GF$
tal que $F$, $G$, $\vphi$ y la identidad $1_{F(\C_1)}\to FG$ definen una \emph{equivalencia} entre
la categoría $\C_1$ y la subcategoría plena $F(\C_1)$ de $\C_2$ (T, 1.2).
\end{env}

\begin{env}[8.1.6]
\label{0.8.1.6}
Aplicamos la Proposición \sref{0.8.1.4} al caso en que $F$ es $h_{X'}$, siendo $X'$ un
objeto cualquiera de $\C$; la aplicación $\beta:\Hom(X,X')\to\Hom(h_X,h_{X'})$ no es otra que la aplicación
$w\mapsto h_w$ definida en \sref{0.8.1.2}; puesto que esta aplicación es \emph{biyectiva}, vemos,
con la terminología de \sref{0.8.1.5}, que:
\end{env}

\begin{proposition}[8.1.7]
\label{0.8.1.7}
El funtor canónico $h:\C\to\CHom(\C\op,\Set)$ es plenamente fiel.
\end{proposition}

\begin{env}[8.1.8]
\label{0.8.1.8}
Sea $F$ un funtor contravariante de $\C$ en $\Set$; decimos que $F$ es \emph{representable} si
existe un objeto $X\in\C$ tal que $F$ es \emph{isomorfo} a $h_X$; de la
Proposición \sref{0.8.1.7} se sigue que los datos de un $X\in\C$ y de un isomorfismo de funtores
$g:h_X\to F$ determinan $X$ salvo isomorfismo único. La Proposición \sref{0.8.1.7} implica entonces
que $h$ define una \emph{equivalencia} entre $\C$ y la subcategoría plena de
$\CHom(\C\op,\Set)$ formada por los \emph{funtores contravariantes representables}. De la
Proposición \sref{0.8.1.4} se sigue que el dato de una transformación natural $g:h_X\to F$ es
equivalente al de un elemento $\xi\in F(X)$; decir que $g$ es un \emph{isomorfismo} equivale
a la siguiente condición sobre $\xi$: \emph{para todo objeto $Y$ de $\C$, la aplicación
$v\mapsto(F(v))(\xi)$ de $\Hom(Y,X)$ en $F(Y)$ es biyectiva}. Cuando $\xi$ satisface esta condición,
decimos que el par $(X,\xi)$ \emph{representa} al funtor representable $F$. Por abuso de lenguaje,
decimos también que el objeto $X\in\C$ representa a $F$ si existe un $\xi\in F(X)$ tal que
$(X,\xi)$ representa a $F$, es decir, si $h_X$ es isomorfo a $F$.

Sean $F$, $F'$ dos funtores contravariantes representables de $\C$ en $\Set$, y sean $h_X\to F$ y
$h_{X'}\to F'$ dos isomorfismos de funtores. Se sigue entonces de \sref{0.8.1.6} que
existe una correspondencia biyectiva canónica entre $\Hom(X,X')$ y el conjunto $\Hom(F,F')$ de
transformaciones naturales $F\to F'$.
\end{env}

\begin{env}[8.1.9]
\label{0.8.1.9}
\textit{Ejemplo I. Límites proyectivos}.
La noción de funtor contravariante representable comprende, en particular, la noción ``dual'' de la noción usual de ``solución de un problema universal''.
\oldpage[0\textsubscript{III}]{8}
Más generalmente, veremos que la noción de \emph{límite proyectivo} es un caso particular de la noción de funtor representable.
Recordemos que, en una categoría $\C$, se define un \emph{sistema proyectivo} mediante el dato de un conjunto preordenado $I$, de una familia $(A_\alpha)_{\alpha\in I}$ de objetos de $\C$ y, para cada par de índices $(\alpha,\beta)$ tal que $\alpha\leq\beta$, de un morfismo $u_{\alpha\beta}:A_\beta\to A_\alpha$.
Un \emph{límite proyectivo} de este sistema en $\C$ consiste en un objeto $B$ de $\C$ (denotado $\varprojlim A_\alpha$) y, para cada $\alpha\in I$, un morfismo $u_\alpha:B\to A_\alpha$ tales que: 1\textsuperscript{o}. $u_\alpha=u_{\alpha\beta}u_\beta$ para $\alpha\leq\beta$; 2\textsuperscript{o}. para todo objeto $X$ de $\C$ y toda familia $(v_\alpha)_{\alpha\in I}$ de morfismos $v_\alpha:X\to A_\alpha$ tal que $v_\alpha=u_{\alpha\beta}v_\beta$ para $\alpha\leq\beta$, existe un único morfismo $v:X\to B$ (denotado $\varprojlim v_\alpha$) tal que $v_\alpha=u_\alpha v$ para todo $\alpha\in I$ (T, 1.8).
Esto puede interpretarse de la manera siguiente: los $u_{\alpha\beta}$ definen canónicamente aplicaciones
\[
  \overline{u}_{\alpha\beta}:\Hom(X,A_\beta)\to\Hom(X,A_\alpha)
\]
que definen un \emph{sistema proyectivo} de conjuntos $(\Hom(X,A_\alpha),\overline{u}_{\alpha\beta})$, y $(v_\alpha)$ es, por definición, un elemento del conjunto $\varprojlim\Hom(X,A_\alpha)$; está claro que $X\mapsto\varprojlim\Hom(X,A_\alpha)$ es un \emph{funtor contravariante} de $\C$ en $\Set$, y la existencia del límite proyectivo $B$ equivale a decir que $(v_\alpha)\mapsto\varprojlim v_\alpha$ es un \emph{isomorfismo} de funtores en $X$
\[
  \varprojlim\Hom(X,A_\alpha)\isoto\Hom(X,B),
  \tag{8.1.9.1}
\]
es decir, que el funtor $X\mapsto\varprojlim\Hom(X,A_\alpha)$ es \emph{representable}.
\end{env}

\begin{env}[8.1.10]
\label{0.8.1.10}
\textit{Ejemplo II. Objetos finales}.
Sea $\C$ una categoría, y sea $\{a\}$ un conjunto unitario.
Consideremos el funtor contravariante $F:\C\to\Set$ que envía cada objeto $X$ de $\C$ al conjunto $\{a\}$, y cada morfismo $X\to X'$ de $\C$ a la única aplicación $\{a\}\to\{a\}$.
Decir que este funtor es \emph{representable} significa que existe un objeto $e\in\C$ tal que, para todo $Y\in\C$, $\Hom(Y,e)=h_e(Y)$ es un \emph{conjunto unitario}; decimos que $e$ es un \emph{objeto final} de $\C$, y está claro que dos objetos finales de $\C$ son isomorfos (lo cual permite definir, en general mediante el axioma de elección, \emph{un} objeto final de $\C$, que denotamos entonces por $e_\C$).
Por ejemplo, en la categoría $\Set$, los objetos finales son los conjuntos unitarios; en la categoría de \emph{álgebras aumentadas} sobre un cuerpo $K$ (donde los morfismos son los homomorfismos de álgebras compatibles con el aumento), $K$ es un objeto final; en la categoría de \emph{$S$-esquemas} \sref[I]{I.2.5.1}, $S$ es un objeto final.
\end{env}

\begin{env}[8.1.11]
\label{0.8.1.11}
Para dos objetos $X$ e $Y$ de una categoría $\C$, pongamos $h_X'(Y)=\Hom(X,Y)$ y, para todo morfismo $u:Y\to Y'$, sea $h_X'(u)$ la aplicación $v\mapsto uv$ de $\Hom(X,Y)$ en $\Hom(X,Y')$; $h_X'$ es entonces un \emph{funtor covariante $\C\to\Set$}, de modo que deducimos, como en \sref{0.8.1.2}, la definición de un funtor covariante canónico $h':\C\op\to\CHom(\C,\Set)$; un funtor \emph{covariante} $F$ de $\C$ en $\Set$, es decir, un objeto de $\CHom(\C,\Set)$, es entonces \emph{representable} si existe un objeto $X\in\C$ (necesariamente único salvo isomorfismo único) tal que $F$ es \emph{isomorfo a $h_X'$}; dejamos al lector el desarrollo de las nociones ``duales'' de las anteriores, que esta vez comprenden la noción de \emph{límite inductivo}, y en particular la noción usual de ``solución de un problema universal''.
\end{env}

\subsection{Estructuras algebraicas en categorías}
\label{subsection:0.8.2}

\begin{env}[8.2.1]
\label{0.8.2.1}
\oldpage[0\textsubscript{III}]{9}
Dados dos funtores contravariantes $F$ y $F'$ de $\C$ en $\Set$, recordemos que, para todo objeto $Y\in\C$, ponemos $(F\times F')(Y)=F(Y)\times F'(Y)$, y para todo morfismo $u:Y\to Y'$ de $\C$, ponemos $(F\times F')(u)=F(u)\times F'(u)$, que es la aplicación $(t,t')\mapsto(F(u)(t),F'(u)(t'))$ de $F(Y')\times F'(Y')$ en $F(Y)\times F'(Y)$; $F\times F':\C\to\Set$ es así un \emph{funtor contravariante} (que no es otro que el \emph{producto} de los objetos $F$ y $F'$ en la categoría $\CHom(\C\op,\Set)$).
Dado un objeto $X\in\C$, llamamos \emph{ley de composición interna} sobre $X$ a una \emph{transformación natural}
\[
\label{0.8.2.1.1}
  \gamma_X:h_X\times h_X\to h_X.
  \tag{8.2.1.1}
\]

En otros términos (T, 1.2), para todo objeto $Y\in\C$, $\gamma_X(Y)$ es una aplicación $h_X(Y)\times h_X(Y)\to h_X(Y)$ (por tanto, por definición, una \emph{ley de composición interna} sobre el conjunto $h_X(Y)$) con la condición de que, para todo morfismo $u:Y\to Y'$ de $\C$, el diagrama
\[
  \xymatrix{
    h_X(Y')\times h_X(Y')\ar[rr]^{h_X(u)\times h_X(u)}\ar[dd]_{\gamma_X(Y')} & &
    h_X(Y)\times h_X(Y)\ar[dd]^{\gamma_X(Y)}\\\\
    h_X(Y')\ar[rr]_{h_X(u)} & &
    h_X(Y)
  }
\]
sea conmutativo; esto implica que, para las leyes de composición $\gamma_X(Y)$ y $\gamma_X(Y')$, $h_X(u)$ es un \emph{homomorfismo} de $h_X(Y')$ en $h_X(Y)$.

De manera análoga, dados dos objetos $Z$ y $X$ de $\C$, llamamos \emph{ley de composición externa sobre $X$, con $Z$ como dominio de operadores} a una transformación natural
\[
\label{0.8.2.1.2}
  \omega_{X,Z}:h_Z\times h_X\to h_X.
  \tag{8.2.1.2}
\]

Vemos como antes que, para todo $Y\in\C$, $\omega_{X,Z}(Y)$ es una ley de composición externa sobre $h_X(Y)$, con $h_Z(Y)$ como dominio de operadores, y tal que, para todo morfismo $u:Y\to Y'$, $h_X(u)$ y $h_Z(u)$ forman un \emph{dihomomorfismo} de $(h_Z(Y'),h_X(Y'))$ en $(h_Z(Y),h_X(Y))$.
\end{env}

\begin{env}[8.2.2]
\label{0.8.2.2}
Sea $X'$ un segundo objeto de $\C$, y supongamos dada una ley de composición interna $\gamma_{X'}$ sobre $X'$; decimos que un morfismo $w:X\to X'$ de $\C$ es un \emph{homomorfismo} para las leyes de composición si, para todo $Y\in\C$, $h_w(Y):h_X(Y)\to h_{X'}(Y)$ es un \emph{homomorfismo} para las leyes de composición $\gamma_X(Y)$ y $\gamma_{X'}(Y)$.
Si $X''$ es un tercer
\oldpage[0\textsubscript{III}]{10}
objeto de $\C$ provisto de una ley de composición interna $\gamma_{X''}$, y $w':X'\to X''$ es un morfismo de $\C$ que es un homomorfismo para $\gamma_{X'}$ y $\gamma_{X''}$, está claro que el morfismo $w'w:X\to X''$ es un homomorfismo para las leyes de composición $\gamma_X$ y $\gamma_{X''}$.
Un isomorfismo $w:X\isoto X'$ de $\C$ se llama \emph{isomorfismo para las leyes de composición $\gamma_X$ y $\gamma_{X'}$} si $w$ es un homomorfismo para estas leyes de composición, y si su morfismo inverso $w^{-1}$ es un homomorfismo para las leyes de composición $\gamma_{X'}$ y $\gamma_X$.

Definimos de manera análoga los \emph{dihomomorfismos} para pares de objetos de $\C$ provistos de leyes de composición externa.
\end{env}

\begin{env}[8.2.3]
\label{0.8.2.3}
Cuando una ley de composición interna $\gamma_X$ sobre un objeto $X\in\C$ es tal que $\gamma_X(Y)$ es una \emph{ley de grupo} sobre $h_X(Y)$ para \emph{todo} $Y\in\C$, decimos que $X$, provisto de esta ley, es un \emph{$\C$-grupo} o un \emph{objeto grupo de $\C$}.
Definimos análogamente los \emph{$\C$-anillos}, \emph{$\C$-módulos}, etc.
\end{env}

\begin{env}[8.2.4]
\label{0.8.2.4}
Supongamos que el \emph{producto} $X\times X$ de un objeto $X\in\C$ por sí mismo existe en $\C$; por definición, tenemos entonces $h_{X\times X}=h_X\times h_X$ salvo isomorfismo canónico, pues es un caso particular del límite proyectivo \sref{0.8.1.9}; una ley de composición interna sobre $X$ puede, por tanto, considerarse como un morfismo funtorial $\gamma_X:h_{X\times X}\to h_X$, y determina así canónicamente \sref{0.8.1.6} un elemento $c_X\in\Hom(X\times X,X)$ tal que $h_{c_X}=\gamma_X$; en este caso, el dato de una ley de composición interna sobre $X$ equivale al dato de un morfismo $X\times X\to X$; cuando $\C$ es la categoría $\Set$, recuperamos la noción clásica de ley de composición interna sobre un conjunto.
Tenemos un resultado análogo para una ley de composición externa cuando el producto $Z\times X$ existe en $\C$.
\end{env}

\begin{env}[8.2.5]
\label{0.8.2.5}
Con la notación anterior, supongamos además que $X\times X\times X$ existe en $\C$; la caracterización del producto como objeto que representa un funtor \sref{0.8.1.9} implica la existencia de isomorfismos canónicos
\[
  (X\times X)\times X\isoto X\times X\times X\isoto X\times(X\times X);
\]
si identificamos canónicamente $X\times X\times X$ con $(X\times X)\times X$, la aplicación $\gamma_X(Y)\times 1_{h_X(Y)}$ se identifica con $h_{c_X\times 1_X}(Y)$ para todo $Y\in\C$.
En consecuencia, es equivalente decir que, para todo $Y\in\C$, la ley interna $\gamma_X(Y)$ es asociativa, o que el diagrama de aplicaciones
\[
  \xymatrix{
    h_X(Y)\times h_X(Y)\times h_X(Y)\ar[rr]^{\gamma_X(Y)\times 1}\ar[dd]_{1\times\gamma_X(Y)} & &
    h_X(Y)\times h_X(Y)\ar[dd]^{\gamma_X(Y)}\\\\
    h_X(Y)\times h_X(Y)\ar[rr]_{\gamma_X(Y)} & &
    h_X(Y)
  }
\]
\oldpage[0\textsubscript{III}]{11}
es conmutativo, o que el diagrama de morfismos
\[
  \xymatrix{
    X\times X\times X\ar[rr]^{c_X\times 1_X}\ar[dd]_{1_X\times c_X} & &
    X\times X\ar[dd]^{c_X}\\\\
    X\times X\ar[rr]_{c_X} & &
    X
  }
\]
es conmutativo.
\end{env}

\begin{env}[8.2.6]
\label{0.8.2.6}
Bajo las hipótesis de \sref{0.8.2.5}, si queremos expresar, para todo $Y\in\C$, que la ley interna $\gamma_X(Y)$ es una \emph{ley de grupo}, es necesario, en primer lugar, que sea asociativa y, en segundo lugar, que exista una aplicación $\alpha_X(Y):h_X(Y)\to h_X(Y)$ con las propiedades de la operación de \emph{inverso} de un grupo; como hemos visto que, para todo morfismo $u:Y\to Y'$ de $\C$, $h_X(u)$ debe ser un homomorfismo de grupos $h_X(Y')\to h_X(Y)$, vemos primero que $\alpha_X:h_X\to h_X$ debe ser una \emph{transformación natural}.
Por otra parte, se pueden expresar las propiedades características del inverso $s\mapsto s^{-1}$ de un grupo $G$ sin hacer intervenir el elemento identidad: basta comprobar que las dos aplicaciones compuestas
\[
  (s,t)\mapsto(s,s^{-1},t)\mapsto(s,s^{-1}t)\mapsto s(s^{-1}t),
\]
\[
  (s,t)\mapsto(s,s^{-1},t)\mapsto(s,ts^{-1})\mapsto(ts^{-1})s
\]
son iguales a la segunda proyección $(s,t)\mapsto t$ de $G\times G$ en $G$.
Por \sref{0.8.1.3}, tenemos $\alpha_X=h_{a_X}$, donde $a_X\in\Hom(X,X)$; la primera condición anterior expresa entonces que el morfismo compuesto
\[
  X\times X\xrightarrow{(1_X,a_X)\times 1_X}X\times X\times X\xrightarrow{1_X\times c_X}X\times X\xrightarrow{c_X}X
\]
es la segunda proyección $X\times X\to X$ en $\C$, y la segunda condición es análoga.
\end{env}

\begin{env}[8.2.7]
\label{0.8.2.7}
Supongamos ahora que existe un \emph{objeto final} $e$ \sref{0.8.1.10} en $\C$.
Supongamos siempre que $\gamma_X(Y)$ es una ley de grupo sobre $h_X(Y)$ para todo $Y\in\C$, y denotemos por $\eta_X(Y)$ el elemento identidad de $\gamma_X(Y)$.
Como, para todo morfismo $u:Y\to Y'$ de $\C$, $h_X(u)$ es un homomorfismo de grupos, tenemos $\eta_X(Y)=(h_X(u))(\eta_X(Y'))$; tomando en particular $Y'=e$, caso en el cual $u$ es el único elemento $\varepsilon$ de $\Hom(Y,e)$, vemos que el elemento $\eta_X(e)$ determina por completo $\eta_X(Y)$ para todo $Y\in\C$.
Pongamos $e_X=\eta_X(X)$, el elemento identidad del grupo $h_X(X)=\Hom(X,X)$; la conmutatividad del diagrama
\[
  \xymatrix{
    h_X(e)\ar[r]^{h_X(\varepsilon)}\ar[d]_{h_{e_X}(e)} &
    h_X(Y)\ar[d]^{h_{e_X}(Y)}\\
    h_X(e)\ar[r]_{h_X(\varepsilon)} &
    h_X(Y)
  }
\]
(cf. \sref{0.8.1.2}) muestra que, sobre el conjunto $h_X(Y)$, la aplicación $h_{e_X}(Y)$ no es otra que
\oldpage[0\textsubscript{III}]{12}
$s\mapsto\eta_X(Y)$, que envía todo elemento al elemento identidad.
Comprobamos entonces que el hecho de que $\eta_X(Y)$ sea el elemento identidad de $\gamma_X(Y)$ para todo $Y\in\C$ equivale a decir que el morfismo compuesto
\[
  X\xrightarrow{(1_X,1_X)}X\times X\xrightarrow{1_X\times e_X}X\times X\xrightarrow{c_X}X,
\]
y su análogo en el que se intercambian $1_X$ y $e_X$, son ambos \emph{iguales} a $1_X$.
\end{env}

\begin{env}[8.2.8]
\label{0.8.2.8}
Naturalmente, se podrían ampliar con facilidad los ejemplos de estructuras algebraicas en categorías.
El ejemplo de los grupos se ha tratado con suficiente detalle, pero en adelante dejaremos por lo general al lector el desarrollo de las nociones análogas para los ejemplos de estructuras algebraicas que encontremos.
\end{env}

\section{Conjuntos construibles}
\label{section:0.9}

\subsection{Conjuntos construibles}
\label{subsection:0.9.1}

\begin{definition}[9.1.1]
\label{0.9.1.1}
Decimos que una aplicación continua $f:X\to Y$ es \emph{cuasicompacta} si, para todo subconjunto abierto cuasicompacto $U$ de $Y$, $f^{-1}(U)$ es cuasicompacto.
Decimos que un subconjunto $Z$ de un espacio topológico $X$ es \emph{retrocompacto} en $X$ si la inyección canónica $Z\to X$ es cuasicompacta; en otros términos, si para todo subconjunto abierto cuasicompacto $U$ de $X$, $U\cap Z$ es cuasicompacto.
\end{definition}

Un subconjunto \emph{cerrado} de $X$ es retrocompacto en $X$, pero un subconjunto cuasicompacto de $X$ no es necesariamente retrocompacto en $X$.
Si $X$ es cuasicompacto, todo subconjunto abierto retrocompacto de $X$ es cuasicompacto.
Está claro que toda unión \emph{finita} de conjuntos retrocompactos en $X$ es retrocompacta en $X$, pues toda unión finita de conjuntos cuasicompactos es cuasicompacta.
Toda intersección finita de abiertos retrocompactos de $X$ es un abierto retrocompacto de $X$.
En un espacio \emph{localmente noetheriano} $X$, todo conjunto cuasicompacto es un subespacio noetheriano y, en consecuencia, \emph{todo subconjunto} de $X$ es retrocompacto en $X$.

\begin{definition}[9.1.2]
\label{0.9.1.2}
Dado un espacio topológico $X$, decimos que un subconjunto de $X$ es \emph{construible} si pertenece al menor conjunto de subconjuntos $\mathfrak{F}$ de $X$ que contiene todos los subconjuntos abiertos retrocompactos de $X$ y es estable por intersecciones finitas y complementarios (lo que implica que $\mathfrak{F}$ también es estable por uniones finitas).
\end{definition}

\begin{proposition}[9.1.3]
\label{0.9.1.3}
Para que un subconjunto de $X$ sea construible, es necesario y suficiente que sea una unión finita de conjuntos de la forma $U\cap\complement{V}$, donde $U$ y $V$ son abiertos retrocompactos de $X$.
\end{proposition}

\begin{proof}
Está claro que la condición es suficiente.
Para ver que es necesaria, consideremos el conjunto $\mathfrak{G}$ de las uniones finitas de conjuntos de la forma $U\cap\complement{V}$, donde $U$ y $V$ son abiertos retrocompactos de $X$; basta ver que el complementario de todo conjunto de $\mathfrak{G}$ pertenece a $\mathfrak{G}$.
Sea $Z=\bigcup_{i\in I}(U_i\cap\complement{V_i})$, donde $I$ es finito y los $U_i$ y $V_i$ son abiertos retrocompactos de $X$; tenemos
$\complement{Z}=\bigcap_{i\in I}(V_i\cup\complement{U_i})$, de modo que $\complement{Z}$ es una unión finita de conjuntos que son intersecciones de cierto número de los $V_i$ y de cierto número de los $\complement{U_i}$, por tanto de la forma
\oldpage[0\textsubscript{III}]{13}
$V\cap\complement{U}$, donde $U$ es la unión de cierto número de los $U_i$ y $V$ es la intersección de cierto número de los $V_i$; pero hemos observado antes que las uniones e intersecciones finitas de abiertos retrocompactos de $X$ son abiertos retrocompactos de $X$, de donde se sigue la conclusión.
\end{proof}

\begin{corollary}[9.1.4]
\label{0.9.1.4}
Todo subconjunto construible de $X$ es retrocompacto en $X$.
\end{corollary}

\begin{proof}
Basta demostrar que, si $U$ y $V$ son abiertos retrocompactos de $X$, entonces $U\cap\complement{V}$ es retrocompacto en $X$; si $W$ es un abierto cuasicompacto de $X$, entonces $W\cap U\cap\complement{V}$ es cerrado en el espacio cuasicompacto $W\cap U$ y, por tanto, es cuasicompacto.
\end{proof}

En particular:
\begin{corollary}[9.1.5]
\label{0.9.1.5}
Para que un subconjunto abierto $U$ de $X$ sea construible, es necesario y suficiente que sea retrocompacto en $X$.
Para que un subconjunto cerrado $F$ de $X$ sea construible, es necesario y suficiente que el abierto $\complement{F}$ sea retrocompacto.
\end{corollary}

\begin{env}[9.1.6]
\label{0.9.1.6}
Un caso importante es aquel en que todo subconjunto abierto cuasicompacto de $X$ es retrocompacto; en otros términos, aquel en que la intersección de dos subconjuntos abiertos cuasicompactos de $X$ es cuasicompacta (cf.~\sref[I]{I.5.5.6}).
Cuando $X$ es además cuasicompacto, esto implica que los subconjuntos abiertos retrocompactos de $X$ coinciden con los subconjuntos abiertos cuasicompactos de $X$, y que los subconjuntos construibles de $X$ son las uniones finitas de conjuntos de la forma $U\cap\complement{V}$, donde $U$ y $V$ son abiertos cuasicompactos.
\end{env}

\begin{corollary}[9.1.7]
\label{0.9.1.7}
Para que un subconjunto de un espacio noetheriano $X$ sea construible, es necesario y suficiente que sea una unión finita de subconjuntos localmente cerrados de $X$.
\end{corollary}

\begin{proposition}[9.1.8]
\label{0.9.1.8}
Sea $X$ un espacio topológico y $U$ un subconjunto abierto de $X$.
\begin{enumerate}
  \item[{\rm(i)}] Si $T$ es un subconjunto construible de $X$, entonces $T\cap U$ es un subconjunto construible de $U$.
  \item[{\rm(ii)}] Supongamos, además, que $U$ es retrocompacto en $X$. Para que un subconjunto $Z$ de $U$ sea construible en $X$, es necesario y suficiente que sea construible en $U$.
\end{enumerate}
\end{proposition}

\begin{proof}
\medskip\noindent
\begin{enumerate}
  \item[(i)] Mediante la Proposición \sref{0.9.1.3}, nos reducimos a demostrar que, si $T$ es un abierto retrocompacto de $X$, entonces $T\cap U$ es un abierto retrocompacto de $U$; en otros términos, para todo abierto cuasicompacto $W\subset U$, $T\cap U\cap W=T\cap W$ es cuasicompacto, lo que se sigue inmediatamente de la hipótesis.
  \item[(ii)] La condición es necesaria por (i), por lo que queda demostrar que es suficiente.
    Por la Proposición \sref{0.9.1.3}, basta considerar el caso en que $Z$ es un abierto retrocompacto \emph{en $U$}, pues entonces se seguirá que $U\setmin Z$ es construible en $X$, y si $Z$ y $Z'$ son dos abiertos retrocompactos \emph{en $U$}, entonces $Z\cap(U\setmin Z')$ será construible en $X$.
    Si $W$ es un abierto cuasicompacto de $X$ y $Z$ un abierto retrocompacto de $U$, tenemos $Z\cap W=Z\cap(W\cap U)$, y por hipótesis $W\cap U$ es un abierto cuasicompacto de $U$; por tanto, $W\cap Z$ es cuasicompacto y, en consecuencia, $Z$ es un abierto retrocompacto de $X$, y \emph{a fortiori} construible en $X$.
\end{enumerate}
\end{proof}

\begin{corollary}[9.1.9]
\label{0.9.1.9}
Sea $X$ un espacio topológico y $(U_i)_{i\in I}$ un recubrimiento finito de $X$ por abiertos retrocompactos de $X$.
Para que un subconjunto $Z$ de $X$ sea construible en $X$, es necesario y suficiente que $Z\cap U_i$ sea construible en $U_i$ para todo $i\in I$.
\end{corollary}

\begin{env}[9.1.10]
\label{0.9.1.10}
Supongamos, en particular, que $X$ es cuasicompacto y que todo punto
\oldpage[0\textsubscript{III}]{14}
de $X$ admite un sistema fundamental de entornos abiertos retrocompactos en $X$ (y \emph{a fortiori} cuasicompactos); entonces la condición para que un subconjunto $Z$ de $X$ sea construible en $X$ es de naturaleza \emph{local}; en otros términos, es necesario y suficiente que, para todo $x\in X$, exista un entorno abierto $V$ de $x$ tal que $V\cap Z$ sea construible en $V$.
En efecto, si se satisface esta condición, entonces para todo $x\in X$ existe un entorno abierto $V$ de $x$ que es \emph{retrocompacto en $X$} y tal que $V\cap Z$ es construible en $V$, por las hipótesis sobre $X$ y por la Proposición \sref{0.9.1.8}[i]; basta entonces recubrir $X$ por un número finito de estos entornos y aplicar el Corolario \sref{0.9.1.9}.
\end{env}

\begin{definition}[9.1.11]
\label{0.9.1.11}
Sea $X$ un espacio topológico.
Decimos que un subconjunto $T$ de $X$ es \emph{localmente construible} en $X$ si, para todo $x\in X$, existe un entorno abierto $V$ de $x$ tal que $T\cap V$ es construible en $V$.
\end{definition}

De la Proposición \sref{0.9.1.8}[i] se sigue que, si $V$ es tal que $V\cap T$ es construible en $V$, entonces para todo abierto $W\subset V$, $W\cap T$ es construible en $W$.
Si $T$ es localmente construible en $X$, entonces, para todo abierto $U$ de $X$, $T\cap U$ es localmente construible en $U$, en virtud de la observación anterior.
La misma observación muestra que el conjunto de subconjuntos localmente construibles de $X$ es estable por uniones finitas e intersecciones finitas; por otra parte, está claro que también es estable por paso al complementario.

\begin{proposition}[9.1.12]
\label{0.9.1.12}
Sea $X$ un espacio topológico.
Todo conjunto construible en $X$ es localmente construible en $X$.
El recíproco es cierto si $X$ es cuasicompacto y si su topología admite una base formada por los conjuntos retrocompactos de $X$.
\end{proposition}

\begin{proof}
La primera afirmación se sigue de la Definición \sref{0.9.1.11}, y la segunda de \sref{0.9.1.10}.
\end{proof}

\begin{corollary}[9.1.13]
\label{0.9.1.13}
Sea $X$ un espacio topológico cuya topología admite una base formada por los conjuntos retrocompactos de $X$.
Entonces todo subconjunto $T$ localmente construible de $X$ es retrocompacto en $X$.
\end{corollary}

\begin{proof}
Sea $U$ un abierto cuasicompacto de $X$; $T\cap U$ es localmente construible en $U$, luego construible en $U$ por la Proposición~\sref{0.9.1.12} y, en consecuencia, cuasicompacto por el Corolario \sref{0.9.1.4}.
\end{proof}

\subsection{Subconjuntos construibles de espacios noetherianos}
\label{subsection:0.9.2}

\begin{env}[9.2.1]
\label{0.9.2.1}
Hemos visto \sref{0.9.1.7} que, en un espacio noetheriano $X$, los subconjuntos construibles de $X$ son las \emph{uniones finitas de subconjuntos localmente cerrados de $X$}.

La imagen inversa de un conjunto construible de $X$ por una aplicación continua de un espacio noetheriano $X'$ en $X$ es construible en $X'$.
Si $Y$ es un subconjunto construible de un espacio noetheriano $X$, entonces los subconjuntos de $Y$ que son construibles como subespacios de $Y$ coinciden con los que son construibles como subespacios de $X$.
\end{env}

\begin{proposition}[9.2.2]
\label{0.9.2.2}
Sea $X$ un espacio noetheriano irreducible, y sea $E$ un subconjunto construible de $X$.
Para que $E$ sea denso en todas partes en $X$, es necesario y suficiente que $E$ contenga un subconjunto abierto no vacío de $X$.
\end{proposition}

\begin{proof}
\oldpage[0\textsubscript{III}]{15}
La condición es evidentemente suficiente, pues todo abierto no vacío es denso en $X$.
Recíprocamente, sea $E=\bigcup_{i=1}^n(U_i\cap F_i)$ un subconjunto construible de $X$, siendo los $U_i$ abiertos no vacíos y los $F_i$ cerrados en $X$; tenemos entonces $\overline{E}\subset\bigcup_i F_i$.
En consecuencia, si $\overline{E}=X$, entonces $X$ es igual a uno de los $F_i$, luego $E\supset U_i$, lo que termina la demostración.
\end{proof}

Cuando $X$ admite un punto genérico $x$ \sref[0]{0.2.1.2}, la condición de la Proposición \sref{0.9.2.2} equivale a la relación $x\in E$.

\begin{proposition}[9.2.3]
\label{0.9.2.3}
Sea $X$ un espacio noetheriano.
Para que un subconjunto $E$ de $X$ sea construible, es necesario y suficiente que, para todo subconjunto cerrado irreducible $Y$ de $X$, $E\cap Y$ sea raro en $Y$ o contenga un subconjunto abierto no vacío de $Y$.
\end{proposition}

\begin{proof}
La necesidad de la condición se sigue de que $E\cap Y$ debe ser un subconjunto construible de $Y$ y de la Proposición \sref{0.9.2.2}, pues un subconjunto no denso de $Y$ es necesariamente raro en el espacio irreducible $Y$ \sref[0]{0.2.1.1}.
Para demostrar que la condición es suficiente, apliquemos el principio de inducción noetheriana \sref[0]{0.2.2.2} al conjunto $\mathfrak{F}$ de subconjuntos cerrados $Y$ de $X$ tales que $Y\cap E$ es construible (respecto de $Y$ o respecto de $X$, lo que es equivalente): podemos suponer así que, para todo subconjunto cerrado $Y\neq X$ de $X$, $E\cap Y$ es construible.
Supongamos primero que $X$ no es irreducible, y sean $X_i$ ($1\leq i\leq m$) sus componentes irreducibles, necesariamente en número finito \sref[0]{0.2.2.5}; por hipótesis, los $E\cap X_i$ son construibles, luego también lo es su unión $E$.
Supongamos ahora que $X$ es irreducible; entonces, por hipótesis, si $E$ es raro, se tiene $\overline{E}\neq X$ y $E=E\cap\overline{E}$ es construible; si $E$ contiene un subconjunto abierto no vacío $U$ de $X$, es la unión de $U$ y $E\cap(X\setmin U)$; pero $X\setmin U$ es un conjunto cerrado distinto de $X$, por lo que $E\cap(X\setmin U)$ es construible; en consecuencia, $E$ mismo es construible, lo que termina la demostración.
\end{proof}

\begin{corollary}[9.2.4]
\label{0.9.2.4}
Sea $X$ un espacio noetheriano y $(E_\alpha)$ una familia filtrante creciente de subconjuntos construibles de $X$, tal que
\begin{enumerate}
  \item[\emph{(1.o)}] $X$ es la unión de la familia $(E_\alpha)$.
  \item[\emph{(2.o)}] Todo subconjunto cerrado irreducible de $X$ está contenido en la adherencia de uno de los $E_\alpha$.
\end{enumerate}

Entonces existe un índice $\alpha$ tal que $X=E_\alpha$.

Cuando todo subconjunto cerrado irreducible de $X$ admite un punto genérico, puede omitirse la hipótesis \emph{(2.o)}.
\end{corollary}

\begin{proof}
Apliquemos el principio de inducción noetheriana \sref[0]{0.2.2.2} al conjunto $\mathfrak{M}$ de subconjuntos cerrados de $X$ contenidos en al menos uno de los $E_\alpha$; podemos suponer así que todo subconjunto cerrado $Y\neq X$ de $X$ está contenido en uno de los $E_\alpha$.
La proposición es evidente si $X$ no es irreducible, porque cada una de las componentes irreducibles $X_i$ de $X$ ($1\leq i\leq m$) está contenida en un $E_{\alpha_i}$, y existe un $E_\alpha$ que contiene todos los $E_{\alpha_i}$.
Supongamos ahora que $X$ es irreducible.
Por hipótesis, existe un $\beta$ tal que $X=\overline{E_\beta}$, de modo que \sref{0.9.2.2} $E_\beta$ contiene un subconjunto abierto no vacío $U$ de $X$.
Pero entonces el conjunto cerrado $X\setmin U$ está contenido en un $E_\gamma$, y basta tomar un $E_\alpha$ que contenga $E_\beta$ y $E_\gamma$.
Cuando todo subconjunto cerrado irreducible $Y$ de $X$
\oldpage[0\textsubscript{III}]{16}
admite un punto genérico $y$, existe $\alpha$ tal que $y\in E_\alpha$, de modo que $Y=\overline{\{y\}}\subset\overline{E_\alpha}$, y la condición (2.o) es, por tanto, consecuencia de (1.o).
\end{proof}

\begin{proposition}[9.2.5]
\label{0.9.2.5}
Sea $X$ un espacio noetheriano, sea $x$ un punto de $X$ y sea $E$ un subconjunto construible de $X$.
Para que $E$ sea un entorno de $x$, es necesario y suficiente que, para todo subconjunto cerrado irreducible $Y$ de $X$ que contenga $x$, $E\cap Y$ sea denso en $Y$ (si existe un punto genérico $y$ de $Y$, esto significa también \sref{0.9.2.2} que $y\in E$).
\end{proposition}

\begin{proof}
La condición es evidentemente necesaria; demostraremos que es suficiente.
Aplicando el principio de inducción noetheriana al conjunto $\mathfrak{M}$ de subconjuntos cerrados $Y$ de $X$ que contienen $x$ y tales que $E\cap Y$ es un entorno de $x$ en $Y$, podemos suponer que todo subconjunto cerrado $Y\neq X$ de $X$ que contiene $x$ pertenece a $\mathfrak{M}$.
Si $X$ no es irreducible, cada una de las componentes irreducibles $X_i$ de $X$ que contienen $x$ es distinta de $X$, luego $E\cap X_i$ es un entorno de $x$ respecto de $X_i$; en consecuencia, $E$ es un entorno de $x$ en la unión de las componentes irreducibles de $X$ que contienen $x$, y como esta unión es un entorno de $x$ en $X$, también lo es $E$.
Si $X$ es irreducible, $E$ es denso en $X$ por hipótesis, por lo que contiene un subconjunto abierto no vacío $U$ de $X$ \sref{0.9.2.2}; la proposición es entonces evidente si $x\in U$; en caso contrario, por hipótesis, $x$ es un punto interior de $E\cap(X\setmin U)$ relativo a $X\setmin U$, de modo que la adherencia de $X\setmin E$ en $X$ no contiene $x$, y el complementario de esta adherencia es un entorno de $x$ contenido en $E$, lo que termina la demostración.
\end{proof}

\begin{corollary}[9.2.6]
\label{0.9.2.6}
Sea $X$ un espacio noetheriano y sea $E$ un subconjunto de $X$.
Para que $E$ sea un abierto de $X$, es necesario y suficiente que, para todo subconjunto cerrado irreducible $Y$ de $X$ que interseca $E$, $E\cap Y$ contenga un subconjunto abierto no vacío de $Y$.
\end{corollary}

\begin{proof}
La condición es evidentemente necesaria; recíprocamente, si se satisface, implica que $E$ es construible por la Proposición \sref{0.9.2.3}.
Además, la Proposición \sref{0.9.2.5} muestra que $E$ es entonces un entorno de cada uno de sus puntos, de donde se sigue la conclusión.
\end{proof}

\subsection{Funciones construibles}
\label{subsection:0.9.3}

\begin{definition}[9.3.1]
\label{0.9.3.1}
Sea $h$ una aplicación de un espacio topológico $X$ en un conjunto $T$.
Decimos que $h$ es \emph{construible} si $h^{-1}(t)$ es construible para todo $t\in T$, y vacío salvo para un número finito de valores de $t$; entonces, para todo subconjunto $S$ de $T$, $h^{-1}(S)$ es construible.
Decimos que $h$ es \emph{localmente construible} si todo $x\in X$ posee un entorno abierto $V$ tal que $h|V$ es construible.
\end{definition}

Toda función construible es localmente construible; el recíproco es cierto cuando $X$ es cuasicompacto y admite una base formada por los abiertos retrocompactos de $X$ (en particular, cuando $X$ es noetheriano).

\begin{proposition}[9.3.2]
\label{0.9.3.2}
Sea $h$ una aplicación de un espacio noetheriano $X$ en un conjunto $T$.
Para que $h$ sea construible, es necesario y suficiente que, para todo subconjunto cerrado irreducible $Y$ de $X$, exista un subconjunto no vacío $U$ de $Y$, abierto respecto de $Y$, sobre el cual $h$ sea constante.
\end{proposition}

\begin{proof}
La condición es necesaria: en efecto, por hipótesis, $h$ solo toma un número finito de valores $t_i$ sobre $Y$, y cada uno de los conjuntos $h^{-1}(t_i)\cap Y$ es construible en $Y$ \sref{0.9.2.1}; como no pueden ser todos subconjuntos raros del espacio $Y$, al menos uno de ellos contiene un abierto no vacío \sref{0.9.2.3}.

\oldpage[0\textsubscript{III}]{17}
Para ver que la condición es suficiente, aplicamos el principio de inducción noetheriana al conjunto $\mathfrak{M}$ de subconjuntos cerrados $Y$ de $X$ tales que la restricción $h|Y$ es construible; podemos suponer así que, para todo subconjunto cerrado $Y\neq X$ de $X$,
$h|Y$ es construible.
Si $X$ no es irreducible, la restricción de $h$ a cada una de las componentes irreducibles (en número finito) $X_i$ de $X$ es construible, y se sigue entonces inmediatamente de la Definición \sref{0.9.3.1} que $h$ es construible.
Si $X$ es irreducible, existe por hipótesis un subconjunto abierto no vacío $U$ de $X$ sobre el cual $h$ es constante; por otra parte, la restricción de $h$ a $X\setmin U$ es construible por hipótesis, y se sigue inmediatamente que $h$ es construible.
\end{proof}

\begin{corollary}[9.3.3]
\label{0.9.3.3}
Sea $X$ un espacio noetheriano en el que todo subconjunto cerrado irreducible admite un punto genérico.
Si $h$ es una aplicación de $X$ en un conjunto $T$ tal que, para todo $t\in T$, $h^{-1}(t)$ es construible, entonces $h$ es construible.
\end{corollary}

\begin{proof}
Si $Y$ es un subconjunto cerrado irreducible de $X$ e $y$ su punto genérico, entonces $Y\cap h^{-1}(h(y))$ es construible y contiene $y$; por tanto, \sref{0.9.2.2} este conjunto contiene un subconjunto abierto no vacío de $Y$, y basta aplicar la Proposición \sref{0.9.3.2}.
\end{proof}

\begin{proposition}[9.3.4]
\label{0.9.3.4}
Sea $X$ un espacio noetheriano en el que todo subconjunto cerrado irreducible admite un punto genérico, y sea $h$ una aplicación construible de $X$ en un conjunto ordenado.
Para que $h$ sea semicontinua superiormente sobre $X$, es necesario y suficiente que, para todo $x\in X$ y toda generalización \sref[0]{0.2.1.2} $x'$ de $x$, se tenga $h(x')\leq h(x)$.
\end{proposition}

\begin{proof}
La función $h$ solo toma un número finito de valores; por tanto, decir que es semicontinua superiormente significa que, para todo $x\in X$, el conjunto $E$ de los $y\in X$ tales que $h(y)\leq h(x)$ es un entorno de $x$.
Por hipótesis, $E$ es un subconjunto construible de $X$; por otra parte, decir que un subconjunto cerrado irreducible $Y$ de $X$ contiene $x$ significa que su punto genérico $y$ es una generalización de $x$; la conclusión se sigue entonces de la Proposición \sref{0.9.2.5}.
\end{proof}

\section{Complemento sobre los módulos planos}
\label{section:0.10}

Para las demostraciones que falten en \hyperref[subsection:0.10.1]{(10.1)} y \hyperref[subsection:0.10.2]{(10.2)}, remitimos al lector a Bourbaki,~\emph{Alg.~comm.}, cap.~II y~III.

\subsection{Relaciones entre módulos planos y módulos libres}
\label{subsection:0.10.1}

\begin{env}[10.1.1]
\label{0.10.1.1}
Sea $A$ un anillo, $\mathfrak{J}$ un ideal de $A$ y $M$ un $A$-módulo;
para todo entero $p\geq0$, tenemos un homomorfismo canónico de $(A/\mathfrak{J})$-módulos
\[
\label{0.10.1.1.1}
  \vphi_p:(M/\mathfrak{J}M)\otimes_{A/\mathfrak{J}}(\mathfrak{J}^p/\mathfrak{J}^{p+1})\to\mathfrak{J}^pM/\mathfrak{J}^{p+1}M,
  \tag{10.1.1.1}
\]
que es evidentemente \emph{sobreyectivo}.
Denotamos por $\gr(A)=\oplus_{p\geq0}\mathfrak{J}^p/\mathfrak{J}^{p+1}$ el anillo graduado asociado a $A$ filtrado por los $\mathfrak{J}^p$, y por $\gr(M)=\oplus_{p\geq0}\mathfrak{J}^pM/\mathfrak{J}^{p+1}M$ el $\gr(A)$-módulo graduado asociado a $M$ filtrado por los $\mathfrak{J}^pM$;
tenemos entonces $\gr_p(A)=\mathfrak{J}^p/\mathfrak{J}^{p+1}$ y $\gr_p(M)=\mathfrak{J}^pM/\mathfrak{J}^{p+1}M$;
los $\vphi$ definen un homomorfismo \emph{sobreyectivo} de $\gr(A)$-módulos graduados
\[
\label{0.10.1.1.2}
  \vphi:\gr_0(M)\otimes_{\gr_0(A)}\gr(A)\to\gr(M).
  \tag{10.1.1.2}
\]
\end{env}

\oldpage[0\textsubscript{III-1}]{18}
\begin{env}[10.1.2]
\label{0.10.1.2}
Supongamos que se satisface \emph{una} de las hipótesis siguientes:
\begin{enumerate}
  \item[(i)] $\mathfrak{J}$ es nilpotente;
  \item[(ii)] $A$ es noetheriano, $\mathfrak{J}$ está contenido en el radical de $A$ y $M$ es de tipo finito.
\end{enumerate}
Entonces las propiedades siguientes son equivalentes.
\begin{enumerate}
  \item[(a)] $M$ es un $A$-módulo libre.
  \item[(b)] $M/\mathfrak{J}M=M\otimes_A(A/\mathfrak{J})$ es un $(A/\mathfrak{J})$-módulo libre y $\Tor_1^A(M,A/\mathfrak{J})=0$.
  \item[(c)] $M/\mathfrak{J}M$ es un $(A/\mathfrak{J})$-módulo libre y el homomorfismo canónico \sref{0.10.1.1.2} es inyectivo (y, por tanto, biyectivo).
\end{enumerate}
\end{env}

\begin{env}[10.1.3]
\label{0.10.1.3}
Supongamos que $A/\mathfrak{J}$ es un \emph{cuerpo} (en otros términos, que $\mathfrak{J}$ es maximal) y que se satisface una de las hipótesis (i) y (ii) de \sref{0.10.1.2}.
Entonces las propiedades siguientes son equivalentes.
\begin{enumerate}
  \item[(a)] $M$ es un $A$-módulo libre.
  \item[(b)] $M$ es un $A$-módulo proyectivo.
  \item[(c)] $M$ es un $A$-módulo plano.
  \item[(d)] $\Tor_1^A(M,A/\mathfrak{J})=0$.
  \item[(e)] El homomorfismo canónico \sref{0.10.1.1.2} es biyectivo.
\end{enumerate}

Este resultado puede aplicarse, en particular, a los dos casos siguientes:
\begin{enumerate}
  \item[(i)] $M$ es un módulo \emph{arbitrario} sobre un anillo local $A$ cuyo ideal maximal $\mathfrak{J}$ es \emph{nilpotente} (por ejemplo, un anillo artiniano local);
  \item[(ii)] $M$ es un módulo \emph{de tipo finito} sobre un anillo \emph{local noetheriano}.
\end{enumerate}
\end{env}

\subsection{Criterios locales de planitud}
\label{subsection:0.10.2}

\begin{env}[10.2.1]
\label{0.10.2.1}
Con las hipótesis y la notación de \sref{0.10.1.1}, consideremos las condiciones siguientes.
\begin{enumerate}
  \item[(a)] $M$ es un $A$-módulo plano.
  \item[(b)] $M/\mathfrak{J}M$ es un $(A/\mathfrak{J})$-módulo plano y $\Tor_1^A(M,A/\mathfrak{J})=0$.
  \item[(c)] $M/\mathfrak{J}M$ es un $(A/\mathfrak{J})$-módulo plano y el homomorfismo canónico \sref{0.10.1.1.2} es biyectivo.
  \item[(d)] Para todo $n\geq1$, $M/\mathfrak{J}^nM$ es un $(A/\mathfrak{J}^n)$-módulo plano.
\end{enumerate}

Entonces tenemos las implicaciones
\[
  \text{(a)}\implies\text{(b)}\implies\text{(c)}\implies\text{(d)},
\]
y, si $\mathfrak{J}$ es \emph{nilpotente}, las cuatro condiciones son \emph{equivalentes}.
Ocurre también así si $A$ es noetheriano y $M$ es \emph{idealmente separado}, es decir, si para todo ideal $\mathfrak{a}$ de $A$, el $A$-módulo $\mathfrak{a}\otimes_A M$ está \emph{separado} para la topología $\mathfrak{J}$-preádica.
\end{env}

\begin{env}[10.2.2]
\label{0.10.2.2}
Sea $A$ un anillo noetheriano, $B$ una $A$-álgebra conmutativa noetheriana, $\mathfrak{J}$ un ideal de $A$ tal que $\mathfrak{J}B$ está contenido en el radical de $B$, y $M$ un $B$-\emph{módulo} de tipo finito.
Entonces, cuando $M$ se considera como $A$-\emph{módulo}, las cuatro condiciones de \sref{0.10.2.1} son \emph{equivalentes}.
\oldpage[0\textsubscript{III-1}]{19}
Este resultado se aplica ante todo al caso en que $A$ y $B$ son anillos noetherianos \emph{locales} y el homomorfismo $A\to B$ es un homomorfismo \emph{local}.
Más precisamente, si $\mathfrak{J}$ es entonces el ideal \emph{maximal} de $A$, podemos suprimir, en las condiciones~\emph{(b)} y \emph{(c)}, la hipótesis de que $M/\mathfrak{J}M$ sea plano, pues se satisface automáticamente, y la condición~\emph{(d)} implica que los módulos $M/\mathfrak{J}^nM$ son \emph{libres} sobre los $A/\mathfrak{J}^n$.
\end{env}

\begin{env}[10.2.3]
\label{0.10.2.3}
Con las hipótesis sobre $A$, $B$, $\mathfrak{J}$ y $M$ del comienzo de \sref{0.10.2.2}, sea $\widehat{A}$ la completación separada de $A$ para la topología $\mathfrak{J}$-preádica, y $\widehat{M}$ la completación separada de $M$ para la topología $\mathfrak{J}B$-preádica.
Para que $M$ sea un $A$-módulo plano, es necesario y suficiente que $\widehat{M}$ sea un $\widehat{A}$-módulo plano.
\end{env}

\begin{env}[10.2.4]
\label{0.10.2.4}
Sea $\rho:A\to B$ un homomorfismo local de anillos noetherianos locales, sea $k$ el cuerpo residual de $A$, y sean $M$ y $N$ dos $B$-módulos de tipo finito, suponiendo $N$ \emph{$A$-plano}.
Sea $u:M\to N$ un $B$-homomorfismo.
Entonces las condiciones siguientes son equivalentes.
\begin{enumerate}
  \item[(a)] $u$ es inyectivo y $\Coker(u)$ es un $A$-módulo plano.
  \item[(b)] $u\otimes1:M\otimes_A k\to N\otimes_A k$ es inyectivo.
\end{enumerate}
\end{env}

\begin{env}[10.2.5]
\label{0.10.2.5}
Sean $\rho:A\to B$ y $\sigma:B\to C$ homomorfismos locales de anillos noetherianos locales, sea $k$ el cuerpo residual de $A$, y sea $M$ un $C$-módulo de tipo finito.
Supongamos que $B$ es un $A$-módulo \emph{plano}.
Entonces las condiciones siguientes son equivalentes.
\begin{enumerate}
  \item[(a)] $M$ es un $B$-módulo plano.
  \item[(b)] $M$ es un $A$-módulo plano y $M\otimes_A k$ es un $(B\otimes_A k)$-módulo plano.
\end{enumerate}
\end{env}

\begin{proposition}[10.2.6]
\label{0.10.2.6}
Sean $A$ y $B$ anillos noetherianos locales, sea $\rho:A\to B$ un homomorfismo local, sea $\mathfrak{J}$ un ideal de $B$ contenido en el ideal maximal, y sea $M$ un $B$-módulo de tipo finito.
Supongamos que, para todo $n\geq0$, $M_n=M/\mathfrak{J}^{n+1}M$ es un $A$-módulo plano.
Entonces $M$ es un $A$-módulo plano.
\end{proposition}

\begin{proof}
Debemos demostrar que, para todo homomorfismo inyectivo $u:N'\to N$ de $A$-módulos de tipo finito, $v=1\otimes u:M\otimes_A N'\to M\otimes_A N$ es inyectivo.
Pero $M\otimes_A N'$ y $M\otimes_A N$ son $B$-módulos de tipo finito y, por tanto, están separados para la topología $\mathfrak{J}$-preádica \sref[0\textsubscript{I}]{0.7.3.5};
basta, pues, demostrar que el homomorfismo $\widehat{v}:\widehat{M\otimes_A N'}\to\widehat{M\otimes_A N}$ de las completaciones separadas es inyectivo.
Pero $\widehat{v}=\varprojlim v_n$, donde $v_n$ es el homomorfismo $1\otimes u:M_n\otimes_A N'\to M_n\otimes_A N$;
como, por hipótesis, $M_n$ es $A$-plano, $v_n$ es inyectivo para todo $n$, y también lo es $\widehat{v}$, pues el funtor $\varprojlim$ es exacto por la izquierda.
\end{proof}

\begin{corollary}[10.2.7]
\label{0.10.2.7}
Sea $A$ un anillo noetheriano, sea $B$ un anillo noetheriano local, sea $\rho:A\to B$ un homomorfismo, sea $f$ un elemento del ideal maximal de $B$, y sea $M$ un $B$-módulo de tipo finito.
Supongamos que la homotecia $f_M:x\to fx$ sobre $M$ es inyectiva y que $M/fM$ es un $A$-módulo plano.
Entonces $M$ es un $A$-módulo plano.
\end{corollary}

\begin{proof}
Sea $M_i=f^iM$ para $i\geq0$;
como $f_M$ es inyectiva, $M_i/M_{i+1}$ es isomorfo a $M/fM$ y, por tanto, $A$-plano para todo $i\geq0$;
la sucesión exacta
\[
  0\to M_i/M_{i+1}\to M/M_{i+1}\to M/M_i\to 0
\]
nos da, por inducción sobre $i$, que $M/M_i$ es \emph{$A$-plano} para todo $i\geq0$ \sref[0\textsubscript{I}]{0.6.1.2};
podemos, pues, aplicar \sref{0.10.2.6}.
También podemos razonar directamente como sigue:
para todo $A$-módulo $N$ de tipo finito, $M\otimes_A N$ es un $B$-módulo de tipo finito;
como $f$ pertenece al radical $\mathfrak{n}$ de $B$, la topología $(f)$-ádica sobre $M\otimes_A N$ es más fina que la topología
\oldpage[0\textsubscript{III-1}]{20}
$\mathfrak{n}$-ádica, y sabemos que esta última está \emph{separada} \sref[0\textsubscript{I}]{0.7.3.5}.
Ahora bien, como $M/M_i$ es $A$-plano, tenemos
\[
  f^i(M\otimes_A N)=\Im(M_i\otimes_A N\to M\otimes_A N)=\Ker(M\otimes_A N\to(M/M_i)\otimes_A N)
\]
por \sref[0\textsubscript{I}]{0.6.1.2}.
Sea, pues, $N$ un $A$-módulo de tipo finito, y sea $N'$ un submódulo de $N$, con inyección canónica $j:N'\to N$; en el diagrama conmutativo
\[
  \xymatrix{
    M\otimes_A N'\ar[r]\ar[d]_{1_M\otimes j} &
    (M/M_i)\otimes_A N'\ar[d]^{1_{M/M_i}\otimes j}\\
    M\otimes_A N\ar[r] &
    (M/M_i)\otimes_A N
  }
\]
$1_{M/M_i}\otimes j$ es inyectivo, porque $M/M_i$ es $A$-plano; concluimos, pues, que
\[
  \Ker(M\otimes_A N'\to M\otimes_A N)\subset\Ker(M\otimes_A N'\to(M/M_i)\otimes_A N')
\]
para todo $i$; como la intersección de estos últimos núcleos es $0$, según hemos visto antes, también lo es el primer núcleo y, por tanto, $M$ es $A$-plano.
\end{proof}

\begin{proposition}[10.2.8]
\label{0.10.2.8}
Sea $A$ un anillo noetheriano reducido y sea $M$ un $A$-módulo de tipo finito.
Supongamos que, para toda $A$-álgebra $B$ que sea un anillo de valoración discreta, $M\otimes_A B$ es un $B$-módulo plano (y, por tanto, libre \sref{0.10.1.3}).
Entonces $M$ es un $A$-módulo plano.
\end{proposition}

\begin{proof}
Sabemos que, para que $M$ sea plano, es necesario y suficiente que $M_\mathfrak{m}$ sea un $A_\mathfrak{m}$-módulo plano para todo ideal maximal $\mathfrak{m}$ de $A$ \sref[0\textsubscript{I}]{0.6.3.3};
podemos, pues, restringirnos al caso en que $A$ es \emph{local} \sref[0\textsubscript{I}]{0.1.2.8}.
Sea $\mathfrak{m}$ el ideal maximal de $A$, sean $\mathfrak{p}_i$ ($1\leq i\leq r$) los ideales primos minimales de $A$, y sea $k$ el cuerpo residual $A/\mathfrak{m}$.
Sabemos \sref[II]{II.7.1.7} que existe, para cada $i$, un anillo de valoración discreta $B_i$ que tiene el mismo cuerpo de fracciones $K_i$ que el dominio íntegro $A/\mathfrak{p}_i$ y que, además, domina a $A/\mathfrak{p}_i$.
Sea $M_i=M\otimes_A B_i$.
Por hipótesis, $M_i$ es libre sobre $B_i$ y, denotando por $k_i$ el cuerpo residual de $B_i$, tenemos
\[
\label{0.10.2.8.1}
  \rg_{k_i}(M_i\otimes_{B_i}k_i)=\rg_{K_i}(M_i\otimes_{B_i}K_i).
  \tag{10.2.8.1}
\]

Pero está claro que el homomorfismo compuesto $A\to A/\mathfrak{p}_i\to B_i$ es local, de modo que $k_i$ es una extensión de $k$, y que tenemos $M_i\otimes_{B_i}k_i = M\otimes_A k_i = (M\otimes_A k)\otimes_k k_i$, así como $M_i\otimes_{B_i}K_i = M\otimes_A K_i$.
La ecuación~\sref{0.10.2.8.1} implica, por tanto, que
\[
  \rg_k(M\otimes_A k)=\rg_{K_i}(M\otimes_A K_i)\quad\mbox{para $1\leq i\leq r$}
\]
y, como $A$ es reducido, sabemos que esta condición implica que $M$ es un $A$-módulo \emph{libre} (Bourbaki, \emph{Alg. comm.}, cap.~II, \textsection~3, n\textsuperscript{o}~2, prop.~7).
\end{proof}

\subsection{Existencia de extensiones planas de anillos locales}
\label{subsection:0.10.3}

\begin{proposition}[10.3.1]
\label{0.10.3.1}
Sea $A$ un anillo noetheriano local, con ideal maximal $\mathfrak{J}$ y cuerpo residual $k=A/\mathfrak{J}$.
Sea $K$ una extensión del cuerpo $k$.
Entonces existe un homomorfismo local de $A$ en un anillo noetheriano local $B$, tal que $B/\mathfrak{J}B$ es $k$-isomorfo a $K$ y tal que $B$ es un $A$-módulo plano.
\end{proposition}

El resto de esta sección está dedicado a demostrar esta proposición paso a paso.
\begin{env}[10.3.1.1]
\label{0.10.3.1.1}
Supongamos primero que $K=k(T)$, donde $T$ es una indeterminada.
En el anillo de polinomios $A'=A[T]$, consideremos el ideal primo $\mathfrak{J}'=\mathfrak{J}A'$, formado por los
\oldpage[0\textsubscript{III-1}]{21}
polinomios cuyos coeficientes pertenecen al ideal $\mathfrak{J}$;
está claro que $A'/\mathfrak{J}'$ es canónicamente isomorfo a $k[T]$.
Mostraremos que el anillo de fracciones $B=A'_{\mathfrak{J}'}$ es el que buscamos (es decir, un anillo que satisface las condiciones de la conclusión de la proposición);
es claramente un anillo noetheriano local, con ideal maximal $\mathfrak{L}=\mathfrak{J}B$.
Además, $B/\mathfrak{L}=(A'/\mathfrak{J}')_{\mathfrak{J}'}=(k[T])_{\mathfrak{J}'}$ es precisamente el cuerpo de fracciones $K$ de $k[T]$.
Finalmente, $B$ es un $A'$-módulo plano y $A'$ es un $A$-módulo libre, luego $B$ es un $A$-módulo plano \sref[0\textsubscript{I}]{0.6.2.1}.
\end{env}

\begin{env}[10.3.1.2]
\label{0.10.3.1.2}
Supongamos ahora que $K=k(t)=k[t]$, donde $t$ es algebraico sobre $k$;
sea $f\in k[T]$ el polinomio minimal de $t$;
existe un polinomio mónico $F\in A[T]$ cuya imagen canónica en $k[T]$ es $f$.
Sea, pues, $A'=A[T]$, y sea $\mathfrak{J}'$ el ideal $\mathfrak{J}A'+(F)$ de $A'$.
Veremos que el anillo cociente $B=A'/(F)$ es el que buscamos.
Ante todo, está claro que $B$ es un $A$-módulo \emph{libre} y, por tanto, plano.
El anillo $A'/\mathfrak{J}'$ es isomorfo a
\[
  (A'/\mathfrak{J}A')/\big((\mathfrak{J}A'+(F))/\mathfrak{J}A'\big)=k[T]/(f)=K;
\]
la imagen $\mathfrak{L}$ de $\mathfrak{J}'$ en $B$ es, por tanto, maximal, y tenemos evidentemente $\mathfrak{L}=\mathfrak{J}B$.
Finalmente, $B$ es un anillo semilocal, porque es un $A$-módulo de tipo finito (Bourbaki, \emph{Alg. comm.}, cap.~IV, \textsection~2, n\textsuperscript{o}~5, cor.~3 de la prop.~9), y sus ideales maximales están en correspondencia biyectiva con los de $B/\mathfrak{J}B$ (\cite[vol.~I, p.~259]{I-13});
los argumentos anteriores demuestran entonces que $B$ es un anillo local.
\end{env}

\begin{lemma}[10.3.1.3]
\label{0.10.3.1.3}
Sea $(A_\lambda,f_{\mu\lambda})$ un sistema inductivo filtrante de anillos locales, tal que los $f_{\mu\lambda}$ son homomorfismos locales;
sea $\mathfrak{m}_\lambda$ el ideal maximal de $A_\lambda$, y sea $K_\lambda=A_\lambda/\mathfrak{m}_\lambda$.
Entonces $A'=\varinjlim A_\lambda$ es un anillo local, con ideal maximal $\mathfrak{m}'=\varinjlim\mathfrak{m}_\lambda$ y cuerpo residual $K=\varinjlim K_\lambda$.
Además, si $\mathfrak{m}_\mu=\mathfrak{m}_\lambda A_\mu$ para $\lambda<\mu$, tenemos $\mathfrak{m}'=\mathfrak{m}_\lambda A'$ para todo $\lambda$.
Si, además, para $\lambda<\mu$, $A_\mu$ es un $A_\lambda$-módulo plano y todos los $A_\lambda$ son noetherianos, entonces $A'$ es noetheriano y es un $A_\lambda$-módulo plano para todo $\lambda$.
\end{lemma}

\begin{proof}
Como, por hipótesis, $f_{\mu\lambda}(\mathfrak{m}_\lambda)\subset\mathfrak{m}_\mu$ para $\lambda<\mu$, los $\mathfrak{m}_\lambda$ forman un sistema inductivo y su límite $\mathfrak{m}'$ es evidentemente un ideal de $A'$.
Además, si $x'\not\in\mathfrak{m}'$, existe un $\lambda$ tal que $x'=f_\lambda(x_\lambda)$ para algún $x_\lambda\in A_\lambda$ (donde $f_\lambda:A_\lambda\to A'$ denota el homomorfismo canónico);
como $x'\not\in\mathfrak{m}'$, necesariamente $x_\lambda\not\in\mathfrak{m}_\lambda$, de modo que $x_\lambda$ admite un inverso $y_\lambda$ en $A_\lambda$, y $y'=f_\lambda(y_\lambda)$ es el inverso de $x'$ en $A'$, lo que demuestra que $A'$ es un anillo local con ideal maximal $\mathfrak{m}'$;
la afirmación sobre $K$ se sigue inmediatamente de que $\varinjlim$ es un funtor exacto.
La hipótesis $\mathfrak{m}_\mu=\mathfrak{m}_\lambda A_\mu$ implica que la aplicación canónica $\mathfrak{m}_\lambda\otimes_{A_\lambda}A_\mu\to\mathfrak{m}_\mu$ es sobreyectiva;
la igualdad $\mathfrak{m}'=\mathfrak{m}_\lambda A'$ se sigue entonces, de nuevo, de que el funtor $\varinjlim$ es exacto y conmuta con el producto tensorial.

Supongamos ahora que, para $\lambda<\mu$, tenemos $\mathfrak{m}_\mu=\mathfrak{m}_\lambda A_\mu$ y que $A_\mu$ es un $A_\lambda$-módulo plano.
Entonces $A'$ es un $A_\lambda$-módulo plano para todo $\lambda$, por \sref[0\textsubscript{I}]{0.6.2.3};
como $A'$ y $A_\lambda$ son anillos locales y $\mathfrak{m}'=\mathfrak{m}_\lambda A'$, $A'$ es incluso un módulo \emph{fielmente plano} sobre $A_\lambda$ \sref[0\textsubscript{I}]{0.6.6.2}.
Supongamos finalmente, además, que los $A_\lambda$ son \emph{noetherianos};
las topologías $\mathfrak{m}_\lambda$-ádicas están entonces separadas \sref[0\textsubscript{I}]{0.7.3.5};
mostraremos ahora que de esto se sigue que la topología $\mathfrak{m}'$-ádica sobre $A'$ está \emph{separada}.
En efecto, si $x'\in A'$ pertenece a todos los $\mathfrak{m}^{'n}$ ($n>0$), entonces es la imagen de algún $x_\mu\in A_\mu$ para cierto índice $\mu$, y como la imagen inversa en $A_\mu$ de $\mathfrak{m}^{'n}=\mathfrak{m}_\mu^n A'$ es $\mathfrak{m}_\mu^n$ \sref[0\textsubscript{I}]{0.6.6.1}, $x_\mu$ pertenece a todos los $\mathfrak{m}_\mu^n$, luego $x_\mu=0$ por hipótesis y, por tanto, $x'=0$.
Denotemos por $\widehat{A}'$ la completación de $A'$ para la topología $\mathfrak{m}'$-ádica;
lo anterior muestra que tenemos $A'\subset\widehat{A}'$.
Mostraremos ahora que $\widehat{A}'$ es \emph{noetheriano} y \emph{$A_\lambda$-plano} para todo $\lambda$;
de esto,
\oldpage[0\textsubscript{III-1}]{22}
se seguirá que $\widehat{A}'$ es $A'$-plano \sref[0\textsubscript{I}]{0.6.2.3} y, como $\mathfrak{m}'\widehat{A}'\neq\widehat{A}'$, que $\widehat{A}'$ es un $A'$-módulo fielmente plano \sref[0\textsubscript{I}]{0.6.6.2}; de ahí se obtiene finalmente que $A'$ es \emph{noetheriano} \sref[0\textsubscript{I}]{0.6.5.2}, lo que terminará la demostración del lema.

Tenemos $\widehat{A}'=\varprojlim_n A'/\mathfrak{m}^{\prime n}$;
por el hecho de que $A'$ es $A_\lambda$-plano, tenemos
\[
  \mathfrak{m}^{\prime n}/\mathfrak{m}^{\prime n+1}=(\mathfrak{m}_\lambda^{n}/\mathfrak{m}_\lambda^{n+1})\otimes_{A_\lambda}A'=(\mathfrak{m}_\lambda^{n}/\mathfrak{m}_\lambda^{n+1})\otimes_{K_\lambda}(K_\lambda\otimes_{A_\lambda}A')=(\mathfrak{m}_\lambda^{n}/\mathfrak{m}_\lambda^{n+1})\otimes_{K_\lambda}K;
\]
como $\mathfrak{m}_\lambda^{n}/\mathfrak{m}_\lambda^{n+1}$ es un $K_\lambda$-espacio vectorial de dimensión finita, $\mathfrak{m}^{\prime n}/\mathfrak{m}^{\prime n+1}$ es un $K$-espacio vectorial de dimensión finita para todo $n\geq 0$.
Se sigue entonces de \sref[0\textsubscript{I}]{0.7.2.12} y \sref[0\textsubscript{I}]{0.7.2.8} que $\widehat{A}'$ es \emph{noetheriano}.
Sabemos además que el ideal maximal de $\widehat{A}'$ es $\mathfrak{m}'\widehat{A}'$, y que $\widehat{A}'/\mathfrak{m}^{\prime n}\widehat{A}'$ es isomorfo a $A'/\mathfrak{m}^{\prime n}$;
como $A'/\mathfrak{m}^{\prime n}=(A_\lambda/\mathfrak{m}_\lambda^n)\otimes_{A_\lambda}A'$, vemos que $A'/\mathfrak{m}^{\prime n}$ es un $(A_\lambda/\mathfrak{m}_\lambda^n)$-módulo plano \sref[0\textsubscript{I}]{0.6.2.1};
el criterio~\sref{0.10.2.2} es, por tanto, aplicable a la $A_\lambda$-álgebra noetheriana $\widehat{A}'$, y muestra que $\widehat{A}'$ es \emph{$A_\lambda$-plano}.
\end{proof}

\begin{env}[10.3.1.4]
\label{0.10.3.1.4}
Tratemos ahora el caso general.
Existe un ordinal $\gamma$ y, para todo ordinal $\lambda\leq\gamma$, un subcuerpo $k_\lambda$ de $K$ que contiene a $k$, tales que (i) para todo $\lambda<\gamma$, $k_{\lambda+1}$ es una extensión de $k_\lambda$ generada por un solo elemento; (ii) para todo ordinal límite $\mu$, $k_\mu=\bigcup_{\lambda<\mu}k_\lambda$; y (iii) $K=k_\gamma$.
En efecto, basta considerar una biyección $\xi\mapsto t_\xi$ del conjunto de ordinales $\xi\leq\beta$ (para algún $\beta$ apropiado) en $K$, y definir $k_\lambda$ por inducción transfinita (para $\lambda\leq\beta$) como la unión de los $k_\mu$ para $\mu<\lambda$ si $\lambda$ es un ordinal límite, y como $k_\nu(t_\xi)$ si $\lambda=\nu+1$, donde $\xi$ es el menor ordinal tal que $t_\xi\not\in k_\nu$;
$\gamma$ es entonces, por definición, el menor ordinal $\leq\beta$ tal que $k_\gamma=K$.

Teniendo esto presente, definiremos, por inducción transfinita, una familia de anillos noetherianos locales $A_\lambda$ para $\lambda\leq\gamma$, y homomorfismos locales $f_{\mu\lambda}:A_\lambda\to A_\mu$ para $\lambda\leq\mu$, que satisfagan las condiciones siguientes:
\begin{enumerate}
  \item[(i)] $(A_\lambda,f_{\mu\lambda})$ es un sistema inductivo y $A_0=A$;
  \item[(ii)] para todo $\lambda$, tenemos un $k$-isomorfismo $A_\lambda/\mathfrak{J}A_\lambda\isoto k_\lambda$;
  \item[(iii)] para $\lambda\leq\mu$, $A_\mu$ es un $A_\lambda$-módulo plano.
\end{enumerate}

Supongamos, pues, definidos los $A_\lambda$ y los $f_{\mu\lambda}$ para $\lambda<\mu<\xi$, y supongamos primero que $\xi=\zeta+1$, de modo que $k_\xi=k_\zeta(t)$.
Si $t$ es trascendente sobre $k_\zeta$, definimos $A_\xi$, siguiendo el procedimiento de \sref{0.10.3.1.1}, como $(A_\zeta[t])_{\mathfrak{J}A_\zeta[t]}$;
la aplicación canónica es $f_{\xi\zeta}$ y, para $\lambda<\zeta$, tomamos $f_{\xi\lambda}=f_{\xi\zeta}\circ f_{\zeta\lambda}$;
la verificación de las condiciones~(i) a (iii) es entonces inmediata, en virtud de lo demostrado en \sref{0.10.3.1.1}.
Supongamos ahora que $t$ es algebraico, y sea $h$ su polinomio minimal en $k_\zeta[T]$, y sea $H$ un polinomio mónico de $A_\zeta[T]$ cuya imagen en $k_\zeta[T]$ es $h$;
tomamos entonces $A_\xi$ igual a $A_\zeta[T]/(H)$, definiendo los $f_{\xi\lambda}$ como antes;
la verificación de las condiciones~(i) a (iii) se sigue entonces de lo demostrado en \sref{0.10.3.1.2}.

Supongamos ahora que $\xi$ no tiene predecesor;
tomamos entonces $A_\xi$ como límite inductivo del sistema inductivo de anillos locales $(A_\lambda,f_{\mu\lambda})$ para $\lambda<\xi$;
definimos $f_{\xi\lambda}$ como la aplicación canónica para $\lambda<\xi$.
El hecho de que $A_\xi$ sea local y noetheriano, de que los $f_{\xi\lambda}$ sean homomorfismos locales y de que las condiciones~(i) a (iii) se satisfagan para $\lambda\leq\xi$
\oldpage[0\textsubscript{III-1}]{23}
se sigue entonces de la hipótesis de inducción y del Lema~\sref{0.10.3.1.3}.
Con esta construcción, está claro que el anillo $B=A_\gamma$ satisface las condiciones de \sref{0.10.3.1}.\qed
\end{env}

Observemos que, por \sref{0.10.2.1}[c], tenemos un isomorfismo canónico
\[
\label{0.10.3.1.5}
  \gr(A)\otimes_k K\simto\gr(B).
  \tag{10.3.1.5}
\]

También podemos reemplazar $B$ por su completación $\mathfrak{J}B$-ádica $\widehat{B}$ sin modificar las conclusiones de \sref{0.10.3.1}, porque $\widehat{B}$ es un $B$-módulo plano \sref[0\textsubscript{I}]{0.7.3.3} y, por tanto, un $A$-módulo plano \sref[0\textsubscript{I}]{0.6.2.1}.

Hemos demostrado también lo siguiente:
\begin{corollary}[10.3.2]
\label{0.10.3.2}
Si $K$ es una extensión de grado finito, podemos suponer que $B$ es una $A$-álgebra finita.
\end{corollary}

\section{Suplemento sobre álgebra homológica}
\label{section:0.11}

\subsection{Repaso de las sucesiones espectrales}
\label{subsection:0.11.1}

\begin{env}[11.1.1]
\label{0.11.1.1}
En lo que sigue utilizamos una noción de sucesión espectral más general que la definida en (T, 2.4); conservando la notación de (T, 2.4), llamamos \emph{sucesión espectral} en una categoría abeliana $\cat{C}$ a un sistema $E$ constituido por las partes siguientes:
\begin{enumerate}
  \item[(a)] Una familia $(E_r^{pq})$ de objetos de $\cat{C}$ definida para $p,q\in\bb{Z}$ y $r\geq 2$.
  \item[(b)] Una familia de morfismos $d_r^{pq}:E_r^{pq}\to E_r^{p+r,q-r+1}$ tales que $d_r^{p+r,q-r+1}d_r^{pq}=0$.
    Ponemos $Z_{r+1}(E_r^{pq})=\Ker(d_r^{pq})$ y $B_{r+1}(E_r^{pq})=\Im(d_r^{p-r,q+r-1})$, de modo que
    \[
      B_{r+1}(E_r^{pq})\subset Z_{r+1}(E_r^{pq})\subset E_r^{pq}.
    \]
  \item[(c)] Una familia de isomorfismos $\alpha_r^{pq}:Z_{r+1}(E_r^{pq})/B_{r+1}(E_r^{pq})\isoto E_{r+1}^{pq}$.

  Definimos entonces, para $k\geq r+1$, por inducción sobre $k$, los subobjetos $B_k(E_r^{pq})$ y $Z_k(E_r^{pq})$ como las imágenes inversas, por el morfismo canónico $E_r^{pq}\to E_r^{pq}/B_{r+1}(E_r^{pq})$, de los subobjetos de este cociente identificados mediante $\alpha_r^{pq}$ con los subobjetos $B_k(E_{r+1}^{pq})$ y $Z_k(E_{r+1}^{pq})$, respectivamente.
Está claro que se tiene entonces, salvo isomorfismo,
  \[
  \label{0.11.1.1.1}
    Z_k(E_r^{pq})/B_k(E_r^{pq})=E_k^{pq}\text{ para }k\geq r+1,
    \tag{11.1.1.1}
  \]
  y, si ponemos $B_r(E_r^{pq})=0$ y $Z_r(E_r^{pq})=E_r^{pq}$, se tienen las relaciones de inclusión
  \[
  \label{0.11.1.1.2}
    0=B_r(E_r^{pq})\subset B_{r+1}(E_r^{pq})\subset B_{r+2}(E_r^{pq})\subset\cdots\subset Z_{r+2}(E_r^{pq})\subset Z_{r+1}(E_r^{pq})\subset Z_r(E_r^{pq})=E_r^{pq}.
    \tag{11.1.1.2}
  \]
  Las otras partes de los datos de $E$ son entonces:
  \item[(d)] Dos subobjetos $B_\infty(E_2^{pq})$ y $Z_\infty(E_2^{pq})$ de $E_2^{pq}$ tales que $B_\infty(E_2^{pq})\subset Z_\infty(E_2^{pq})$ y, para todo $k\geq 2$,
    \[
      B_k(E_2^{pq})\subset B_\infty(E_2^{pq})\text{ y }Z_\infty(E_2^{pq})\subset Z_k(E_2^{pq}).
    \]
    Ponemos
    \[
    \label{0.11.1.1.3}
      E_\infty^{pq}=Z_\infty(E_2^{pq})/B_\infty(E_2^{pq}).
      \tag{11.1.1.3}
    \]
  \item[(e)]
\oldpage[0\textsubscript{III}]{24}
    Una familia $(E^n)$ de objetos de $\cat{C}$, cada uno provisto de una \emph{filtración decreciente $(F^p(E^n))_{p\in\bb{Z}}$}.
    Como de costumbre, denotamos por $\gr(E^n)$ el objeto graduado asociado al objeto filtrado $E^n$, suma directa de los $\gr_p(E^n)=F^p(E^n)/F^{p+1}(E^n)$.
  \item[(f)] Para todo par $(p,q)\in\bb{Z}\times\bb{Z}$, un isomorfismo $\beta^{pq}:E_\infty^{pq}\isoto\gr_p(E^{p+q})$.
\end{enumerate}

La familia $(E^n)$, sin las filtraciones, se llama el \emph{término límite} (o \emph{límite}) de la sucesión espectral $E$.

Supongamos que la categoría $\cat{C}$ admite sumas directas infinitas, o que para todo $r\geq 2$ y todo $n\in\bb{Z}$ existe solo un número finito de pares $(p,q)$ tales que $p+q=n$ y $E_r^{pq}\neq 0$ (basta que esto se cumpla para $r=2$).
Entonces están definidos los $E_r^{(n)}=\sum_{p+q=n}E_r^{pq}$ y, si denotamos por $d_r^{(n)}$ el morfismo $E_r^{(n)}\to E_r^{(n+1)}$ cuya restricción a $E_r^{pq}$ es $d_r^{pq}$ para todo par $(p,q)$ tal que $p+q=n$, entonces $d_r^{(n+1)}\circ d_r^{(n)}=0$; en otras palabras, $(E_r^{(n)})_{n\in\bb{Z}}$ es un \emph{complejo $E_r^{(\bullet)}$} en $\cat{C}$, con diferencial de grado $+1$, y de (c) se deduce que $\HH^n(E_r^{(\bullet)})$ es \emph{isomorfo a $E_{r+1}^{(n)}$} para todo $r\geq 2$.
\end{env}

\begin{env}[11.1.2]
\label{0.11.1.2}
Un \emph{morfismo $u:E\to E'$} de una sucesión espectral $E$ en una sucesión espectral $E'=(E_r^{\prime pq},E^{\prime n})$ consiste en sistemas de morfismos $u_r^{pq}:E_r^{pq}\to E_r^{\prime pq}$ y $u^n:E^n\to E^{\prime n}$, siendo los $u^n$ compatibles con las filtraciones de $E^n$ y $E^{\prime n}$, y tales que los diagramas
\[
  \xymatrix{
    E_r^{pq}\ar[r]^-{d_r^{pq}}\ar[d]_{u_r^{pq}} &
    E_r^{p+r,q-r+1}\ar[d]^{u_r^{p+r,q-r+1}}\\
    E_r^{\prime pq}\ar[r]_{d_r^{\prime pq}} &
    E_r^{\prime p+r,q-r+1}
  }
\]
sean conmutativos; además, pasando a los cocientes, $u_r^{pq}$ da un morfismo $\overline{u}_r^{pq}:Z_{r+1}(E_r^{pq})/B_{r+1}(E_r^{pq})\to Z_{r+1}(E_r^{\prime pq})/B_{r+1}(E_r^{\prime pq})$ y debe tenerse $\alpha_r^{\prime pq}\circ\overline{u}_r^{pq}=u_{r+1}^{pq}\circ\alpha_r^{pq}$; finalmente, debe tenerse $u_2^{pq}(B_\infty(E_2^{pq}))\subset B_\infty(E_2^{\prime pq})$ y $u_2^{pq}(Z_\infty(E_2^{pq}))\subset Z_\infty(E_2^{\prime pq})$; pasando a los cocientes, $u_2^{pq}$ da entonces un morfismo $u_\infty^{pq}:E_\infty^{pq}\to E_\infty^{\prime pq}$, y el diagrama
\[
  \xymatrix{
    E_\infty^{pq}\ar[r]^-{u_\infty^{pq}}\ar[d]_{\beta^{pq}} &
    E_\infty^{\prime pq}\ar[d]^{\beta^{\prime pq}}\\
    \gr_p(E^{p+q})\ar[r]_{\gr_p(u^{p+q})} &
    \gr_p(E^{\prime p+q})
  }
\]
debe ser conmutativo.

Las definiciones anteriores muestran, por inducción sobre $r$, que si los $u_2^{pq}$ son \emph{isomorfismos}, también lo son los $u_r^{pq}$ para $r\geq 2$; \emph{si además sabemos que $u_2^{pq}(B_\infty(E_2^{pq}))=B_\infty(E_2^{\prime pq})$ y $u_2^{pq}(Z_\infty(E_2^{pq}))=Z_\infty(E_2^{\prime pq})$ y que los $u^n$ son isomorfismos}, entonces podemos concluir que $u$ es un \emph{isomorfismo}.
\end{env}

\begin{env}[11.1.3]
\label{0.11.1.3}
\oldpage[0\textsubscript{III}]{25}
Recordemos que si $(F^p(X))_{p\in\bb{Z}}$ es una \emph{filtración} (decreciente) de un objeto $X\in\cat{C}$, se dice que esta filtración es \emph{separada} si $\inf(F^p(X))=0$, \emph{discreta} si existe un $p$ tal que $F^p(X)=0$, \emph{exhaustiva} (o \emph{coseparada}) si $\sup(F^p(X))=X$, y \emph{codiscreta} si existe un $p$ tal que $F^p(X)=X$.

Se dice que una sucesión espectral $E=(E_r^{pq},E^n)$ es \emph{débilmente convergente} si se tiene $B_\infty(E_2^{pq})=\sup_k(B_k(E_2^{pq}))$ y $Z_\infty(E_2^{pq})=\inf_k(Z_k(E_2^{pq}))$ (en otras palabras, los objetos $B_\infty(E_2^{pq})$ y $Z_\infty(E_2^{pq})$ están determinados por los datos (a)--(c) de la sucesión espectral $E$).
Se dice que la sucesión espectral $E$ es \emph{regular} si es débilmente convergente y si, además:
\begin{enumerate}
  \item[$1^\circ$] Para todo par $(p,q)$, la sucesión decreciente $(Z_k(E_2^{pq}))_{k\geq 2}$ es \emph{estable}; la hipótesis de que $E$ es débilmente convergente implica entonces que $Z_\infty(E_2^{pq})=Z_k(E_2^{pq})$ para $k$ suficientemente grande (dependiendo de $p$ y $q$).
  \item[$2^\circ$] Para todo $n$, la filtración $(F^p(E^n))_{p\in\bb{Z}}$ de $E^n$ es \emph{discreta} y \emph{exhaustiva}.
\end{enumerate}

Se dice que la sucesión espectral $E$ es \emph{corregular} si es débilmente convergente y si, además:
\begin{enumerate}
  \item[$3^\circ$] Para todo par $(p,q)$, la sucesión creciente $(B_k(E_2^{pq}))_{k\geq 2}$ es \emph{estable}, lo que implica que $B_\infty(E_2^{pq})=B_k(E_2^{pq})$ y, en consecuencia, $E_\infty^{pq}=\inf E_k^{pq}$.
  \item[$4^\circ$] Para todo $n$, la filtración de $E^n$ es \emph{codiscreta}.
\end{enumerate}

Finalmente, se dice que $E$ es \emph{birregular} si es a la vez regular y corregular, es decir, si se cumplen las condiciones siguientes:
\begin{enumerate}
  \item[(a)] Para todo par $(p,q)$, las sucesiones $(B_k(E_2^{pq}))_{k\geq 2}$ y $(Z_k(E_2^{pq}))_{k\geq 2}$ son \emph{estables} y se tiene $B_\infty(E_2^{pq})=B_k(E_2^{pq})$ y $Z_\infty(E_2^{pq})=Z_k(E_2^{pq})$ para $k$ suficientemente grande (lo que implica que $E_\infty^{pq}=E_k^{pq}$).
  \item[(b)] Para todo $n$, la filtración $(F^p(E^n))_{p\in\bb{Z}}$ es \emph{discreta} y \emph{codiscreta} (también se dice \emph{finita}).
\end{enumerate}

Las sucesiones espectrales definidas en (T, 2.4) son, por tanto, sucesiones espectrales birregulares.
\end{env}

\begin{env}[11.1.4]
\label{0.11.1.4}
Supongamos que en la categoría $\cat{C}$ existen los límites inductivos filtrantes y que el funtor $\varinjlim$ es \emph{exacto} (lo que equivale a decir que se satisface el axioma (AB~5) de (T, 1.5) (cf.~T, 1.8)).
La condición de que la filtración $(F^p(X))_{p\in\bb{Z}}$ de un objeto $X\in\cat{C}$ sea exhaustiva se expresa entonces por $\varinjlim_{p\to-\infty}F^p(X)=X$.
Si una sucesión espectral $E$ es débilmente convergente, se tiene $B_\infty(E_2^{pq})=\varinjlim_{k\to\infty}B_k(E_2^{pq})$; si además $u:E\to E'$ es un morfismo de $E$ en una sucesión espectral débilmente convergente $E'$ en $\cat{C}$, entonces se tiene $u_2^{pq}(B_\infty(E_2^{pq}))=B_\infty(E_2^{\prime pq})$, por la exactitud de $\varinjlim$.
Además:
\end{env}

\begin{proposition}[11.1.5]
\label{0.11.1.5}
Sea $\cat{C}$ una categoría abeliana en la que los límites inductivos filtrantes son exactos, sean $E$ y $E'$ dos sucesiones espectrales regulares en $\cat{C}$, y sea $u:E\to E'$ un morfismo de sucesiones espectrales.
Si los $u_2^{pq}$ son isomorfismos, entonces $u$ también lo es.
\end{proposition}

\begin{proof}
Ya sabemos por \sref{0.11.1.2} que los $u_r^{pq}$ son isomorfismos y que
\[
  u_2^{pq}(B_\infty(E_2^{pq}))=B_\infty(E_2^{\prime pq});
\]
\oldpage[0\textsubscript{III}]{26}
la hipótesis de que $E$ y $E'$ son regulares implica también que $u_2^{pq}(Z_\infty(E_2^{pq}))=Z_\infty(E_2^{\prime pq})$ y, como $u_2^{pq}$ es un isomorfismo, también lo es $u_\infty^{pq}$; concluimos así que $\gr_p(u^{p+q})$ es asimismo un isomorfismo.
Pero, como las filtraciones de los $E^n$ y los $E^{\prime n}$ son discretas y exhaustivas, esto implica que los $u^n$ son también isomorfismos (Bourbaki, \emph{Alg. comm.}, cap.~III, \textsection2, n\textsuperscript{o}8, th.~1).
\end{proof}

\begin{env}[11.1.6]
\label{0.11.1.6}
De (\hyperref[0.11.1.1.2]{11.1.1.2}) y de la definición (\hyperref[0.11.1.1.3]{11.1.1.3}) se deduce que si, para una sucesión espectral $E$, se tiene $E_r^{pq}=0$, entonces $E_k^{pq}=0$ para $k\geq r$ y $E_\infty^{pq}=0$.
Se dice que una sucesión espectral \emph{degenera} si existe un entero $r\geq 2$ y, para todo entero $n\in\bb{Z}$, un entero $q(n)$ tal que \emph{$E_r^{n-q,q}=0$ para todo $q\neq q(n)$}.
De la observación anterior deducimos primero que también se tiene $E_k^{n-q,q}=0$ para $k\geq r$ (incluido $k=\infty$) y $q\neq q(n)$.
Además, la definición de $E_{r+1}^{pq}$ muestra que se tiene $E_{r+1}^{n-q(n),q(n)}=E_r^{n-q(n),q(n)}$; si $E$ es \emph{débilmente convergente}, se tiene también $E_\infty^{n-q(n),q(n)}=E_r^{n-q(n),q(n)}$; en otras palabras, para todo $n\in\bb{Z}$, $\gr_p(E^n)=0$ para $p\neq n-q(n)$ y $\gr_{n-q(n)}(E^n)=E_r^{n-q(n),q(n)}$.
Si además la filtración de $E^n$ es \emph{discreta} y \emph{exhaustiva}, entonces la sucesión espectral $E$ es \emph{regular}, y se tiene $E^n=E_r^{n-q(n),q(n)}$ salvo isomorfismo.
\end{env}

\begin{env}[11.1.7]
\label{0.11.1.7}
Supongamos que en la categoría $\cat{C}$ existen los límites inductivos filtrantes y son exactos, y sea $(E_\lambda,u_{\mu\lambda})$ un sistema inductivo (sobre un conjunto filtrante de índices) de sucesiones espectrales en $\cat{C}$.
Entonces el \emph{límite inductivo} de este sistema inductivo existe en la categoría aditiva de las sucesiones espectrales de objetos de $\cat{C}$: para verlo, basta definir $E_r^{pq}$, $d_r^{pq}$, $\alpha_r^{pq}$, $B_\infty(E_2^{pq})$, $Z_\infty(E_2^{pq})$, $E^n$, $F^p(E^n)$ y $\beta^{pq}$ como los límites inductivos respectivos de los $E_{r,\lambda}^{pq}$, $d_{r,\lambda}^{pq}$, $\alpha_{r,\lambda}^{pq}$, $B_\infty(E_{2,\lambda}^{pq})$, $Z_\infty(E_{2,\lambda}^{pq})$, $E_\lambda^n$, $F^p(E_\lambda^n)$ y $\beta_\lambda^{pq}$; la verificación de las condiciones de \sref{0.11.1.1} se deduce de la exactitud del funtor $\varinjlim$ sobre $\cat{C}$.
\end{env}

\begin{remark}[11.1.8]
\label{0.11.1.8}
Supongamos que la categoría $\cat{C}$ es la categoría de los $A$-módulos sobre un anillo $A$ \emph{noetheriano} (resp. un anillo $A$).
Las definiciones de \sref{0.11.1.1} muestran entonces que si, para un $r$ dado, los $E_r^{pq}$ son $A$-módulos \emph{de tipo finito} (resp. \emph{de longitud finita}), también lo son todos los módulos $E_s^{pq}$ para $s\geq r$ y, por tanto, también $E_\infty^{pq}$.
Si además la filtración del término límite $(E^n)$ es \emph{discreta} y \emph{codiscreta} para todo $n$, concluimos que cada $E^n$ es también un $A$-módulo \emph{de tipo finito} (resp. \emph{de longitud finita}).
\end{remark}

\begin{env}[11.1.9]
\label{0.11.1.9}
Habremos de considerar condiciones que aseguren que una sucesión espectral $E$ sea birregular de manera ``uniforme'' en $p+q=n$.
Utilizaremos entonces el lema siguiente:
\end{env}

\begin{lemma}[11.1.10]
\label{0.11.1.10}
Sea $(E_r^{pq})$ una familia de objetos de $\cat{C}$ relacionados por los datos \emph{(a)}, \emph{(b)} y \emph{(c)} de \sref{0.11.1.1}.
Para un entero fijo $n$, las propiedades siguientes son equivalentes:
\begin{enumerate}
  \item[{\rm(a)}] Existe un entero $r(n)$ tal que, para $r\geq r(n)$, $p+q=n$ o $p+q=n-1$, los morfismos $d_r^{pq}$ son nulos.
  \item[{\rm(b)}] Existe un entero $r(n)$ tal que, para $p+q=n$ o $p+q=n+1$, se tiene $B_r(E_2^{pq})=B_s(E_2^{pq})$ para $s\geq r\geq r(n)$.
  \item[{\rm(c)}] Existe un entero $r(n)$ tal que, para $p+q=n$ o $p+q=n-1$, se tiene $Z_r(E_2^{pq})=Z_s(E_2^{pq})$ para $s\geq r\geq r(n)$.
  \item[{\rm(d)}] Existe un entero $r(n)$ tal que, para $p+q=n$, se tiene $B_r(E_2^{pq})=B_s(E_2^{pq})$ y $Z_r(E_2^{pq})=Z_s(E_2^{pq})$ para $s\geq r\geq r(n)$.
\end{enumerate}
\end{lemma}

\begin{proof}
\oldpage[0\textsubscript{III}]{27}
Según las condiciones (a), (b) y (c) de \sref{0.11.1.1}, decir que $Z_{r+1}(E_2^{pq})=Z_r(E_2^{pq})$ equivale a decir que $d_r^{pq}=0$, y decir que $B_r(E_2^{p+r,q-r+1})=B_{r+1}(E_2^{p+r,q-r+1})$ equivale a decir que $d_r^{pq}=0$; el lema se deduce inmediatamente de esta observación.
\end{proof}

\subsection{La sucesión espectral de un complejo filtrado}
\label{subsection:0.11.2}

\begin{env}[11.2.1]
\label{0.11.2.1}
Sea $\cat{C}$ una categoría abeliana; convenimos en denotar mediante notaciones como $K^\bullet$ los \emph{complejos $(K^i)_{i\in\bb{Z}}$} de objetos de $\cat{C}$ cuyo diferencial es de grado $+1$, y mediante notaciones como $K_\bullet$ los complejos $(K_i)_{i\in\bb{Z}}$ de objetos de $\cat{C}$ cuyo diferencial es de grado $-1$.
A todo complejo $K^\bullet=(K^i)$ cuyo diferencial $d$ es de grado $+1$ podemos asociarle un complejo $K_\bullet'=(K_i')$ poniendo $K_i'=K^{-i}$, siendo el diferencial $K_i'\to K_{i-1}'$ el operador $d:K^{-i}\to K^{-i+1}$; y \emph{viceversa}, lo que permitirá, según las circunstancias, considerar uno u otro de los tipos de complejos y traducir todo resultado de uno de los tipos en resultados para el otro.
Denotamos análogamente mediante notaciones como $K^{\bullet\bullet}=(K^{ij})$ (resp. $K_{\bullet\bullet}=(K_{ij})$) los \emph{bicomplejos} (o \emph{complejos dobles}) de objetos de $\cat{C}$ en los que los \emph{dos} diferenciales son de grado $+1$ (resp. $-1$); todavía se puede pasar de un tipo al otro cambiando los signos de los índices, y tenemos notaciones y observaciones análogas para los multicomplejos cualesquiera.
Las notaciones $K^\bullet$ y $K_\bullet$ se utilizarán también para \emph{objetos $\bb{Z}$-graduados} de $\cat{C}$, que no son necesariamente complejos (pueden considerarse como tales para los diferenciales \emph{nulos}); por ejemplo, escribimos $\HH^\bullet(K^\bullet)=(\HH^i(K^\bullet))_{i\in\bb{Z}}$ para la \emph{cohomología} de un complejo $K^\bullet$ cuyo diferencial es de grado $+1$, y $\HH_\bullet(K_\bullet)=(\HH_i(K_\bullet))_{i\in\bb{Z}}$ para la \emph{homología} de un complejo $K_\bullet$ cuyo diferencial es de grado $-1$; cuando se pasa de $K^\bullet$ a $K_\bullet'$ por el procedimiento descrito anteriormente, se tiene $\HH_i(K_\bullet')=\HH^{-i}(K^\bullet)$.

Recordemos en este caso que, para un complejo $K^\bullet$ (resp. $K_\bullet$), escribiremos en general $Z^i(K^\bullet)=\Ker(K^i\to K^{i+1})$ (``objeto de cociclos'') y $B^i(K^\bullet)=\Im(K^{i-1}\to K^i)$ (``objeto de cobordes'') (resp. $Z_i(K_\bullet)=\Ker(K_i\to K_{i-1})$ (``objeto de ciclos'') y $B_i(K_\bullet)=\Im(K_{i+1}\to K_i)$ (``objeto de bordes'')), de modo que $\HH^i(K^\bullet)=Z^i(K^\bullet)/B^i(K^\bullet)$ (resp. $\HH_i(K_\bullet)=Z_i(K_\bullet)/B_i(K_\bullet)$).

Si $K^\bullet=(K^i)$ (resp. $K_\bullet=(K_i)$) es un complejo en $\cat{C}$ y $T:\cat{C}\to\cat{C}'$ un funtor de $\cat{C}$ en una categoría abeliana $\cat{C}'$, denotamos por $T(K^\bullet)$ (resp. $T(K_\bullet)$) el complejo $(T(K^i))$ (resp. $(T(K_i))$) en $\cat{C}'$.

No repasaremos la definición de los \emph{$\partial$-funtores} (T, 2.1), salvo para observar que se dice \emph{también} $\partial$-funtor en lugar de $\partial^*$-funtor cuando el morfismo $\partial$ disminuye el grado en una unidad, aclarando el contexto este punto si puede haber confusión.

Finalmente, se dice que un \emph{objeto graduado $(A_i)_{i\in\bb{Z}}$} de $\cat{C}$ está \emph{acotado inferiormente} (resp. \emph{superiormente}) si existe un $i_0$ tal que $A_i=0$ para $i<i_0$ (resp. $i>i_0$).
\end{env}

\begin{env}[11.2.2]
\label{0.11.2.2}
Sea $K^\bullet$ un complejo en $\cat{C}$ cuyo diferencial $d$ es de grado $+1$, y supongamos que está provisto de una \emph{filtración $F(K^\bullet)=(F^p(K^\bullet))_{p\in\bb{Z}}$} formada por subobjetos \emph{graduados}
\oldpage[0\textsubscript{III}]{28}
de $K^\bullet$, es decir, $F^p(K^\bullet)=(K^i\cap F^p(K^\bullet))_{i\in\bb{Z}}$; suponemos además que $d(F^p(K^\bullet))\subset F^p(K^\bullet)$ para todo $p\in\bb{Z}$.
Recordemos rápidamente cómo se define \emph{funtorialmente} una sucesión espectral $E(K^\bullet)$ a partir de $K^\bullet$ (M,~XV,~4 y G,~I,~4.3).
Para $r\geq 2$, el morfismo canónico $F^p(K^\bullet)/F^{p+r}(K^\bullet)\to F^p(K^\bullet)/F^{p+1}(K^\bullet)$ define un morfismo en cohomología
\[
  \HH^{p+q}(F^p(K^\bullet)/F^{p+r}(K^\bullet))\to\HH^{p+q}(F^p(K^\bullet)/F^{p+1}(K^\bullet)).
\]

Denotamos por $Z_r^{pq}(K^\bullet)$ la imagen de este morfismo.
Análogamente, de la sucesión exacta
\[
  0\to F^p(K^\bullet)/F^{p+1}(K^\bullet)\to F^{p-r+1}(K^\bullet)/F^{p+1}(K^\bullet)\to F^{p-r+1}(K^\bullet)/F^p(K^\bullet)\to 0,
\]
deducimos de la sucesión exacta de cohomología un morfismo
\[
  \HH^{p+q-1}(F^{p-r+1}(K^\bullet)/F^p(K^\bullet))\to\HH^{p+q}(F^p(K^\bullet)/F^{p+1}(K^\bullet)),
\]
y denotamos por $B_r^{pq}(K^\bullet)$ la imagen de este morfismo; se demuestra que $B_r^{pq}(K^\bullet)\subset Z_r^{pq}(K^\bullet)$ y se toma $E_r^{pq}(K^\bullet)=Z_r^{pq}(K^\bullet)/B_r^{pq}(K^\bullet)$; no precisaremos la definición de los $d_r^{pq}$ ni de los $\alpha_r^{pq}$.

Observemos aquí que todos los $Z_r^{pq}(K^\bullet)$ y $B_r^{pq}(K^\bullet)$, para $p$ y $q$ fijos, son subobjetos del mismo objeto $\HH^{p+q}(F^p(K^\bullet)/F^{p+1}(K^\bullet))$, que denotamos por $Z_1^{pq}(K^\bullet)$; ponemos $B_1^{pq}(K^\bullet)=0$, de modo que las definiciones anteriores de $Z_r^{pq}(K^\bullet)$ y $B_r^{pq}(K^\bullet)$ valen también para $r=1$; ponemos asimismo $E_1^{pq}(K^\bullet)=Z_1^{pq}(K^\bullet)$.
Definimos $d_1^{pq}$ y $\alpha_1^{pq}$ de modo que se satisfagan las condiciones de \sref{0.11.1.1} para $r=1$.
Por otra parte, definimos los subobjetos $Z_\infty^{pq}(K^\bullet)$ como la imagen del morfismo
\[
  \HH^{p+q}(F^p(K^\bullet))\to\HH^{p+q}(F^p(K^\bullet)/F^{p+1}(K^\bullet))=E_1^{pq}(K^\bullet),
\]
y $B_\infty^{pq}(K^\bullet)$ como la imagen del morfismo
\[
  \HH^{p+q-1}(K^\bullet/F^p(K^\bullet))\to\HH^{p+q}(F^p(K^\bullet)/F^{p+1}(K^\bullet))=E_1^{pq}(K^\bullet),
\]
inducido como antes por la sucesión exacta de cohomología.
Definimos $Z_\infty(E_2^{pq}(K^\bullet))$ y $B_\infty(E_2^{pq}(K^\bullet))$ como las imágenes canónicas en $E_2^{pq}(K^\bullet)$ de $Z_\infty^{pq}(K^\bullet)$ y $B_\infty^{pq}(K^\bullet)$.

Finalmente, denotamos por $F^p(\HH^n(K^\bullet))$ la imagen en $\HH^n(K^\bullet)$ del morfismo $\HH^n(F^p(K^\bullet))\to\HH^n(K^\bullet)$ inducido por la inyección canónica $F^p(K^\bullet)\to K^\bullet$; por la sucesión exacta de cohomología, este es también el núcleo del morfismo $\HH^n(K^\bullet)\to\HH^n(K^\bullet/F^p(K^\bullet))$.
Esto define una filtración sobre $E^n(K^\bullet)=\HH^n(K^\bullet)$; no daremos aquí la definición de los isomorfismos $\beta^{pq}$.
\end{env}

\begin{env}[11.2.3]
\label{0.11.2.3}
El carácter \emph{funtorial} de $E(K^\bullet)$ se entiende del modo siguiente: dados dos complejos \emph{filtrados} $K^\bullet$ y $K^{\prime\bullet}$ en $\cat{C}$ y un morfismo de complejos $u:K^\bullet\to K^{\prime\bullet}$ \emph{compatible con las filtraciones}, se inducen de manera evidente los morfismos $u_r^{pq}$ (para $r\geq 1$) y $u^n$, y se demuestra que estos morfismos son compatibles con los $d_r^{pq}$, $\alpha_r^{pq}$ y $\beta^{pq}$ en el sentido de \sref{0.11.1.2}, y definen así un morfismo bien definido $E(u):E(K^\bullet)\to E(K^{\prime\bullet})$ de sucesiones espectrales.
Además, se demuestra que si $u$ y $v$ son morfismos $K^\bullet\to K^{\prime\bullet}$ del tipo anterior, \emph{homótopos de orden $\leq k$}, entonces $u_r^{pq}=v_r^{pq}$ para $r>k$ y $u^n=v^n$ para todo $n$ (M,~XV,~3.1).
\end{env}

\begin{env}[11.2.4]
\label{0.11.2.4}
\oldpage[0\textsubscript{III}]{29}
Supongamos que los límites inductivos filtrantes en $\cat{C}$ son exactos.
Entonces, si la filtración $(F^p(K^\bullet))$ de $K^\bullet$ es \emph{exhaustiva}, también lo es la filtración $(F^p(\HH^n(K^\bullet)))$ para todo $n$, pues por hipótesis se tiene $K^\bullet=\varinjlim_{p\to-\infty}F^p(K^\bullet)$ y la hipótesis sobre $\cat{C}$ implica que la cohomología conmuta con los límites inductivos.
Además, por la misma razón, se tiene $B_\infty(E_2^{pq}(K^\bullet))=\sup_k B_k(E_2^{pq}(K^\bullet))$.
Se dice que la filtración $(F^p(K^\bullet))$ de $K^\bullet$ es \emph{regular} si para todo $n$ existe un entero $u(n)$ tal que $\HH^n(F^p(K^\bullet))=0$ para $p>u(n)$.
Este es, en particular, el caso cuando la filtración de $K^\bullet$ es \emph{discreta}.
Cuando la filtración de $K^\bullet$ es regular y exhaustiva, y los límites inductivos filtrantes son exactos en $\cat{C}$, se tiene (M,~XV,~4) que la sucesión espectral $E(K^\bullet)$ es \emph{regular}.
\end{env}

\subsection{Las sucesiones espectrales de un bicomplejo}
\label{subsection:0.11.3}

\begin{env}[11.3.1]
\label{0.11.3.1}
En cuanto a las convenciones para los bicomplejos, seguimos las de~(T,~2.4) antes que las de~(M), suponiéndose así que los dos diferenciales $d'$ y $d''$ (de grado~$+1$) de un tal bicomplejo $K^{\bullet\bullet}=(K^{ij})$ son \emph{permutables}.
Supongamos que se satisface \emph{una} de las dos condiciones siguientes:
\begin{enumerate}
   \item Existen \emph{sumas directas infinitas} en $\cat{C}$;
   \item Para todo $n\in\bb{Z}$, solo hay un número \emph{finito} de pares $(p,q)$ tales que $p+q=n$ y $K^{pq}\neq 0$.
 \end{enumerate}
Entonces el bicomplejo $K^{\bullet\bullet}$ define un \emph{complejo} (simple) $(K^{\prime n})_{n\in\bb{Z}}$ con $K^{\prime n}=\sum_{i+j=n}K^{ij}$, viniendo dado el diferencial $d$ (de grado~$+1$) de este complejo por $dx=d'x+(-1)^i d''x$ para $x\in K^{ij}$.
Cuando hablemos más adelante del \emph{complejo (simple) definido por un bicomplejo $K^{\bullet\bullet}$}, se sobreentenderá siempre que se satisface una de las condiciones anteriores.
Adoptamos las convenciones análogas para los multicomplejos.

Denotamos por $K^{i,\bullet}$ (resp.~$K^{\bullet,j}$) el complejo simple $(K^{ij})_{j\in\bb{Z}}$ (resp.~$(K^{ij})_{i\in\bb{Z}}$), y por $Z_\mathrm{II}^p(K^{i,\bullet})$, $B_\mathrm{II}^p(K^{i,\bullet})$, $\HH_\mathrm{II}^p(K^{i,\bullet})$ (resp.~$Z_\mathrm{I}^p(K^{\bullet,j})$, $B_\mathrm{I}^p(K^{\bullet,j})$, $\HH_\mathrm{I}^p(K^{\bullet,j})$) sus objetos de cociclos, de cobordes y de cohomología de grado $p$, respectivamente;
el diferencial $d':K^{i,\bullet}\to K^{i+1,\bullet}$ es un morfismo de complejos, que da así un operador sobre los cociclos, los cobordes y la cohomología,
\begin{align*}
  d'&:Z_\mathrm{II}^p(K^{i,\bullet})\to Z_\mathrm{II}^p(K^{i+1,\bullet})
\\d'&:B_\mathrm{II}^p(K^{i,\bullet})\to B_\mathrm{II}^p(K^{i+1,\bullet})
\\d'&:\HH_\mathrm{II}^p(K^{i,\bullet})\to\HH_\mathrm{II}^p(K^{i+1,\bullet})
\end{align*}
y está claro que, para estos operadores, $(Z_\mathrm{II}^p(K^{i,\bullet}))_{i\in\bb{Z}}$, $(B_\mathrm{II}^p(K^{i,\bullet}))_{i\in\bb{Z}}$ y $(\HH_\mathrm{II}^p(K^{i,\bullet}))_{i\in\bb{Z}}$ son complejos;
denotamos el complejo $(\HH_\mathrm{II}^p(K^{i,\bullet}))_{i\in\bb{Z}}$ por $\HH_\mathrm{II}^p(K^{\bullet\bullet})$, y sus objetos de cociclos, cobordes y cohomología de grado $q$ por $Z_\mathrm{I}^q(\HH_\mathrm{II}^p(K^{\bullet\bullet}))$, $B_\mathrm{I}^q(\HH_\mathrm{II}^p(K^{\bullet\bullet}))$ y $\HH_\mathrm{I}^q(\HH_\mathrm{II}^p(K^{\bullet\bullet}))$.
Definimos análogamente los complejos $\HH_\mathrm{I}^p(K^{\bullet\bullet})$ y sus objetos de cohomología $\HH_\mathrm{II}^q(\HH_\mathrm{I}^p(K^{\bullet\bullet}))$.
Recordemos, no obstante, que $\HH^n(K^{\bullet\bullet})$ denota el objeto de grado $n$ de la cohomología del complejo \emph{(simple)} definido por $K^{\bullet\bullet}$.
\end{env}

\begin{env}[11.3.2]
\label{0.11.3.2}
Sobre el complejo definido por un bicomplejo $K^{\bullet\bullet}$ pueden considerarse dos filtraciones canónicas, $(F_\mathrm{I}^p(K^{\bullet\bullet}))$ y $(F_\mathrm{II}^p(K^{\bullet\bullet}))$, dadas por
\[
\label{0.11.3.2.1}
  F_\mathrm{I}^p(K^{\bullet\bullet})
  = \left(\sum_{i+j=n,i\geq p}K^{ij}\right)_{n\in\bb{Z}}
  \quad\text{y}\quad
  F_\mathrm{II}^p(K^{\bullet\bullet})
  = \left(\sum_{i+j=n,j\geq p}K^{ij}\right)_{n\in\bb{Z}}
  \tag{11.3.2.1}
\]
\oldpage[0\textsubscript{III}]{30}
que, por definición, son subobjetos graduados del complejo (simple) definido por $K^{\bullet\bullet}$ y hacen así de este complejo un complejo filtrado;
además, está claro que estas filtraciones son \emph{exhaustivas} y \emph{separadas}.

A cada una de estas filtraciones le corresponde una sucesión espectral~\sref{0.11.2.2};
denotamos por $'E(K^{\bullet\bullet})$ y $''E(K^{\bullet\bullet})$ las sucesiones espectrales correspondientes a $(F_\mathrm{I}^p(K^{\bullet\bullet}))$ y $(F_\mathrm{II}^p(K^{\bullet\bullet}))$, respectivamente, llamadas \emph{sucesiones espectrales del bicomplejo $K^{\bullet\bullet}$}, y que tienen ambas por término límite la cohomología $(\HH^n(K^{\bullet\bullet}))$.
Se demuestra además~(M,~XV,~6) que se tiene
\[
\label{0.11.3.2.2}
  'E_2^{pq}(K^{\bullet\bullet})
  = \HH_\mathrm{I}^p(\HH_\mathrm{II}^q(K^{\bullet\bullet}))
  \quad\text{y}\quad
  ''E_2^{pq}(K^{\bullet\bullet})
  = \HH_\mathrm{II}^q(\HH_\mathrm{I}^p(K^{\bullet\bullet})).
  \tag{11.3.2.2}
\]

Todo morfismo $u:K^{\bullet\bullet}\to K^{\prime\bullet\bullet}$ de bicomplejos es \emph{ipso facto} compatible con las filtraciones del mismo tipo de $K^{\bullet\bullet}$ y $K^{\prime\bullet\bullet}$, y define así un morfismo para cada una de las dos sucesiones espectrales;
además, dos morfismos \emph{homótopos} definen una homotopía \emph{de orden $\leq 1$} de los complejos (simples) filtrados correspondientes, y por tanto el \emph{mismo} morfismo para cada una de las dos sucesiones espectrales~(M,~XV,~6.1).
\end{env}

\begin{proposition}[11.3.3]
\label{0.11.3.3}
Sea $K^{\bullet\bullet}=(K^{ij})$ un bicomplejo en una categoría abeliana $\cat{C}$.
\begin{enumerate}
  \item[{\rm(i)}] Si existen $i_0$ y $j_0$ tales que $K^{ij}=0$ para $i<i_0$ o $j<j_0$ (resp.~$i>i_0$ o $j>j_0$), entonces las dos sucesiones espectrales $'E(K^{\bullet\bullet})$ y $''E(K^{\bullet\bullet})$ son birregulares.
  \item[{\rm(ii)}] Si existen $i_0$ e $i_1$ tales que $K^{ij}=0$ para $i<i_0$ o $i>i_1$ (resp.~si existen $j_0$ y $j_1$ tales que $K^{ij}=0$ para $j<j_0$ o $j>j_1$), entonces las dos sucesiones espectrales $'E(K^{\bullet\bullet})$ y $''E(K^{\bullet\bullet})$ son birregulares.
\end{enumerate}
Supongamos además que los límites inductivos filtrantes existen y son exactos en $\cat{C}$. Entonces:
\begin{enumerate}
  \item[{\rm(iii)}] Si existe $i_0$ tal que $K^{ij}=0$ para $i>i_0$ (resp.~si existe $j_0$ tal que $K^{ij}=0$ para $j<j_0$), entonces la sucesión espectral $'E(K^{\bullet\bullet})$ es regular.
  \item[{\rm(iv)}] Si existe $i_0$ tal que $K^{ij}=0$ para $i<i_0$ (resp.~si existe $j_0$ tal que $K^{ij}=0$ para $j>j_0$), entonces la sucesión espectral $''E(K^{\bullet\bullet})$ es regular.
\end{enumerate}
\end{proposition}

\begin{proof}
La proposición se deduce inmediatamente de las definiciones~\sref{0.11.1.3} y de~\sref{0.11.2.4}, así como de las observaciones siguientes relativas a la filtración $F_\mathrm{I}$ (y de las observaciones análogas que se deducen para $F_\mathrm{II}$ intercambiando los papeles de los dos índices en $K^{\bullet\bullet}$):
\begin{enumerate}
  \item[{$1^{\circ}$}] Si existe $i_0$ tal que $K^{ij}=0$ para $i>i_0$, entonces la filtración $F_\mathrm{I}(K^{\bullet\bullet})$ es \emph{discreta}.
  \item[{$2^{\circ}$}] Si existe $i_0$ tal que $K^{ij}=0$ para $i<i_0$, entonces la filtración $F_\mathrm{I}(K^{\bullet\bullet})$ es \emph{codiscreta}.
    Se deduce inmediatamente que ocurre lo mismo para la filtración correspondiente $F_\mathrm{I}(\HH^n(K^{\bullet\bullet}))$ para todo $n$.
    Además, la definición de los $B_{r}^{pq}$ correspondientes a la filtración $F_\mathrm{I}(K^{\bullet\bullet})$ \sref{0.11.2.2} muestra que, para todo par $(p,q)$, la sucesión $(B_{r}^{pq})_{r\geq 2}$ es estable.
  \item[{$3^{\circ}$}] Si existe $j_0$ tal que $K^{ij}=0$ para $j<j_0$, entonces se tiene
    \[
      F_\mathrm{I}^{p+r}(K^{\bullet\bullet})
      \cap (\sum_{i+j=n}K^{ij})
      = 0
    \]
    desde que $p+r+j_{0}>n$, luego $Z_{r}^{pq} = Z_\infty(E_2^{pq})$ para $r>q-j_{0}+1$.
    Por otra parte, $\HH^{n}(F_\mathrm{I}^{p}(K^{\bullet\bullet})) = 0$ para $p>n-j_{0}+1$.
  \item[{$4^{\circ}$}] Si existe $j_0$ tal que $K^{ij}=0$ para $j>j_0$, entonces se tiene
    \[
      F_\mathrm{I}^{p-r+1}(K^{\bullet\bullet})
      \cap (\sum_{i+j=n}K^{ij})
      = \sum_{i+j=n}K^{ij}
    \]
    \oldpage[0\textsubscript{III}]{31}
    desde que $p-r+1+j_{0}<n$, luego $B_{r}^{pq} = B_\infty(E_2^{pq})$ para $r>j_{0}-q+1$.
    Por otra parte, $\HH^{n}(F_\mathrm{I}^{p}(K^{\bullet\bullet})) = \HH^{n}(K^{\bullet\bullet})$ para $p+j_{0}<n-1$.
\end{enumerate}
\end{proof}

\begin{env}[11.3.4]
\label{0.11.3.4}
Supongamos que el bicomplejo $K^{\bullet\bullet}=(K^{ij})$ es tal que $K^{ij}=0$ para $i<0$ o $j<0$.
Sabemos que se puede definir, para todo $p\in\bb{Z}$, un ``homomorfismo de borde'' canónico.
\[
\label{0.11.3.4.1}
  'E_2^{p0}(K^{\bullet\bullet})
  \rightarrow \HH^{p}(K^{\bullet\bullet})
\tag{11.3.4.1}
\]
(M,~XV,~6).
Recordemos que esto se debe, por una parte, a la sucesión espectral
$'E(K^{\bullet\bullet})$, pues $Z_r^{p0}=Z_{\mathrm{I}}^{p}(Z_{\mathrm{II}}^{0}(K^{\bullet\bullet}))$ para $2\leq r \leq +\infty$, y, por otra parte, a que $\HH^{p}(F_\mathrm{I}^{p+1}(K^{\bullet\bullet})) = 0$, de modo que el isomorfismo $\beta^{p0}:\ 'E_\infty^{p0} \simto \HH^{p}(F_\mathrm{I}^p)/\HH^{p}(F_\mathrm{I}^{p+1})$ da un homomorfismo $'E_\infty^{p0} \rightarrow \HH^{p}(F_\mathrm{I}^p(K^{\bullet\bullet})) \rightarrow \HH^{p}(K^{\bullet\bullet})$.
La igualdad de todos los $Z_r^{p0}$ permite definir homomorfismos canónicos $'E_r^{p0} \rightarrow 'E_s^{p0}$ para $r \leq s$ y, en particular, un homomorfismo $'E_2^{p0} \rightarrow 'E_\infty^{p0}$, de donde, por composición, el homomorfismo de borde $'E_2^{p0}(K^{\bullet\bullet}) \rightarrow \HH^{p}(K^{\bullet\bullet})$.
Además, se verifica inmediatamente que el homomorfismo de borde así definido envía la clase módulo $B_2^{p0}$ de un elemento $z\in Z_\mathrm{II}^0(K^{\bullet\bullet})\subset K^{p0}$ tal que $d'z=0$ en la clase de $z$ módulo $B_\infty^{p0}$ en $'E_\infty^{p0}$, y después en la clase de cohomología de $z$ en $\HH^p(K^{\bullet\bullet})$.
Vemos, pues, finalmente que el \emph{homomorfismo de borde} \sref{0.11.3.4.1} \emph{proviene, al pasar a la cohomología, de la inyección canónica} $Z_\mathrm{II}^0(K^{\bullet\bullet}) \rightarrow K^{\bullet\bullet}$ (donde $K^{\bullet\bullet}$ se considera como complejo simple).
Interpretamos naturalmente del mismo modo el homomorfismo de borde
\[
\label{0.11.3.4.2}
  ''E_2^{p0}(K^{\bullet\bullet})
  \rightarrow \HH^{p}(K^{\bullet\bullet})
\tag{11.3.4.2}
\]
como proveniente de la inyección canónica $Z_\mathrm{I}^0(K^{\bullet\bullet}) \rightarrow K^{\bullet\bullet}$.
\end{env}

\begin{env}[11.3.5]
\label{0.11.3.5}
Sea ahora $K_{\bullet\bullet} = (K_{ij})$ un bicomplejo en $\cat{C}$ cuyos dos operadores diferenciales son de grado~$-1$.
Escribimos entonces $K_{i,\bullet}$ (resp. $K_{\bullet,j}$) para designar el complejo simple $(K_{ij})_{j\in\bb{Z}}$ (resp. $(K_{ij})_{i\in\bb{Z}}$), $\HH_{p}^\mathrm{II}(K_{i, \bullet})$ (resp. $\HH_{p}^\mathrm{I}(K_{\bullet, j})$) para designar el objeto de homología de grado $p$, $\HH_{p}^\mathrm{II}(K_{\bullet\bullet})$ (resp. $\HH_{p}^\mathrm{I}(K_{\bullet\bullet})$) para designar el complejo $(\HH_{p}^\mathrm{II}(K_{i, \bullet}))_{i \in\bb{Z}}$ (resp. $(\HH_{p}^\mathrm{I}(K_{\bullet, j}))_{j \in\bb{Z}}$), y $\HH_q^\mathrm{I}(\HH_{p}^\mathrm{II}(K_{\bullet\bullet}))$ (resp. $\HH_q^\mathrm{II}(\HH_{p}^\mathrm{I}(K_{\bullet\bullet}))$) para designar el objeto de homología de grado $q$;
utilizamos notaciones análogas para los objetos de ciclos y los objetos de bordes;
finalmente, denotaremos por $\HH_n(K_{\bullet\bullet})$ el objeto de homología de grado $n$ (cuando exista) del complejo simple (con un operador diferencial de grado~$-1$) definido por $K_{\bullet\bullet}$.
Sea $K'^{\bullet\bullet}=(K'^{ij})$, con $K'^{ij}=K_{-i,-j}$, el bicomplejo con operadores diferenciales de grado~$+1$ asociado a $K_{\bullet\bullet}$.
Por definición, las \emph{sucesiones espectrales} de $K_{\bullet\bullet}$ son las de $K'^{\bullet\bullet}$, que denotamos por $'E(K_{\bullet\bullet})$ y $''E(K_{\bullet\bullet})$, cambiando, sin embargo, la notación al poner
\[
  'E_{pq}^{r}(K_{\bullet\bullet})
  = 'E_{r}^{-p,-q}(K'^{\bullet\bullet})
  \quad\text{y}\quad
  ''E_{pq}^{r}(K_{\bullet\bullet})
  = ''E_{r}^{-p,-q}(K'^{\bullet\bullet}), 
\]
para $2\leq r \leq \infty$.
Con esta notación, se tiene
\[
  'E_{pq}^{2}(K_{\bullet\bullet})
  = \HH_{p}^\mathrm{I}(\HH_{q}^\mathrm{II}(K_{\bullet\bullet}))
  \quad\text{y}\quad
  ''E_{pq}^{2}(K_{\bullet\bullet})
  = \HH_{q}^\mathrm{II}(\HH_{p}^\mathrm{I}(K_{\bullet\bullet})).
\]
Para evitar errores de signo, será preferible, en general, volver al complejo $K'^{\bullet\bullet}$ para las relaciones entre estas sucesiones espectrales y sus convergencias.
Señalemos, no obstante, los criterios correspondientes a \sref{0.11.3.3}:
\end{env}

\begin{env}[11.3.6]
\label{0.11.3.6}
Las sucesiones espectrales $'E(K_{\bullet\bullet})$ y $''E(K_{\bullet\bullet})$ son \emph{birregulares} en los casos siguientes:
\begin{enumerate}
  \item[(a)] Existen $i_0$ y $j_0$ tales que $K_{ij}=0$ para $i>i_0$ o $j>j_0$ (resp. para $i<i_0$
    \oldpage[0\textsubscript{III}]{32}
    o $j<j_0$);
  \item[(b)] Existen $i_0$ e $i_1$ tales que $K_{ij}=0$ para $i<i_0$ y $i>i_1$;
  \item[(c)] Existen $j_0$ y $j_1$ tales que $K_{ij}=0$ para $j<j_0$ y $j>j_1$.
\end{enumerate}

La sucesión $'E(K_{\bullet\bullet})$ es \emph{regular} si existe $i_0$ tal que $K_{ij}=0$ para $i<i_0$, o si existe $j_0$ tal que $K_{ij}=0$ para $j>j_0$.

La sucesión $''E(K_{\bullet\bullet})$ es \emph{regular} si existe $i_0$ tal que $K_{ij}=0$ para $i>i_0$, o si existe $j_0$ tal que $K_{ij}=0$ para $j<j_0$.
\end{env}


\subsection{Hipercohomología de un funtor con respecto a un complejo $K^{\bullet}$}
\label{subsection:0.11.4}

\begin{env}[11.4.1]
\label{0.11.4.1}
Sea $\cat{C}$ una categoría abeliana.
Recordemos que se ha definido una \emph{resolución por la derecha} (o \emph{resolución cohomológica}) de un objeto $A$ de $\cat{C}$ como un complejo de objetos de $\cat{C}$, cuyo operador diferencial es de grado $+1$,
\[
  0 \rightarrow L^0 \rightarrow L^1 \rightarrow L^2 \rightarrow \ldots
\]
provisto de un morfismo $\varepsilon: A \rightarrow L^0$, llamado la \emph{aumentación} de la resolución, que puede considerarse como un morfismo de complejos
\[
  \xymatrix{
    0 \ar[r]
    & A \ar[r] \ar[d]_{\varepsilon}
    & 0 \ar[r] \ar[d]
    & 0 \ar[r] \ar[d]
    & \ldots
  \\0 \ar[r]
    & L^0 \ar[r]
    & L^1 \ar[r]
    & L^2 \ar[r]
    & \ldots
  }
\]
tal que la sucesión
\[
  0 \rightarrow A \xrightarrow{\varepsilon} L^0 \rightarrow L^1 \longrightarrow \ldots
\]
sea exacta.
Análogamente, una \emph{resolución por la izquierda} (o \emph{resolución homológica}) de $A$ es un complejo $0 \leftarrow L_0 \leftarrow L_1 \leftarrow L_2 \leftarrow \ldots$ de objetos de $\cat{C}$, cuyo operador diferencial es de grado~$-1$, provisto de una \emph{aumentación} $L_0 \xrightarrow{\varepsilon} A$ tal que la sucesión
\[
  0 \leftarrow A \xleftarrow{\varepsilon} L_0 \leftarrow L_1 \leftarrow \ldots
\]
sea exacta.

Cuando la resolución por la derecha $(L^i)_{i\geq 0}$ de un objeto $A$ es tal que $L^i=0$ para $i\geq n+1$, se dice que esta resolución es \emph{de longitud $\leq n$}.
Se define análogamente una resolución por la izquierda de longitud $\leq n$.
Una resolución que es de longitud $\leq n$ para algún entero $n$ se dice \emph{finita}.

Una resolución de $A$ se llama \emph{proyectiva} (resp. \emph{inyectiva}) si los objetos de $\cat{C}$ distintos de $A$ que la componen son \emph{proyectivos} (resp. \emph{inyectivos}).
Cuando $\cat{C}$ es la categoría de los módulos (a izquierda, por ejemplo) sobre un anillo, se dice que una resolución de $A$ es \emph{plana} (resp. \emph{libre}) cuando los módulos distintos de $A$ que la componen son \emph{planos} (resp. \emph{libres}).
\end{env}

\begin{env}[11.4.2]
\label{0.11.4.2}
Sea $K^\bullet = (K^i)_{i\in\bb{Z}}$ un complejo de objetos de $\cat{C}$, cuyo operador diferencial es de grado~$+1$.
Se define una \emph{resolución de Cartan--Eilenberg por la derecha} de $K^\bullet$ como un par formado por un bicomplejo $L^{\bullet\bullet} = (L^{ij})$ con operadores diferenciales de grado~$+1$ y tal que $L^{ij}=0$ para $j<0$, y un morfismo de complejos simples $\varepsilon: K^\bullet \rightarrow L^{\bullet, 0}$, de modo que se cumplan las condiciones siguientes:
\oldpage[0\textsubscript{III}]{33}
\begin{enumerate}
  \item[(i)] Para cada índice $i$, las sucesiones
    \[
      \begin{aligned}
        0 \rightarrow
        & K^i
        \xrightarrow{\varepsilon} L^{i0}
        \rightarrow L^{i1}
        \rightarrow \ldots
      \\0 \rightarrow
        & B^i(K^\bullet)
        \xrightarrow{\varepsilon} B_\mathrm{I}^i(L^{\bullet,0})
        \rightarrow B_\mathrm{I}^i(L^{\bullet,1})
        \rightarrow \ldots
      \\0 \rightarrow
        & Z^i(K^\bullet)
        \xrightarrow{\varepsilon} Z_\mathrm{I}^i(L^{\bullet,0})
        \rightarrow Z_\mathrm{I}^i(L^{\bullet,1})
        \rightarrow \ldots
      \\0 \rightarrow
        & \HH^i(K^\bullet)
        \xrightarrow{\varepsilon} \HH_\mathrm{I}^i(L^{\bullet,0})
        \rightarrow \HH_\mathrm{I}^i(L^{\bullet,1})
        \rightarrow \ldots
      \end{aligned}
    \]
    son exactas o, en otras palabras, $(L^{i,\bullet}), (B_\mathrm{I}^i(L^{\bullet\bullet})), (Z_\mathrm{I}^i(L^{\bullet\bullet}))$ y $(\HH_\mathrm{I}^i(L^{\bullet\bullet}))$ son \emph{resoluciones} de $K^i, B^i(K^\bullet), Z^i(K^\bullet)$ y $\HH^i(K^\bullet)$, respectivamente.
  \item[(ii)] Para cada $j$, el complejo simple $L^{\bullet, j}$ es \emph{escindido} o, en otras palabras, las sucesiones exactas
    \[
      \label{0.11.4.2.1}
      0 \rightarrow B_\mathrm{I}^i(L^{\bullet,j}) \rightarrow Z_\mathrm{I}^i(L^{\bullet,j}) \rightarrow \HH_\mathrm{I}^i(L^{\bullet,j}) \rightarrow 0
      \tag{11.4.2.1}
    \]
    \[
      \label{0.11.4.2.2}
      0 \rightarrow Z_\mathrm{I}^i(L^{\bullet,j}) \rightarrow L^{ij} \rightarrow B_\mathrm{I}^{i+1}(L^{\bullet,j}) \rightarrow 0
      \tag{11.4.2.2}
    \]
    son \emph{escindidas}.
\end{enumerate}

Se demuestra (M,~XVII,~1.2) que si todo objeto de $\cat{C}$ es subobjeto de un objeto inyectivo, entonces todo complejo $K^\bullet$ de $\cat{C}$ admite una resolución de Cartan--Eilenberg \emph{inyectiva}, es decir, formada por objetos inyectivos $L^{ij}$ (la condición (ii) anterior implica entonces que los $(B_\mathrm{I}^i(L^{\bullet,j})), (Z_\mathrm{I}^i(L^{\bullet,j}))$ y $(\HH_\mathrm{I}^i(L^{\bullet,j}))$ están también formados por objetos inyectivos).
Además, para todo morfismo $f: K^\bullet \rightarrow K'^\bullet$ de complejos de $\cat{C}$,
toda resolución de Cartan--Eilenberg $L^{\bullet\bullet}$ de $K^\bullet$ y toda resolución de Cartan--Eilenberg \emph{inyectiva} $L'^{\bullet\bullet}$ de $K'^\bullet$, existe un morfismo de bicomplejos $F: L^{\bullet\bullet} \rightarrow L'^{\bullet\bullet}$ compatible con $f$ y las aumentaciones, y si $f$ y $g$ son morfismos homótopos de $K^\bullet$ en $K'^\bullet$, entonces los morfismos correspondientes de $L^{\bullet\bullet}$ en $L'^{\bullet\bullet}$ son también homótopos \emph{(loc. cit.)}.

Cuando $K^{\bullet}$ está \emph{acotado inferiormente} (resp. \emph{superiormente}), puede tomarse $L^{\bullet\bullet}$ tal que $L^{ij}=0$ para $i<i_0$ (resp. $i>i_0$) si $K^i=0$ para $i<i_0$ (resp. $i>i_0$) (M,~XVII,~1.3).

Supongamos, por otra parte, que existe un entero $n$ tal que todo objeto de $\cat{C}$ admite una \emph{resolución inyectiva de longitud $\leq n$};
entonces puede suponerse que $L^{ij}=0$ para $j>n$ (M,~XVII,~1.4).
\end{env}


\begin{env}[11.4.3]
\label{0.11.4.3}
Sea ahora $T$ un \emph{funtor covariante aditivo} de $\cat{C}$ en una categoría abeliana $\cat{C'}$.
Dado un complejo $K^\bullet$ de $\cat{C}$ y una resolución de Cartan--Eilenberg \emph{inyectiva} $L^{\bullet\bullet}$ de $K^{\bullet}$, supongamos que existe el complejo (simple) definido por el bicomplejo $T(L^{\bullet\bullet})$ \sref{0.11.3.1};
entonces las dos sucesiones espectrales $'E(T(L^{\bullet\bullet}))$ y $''E(T(L^{\bullet\bullet}))$ de este bicomplejo se llaman las \emph{sucesiones espectrales de hipercohomología} de $T$ con respecto a $K^\bullet$;
por \sref{0.11.4.2} y \sref{0.11.3.2}, dependen únicamente del complejo $K^\bullet$ y no de la elección de la resolución de Cartan--Eilenberg inyectiva $L^{\bullet\bullet}$;
además, dependen \emph{funtorialmente} de $K^\bullet$.
Tienen el mismo término límite $\HH^\bullet(T(L^{\bullet\bullet}))$, llamado también la \emph{hipercohomología} de $T$ con respecto a $K^\bullet$, y denotado por $\RR^\bullet T(K^\bullet)$.
Se demuestra que los términos $E_2$ de las dos sucesiones espectrales anteriores vienen dados por
\[
\label{0.11.4.3.1}
  'E_2^{pq}
  = \HH^p(\RR^q T(K^\bullet))
\tag{11.4.3.1}
\]
\[
\label{0.11.4.3.2}
  ''E_2^{pq}
  = \RR^p T(\HH^q(K^\bullet))
\tag{11.4.3.2}
\]
\oldpage[0\textsubscript{III}]{34}
donde $\RR^p T$ denota, como de costumbre, el \emph{funtor derivado} de grado $p$ de $T$ para $p\in\bb{Z}$;
$\RR^q T(K^\bullet)$ denota el complejo $(\RR^q T(K^i))_{i\in\bb{Z}}$.
\phantomsection\label{0.11.4.4}
Salvo mención explícita en contrario, supondremos en lo sucesivo que todo objeto de $\cat{C}$ es subobjeto de un objeto inyectivo de $\cat{C}$, de modo que existen resoluciones de Cartan--Eilenberg inyectivas para todo complejo de $\cat{C}$.
Como $L^{ij}=0$ para $j<0$, los criterios de \sref{0.11.3.3} muestran que las \emph{dos} sucesiones espectrales de hipercohomología de $T$ con respecto a $K^\bullet$ existen y son \emph{birregulares} en cada uno de los dos casos siguientes:
\begin{enumerate}
  \item $K^\bullet$ está acotado inferiormente;
  \item Todo objeto de $\cat{C}$ admite una resolución inyectiva de longitud a lo sumo $n$, para algún entero $n$ independiente del objeto considerado.
\end{enumerate}
En efecto, puede suponerse, en el primer caso, por \sref{0.11.4.2}, que existe $i_0$ tal que $L^{ij}=0$ para $i<i_0$, y, en el segundo caso, que existe $j_1$ tal que $L^{ij}=0$ para $j>j_1$;
en cada uno de estos dos casos, está además claro que, para cada $n$ dado, solo existe un número finito de pares $(i,j)$ tales que $L^{ij}\neq 0$ e $i+j=n$, lo que demuestra nuestras aserciones.

Si suponemos que los límites inductivos filtrantes existen en $\cat{C}'$ y son exactos (lo que implica, en particular, la existencia de sumas directas infinitas en $\cat{C}'$), entonces existe el complejo definido por el bicomplejo $T(L^{\bullet\bullet})$, y el criterio \sref{0.11.3.3} muestra que la sucesión $'E(T(L^{\bullet\bullet}))$ es siempre \emph{regular}.

Si $K^\bullet$ es un complejo tal que todos los $K^i$ son nulos salvo un único $K^{i_0}$, entonces
$\RR^n T(K^\bullet)$ es isomorfo a $\RR^{n-i_0} T(K^{i_0})$, como se deduce inmediatamente de las definiciones tomando una resolución de Cartan--Eilenberg $L^{\bullet\bullet}$ tal que $L^{ij}=0$ para $i\neq i_0$.

Si $K^\bullet$ y $K'^\bullet$ son dos complejos de $\cat{C}$, y $f$ y $g$ son morfismos homótopos de $K^\bullet$ en $K'^\bullet$, entonces los morfismos $\RR^\bullet T(K^\bullet) \rightarrow \RR^\bullet T(K'^\bullet)$ inducidos por $f$ y $g$ son idénticos, y lo mismo ocurre con los morfismos de las sucesiones espectrales de hipercohomología.
\end{env}

\begin{proposition}[11.4.5]
\label{0.11.4.5}
Supongamos que los límites inductivos filtrantes existen en $\cat{C}'$ y son exactos.
Si $\RR^n T(K^i) = 0$ para todo $n>0$ y todo $i\in\bb{Z}$, entonces se tienen isomorfismos funtoriales
\[
\label{0.11.4.5.1}
  \RR^i T(K^\bullet) \simto \HH^i(T(K^\bullet))
\tag{11.4.5.1}
\]
para todo $i\in\bb{Z}$.
\end{proposition}

\begin{proof}
En efecto, los únicos términos $E_2$ no nulos de la primera sucesión espectral \sref{0.11.4.3.1} son entonces $'E_2^{p0} = \HH^p(T(K^\bullet))$;
en otras palabras, esta sucesión es \emph{degenerada};
como es regular \sref{0.11.4.4}, la conclusión se deduce de \sref{0.11.1.6}.
\end{proof}

\begin{env}[11.4.6]
\label{0.11.4.6}
Consideremos ahora, por ejemplo, un bifuntor covariante $(M, N) \rightarrow T(M, N)$ de $\cat{C} \times \cat{C}'$ en $\cat{C}''$, donde $\cat{C}$, $\cat{C}'$ y $\cat{C}''$ son tres categorías abelianas;
suponemos, por simplicidad, que $T$ es aditivo en cada uno de sus argumentos y, además, que todo objeto de $\cat{C}$ y todo objeto de $\cat{C}'$ son subobjetos de un objeto inyectivo,
y que los límites inductivos filtrantes existen en $\cat{C}''$ y son exactos.
Se define entonces la \emph{hipercohomología} de $T$ con respecto a los complejos $K^\bullet$ y $K'^\bullet$ de $\cat{C}$ y $\cat{C}'$, respectivamente, con operadores diferenciales de grado~$+1$, reemplazando $K^\bullet$ (resp. $K'^\bullet$) por una resolución de Cartan--Eilenberg inyectiva $L^{\bullet\bullet}$ (resp. $L'^{\bullet\bullet}$);
entonces $T(L^{\bullet\bullet}, L'^{\bullet\bullet})$ es un cuadricomplejo de $\cat{C}''$, que consideramos como un \emph{bicomplejo} de $\cat{C}''$, siendo los grados de $T(L^{ij}, L'^{hk})$ los enteros $i+h$ y $j+k$.
La \emph{hipercohomología} de $T$ con respecto a $K^\bullet$ y $K'^\bullet$ es, por definición, la cohomología $\HH^\bullet(T(L^{\bullet\bullet}, L'^{\bullet\bullet}))$ de este bicomplejo (en otras palabras, la del complejo simple asociado) y se denota por $\RR^\bullet T(K^\bullet, K'^\bullet)$;
\oldpage[0\textsubscript{III}]{35}
es el término límite de dos sucesiones espectrales cuyos términos $E_2$ vienen dados por
\[
\label{0.11.4.6.1}
  'E_2^{pq}
  = \HH^p(\RR^q T(K^\bullet, K'^\bullet))
\tag{11.4.6.1}
\]
\[
\label{0.11.4.6.2}
  ''E_2^{pq}
  = \sum_{q'+q''=q}\RR^p T(\HH^{q'}(K^\bullet),\HH^{q''}(K'^\bullet))
\tag{11.4.6.2}
\]
(cf.~M,~XVII,~2).

Aquí $\RR^q T(K^\bullet, K'^\bullet)$ es el bicomplejo $(\RR^q T(K^i, K'^j))_{(i,j)\in \bb{Z}\times \bb{Z}}$, y el segundo miembro de \sref{0.11.4.6.1} es su cohomología cuando se lo considera como complejo simple.

Además, la primera sucesión espectral es siempre regular, y las dos sucesiones espectrales son birregulares cuando existe $n$ tal que todo objeto de $\cat{C}$ y todo objeto de $\cat{C}'$ admiten una resolución inyectiva de longitud $\leq n$, o cuando $K^\bullet$ y $K'^\bullet$ están acotados inferiormente;
en este último caso, puede omitirse además la hipótesis de que existan los límites inductivos en $\cat{C}''$.

Si $K_1^\bullet$ y ${K'_1}^\bullet$ son otros dos complejos de $\cat{C}$ y $\cat{C}'$, respectivamente, $f$ y $g$ son morfismos \emph{homótopos} de $K^\bullet$ en $K_1^\bullet$,
y $f'$ y $g'$ son morfismos \emph{homótopos} de $K'^\bullet$ en ${K'_1}^\bullet$, entonces los morfismos $\RR^\bullet  T(K^\bullet, K'^\bullet) \rightarrow \RR^\bullet  T(K_1^\bullet, {K'_1}^\bullet)$ inducidos por $f$ y $f'$, por una parte, y por $g$ y $g'$, por otra, son idénticos, y lo mismo ocurre con los morfismos de las sucesiones espectrales de hipercohomología.

Se generaliza fácilmente a todo multifuntor covariante aditivo.
\end{env}

\begin{proposition}[11.4.7]
\label{0.11.4.7}
Supongamos que, para todo objeto inyectivo $I$ de $\cat{C}$ (resp. $I'$ de $\cat{C}'$), $A'\mapsto T(I, A')$ (resp. $A\mapsto T(A, I')$) es un funtor exacto.
Entonces, con la notación de \sref{0.11.4.6}, se tiene isomorfismos canónicos
\[
  \label{0.11.4.7.1} \RR^\bullet T(K^\bullet, K'^\bullet) \simto \HH^\bullet(T(L^{\bullet\bullet}, K'^\bullet)) \simto \HH^\bullet(T(K^\bullet, L'^{\bullet\bullet})) \tag{11.4.7.1} 
\]
donde los dos últimos términos son la cohomología de los complejos simples definidos por los tricomplejos $T(L^{\bullet\bullet}, K'^\bullet)$ y $T(K^\bullet, L'^{\bullet\bullet})$, respectivamente.
\end{proposition}

\begin{proof}
Definamos, por ejemplo, el primero de estos isomorfismos.
El cuadricomplejo $T(L^{\bullet\bullet}, L'^{\bullet\bullet})$ puede considerarse como un bicomplejo, donde los grados de $T(L^{ij}, L'^{hk})$ son los enteros $i+j+h$ y $k$.
Como, para cada $h$, ${L'}^{h,\bullet}$ es una \emph{resolución} de ${K'}^h$, se tiene, para este bicomplejo, en virtud de las hipótesis sobre $T$, $\HH_\mathrm{II}^q(T(L^{\bullet\bullet}, L'^{\bullet\bullet}))=0$ para $q\neq 0$, y $\HH_\mathrm{II}^0(T(L^{\bullet\bullet}, L'^{\bullet\bullet})) = T(L^{\bullet\bullet}, K'^\bullet)$;
la primera sucesión espectral de este bicomplejo es, pues, \emph{degenerada};
como $L'^{hk}=0$ para $k<0$, esta sucesión es también regular \sref{0.11.3.3}, y la conclusión se deduce entonces de \sref{0.11.1.6}.
\end{proof}

Se tienen resultados análogos para un multifuntor covariante de un número cualquiera $n$ de argumentos: en el cálculo de la hipercohomología no es necesario reemplazar \emph{todos} los complejos
por una resolución de Cartan--Eilenberg, sino únicamente $n-1$ de ellos, siempre que, al fijar $n-1$ argumentos arbitrarios tomándolos como objetos \emph{inyectivos}, el funtor covariante en el argumento restante sea \emph{exacto}.


\subsection{Paso al límite inductivo en la hipercohomología}
\label{subsection:0.11.5}

\begin{lemma}[11.5.1]
\label{0.11.5.1}
Sea $K^\bullet=(K^i)_{i\in\bb{Z}}$ un complejo de $\cat{C}$ y, para todo entero $r\in\bb{Z}$, sea $K^\bullet_{(r)}$ el complejo tal que $K^i_{(r)}=0$ para $i<r$ y $K^i_{(r)}=K^i$ para $i\geq r$.
Sea $T$ un funtor covariante aditivo
\oldpage[0\textsubscript{III}]{36}
de $\cat{C}$ en $\cat{C}'$ que conmuta con los límites inductivos; suponemos que los límites inductivos filtrantes existen y son exactos en $\cat{C}$ y $\cat{C}'$.
Entonces $\RR^\bullet T(K^\bullet)$ es isomorfo al límite inductivo
\[
  \varinjlim_{r\to-\infty}\RR^\bullet T(K^\bullet_{(r)}).
\]
\end{lemma}

Una resolución de Cartan--Eilenberg inyectiva de $K^\bullet$ se construye eligiendo arbitrariamente, para cada $i$, una resolución inyectiva $(X_{\mathrm{B}}^{ij})_{j\geq 0}$ de $\mathrm{B}^i(K^\bullet)$ y una resolución inyectiva $(X_{\mathrm{H}}^{ij})_{j\geq 0}$ de $\HH^i(K^\bullet)$.
Hecho esto, la construcción muestra que la resolución inyectiva $(L^{ij})_{j\geq 0}$ de $K^i$ y los operadores diferenciales $L^{i,j}\to L^{i+1,j}$ dependen únicamente de las resoluciones $(X_{\mathrm{B}}^{i,\bullet})$, $(X_{\mathrm{H}}^{i,\bullet})$ y $(X_{\mathrm{B}}^{i+1,\bullet})$ (M,~XVII,~1.2).

Está claro ahora que $\mathrm{B}^i(K^\bullet_{(r)})=\mathrm{B}^i(K^\bullet)$ y $\HH^i(K^\bullet_{(r)})=\HH^i(K^\bullet)$ para $i\geq r+1$.
Por otra parte, para todo $r$ existe una inyección canónica
\[
  \varphi^\bullet_{r-1,r}:K^\bullet_{(r)}\longrightarrow K^\bullet_{(r-1)},
\]
donde $\varphi^i_{r-1,r}$ es la identidad para $i\neq r-1$.
La observación anterior muestra que, si $L^{\bullet\bullet}=(L^{ij})$ es una resolución de Cartan--Eilenberg inyectiva de $K^\bullet$, entonces para todo $r$ se puede definir una resolución de Cartan--Eilenberg inyectiva $L^{\bullet\bullet}_{(r)}=(L^{ij}_{(r)})$ de $K^\bullet_{(r)}$ tal que $L^{ij}_{(r)}=0$ para $i<r$ y $L^{ij}_{(r)}=L^{ij}$ para $i\geq r+1$.
También se puede definir un morfismo de bicomplejos
\[
  \Phi^{\bullet\bullet}_{r-1,r}:L^{\bullet\bullet}_{(r)}
  \longrightarrow L^{\bullet\bullet}_{(r-1)}
\]
correspondiente a $\varphi^\bullet_{r-1,r}$; la construcción de este morfismo (\emph{loc.\ cit.}) muestra de nuevo que puede elegirse de modo que $\Phi^{ij}_{r-1,r}$ sea la identidad para $i\neq r$ e $i\neq r-1$.

Hemos definido así un sistema inductivo $(L^{\bullet\bullet}_{(r)})$ de bicomplejos de $\cat{C}$ cuyo límite inductivo, cuando $r$ tiende a $-\infty$, es evidentemente $L^{\bullet\bullet}$.
Como las sumas directas conmutan con los límites inductivos, el complejo simple asociado a $L^{\bullet\bullet}$ es asimismo el límite inductivo de los complejos simples asociados a los $L^{\bullet\bullet}_{(r)}$.
Por hipótesis, $T$ conmuta con los límites inductivos, al igual que la cohomología, porque el funtor $\varinjlim$ es exacto.
Por consiguiente, salvo isomorfismo,
\[
  \HH^\bullet(T(L^{\bullet\bullet}))
  =\varinjlim_{r\to-\infty}\HH^\bullet(T(L^{\bullet\bullet}_{(r)})).
\]

El Lema~\sref{0.11.5.1} permite, mediante el paso al límite inductivo, extender a complejos arbitrarios $K^\bullet$ propiedades que se cumplen para complejos \emph{acotados inferiormente}.
Como primer ejemplo, demostramos lo siguiente.

\begin{proposition}[11.5.2]
\label{0.11.5.2}
Bajo las hipótesis de \sref{0.11.5.1} relativas a $\cat{C}$, $\cat{C}'$ y $T$, el funtor $\RR^\bullet T(K^\bullet)$ es un funtor cohomológico sobre la categoría abeliana de los complejos de $\cat{C}$.
\end{proposition}

\begin{proof}
Mostremos primero que basta tratar los complejos acotados inferiormente.
Dada una sucesión exacta
\[
  0\longrightarrow K'^\bullet\longrightarrow K^\bullet
  \longrightarrow K''^\bullet\longrightarrow 0
\]
de complejos, para todo $r$ induce una sucesión exacta
\[
  0\longrightarrow K'^\bullet_{(r)}\longrightarrow K^\bullet_{(r)}
  \longrightarrow K''^\bullet_{(r)}\longrightarrow 0.
\]
Por hipótesis, esta última da una sucesión exacta
\[
  \cdots\longrightarrow\RR^nT(K'^\bullet_{(r)})
  \longrightarrow\RR^nT(K^\bullet_{(r)})
  \longrightarrow\RR^nT(K''^\bullet_{(r)})
  \xrightarrow{\partial}\RR^{n+1}T(K'^\bullet_{(r)})
  \longrightarrow\cdots .
\]
Estas sucesiones exactas forman un sistema inductivo.
El Lema~\sref{0.11.5.1} y la exactitud de $\varinjlim$ dan, por tanto, una sucesión exacta
\[
  \cdots\longrightarrow\RR^nT(K'^\bullet)
  \longrightarrow\RR^nT(K^\bullet)
  \longrightarrow\RR^nT(K''^\bullet)
  \xrightarrow{\partial}\RR^{n+1}T(K'^\bullet)
  \longrightarrow\cdots .
\]
\end{proof}

Para tratar los complejos acotados inferiormente, basta considerar aquellos para los cuales $K^i=0$ cuando $i<0$; estos forman manifiestamente una categoría abeliana, que denotamos por $\cat{K}$.

\begin{lemma}[11.5.2.1]
\label{0.11.5.2.1}
En $\cat{K}$, sea $\mathfrak{J}$ la clase de los complejos $Q^\bullet=(Q^i)_{i\geq0}$ que poseen las
\oldpage[0\textsubscript{III}]{37}
propiedades siguientes:
\begin{enumerate}
  \item[{\rm(1)}] todo $Q^i$ es un objeto inyectivo de $\cat{C}$;
  \item[{\rm(2)}] para todo $i\geq0$, $\mathrm{Z}^i(Q^\bullet)=\mathrm{B}^i(Q^\bullet)$, y $\mathrm{Z}^i(Q^\bullet)$ es un sumando directo de $Q^i$.
\end{enumerate}
Entonces:
\begin{enumerate}
  \item[{\rm(i)}] todo $Q^\bullet\in\mathfrak{J}$ es un objeto inyectivo de $\cat{K}$;
  \item[{\rm(ii)}] todo objeto de $\cat{K}$ es isomorfo a un subcomplejo de un complejo perteneciente a $\mathfrak{J}$.
\end{enumerate}
\end{lemma}

\begin{proof}
{\rm(i)}
Sea $A^\bullet=(A^i)$ un objeto de $\cat{K}$, sea $A'^\bullet=(A'^i)$ un subobjeto de $A^\bullet$, sea $Q^\bullet=(Q^i)$ perteneciente a $\mathfrak{J}$, y supongamos dado un morfismo
\[
  f=(f^i):A'^\bullet\longrightarrow Q^\bullet
\]
.
Debemos extenderlo a un morfismo $g=(g^i):A^\bullet\to Q^\bullet$.
Por simplicidad, utilizamos el lenguaje de la categoría de módulos (cf.~[27]).

Identifiquemos $Q^i$ con $\mathrm{B}^i(Q^\bullet)\oplus\mathrm{B}^{i+1}(Q^\bullet)$.
Procedemos por inducción sobre $i$.
Supongamos que $g^j$ ha sido definido para $j<i$, compatiblemente con los diferenciales $d^j:A^j\to A^{j+1}$ y $d'^j:Q^j\to Q^{j+1}$ para $j<i-1$, y que además:
\begin{enumerate}
  \item[{\rm(1)}] $g^{i-1}(\mathrm{Z}^{i-1}(A^\bullet))\subset\mathrm{Z}^{i-1}(Q^\bullet)$;
  \item[{\rm(2)}] si $C^j=(d^j)^{-1}(A'^{j+1})$ para todo $j$, entonces $d'^{i-1}\circ g^{i-1}$ coincide con $f^i\circ d^{i-1}$ sobre $C^{i-1}$.
\end{enumerate}
Componiendo $f^i:A'^i\to Q^i$ con las dos proyecciones, obtenemos morfismos
\[
  f'^i:A'^i\longrightarrow\mathrm{B}^i(Q^\bullet),
  \qquad
  f''^i:A'^i\longrightarrow\mathrm{B}^{i+1}(Q^\bullet).
\]
Como $d'^{i-1}\circ g^{i-1}$ envía $A^{i-1}$ en $\mathrm{B}^i(Q^\bullet)$ y se anula sobre $\mathrm{Z}^{i-1}(A^\bullet)$, define un morfismo
\[
  h^i:\mathrm{B}^i(A^\bullet)\longrightarrow\mathrm{B}^i(Q^\bullet).
\]
La condición~{\rm(2)} dice que $h^i$ coincide con $f'^i$ sobre $\mathrm{B}^i(A^\bullet)\cap A'^i$.
Como $\mathrm{B}^i(Q^\bullet)$ es un sumando directo de $Q^i$, es inyectivo; existe, por tanto, un morfismo
\[
  g'^i:A^i\longrightarrow\mathrm{B}^i(Q^\bullet)
\]
que coincide con $h^i$ sobre $\mathrm{B}^i(A^\bullet)$ y con $f'^i$ sobre $A'^i$.

Por otra parte, consideremos el morfismo
\[
  f'^{i+1}\circ d^i:C^i\longrightarrow\mathrm{B}^{i+1}(Q^\bullet),
\]
que se anula sobre $\mathrm{Z}^i(A^\bullet)$.
Como $\mathrm{B}^{i+1}(Q^\bullet)$ es inyectivo, existe un morfismo
\[
  g''^i:A^i\longrightarrow\mathrm{B}^{i+1}(Q^\bullet)
\]
que coincide con $f'^{i+1}\circ d^i$ sobre $C^i$ y con cero sobre $\mathrm{Z}^i(A^\bullet)$.
Tomando $g^i=g'^i+g''^i$ puede continuarse la inducción.

\smallskip
\noindent{\rm(ii)}
Para sumergir $A^\bullet=(A^i)$ en un complejo perteneciente a $\mathfrak{J}$, elijamos para todo $i\geq1$ un objeto inyectivo $Q'^i$ de $\cat{C}$ y una inyección $f'^i:A^i\to Q'^i$.
Pongamos
\[
  Q^i=0\quad(i<0),\qquad Q^0=Q'^1,\qquad
  Q^i=Q'^i\oplus Q'^{i+1}\quad(i\geq1),
\]
con el diferencial evidente.
Para todo $i$, pongamos
\[
  f^i=f'^i+(f'^{i+1}\circ d^i),
\]
donde $f'^i=0$ para $i\leq0$.
Es inmediato que todo $f^i$ es inyectivo y que los $f^i$ conmutan con los diferenciales.
\end{proof}

\begin{corollary}[11.5.2.2]
\label{0.11.5.2.2}
Todo objeto $K^\bullet$ de $\cat{K}$ admite una resolución por la derecha formada por objetos de $\mathfrak{J}$.
Si $L^{\bullet\bullet}$ es una tal resolución, entonces para toda resolución $L'^{\bullet\bullet}$ de $K^\bullet$ formada por objetos de $\cat{K}$ existe un morfismo de bicomplejos
\[
  F:L'^{\bullet\bullet}\longrightarrow L^{\bullet\bullet}
\]
compatible con las aumentaciones, y dos morfismos cualesquiera $F,F'$ de este tipo son homótopos.
\end{corollary}

\begin{proof}
Esto es precisamente (M, V, 1.1~a)) aplicado a la categoría abeliana $\cat{K}$.
\end{proof}

\begin{env}[11.5.2.3]
\label{0.11.5.2.3}
Establecidos estos preliminares, sea $L'^{\bullet\bullet}$ una resolución de Cartan--Eilenberg inyectiva de $K^\bullet$ y sea $L^{\bullet\bullet}$ una resolución de $K^\bullet$ formada por objetos de $\mathfrak{J}$.
Mostremos que existe un isomorfismo
\[
  \HH^\bullet(T(L'^{\bullet\bullet}))
  \simto
  \HH^\bullet(T(L^{\bullet\bullet})).
\]
El Corolario~\sref{0.11.5.2.2} da un morfismo de bicomplejos
\[
  T(L'^{\bullet\bullet})\longrightarrow T(L^{\bullet\bullet}),
\]
y, por tanto, un morfismo
\[
  'E(T(L'^{\bullet\bullet}))\longrightarrow{}'E(T(L^{\bullet\bullet}))
\]
entre sus primeras sucesiones espectrales.
Por \sref{0.11.3.3}, estas sucesiones espectrales son regulares, de modo que por \sref{0.11.1.5} basta comprobar que este morfismo es un isomorfismo sobre los términos $E_2$.
Equivalentemente, hay que demostrar que
\[
  \HH_\mathrm{II}^q(T(L^{i,\bullet}))=\RR^qT(K^i).
\]
Como $L^{i,\bullet}$ es una resolución por la derecha de $K^i$, esto reduce la cuestión al lema siguiente.
\end{env}

\oldpage[0\textsubscript{III}]{38}
\begin{lemma}[11.5.2.4]
\label{0.11.5.2.4}
Sea $L^\bullet=(L^i)_{i\in\bb{Z}}$ una resolución por la derecha de un objeto $A$ de $\cat{C}$ tal que $\RR^nT(L^i)=0$ para todo $i\in\bb{Z}$ y todo $n>0$.
Entonces
\[
  \HH^\bullet(T(L^\bullet))=\RR^\bullet T(A).
\]
\end{lemma}

\begin{proof}
Este es un caso particular de (T, 2.5.2).
\end{proof}

\begin{env}[11.5.2.5]
\label{0.11.5.2.5}
La demostración de \sref{0.11.5.2} se deduce ahora inmediatamente, puesto que $L'^{\bullet\bullet}$ es una resolución inyectiva de $K^\bullet$ en la categoría abeliana $\cat{K}$.
En otras palabras, el funtor
\[
  K^\bullet\longmapsto\HH^\bullet(T(L'^{\bullet\bullet}))
\]
es precisamente el funtor cohomológico derivado por la derecha de $T$ en la categoría $\cat{K}$ (T, 2.3).
\end{env}

\begin{proposition}[11.5.3]
\label{0.11.5.3}
Bajo las hipótesis de \sref{0.11.5.1} relativas a $\cat{C}$, $\cat{C}'$ y $T$, sea $L^{\bullet\bullet}=(L^{ij})$ un bicomplejo tal que $L^{ij}=0$ para $j<0$ y tal que, para todo $i$, $L^{i,\bullet}$ es una resolución de $K^i$.
Supongamos además que $\RR^nT(L^{ij})=0$ para todo par $(i,j)$ y todo $n>0$.
Entonces existe un isomorfismo funtorial
\[
\label{0.11.5.3.1}
\tag{11.5.3.1}
  \RR^\bullet T(K^\bullet)
  \simto
  \HH^\bullet(T(L^{\bullet\bullet})).
\]
\end{proposition}

\begin{proof}
Sea $L^{\bullet\bullet}_{(r)}=(L^{ij}_{(r)})$ el bicomplejo definido por
\[
  L^{ij}_{(r)}=0\quad(i<r),\qquad
  L^{ij}_{(r)}=L^{ij}\quad(i\geq r).
\]
Es inmediato que $L^{\bullet\bullet}$ es el límite inductivo de los $L^{\bullet\bullet}_{(r)}$ cuando $r$ tiende a $-\infty$.
Por la hipótesis sobre $T$ y por \sref{0.11.5.1}, basta, por tanto, demostrar la proposición cuando $K^\bullet$ está acotado inferiormente, por ejemplo cuando $K^i=0$ para $i<0$, con $L^{ij}=0$ para $i<0$.
Sea $L'^{\bullet\bullet}=(L'^{ij})$ una resolución por la derecha de $K^\bullet$ formada por objetos de $\mathfrak{J}$, como en \sref{0.11.5.2.2}.
Existe un morfismo de bicomplejos
\[
  L^{\bullet\bullet}\longrightarrow L'^{\bullet\bullet}
\]
compatible con las aumentaciones y, por tanto, un morfismo
\[
  'E(T(L^{\bullet\bullet}))\longrightarrow{}'E(T(L'^{\bullet\bullet}))
\]
entre las primeras sucesiones espectrales.
El Lema~\sref{0.11.5.2.4} muestra, exactamente como en \sref{0.11.5.2.3}, que este morfismo es un isomorfismo, lo que demuestra la aserción.
\end{proof}

\begin{remark}[11.5.4]
\label{0.11.5.4}
Los razonamientos anteriores muestran que las conclusiones de \sref{0.11.5.2} y \sref{0.11.5.3} se cumplen en la categoría de los complejos $K^\bullet$ acotados inferiormente sin suponer que $T$ conmuta con los límites inductivos filtrantes.
Además, si se considera únicamente la categoría $\cat{K}$ de los complejos $K^\bullet$ tales que $K^i=0$ para $i<0$, la caracterización de $\RR^\bullet T(K^\bullet)$ como el sistema de funtores derivados por la derecha de $T$ en $\cat{K}$ muestra que este funtor cohomológico es \emph{universal} (T, 2.3).

Otro caso en el que \sref{0.11.5.2} se cumple sin hipótesis adicional alguna sobre $T$ es aquel en que existe un entero $m>0$ tal que todo objeto de $\cat{C}$ admite una resolución inyectiva de longitud a lo sumo $m$.
En efecto, en la demostración de \sref{0.11.5.1}, todas las resoluciones inyectivas de objetos de $\cat{C}$ que aparecen pueden elegirse de longitud a lo sumo $m$.
Se deduce inmediatamente que, para un grado total fijo $n$ y para $-r$ suficientemente grande, los términos de grado total $n$ en $T(L^{\bullet\bullet}_{(r)})$ son iguales a los de $T(L^{\bullet\bullet})$ y son finitos en número.
Por consiguiente,
\[
  \HH^n(T(L^{\bullet\bullet}))
  =\HH^n(T(L^{\bullet\bullet}_{(r)}))
\]
para $-r$ suficientemente grande.
Con la notación de \sref{0.11.5.2}, se tiene análogamente
\[
  \RR^nT(K^\bullet_{(r)})=\RR^nT(K^\bullet)
\]
para $-r$ suficientemente grande (dependiendo de $n$), y análogamente para $K'^\bullet$ y $K''^\bullet$, lo que demuestra la conclusión.
Del mismo modo, \sref{0.11.5.3} se cumple sin condición adicional sobre $T$ cuando $\cat{C}$ satisface la hipótesis anterior y se toman las resoluciones $L^{i,\bullet}$ de longitud a lo sumo $m$.
\end{remark}

\begin{env}[11.5.5]
\label{0.11.5.5}
El resultado de \sref{0.11.5.2} se extiende a los multifuntores covariantes.
Consideremos, por ejemplo, la situación de \sref{0.11.4.6}, en la que se supone que los límites inductivos filtrantes existen y son exactos en $\cat{C}$, $\cat{C}'$ y $\cat{C}''$,
\oldpage[0\textsubscript{III}]{39}
y se supone que $T$ conmuta con los límites inductivos.
Entonces $\RR^\bullet T(K^\bullet,K'^\bullet)$ es un bifuntor cohomológico en cada uno de los complejos $K^\bullet,K'^\bullet$.
Para verlo, se reduce, como en \sref{0.11.5.2}, al caso en que $K^\bullet$ y $K'^\bullet$ están acotados inferiormente.
Eligiendo resoluciones inyectivas de $K^\bullet$ y $K'^\bullet$ del tipo descrito en \sref{0.11.5.2.2}, se queda reducido a la propiedad general demostrada en (M, V, 4.1).
\end{env}

\begin{env}[11.5.6]
\label{0.11.5.6}
Análogamente, los resultados de \sref{0.11.4.7} y \sref{0.11.5.3} se extienden como sigue.
Bajo las hipótesis de \sref{0.11.5.5}, supongamos dados dos bicomplejos
\[
  L^{\bullet\bullet}=(L^{ij}),\qquad
  L'^{\bullet\bullet}=(L'^{ij})
\]
tales que $L^{ij}=L'^{ij}=0$ para $j<0$, tales que, para todo $i$, $L^{i,\bullet}$ es una resolución de $K^i$ y $L'^{i,\bullet}$ es una resolución de $K'^i$, y tales que
\[
  \RR^nT(L^{ij},L'^{hk})=0
\]
para $n>0$ y todo sistema de índices $(i,j,h,k)$.
Entonces existe un isomorfismo, funtorial en $K^\bullet$ y $K'^\bullet$,
\[
\label{0.11.5.6.1}
\tag{11.5.6.1}
  \RR^\bullet T(K^\bullet,K'^\bullet)
  \simto
  \HH^\bullet(T(L^{\bullet\bullet},L'^{\bullet\bullet})).
\]
Esto se demuestra como en \sref{0.11.5.3}, reduciéndose al caso en que $K^\bullet$ y $K'^\bullet$ están acotados inferiormente.

Supongamos además que, para todo par $(i,j)$ y todo par $(h,k)$, los funtores
\[
  A\longmapsto T(A,L'^{hk}),
  \qquad
  A'\longmapsto T(L^{ij},A')
\]
son exactos en $\cat{C}$ y $\cat{C}'$, respectivamente.
Entonces existen también isomorfismos funtoriales
\[
\label{0.11.5.6.2}
\tag{11.5.6.2}
  \RR^\bullet T(K^\bullet,K'^\bullet)
  \simto
  \HH^\bullet(T(L^{\bullet\bullet},K'^\bullet))
  \simto
  \HH^\bullet(T(K^\bullet,L'^{\bullet\bullet})).
\]
La demostración es análoga a la de \sref{0.11.4.7}.
\end{env}

\begin{env}[11.5.7]
\label{0.11.5.7}
Los resultados de \sref{0.11.5.5} y \sref{0.11.5.6} se cumplen también sin suponer que $T$ conmuta con los límites inductivos, con tal de limitarse a complejos $K^\bullet,K'^\bullet$ acotados inferiormente, o de suponer que todo objeto de $\cat{C}$ (respectivamente, de $\cat{C}'$) admite una resolución inyectiva de longitud acotada y, en \sref{0.11.5.6}, que los segundos grados de $L^{\bullet\bullet}$ y $L'^{\bullet\bullet}$ están acotados superiormente.
\end{env}

\subsection{Hiperhomología de un funtor con respecto a un complejo $K_{\bullet}$}
\label{subsection:0.11.6}

\begin{env}[11.6.1]
\label{0.11.6.1}
Sea $\cat{C}$ una categoría abeliana y sea $K_\bullet=(K_i)_{i\in\bb{Z}}$ un complejo de objetos de $\cat{C}$ cuyo diferencial es de grado $-1$.
Una \emph{resolución de Cartan--Eilenberg por la izquierda} de $K_\bullet$ consiste en un bicomplejo $L_{\bullet\bullet}=(L_{ij})$, cuyos diferenciales son de grado $-1$ y para el cual $L_{ij}=0$ cuando $j<0$, junto con un morfismo de complejos simples
\[
  \varepsilon:L_{\bullet,0}\longrightarrow K_\bullet,
\]
tal que se satisfacen las condiciones obtenidas de las de \sref{0.11.4.2} ``invirtiendo las flechas''.
Si todo objeto de $\cat{C}$ es cociente de un objeto proyectivo, entonces todo complejo $K_\bullet$ de $\cat{C}$ admite una \emph{resolución de Cartan--Eilenberg proyectiva}, es decir, formada por objetos proyectivos $L_{ij}$; posee propiedades funtoriales análogas a las de \sref{0.11.4.2}.
Además, si $K_\bullet$ está acotado inferiormente (respectivamente, superiormente), puede exigirse $L_{ij}=0$ para $i<i_0$ (respectivamente, $i>i_0$) siempre que $K_i=0$ para $i<i_0$ (respectivamente, $i>i_0$).
Si todo objeto de $\cat{C}$ admite una resolución proyectiva de longitud a lo sumo $n$, puede exigirse $L_{ij}=0$ para $j>n$.
\end{env}

\begin{env}[11.6.2]
\label{0.11.6.2}
\oldpage[0\textsubscript{III}]{40}
Supongamos que $T$ es un funtor covariante aditivo de $\cat{C}$ en una categoría abeliana $\cat{C}'$.
La \emph{hiperhomología} $\LL_\bullet T(K_\bullet)$ y las \emph{sucesiones espectrales de hiperhomología} de $T$ con respecto a un complejo $K_\bullet$ de $\cat{C}$ (cuando existen) se definen a partir de \sref{0.11.4.3}, de nuevo ``invirtiendo las flechas''.
Los términos $E^2$ de las dos sucesiones espectrales así obtenidas son
\[
\label{0.11.6.2.1}
\tag{11.6.2.1}
  {}'E^2_{pq}
  =\HH_p(\LL_qT(K_\bullet))
\]
y
\[
\label{0.11.6.2.2}
\tag{11.6.2.2}
  {}''E^2_{pq}
  =\LL_pT(\HH_q(K_\bullet)),
\]
donde $\LL_pT$ denota el funtor derivado por la izquierda de grado $p$ de $T$ para $p\geq0$, y cero para $p<0$, mientras que $\LL_qT(K_\bullet)$ denota el complejo $(\LL_qT(K_i))_{i\in\bb{Z}}$.

Las propiedades de la hiperhomología no se deducen todas de las de la hipercohomología con solo ``invertir las flechas'' (a menos que se impongan a $\cat{C}'$ hipótesis adicionales de tipo (AB~5*) de (T, 1.5)).
La razón son las condiciones de regularidad de las dos sucesiones espectrales anteriores, a las que ahora hay que aplicar los criterios de \sref{0.11.3.4}.
Estos criterios muestran que, si los límites inductivos filtrantes existen y son exactos en $\cat{C}'$, entonces existe el complejo definido por el bicomplejo $T(L_{\bullet\bullet})$ y la \emph{segunda} sucesión espectral ${}''E(T(L_{\bullet\bullet}))$ es regular.
Si $K_\bullet$ está acotado inferiormente o existe un entero $n$ tal que todo objeto de $\cat{C}$ admite una resolución proyectiva de longitud a lo sumo $n$, entonces las \emph{dos} sucesiones espectrales de hiperhomología existen (sin hipótesis alguna sobre $\cat{C}'$) y son birregulares.
\end{env}

\begin{proposition}[11.6.3]
\label{0.11.6.3}
Sean $\cat{C}$ y $\cat{C}'$ dos categorías abelianas, y sea $T$ un funtor covariante aditivo de $\cat{C}$ en $\cat{C}'$.
Entonces:
\begin{enumerate}
  \item[{\rm(i)}] \emph{La hiperhomología $\LL_\bullet T(K_\bullet)$ es un funtor homológico sobre la categoría abeliana de los complejos de $\cat{C}$ acotados inferiormente.}

  \item[{\rm(ii)}] \emph{Sea $K_\bullet$ un complejo de $\cat{C}$ acotado inferiormente.
  Si $\LL_nT(K_i)=0$ para todo $n>0$ y todo $i\in\bb{Z}$, existen isomorfismos funtoriales}
  \[
  \label{0.11.6.3.1}
  \tag{11.6.3.1}
    \LL_iT(K_\bullet)
    \simto
    \HH_i(T(K_\bullet))
    \qquad(i\in\bb{Z}).
  \]

  \item[{\rm(iii)}] \emph{Sea $K_\bullet$ un complejo de $\cat{C}$ acotado inferiormente.
  Sea $L_{\bullet\bullet}=(L_{ij})$ un bicomplejo tal que $L_{ij}=0$ para $j<0$ y, para todo $i$, $L_{i,\bullet}$ es una resolución de $K_i$.
  Supongamos además que $\LL_nT(L_{ij})=0$ para todo par $(i,j)$ y todo $n>0$.
  Entonces existe un isomorfismo funtorial}
  \[
  \label{0.11.6.3.2}
  \tag{11.6.3.2}
    \LL_\bullet T(K_\bullet)
    \simto
    \HH_\bullet(T(L_{\bullet\bullet})).
  \]
\end{enumerate}

Las demostraciones proceden como las de \sref{0.11.5.2}, \sref{0.11.4.5} y \sref{0.11.5.3} en el caso de complejos acotados inferiormente.
Dejamos al lector los detalles de estos razonamientos.
\end{proposition}

\begin{env}[11.6.4]
\label{0.11.6.4}
Se tienen resultados enteramente análogos para los multifuntores covariantes aditivos en cada argumento.
Por ejemplo, para un bifuntor $T$, las dos sucesiones espectrales de hiperhomología tienen términos $E^2$
\oldpage[0\textsubscript{III}]{41}
\[
\label{0.11.6.4.1}
\tag{11.6.4.1}
  {}'E^2_{pq}
  =\HH_p(\LL_qT(K_\bullet,K'_\bullet))
\]
y
\[
\label{0.11.6.4.2}
\tag{11.6.4.2}
  {}''E^2_{pq}
  =
  \sum_{q'+q''=q}
  \LL_pT(\HH_{q'}(K_\bullet),\HH_{q''}(K'_\bullet)).
\]
También aquí es la \emph{segunda} sucesión espectral la que es regular.
Las dos sucesiones son birregulares cuando $K_\bullet$ y $K'_\bullet$ están acotados inferiormente, o cuando los objetos de las categorías abelianas consideradas admiten resoluciones proyectivas de una longitud fija.

Además:
\end{env}

\begin{proposition}[11.6.5]
\label{0.11.6.5}
Sean $\cat{C}$, $\cat{C}'$ y $\cat{C}''$ tres categorías abelianas, y sea $T$ un bifuntor covariante biaditivo de $\cat{C}\times\cat{C}'$ en $\cat{C}''$.
\begin{enumerate}
  \item[{\rm(i)}] \emph{$\LL_\bullet T(K_\bullet,K'_\bullet)$ es un bifuntor homológico en cada uno de los complejos acotados inferiormente $K_\bullet,K'_\bullet$, formados respectivamente por objetos de $\cat{C}$ y de $\cat{C}'$.}

  \item[{\rm(ii)}] \emph{Supongamos que $K_\bullet$ y $K'_\bullet$ están acotados inferiormente.
  Sean $L_{\bullet\bullet}=(L_{ij})$ y $L'_{\bullet\bullet}=(L'_{ij})$ dos bicomplejos tales que $L_{ij}=L'_{ij}=0$ para $j<0$ y tales que, para todo $i$, $L_{i,\bullet}$ es una resolución de $K_i$ y $L'_{i,\bullet}$ es una resolución de $K'_i$.
  Supongamos finalmente que $\LL_nT(L_{ij},L'_{hk})=0$ para todo $n>0$ y todo sistema de índices $(i,j,h,k)$.
  Entonces existe un isomorfismo funtorial}
  \[
  \label{0.11.6.5.1}
  \tag{11.6.5.1}
    \LL_\bullet T(K_\bullet,K'_\bullet)
    \simto
    \HH_\bullet(T(L_{\bullet\bullet},L'_{\bullet\bullet})).
  \]

  \item[{\rm(iii)}] \emph{Supongamos además que, para todo par $(i,j)$ y todo par $(h,k)$, los funtores
  \[
    A\longmapsto T(A,L'_{hk}),
    \qquad
    A'\longmapsto T(L_{ij},A')
  \]
  son exactos en $\cat{C}$ y $\cat{C}'$, respectivamente.
  Entonces existen isomorfismos funtoriales}
  \[
  \label{0.11.6.5.2}
  \tag{11.6.5.2}
    \LL_\bullet T(K_\bullet,K'_\bullet)
    \simto
    \HH_\bullet(T(L_{\bullet\bullet},K'_\bullet))
    \simto
    \HH_\bullet(T(K_\bullet,L'_{\bullet\bullet})).
  \]
\end{enumerate}

Las demostraciones son análogas a las de \sref{0.11.5.5} y \sref{0.11.5.6}.
\end{proposition}

\subsection{Hiperhomología de un funtor con respecto a un bicomplejo $K_{\bullet\bullet}$}
\label{subsection:0.11.7}

\begin{env}[11.7.1]
\label{0.11.7.1}
Sea $\cat{C}$ una categoría abeliana en la que todo objeto es cociente de un
objeto proyectivo.
Consideremos un bicomplejo $K_{\bullet\bullet}=(K_{ij})$ de objetos de $\cat{C}$ cuyos
dos grados están acotados inferiormente.
Podemos limitarnos siempre al caso en que $K_{ij}=0$ para $i<0$ o
$j<0$, y así lo haremos en lo sucesivo.
Podemos considerar $K_{\bullet\bullet}$ como un complejo ordinario formado por los
objetos
\[
  K_{i,\bullet}=(K_{ij})_{j\geq 0}
\]
de la categoría abeliana $\cat{K}$ de los complejos no negativos de objetos de
$\cat{C}$.
Del Lema~\sref{0.11.5.2.1} (o, más precisamente, del lema
``dual'' obtenido ``invirtiendo las flechas'') y de (M,~V,~2.2) se deduce
que $K_{\bullet\bullet}$ admite una \emph{resolución de Cartan--Eilenberg
proyectiva} en la categoría $\cat{K}$.
Una tal resolución es un tricomplejo
$M_{\bullet\bullet\bullet}=(M_{ijk})$ en $\cat{C}$, todos cuyos grados son
no negativos y cuyos objetos son proyectivos, tal que, para todo $i$,
\[
  M_{i,\bullet\bullet},\qquad
  \mathrm{B}^{\mathrm{I}}_i(M_{\bullet\bullet\bullet}),\qquad
  \mathrm{Z}^{\mathrm{I}}_i(M_{\bullet\bullet\bullet}),\qquad
  \HH^{\mathrm{I}}_i(M_{\bullet\bullet\bullet})
\]
son resoluciones proyectivas en $\cat{K}$, respectivamente, de
\[
  K_{i,\bullet},\qquad
  \mathrm{B}^{\mathrm{I}}_i(K_{\bullet\bullet}),\qquad
  \mathrm{Z}^{\mathrm{I}}_i(K_{\bullet\bullet}),\qquad
  \HH^{\mathrm{I}}_i(K_{\bullet\bullet}).
\]
En particular, para todo par $(i,j)$, $M_{ij,\bullet}$ es una resolución
proyectiva de $K_{ij}$ en $\cat{C}$.
\end{env}

\begin{proposition}[11.7.2]
\label{0.11.7.2}
Sea $T$ un funtor covariante aditivo de $\cat{C}$ en una categoría
abeliana $\cat{C}'$.
Con la notación de \sref{0.11.7.1}, la homología
\[
  \HH_\bullet\bigl(T(M_{\bullet\bullet\bullet})\bigr)
\]
del complejo simple asociado al tricomplejo
$T(M_{\bullet\bullet\bullet})$ es canónicamente isomorfa a la
hiperhomología $\LL_\bullet T(K_{\bullet\bullet})$ del
\oldpage[0\textsubscript{III}]{42}
complejo simple asociado a $K_{\bullet\bullet}$
\sref{0.11.6.2}.
Es también el término límite común de seis sucesiones espectrales birregulares denotadas
${}^{(t)}\mathrm{E}$, donde $t=a,b,a',b',c$ o $d$, cuyos
términos $\mathrm{E}^2$ son
\[
\begin{aligned}
  {}^{(a)}\mathrm{E}^2_{pq}
    &= \LL_pT\bigl(\HH_q(K_{\bullet\bullet})\bigr),\\
  {}^{(b)}\mathrm{E}^2_{pq}
    &= \HH_p\bigl(\LL^{\mathrm{II}}_qT(K_{\bullet\bullet})\bigr),\\
  {}^{(a')}\mathrm{E}^2_{pq}
    &= \LL_pT\bigl(\HH^{\mathrm{I}}_q(K_{\bullet\bullet})\bigr),\\
  {}^{(b')}\mathrm{E}^2_{pq}
    &= \HH_p\bigl(\LL_qT(K_{\bullet\bullet})\bigr),\\
  {}^{(c)}\mathrm{E}^2_{pq}
    &= \LL_pT\bigl(\HH^{\mathrm{II}}_q(K_{\bullet\bullet})\bigr),\\
  {}^{(d)}\mathrm{E}^2_{pq}
    &= \HH_p\bigl(\LL^{\mathrm{I}}_qT(K_{\bullet\bullet})\bigr).
\end{aligned}
\]
\end{proposition}

Recordemos que, para un complejo $A_\bullet=(A_i)$, utilizamos la notación
$F(A_\bullet)$ para el complejo cuyos objetos son los $F(A_i)$.
Por ejemplo, $\LL^{\mathrm{II}}_qT(K_{\bullet\bullet})$ denota el complejo
\[
  \bigl(\LL^{\mathrm{II}}_qT(K_{i,\bullet})\bigr)_{i\geq 0},
\]
donde $\LL^{\mathrm{II}}_qT(K_{i,\bullet})$ es la hiperhomología de índice
$q$ del funtor $T$ relativa al complejo simple $K_{i,\bullet}$.

\begin{proof}
Sea $L_\bullet$ el complejo simple asociado a
$K_{\bullet\bullet}$, de modo que
\[
  L_i=\bigoplus_{r+s=i}K_{rs},
\]
y pongamos
\[
  N_{ij}=\bigoplus_{r+s=i}M_{rsj}.
\]
Está claro que, para todo $i$, $N_{i,\bullet}$ es una resolución proyectiva
de $L_i$ en $\cat{C}$.
De \sref{0.11.6.3} y \sref{0.11.6.4} se deduce que existe un
isomorfismo funtorial
\[
  \LL_\bullet T(L_\bullet)
  \isoto
  \HH_\bullet\bigl(T(N_{\bullet\bullet})\bigr).
\]
Como el complejo simple asociado al bicomplejo
$T(N_{\bullet\bullet})$ es también el complejo simple asociado al
tricomplejo $T(M_{\bullet\bullet\bullet})$, esto demuestra la primera aserción.

Además, $\LL_\bullet T(L_\bullet)$ es el término límite de las dos
sucesiones espectrales de hiperhomología \sref{0.11.6.2} de $T$ relativas al
complejo simple $L_\bullet$.
Estas son precisamente ${}^{(b')}\mathrm{E}$ y ${}^{(a)}\mathrm{E}$,
respectivamente.

Consideremos ahora $M_{\bullet\bullet\bullet}$ como un bicomplejo
$U_{\bullet\bullet}$, donde
\[
  U_{ij}=\bigoplus_{r+s=j}M_{irs}.
\]
La homología $\HH_\bullet(T(M_{\bullet\bullet\bullet}))$ es también la
término límite de las dos sucesiones espectrales del bicomplejo
$T(U_{\bullet\bullet})$.
Para todo $i$, el bicomplejo $M_{i,\bullet\bullet}$ satisface las
hipótesis de \sref{0.11.6.3} relativas al complejo simple
$K_{i,\bullet}$.
Por consiguiente,
\[
  \HH^{\mathrm{II}}_q\bigl(T(U_{i,\bullet})\bigr)
  =
  \LL^{\mathrm{II}}_qT(K_{i,\bullet}),
\]
y la primera sucesión espectral de $T(U_{\bullet\bullet})$ es precisamente
${}^{(b)}\mathrm{E}$.

Por otra parte, para todo $r$, $M_{\bullet,r,\bullet}$ es una
resolución de Cartan--Eilenberg del complejo simple $K_{\bullet,r}$.
El cálculo efectuado en (M,~XV,~2) muestra que
\[
  \HH^{\mathrm{I}}_q\bigl(T(M_{\bullet,rs})\bigr)
  =
  T\bigl(\HH^{\mathrm{I}}_q(M_{\bullet,rs})\bigr),
\]
y, por tanto,
\[
  \HH^{\mathrm{I}}_q\bigl(T(U_{\bullet,j})\bigr)
  =
  \bigoplus_{r+s=j}
  T\bigl(\HH^{\mathrm{I}}_q(M_{\bullet,rs})\bigr).
\]
En otras palabras, el complejo simple
$\HH^{\mathrm{I}}_q(T(U_{\bullet\bullet}))$ es el complejo simple
asociado al bicomplejo
$T(\HH^{\mathrm{I}}_q(M_{\bullet\bullet\bullet}))$.
Para todo $q$, $\HH^{\mathrm{I}}_q(M_{\bullet\bullet\bullet})$ es una
resolución proyectiva en $\cat{K}$ del complejo simple
$\HH^{\mathrm{I}}_q(K_{\bullet\bullet})$.
Aplicando \sref{0.11.6.3}, obtenemos
\[
  \HH^{\mathrm{II}}_p
  \bigl(T(\HH^{\mathrm{I}}_q(M_{\bullet\bullet\bullet}))\bigr)
  =
  \LL_pT\bigl(\HH^{\mathrm{I}}_q(K_{\bullet\bullet})\bigr),
\]
y, por tanto, la sucesión espectral ${}^{(a')}\mathrm{E}$.
Finalmente, las sucesiones espectrales ${}^{(c)}\mathrm{E}$ y
${}^{(d)}\mathrm{E}$ se obtienen intercambiando los papeles de los
índices $i$ y $j$ en la definición del tricomplejo
$M_{\bullet\bullet\bullet}$ y aplicando los razonamientos anteriores al
tricomplejo resultante.
\end{proof}

Llamamos a $\LL_\bullet T(K_{\bullet\bullet})$ la \emph{hiperhomología} de
$T$ relativa al bicomplejo $K_{\bullet\bullet}$.

\begin{env}[11.7.3]
\label{0.11.7.3}
\emph{Observaciones}.
\begin{enumerate}
  \item[{\rm(i)}]
  De \sref{0.11.6.3} se deduce que
  $\LL_\bullet T(K_{\bullet\bullet})$ es un funtor homológico sobre la
  categoría de los bicomplejos $K_{\bullet\bullet}$ en $\cat{C}$ que están
  acotados inferiormente en cada grado.

\oldpage[0\textsubscript{III}]{43}
  \item[{\rm(ii)}]
  Sea $M_{\bullet\bullet\bullet}$ un tricomplejo en $\cat{C}$ tal que,
  para todo par $(r,s)$, $M_{rs,\bullet}$ sea una resolución de $K_{rs}$,
  y tal que
  \[
    \LL_nT(M_{ijk})=0
  \]
  para toda terna $(i,j,k)$ y todo $n>0$.
  Entonces existe un isomorfismo
  \[
    \LL_\bullet T(K_{\bullet\bullet})
    \isoto
    \HH_\bullet\bigl(T(M_{\bullet\bullet\bullet})\bigr).
  \]
  En efecto, con la notación utilizada en la demostración de \sref{0.11.7.2},
  $N_{i,\bullet}$ es una resolución de $L_i$ tal que
  $\LL_nT(N_{ij})=0$ para todo $n>0$ y todo par $(i,j)$, por lo que
  basta aplicar \sref{0.11.6.3},~(iii).

  \item[{\rm(iii)}]
  Los resultados de \sref{0.11.7.2} se generalizan inmediatamente a los
  multifuntores covariantes.
  Por ejemplo, sea $\cat{C}'$ una segunda categoría abeliana en la que todo
  objeto sea cociente de un objeto proyectivo, sea
  $K'_{\bullet\bullet}$ un bicomplejo en $\cat{C}'$ cuyos dos grados estén
  acotados inferiormente, y sea $T$ un bifuntor covariante aditivo de
  $\cat{C}\times\cat{C}'$ en una categoría abeliana $\cat{C}''$.
  Si $L_\bullet$ y $L'_\bullet$ son los complejos simples asociados a
  $K_{\bullet\bullet}$ y $K'_{\bullet\bullet}$, respectivamente, se define
  la hiperhomología de $T$ relativa a los dos bicomplejos
  $K_{\bullet\bullet}$ y $K'_{\bullet\bullet}$ como
  \[
    \LL_\bullet T(L_\bullet,L'_\bullet).
  \]
  Aplicando \sref{0.11.6.4} y \sref{0.11.6.5}, se obtienen, como en
  \sref{0.11.7.2}, seis sucesiones espectrales que convergen a esta hiperhomología;
  dejamos al lector su formulación explícita.
\end{enumerate}
\end{env}

\subsection{Suplemento sobre la cohomología de los complejos simpliciales}
\label{subsection:0.11.8}

\begin{env}[11.8.1]
\label{0.11.8.1}
Sea $A$ un conjunto finito, y sea $\Sigma(A)$ el conjunto de las sucesiones
finitas
\[
  \sigma=(\alpha_0,\ldots,\alpha_h)
\]
de elementos de $A$ (los ``símplices'' de $A$); pongamos
$|\sigma|=\{\alpha_0,\ldots,\alpha_h\}$.
Recordemos que el complejo de cadenas $\mathrm{C}_\bullet(A)$ es el grupo abeliano
libre graduado engendrado por los elementos de $\Sigma(A)$, con
$(\alpha_0,\ldots,\alpha_h)$ en grado $h$, y con diferencial
\[
  d(\alpha_0,\ldots,\alpha_h)
  =
  \sum_{j=0}^{h}(-1)^j
  (\alpha_0,\ldots,\widehat{\alpha_j},\ldots,\alpha_h).
\]
El subgrupo $\mathrm{D}_\bullet(A)$ de $\mathrm{C}_\bullet(A)$ engendrado
por las cadenas $\sigma=(\alpha_0,\ldots,\alpha_h)$ para las cuales dos de los
$\alpha_i$ son iguales, y por las cadenas
\[
  \pi(\sigma)-\varepsilon_\pi\sigma,
\]
donde, para toda permutación $\pi$,
\[
  \pi(\sigma)
  =
  (\alpha_{\pi(0)},\ldots,\alpha_{\pi(h)})
\]
y $\varepsilon_\pi$ es el signo de $\pi$, es un subcomplejo de
$\mathrm{C}_\bullet(A)$.
Sus elementos se llaman \emph{cadenas degeneradas}.
Pongamos
\[
  \mathrm{L}_\bullet(A)
  =
  \mathrm{C}_\bullet(A)/\mathrm{D}_\bullet(A);
\]
existe un homomorfismo natural de complejos
\[
  p:\mathrm{C}_\bullet(A)\longrightarrow\mathrm{L}_\bullet(A).
\]

Recíprocamente, definamos un homomorfismo de complejos
\[
  j:\mathrm{L}_\bullet(A)\longrightarrow\mathrm{C}_\bullet(A)
\]
como sigue.
Elijamos un orden total sobre $A$.
A la clase módulo $\mathrm{D}_\bullet(A)$ de un símplex
$\sigma=(\alpha_0,\ldots,\alpha_h)$, asociemos $0$ si dos de los
$\alpha_i$ son iguales y, en caso contrario, asociemos la sucesión
$(\beta_0,\ldots,\beta_h)$ obtenida ordenando los $\alpha_i$ en
orden creciente.
Está claro que $p\circ j$ es la identidad de
$\mathrm{L}_\bullet(A)$.
\end{env}

\begin{env}[11.8.2]
\label{0.11.8.2}
Sea $B$ un segundo conjunto finito.
Si $d'$ y $d''$ son los diferenciales de
$\mathrm{C}_\bullet(A)$ y $\mathrm{C}_\bullet(B)$, el complejo producto
tensorial
\[
  \mathrm{C}_\bullet(A)\otimes\mathrm{C}_\bullet(B)
\]
puede considerarse como el grupo abeliano libre engendrado por los elementos de
$\Sigma(A)\times\Sigma(B)$, con diferencial
\[
  d(\sigma,\tau)
  =
  (d'\sigma,\tau)+(-1)^{h+1}(\sigma,d''\tau)
  \qquad\text{si }\operatorname{Card}|\sigma|=h+1.
\]

Los homomorfismos naturales
$\mathrm{C}_\bullet(A)\to\mathrm{L}_\bullet(A)$ y
$\mathrm{C}_\bullet(B)\to\mathrm{L}_\bullet(B)$ definen un homomorfismo
\[
  p:
  \mathrm{C}_\bullet(A)\otimes\mathrm{C}_\bullet(B)
  \longrightarrow
  \mathrm{L}_\bullet(A)\otimes\mathrm{L}_\bullet(B),
\]
siendo este último producto tensorial isomorfo a
\[
  \frac{\mathrm{C}_\bullet(A)\otimes\mathrm{C}_\bullet(B)}
  {\mathrm{C}_\bullet(A)\otimes\mathrm{D}_\bullet(B)
   +\mathrm{D}_\bullet(A)\otimes\mathrm{C}_\bullet(B)}.
\]
Análogamente, utilizando los homomorfismos
$\mathrm{L}_\bullet(A)\to\mathrm{C}_\bullet(A)$ y
$\mathrm{L}_\bullet(B)\to\mathrm{C}_\bullet(B)$ definidos en
\sref{0.11.8.1}, con órdenes totales sobre $A$ y $B$, se obtiene un
homomorfismo
\[
  j:
  \mathrm{L}_\bullet(A)\otimes\mathrm{L}_\bullet(B)
  \longrightarrow
  \mathrm{C}_\bullet(A)\otimes\mathrm{C}_\bullet(B)
\]
tal que $p\circ j$ es la identidad.
\end{env}

\oldpage[0\textsubscript{III}]{44}
Con esta notación:
\begin{proposition}[11.8.3]
\label{0.11.8.3}
Existe una homotopía
\[
  h:
  \mathrm{C}_\bullet(A)\otimes\mathrm{C}_\bullet(B)
  \longrightarrow
  \mathrm{C}_\bullet(A)\otimes\mathrm{C}_\bullet(B)
\]
tal que $h(\sigma,\tau)$ es una combinación lineal de pares de símplices
$(\sigma_i,\tau_i)$ que satisfacen
$|\sigma_i|\subset|\sigma|$ y $|\tau_i|\subset|\tau|$, y tal que,
si $f=j\circ p$, entonces
\[
  f-1=h\circ d+d\circ h.
  \tag{11.8.3.1}\label{0.11.8.3.1}
\]
\end{proposition}

\begin{proof}
Basta definir $h$ sobre cada par $(\sigma,\tau)$ de símplices, por
inducción sobre la suma de los grados de $\sigma$ y $\tau$; cuando esa suma
es $0$, puede tomarse $h=0$.
Pongamos
\[
  \omega
  =
  f(\sigma,\tau)-(\sigma,\tau)-h\bigl(d(\sigma,\tau)\bigr).
\]
Por la hipótesis de inducción y la definición de $d$, se tiene
\[
  \omega
  \in
  \mathrm{C}_\bullet(|\sigma|)
  \otimes
  \mathrm{C}_\bullet(|\tau|).
\]
Además,
\[
  d\omega
  =
  f\bigl(d(\sigma,\tau)\bigr)
  -d(\sigma,\tau)
  -d\bigl(h(d(\sigma,\tau))\bigr)
  =
  h\bigl(d(d(\sigma,\tau))\bigr)
  =
  0
\]
por \eqref{0.11.8.3.1} y la hipótesis de inducción.
Ahora bien,
\[
  \HH_q\bigl(\mathrm{C}_\bullet(A)\bigr)=0
  \qquad(q>0)
\]
(G, I, 3.7.4), y, por consiguiente,
\[
  \HH_q\bigl(
    \mathrm{C}_\bullet(A)\otimes\mathrm{C}_\bullet(B)
  \bigr)=0
  \qquad(q>0)
\]
por la fórmula de K\"unneth (G, I, 5.5.2).
Aplicando esta observación con $A$ reemplazado por $|\sigma|$ y $B$ por
$|\tau|$, se encuentra
\[
  \omega'
  \in
  \mathrm{C}_\bullet(|\sigma|)
  \otimes
  \mathrm{C}_\bullet(|\tau|)
\]
tal que $\omega=d\omega'$.
Tomar $h(\sigma,\tau)=\omega'$ verifica \eqref{0.11.8.3.1} para el par
$(\sigma,\tau)$ y permite continuar la inducción.
\end{proof}

\begin{env}[11.8.4]
\label{0.11.8.4}
Con la notación de \sref{0.11.8.2}, escribamos
\[
  (\sigma,\tau)\leq(\sigma',\tau')
\]
si $|\sigma|\subset|\sigma'|$ y $|\tau|\subset|\tau'|$.
Un \emph{sistema de coeficientes} $\sh{S}$ sobre
$\Sigma(A)\times\Sigma(B)$ consiste en una familia
$(\Gamma_{\sigma,\tau})$ de grupos abelianos, donde
$\Gamma_{\sigma,\tau}$ depende únicamente de los conjuntos $|\sigma|$ y
$|\tau|$, junto con homomorfismos
\[
  \Gamma_{\sigma,\tau}
  \longrightarrow
  \Gamma_{\sigma',\tau'}
  \qquad
  \bigl((\sigma,\tau)\leq(\sigma',\tau')\bigr)
\]
que forman un sistema inductivo para este preorden.

Definamos un complejo de cocadenas $\mathrm{C}^\bullet(A,B;\sh{S})$ cuyos elementos
son las familias
\[
  \lambda=\bigl(\lambda(\sigma,\tau)\bigr),
  \qquad
  \lambda(\sigma,\tau)\in\Gamma_{\sigma,\tau},
\]
cuando $(\sigma,\tau)$ recorre $\Sigma(A)\times\Sigma(B)$.
Su diferencial se define como sigue.
Si
\[
  d(\sigma,\tau)=\sum_i\pm(\sigma_i,\tau_i),
\]
entonces $|\sigma_i|\subset|\sigma|$ y $|\tau_i|\subset|\tau|$ para todo
$i$, y se pone
\[
  d\lambda(\sigma,\tau)
  =
  \sum_i\pm\lambda_i(\sigma_i,\tau_i),
\]
donde $\lambda_i(\sigma_i,\tau_i)$ denota la imagen canónica de
$\lambda(\sigma_i,\tau_i)$ en $\Gamma_{\sigma,\tau}$.

Una cocadena $\lambda\in\mathrm{C}^\bullet(A,B;\sh{S})$ se llama
\emph{bialternada} si $\lambda(\sigma,\tau)=0$ siempre que $\sigma$
o $\tau$ tenga dos términos iguales, y si
\[
  \lambda(\pi(\sigma),\tau)
  =
  \varepsilon_\pi\lambda(\sigma,\tau),
  \qquad
  \lambda(\sigma,\pi'(\tau))
  =
  \varepsilon_{\pi'}\lambda(\sigma,\tau)
\]
para permutaciones arbitrarias $\pi$ y $\pi'$ de los índices.
Estas cocadenas engendran manifiestamente un subcomplejo
$\mathrm{L}^\bullet(A,B;\sh{S})$ de
$\mathrm{C}^\bullet(A,B;\sh{S})$.
\end{env}

\begin{proposition}[11.8.5]
\label{0.11.8.5}
La inyección canónica
\[
  \mathrm{L}^\bullet(A,B;\sh{S})
  \longrightarrow
  \mathrm{C}^\bullet(A,B;\sh{S})
\]
induce un isomorfismo en cohomología.
\end{proposition}

\begin{proof}
Si $p$ y $j$ tienen los significados asignados en \sref{0.11.8.2}, entonces
las aplicaciones
\[
  {}'\!p:\lambda\longmapsto\lambda\circ p,
  \qquad
  {}'\!j:\lambda\longmapsto\lambda\circ j
\]
están definidas sobre $\mathrm{L}^\bullet(A,B;\sh{S})$ y
$\mathrm{C}^\bullet(A,B;\sh{S})$, respectivamente; la primera es precisamente
la inyección canónica.
Como ${}'\!j\circ{}'\!p$ es la identidad, solo queda demostrar
que ${}'\!p\circ{}'\!j$ es homótopa a la identidad.
Por \sref{0.11.8.3}, la aplicación
\[
  {}'\!h:\lambda\longmapsto\lambda\circ h
\]
está definida sobre $\mathrm{C}^\bullet(A,B;\sh{S})$, de modo que al trasponer
\eqref{0.11.8.3.1} se obtiene el resultado deseado.
\end{proof}

\begin{env}[11.8.6]
\label{0.11.8.6}
\oldpage[0\textsubscript{III}]{45}
La Proposición~\sref{0.11.8.5} reduce el cálculo de la cohomología de
$\mathrm{L}^\bullet(A,B;\sh{S})$ al de
$\mathrm{C}^\bullet(A,B;\sh{S})$.
Por otra parte, por el teorema de Eilenberg--Zilber
(G, I, 3.10.2), esta última cohomología es canónicamente isomorfa a la
cohomología obtenida a partir del complejo de cadenas
$\mathrm{P}_\bullet(A,B)$ de combinaciones lineales de pares
$(\sigma,\tau)\in\Sigma(A)\times\Sigma(B)$ para los cuales $\sigma$ y
$\tau$ tienen el mismo grado.
El diferencial de este complejo es
\[
  d(\sigma,\tau)
  =
  \sum_{j,k}(-1)^{j+k}(\sigma_j,\tau_k)
\]
cuando
\[
  d\sigma=\sum_j(-1)^j\sigma_j,
  \qquad
  d\tau=\sum_k(-1)^k\tau_k.
\]
Existen dos homomorfismos canónicos de complejos
\[
  f:\mathrm{P}_\bullet(A,B)
  \longrightarrow
  \mathrm{C}_\bullet(A)\otimes\mathrm{C}_\bullet(B),
  \qquad
  g:\mathrm{C}_\bullet(A)\otimes\mathrm{C}_\bullet(B)
  \longrightarrow
  \mathrm{P}_\bullet(A,B),
\]
y se demuestra allí que existen homotopías $h$ y $h'$ tales que
\[
  f\circ g-1=d\circ h+h\circ d,
  \qquad
  g\circ f-1=d\circ h'+h'\circ d.
\]
Además,
\[
  f(\sigma,\tau)
  \in
  \mathrm{C}_\bullet(|\sigma|)
  \otimes
  \mathrm{C}_\bullet(|\tau|),
  \qquad
  g(\sigma,\tau)
  \in
  \mathrm{P}_\bullet(|\sigma|,|\tau|),
\]
y las homotopías pueden elegirse de modo que
\[
  h(\sigma,\tau)
  \in
  \mathrm{C}_\bullet(|\sigma|)
  \otimes
  \mathrm{C}_\bullet(|\tau|),
  \qquad
  h'(\sigma,\tau)
  \in
  \mathrm{P}_\bullet(|\sigma|,|\tau|).
\]
En efecto, $h(\sigma,\tau)$ y $h'(\sigma,\tau)$ pueden definirse por
inducción sobre la suma de los grados de $\sigma$ y $\tau$, mientras que
\[
  \HH_q\bigl(
    \mathrm{C}_\bullet(|\sigma|)
    \otimes
    \mathrm{C}_\bullet(|\tau|)
  \bigr)
  =
  \HH_q\bigl(\mathrm{P}_\bullet(|\sigma|,|\tau|)\bigr)
  =
  0
  \qquad(q>0)
\]
(\emph{loc.\ cit.}); el razonamiento es el mismo que en
\sref{0.11.8.3}.

Definamos ahora $\mathrm{P}^\bullet(A,B;\sh{S})$ como el conjunto de familias
\[
  \lambda=\bigl(\lambda(\sigma,\tau)\bigr),
  \qquad
  \lambda(\sigma,\tau)\in\Gamma_{\sigma,\tau},
\]
donde $(\sigma,\tau)$ recorre los pares de símplices del mismo grado.
Cuando
\[
  d\sigma=\sum_j(-1)^j\sigma_j,
  \qquad
  d\tau=\sum_k(-1)^k\tau_k,
\]
la fórmula
\[
  d\lambda(\sigma,\tau)
  =
  \sum_{j,k}(-1)^{j+k}\lambda(\sigma_j,\tau_k)
\]
está definida y da el diferencial de
$\mathrm{P}^\bullet(A,B;\sh{S})$.
Las aplicaciones
\[
  {}'\!f:\lambda\longmapsto\lambda\circ f,\quad
  {}'\!g:\lambda\longmapsto\lambda\circ g,\quad
  {}'\!h:\lambda\longmapsto\lambda\circ h,\quad
  {}'\!h':\lambda\longmapsto\lambda\circ h'
\]
están todas definidas por las propiedades de soporte anteriores.
Así, ${}'\!f\circ{}'\!g$ y ${}'\!g\circ{}'\!f$ son homótopas a la
identidad, dando el isomorfismo requerido entre la cohomología de
$\mathrm{C}^\bullet(A,B;\sh{S})$ y la de
$\mathrm{P}^\bullet(A,B;\sh{S})$.
\end{env}

\begin{env}[11.8.7]
\label{0.11.8.7}
\emph{Observación}.
El razonamiento de \sref{0.11.8.3}, aplicado a
$\mathrm{C}_\bullet(A)$ y $\mathrm{L}_\bullet(A)$, muestra que si
$j$ y $p$ se definen como en \sref{0.11.8.1}, entonces
$f=j\circ p$ satisface de nuevo \eqref{0.11.8.3.1}, con
$|h(\sigma)|\subset|\sigma|$.
Como en \sref{0.11.8.5}, se obtiene un isomorfismo de la cohomología de
$\mathrm{L}^\bullet(A;\sh{S})$ en la de
$\mathrm{C}^\bullet(A;\sh{S})$, definiéndose los dos complejos de manera
evidente.
Este es el resultado cuya demostración se esboza en (G, I, 3.8.1).
\end{env}

\begin{env}[11.8.8]
\label{0.11.8.8}
Volvamos ahora a la notación y las hipótesis de \sref{0.11.8.2}, y sea
\[
  \sh{S}^\bullet=(\sh{S}^k)
\]
un complejo de sistemas de coeficientes sobre
$\Sigma(A)\times\Sigma(B)$.
Para cada $(\sigma,\tau)$, los grupos $\Gamma^k_{\sigma,\tau}$
forman, por tanto, un complejo de grupos abelianos $(k\in\bb{Z})$, y para
$(\sigma,\tau)\leq(\sigma',\tau')$ existen diagramas conmutativos
\[
  \xymatrix{
    \Gamma^k_{\sigma,\tau}\ar[r]\ar[d] &
    \Gamma^{k+1}_{\sigma,\tau}\ar[d]\\
    \Gamma^k_{\sigma',\tau'}\ar[r] &
    \Gamma^{k+1}_{\sigma',\tau'}.
  }
\]
\oldpage[0\textsubscript{III}]{46}
Se deduce inmediatamente que
\[
  \mathrm{C}^\bullet(A,B;\sh{S}^\bullet)
  =
  \bigl(
    \mathrm{C}^h(A,B;\sh{S}^k)
  \bigr)_{(h,k)\in\bb{Z}\times\bb{Z}}
\]
es un bicomplejo de grupos abelianos, y que
\[
  \mathrm{L}^\bullet(A,B;\sh{S}^\bullet)
  =
  \bigl(
    \mathrm{L}^h(A,B;\sh{S}^k)
  \bigr)
\]
es un subbicomplejo.
\end{env}

\begin{proposition}[11.8.9]
\label{0.11.8.9}
La inyección canónica
\[
  \mathrm{L}^\bullet(A,B;\sh{S}^\bullet)
  \longrightarrow
  \mathrm{C}^\bullet(A,B;\sh{S}^\bullet)
\]
induce un isomorfismo en la cohomología de estos dos bicomplejos.
\end{proposition}

\begin{proof}
Por brevedad, pongamos
\[
  \mathrm{C}^{\bullet\bullet}
  =
  \mathrm{C}^\bullet(A,B;\sh{S}^\bullet),
  \qquad
  \mathrm{L}^{\bullet\bullet}
  =
  \mathrm{L}^\bullet(A,B;\sh{S}^\bullet).
\]
Como $\mathrm{C}^{hk}=\mathrm{L}^{hk}=0$ para $h<0$, las segundas
sucesiones espectrales de estos bicomplejos son regulares
\sref{0.11.3.3}.
El homomorfismo
$\mathrm{L}^{\bullet\bullet}\to\mathrm{C}^{\bullet\bullet}$ da, por tanto,
un morfismo de sucesiones espectrales
\[
  {}''\mathrm{E}(\mathrm{L}^{\bullet\bullet})
  \longrightarrow
  {}''\mathrm{E}(\mathrm{C}^{\bullet\bullet})
\]
que, sobre los términos $\mathrm{E}_2$, es
\[
  \HH^p_{\mathrm{II}}
  \bigl(\HH^q_{\mathrm{I}}(\mathrm{L}^{\bullet\bullet})\bigr)
  \longrightarrow
  \HH^p_{\mathrm{II}}
  \bigl(\HH^q_{\mathrm{I}}(\mathrm{C}^{\bullet\bullet})\bigr).
  \tag{11.8.9.1}\label{0.11.8.9.1}
\]
Sin embargo, para todo $k\in\bb{Z}$, \sref{0.11.8.5} muestra que
\[
  \HH^q_{\mathrm{I}}(\mathrm{L}^{\bullet,k})
  \longrightarrow
  \HH^q_{\mathrm{I}}(\mathrm{C}^{\bullet,k})
\]
es biyectivo.
La conclusión se deduce de \sref{0.11.1.5}.
\end{proof}

\begin{env}[11.8.10]
\label{0.11.8.10}
Análogamente, con la notación de \sref{0.11.8.6}, existe un
homomorfismo canónico de bicomplejos
\[
  \mathrm{C}^\bullet(A,B;\sh{S}^\bullet)
  \longrightarrow
  \mathrm{P}^\bullet(A,B;\sh{S}^\bullet),
\]
con la notación evidente.
El razonamiento de \sref{0.11.8.9}, utilizando ahora \sref{0.11.8.6}, muestra que
este homomorfismo induce asimismo un isomorfismo en cohomología.
\end{env}

\subsection{Un lema sobre los complejos de tipo finito}
\label{subsection:0.11.9}

\begin{proposition}[11.9.1]
\label{0.11.9.1}
Sea $\cat{C}$ una categoría abeliana, y sean $K'$ y $K''$ subconjuntos
del conjunto de objetos de $\cat{C}$, con $K''\subset K'$, que satisfacen las
condiciones siguientes:
\begin{enumerate}
  \item[{\rm(i)}]
  \emph{Para todo $A'\in K'$ y todo epimorfismo $u:A\to A'$ en
  $\cat{C}$, existen un objeto $B\in K''$ y un morfismo $v:B\to A$
  tales que $u\circ v$ es un epimorfismo.}

  \item[{\rm(ii)}]
  \emph{Para todo par $A,B\in K'$ y todo epimorfismo $u:A\to B$,
  el objeto $\Ker(u)$ pertenece a $K'$.}

  \item[{\rm(iii)}]
  \emph{El producto de dos objetos de $K''$ pertenece a $K''$.}
\end{enumerate}

\emph{Sea $P_\bullet=(P_i)_{i\in\bb{Z}}$ un complejo en $\cat{C}$ tal
que $\HH_i(P_\bullet)\in K'$ para todo $i$, y supongamos que existe un
entero $d$ para el cual $\HH_i(P_\bullet)=0$ cuando $i<d$.
Entonces existen un complejo $Q_\bullet=(Q_i)_{i\in\bb{Z}}$ en $\cat{C}$
tal que $Q_i\in K''$ para todo $i$ y $Q_i=0$ para $i<d$, y un
homomorfismo de complejos
\[
  u:Q_\bullet\longrightarrow P_\bullet
\]
cuyo morfismo inducido
\[
  \HH_\bullet(Q_\bullet)\longrightarrow\HH_\bullet(P_\bullet)
\]
es un isomorfismo.}
\end{proposition}

\begin{proof}
Demostremos primero la consecuencia siguiente de la condición~(i).

\medskip\noindent
\emph{(i bis)}
\emph{Sea $u:C\to B$ un epimorfismo en $\cat{C}$, sea $A$ un
objeto de $K'$, y sea $v:A\to B$ un morfismo en $\cat{C}$.
Entonces existen un objeto $D\in K''$, un epimorfismo $u':D\to A$ y un
morfismo $v':D\to C$ tales que el diagrama}
\[
  \xymatrix{
    D\ar[r]^{u'}\ar[d]_{v'} & A\ar[d]^{v}\\
    C\ar[r]_{u} & B
  }
  \tag{11.9.1.1}\label{0.11.9.1.1}
\]
\emph{es conmutativo.}

Formemos el producto fibrado $C\times_B A$ en $\cat{C}$ y sean
\[
  p:C\times_B A\longrightarrow C,
  \qquad
  q:C\times_B A\longrightarrow A
\]
las proyecciones canónicas, de modo que
\oldpage[0\textsubscript{III}]{47}
\[
  \xymatrix{
    C\times_B A\ar[r]^{q}\ar[d]_{p} & A\ar[d]^{v}\\
    C\ar[r]_{u} & B
  }
  \tag{11.9.1.2}\label{0.11.9.1.2}
\]
es conmutativo.
El conúcleo de $q$ es el cociente de $A$ por $v^{-1}(u(C))$
([27, p.~1-12]).
Como $u$ es un epimorfismo, $u(C)=B$ y
$v^{-1}(u(C))=A$; por tanto, $q$ es un epimorfismo.
Aplicando~(i) a $q:C\times_B A\to A$, se obtienen un objeto
$D\in K''$ y un morfismo $w:D\to C\times_B A$ tales que $q\circ w$ es
un epimorfismo.
Basta tomar $u'=q\circ w$ y $v'=p\circ w$.

\medskip
Demostremos ahora la proposición por inducción.
Supongamos, para algún $i\geq d-1$, que para todo $j\leq i$ hemos
construido objetos $Q_j$, diferenciales
\[
  d^Q_j:Q_j\longrightarrow Q_{j-1}
\]
con $Q_j=0$ para $j<d$ y
$d^Q_{j-1}\circ d^Q_j=0$, y morfismos
\[
  u_j:Q_j\longrightarrow P_j
\]
tales que
\[
  d^P_j\circ u_j=u_{j-1}\circ d^Q_j
  \qquad(j\leq i).
\]
Supongamos además que se cumplen las condiciones siguientes:
\begin{enumerate}
  \item[$(\mathrm{I}_i)$]
  $Q_j\in K''$ para $j\leq i$, y
  $\mathrm{B}_j(Q_\bullet)\in K'$ para $j<i$.

  \item[$(\mathrm{II}_i)$]
  Para $j<i$, el homomorfismo
  \[
    \HH_j(Q_\bullet)\longrightarrow\HH_j(P_\bullet)
  \]
  inducido por $(u_k)_{k\leq i}$ es un isomorfismo.

  \item[$(\mathrm{III}_i)$]
  El compuesto
  \[
    v_i:
    \mathrm{Z}_i(Q_\bullet)
    \longrightarrow
    \mathrm{Z}_i(P_\bullet)
    \longrightarrow
    \HH_i(P_\bullet),
  \]
  donde la primera flecha es la restricción de $u_i$ y la segunda es
  canónica, es un epimorfismo.
\end{enumerate}

Por~(ii), $\mathrm{Z}_i(Q_\bullet)$ pertenece a $K'$: es el núcleo del
epimorfismo
\[
  Q_i\longrightarrow\mathrm{B}_{i-1}(Q_\bullet),
\]
y aquí se utiliza $(\mathrm{I}_i)$.
La condición~(ii) y $(\mathrm{III}_i)$ muestran también que
\[
  N_i=\Ker(v_i)
\]
pertenece a $K'$.
Por (i bis), existen un objeto $Q'_{i+1}\in K''$, un epimorfismo
\[
  d'_{i+1}:Q'_{i+1}\longrightarrow N_i,
\]
y un morfismo $u'_{i+1}:Q'_{i+1}\to P_{i+1}$ tales que
\[
  \xymatrix{
    Q'_{i+1}\ar[r]^{d'_{i+1}}\ar[d]_{u'_{i+1}} &
    N_i\ar[d]^{u_i}\\
    P_{i+1}\ar[r]_{d^P_{i+1}} &
    \mathrm{B}_i(P_\bullet)
  }
  \tag{11.9.1.3}\label{0.11.9.1.3}
\]
es conmutativo.

El morfismo canónico
\[
  \mathrm{Z}_{i+1}(P_\bullet)
  \longrightarrow
  \HH_{i+1}(P_\bullet)
\]
es un epimorfismo, y $\HH_{i+1}(P_\bullet)\in K'$.
La condición~(i) proporciona, por tanto, un objeto $Q''_{i+1}\in K''$ y un
morfismo
\[
  u''_{i+1}:Q''_{i+1}\longrightarrow
  \mathrm{Z}_{i+1}(P_\bullet)
\]
tal que el compuesto
\[
  Q''_{i+1}
  \longrightarrow
  \mathrm{Z}_{i+1}(P_\bullet)
  \longrightarrow
  \HH_{i+1}(P_\bullet)
\]
es un epimorfismo.
Si $d''_{i+1}:Q''_{i+1}\to N_i$ es el morfismo nulo, entonces
\[
  \xymatrix{
    Q''_{i+1}\ar[r]^{d''_{i+1}}\ar[d]_{u''_{i+1}} &
    N_i\ar[d]^{u_i}\\
    \mathrm{Z}_{i+1}(P_\bullet)\ar[r]_{d^P_{i+1}} &
    P_i
  }
  \tag{11.9.1.4}\label{0.11.9.1.4}
\]
es conmutativo; la flecha horizontal inferior es nula.

\oldpage[0\textsubscript{III}]{48}
Pongamos
\[
  Q_{i+1}=Q'_{i+1}\times Q''_{i+1},
  \qquad
  d^Q_{i+1}=d'_{i+1}+d''_{i+1},
  \qquad
  u_{i+1}=u'_{i+1}+u''_{i+1}.
\]
Por~(iii), $Q_{i+1}$ pertenece a $K''$.
Como
\[
  d^Q_{i+1}(Q_{i+1})
  =
  d'_{i+1}(Q'_{i+1})
  =
  N_i
  \subset
  \mathrm{Z}_i(Q_\bullet),
\]
se tiene
$d^Q_i\circ d^Q_{i+1}=0$ y
$\mathrm{B}_i(Q_\bullet)=N_i$; esto demuestra
$(\mathrm{I}_{i+1})$.
La conmutatividad de \eqref{0.11.9.1.3} y
\eqref{0.11.9.1.4} da
\[
  d^P_{i+1}\circ u_{i+1}
  =
  u_i\circ d^Q_{i+1}.
\]
Por la definición de $N_i$, el morfismo
\[
  \HH_i(Q_\bullet)
  =
  \mathrm{Z}_i(Q_\bullet)/N_i
  \longrightarrow
  \HH_i(P_\bullet)
\]
inducido por los $u_k$ para $k\leq i+1$ es un isomorfismo, puesto que $v_i$ es
un epimorfismo; por tanto, se cumple $(\mathrm{II}_{i+1})$.
Finalmente, el factor $Q''_{i+1}$ está contenido en
$\mathrm{Z}_{i+1}(Q_\bullet)$.
La elección de $u''_{i+1}$ muestra que
\[
  v_{i+1}:
  \mathrm{Z}_{i+1}(Q_\bullet)
  \longrightarrow
  \mathrm{Z}_{i+1}(P_\bullet)
  \longrightarrow
  \HH_{i+1}(P_\bullet)
\]
es un epimorfismo, porque su restricción a $Q''_{i+1}$ ya lo
es.
Así se cumple $(\mathrm{III}_{i+1})$, la inducción continúa y se deduce la
proposición.
\end{proof}

\begin{corollary}[11.9.2]
\label{0.11.9.2}
Sea $A$ un anillo noetheriano, no necesariamente conmutativo, y sea
$P_\bullet=(P_i)_{i\in\bb{Z}}$ un complejo de $A$-módulos a derecha.
Supongamos que todo $\HH_i(P_\bullet)$ es un $A$-módulo finitamente
engendrado y que existe un entero $d$ tal que
$\HH_i(P_\bullet)=0$ para $i<d$.
Entonces existen un complejo $Q_\bullet=(Q_i)_{i\in\bb{Z}}$ de
$A$-módulos libres de rango finito, con $Q_i=0$ para $i<d$, y un homomorfismo
de complejos
\[
  u:Q_\bullet\longrightarrow P_\bullet
\]
cuyo morfismo inducido
\[
  \HH_\bullet(Q_\bullet)\longrightarrow\HH_\bullet(P_\bullet)
\]
es un isomorfismo.
\end{corollary}

\begin{proof}
Apliquemos \sref{0.11.9.1} tomando como $\cat{C}$ la categoría de los
$A$-módulos a derecha, como $K'$ el conjunto de los $A$-módulos finitamente engendrados y
como $K''$ el conjunto de los $A$-módulos libres de rango finito.
Las condiciones~(i), (ii) y~(iii) de \sref{0.11.9.1} se deducen inmediatamente
de la hipótesis de que $A$ es noetheriano.
\end{proof}

\begin{env}[11.9.3]
\label{0.11.9.3}
\emph{Observaciones}.
\begin{enumerate}
  \item[{\rm(i)}]
  Bajo las hipótesis de \sref{0.11.9.2}, supongamos además que los
  $P_i$ son $A$-módulos planos a derecha.
  Entonces, para todo $A$-módulo a izquierda $M$, el homomorfismo de complejos
  \[
    u\otimes 1:
    Q_\bullet\otimes_A M
    \longrightarrow
    P_\bullet\otimes_A M
  \]
  induce de nuevo un isomorfismo
  \[
    \HH_\bullet(Q_\bullet\otimes_A M)
    \isoto
    \HH_\bullet(P_\bullet\otimes_A M),
  \]
  como se verá en el Capítulo~III.

  \item[{\rm(ii)}]
  La conclusión de \sref{0.11.9.2} no tiene por qué cumplirse cuando $A$ no es
  noetheriano.
  En efecto, aplicarla a un complejo que solo tuviese un término no nulo
  implicaría que todo $A$-módulo por la derecha finitamente engendrado admite una resolución por
  módulos libres finitamente engendrados, lo que es falso en general
(Bourbaki, \emph{Alg\`ebre commutative}, cap.~I, \S2, ejerc.~6).

  Sin embargo, es posible reemplazar la hipótesis de que $A$ es
  noetheriano por la suposición de que los $\HH_i(P_\bullet)$ admiten
  $\infty$-presentaciones finitas (cf.\ Capítulo~IV).
\end{enumerate}
\end{env}

\subsection{Característica de Euler--Poincar\'e de un complejo de módulos de longitud finita}
\label{subsection:0.11.10}

\begin{env}[11.10.1]
\label{0.11.10.1}
Sea $A$ un anillo, no necesariamente conmutativo, y sea
\[
  M^\bullet:
  0\longrightarrow M^0\longrightarrow M^1
  \longrightarrow\cdots\longrightarrow M^n\longrightarrow 0
  \tag{11.10.1.1}\label{0.11.10.1.1}
\]
un complejo de $A$-módulos a izquierda de longitud finita.
La \emph{característica de Euler--Poincar\'e} de este complejo es el
entero
\oldpage[0\textsubscript{III}]{49}
\[
  \chi(M^\bullet)
  =
  \sum_{i=0}^{n}(-1)^i\operatorname{length}(M^i).
  \tag{11.10.1.2}\label{0.11.10.1.2}
\]
\end{env}

\begin{proposition}[11.10.2]
\label{0.11.10.2}
Para todo complejo finito $M^\bullet$ de $A$-módulos a izquierda de longitud
finita,
\[
  \chi(M^\bullet)=\chi\bigl(\HH^\bullet(M^\bullet)\bigr),
\]
donde $\HH^\bullet(M^\bullet)$ se considera como un complejo con diferencial
nulo.
En particular, si la sucesión \eqref{0.11.10.1.1} es exacta, entonces
$\chi(M^\bullet)=0$.
\end{proposition}

\begin{proof}
Por brevedad, escribamos
\[
  \mathrm{B}^i=\mathrm{B}^i(M^\bullet),
  \qquad
  \mathrm{Z}^i=\mathrm{Z}^i(M^\bullet),
  \qquad
  \HH^i=\HH^i(M^\bullet)=\mathrm{Z}^i/\mathrm{B}^i.
\]
Los módulos $\mathrm{B}^i$, $\mathrm{Z}^i$ y $\HH^i$ tienen todos
longitud finita.
Las sucesiones exactas
\[
  0\longrightarrow\mathrm{B}^i
  \longrightarrow\mathrm{Z}^i
  \longrightarrow\HH^i
  \longrightarrow0
\]
y
\[
  0\longrightarrow\mathrm{Z}^i
  \longrightarrow M^i
  \longrightarrow\mathrm{B}^{i+1}
  \longrightarrow0
\]
dan
\[
  \operatorname{length}(\mathrm{Z}^i)
  =
  \operatorname{length}(\HH^i)
  +
  \operatorname{length}(\mathrm{B}^i)
\]
y
\[
  \operatorname{length}(M^i)
  =
  \operatorname{length}(\mathrm{Z}^i)
  +
  \operatorname{length}(\mathrm{B}^{i+1}).
\]
Por consiguiente,
\[
  \operatorname{length}(M^i)-\operatorname{length}(\HH^i)
  =
  \operatorname{length}(\mathrm{B}^{i+1})
  +
  \operatorname{length}(\mathrm{B}^i).
\]
Multipliquemos esta igualdad por $(-1)^i$ y sumemos para $0\leq i\leq n$.
Como $\mathrm{B}^0=\mathrm{B}^{n+1}=0$, se deduce la igualdad deseada.
\end{proof}

\begin{corollary}[11.10.3]
\label{0.11.10.3}
Sea $\mathrm{E}=(\mathrm{E}_r^{pq})$ una sucesión espectral en la
categoría de los módulos sobre un anillo $A$.
Supongamos que los $\mathrm{E}_2^{pq}$ son $A$-módulos de longitud finita y
que $\mathrm{E}_2^{pq}\neq0$ solo para un número finito de pares $(p,q)$.
Entonces las características de Euler--Poincar\'e
\[
  \chi\bigl(\mathrm{E}_r^{(\bullet)}\bigr)
\]
de todos los complejos
\[
  \mathrm{E}_r^{(\bullet)}
  =
  \bigl(\mathrm{E}_r^{(n)}\bigr)_{n\in\bb{Z}}
\]
de \sref{0.11.1.1} son iguales.
Si, además, $\mathrm{E}$ es débilmente convergente y se pone
\[
  \mathrm{E}_\infty^{(n)}
  =
  \bigoplus_{p+q=n}\mathrm{E}_\infty^{pq}
  \qquad(n\in\bb{Z}),
\]
entonces
\[
  \chi\bigl(\mathrm{E}_\infty^{(\bullet)}\bigr)
  =
  \chi\bigl(\mathrm{E}_2^{(\bullet)}\bigr),
\]
donde
\[
  \mathrm{E}_\infty^{(\bullet)}
  =
  \bigl(\mathrm{E}_\infty^{(n)}\bigr)_{n\in\bb{Z}}
\]
se considera como un complejo con diferencial nulo.
\end{corollary}

\begin{proof}
Observemos primero que si $\mathrm{E}_2^{pq}=0$, entonces
$\mathrm{E}_r^{pq}=0$ para $2\leq r\leq+\infty$.
Así, todo complejo $\mathrm{E}_r^{(\bullet)}$ es finito y está formado por
$A$-módulos de longitud finita.
Para todo $r$ finito, la igualdad
\[
  \chi\bigl(\mathrm{E}_r^{(\bullet)}\bigr)
  =
  \chi\bigl(\mathrm{E}_{r+1}^{(\bullet)}\bigr)
\]
se deduce, por tanto, de \sref{0.11.10.2} y del isomorfismo
\[
  \HH^\bullet\bigl(\mathrm{E}_r^{(\bullet)}\bigr)
  \simeq
  \mathrm{E}_{r+1}^{(\bullet)}
\]
de complejos con diferencial nulo.

Como los $\mathrm{E}_r^{pq}$ tienen longitud finita, para cada par
$(p,q)$ las sucesiones
\[
  \bigl(\mathrm{B}_k(\mathrm{E}_2^{pq})\bigr)_{k\geq2},
  \qquad
  \bigl(\mathrm{Z}_k(\mathrm{E}_2^{pq})\bigr)_{k\geq2}
\]
son estacionarias.
La convergencia débil, junto con el hecho de que
$\mathrm{E}_2^{pq}=0$ salvo para un número finito de pares $(p,q)$, implica, por tanto,
que algún entero $r\geq2$ satisface
\[
  \mathrm{E}_\infty^{pq}=\mathrm{E}_r^{pq}
\]
para todo par $(p,q)$.
Se deduce la aserción relativa a
$\chi(\mathrm{E}_\infty^{(\bullet)})$.
\end{proof}

\section{Suplemento sobre la cohomología de haces}
\label{section:0.12}

\subsection{Cohomología de haces de módulos sobre espacios anillados}
\label{subsection:0.12.1}

\begin{env}[12.1.1]
\label{0.12.1.1}
\oldpage[0\textsubscript{III}]{49}
Sea $(X,\sh{O}_X)$ un espacio anillado.
Recordemos que, para todo $\sh{O}_X$-módulo $\sh{F}$, se define la
cohomología $\HH^\bullet(X,\sh{F})$, que es un funtor cohomológico universal
(T, 2.2) de la categoría $\cat{C}(X)$ de
$\sh{O}_X$-módulos en la categoría de grupos abelianos; es el funtor derivado
del funtor exacto por la izquierda
\[
  \sh{F}\longmapsto\Gamma(X,\sh{F}).
\]
El funtor $\HH^\bullet(X,\sh{F})$ es isomorfo
\oldpage[0\textsubscript{III}]{50}
a la restricción a la categoría $\cat{C}(X)$ del funtor cohomológico
definido del mismo modo sobre la categoría de haces de grupos abelianos
en $X$ (G, II, 7.2.1).
\end{env}

\begin{env}[12.1.2]
\label{0.12.1.2}
Pongamos $A=\Gamma(X,\sh{O}_X)$.
Como todo elemento de $A$ define un endomorfismo del grupo abeliano
$\Gamma(X,\sh{F})$, define por funtorialidad un endomorfismo del
$\delta$-funtor $\HH^\bullet(X,\sh{F})$.
Estos endomorfismos dotan a cada $\HH^p(X,\sh{F})$ de una
estructura de $A$-módulo, y el operador de conexión $\partial$ es
$A$-lineal.
Además, para cualesquiera dos enteros no negativos $p,q$ y cualesquiera dos
$\sh{O}_X$-módulos $\sh{F},\sh{G}$, existe un homomorfismo de
$A$-módulos, llamado \emph{producto cup},
\[
\label{0.12.1.2.1}
\tag{12.1.2.1}
  \HH^p(X,\sh{F})\otimes_A\HH^q(X,\sh{G})
  \longrightarrow
  \HH^{p+q}\bigl(X,\sh{F}\otimes_{\sh{O}_X}\sh{G}\bigr)
\]
(G, II, 6.6).
Estos homomorfismos hacen de la suma directa $S$ de los
$\HH^p(X,\sh{O}_X)$, para $p\geq0$, una $A$-álgebra graduada
anticomutativa, y de la suma directa de los $\HH^p(X,\sh{F})$ un
$S$-módulo graduado.

Para todo recubrimiento abierto $\mathfrak{U}$ de $X$, denotaremos siempre
por $\check{\mathrm{C}}^\bullet(\mathfrak{U},\sh{F})$ (contrariamente a
(G, II, 5.1)) el complejo de cocadenas alternadas del nervio de
$\mathfrak{U}$ con valores en el sistema de coeficientes
$\Gamma(U_\alpha,\sh{F})$.
Está claro que
$\check{\mathrm{C}}^\bullet(\mathfrak{U},\sh{F})$ es un
$A$-módulo graduado, de modo que los grupos de cohomología
$\check{\HH}^p(\mathfrak{U},\sh{F})$ de este complejo están dotados de
una estructura de $A$-módulo.
Además, las aplicaciones canónicas
\[
  \check{\HH}^p(\mathfrak{U},\sh{F})
  \longrightarrow
  \HH^p(X,\sh{F})
\]
(G, II, 5.4) son $A$-lineales.
\end{env}

\begin{env}[12.1.3]
\label{0.12.1.3}
Sea $(X',\sh{O}_{X'})$ un segundo espacio anillado, y sea
$f=(\psi,\theta)$ un morfismo de $X'$ en $X$.
Pongamos $A'=\Gamma(X',\sh{O}_{X'})$; el $\psi$-morfismo $\theta$
define canónicamente un homomorfismo de anillos $A\to A'$.
Sea $\sh{F}$ un $\sh{O}_X$-módulo y $\sh{F}'$ un
$\sh{O}_{X'}$-módulo.
Para todo $f$-morfismo $u:\sh{F}\to\sh{F}'$
\sref[0\textsubscript{I}]{0.4.4.1}, veremos que se puede definir,
para todo $p\geq0$, un dihomomorfismo
\[
\label{0.12.1.3.1}
\tag{12.1.3.1}
  u_p:\HH^p(X,\sh{F})\longrightarrow\HH^p(X',\sh{F}').
\]

En efecto, como $\psi^*$ es exacto en la categoría de haces de grupos
abelianos sobre $X$, el funtor
\[
  \sh{F}\longmapsto\HH^\bullet(X',\psi^*(\sh{F}))
\]
es un $\delta$-funtor sobre dicha categoría, y se sabe que existe un
homomorfismo canónico de $\delta$-funtores
\[
\label{0.12.1.3.2}
\tag{12.1.3.2}
  \HH^\bullet(X,\sh{F})
  \longrightarrow
  \HH^\bullet(X',\psi^*(\sh{F}))
\]
determinado de manera única por la condición de que, en grado $0$, se reduce
al homomorfismo canónico
\[
  \Gamma(X,\sh{F})
  \longrightarrow
  \Gamma(X',\psi^*(\sh{F}))
\]
(T, 3.2.2).
Además, todo elemento de $A$ determina un endomorfismo $\mu$ de
$\Gamma(X,\sh{F})$ y un endomorfismo $\mu'$ de
$\Gamma(X',\psi^*(\sh{F}))$ tal que el diagrama
\[
\label{0.12.1.3.3}
\tag{12.1.3.3}
  \xymatrix{
    \Gamma(X,\sh{F})\ar[r]\ar[d]_\mu
      & \Gamma(X',\psi^*(\sh{F}))\ar[d]^{\mu'}\\
    \Gamma(X,\sh{F})\ar[r]
      & \Gamma(X',\psi^*(\sh{F}))
  }
\]
es conmutativo.
\oldpage[0\textsubscript{III}]{51}
Por la propiedad de unicidad para prolongar morfismos de funtores
cohomológicos universales (T, 2.2), se obtienen prolongaciones únicas de
$\mu$ y $\mu'$ a la cohomología que hacen conmutativos los diagramas
análogos a \sref{0.12.1.3.3}.
Así, \sref{0.12.1.3.2} es un homomorfismo de $A$-módulos.

Observemos ahora que
\[
  f^*(\sh{F})
  =
  \psi^*(\sh{F})
  \otimes_{\psi^*(\sh{O}_X)}
  \sh{O}_{X'},
\]
y, por tanto, que existe un dihomomorfismo canónico
\[
  \psi^*(\sh{F})\longrightarrow f^*(\sh{F})
\]
del $\psi^*(\sh{O}_X)$-módulo $\psi^*(\sh{F})$ en el
$\sh{O}_{X'}$-módulo $f^*(\sh{F})$.
Por funtorialidad se obtiene, pues, un dihomomorfismo funtorial
\[
\label{0.12.1.3.4}
\tag{12.1.3.4}
  \HH^p(X',\psi^*(\sh{F}))
  \longrightarrow
  \HH^p(X',f^*(\sh{F})),
\]
siendo $A$ y $A'$ los anillos correspondientes.
Componiendo este dihomomorfismo con \sref{0.12.1.3.2}, se obtiene un
dihomomorfismo canónico, funtorial en $\sh{F}$,
\[
\label{0.12.1.3.5}
\tag{12.1.3.5}
  \theta_p:\HH^p(X,\sh{F})
  \longrightarrow
  \HH^p(X',f^*(\sh{F})).
\]

Finalmente, por funtorialidad, el homomorfismo
$u^\sharp:f^*(\sh{F})\to\sh{F}'$ da un homomorfismo de
$A'$-módulos
\[
  \HH^p(X',f^*(\sh{F}))
  \longrightarrow
  \HH^p(X',\sh{F}'),
\]
que, compuesto con \sref{0.12.1.3.5}, da
\sref{0.12.1.3.1}.

Sea $f'=(\psi',\theta'):X''\to X'$ un segundo morfismo de espacios
anillados, y sea $f''=f\circ f'$ el morfismo compuesto.
Teniendo en cuenta la compatibilidad del funtor $\psi'^*$ con los
productos tensoriales \sref[0\textsubscript{I}]{0.4.3.3}, se comprueba
inmediatamente que el compuesto del dihomomorfismo
\[
  \HH^p(X',f^*(\sh{F}))
  \longrightarrow
  \HH^p\bigl(X'',f'^*(f^*(\sh{F}))\bigr)
\]
con \sref{0.12.1.3.5} es el dihomomorfismo correspondiente
\[
  \HH^p(X,\sh{F})
  \longrightarrow
  \HH^p(X'',f''^*(\sh{F})).
\]
\end{env}

\begin{env}[12.1.4]
\label{0.12.1.4}
Puede darse una definición directa del homomorfismo
\sref{0.12.1.3.2} como sigue.
Tomemos una resolución inyectiva
$\sh{L}^\bullet=(\sh{L}^i)$ de $\sh{F}$ constituida por haces de
grupos abelianos sobre $X$.
Como el funtor $\psi^*$ es exacto,
$\psi^*(\sh{L}^\bullet)$ es una resolución de $\psi^*(\sh{F})$
constituida por haces sobre $X'$.
Si $\sh{L}'^\bullet=(\sh{L}'^i)$ es una resolución inyectiva de
$\psi^*(\sh{F})$ en la categoría de haces de grupos abelianos sobre
$X'$, existe un morfismo
\[
  \psi^*(\sh{L}^\bullet)\longrightarrow\sh{L}'^\bullet
\]
de complejos de haces de grupos abelianos, compatible con los
aumentos (M, V, 1.1.a), y determinado salvo homotopía.
Se obtienen así homomorfismos
\[
  \Gamma(X,\sh{L}^\bullet)
  \longrightarrow
  \Gamma(X',\psi^*(\sh{L}^\bullet))
  \longrightarrow
  \Gamma(X',\sh{L}'^\bullet)
\]
de complejos de grupos abelianos.
Su compuesto, al pasar a la cohomología, da un morfismo de
$\delta$-funtores
\[
  \HH^\bullet(X,\sh{F})
  \longrightarrow
  \HH^\bullet(X',\psi^*(\sh{F}));
\]
como coincide con \sref{0.12.1.3.2} en grado $0$, es idéntico a él
(T, 2.2).

Consideremos ahora un recubrimiento abierto
$\mathfrak{U}=(U_\alpha)$ de $X$, y sea
$\mathfrak{U}'=(f^{-1}(U_\alpha))$ su recubrimiento abierto imagen inversa
de $X'$.
Para todo subconjunto abierto $V$ de $X$, los homomorfismos canónicos
\[
  \Gamma(V,\sh{F})
  \longrightarrow
  \Gamma(f^{-1}(V),f^*(\sh{F}))
\]
definen inmediatamente (cf.\ G, II, 5.1) un homomorfismo de complejos
\[
  \check{\mathrm{C}}^\bullet(\mathfrak{U},\sh{F})
  \longrightarrow
  \check{\mathrm{C}}^\bullet(\mathfrak{U}',f^*(\sh{F})),
\]
y, por tanto, homomorfismos canónicos
\[
\label{0.12.1.4.1}
\tag{12.1.4.1}
  \theta_p:
  \check{\HH}^p(\mathfrak{U},\sh{F})
  \longrightarrow
  \check{\HH}^p(\mathfrak{U}',f^*(\sh{F})).
\]

\oldpage[0\textsubscript{III}]{52}
Además, existen diagramas conmutativos
\[
\label{0.12.1.4.2}
\tag{12.1.4.2}
  \xymatrix{
    \check{\HH}^p(\mathfrak{U},\sh{F})
      \ar[r]^-{\theta_p}\ar[d]
      & \check{\HH}^p(\mathfrak{U}',f^*(\sh{F}))\ar[d]\\
    \HH^p(X,\sh{F})
      \ar[r]_-{\theta_p}
      & \HH^p(X',f^*(\sh{F})),
  }
\]
donde las flechas verticales son los homomorfismos canónicos de
(G, II, 5.2).
Para demostrar la conmutatividad de \sref{0.12.1.4.2}, consideremos el
complejo de haces de cocadenas alternadas de $\sh{F}$ relativo a
$\mathfrak{U}$,
$\sh{C}^\bullet(\mathfrak{U},\sh{F})$, para el cual
\[
  \Gamma\bigl(X,\sh{C}^\bullet(\mathfrak{U},\sh{F})\bigr)
  =
  \check{\mathrm{C}}^\bullet(\mathfrak{U},\sh{F})
\]
(G, II, 5.2).
Los homomorfismos canónicos
\[
  \Gamma(V,\sh{F})
  \longrightarrow
  \Gamma(f^{-1}(V),\psi^*(\sh{F}))
\]
definen entonces un $\psi$-morfismo
\[
  \sh{C}^\bullet(\mathfrak{U},\sh{F})
  \longrightarrow
  \sh{C}^\bullet(\mathfrak{U}',\psi^*(\sh{F})),
\]
y, con la notación precedente, existe un diagrama conmutativo
\[
  \xymatrix{
    \Gamma\bigl(X,\sh{C}^\bullet(\mathfrak{U},\sh{F})\bigr)
      \ar[r]\ar[d]
      & \Gamma\bigl(X',
          \sh{C}^\bullet(\mathfrak{U}',\psi^*(\sh{F}))\bigr)\ar[d]\\
    \Gamma(X,\sh{L}^\bullet)\ar[r]
      & \Gamma(X',\psi^*(\sh{L}^\bullet))\ar[r]
      & \Gamma(X',\sh{L}'^\bullet).
  }
\]
Al pasar a la cohomología se obtienen diagramas conmutativos
\[
  \xymatrix{
    \check{\HH}^p(\mathfrak{U},\sh{F})\ar[r]\ar[d]
      & \check{\HH}^p(\mathfrak{U}',\psi^*(\sh{F}))\ar[d]\\
    \HH^p(X,\sh{F})\ar[r]
      & \HH^p(X',\psi^*(\sh{F})),
  }
\]
donde las flechas verticales son los homomorfismos canónicos de
(G, II, 5.2).
Solo queda combinar estos diagramas con los diagramas conmutativos
\[
  \xymatrix{
    \check{\HH}^p(\mathfrak{U}',\psi^*(\sh{F}))\ar[r]\ar[d]
      & \check{\HH}^p(\mathfrak{U}',f^*(\sh{F}))\ar[d]\\
    \HH^p(X',\psi^*(\sh{F}))\ar[r]
      & \HH^p(X',f^*(\sh{F})),
  }
\]
\oldpage[0\textsubscript{III}]{53}
que provienen del homomorfismo
$\psi^*(\sh{F})\to f^*(\sh{F})$ y de la funtorialidad de los
homomorfismos canónicos de (G, II, 5.2), para obtener la conmutatividad
de \sref{0.12.1.4.2}.

Observemos que si
$A=\Gamma(X,\sh{O}_X)$ y $A'=\Gamma(X',\sh{O}_{X'})$, entonces
\sref{0.12.1.4.1} es un dihomomorfismo de módulos correspondiente a
los anillos $A$ y $A'$.
Existe para \sref{0.12.1.4.1} una propiedad de transitividad respecto
de un compuesto de dos morfismos, análoga a la de
\sref{0.12.1.3.5}.
Finalmente, en las definiciones precedentes, en lugar de una resolución
inyectiva $\sh{L}^\bullet$ de $\sh{F}$, se podría haber partido igualmente
de una resolución tal que
\[
  \HH^p(X,\sh{L}^i)=0
  \qquad\text{para todo $i$ y todo $p>0$}
\]
(G, II, 4.7.1).
\end{env}

\begin{env}[12.1.5]
\label{0.12.1.5}
Sean $\sh{F},\sh{G},\sh{H}$ tres $\sh{O}_X$-módulos, y consideremos
un $\sh{O}_X$-homomorfismo
\[
  u:\sh{F}\otimes_{\sh{O}_X}\sh{G}\longrightarrow\sh{H},
\]
que da, en cohomología, los homomorfismos
\[
\label{0.12.1.5.1}
\tag{12.1.5.1}
  \HH^p(X,\sh{F})\otimes_A\HH^q(X,\sh{G})
  \longrightarrow
  \HH^{p+q}(X,\sh{H})
\]
deducidos del producto cup \sref{0.12.1.2.1}.
Demostraremos que, bajo las hipótesis y con la notación de
\sref{0.12.1.3}, los diagramas
\[
\label{0.12.1.5.2}
\tag{12.1.5.2}
  \xymatrix@C=10pt{
    \HH^p(X,\sh{F})\otimes_A\HH^q(X,\sh{G})
      \ar[r]\ar[d]
      & \HH^{p+q}(X,\sh{H})\ar[d]\\
    \HH^p(X',f^*(\sh{F}))\otimes_{A'}
      \HH^q(X',f^*(\sh{G}))
      \ar[r]
      & \HH^{p+q}(X',f^*(\sh{H}))
  }
\]
son conmutativos, donde las flechas verticales provienen de los
homomorfismos canónicos \sref{0.12.1.3.5}.

Recordemos para ello que \sref{0.12.1.5.1} puede obtenerse a partir de
las resoluciones canónicas (G, II, 4.3)
$\sh{L}^\bullet,\sh{M}^\bullet,\sh{N}^\bullet$ de
$\sh{F},\sh{G},\sh{H}$ respectivamente, que están constituidas por
$\sh{O}_X$-módulos, y de la aplicación lineal
\[
  \sh{L}^\bullet\otimes_{\sh{O}_X}\sh{M}^\bullet
  \longrightarrow
  \sh{N}^\bullet
\]
de complejos de $\sh{O}_X$-módulos correspondiente a $u$.
Esta aplicación proporciona un homomorfismo de complejos de $A$-módulos
\[
  \Gamma(X,\sh{L}^\bullet)\otimes_A\Gamma(X,\sh{M}^\bullet)
  \longrightarrow
  \Gamma(X,\sh{N}^\bullet),
\]
y, por tanto, al pasar a la cohomología, homomorfismos
\[
  \HH^p\bigl(\Gamma(X,\sh{L}^\bullet)\bigr)
  \otimes_A
  \HH^q\bigl(\Gamma(X,\sh{M}^\bullet)\bigr)
  \longrightarrow
  \HH^{p+q}\bigl(\Gamma(X,\sh{N}^\bullet)\bigr)
\]
(G, II, 6.6).
Existe claramente un diagrama conmutativo
\[
\label{0.12.1.5.3}
\tag{12.1.5.3}
  \xymatrix@C=8pt{
    \Gamma(X,\sh{L}^\bullet)\otimes_A
      \Gamma(X,\sh{M}^\bullet)
      \ar[r]\ar[d]
      & \Gamma(X,\sh{N}^\bullet)\ar[d]\\
    \Gamma(X',\psi^*(\sh{L}^\bullet))
      \otimes_{\Gamma(X',\psi^*(\sh{O}_X))}
      \Gamma(X',\psi^*(\sh{M}^\bullet))
      \ar[r]
      & \Gamma(X',\psi^*(\sh{N}^\bullet)).
  }
\]
\oldpage[0\textsubscript{III}]{54}
Al pasar a la cohomología se obtienen diagramas conmutativos
\[
\label{0.12.1.5.4}
\tag{12.1.5.4}
  \xymatrix@C=7pt{
    \HH^p(X,\sh{F})\otimes_A\HH^q(X,\sh{G})
      \ar[r]\ar[d]
      & \HH^{p+q}(X,\sh{H})\ar[d]\\
    \HH^p\bigl(\Gamma(X',\psi^*(\sh{L}^\bullet))\bigr)
      \otimes_{\Gamma(X',\psi^*(\sh{O}_X))}
      \HH^q\bigl(\Gamma(X',\psi^*(\sh{M}^\bullet))\bigr)
      \ar[r]
      & \HH^{p+q}\bigl(\Gamma(X',\psi^*(\sh{N}^\bullet))\bigr).
  }
\]
Pero como
$\psi^*(\sh{L}^\bullet)$,
$\psi^*(\sh{M}^\bullet)$, y
$\psi^*(\sh{N}^\bullet)$
son resoluciones de
$\psi^*(\sh{F})$,
$\psi^*(\sh{G})$, y
$\psi^*(\sh{H})$ respectivamente, existe un diagrama conmutativo
(G, II, 6.6.1)
\[
\label{0.12.1.5.5}
\tag{12.1.5.5}
  \xymatrix@C=7pt{
    \HH^p\bigl(\Gamma(X',\psi^*(\sh{L}^\bullet))\bigr)
      \otimes_{\Gamma(X',\psi^*(\sh{O}_X))}
      \HH^q\bigl(\Gamma(X',\psi^*(\sh{M}^\bullet))\bigr)
      \ar[r]\ar[d]
      & \HH^{p+q}\bigl(\Gamma(X',\psi^*(\sh{N}^\bullet))\bigr)
        \ar[d]\\
    \HH^p(X',\psi^*(\sh{F}))
      \otimes_{\Gamma(X',\psi^*(\sh{O}_X))}
      \HH^q(X',\psi^*(\sh{G}))
      \ar[r]
      & \HH^{p+q}(X',\psi^*(\sh{H})).
  }
\]
Finalmente, por funtorialidad, existe un diagrama conmutativo
\[
\label{0.12.1.5.6}
\tag{12.1.5.6}
  \xymatrix@C=7pt{
    \HH^p(X',\psi^*(\sh{F}))
      \otimes_{\Gamma(X',\psi^*(\sh{O}_X))}
      \HH^q(X',\psi^*(\sh{G}))
      \ar[r]\ar[d]
      & \HH^{p+q}(X',\psi^*(\sh{H}))\ar[d]\\
    \HH^p(X',f^*(\sh{F}))\otimes_{A'}
      \HH^q(X',f^*(\sh{G}))
      \ar[r]
      & \HH^{p+q}(X',f^*(\sh{H})).
  }
\]
Combinando los tres diagramas \sref{0.12.1.5.4},
\sref{0.12.1.5.5} y \sref{0.12.1.5.6}, se obtiene el diagrama
conmutativo deseado \sref{0.12.1.5.2}.
\end{env}

\begin{remark}[12.1.6]
\label{0.12.1.6}
\oldpage[0\textsubscript{III}]{55}
Con la notación de \sref{0.12.1.3}, supongamos que existe un
diagrama conmutativo
\[
\label{0.12.1.6.1}
\tag{12.1.6.1}
  \xymatrix{
    0\ar[r] & \sh{F}\ar[r]^r\ar[d]_u
      & \sh{G}\ar[r]^s\ar[d]_v
      & \sh{H}\ar[r]\ar[d]_w & 0\\
    0\ar[r] & \sh{F}'\ar[r]_{r'}
      & \sh{G}'\ar[r]_{s'}
      & \sh{H}'\ar[r] & 0,
  }
\]
donde $r,s$ son homomorfismos de $\sh{O}_X$-módulos,
$r',s'$ son homomorfismos de $\sh{O}_{X'}$-módulos,
$u,v,w$ son $f$-morfismos, y las filas son exactas.
Se sigue que el diagrama
\[
\label{0.12.1.6.2}
\tag{12.1.6.2}
  \xymatrix@C=6pt{
    \cdots\ar[r]
      & \HH^p(X,\sh{F})\ar[r]\ar[d]_{u_p}
      & \HH^p(X,\sh{G})\ar[r]\ar[d]_{v_p}
      & \HH^p(X,\sh{H})\ar[r]^-\partial\ar[d]_{w_p}
      & \HH^{p+1}(X,\sh{F})\ar[r]\ar[d]_{u_{p+1}}
      & \cdots\\
    \cdots\ar[r]
      & \HH^p(X',\sh{F}')\ar[r]
      & \HH^p(X',\sh{G}')\ar[r]
      & \HH^p(X',\sh{H}')\ar[r]_-\partial
      & \HH^{p+1}(X',\sh{F}')\ar[r]
      & \cdots
  }
\]
es conmutativo.
En efecto, \sref{0.12.1.6.1} se factoriza como
\[
  \xymatrix@R=14pt{
    0\ar[r] & \sh{F}\ar[r]\ar[d]
      & \sh{G}\ar[r]\ar[d]
      & \sh{H}\ar[r]\ar[d] & 0\\
    0\ar[r] & \psi^*(\sh{F})\ar[r]\ar[d]
      & \psi^*(\sh{G})\ar[r]\ar[d]
      & \psi^*(\sh{H})\ar[r]\ar[d] & 0\\
    0\ar[r] & \sh{F}'\ar[r]
      & \sh{G}'\ar[r]
      & \sh{H}'\ar[r] & 0,
  }
\]
donde la fila intermedia es exacta
\sref[0\textsubscript{I}]{0.3.7.2}.
Basta entonces utilizar el hecho de que \sref{0.12.1.3.2} es un
homomorfismo de $\delta$-funtores y que los
$\HH^p(X',\sh{F}')$ forman un $\delta$-funtor en $\sh{F}'$.
\end{remark}

\begin{env}[12.1.7]
\label{0.12.1.7}
Con las hipótesis y la notación de \sref{0.12.1.3}, consideremos ahora el
caso
\[
  \sh{F}=f_*(\sh{F}')=\psi_*(\sh{F}').
\]
Veremos que el dihomomorfismo definido en \sref{0.12.1.3},
\[
  \HH^p(X,f_*(\sh{F}'))\longrightarrow\HH^p(X',\sh{F}'),
  \tag{12.1.7.1}
  \label{0.12.1.7.1}
\]
puede obtenerse (salvo un automorfismo de $\HH^p(X',\sh{F}')$) como un
homomorfismo de borde de una sucesión espectral para el funtor compuesto
\[
  \sh{F}'\leadsto\Gamma(X,\psi_*(\sh{F}'))
\]
(T, 2.4).
La descripción de \sref{0.12.1.7.1} dada en \sref{0.12.1.4} muestra
aquí que este homomorfismo puede obtenerse como sigue.
Tomemos resoluciones inyectivas $\sh{L}^\bullet$ y
${\sh{L}'}^\bullet$ de $\psi_*(\sh{F}')$ y $\sh{F}'$,
respectivamente, y después un homomorfismo de complejos
\[
  v:\psi^*(\sh{L}^\bullet)\longrightarrow{\sh{L}'}^\bullet
\]
``por encima'' del homomorfismo canónico
$\psi^*(\psi_*(\sh{F}'))\to\sh{F}'$.
\oldpage[0\textsubscript{III}]{56}
Observemos a continuación que
\[
  \Gamma(X',{\sh{L}'}^\bullet)
  =\Gamma(X,\psi_*({\sh{L}'}^\bullet))
\]
y que el homomorfismo compuesto
\[
  \Gamma(X,\sh{L}^\bullet)
  \longrightarrow\Gamma(X',\psi^*(\sh{L}^\bullet))
  \xrightarrow{\Gamma(v)}\Gamma(X',{\sh{L}'}^\bullet)
\]
no es otro que
\[
  \Gamma(v^\flat):
  \Gamma(X,\sh{L}^\bullet)
  \longrightarrow\Gamma(X,\psi_*({\sh{L}'}^\bullet))
  \tag{12.1.7.2}
  \label{0.12.1.7.2}
\]
\sref[0\textsubscript{I}]{0.3.7.1}; y \sref{0.12.1.7.1} se
obtiene pasando a la cohomología en \sref{0.12.1.7.2}.

Por otra parte, las sucesiones espectrales para el funtor compuesto
\[
  \sh{F}'\leadsto\Gamma(X,\psi_*(\sh{F}'))
\]
se obtienen tomando una resolución inyectiva de Cartan--Eilenberg
\[
  \sh{M}^{\bullet\bullet}=(\sh{M}^{ij})
\]
del complejo $\psi_*({\sh{L}'}^\bullet)$ en la categoría de
haces de grupos abelianos sobre $X$.
Las sucesiones espectrales en cuestión son las del bicomplejo
$\Gamma(X,\sh{M}^{\bullet\bullet})$; son birregulares, pues
$\sh{M}^{ij}=0$ para $i<0$ o $j<0$.
Ahora bien, la primera sucesión espectral de este bicomplejo degenera, pues
los haces $\psi_*({\sh{L}'}^i)$ son flácidos (G, II, 3.1.1), y
por tanto
\[
  \HH_\mathrm{II}^q
    \bigl(\Gamma(X,\sh{M}^{i,\bullet})\bigr)
  =\HH^q\bigl(\psi_*({\sh{L}'}^i)\bigr)=0
  \qquad(q>0)
\]
(G, II, 4.4.3).
Por consiguiente, existen homomorfismos de borde biyectivos
\sref{0.11.1.6}
\[
  {}'E_2^{i0}
  =\HH^i\!\left(
      \HH_\mathrm{II}^0
        \bigl(\Gamma(X,\sh{M}^{\bullet\bullet})\bigr)
    \right)
  \longrightarrow
  \HH^i\bigl(\Gamma(X,\sh{M}^{\bullet\bullet})\bigr),
  \tag{12.1.7.3}
  \label{0.12.1.7.3}
\]
y sabemos por \sref{0.11.3.4} que este homomorfismo proviene, al pasar
a la cohomología, del aumento
\[
  \Gamma(X,\psi_*({\sh{L}'}^\bullet))
  \longrightarrow\Gamma(X,\sh{M}^{\bullet\bullet}),
  \tag{12.1.7.4}
  \label{0.12.1.7.4}
\]
que a su vez proviene del aumento
\[
  \eta:\psi_*({\sh{L}'}^\bullet)
  \longrightarrow\sh{M}^{\bullet\bullet}.
\]

Para la segunda sucesión espectral, por otra parte, existen homomorfismos
de borde
\[
  {}''E_2^{i0}
  =\HH^i\!\left(
      \HH_\mathrm{I}^0
        \bigl(\Gamma(X,\sh{M}^{\bullet\bullet})\bigr)
    \right)
  \longrightarrow
  \HH^i\bigl(\Gamma(X,\sh{M}^{\bullet\bullet})\bigr),
  \tag{12.1.7.5}
  \label{0.12.1.7.5}
\]
que provienen, por \sref{0.11.3.4}, al pasar a la cohomología, del
homomorfismo de complejos
\[
  Z_\mathrm{I}^0
    \bigl(\Gamma(X,\sh{M}^{\bullet\bullet})\bigr)
  \longrightarrow\Gamma(X,\sh{M}^{\bullet\bullet}).
\]
Como $\psi_*$ es exacto por la izquierda, la sucesión
\[
  0\longrightarrow\psi_*(\sh{F}')
  \longrightarrow\psi_*({\sh{L}'}^0)
  \longrightarrow\psi_*({\sh{L}'}^1)
\]
es exacta.
Por la definición de una resolución de Cartan--Eilenberg
\sref{0.11.4.2}, podemos tomar, por tanto,
\[
  B_\mathrm{I}^0(\sh{M}^{\bullet\bullet})=0,
  \qquad
  Z_\mathrm{I}^0(\sh{M}^{\bullet\bullet})=\sh{L}^\bullet.
\]
Como el diagrama
\[
  \xymatrix@C=50pt@R=24pt{
    \sh{L}^0\ar[r]^{i^0}
      & \sh{M}^{00}\\
    \psi_*(\sh{F}')\ar[u]^\varepsilon
      \ar[r]_{\varepsilon'}\ar[ur]^{\varepsilon''}
      & \psi_*({\sh{L}'}^0)\ar[u]_{\eta^0}
  }
\]
es conmutativo, la inyección de complejos
$i:\sh{L}^\bullet\to\sh{M}^{\bullet\bullet}$ es compatible con los
aumentos $\varepsilon$ y $\varepsilon''$.
Existen, pues, dos homomorfismos de complejos de $\sh{L}^\bullet$
en $\sh{M}^{\bullet\bullet}$,
\[
  \xymatrix@C=52pt@R=22pt{
    \sh{L}^\bullet\ar[r]^i\ar[dr]_{v^\flat}
      & \sh{M}^{\bullet\bullet}\\
    & \psi_*({\sh{L}'}^\bullet)\ar[u]_\eta
  }
\]
\oldpage[0\textsubscript{III}]{57}
compatibles con los aumentos $\varepsilon$ y
$\varepsilon''$.
Como $\sh{L}^\bullet$ es una resolución inyectiva y
$\sh{M}^{\bullet\bullet}$ está constituido por haces inyectivos, se sigue
de (M, V, 1.1(a)) que estos dos homomorfismos son homotópicos.
Lo mismo vale, por tanto, para los dos homomorfismos correspondientes
\[
  \Gamma(X,\sh{L}^\bullet)
  \longrightarrow\Gamma(X,\sh{M}^{\bullet\bullet}),
\]
y, al pasar a la cohomología, dan por consiguiente el mismo
homomorfismo.
En otras palabras, hemos demostrado que el homomorfismo de borde
\sref{0.12.1.7.5}, escrito
\[
  \HH^p(X,\psi_*(\sh{F}'))
  \longrightarrow\HH^p\bigl(\Gamma(X,\sh{M}^{\bullet\bullet})\bigr),
\]
es el compuesto de \sref{0.12.1.7.1} y
\sref{0.12.1.7.3}, estando este último escrito
\[
  \HH^p(X',\sh{F}')
  \longrightarrow\HH^p\bigl(\Gamma(X,\sh{M}^{\bullet\bullet})\bigr)
\]
y habiéndose visto que es un isomorfismo.
Esto demuestra la aserción.
\end{env}

\subsection{Imágenes directas superiores}
\label{subsection:0.12.2}

\begin{env}[12.2.1]
\label{0.12.2.1}
Sean $(X,\sh{O}_X)$ y $(Y,\sh{O}_Y)$ dos espacios anillados, y sea
$f=(\psi,\theta)$ un morfismo de $X$ en $Y$.  Define el funtor imagen
directa
\[
  f_*:\cat{C}(X)\longrightarrow\cat{C}(Y),
\]
que es además idéntico a la restricción a $\cat{C}(X)$ del
funtor $\psi_*$ definido sobre la categoría de haces de grupos abelianos
en $X$.
Este último funtor es aditivo y exacto por la izquierda; y como todo haz de
grupos abelianos sobre $X$ es isomorfo a un subhaz de un haz inyectivo
de grupos abelianos, se definen los funtores derivados por la derecha
\[
  \sh{F}\longmapsto\RR^p\psi_*(\sh{F})
\]
de $\psi_*$.  Los $\RR^p\psi_*(\sh{F})$ son haces de grupos abelianos
sobre $Y$, y los $\RR^p\psi_*$ forman un funtor cohomológico universal
(T, 2.3).

Además, $\RR^p\psi_*(\sh{F})$ es el haz asociado al
prehaz
\[
  V\longmapsto\HH^p(f^{-1}(V),\sh{F})
\]
(T, 3.7.2).
Supongamos ahora que $\sh{F}$ es un $\sh{O}_X$-módulo.
Entonces $\HH^p(f^{-1}(V),\sh{F})$ está dotado naturalmente de la estructura
de $\Gamma(f^{-1}(V),\sh{O}_X)$-módulo, y por tanto de la de
$\Gamma(V,\psi_*(\sh{O}_X))$-módulo; el homomorfismo
\[
  \theta:\sh{O}_Y\longrightarrow\psi_*(\sh{O}_X)
\]
lo dota entonces de la estructura de
$\Gamma(V,\sh{O}_Y)$-módulo.
Para las estructuras así definidas, la restricción de un subconjunto abierto
$V$ a un subconjunto abierto $V'\subset V$ define un dihomomorfismo.  Esto
define, por tanto, sobre todo $\RR^p\psi_*(\sh{F})$ la estructura de un
$\sh{O}_Y$-módulo.  Denotamos este $\sh{O}_Y$-módulo por
$\RR^p f_*(\sh{F})$; así, $\RR^p f_*$ queda definido como un funtor
aditivo de $\cat{C}(X)$ en $\cat{C}(Y)$.

Además, los $\RR^p f_*$ forman un $\delta$-funtor.  En efecto, si
\[
  0\longrightarrow\sh{F}'\longrightarrow\sh{F}
  \longrightarrow\sh{F}''\longrightarrow0
\]
es una sucesión exacta de $\sh{O}_X$-módulos, la descripción precedente
de los $\RR^p\psi_*$ y de la estructura de $\sh{O}_Y$-módulo sobre
$\RR^p\psi_*(\sh{F})$ muestra inmediatamente que el homomorfismo de conexión
\[
  \partial:\RR^p\psi_*(\sh{F}'')
  \longrightarrow\RR^{p+1}\psi_*(\sh{F}')
\]
es un homomorfismo de $\sh{O}_Y$-módulos.
Finalmente, los $\RR^p f_*$ se identifican con los funtores derivados por la
derecha de $f_*$.  Todo $\sh{O}_X$-módulo admite una resolución inyectiva por
$\sh{O}_X$-módulos, y tal resolución está constituida por haces flácidos
de grupos abelianos (G, II, 7.1); puede, por tanto, utilizarse para calcular
los $\RR^p\psi_*(\sh{F})$, pues $\RR^n\psi_*(\sh{G})=0$ para $n\geq1$
y todo haz flácido $\sh{G}$ (T, 2.4.1, Observación~3, y el
 corolario de la Proposición~\EGAUnavailableHyperref{0.3.3.2}{3.3.2}).
Se sigue que los $\RR^p f_*$ forman un funtor cohomológico universal
de $\cat{C}(X)$ en $\cat{C}(Y)$ (T, 2.3).
\end{env}

\begin{env}[12.2.2]
\label{0.12.2.2}
Sean $\sh{F}$ y $\sh{G}$ dos $\sh{O}_X$-módulos.
Con la notación de \sref{0.12.2.1}, para todo subconjunto abierto $V$ de
$Y$ existe el homomorfismo de producto cup \sref{0.12.1.2.1}
\[
  \begin{aligned}
  &\HH^p(f^{-1}(V),\sh{F})
    \otimes_{\Gamma(f^{-1}(V),\sh{O}_X)}
    \HH^q(f^{-1}(V),\sh{G})\\
  &\hspace{55mm}\longrightarrow
    \HH^{p+q}\bigl(f^{-1}(V),
      \sh{F}\otimes_{\sh{O}_X}\sh{G}\bigr).
  \end{aligned}
\]
\oldpage[0\textsubscript{III}]{58}
Se sigue inmediatamente de la definición del producto cup
(G, II, 6.6) que estos homomorfismos conmutan con la restricción de
$V$ a un subespacio abierto $V'$ de $V$.
Por otra parte, $\theta$ da un homomorfismo de anillos
\[
  \Gamma(V,\sh{O}_Y)
  \longrightarrow\Gamma(V,\psi_*(\sh{O}_X))
  =\Gamma(f^{-1}(V),\sh{O}_X),
\]
y, por tanto, un homomorfismo canónico de productos tensoriales
\[
  \begin{aligned}
  &\HH^p(f^{-1}(V),\sh{F})
    \otimes_{\Gamma(V,\sh{O}_Y)}
    \HH^q(f^{-1}(V),\sh{G})\\
  &\quad\longrightarrow
    \HH^p(f^{-1}(V),\sh{F})
    \otimes_{\Gamma(f^{-1}(V),\sh{O}_X)}
    \HH^q(f^{-1}(V),\sh{G}),
  \end{aligned}
\]
que es asimismo compatible con la restricción de $V$ a $V'$.
Por composición se obtiene así un homomorfismo de
$\Gamma(V,\sh{O}_Y)$-módulos.  Para los haces asociados a los
prehaces en cuestión, define un homomorfismo canónico,
funtorial en $\sh{F}$ y $\sh{G}$,
\[
  \RR^p f_*(\sh{F})\otimes_{\sh{O}_Y}\RR^q f_*(\sh{G})
  \longrightarrow
  \RR^{p+q}f_*
    \bigl(\sh{F}\otimes_{\sh{O}_X}\sh{G}\bigr).
  \tag{12.2.2.1}
  \label{0.12.2.2.1}
\]
Para $p=q=0$, este homomorfismo se reduce a
\sref[0\textsubscript{I}]{0.4.2.2.1}.
\end{env}

\begin{proposition}[12.2.3]
\label{0.12.2.3}
Para todo $\sh{O}_X$-módulo $\sh{F}$ y todo $\sh{O}_Y$-módulo
$\sh{L}$ localmente libre de rango finito, existen isomorfismos canónicos
funtoriales
\[
  \RR^p f_*(\sh{F})\otimes_{\sh{O}_Y}\sh{L}
  \simto
  \RR^p f_*
    \bigl(\sh{F}\otimes_{\sh{O}_X}f^*(\sh{L})\bigr).
  \tag{12.2.3.1}
  \label{0.12.2.3.1}
\]
\end{proposition}

\begin{proof}
El homomorfismo \sref{0.12.2.3.1} se obtiene componiendo el
siguiente caso particular de \sref{0.12.2.2.1},
\[
  \RR^p f_*(\sh{F})
    \otimes_{\sh{O}_Y}f_*\bigl(f^*(\sh{L})\bigr)
  \longrightarrow
  \RR^p f_*
    \bigl(\sh{F}\otimes_{\sh{O}_X}f^*(\sh{L})\bigr),
  \tag{12.2.3.2}
  \label{0.12.2.3.2}
\]
con el homomorfismo del primer miembro de
\sref{0.12.2.3.1} al de \sref{0.12.2.3.2} inducido por el
homomorfismo canónico \sref[0\textsubscript{I}]{0.4.4.3.2}.
Para comprobar que \sref{0.12.2.3.1} es un isomorfismo cuando $\sh{L}$ es
localmente libre, puede reducirse inmediatamente al caso
$\sh{L}=\sh{O}_Y$, pues la cuestión es local sobre $Y$ y los
funtores en cuestión son aditivos en $\sh{L}$.
En ese caso, en vista de la definición de \sref{0.12.2.2.1}, la
proposición se reduce a comprobar que el homomorfismo correspondiente
de prehaces es biyectivo; esto se sigue inmediatamente de
$f^*(\sh{O}_Y)=\sh{O}_X$.
\end{proof}

\begin{env}[12.2.4]
\label{0.12.2.4}
Sea $(Z,\sh{O}_Z)$ un tercer espacio anillado, y sea
$g:Y\to Z$ un morfismo de espacios anillados.
Sabemos (G, II, 7.1 y 3.1.1) que, para todo
$\sh{O}_X$-módulo inyectivo $\sh{G}$, $f_*(\sh{G})$ es un haz flácido de
grupos abelianos.  Por consiguiente, por \sref{0.12.2.1},
\[
  \RR^p g_*\bigl(f_*(\sh{G})\bigr)=0
  \qquad(p>0).
\]
Se sigue (T, 2.4.1) que la sucesión espectral de Leray para funtores
compuestos se aplica al compuesto $g_*\circ f_*$.  Existe una
sucesión espectral birregular cuyo término límite es el funtor
$\RR^\bullet h_*$, donde $h=g\circ f$, y cuyo término $E_2$ es
\[
  E_2^{pq}=\RR^p g_*\bigl(\RR^q f_*(\sh{F})\bigr).
  \tag{12.2.4.1}
  \label{0.12.2.4.1}
\]
\end{env}

\begin{env}[12.2.5]
\label{0.12.2.5}
Bajo las condiciones de \sref{0.12.2.4}, definiremos directamente
homomorfismos canónicos de $\sh{O}_Z$-módulos
\[
  \RR^n g_*\bigl(f_*(\sh{F})\bigr)
  \longrightarrow\RR^n h_*(\sh{F})
  \tag{12.2.5.1}
  \label{0.12.2.5.1}
\]
y
\[
  \RR^n h_*(\sh{F})
  \longrightarrow g_*\bigl(\RR^n f_*(\sh{F})\bigr),
  \tag{12.2.5.2}
  \label{0.12.2.5.2}
\]
\oldpage[0\textsubscript{III}]{59}
que pueden identificarse con los ``homomorfismos de borde'' de la sucesión
espectral de Leray (cf.\ \sref{0.12.1.7}).
Basta trabajar con los prehaces a los que están asociados los haces de
imágenes directas superiores \sref{0.12.2.1}.
Para ello, sea $W$ un subconjunto abierto cualquiera de $Z$, y consideremos
su imagen inversa $g^{-1}(W)$ en $Y$.  Existe un
dihomomorfismo canónico
\[
  \begin{aligned}
  \HH^n\bigl(g^{-1}(W),f_*(\sh{F})\bigr)
  &\longrightarrow
  \HH^n\bigl(f^{-1}(g^{-1}(W)),
    f^*(f_*(\sh{F}))\bigr),
  \end{aligned}
  \tag{12.2.5.3}
  \label{0.12.2.5.3}
\]
siendo los anillos correspondientes
$\Gamma(g^{-1}(W),\sh{O}_Y)$ y
$\Gamma(h^{-1}(W),\sh{O}_X)$.
Por otra parte, el homomorfismo canónico
\sref[0\textsubscript{I}]{0.4.4.3.3} da, por funtorialidad,
homomorfismos canónicos
\[
  \HH^n\bigl(h^{-1}(W),f^*(f_*(\sh{F}))\bigr)
  \longrightarrow
  \HH^n\bigl(h^{-1}(W),\sh{F}\bigr),
  \tag{12.2.5.4}
  \label{0.12.2.5.4}
\]
que son homomorfismos de
$\Gamma(h^{-1}(W),\sh{O}_X)$-módulos.
Teniendo en cuenta el homomorfismo de anillos
\[
  \Gamma(W,\sh{O}_Z)
  \longrightarrow\Gamma(g^{-1}(W),\sh{O}_Y),
\]
vemos que, componiendo \sref{0.12.2.5.3} con
\sref{0.12.2.5.4}, se obtiene un homomorfismo de prehaces y, por tanto, el
homomorfismo de haces \sref{0.12.2.5.1}.

La definición de \sref{0.12.2.5.2} es aún más sencilla.
Por definición, $\RR^n h_*(\sh{F})$ está asociado al prehaz
\[
  W\longmapsto
  \HH^n\bigl(f^{-1}(g^{-1}(W)),\sh{F}\bigr),
\]
mientras que $\RR^n f_*(\sh{F})$ está asociado al prehaz
\[
  V\longmapsto\HH^n(f^{-1}(V),\sh{F}).
\]
Existe, por tanto, un homomorfismo canónico
\[
  \HH^n\bigl(f^{-1}(g^{-1}(W)),\sh{F}\bigr)
  \longrightarrow
  \Gamma\bigl(g^{-1}(W),\RR^n f_*(\sh{F})\bigr),
\]
y es inmediato que estos homomorfismos definen un homomorfismo de
prehaces, que a su vez define \sref{0.12.2.5.2}.
\end{env}

\begin{env}[12.2.6]
\label{0.12.2.6}
Bajo las hipótesis de \sref{0.12.2.4}, sean $\sh{F}$, $\sh{G}$
y $\sh{K}$ tres $\sh{O}_X$-módulos, y sea
\[
  u:\sh{F}\otimes_{\sh{O}_X}\sh{G}\longrightarrow\sh{K}
\]
un $\sh{O}_X$-homomorfismo.
Entonces los diagramas siguientes son conmutativos:
\[
  \xymatrix@C=10pt@R=24pt{
    \RR^p g_*(f_*(\sh{F}))
      \otimes_{\sh{O}_Z}\RR^q g_*(f_*(\sh{G}))
      \ar[r]\ar[d]
      & \RR^{p+q}g_*(f_*(\sh{K}))\ar[d]\\
    \RR^p h_*(\sh{F})
      \otimes_{\sh{O}_Z}\RR^q h_*(\sh{G})
      \ar[r]
      & \RR^{p+q}h_*(\sh{K}),
  }
  \tag{12.2.6.1}
  \label{0.12.2.6.1}
\]
y
\[
  \xymatrix@C=10pt@R=24pt{
    \RR^p h_*(\sh{F})
      \otimes_{\sh{O}_Z}\RR^q h_*(\sh{G})
      \ar[r]\ar[d]
      & \RR^{p+q}h_*(\sh{K})\ar[d]\\
    g_*(\RR^p f_*(\sh{F}))
      \otimes_{\sh{O}_Z}g_*(\RR^q f_*(\sh{G}))
      \ar[r]
      & g_*(\RR^{p+q}f_*(\sh{K})).
  }
  \tag{12.2.6.2}
  \label{0.12.2.6.2}
\]
\oldpage[0\textsubscript{III}]{60}
Aquí las flechas horizontales provienen de \sref{0.12.2.2.1} (la última
combinada con \sref[0\textsubscript{I}]{0.4.2.2.1}), y las flechas
verticales de \sref{0.12.2.5.1} y \sref{0.12.2.5.2}, respectivamente.

En efecto, basta verificar la aserción para los homomorfismos de prehaces
correspondientes.  Volviendo a las definiciones de
\sref{0.12.2.2} y \sref{0.12.2.5}, se reduce inmediatamente, para
\sref{0.12.2.6.1}, a los diagramas conmutativos
\sref{0.12.1.5.2}; la verificación es aún más sencilla para
\sref{0.12.2.6.2}.
\end{env}

\subsection{Suplementos sobre los funtores Ext de haces}
\label{subsection:0.12.3}

\begin{env}[12.3.1]
\label{0.12.3.1}
Sea $(X,\sh{O}_X)$ un espacio anillado.
No recordaremos la definición ni las propiedades principales de los
bifuntores
\[
  \Ext^p_{\sh{O}_X}(X;\sh{F},\sh{G}),
  \qquad
  \shExt^p_{\sh{O}_X}(\sh{F},\sh{G}),
\]
el primero con valores en la categoría de
$\Gamma(X,\sh{O}_X)$-módulos y el segundo en la categoría de
$\sh{O}_X$-módulos, ni la sucesión espectral birregular
$E(\sh{F},\sh{G})$ que los relaciona (T, 4.2 y G, II, 7.3).
\end{env}

\begin{env}[12.3.2]
\label{0.12.3.2}
Al igual que en (M, XIV, 1), se define la noción de una
\emph{extensión} de un $\sh{O}_X$-módulo $\sh{F}$ por un
$\sh{O}_X$-módulo $\sh{G}$ y la ley de composición sobre las clases de
extensiones equivalentes: los argumentos utilizados para los módulos se
trasladan manifiestamente a una categoría abeliana arbitraria.
La segunda demostración de (M, XIV, 1.1), que solo utiliza la existencia de
inmersiones en objetos inyectivos, sigue siendo, por tanto, válida en la
categoría de $\sh{O}_X$-módulos.  Muestra que
\[
  \Ext^1_{\sh{O}_X}(X;\sh{F},\sh{G})
\]
se identifica canónicamente con el grupo abeliano de clases de
extensiones de $\sh{F}$ por $\sh{G}$.
\end{env}

\begin{proposition}[12.3.3]
\label{0.12.3.3}
Sea $(X,\sh{O}_X)$ un espacio anillado cuyo haz de anillos
$\sh{O}_X$ es coherente.
Entonces, para todo par de $\sh{O}_X$-módulos coherentes $\sh{F}$ y
$\sh{G}$ y todo $p\geq0$,
\[
  \shExt^p_{\sh{O}_X}(\sh{F},\sh{G})
\]
es un $\sh{O}_X$-módulo coherente.
\end{proposition}

\begin{proof}
Los $\shExt^p_{\sh{O}_X}(\sh{F},\sh{G})$ forman un funtor cohomológico
contravariante en $\sh{F}$.
Como $\sh{F}$ es coherente, para todo $p$ y todo punto $x\in X$
existen un entorno abierto $U$ de $x$ en $X$ y una sucesión exacta
de $(\sh{O}_X|U)$-módulos
\[
  0\longrightarrow\sh{R}\longrightarrow\sh{L}_{p-1}
  \longrightarrow\cdots\longrightarrow\sh{L}_0
  \longrightarrow\sh{F}|U\longrightarrow0,
\]
donde cada $\sh{L}_i$ ($0\leq i\leq p-1$) es isomorfo a
$\sh{O}_X^{n_i}|U$, y $\sh{R}$ es coherente.
Esto se sigue por inducción sobre $p$ de
\sref[0\textsubscript{I}]{0.5.3.2} y
\sref[0\textsubscript{I}]{0.5.3.4}, pues $\sh{O}_X$ es coherente.

Además, para $p\geq1$,
\[
  \shExt^p_{\sh{O}_X|U}(\sh{L}|U,\sh{G}|U)=0
\]
para todo $\sh{O}_X$-módulo $\sh{L}$ tal que $\sh{L}|U$ sea
isomorfo a $\sh{O}_X^n|U$ (T, 4.2.3).
El argumento de (M, V, 7.2) se aplica, por tanto, al funtor cohomológico
contravariante
\[
  \sh{F}\longmapsto
  \shHom_{\sh{O}_X|U}(\sh{F}|U,\sh{G}|U)
\]
y da una sucesión exacta
\[
  \begin{aligned}
  \shHom_{\sh{O}_X|U}(\sh{L}_{p-1},\sh{G}|U)
  &\longrightarrow
  \shHom_{\sh{O}_X|U}(\sh{R},\sh{G}|U)\\
  &\longrightarrow
  \shExt^p_{\sh{O}_X|U}(\sh{F}|U,\sh{G}|U)
  \longrightarrow0.
  \end{aligned}
\]
Los dos primeros términos son $(\sh{O}_X|U)$-módulos coherentes
\sref[0\textsubscript{I}]{0.5.3.5}; por tanto, también lo es el tercero
\sref[0\textsubscript{I}]{0.5.3.4}.
\end{proof}

\oldpage[0\textsubscript{III}]{61}
\begin{proposition}[12.3.4]
\label{0.12.3.4}
Sea $f:X\to Y$ un morfismo plano de espacios anillados, y sean
$\sh{F}$ y $\sh{G}$ dos $\sh{O}_Y$-módulos.
\begin{enumerate}
  \item[\rm(i)] Existe un homomorfismo de bifuntores cohomológicos
  \[
    f^*\bigl(\shExt^\bullet_{\sh{O}_Y}(\sh{F},\sh{G})\bigr)
    \longrightarrow
    \shExt^\bullet_{\sh{O}_X}
      \bigl(f^*(\sh{F}),f^*(\sh{G})\bigr),
    \tag{12.3.4.1}
    \label{0.12.3.4.1}
  \]
  que en grado $0$ se reduce al homomorfismo canónico
  \sref[0\textsubscript{I}]{0.4.4.6}.

  \item[\rm(ii)] Existe un morfismo canónico de sucesiones espectrales
  \[
    E(\sh{F},\sh{G})
    \longrightarrow E\bigl(f^*(\sh{F}),f^*(\sh{G})\bigr),
    \tag{12.3.4.2}
    \label{0.12.3.4.2}
  \]
  cuya aplicación sobre los términos $E_2$ viene dada por los homomorfismos
  \[
    \begin{aligned}
    \HH^p\bigl(Y,\shExt^q_{\sh{O}_Y}(\sh{F},\sh{G})\bigr)
    \longrightarrow
    \HH^p\bigl(X,\shExt^q_{\sh{O}_X}
      (f^*(\sh{F}),f^*(\sh{G}))\bigr),
    \end{aligned}
    \tag{12.3.4.3}
    \label{0.12.3.4.3}
  \]
  deducidos de \sref{0.12.3.4.1} y \sref{0.12.1.3.1}.
\end{enumerate}
\end{proposition}

\begin{proof}
\begin{enumerate}
  \item[\rm(i)]
  Como $f^*$ es exacto sobre la categoría de $\sh{O}_Y$-módulos
  \sref[0\textsubscript{I}]{0.6.7.2}, los funtores
  \[
    \sh{G}\longmapsto
      f^*\bigl(\shHom_{\sh{O}_Y}(\sh{F},\sh{G})\bigr),
    \qquad
    \sh{G}\longmapsto
      \shHom_{\sh{O}_X}\bigl(f^*(\sh{F}),f^*(\sh{G})\bigr)
  \]
  son exactos por la izquierda.  El homomorfismo canónico
  \sref[0\textsubscript{I}]{0.4.4.6} induce, por tanto, canónicamente un
  homomorfismo entre sus funtores derivados.

  Para calcular estos funtores derivados, tomemos una resolución inyectiva
  $\sh{L}^\bullet=(\sh{L}^i)$ de $\sh{G}$.  Se tienen entonces
  homomorfismos
  \[
    \begin{aligned}
    \sh{H}^p\!\left(
      f^*\bigl(\shHom_{\sh{O}_Y}
        (\sh{F},\sh{L}^\bullet)\bigr)\right)
    \longrightarrow
    \sh{H}^p\!\left(
      \shHom_{\sh{O}_X}
        (f^*(\sh{F}),f^*(\sh{L}^\bullet))\right)
    \end{aligned}
  \]
  entre los haces de cohomología de complejos.
  Por la exactitud de $f^*$, el primer miembro es
  \[
    f^*\bigl(\shExt^p_{\sh{O}_Y}(\sh{F},\sh{G})\bigr).
  \]
  La exactitud de $f^*$ implica también que
  $f^*(\sh{L}^\bullet)$ es una resolución de $f^*(\sh{G})$.
  Si ${\sh{L}'}^\bullet=({\sh{L}'}^i)$ es una resolución inyectiva
  de $f^*(\sh{G})$ en la categoría de $\sh{O}_X$-módulos, existe un
  homomorfismo de complejos
  \[
    f^*(\sh{L}^\bullet)\longrightarrow{\sh{L}'}^\bullet,
  \]
  determinado salvo homotopía, y por tanto un homomorfismo bien definido en
  cohomología.  Su compuesto con el homomorfismo precedente es
  \sref{0.12.3.4.1}.

  \item[\rm(ii)]
  Con la notación precedente, existe un homomorfismo de complejos de
  $\sh{O}_X$-módulos
  \[
    f^*\bigl(\shHom_{\sh{O}_Y}(\sh{F},\sh{L}^\bullet)\bigr)
    \longrightarrow
    \shHom_{\sh{O}_X}
      \bigl(f^*(\sh{F}),{\sh{L}'}^\bullet\bigr).
  \]
  Sea $\sh{M}^{\bullet\bullet}$ una resolución inyectiva de
  Cartan--Eilenberg de
  $\shHom_{\sh{O}_Y}(\sh{F},\sh{L}^\bullet)$ en la categoría de
  $\sh{O}_Y$-módulos.  Como $f^*$ es exacto,
  $f^*(\sh{M}^{\bullet\bullet})$ es una resolución de Cartan--Eilenberg de
  $f^*(\shHom_{\sh{O}_Y}(\sh{F},\sh{L}^\bullet))$.
  Si ${\sh{M}'}^{\bullet\bullet}$ es una resolución inyectiva de
  Cartan--Eilenberg de
  $\shHom_{\sh{O}_X}(f^*(\sh{F}),{\sh{L}'}^\bullet)$, existe
  \sref{0.11.4.2} un homomorfismo, determinado salvo homotopía,
  \[
    f^*(\sh{M}^{\bullet\bullet})
    \longrightarrow{\sh{M}'}^{\bullet\bullet},
  \]
  compatible con el homomorfismo precedente; en otras palabras, existe
  un $f$-morfismo
  $\sh{M}^{\bullet\bullet}\to{\sh{M}'}^{\bullet\bullet}$ de
  de bicomplejos de haces.
  Induce un dihomomorfismo
  \[
    \Gamma(Y,\sh{M}^{\bullet\bullet})
    \longrightarrow\Gamma(X,{\sh{M}'}^{\bullet\bullet})
  \]
  de bicomplejos de módulos, determinado salvo homotopía, y por tanto
  un morfismo bien definido de sucesiones espectrales
  \sref{0.11.3.2}.  Este es precisamente el morfismo deseado
  \sref{0.12.3.4.2}, y la descripción
  \sref{0.12.3.4.3} se sigue inmediatamente de las definiciones.
\end{enumerate}
\end{proof}

\begin{proposition}[12.3.5]
\label{0.12.3.5}
Bajo las hipótesis de \sref{0.12.3.4}, supongamos además que el
haz de anillos $\sh{O}_Y$ es coherente.
Entonces, para todo $\sh{O}_Y$-módulo coherente $\sh{F}$, los
homomorfismos canónicos \sref{0.12.3.4.1} son biyectivos.
\end{proposition}

\begin{proof}
\oldpage[0\textsubscript{III}]{62}
Como la cuestión es local sobre $Y$, podemos suponer que existe una
sucesión exacta
\[
  0\longrightarrow\sh{R}\longrightarrow\sh{O}_Y^n
  \longrightarrow\sh{F}\longrightarrow0.
\]
Entonces $\sh{R}$ es también un $\sh{O}_Y$-módulo coherente
\sref[0\textsubscript{I}]{0.5.3.4}.
Para demostrar que los homomorfismos
\[
  f^*\bigl(\shExt^p_{\sh{O}_Y}(\sh{F},\sh{G})\bigr)
  \longrightarrow
  \shExt^p_{\sh{O}_X}
    \bigl(f^*(\sh{F}),f^*(\sh{G})\bigr)
\]
son biyectivos, razonamos por inducción sobre $p$.
Para $p=0$, la aserción es
\sref[0\textsubscript{I}]{0.6.7.6.1}.
Existe un diagrama conmutativo con filas exactas
\[
{\tiny
  \xymatrix@C=2pt@R=24pt{
  f^*\!\left(\shExt^{p-1}_{\sh{O}_Y}(\sh{O}_Y^n,\sh{G})\right)
    \ar[r]\ar[d]
  & f^*\!\left(\shExt^{p-1}_{\sh{O}_Y}(\sh{R},\sh{G})\right)
    \ar[r]^\partial\ar[d]
  & f^*\!\left(\shExt^p_{\sh{O}_Y}(\sh{F},\sh{G})\right)
    \ar[r]\ar[d]
  & f^*\!\left(\shExt^p_{\sh{O}_Y}(\sh{O}_Y^n,\sh{G})\right)
    \ar[d]\\
  \shExt^{p-1}_{\sh{O}_X}(\sh{O}_X^n,f^*(\sh{G}))
    \ar[r]
  & \shExt^{p-1}_{\sh{O}_X}(f^*(\sh{R}),f^*(\sh{G}))
    \ar[r]^\partial
  & \shExt^p_{\sh{O}_X}(f^*(\sh{F}),f^*(\sh{G}))
    \ar[r]
  & \shExt^p_{\sh{O}_X}(\sh{O}_X^n,f^*(\sh{G})).
  }}
\]
Aquí hemos utilizado $f^*(\sh{O}_Y)=\sh{O}_X$; como $f^*$ es exacto,
ambas filas son exactas.
Además,
\[
  \shExt^p_{\sh{O}_Y}(\sh{O}_Y^n,\sh{G})=0,
  \qquad
  \shExt^p_{\sh{O}_X}(\sh{O}_X^n,f^*(\sh{G}))=0
  \qquad(p>0)
\]
(T, 4.2.3).
Por la hipótesis de inducción, las dos primeras flechas verticales del
diagrama precedente son isomorfismos, mientras que los términos de la
derecha son nulos.  Se sigue que
\[
  f^*\bigl(\shExt^p_{\sh{O}_Y}(\sh{F},\sh{G})\bigr)
  \longrightarrow
  \shExt^p_{\sh{O}_X}
    \bigl(f^*(\sh{F}),f^*(\sh{G})\bigr)
\]
es un isomorfismo.
\end{proof}

\subsection{Hipercohomología del funtor imagen directa}
\label{subsection:0.12.4}

\begin{env}[12.4.1]
\label{0.12.4.1}
Sean $(X,\sh{O}_X)$ y $(Y,\sh{O}_Y)$ dos espacios anillados, y sea
$f:X\longrightarrow Y$ un morfismo de espacios anillados.
Se puede formar la hipercohomología de $f_*$ respecto de un complejo
arbitrario de $\sh{O}_X$-módulos
$\sh{K}^\bullet=(\sh{K}^i)_{i\in\bb{Z}}$
\sref{0.11.4.4}, pues los límites inductivos filtrantes existen y son exactos
en la categoría abeliana de $\sh{O}_Y$-módulos (T, 3.1.1).
Los $\sh{O}_Y$-módulos de hipercohomología
$\RR^p f_*(\sh{K}^\bullet)$ se denotarán también
$\sh{H}^p(f,\sh{K}^\bullet)$, o simplemente
$\sh{H}^p_f(\sh{K}^\bullet)$.
Recordemos que $\sh{H}^\bullet(f,\sh{K}^\bullet)$ es la cohomología del
bicomplejo de $\sh{O}_Y$-módulos $f_*(\sh{L}^{\bullet\bullet})$,
donde $\sh{L}^{\bullet\bullet}$ es una resolución inyectiva de
Cartan--Eilenberg de $\sh{K}^\bullet$ en la categoría de
$\sh{O}_X$-módulos.
Es el término límite de dos sucesiones espectrales cuyos términos $E_2$
vienen dados por
\begin{equation}
\label{0.12.4.1.1}
\tag{12.4.1.1}
  {}'E_2^{pq}
  =\sh{H}^p\bigl(\sh{H}^q(f,\sh{K}^\bullet)\bigr)
\end{equation}
y
\begin{equation}
\label{0.12.4.1.2}
\tag{12.4.1.2}
  {}''E_2^{pq}
  =\sh{H}^p\bigl(f,\sh{H}^q(\sh{K}^\bullet)\bigr)
  \quad
  \bigl(=\RR^p f_*(\sh{H}^q(\sh{K}^\bullet))\bigr).
\end{equation}
En estas fórmulas hemos adoptado la notación general $T(A^\bullet)$
para el transformado de un complejo por un funtor \sref{0.11.2.1}, y, para
un $\sh{O}_X$-módulo $\sh{F}$, escribimos
$\sh{H}^p(f,\sh{F})$ en lugar de $\RR^p f_*(\sh{F})$.
Recordemos también que la primera sucesión espectral es siempre regular.
Ambas sucesiones espectrales son birregulares cuando $\sh{K}^\bullet$ está
acotado inferiormente,
\oldpage[0\textsubscript{III}]{63}
o cuando existe un entero $m$ tal que todo $\sh{O}_X$-módulo
admite una resolución flácida de longitud $\leq m$ \sref{0.11.4.4}.
\end{env}

\begin{env}[12.4.2]
\label{0.12.4.2}
Denotaremos análogamente por $\bb{H}^\bullet(X,\sh{K}^\bullet)$ la
hipercohomología del funtor $\Gamma$ respecto de un complejo
$\sh{K}^\bullet$ de $\sh{O}_X$-módulos.
Así, los $\bb{H}^p(X,\sh{K}^\bullet)$ son
$\Gamma(X,\sh{O}_X)$-módulos.
Además, $\bb{H}^\bullet(X,\sh{K}^\bullet)$ puede considerarse como un
caso particular de $\sh{H}^\bullet(f,\sh{K}^\bullet)$, donde $f$ es un
morfismo de $(X,\sh{O}_X)$ a un espacio anillado constituido por un punto
dotado del anillo $\Gamma(X,\sh{O}_X)$.

Para todo subconjunto abierto $V$ de $X$, escribiremos
$\bb{H}^\bullet(V,\sh{K}^\bullet)$ en lugar de
$\bb{H}^\bullet(V,\sh{K}^\bullet|V)$.
\end{env}

\begin{proposition}[12.4.3]
\label{0.12.4.3}
Para todo entero $p\in\bb{Z}$, el $\sh{O}_Y$-módulo
$\sh{H}^p(f,\sh{K}^\bullet)$ es canónicamente isomorfo al haz
asociado al prehaz
\[
  U\longmapsto
  \bb{H}^p\bigl(f^{-1}(U),\sh{K}^\bullet\bigr)
\]
sobre $Y$.
\end{proposition}

\begin{proof}
Con la notación de \sref{0.12.4.1}, el haz de cohomología
$\sh{H}^p(f_*(\sh{L}^{\bullet\bullet}))$ está asociado al
prehaz
\[
  U\longmapsto
  \HH^p\bigl(\Gamma(U,f_*(\sh{L}^{\bullet\bullet}))\bigr)
  =
  \HH^p\bigl(\Gamma(f^{-1}(U),\sh{L}^{\bullet\bullet})\bigr).
\]
Pero está claro que
$\sh{L}^{\bullet\bullet}|f^{-1}(U)$ es una resolución inyectiva de
Cartan--Eilenberg de $\sh{K}^\bullet|f^{-1}(U)$ (T, 3.1.3).
Por consiguiente,
\[
  \HH^p\bigl(\Gamma(f^{-1}(U),\sh{L}^{\bullet\bullet})\bigr)
  =
  \bb{H}^p\bigl(f^{-1}(U),\sh{K}^\bullet\bigr)
\]
por definición.
\end{proof}

\begin{proposition}[12.4.4]
\label{0.12.4.4}
La hipercohomología $\sh{H}^\bullet(f,\sh{K}^\bullet)$ es un
funtor cohomológico de $\sh{K}^\bullet$ en cada uno de los casos
siguientes:
\begin{enumerate}
  \item[\rm{a)}] $\sh{K}^\bullet$ varía en la categoría de complejos
    acotados inferiormente;
  \item[\rm{b)}] existe un entero $m$ tal que todo
    $\sh{O}_X$-módulo admite una resolución flácida de longitud $\leq m$;
  \item[\rm{c)}] $X$ es un espacio noetheriano.
\end{enumerate}
\end{proposition}

\begin{proof}
Los casos~a) y~b) son casos particulares de \sref{0.11.5.4}.
El caso~c), por otra parte, se sigue de \sref{0.11.5.2}, pues en
este caso el funtor $f_*$ conmuta con los límites inductivos
(G, II, 3.10.1).
\end{proof}

\begin{env}[12.4.5]
\label{0.12.4.5}
Sea ahora $\mathfrak{U}=(U_\alpha)$ un recubrimiento abierto de $X$.
Para todo complejo de \emph{prehaces}
$\sh{K}^\bullet=(\sh{K}^j)$ sobre $X$, consideremos el bicomplejo
$\check{\mathrm{C}}^\bullet(\mathfrak{U},\sh{K}^\bullet)$ cuya
componente de bigrado $(i,j)$ es
$\check{\mathrm{C}}^i(\mathfrak{U},\sh{K}^j)$, el grupo de
$i$-cocadenas alternadas del nervio de $\mathfrak{U}$ con valores
en $\sh{K}^j$ (G, II, 5.1).
Llamamos a la cohomología de este bicomplejo la
\emph{hipercohomología del recubrimiento $\mathfrak{U}$ con
coeficientes en $\sh{K}^\bullet$}, y la denotamos por
\[
  \bb{H}^\bullet(\mathfrak{U},\sh{K}^\bullet)
  =
  \HH^\bullet
  \bigl(\check{\mathrm{C}}^\bullet
    (\mathfrak{U},\sh{K}^\bullet)\bigr).
\]
La sucesión espectral de Leray de un recubrimiento
(T, 3.8.1 y G, II, 5.9.1) se generaliza a la hipercohomología como sigue.
\end{env}

\begin{proposition}[12.4.6]
\label{0.12.4.6}
Sea $\sh{K}^\bullet=(\sh{K}^i)$ un complejo de
$\sh{O}_X$-módulos.
Existe un funtor espectral regular de $\sh{K}^\bullet$ con
término límite $\bb{H}^\bullet(X,\sh{K}^\bullet)$ y término $E_2$
\begin{equation}
\label{0.12.4.6.1}
\tag{12.4.6.1}
  E_2^{pq}
  =
  \CHH^p\bigl(\mathfrak{U},h^q(\sh{K}^\bullet)\bigr),
\end{equation}
donde $h^q(\sh{K}^\bullet)$ denota el complejo de prehaces
\[
  V\longmapsto\HH^q(V,\sh{K}^\bullet)
\]
sobre $X$.
Esta sucesión espectral es birregular si $\sh{K}^\bullet$ está acotado
inferiormente.
\end{proposition}

\begin{proof}
Sea $\sh{L}^{\bullet\bullet}$ una resolución inyectiva de
Cartan--Eilenberg de $\sh{K}^\bullet$, y consideremos el tricomplejo
\[
  \check{\mathrm{C}}^\bullet
  (\mathfrak{U},\sh{L}^{\bullet\bullet})
  =
  \bigl(\check{\mathrm{C}}^i
    (\mathfrak{U},\sh{L}^{jk})\bigr).
\]
Consideremos primero este tricomplejo como un bicomplejo en los grados $i$ y
$j+k$.
Como $i$ solo toma valores no negativos, la segunda sucesión espectral
de este bicomplejo es regular \sref{0.11.3.3} y degenerada, porque
\[
  \CHH^q(\mathfrak{U},\sh{L}^{jk})=0
  \qquad(q>0),
\]
siendo los $\sh{O}_X$-módulos $\sh{L}^{jk}$ haces flácidos
\oldpage[0\textsubscript{III}]{64}
(G, II, 5.2.3).
Se sigue que existe un isomorfismo canónico
\[
  \HH^n\bigl(
    \check{\mathrm{C}}^\bullet
      (\mathfrak{U},\sh{L}^{\bullet\bullet})
  \bigr)
  \isoto
  \HH^n\bigl(\Gamma(X,\sh{L}^{\bullet\bullet})\bigr)
\]
\sref{0.11.1.6}, en virtud de (G, II, 5.2.2), y por tanto, por la
definición \sref{0.12.4.2}, un isomorfismo
\[
  \HH^n\bigl(
    \check{\mathrm{C}}^\bullet
      (\mathfrak{U},\sh{L}^{\bullet\bullet})
  \bigr)
  \isoto
  \bb{H}^n(X,\sh{K}^\bullet).
\]

Consideremos ahora
$\check{\mathrm{C}}^\bullet
(\mathfrak{U},\sh{L}^{\bullet\bullet})$
como un bicomplejo en los grados $i+j$ y $k$.
Como $k$ solo toma valores no negativos, la primera sucesión espectral
de este bicomplejo es siempre regular.
Es birregular si $\sh{L}^{jk}=0$ para $j<j_0$, es decir, cuando
$\sh{K}^\bullet$ está acotado inferiormente \sref{0.11.3.3}.
Esta es la sucesión espectral deseada.
En efecto, para todo $j$, $\sh{L}^{j,\bullet}$ es una resolución
inyectiva de $\sh{K}^j$; por consiguiente, cuando $i$ varía,
\[
  \HH^q\bigl(
    \check{\mathrm{C}}^i
      (\mathfrak{U},\sh{L}^{j,\bullet})
  \bigr)
\]
es precisamente el complejo de cocadenas
$\check{\mathrm{C}}^i(\mathfrak{U},h^q(\sh{K}^j))$.
Esto demuestra la proposición.
\end{proof}

\begin{corollary}[12.4.7]
\label{0.12.4.7}
Si, para todo símplex $\sigma$ del nervio de $\mathfrak{U}$ y
todo entero $i$, se tiene
\[
  \HH^q(U_\sigma,\sh{K}^i)=0
  \qquad(q>0),
\]
entonces existe un isomorfismo canónico
\begin{equation}
\label{0.12.4.7.1}
\tag{12.4.7.1}
  \bb{H}^\bullet(\mathfrak{U},\sh{K}^\bullet)
  \isoto
  \bb{H}^\bullet(X,\sh{K}^\bullet).
\end{equation}
\end{corollary}

\begin{proof}
La hipótesis implica que
$\check{\mathrm{C}}^\bullet
(\mathfrak{U},h^q(\sh{K}^j))=0$ para $q>0$, y por tanto que
$E_2^{pq}=0$ para $q>0$.
Como la sucesión espectral \sref{0.12.4.6.1} es degenerada y
regular, la conclusión se sigue de la definición
\sref{0.12.4.5} de
$\bb{H}^\bullet(\mathfrak{U},\sh{K}^\bullet)$ y de
\sref{0.11.1.6}.
\end{proof}

\begin{env}[12.4.8]
\label{0.12.4.8}
Sea $(X',\sh{O}_{X'})$ un segundo espacio anillado, y sea
$f=(\psi,\theta)$ un morfismo de $X'$ en $X$.
Por el mismo método que en \sref{0.12.1.3} y \sref{0.12.1.4}, se
define un dihomomorfismo para la hipercohomología de un complejo
$\sh{K}^\bullet$ de $\sh{O}_X$-módulos:
\begin{equation}
\label{0.12.4.8.1}
\tag{12.4.8.1}
  \bb{H}^p(X,\sh{K}^\bullet)
  \longrightarrow
  \bb{H}^p(X',f^*(\sh{K}^\bullet)).
\end{equation}
Partiendo de una resolución inyectiva de Cartan--Eilenberg
$\sh{L}^{\bullet\bullet}$ de $\sh{K}^\bullet$, observemos que, como
$\psi^*$ es exacto, $\psi^*(\sh{L}^{\bullet\bullet})$ es una
resolución de Cartan--Eilenberg de $\psi^*(\sh{K}^\bullet)$ en la
categoría de $\psi^*(\sh{O}_X)$-módulos.
Existe, por tanto, un morfismo
\[
  \psi^*(\sh{L}^{\bullet\bullet})
  \longrightarrow
  {\sh{L}'}^{\bullet\bullet},
\]
donde ${\sh{L}'}^{\bullet\bullet}$ es una resolución inyectiva de
Cartan--Eilenberg de $\psi^*(\sh{K}^\bullet)$.
Induce un morfismo en cohomología
\[
  \bb{H}^\bullet(X,\sh{K}^\bullet)
  \longrightarrow
  \bb{H}^\bullet
  (X',\psi^*(\sh{K}^\bullet)).
\]
Componiéndolo con el morfismo inducido por funtorialidad a partir de
$\psi^*(\sh{K}^\bullet)\longrightarrow f^*(\sh{K}^\bullet)$ se obtiene
el morfismo deseado \sref{0.12.4.8.1}.

Partiendo de \sref{0.12.4.8.1} y \sref{0.12.4.3}, y razonando como
en \sref{0.12.2.5}, se pueden definir entonces, para dos morfismos de
espacios anillados
\[
  f:X\longrightarrow Y,
  \qquad
  g:Y\longrightarrow Z,
\]
con $h=g\circ f$, homomorfismos para la hipercohomología de un complejo
$\sh{K}^\bullet$ de $\sh{O}_X$-módulos:
\begin{equation}
\label{0.12.4.8.2}
\tag{12.4.8.2}
  \sh{H}^n(g,f_*(\sh{K}^\bullet))
  \longrightarrow
  \sh{H}^n(h,\sh{K}^\bullet),
\end{equation}
y
\begin{equation}
\label{0.12.4.8.3}
\tag{12.4.8.3}
  \sh{H}^n(h,\sh{K}^\bullet)
  \longrightarrow
  g_*\bigl(\sh{H}^n(f,\sh{K}^\bullet)\bigr).
\end{equation}
Dejamos al lector los detalles de las definiciones.
\end{env}

\section{Límites proyectivos en álgebra homológica}
\label{section:0.13}

\subsection{La condición de Mittag-Leffler}
\label{subsection:0.13.1}

\begin{env}[13.1.1]
\label{0.13.1.1}
Sea $\cat{C}$ una categoría abeliana en la que existan los productos
infinitos (axioma AB~3* de T, 1.5).
Entonces existe el ínfimo de una familia de subobjetos de un objeto de
$\cat{C}$, y todo sistema proyectivo de objetos de $\cat{C}$ posee un
límite proyectivo; este último es un funtor exacto por la izquierda del
sistema proyectivo en cuestión (T, 1.8).
\oldpage[0\textsubscript{III}]{65}
Sea $(A_\alpha,f_{\alpha\beta})$ un sistema proyectivo de objetos de
$\cat{C}$ cuyo conjunto de índices $\mathrm{I}$ está filtrado por la derecha.
Pongamos
\[
  A=\varprojlim A_\alpha,
\]
y, para cada $\alpha\in\mathrm{I}$, sea
$f_\alpha:A\longrightarrow A_\alpha$ el morfismo canónico.
Para todo $\alpha\in\mathrm{I}$, los subobjetos
$f_{\alpha\beta}(A_\beta)$, para $\alpha\leq\beta$, forman una familia
filtrada decreciente de subobjetos de $A_\alpha$.
El subobjeto
\[
  A'_\alpha
  =
  \inf_{\beta\geq\alpha}f_{\alpha\beta}(A_\beta)
\]
se llama el subobjeto de las \emph{imágenes universales} en $A_\alpha$.
Está claro que
\[
  f_\alpha(A)\subset A'_\alpha,
  \qquad
  f_{\alpha\beta}(A'_\beta)\subset A'_\alpha
  \quad(\alpha\leq\beta).
\]
Así,
$(A'_\alpha,f_{\alpha\beta}|A'_\beta)$ es un sistema proyectivo y
\[
  A=\varprojlim A'_\alpha.
\]
\end{env}

\begin{env}[13.1.2]
\label{0.13.1.2}
Para un sistema proyectivo $(A_\alpha,f_{\alpha\beta})$ en $\cat{C}$,
se llama \emph{condición de Mittag-Leffler} a la siguiente:
\begin{itemize}
  \item[\textbf{(ML)}]
  \emph{Para todo índice $\alpha$, existe un $\beta\geq\alpha$ tal
  que, para todo $\gamma\geq\beta$,}
  \[
    f_{\alpha\gamma}(A_\gamma)
    =
    f_{\alpha\beta}(A_\beta).
  \]
\end{itemize}

Está claro que la condición (ML) se cumple si todos los
$f_{\alpha\beta}$ son epimorfismos.
Recíprocamente, supongamos que se cumple (ML), y para cada
$\alpha\in\mathrm{I}$ sea $A'_\alpha$ el subobjeto de las imágenes
universales en $A_\alpha$.
Entonces, para $\alpha\leq\beta$, la restricción de
$f_{\alpha\beta}$ a $A'_\beta$ es un \emph{epimorfismo}
\[
  A'_\beta\longrightarrow A'_\alpha.
\]
En efecto, elijamos $\gamma\geq\beta$ tal que
\[
  f_{\beta\delta}(A_\delta)
  =
  f_{\beta\gamma}(A_\gamma)
  \qquad(\delta\geq\gamma).
\]
Entonces $A'_\beta=f_{\beta\gamma}(A_\gamma)$, y por tanto
\[
  f_{\alpha\beta}(A'_\beta)
  =
  f_{\alpha\gamma}(A_\gamma),
  \qquad
  A'_\alpha
  =
  f_{\alpha\gamma}(A_\gamma)
  =
  f_{\alpha\beta}(A'_\beta).
\]

La condición (ML) se cumple también cuando los objetos $A_\alpha$ son
\emph{artinianos} en $\cat{C}$, es decir, cuando toda familia de
subobjetos de $A_\alpha$ posee un elemento minimal.
En efecto, un elemento minimal de la familia filtrada decreciente
$\bigl(f_{\alpha\beta}(A_\beta)\bigr)$ de subobjetos de $A_\alpha$
es necesariamente el menor de esos subobjetos.
\end{env}

\begin{env}[13.1.3]
\label{0.13.1.3}
\emph{Observación} (13.1.3).
La condición (ML) puede formularse también cuando $\cat{C}$ es, por
ejemplo, la categoría de conjuntos.
Se puede definir de nuevo el subconjunto de las imágenes universales en
$A_\alpha$, y las observaciones hechas en \sref{0.13.1.1} y
\sref{0.13.1.2} siguen siendo válidas.
\end{env}

\subsection{La condición de Mittag-Leffler para grupos abelianos}
\label{subsection:0.13.2}

\begin{proposition}[13.2.1]
\label{0.13.2.1}
Sea
\[
  0\longrightarrow A_\alpha
  \xrightarrow{u_\alpha}B_\alpha
  \xrightarrow{v_\alpha}C_\alpha
  \longrightarrow0
\]
una sucesión exacta de sistemas proyectivos de grupos abelianos, relativos
al mismo conjunto de índices filtrado $\mathrm{I}$.
\begin{enumerate}
  \item[{\rm(i)}]
  Si $(B_\alpha)$ satisface {\rm(ML)}, entonces también $(C_\alpha)$.
  \item[{\rm(ii)}]
  Si $(A_\alpha)$ y $(C_\alpha)$ satisfacen {\rm(ML)}, entonces también
  $(B_\alpha)$.
\end{enumerate}
\end{proposition}

\begin{proof}
Sean $(f_{\alpha\beta})$, $(g_{\alpha\beta})$ y
$(h_{\alpha\beta})$ los sistemas de homomorfismos que definen
$(A_\alpha)$, $(B_\alpha)$ y $(C_\alpha)$, respectivamente.

(i) Supongamos que
\[
  g_{\alpha\beta}(B_\beta)
  =
  g_{\alpha\lambda}(B_\lambda)
  \qquad(\lambda\geq\beta).
\]
Como $v_\beta$ y $v_\lambda$ son sobreyectivos,
\[
  h_{\alpha\beta}(C_\beta)
  =
  v_\alpha\bigl(g_{\alpha\beta}(B_\beta)\bigr)
  =
  v_\alpha\bigl(g_{\alpha\lambda}(B_\lambda)\bigr)
  =
  h_{\alpha\lambda}(C_\lambda)
\]
para $\lambda\geq\beta$.

(ii) Sea $\alpha\in\mathrm{I}$.
Elijamos $\beta\geq\alpha$ tal que
\[
  f_{\alpha\beta}(A_\beta)
  =
  f_{\alpha\lambda}(A_\lambda)
  \qquad(\lambda\geq\beta),
\]
y elijamos después $\gamma\geq\beta$ tal que
\[
  h_{\beta\gamma}(C_\gamma)
  =
  h_{\beta\lambda}(C_\lambda)
  \qquad(\lambda\geq\gamma).
\]
Sea $y_\alpha\in g_{\alpha\gamma}(B_\gamma)$, escribamos
$y_\alpha=g_{\alpha\gamma}(y_\gamma)$ con $y_\gamma\in B_\gamma$,
y pongamos $y_\beta=g_{\beta\gamma}(y_\gamma)$.
Entonces
\[
  v_\beta(y_\beta)
  =
  h_{\beta\gamma}\bigl(v_\gamma(y_\gamma)\bigr).
\]
Para todo $\lambda\geq\gamma$, existe, por hipótesis, un
$y_\lambda\in B_\lambda$ tal que
\[
  h_{\beta\gamma}\bigl(v_\gamma(y_\gamma)\bigr)
  =
  h_{\beta\lambda}\bigl(v_\lambda(y_\lambda)\bigr).
\]
De aquí que
$v_\beta(y_\beta-g_{\beta\lambda}(y_\lambda))=0$, y por tanto
existe un $x_\beta\in A_\beta$
\oldpage[0\textsubscript{III}]{66}
tal que
\[
  y_\beta
  =
  g_{\beta\lambda}(y_\lambda)+u_\beta(x_\beta).
\]
Se sigue que
\[
  y_\alpha
  =
  g_{\alpha\lambda}(y_\lambda)
  +
  u_\alpha\bigl(f_{\alpha\beta}(x_\beta)\bigr).
\]
Como $\lambda\geq\beta$, existe un $x_\lambda\in A_\lambda$ tal
que
\[
  f_{\alpha\beta}(x_\beta)=f_{\alpha\lambda}(x_\lambda).
\]
Por consiguiente,
\[
  y_\alpha
  =
  g_{\alpha\lambda}
  \bigl(y_\lambda+u_\lambda(x_\lambda)\bigr)
  \in g_{\alpha\lambda}(B_\lambda),
\]
lo que demuestra la aserción.
\end{proof}

\begin{proposition}[13.2.2]
\label{0.13.2.2}
Sea $\mathrm{I}$ un conjunto ordenado filtrado que posee un subconjunto
cofinal numerable, y sea
\[
  0\longrightarrow A_\alpha
  \xrightarrow{u_\alpha}B_\alpha
  \xrightarrow{v_\alpha}C_\alpha
  \longrightarrow0
\]
una sucesión exacta de sistemas proyectivos de grupos abelianos indexados
por $\mathrm{I}$.
Si $(A_\alpha)$ satisface {\rm(ML)}, entonces la sucesión
\[
  0\longrightarrow\varprojlim A_\alpha
  \longrightarrow\varprojlim B_\alpha
  \longrightarrow\varprojlim C_\alpha
  \longrightarrow0
\]
es exacta.
\end{proposition}

\begin{proof}
Solo queda demostrar que el homomorfismo
\[
  v=\varprojlim v_\alpha:
  \varprojlim B_\alpha\longrightarrow\varprojlim C_\alpha
\]
es sobreyectivo.
Sea $z=(z_\alpha)\in\varprojlim C_\alpha$, y pongamos
\[
  E_\alpha=v_\alpha^{-1}(z_\alpha).
\]
Los $E_\alpha$ forman manifiestamente un sistema proyectivo de conjuntos no
vacíos mediante las restricciones de los homomorfismos
$g_{\alpha\beta}:B_\beta\longrightarrow B_\alpha$.

Demostremos que este sistema proyectivo satisface (ML).
Identificando $A_\alpha$ con un subgrupo de $B_\alpha$ por medio de
$u_\alpha$, elijamos, para todo $\alpha\in\mathrm{I}$, un índice
$\beta\geq\alpha$ tal que
\[
  g_{\alpha\beta}(A_\beta)
  =
  g_{\alpha\lambda}(A_\lambda)
  \qquad(\lambda\geq\beta).
\]
Afirmamos que también
\[
  g_{\alpha\beta}(E_\beta)
  =
  g_{\alpha\lambda}(E_\lambda)
  \qquad(\lambda\geq\beta).
\]
En efecto, tomemos $y_\lambda\in E_\lambda$ y pongamos
\[
  y_\beta=g_{\beta\lambda}(y_\lambda),
  \qquad
  y_\alpha=g_{\alpha\lambda}(y_\lambda).
\]
Sea $y'_\alpha=g_{\alpha\beta}(y'_\beta)$ para algún
$y'_\beta\in E_\beta$.
Entonces $y'_\beta-y_\beta=x_\beta\in A_\beta$, y por hipótesis existe
un $x_\lambda\in A_\lambda$ tal que
\[
  g_{\alpha\beta}(x_\beta)
  =
  g_{\alpha\lambda}(x_\lambda).
\]
Así,
\[
  y'_\alpha
  =
  g_{\alpha\beta}(y_\beta)+g_{\alpha\beta}(x_\beta)
  =
  g_{\alpha\lambda}(y_\lambda+x_\lambda)
  \in g_{\alpha\lambda}(E_\lambda),
\]
lo que demuestra la afirmación.

Se sabe (Bourbaki, \emph{Top.\ g\'{e}n.}, cap.~II,
\textsection3, teorema~1) que, bajo las hipótesis impuestas a
$\mathrm{I}$, un sistema proyectivo de conjuntos no vacíos que satisface (ML)
posee un límite proyectivo no vacío.
Existe, por tanto, un punto
$y=(y_\alpha)\in\varprojlim E_\alpha$.
Como $v_\alpha(y_\alpha)=z_\alpha$ para todo $\alpha$, se tiene
$z=v(y)$.
\end{proof}

\begin{proposition}[13.2.3]
\label{0.13.2.3}
Bajo las hipótesis sobre $\mathrm{I}$ de \sref{0.13.2.2}, sea
$(\mathrm{K}^\bullet_\alpha)_{\alpha\in\mathrm{I}}$ un sistema proyectivo
de complejos de grupos abelianos
\[
  \mathrm{K}^\bullet_\alpha
  =
  (\mathrm{K}^n_\alpha)_{n\in\bb{Z}},
\]
cuyas diferenciales tienen grado $+1$.
Para todo $n$, existe un homomorfismo canónico funtorial
\begin{equation}
\label{0.13.2.3.1}
\tag{13.2.3.1}
  h_n:
  \HH^n\bigl(\varprojlim\mathrm{K}^\bullet_\alpha\bigr)
  \longrightarrow
  \varprojlim\HH^n(\mathrm{K}^\bullet_\alpha).
\end{equation}
Si, para todo grado $n$, el sistema proyectivo de grupos abelianos
$(\mathrm{K}^n_\alpha)_{\alpha\in\mathrm{I}}$ satisface {\rm(ML)},
entonces todo $h_n$ es sobreyectivo.
Si, además, el sistema proyectivo
$(\HH^{n-1}(\mathrm{K}^\bullet_\alpha))_{\alpha\in\mathrm{I}}$
satisface {\rm(ML)}, entonces $h_n$ es biyectivo.
\end{proposition}

\begin{proof}
Pongamos
\[
  \mathrm{K}^n=\varprojlim\mathrm{K}^n_\alpha
  \qquad(n\in\bb{Z}).
\]
Los homomorfismos $h_n$ están definidos por los diagramas conmutativos
\[
  \xymatrix{
    \cdots\ar[r]
      & \mathrm{K}^{n-1}\ar[r]\ar[d]
      & \mathrm{K}^n\ar[r]\ar[d]
      & \mathrm{K}^{n+1}\ar[r]\ar[d]
      & \cdots\\
    \cdots\ar[r]
      & \mathrm{K}^{n-1}_\alpha\ar[r]
      & \mathrm{K}^n_\alpha\ar[r]
      & \mathrm{K}^{n+1}_\alpha\ar[r]
      & \cdots,
  }
\]
\oldpage[0\textsubscript{III}]{67}
siendo las diferenciales en $\mathrm{K}^\bullet$ los límites proyectivos
de las diferenciales correspondientes en los
$\mathrm{K}^\bullet_\alpha$.
Consideremos las sucesiones exactas
\[
  (*_n)\qquad
  0\longrightarrow
  \mathrm{B}^n(\mathrm{K}^\bullet_\alpha)
  \longrightarrow
  \mathrm{Z}^n(\mathrm{K}^\bullet_\alpha)
  \longrightarrow
  \HH^n(\mathrm{K}^\bullet_\alpha)
  \longrightarrow0
\]
y
\[
  (**_n)\qquad
  0\longrightarrow
  \mathrm{Z}^{n-1}(\mathrm{K}^\bullet_\alpha)
  \longrightarrow
  \mathrm{K}^{n-1}_\alpha
  \longrightarrow
  \mathrm{B}^n(\mathrm{K}^\bullet_\alpha)
  \longrightarrow0.
\]
La hipótesis y la Proposición~\sref{0.13.2.1}~(i) muestran que el
sistema proyectivo
$(\mathrm{B}^n(\mathrm{K}^\bullet_\alpha))_{\alpha\in\mathrm{I}}$
satisface (ML) para todo $n$.
Se sigue de \sref{0.13.2.2} que
\[
  (***_n)\qquad
  0\longrightarrow
  \varprojlim_\alpha\mathrm{B}^n(\mathrm{K}^\bullet_\alpha)
  \longrightarrow
  \varprojlim_\alpha\mathrm{Z}^n(\mathrm{K}^\bullet_\alpha)
  \longrightarrow
  \varprojlim_\alpha\HH^n(\mathrm{K}^\bullet_\alpha)
  \longrightarrow0
\]
es exacta.
Ahora bien,
$\varprojlim_\alpha\mathrm{B}^n(\mathrm{K}^\bullet_\alpha)$
se identifica con un subgrupo de $\mathrm{K}^{n}$ que contiene
$\mathrm{B}^n(\mathrm{K}^\bullet)$, mientras que
$\varprojlim_\alpha\mathrm{Z}^n(\mathrm{K}^\bullet_\alpha)$
se identifica con un subgrupo de $\mathrm{Z}^n(\mathrm{K}^\bullet)$.
Por consiguiente, $h_n$ es sobreyectivo.

Supongamos ahora, además, que
$(\HH^{n-1}(\mathrm{K}^\bullet_\alpha))_{\alpha\in\mathrm{I}}$
satisface (ML).
Las sucesiones exactas $(*_{n-1})$ y la
Proposición~\sref{0.13.2.1}~(ii) muestran que
$(\mathrm{Z}^{n-1}(\mathrm{K}^\bullet_\alpha))_{\alpha\in\mathrm{I}}$
satisface (ML).
Aplicando \sref{0.13.2.2} a las sucesiones exactas $(**_n)$ se obtiene una
sucesión exacta
\[
  0\longrightarrow
  \varprojlim_\alpha
    \mathrm{Z}^{n-1}(\mathrm{K}^\bullet_\alpha)
  \longrightarrow
  \mathrm{K}^{n-1}
  \xrightarrow{u}
  \varprojlim_\alpha
    \mathrm{B}^n(\mathrm{K}^\bullet_\alpha)
  \longrightarrow0.
\]
Como
$\varprojlim_\alpha\mathrm{B}^n(\mathrm{K}^\bullet_\alpha)
\supset\mathrm{B}^n(\mathrm{K}^\bullet)$
y el compuesto de $u$ con la inyección
\[
  \varprojlim_\alpha\mathrm{B}^n(\mathrm{K}^\bullet_\alpha)
  \longrightarrow\mathrm{K}^n
\]
es la diferencial
$\mathrm{K}^{n-1}\longrightarrow\mathrm{K}^n$, la sobreyectividad de
$u$ implica
\[
  \varprojlim_\alpha\mathrm{B}^n(\mathrm{K}^\bullet_\alpha)
  =
  \mathrm{B}^n(\mathrm{K}^\bullet).
\]
Así, $h_n$ es inyectivo.
\end{proof}

\begin{env}[13.2.4]
\label{0.13.2.4}
\emph{Observaciones} (13.2.4).
\begin{enumerate}
  \item[{\rm(i)}]
  El argumento de \sref{0.13.2.2}
  (cf.\ Bourbaki, \emph{loc.\ cit.}) muestra que la conclusión de
  esa proposición sigue siendo válida si se supone únicamente que los
  $A_\alpha$ pueden dotarse de estructuras de espacios metrizables completos
  en los que las traslaciones son homeomorfismos, que las aplicaciones
  $f_{\alpha\beta}:A_\beta\longrightarrow A_\alpha$ que definen el
  sistema proyectivo son uniformemente continuas para las métricas en
  cuestión, y que $(A_\alpha)$ satisface la condición
  \begin{itemize}
    \item[\textbf{(ML')}]
    \emph{Para todo índice $\alpha$, existe un $\beta\geq\alpha$ tal
    que, para todo $\gamma\geq\beta$,
    $f_{\alpha\gamma}(A_\gamma)$ es denso en
    $f_{\alpha\beta}(A_\beta)$.}
  \end{itemize}

  Esto da un suplemento análogo a \sref{0.13.2.3}.
  Supongamos que $\mathrm{K}^n_\alpha=0$ para $n<0$ y todo $\alpha$;
  que
  $(\mathrm{K}^n_\alpha)_{\alpha\in\mathrm{I}}$ satisface (ML) para
  $n\geq0$; y que
  $A_\alpha=\HH^0(\mathrm{K}^\bullet_\alpha)$ puede dotarse de
  estructuras de espacio métrico que tengan las propiedades recién enunciadas.
  Entonces las conclusiones de \sref{0.13.2.3} no cambian para
  $n\geq2$, y $h_1$ es además \emph{biyectivo}.
  En efecto, el argumento de \sref{0.13.2.2} muestra también que
  $(\mathrm{B}^1(\mathrm{K}^\bullet_\alpha))_{\alpha\in\mathrm{I}}$
  satisface (ML), que la sucesión $(***_1)$ es exacta y, finalmente,
  que
  \[
    \varprojlim_\alpha
    \mathrm{B}^1(\mathrm{K}^\bullet_\alpha)
    =
    \mathrm{B}^1(\mathrm{K}^\bullet).
  \]
  Esto demuestra, entre otras cosas, las aserciones de (T, 3.10.2).

  \item[{\rm(ii)}]
  Se pueden introducir los funtores derivados por la derecha
  $\varprojlim^{(n)}$ del funtor $\varprojlim$, y obtener enunciados más
  completos que los precedentes~\cite{III-28}.
\end{enumerate}
\end{env}

\subsection{Aplicación: cohomología de un límite proyectivo de haces}
\label{subsection:0.13.3}

\begin{proposition}[13.3.1]
\label{0.13.3.1}
\oldpage[0\textsubscript{III}]{68}
Sea $X$ un espacio topológico, sea
$(\sh{F}_k)_{k\in\bb{N}}$ un sistema proyectivo de haces de
grupos abelianos sobre $X$, y pongamos
$\sh{F}=\varprojlim_k\sh{F}_k$.
Supongamos que se cumplen las condiciones siguientes:
\begin{enumerate}
  \item[{\rm(i)}]
  Existe una base $\mathfrak{B}$ de la topología de $X$ tal que,
  para todo $U\in\mathfrak{B}$ y todo $i\geq0$, el sistema
  proyectivo
  $(\HH^i(U,\sh{F}_k))_{k\in\bb{N}}$ satisface (ML).

  \item[{\rm(ii)}]
  Para todo $x\in X$ y todo $i>0$, se tiene
  \[
    \varinjlim_U\bigl(\varprojlim_k
    \HH^i(U,\sh{F}_k)\bigr)=0,
  \]
  donde $U$ recorre los entornos de $x$ pertenecientes a
  $\mathfrak{B}$.

  \item[{\rm(iii)}]
  Los homomorfismos
  $u_{hk}:\sh{F}_k\longrightarrow\sh{F}_h$ $(h\leq k)$ que definen el
  sistema proyectivo $(\sh{F}_k)$ son sobreyectivos.
\end{enumerate}
Bajo estas condiciones, para todo $i>0$, el homomorfismo canónico
\[
  h_i:\HH^i(X,\sh{F})
  \longrightarrow
  \varprojlim_k\HH^i(X,\sh{F}_k)
\]
es sobreyectivo. Si, además, para un valor dado de $i$ el sistema
proyectivo
$(\HH^{i-1}(X,\sh{F}_k))_{k\in\bb{N}}$ satisface (ML), entonces $h_i$ es
biyectivo.
\end{proposition}

\begin{proof}
\emph{a)}
Supongamos primero que los $\sh{F}_k$, así como los núcleos
$\sh{N}_{hk}$ de los $u_{hk}$, son flácidos; mostraremos entonces que
la condición~{\rm(iii)} del enunciado implica
$\HH^i(X,\sh{F})=0$ para $i>0$.
Basta demostrar que, para todo subconjunto abierto $U$ de $X$ y todo
recubrimiento $\mathfrak{U}$ de $U$ por subconjuntos abiertos de $U$, se tiene
\[
  \CHH^i(\mathfrak{U},\sh{F})=0
  \qquad (i>0).
\]
En efecto, se sigue entonces primero que, para la cohomología de \v{C}ech,
$\check{\HH}^{\,i}(U,\sh{F})=0$ para todo $i>0$, y por tanto, por
(G, II, 5.9.2) aplicado al conjunto de todos los subconjuntos abiertos de
$X$, que $\HH^i(X,\sh{F})=0$ para todo $i>0$.

Como los $\sh{F}_k$ son flácidos, se tiene
$\CHH^i(\mathfrak{U},\sh{F}_k)=0$ para todo $i>0$
(G, II, 5.2.3).
Para cada $k$, consideremos el complejo
$\check{\mathrm{C}}^\bullet(\mathfrak{U},\sh{F}_k)$ de cocadenas
alternadas sobre el nervio del recubrimiento $\mathfrak{U}$ (G, II, 5.1);
estos complejos forman manifiestamente un sistema proyectivo de complejos de
grupos abelianos.
Demostremos que todas las aplicaciones
\[
  \check{\mathrm{C}}^\bullet(\mathfrak{U},\sh{F}_k)
  \longrightarrow
  \check{\mathrm{C}}^\bullet(\mathfrak{U},\sh{F}_h)
  \qquad(h\leq k)
\]
son sobreyectivas.
Por definición, basta mostrar que, para todo subconjunto abierto $V$
de $X$, la aplicación
$\Gamma(V,\sh{F}_k)\longrightarrow\Gamma(V,\sh{F}_h)$ es
es sobreyectiva. Pero la sucesión
\[
  0\longrightarrow\sh{N}_{hk}\longrightarrow\sh{F}_k
  \longrightarrow\sh{F}_h\longrightarrow0
\]
es exacta por hipótesis y da, por tanto, la sucesión exacta de
cohomología
\[
  \Gamma(V,\sh{F}_k)\longrightarrow\Gamma(V,\sh{F}_h)
  \longrightarrow\HH^1(V,\sh{N}_{hk})=0,
\]
pues $\sh{N}_{hk}$ es flácido.

El sistema proyectivo
$(\check{\mathrm{C}}^\bullet(\mathfrak{U},\sh{F}_k))_{k\in\bb{N}}$
satisface, por tanto, (ML), y lo mismo vale para
$(\CHH^i(\mathfrak{U},\sh{F}_k))_{k\in\bb{N}}$ para todo $i\geq0$:
para $i>0$ esto es trivial, mientras que para $i=0$ se sigue de lo
precedente, pues
$\CHH^0(\mathfrak{U},\sh{F}_k)=\Gamma(U,\sh{F}_k)$
(G, II, 5.2.2).
Podemos, por tanto, aplicar \sref{0.13.2.3}, que da
\[
  \CHH^i(\mathfrak{U},\sh{F})
  =
  \varprojlim_k\CHH^i(\mathfrak{U},\sh{F}_k)
  =0
  \qquad(i>0).
\]

\emph{b)}
Pasemos ahora al caso general. Para cada $k\in\bb{N}$, consideremos la
\emph{resolución canónica}
\[
  \sh{C}^\bullet(X,\sh{F}_k)
  =
  \bigl(\sh{C}^i(X,\sh{F}_k)\bigr)_{i\geq0}
\]
de $\sh{F}_k$ por haces flácidos (G, II, 4.3).
Para todo $i\geq0$, está claro que
$(\sh{C}^i(X,\sh{F}_k))_{k\in\bb{N}}$ es un sistema proyectivo de
haces flácidos; demostremos que satisface las condiciones de la
parte~\emph{a)}.
\oldpage[0\textsubscript{III}]{69}
En efecto, si $\sh{N}_{hk}$ denota el núcleo de $u_{hk}$ para $h\leq k$,
entonces la sucesión
\[
  0\longrightarrow\sh{N}_{hk}\longrightarrow\sh{F}_k
  \longrightarrow\sh{F}_h\longrightarrow0
\]
es exacta por~{\rm(iii)}, y nuestra aserción se sigue de la exactitud
del funtor
$\sh{A}\longmapsto\sh{C}^i(X,\sh{A})$ (G, II, 4.3).

Pongamos
\[
  \sh{G}^i=\varprojlim_k\sh{C}^i(X,\sh{F}_k).
\]
Por la parte~\emph{a)}, se tiene
$\HH^j(X,\sh{G}^i)=0$ para $j>0$ e $i\geq0$.
Demostraremos que
$\sh{G}^\bullet=(\sh{G}^i)_{i\geq0}$ es una resolución del haz
$\sh{F}$. Como $\HH^j(X,\sh{G}^i)=0$ para $j>0$, la cohomología
$\HH^\bullet(X,\sh{F})$ será entonces igual a
$\HH^\bullet(\Gamma(X,\sh{G}^\bullet))$ (G, II, 4.7.1).

Tomando el límite proyectivo de las sucesiones exactas
\[
  0\longrightarrow\sh{F}_k
  \longrightarrow\sh{C}^0(X,\sh{F}_k)
  \longrightarrow\sh{C}^1(X,\sh{F}_k)
  \longrightarrow\cdots
\]
se obtiene manifiestamente un complejo de haces de grupos abelianos
\[
  0\longrightarrow\sh{F}
  \longrightarrow\sh{G}^0
  \longrightarrow\sh{G}^1
  \longrightarrow\cdots.
\]
Para demostrar nuestra aserción, queda establecer que
$\sh{H}^i(\sh{G}^\bullet)=0$ para $i>0$.
Este haz está generado por el prehaz
\[
  U\longmapsto\HH^i\bigl(\Gamma(U,\sh{G}^\bullet)\bigr)
\]
(G, II, 4.1).
Ahora bien, el complejo $\Gamma(U,\sh{G}^\bullet)$ es el límite proyectivo
del sistema proyectivo de complejos de grupos abelianos
\[
  \bigl(\Gamma(U,\sh{C}^\bullet(X,\sh{F}_k))\bigr)_{k\in\bb{N}}
\]
(0\textsubscript{I}, \EGAUnavailableHyperref{0.3.2.6}{3.2.6}).
Vimos en la parte~\emph{a)} que, para todo $i\geq0$, las aplicaciones
\[
  \Gamma(U,\sh{C}^i(X,\sh{F}_k))
  \longrightarrow
  \Gamma(U,\sh{C}^i(X,\sh{F}_h))
  \qquad(h\leq k)
\]
son sobreyectivas. Por otra parte,
\[
  \HH^i(U,\sh{F}_k)
  =
  \HH^i\bigl(\Gamma(U,\sh{C}^\bullet(X,\sh{F}_k))\bigr),
\]
pues la resolución canónica $\sh{C}^\bullet(U,\sh{F}_k|U)$ es
inducida sobre $U$ por $\sh{C}^\bullet(X,\sh{F}_k)$.
Por la hipótesis~{\rm(i)}, para todo $U\in\mathfrak{B}$ podemos aplicar
\sref{0.13.2.3} a este sistema proyectivo de complejos, y por tanto
\[
  \HH^i\bigl(\Gamma(U,\sh{G}^\bullet)\bigr)
  =
  \varprojlim_k\HH^i(U,\sh{F}_k)
  \qquad(i\geq0).
\]
La hipótesis~{\rm(ii)} muestra entonces, por definición, que los haces
$\sh{H}^i(\sh{G}^\bullet)$ se anulan para $i>0$.

Así, para todo $i\geq0$,
\[
  \HH^i(X,\sh{F})
  =
  \HH^i\bigl(\Gamma(X,\sh{G}^\bullet)\bigr),
  \qquad
  \Gamma(X,\sh{G}^\bullet)
  =
  \varprojlim_k\Gamma(X,\sh{C}^\bullet(X,\sh{F}_k)).
\]
Acabamos de observar que todas las aplicaciones
\[
  \Gamma(X,\sh{C}^i(X,\sh{F}_k))
  \longrightarrow
  \Gamma(X,\sh{C}^i(X,\sh{F}_h))
  \qquad(h\leq k)
\]
son sobreyectivas; la conclusión se sigue, por tanto, una vez más de
\sref{0.13.2.3}.
\end{proof}

\begin{env}[13.3.2]
\label{0.13.3.2}
\emph{Observaciones} (13.3.2).
\begin{enumerate}
  \item[{\rm(i)}]
  El enunciado de \sref{0.13.3.1} solo tiene interés para $i>0$,
  pues, para $i=0$, $h_i$ es siempre un isomorfismo sin hipótesis
  alguna (0\textsubscript{I}, \EGAUnavailableHyperref{0.3.2.6}{3.2.6}).

  \item[{\rm(ii)}]
  Las condiciones~{\rm(i)} y~{\rm(ii)} de \sref{0.13.3.1} se satisfacen,
  en particular, si
  $\HH^i(U,\sh{F}_k)=0$ para todo $k$, todo $i>0$ y todo
  $U\in\mathfrak{B}$, y si, para $U\in\mathfrak{B}$, las aplicaciones
  \[
    \Gamma(U,\sh{F}_k)\longrightarrow\Gamma(U,\sh{F}_h)
  \]
  son sobreyectivas. Este es el caso más frecuente en el que se aplica
  \sref{0.13.3.1}.
\end{enumerate}
\end{env}

\subsection{Condición de Mittag-Leffler y objetos graduados asociados a sistemas proyectivos}
\label{subsection:0.13.4}
\label{0.13.4}

\begin{env}[13.4.1]
\label{0.13.4.1}
Sea
$\bb{A}=(A_k,u_{kh})_{k\in\bb{Z}}$ un sistema proyectivo en una
categoría abeliana $\cat{C}$.
Decimos que está \emph{acotado inferiormente} si existe un $k_0$ tal que
$A_k=0$ para $k<k_0$.
Sobre cada $A_k$ definimos una \emph{filtración}
$(\mathrm{F}^p(A_k))_{p\in\bb{Z}}$ por
\begin{equation}
\label{0.13.4.1.1}
\tag{13.4.1.1}
  \mathrm{F}^p(A_k)=
  \begin{cases}
    \Ker(A_k\longrightarrow A_{p-1}) & \text{para }p\leq k+1,\\
    0 & \text{para }p\geq k+1.
  \end{cases}
\end{equation}

\oldpage[0\textsubscript{III}]{70}
Por hipótesis,
$\mathrm{F}^{k_0}(A_k)=A_k$ y
$\mathrm{F}^{k+1}(A_k)=0$; en otras palabras, la filtración considerada
es finita (\sref{0.11.1.3}).
Los objetos graduados asociados a esta filtración son, por tanto,
\[
  \mathrm{gr}^p(A_k)
  =
  \Ker(A_k\longrightarrow A_{p-1})
  \big/
  \Ker(A_k\longrightarrow A_p).
\]
Por consiguiente, $\mathrm{gr}^p(A_k)$ es isomorfo a la imagen por
$A_k\longrightarrow A_p$ de
$\Ker(A_k\longrightarrow A_{p-1})$.
Por la transitividad de los morfismos que definen un sistema proyectivo,
se tiene, por tanto,
\begin{equation}
\label{0.13.4.1.2}
\tag{13.4.1.2}
  \mathrm{gr}^p(A_k)
  =
  \Ker(A_p\longrightarrow A_{p-1})
  \cap
  \Im(A_k\longrightarrow A_p).
\end{equation}
Pero \sref{0.13.4.1.1} da
$\Ker(A_p\longrightarrow A_{p-1})=\mathrm{gr}^p(A_p)$, y por tanto
\begin{equation}
\label{0.13.4.1.3}
\tag{13.4.1.3}
  \mathrm{gr}^p(A_k)
  =
  \mathrm{gr}^p(A_p)
  \cap
  \Im(A_k\longrightarrow A_p).
\end{equation}

Las definiciones precedentes muestran también que, para $k\leq h$,
\[
  u_{kh}\bigl(\mathrm{F}^p(A_h)\bigr)
  \subset
  \mathrm{F}^p(A_k),
\]
y, por consiguiente, para todo $p\in\bb{Z}$, los
$\mathrm{gr}^p(u_{kh})$ definen un sistema proyectivo
$(\mathrm{gr}^p(A_k))_{k\in\bb{Z}}$.
\end{env}

\begin{env}[13.4.2]
\label{0.13.4.2}
Decimos que el sistema proyectivo $\bb{A}$ es
\emph{esencialmente constante} si los morfismos
$A_{k+1}\longrightarrow A_k$ son isomorfismos para todo $k$
suficientemente grande.
Decimos que $\bb{A}$ es \emph{estricto} si los morfismos
$A_i\longrightarrow A_j$ $(j\leq i)$ son epimorfismos.

Cuando $\bb{A}$ es estricto, se sigue de \sref{0.13.4.1.3} que, para
$p\leq k\leq h$, el morfismo canónico
\[
  \mathrm{gr}^p(A_h)\longrightarrow\mathrm{gr}^p(A_k)
\]
es un isomorfismo; en otras palabras, el sistema proyectivo
$(\mathrm{gr}^p(A_k))_{k\in\bb{Z}}$ es esencialmente constante.
La sucesión de objetos $\mathrm{gr}^p(A_p)$, identificada para todo
$p$ con $\varprojlim_k\mathrm{gr}^p(A_k)$, se denota entonces
$\mathrm{gr}^\bullet(\bb{A})$ y se llama el
\emph{objeto graduado asociado al sistema proyectivo estricto}
$\bb{A}=(A_k)$.

Supongamos ahora que el sistema proyectivo $\bb{A}$, acotado
inferiormente, satisface (ML).
Por \sref{0.13.1.2}, el sistema proyectivo
$\bb{A}'=(A'_k)$ de objetos de ``imagen universal'' es estricto; también está
acotado inferiormente.
El objeto graduado $\mathrm{gr}^\bullet(\bb{A}')$ asociado a
$\bb{A}'$ se llama asimismo el objeto graduado
\emph{asociado a} $\bb{A}$ y se denota
$\mathrm{gr}^\bullet(\bb{A})$.
\end{env}

\begin{proposition}[13.4.3]
\label{0.13.4.3}
Sea $\bb{A}=(A_k)_{k\in\bb{Z}}$ un sistema proyectivo acotado
inferiormente en una categoría abeliana $\cat{C}$.
Las condiciones siguientes son equivalentes:
\begin{enumerate}
  \item[{\rm a)}] $\bb{A}$ satisface (ML).
  \item[{\rm b)}] Para todo $p\in\bb{Z}$, el sistema proyectivo
  $(\mathrm{gr}^p(A_k))_{k\in\bb{Z}}$ es esencialmente constante.
\end{enumerate}
Además, cuando se cumplen estas condiciones, existe, para todo
$p\in\bb{Z}$, un isomorfismo canónico
\begin{equation}
\label{0.13.4.3.1}
\tag{13.4.3.1}
  \mathrm{gr}^p(\bb{A})
  \isoto
  \varprojlim_k\mathrm{gr}^p(A_k).
\end{equation}
\end{proposition}

\begin{proof}
Se sigue inmediatamente de \sref{0.13.4.1.2} que \emph{a)} implica
\emph{b)}. La misma fórmula, aplicada al sistema proyectivo
$\bb{A}'$ con la notación de \sref{0.13.4.2}, da
\sref{0.13.4.3.1} por definición.

Para $k\leq h$, pongamos
\[
  A_{kh}=\Im(A_h\longrightarrow A_k).
\]
Si $k\leq h\leq j$, entonces
$A_{kj}\subset A_{kh}\subset A_k$.
Dotemos a $A_{kh}$ de la filtración inducida por
$(\mathrm{F}^p(A_k))$.
La transitividad de los morfismos que definen $\bb{A}$ muestra
inmediatamente que esta es también la filtración cociente de
$(\mathrm{F}^p(A_h))$; por consiguiente,
\begin{equation}
\label{0.13.4.3.2}
\tag{13.4.3.2}
  \mathrm{gr}^p(A_{kh})
  =
  \Im\bigl(
    \mathrm{gr}^p(A_h)\longrightarrow\mathrm{gr}^p(A_k)
  \bigr).
\end{equation}

\oldpage[0\textsubscript{III}]{71}
Supongamos ahora que se cumple \emph{b)}.
Para todo $p\in\bb{Z}$ y todo $k\geq p$, existe un entero
$\mathrm{L}(p,k)$ tal que el miembro derecho de
\sref{0.13.4.3.2} es constante para $h\geq\mathrm{L}(p,k)$.
Como $\mathrm{gr}^p(A_k)=0$ para $p<k_0$, para $k$ fijo solo
un número finito de los enteros $\mathrm{L}(p,k)$ puede ser no nulo cuando
$p$ recorre los enteros $\leq k$.
Sea $\mathrm{L}(k)=m$ el mayor de estos enteros.
Para todo $h\geq m$, se tiene
$A_{kh}\subset A_{km}$, y, por la definición de $m$, la inyección
canónica
$A_{kh}\longrightarrow A_{km}$ induce un isomorfismo
\[
  \mathrm{gr}^\bullet(A_{kh})
  \simto
  \mathrm{gr}^\bullet(A_{km}).
\]
Como las filtraciones son finitas, esta inyección es ella misma biyectiva
(Bourbaki, \emph{Alg.\ comm.}, cap.~III, \textsection2,
n\textsuperscript{o}~8, teorema~1).
Esto demuestra que $\bb{A}$ satisface (ML).
\end{proof}

\begin{env}[13.4.4]
\label{0.13.4.4}
Supongamos que el límite proyectivo
$A=\varprojlim_k A_k$ existe en $\cat{C}$.
En las definiciones de \sref{0.13.4.1}, se puede entonces reemplazar $A_k$
por $A$, y la filtración así definida sobre $A$ sigue satisfaciendo
\begin{equation}
\label{0.13.4.4.1}
\tag{13.4.4.1}
  \mathrm{gr}^p(A)
  =
  \mathrm{gr}^p(A_p)
  \cap
  \Im(A\longrightarrow A_p).
\end{equation}
\end{env}

\begin{corollary}[13.4.5]
\label{0.13.4.5}
Supongamos que $\cat{C}$ es la categoría de grupos abelianos.
Si el sistema proyectivo $\bb{A}$ satisface (ML) y
$A=\varprojlim_k A_k$, entonces para todo $p\in\bb{Z}$ existe un
isomorfismo canónico
\begin{equation}
\label{0.13.4.5.1}
\tag{13.4.5.1}
  \mathrm{gr}^p(A)
  \isoto
  \varprojlim_k\mathrm{gr}^p(A_k).
\end{equation}
\end{corollary}

\begin{proof}
En efecto,
\[
  \Im(A_k\longrightarrow A_p)
  =
  \Im(A\longrightarrow A_p)
\]
para todo $k$ suficientemente grande
(Bourbaki, \emph{Top.\ g\'{e}n.}, cap.~II,
3\textsuperscript{e}~\'{e}d., \textsection3,
n\textsuperscript{o}~5, teorema~1), y la conclusión se sigue de
\sref{0.13.4.1.3} y \sref{0.13.4.4.1}.
\end{proof}

\subsection{Límites proyectivos de sucesiones espectrales de complejos filtrados}
\label{subsection:0.13.5}

\begin{env}[13.5.1]
\label{0.13.5.1}
Sea $\cat{C}$ una categoría abeliana y sea $X^\bullet$ un complejo
de objetos de $\cat{C}$ dotado de una filtración
$(\mathrm{F}^p(X^\bullet))_{p\in\bb{Z}}$ tal que
$\mathrm{F}^{p_0}(X^\bullet)=X^\bullet$ para algún índice $p_0$.
Para todo $k\in\bb{Z}$, consideremos el complejo
\[
  X^\bullet_k=X^\bullet/\mathrm{F}^{k+1}(X^\bullet);
\]
está dotado canónicamente de la filtración constituida por
\[
  \mathrm{F}^p(X^\bullet_k)
  =
  \mathrm{F}^p(X^\bullet)/\mathrm{F}^{k+1}(X^\bullet)
  \quad\text{para }p\leq k,
\]
y $\mathrm{F}^p(X^\bullet_k)=0$ para $p\geq k+1$.
Además, existen morfismos canónicos
$X^\bullet_{k+1}\longrightarrow X^\bullet_k$, que hacen de
$\bb{X}^\bullet=(X^\bullet_k)_{k\in\bb{Z}}$ un \emph{sistema proyectivo
de complejos filtrados de objetos de $\cat{C}$}.
Observemos que este sistema proyectivo es \emph{estricto} y que
$X^\bullet_k=0$ para $k<p_0$.
\end{env}

\begin{env}[13.5.2]
\label{0.13.5.2}
Más generalmente, consideremos un sistema proyectivo estricto acotado
inferiormente
$\bb{X}^\bullet=(X^\bullet_k)_{k\in\bb{Z}}$ de complejos de objetos de
$\cat{C}$, y dotemos cada $X^\bullet_k$ de la filtración definida en
\sref{0.13.4.1} (aplicada en la categoría abeliana de complejos acotados
inferiormente de objetos de $\cat{C}$).
Los morfismos
$X^\bullet_k\longrightarrow X^\bullet_p$ ($p\leq k$) se convierten entonces
en morfismos de complejos \emph{filtrados} con filtraciones finitas.
La funtorialidad de las sucesiones espectrales de complejos filtrados
(\sref{0.11.2.3}) muestra que los morfismos de transición de
$\bb{X}^\bullet$ inducen morfismos que hacen de
\[
  \mathrm{E}(\bb{X}^\bullet)
  =
  \bigl(\mathrm{E}(X^\bullet_k)\bigr)_{k\in\bb{Z}}
\]
un \emph{sistema proyectivo de sucesiones espectrales}.
\end{env}

\begin{lemma}[13.5.3]
\label{0.13.5.3}
Supongamos que el sistema proyectivo
$\bb{X}^\bullet=(X^\bullet_k)_{k\in\bb{Z}}$ de complejos filtrados se
obtiene como en \sref{0.13.5.2}. Entonces:
\begin{enumerate}
  \item[{\rm a)}] Para $r\geq p-p_0$, se tiene
  \[
    \mathrm{B}_r^{pq}(X^\bullet_k)
    =
    \mathrm{B}_\infty^{pq}(X^\bullet_k)
  \]
  para todo $k\in\bb{Z}$.
  \item[{\rm b)}] Para $k+1\leq p+r$, se tiene
  \[
    \mathrm{Z}_r^{pq}(X^\bullet_k)
    =
    \mathrm{Z}_\infty^{pq}(X^\bullet_k).
  \]
  \item[{\rm c)}] Para $k+1\geq p+r$, los morfismos
  \[
    \mathrm{Z}_r^{pq}(X^\bullet_h)
    \longrightarrow
    \mathrm{Z}_r^{pq}(X^\bullet_k)
    \quad\text{y}\quad
    \mathrm{B}_r^{pq}(X^\bullet_h)
    \longrightarrow
    \mathrm{B}_r^{pq}(X^\bullet_k)
  \]
  son isomorfismos para todo $h\geq k$.
\end{enumerate}
\end{lemma}

\oldpage[0\textsubscript{III}]{72}
Estas tres propiedades se siguen inmediatamente de las definiciones de
\sref{0.11.2.2}, teniendo en cuenta que
$\mathrm{F}^{p-r+1}(X^\bullet_k)=X^\bullet_k$ cuando $p-r<p_0$.

\begin{env}[13.5.4]
\label{0.13.5.4}
Supongamos que se cumplen las hipótesis de \sref{0.13.5.3}.
Entonces, para $p$, $q$ y $r$ fijos (con $r$ finito), los sistemas
proyectivos
\[
  \bigl(\mathrm{Z}_r^{pq}(X^\bullet_k)\bigr)_{k\in\bb{Z}},
  \qquad
  \bigl(\mathrm{B}_r^{pq}(X^\bullet_k)\bigr)_{k\in\bb{Z}},
  \qquad
  \bigl(\mathrm{E}_r^{pq}(X^\bullet_k)\bigr)_{k\in\bb{Z}}
\]
son \emph{esencialmente constantes}; sus respectivos límites proyectivos
se denotarán por
$\mathrm{Z}_r^{pq}(\bb{X}^\bullet)$,
$\mathrm{B}_r^{pq}(\bb{X}^\bullet)$, y
\[
  \mathrm{E}_r^{pq}(\bb{X}^\bullet)
  =
  \mathrm{Z}_r^{pq}(\bb{X}^\bullet)/
  \mathrm{B}_r^{pq}(\bb{X}^\bullet).
\]
Los objetos $\mathrm{Z}_r^{pq}(\bb{X}^\bullet)$ y
$\mathrm{B}_r^{pq}(\bb{X}^\bullet)$ se identifican canónicamente con
subobjetos de $\mathrm{E}_1^{pq}(\bb{X}^\bullet)$.
La definición de los $d_r^{pq}$ (M, XV, 1) muestra que estos
morfismos, tomados respecto de los $X^\bullet_k$, son también esencialmente
constantes, y definen, por tanto, morfismos
\begin{equation}
\label{0.13.5.4.1}
\tag{13.5.4.1}
  d_r^{pq}:
  \mathrm{E}_r^{pq}(\bb{X}^\bullet)
  \longrightarrow
  \mathrm{E}_r^{p+r,q-r+1}(\bb{X}^\bullet)
\end{equation}
tal que
$d_r^{p+r,q-r+1}\circ d_r^{pq}=0$.
Además, existen isomorfismos canónicos de
$\Ker(d_r^{pq})$ sobre
\[
  \mathrm{Z}_{r+1}^{pq}(\bb{X}^\bullet)/
  \mathrm{B}_r^{pq}(\bb{X}^\bullet)
\]
y de $\Im(d_r^{pq})$ sobre
\[
  \mathrm{B}_{r+1}^{p+r,q-r+1}(\bb{X}^\bullet)/
  \mathrm{B}_r^{p+r,q-r+1}(\bb{X}^\bullet).
\]
\end{env}

\begin{lemma}[13.5.5]
\label{0.13.5.5}
Bajo las hipótesis de \sref{0.13.5.3}, para
$s\geq r>p-p_0$ existe un monomorfismo canónico
\begin{equation}
\label{0.13.5.5.1}
\tag{13.5.5.1}
  i:
  \mathrm{E}_s^{pq}(\bb{X}^\bullet)
  \longrightarrow
  \mathrm{E}_r^{pq}(\bb{X}^\bullet)
\end{equation}
y un isomorfismo canónico
\begin{equation}
\label{0.13.5.5.2}
\tag{13.5.5.2}
  j_r:
  \mathrm{E}_r^{pq}(\bb{X}^\bullet)
  \simto
  \mathrm{E}_\infty^{pq}(X^\bullet_{p+r-1})
\end{equation}
tal que el diagrama
\begin{equation}
\label{0.13.5.5.3}
\tag{13.5.5.3}
  \xymatrix{
    \mathrm{E}_s^{pq}(\bb{X}^\bullet)
      \ar[r]^{j_s}\ar[d]_i
    &
    \mathrm{E}_\infty^{pq}(X^\bullet_{p+s-1})
      \ar[d]
    \\
    \mathrm{E}_r^{pq}(\bb{X}^\bullet)
      \ar[r]^{j_r}
    &
    \mathrm{E}_\infty^{pq}(X^\bullet_{p+r-1})
  }
\end{equation}
es conmutativo, donde la flecha vertical derecha está inducida por el
morfismo
$X^\bullet_{p+s-1}\longrightarrow X^\bullet_{p+r-1}$.
\end{lemma}

\begin{proof}
La existencia de $i$ se sigue de
\[
  \mathrm{B}_r^{pq}(X^\bullet_k)
  =
  \mathrm{B}_\infty^{pq}(X^\bullet_k)
  \qquad (r>p-p_0)
\]
(\sref{0.13.5.3},~a)).
Para $k+1\leq p+r$, se tiene
\[
  \mathrm{Z}_r^{pq}(X^\bullet_k)
  =
  \mathrm{Z}_\infty^{pq}(X^\bullet_k)
\]
(\sref{0.13.5.3},~b)); en particular,
\[
  \mathrm{Z}_r^{pq}(X^\bullet_{p+r-1})
  =
  \mathrm{Z}_\infty^{pq}(X^\bullet_{p+r-1}).
\]
Por otra parte,
$\mathrm{Z}_r^{pq}(X^\bullet_{p+r-1})$ y
$\mathrm{B}_r^{pq}(X^\bullet_{p+r-1})$ se identifican canónicamente con
$\mathrm{Z}_r^{pq}(\bb{X}^\bullet)$ y
$\mathrm{B}_r^{pq}(\bb{X}^\bullet)$ por
\sref{0.13.5.3},~c).
Esto demuestra la existencia de $j_r$ y la conmutatividad de
\sref{0.13.5.5.3}.
\end{proof}

\begin{corollary}[13.5.6]
\label{0.13.5.6}
Bajo las hipótesis de \sref{0.13.5.3}, si uno de los límites
proyectivos
\[
  \varprojlim_r\mathrm{E}_r^{pq}(\bb{X}^\bullet),
  \qquad
  \varprojlim_k\mathrm{E}_\infty^{pq}(X^\bullet_k)
\]
existe, entonces también existe el otro, y hay un isomorfismo canónico
\begin{equation}
\label{0.13.5.6.1}
\tag{13.5.6.1}
  j_\infty:
  \varprojlim_r\mathrm{E}_r^{pq}(\bb{X}^\bullet)
  \simto
  \varprojlim_k\mathrm{E}_\infty^{pq}(X^\bullet_k).
\end{equation}
Además, el sistema proyectivo
$(\mathrm{E}_r^{pq}(\bb{X}^\bullet))_{r\in\bb{Z}}$ es esencialmente
constante (\sref{0.13.4.2}) si y solo si el sistema proyectivo
$(\mathrm{E}_\infty^{pq}(X^\bullet_k))_{k\in\bb{Z}}$ es esencialmente
constante.
\end{corollary}

\begin{env}[13.5.7]
\label{0.13.5.7}
Denotamos por
$\mathrm{B}_\infty^{pq}(\bb{X}^\bullet)$ y
$\mathrm{Z}_\infty^{pq}(\bb{X}^\bullet)$ los subobjetos de
$\mathrm{E}_1^{pq}(\bb{X}^\bullet)$ iguales respectivamente a
$\mathrm{B}_r^{pq}(\bb{X}^\bullet)$ para $r>p-p_0$ y a
$\inf_r\mathrm{Z}_r^{pq}(\bb{X}^\bullet)$, cuando este último existe.
Así,
\[
  \varprojlim_r\mathrm{E}_r^{pq}(\bb{X}^\bullet)
\]
\oldpage[0\textsubscript{III}]{73}
se identifica canónicamente con
\[
  \mathrm{E}_\infty^{pq}(\bb{X}^\bullet)
  =
  \mathrm{Z}_\infty^{pq}(\bb{X}^\bullet)/
  \mathrm{B}_\infty^{pq}(\bb{X}^\bullet).
\]
Observemos que los objetos
$\mathrm{Z}_r^{pq}(\bb{X}^\bullet)$,
$\mathrm{B}_r^{pq}(\bb{X}^\bullet)$,
$\mathrm{E}_r^{pq}(\bb{X}^\bullet)$
($1\leq r\leq+\infty$), y los morfismos $d_r^{pq}$ dependen
\emph{funtorialmente} del sistema proyectivo $\bb{X}^\bullet$, con las
restricciones de \sref{0.13.5.5}, y los morfismos definidos en
\sref{0.13.5.5} y \sref{0.13.5.6} son funtoriales.
\end{env}

\subsection{Sucesión espectral de un funtor relativa a un objeto dotado de una filtración finita}
\label{subsection:0.13.6}

\begin{env}[13.6.1]
\label{0.13.6.1}
Sean $\cat{C}$ y $\cat{C}'$ dos categorías abelianas, y sea
$\mathrm{T}:\cat{C}\longrightarrow\cat{C}'$ un funtor covariante
aditivo.
Supongamos que todo objeto de $\cat{C}$ es isomorfo a un subobjeto de
un objeto inyectivo, de modo que existan los funtores derivados por la derecha
$\RR^p\mathrm{T}$ ($p\geq0$).
\end{env}

\begin{lemma}[13.6.2]
\label{0.13.6.2}
Sea $A$ un objeto de $\cat{C}$ dotado de una filtración finita
$(\mathrm{F}^i(A))_{i\in\bb{Z}}$.
Existe una resolución inyectiva
$X^\bullet=(X^j)_{j\geq0}$ de $A$, dotada de una filtración finita
$(\mathrm{F}^i(X^\bullet))_{i\in\bb{Z}}$, tal que
\[
  \mathrm{F}^i(A)=A
  \quad\text{(respectivamente, }\mathrm{F}^i(A)=0\text{)}
\]
implica
\[
  \mathrm{F}^i(X^\bullet)=X^\bullet
  \quad\text{(respectivamente, }\mathrm{F}^i(X^\bullet)=0\text{)},
\]
y tal que, para todo $i\in\bb{Z}$,
$\mathrm{F}^i(X^\bullet)$ es una resolución inyectiva de
$\mathrm{F}^i(A)$.
\end{lemma}

\begin{proof}
Sea $p$ (respectivamente, $q>p$) el mayor índice tal que
$\mathrm{F}^p(A)=A$ (respectivamente, el menor índice tal que
$\mathrm{F}^q(A)=0$).
Razonamos por inducción sobre $q-p$, siendo inmediato el lema cuando
$q-p=1$.
Supongamos construida una resolución inyectiva ${X''}^\bullet$ de
$A/\mathrm{F}^{q-1}(A)$ que tenga las propiedades requeridas.
Consideremos la sucesión exacta
\[
  0\longrightarrow\mathrm{F}^{q-1}(A)
  \longrightarrow A
  \longrightarrow A/\mathrm{F}^{q-1}(A)
  \longrightarrow0,
\]
tomemos una resolución inyectiva ${X'}^\bullet$ de
$\mathrm{F}^{q-1}(A)$, y elijamos después una resolución inyectiva
$X^\bullet$ de $A$ de modo que se obtenga una sucesión exacta
\[
  0\longrightarrow {X'}^\bullet
  \longrightarrow X^\bullet
  \longrightarrow {X''}^\bullet
  \longrightarrow0
\]
compatible con la sucesión precedente (M, V, 2.2).
Está claro que $X^\bullet$ tiene las propiedades requeridas.
\end{proof}

\begin{corollary}[13.6.3]
\label{0.13.6.3}
Sea $B$ un segundo objeto de $\cat{C}$, dotado de una filtración
finita $(\mathrm{F}^i(B))_{i\in\bb{Z}}$, sea $s$ un entero,
y sea $u:A\longrightarrow B$ un morfismo tal que
\[
  u(\mathrm{F}^i(A))\subset\mathrm{F}^{i+s}(B)
\]
para todo $i\in\bb{Z}$.
Si $Y^\bullet=(Y^j)_{j\geq0}$ es una resolución inyectiva de $B$
dotada de una filtración
$(\mathrm{F}^i(Y^\bullet))_{i\in\bb{Z}}$ que tenga las propiedades
enunciadas en \sref{0.13.6.2}, entonces existe un morfismo
$v:X^\bullet\longrightarrow Y^\bullet$, compatible con $u$, tal
que
\[
  v(\mathrm{F}^i(X^\bullet))
  \subset
  \mathrm{F}^{i+s}(Y^\bullet)
\]
para todo $i\in\bb{Z}$.
Además, cualesquiera dos de tales morfismos $v$ y $v'$ son homotópicos.
\end{corollary}

Esto se sigue inmediatamente, por inducción sobre $q-p$, de la construcción
precedente y de (M, V, 2.3).

\begin{env}[13.6.4]
\label{0.13.6.4}
Bajo las hipótesis de \sref{0.13.6.2}, consideremos el complejo
$\mathrm{T}(X^\bullet)$ en $\cat{C}'$.
Está filtrado por los complejos
$\mathrm{T}(\mathrm{F}^i(X^\bullet))$, pues
$\mathrm{F}^i(X^\bullet)$ es un sumando directo de $X^\bullet$.
Se sigue de \sref{0.13.6.3} que la \emph{sucesión espectral} de
este complejo filtrado depende, salvo isomorfismo, únicamente del
objeto filtrado $A$.
Su término límite es la cohomología $\RR^\bullet\mathrm{T}(A)$, dotada
de la filtración
\begin{equation}
\label{0.13.6.4.1}
\tag{13.6.4.1}
  \begin{split}
  \mathrm{F}^p(\RR^n\mathrm{T}(A))
  &=
  \Im\bigl(
    \RR^n\mathrm{T}(\mathrm{F}^p(A))
    \longrightarrow
    \RR^n\mathrm{T}(A)
  \bigr)
  \\
  &=
  \Ker\bigl(
    \RR^n\mathrm{T}(A)
    \longrightarrow
    \RR^n\mathrm{T}(A/\mathrm{F}^p(A))
  \bigr),
  \end{split}
\end{equation}
por \sref{0.11.2.2}, y su término $\mathrm{E}_1$ es
\begin{equation}
\label{0.13.6.4.2}
\tag{13.6.4.2}
  \mathrm{E}_1^{pq}
  =
  \RR^{p+q}\mathrm{T}(\mathrm{gr}^p(A)),
\end{equation}
donde, como de costumbre,
$\mathrm{gr}^p(A)=\mathrm{F}^p(A)/\mathrm{F}^{p+1}(A)$.
Se desprende de \sref{0.11.2.2} que la filtración del término límite
es finita y que, para $p$ y $q$ fijos, las sucesiones
\oldpage[0\textsubscript{III}]{74}
\[
  \mathrm{B}_r^{pq}(A)
  =
  \mathrm{B}_r^{pq}(\mathrm{T}(X^\bullet))
  \quad\text{y}\quad
  \mathrm{Z}_r^{pq}(A)
  =
  \mathrm{Z}_r^{pq}(\mathrm{T}(X^\bullet))
\]
son estacionarias.
La sucesión espectral precedente es, por tanto, \emph{birregular}
(\sref{0.11.1.3}).
La denotamos por
$\mathrm{E}(A)=(\mathrm{E}_r^{pq}(A))$ y la llamamos la
\emph{sucesión espectral del funtor $\mathrm{T}$ relativa al
objeto filtrado $A$}.
\end{env}

\begin{env}[13.6.5]
\label{0.13.6.5}
Supongamos ahora que se cumplen las hipótesis de \sref{0.13.6.3}, y
conservemos su notación.
Como $\mathrm{F}^i(X^\bullet)$ (respectivamente,
$\mathrm{F}^i(Y^\bullet)$) es un sumando directo de $X^\bullet$
(respectivamente, $Y^\bullet$), se tiene
\[
  (\mathrm{T}v)\bigl(\mathrm{T}(\mathrm{F}^i(X^\bullet))\bigr)
  \subset
  \mathrm{T}(\mathrm{F}^{i+s}(Y^\bullet))
\]
para todo $i\in\bb{Z}$.
Las definiciones de \sref{0.11.2.2} muestran entonces que, para
$1\leq r\leq+\infty$, $\mathrm{T}v$ define morfismos
\[
  \mathrm{B}_r^{pq}(\mathrm{T}(X^\bullet))
  \longrightarrow
  \mathrm{B}_r^{p+s,q-s}(\mathrm{T}(Y^\bullet))
\]
y
\[
  \mathrm{Z}_r^{pq}(\mathrm{T}(X^\bullet))
  \longrightarrow
  \mathrm{Z}_r^{p+s,q-s}(\mathrm{T}(Y^\bullet)),
\]
y, por tanto, un morfismo
\[
  w_r:
  \mathrm{E}_r^{pq}(A)
  \longrightarrow
  \mathrm{E}_r^{p+s,q-s}(B).
\]
Análogamente, sobre los términos límite se tienen morfismos
\[
  u_n:
  \RR^n\mathrm{T}(A)
  \longrightarrow
  \RR^n\mathrm{T}(B)
\]
tales que
\[
  u_n\bigl(\mathrm{F}^p(\RR^n\mathrm{T}(A))\bigr)
  \subset
  \mathrm{F}^{p+s}(\RR^n\mathrm{T}(B)).
\]
La definición de los $d_r^{pq}$ (M, XV, 1) muestra también que los
diagramas
\[
  \xymatrix{
    \mathrm{E}_r^{pq}(A)
      \ar[r]^{d_r^{pq}}\ar[d]_{w_r}
    &
    \mathrm{E}_r^{p+r,q-r+1}(A)
      \ar[d]^{w_r}
    \\
    \mathrm{E}_r^{p+s,q-s}(B)
      \ar[r]_{d_r^{p+s,q-s}}
    &
    \mathrm{E}_r^{p+r+s,q-r-s+1}(B)
  }
\]
son conmutativos.
Se sigue que existe un diagrama conmutativo análogo para los
isomorfismos $\alpha_r^{pq}$, cuya explicitación dejamos al lector.
Finalmente (\emph{loc.\ cit.}), se tienen también diagramas conmutativos
sobre los términos límite:
\[
  \xymatrix{
    \mathrm{E}_\infty^{pq}(A)
      \ar[r]^{\beta^{pq}}\ar[d]_{w_\infty}
    &
    \mathrm{gr}^p(\RR^{p+q}\mathrm{T}(A))
      \ar[d]^{u_{p+q}}
    \\
    \mathrm{E}_\infty^{p+s,q-s}(B)
      \ar[r]_{\beta^{p+s,q-s}}
    &
    \mathrm{gr}^{p+s}(\RR^{p+q}\mathrm{T}(B)).
  }
\]
\end{env}

\begin{env}[13.6.6]
\label{0.13.6.6}
Supongamos, en particular, que existe un anillo $S$, dotado de una
\emph{filtración} $(\mathrm{F}^i(S))_{i\in\bb{Z}}$, y un homomorfismo
de anillos
\begin{equation}
\label{0.13.6.6.1}
\tag{13.6.6.1}
  h:S\longrightarrow\Hom_{\cat{C}}(A,A)
\end{equation}
\oldpage[0\textsubscript{III}]{75}
tal que, para todo $t\in\mathrm{F}^i(S)$,
\[
  h_t(\mathrm{F}^j(A))
  \subset
  \mathrm{F}^{i+j}(A)
\]
para todo par $i,j$.
Para abreviar, decimos entonces que $A$ está dotado de la estructura de un
\emph{$S$--$\cat{C}$-módulo filtrado} sobre el anillo filtrado $S$.

Al pasar a los objetos graduados asociados, todo $h_t$, para
$t\in\mathrm{F}^i(S)$, define un endomorfismo graduado
$\overline{h}_t$ de $\mathrm{gr}^\bullet(A)$, homogéneo de grado
$i$.
Además, este endomorfismo depende únicamente de la clase de $t$ en
$\mathrm{gr}^i(S)$, y se obtiene así un homomorfismo de
\emph{anillos graduados}
\[
  \overline{h}:
  \mathrm{gr}^\bullet(S)
  \longrightarrow
  \Hom_{\cat{C}}^\bullet
  \bigl(\mathrm{gr}^\bullet(A),\mathrm{gr}^\bullet(A)\bigr),
\]
donde el miembro derecho es el anillo de endomorfismos graduados de
$\mathrm{gr}^\bullet(A)$.
Decimos que $\mathrm{gr}^\bullet(A)$ está dotado de la estructura
de un \emph{$\mathrm{gr}^\bullet(S)$--$\cat{C}$-módulo graduado}.

Se sigue de \sref{0.13.6.5} que, para
$1\leq r\leq+\infty$, todo
$\overline{t}\in\mathrm{gr}^i(S)$ define canónicamente, sobre los
objetos bigraduados
\[
  \bigl(\mathrm{B}_r^{pq}(A)\bigr)_{(p,q)\in\bb{Z}\times\bb{Z}},
  \qquad
  \bigl(\mathrm{Z}_r^{pq}(A)\bigr)_{(p,q)\in\bb{Z}\times\bb{Z}},
\]
y
\[
  \mathrm{E}_r(A)
  =
  \bigl(\mathrm{E}_r^{pq}(A)\bigr)_{(p,q)\in\bb{Z}\times\bb{Z}},
\]
endomorfismos bigraduados de grado $(i,-i)$.
Sobre $\mathrm{E}_r(A)$, para $r$ finito, este endomorfismo conmuta con
el endomorfismo bigraduado definido por los $d_r^{pq}$.
Estos endomorfismos satisfacen las condiciones usuales de asociatividad y
distributividad respecto de la adición en
$\mathrm{gr}^\bullet(S)$ y en los objetos bigraduados en cuestión.
Para abreviar, decimos que estos últimos son
\emph{$\mathrm{gr}^\bullet(S)$--$\cat{C}'$-módulos bigraduados};
es inmediato que los $\alpha_r^{pq}$ definen un isomorfismo para
este tipo de estructura.

Para todo entero $n$, sea $\mathrm{B}_r^{(n)}(A)$
(respectivamente, $\mathrm{Z}_r^{(n)}(A)$ y
$\mathrm{E}_r^{(n)}(A)$) el subobjeto graduado de
$\mathrm{B}_r^{\bullet,\bullet}(A)$
(respectivamente, $\mathrm{Z}_r^{\bullet,\bullet}(A)$ y
$\mathrm{E}_r^{\bullet,\bullet}(A)$) constituido por los términos para
los que $p+q=n$ ($1\leq r\leq+\infty$).
Es inmediato que estos son
\emph{$\mathrm{gr}^\bullet(S)$--$\cat{C}'$-módulos graduados}.

Finalmente, todo $\overline{t}\in\mathrm{gr}^i(S)$ define, para todo
$n$, un endomorfismo graduado de grado $i$ sobre el objeto graduado
$\mathrm{gr}^\bullet(\RR^n\mathrm{T}(A))$, que adquiere así la
estructura de un
$\mathrm{gr}^\bullet(S)$--$\cat{C}'$-módulo graduado.
Para $p+q=n$, los $\beta^{pq}$ definen un isomorfismo
\[
  \mathrm{E}^{(n)}(A)
  \isoto
  \mathrm{gr}^\bullet(\RR^n\mathrm{T}(A))
\]
para este tipo de estructura.

Observemos que, cuando $\cat{C}'$ es la categoría de grupos abelianos, las
estructuras de $S$--$\cat{C}'$-módulo (respectivamente, de
$\mathrm{gr}^\bullet(S)$--$\cat{C}'$-módulo graduado o bigraduado) son
simplemente las estructuras usuales de $S$-módulo (respectivamente, de
$\mathrm{gr}^\bullet(S)$-módulo graduado o bigraduado).
\end{env}

\subsection{Funtores derivados de un límite proyectivo de argumentos}
\label{subsection:0.13.7}
\label{0.13.7}

\begin{env}[13.7.1]
\label{0.13.7.1}
Sean $\cat{C}$ y $\cat{C}'$ dos categorías abelianas, y supongamos
que todo objeto de $\cat{C}$ es isomorfo a un subobjeto de un
objeto inyectivo.
Sea
\[
  \mathrm{T}:\cat{C}\longrightarrow\cat{C}'
\]
un funtor covariante aditivo.
Consideremos un sistema proyectivo estricto
$\bb{A}=(A_k)_{k\in\bb{Z}}$ en $\cat{C}$, acotado inferiormente; más
precisamente, supongamos que $A_k=0$ para $k<k_0$.
A este sistema asociamos canónicamente, sobre cada $A_k$, la
filtración $(\mathrm{F}^p(A_k))_{p\in\bb{Z}}$ dada por
\sref{0.13.4.1.1}.
Como el sistema proyectivo es estricto, los morfismos canónicos
\begin{equation}
\label{0.13.7.1.1}
\tag{13.7.1.1}
  \mathrm{F}^i(A_h)/\mathrm{F}^j(A_h)
  \longrightarrow
  \mathrm{F}^i(A_k)/\mathrm{F}^j(A_k)
  \qquad(h\geq k)
\end{equation}
son isomorfismos para $i\leq j\leq k+1$.
Recordemos también que
$\mathrm{F}^{k_0}(A_k)=A_k$ y
$\mathrm{F}^{k+1}(A_k)=0$ para todo $k$.
\end{env}

\begin{env}[13.7.2]
\label{0.13.7.2}
Para cada $k$, construyamos una resolución inyectiva
$\mathrm{X}_k^\bullet=(\mathrm{X}_k^j)_{j\geq0}$ de $A_k$ que tenga las
propiedades de \sref{0.13.6.2}.
Los morfismos canónicos $A_{k+1}\longrightarrow A_k$ nos permiten, por
\sref{0.13.6.3}, definir para cada $k$ un morfismo de complejos
\[
  \mathrm{X}_{k+1}^\bullet\longrightarrow\mathrm{X}_k^\bullet
\]
compatible con las filtraciones y que hace de
$\bb{X}^\bullet=(\mathrm{X}_k^\bullet)_{k\in\bb{Z}}$ un sistema
proyectivo de complejos.
\oldpage[0\textsubscript{III}]{76}
Podemos suponer además que este sistema proyectivo es estricto.
En efecto, por \sref{0.13.7.1.1},
\[
  A_k
  \simeq
  A_{k+1}/\mathrm{F}^{k+1}(A_{k+1}).
\]
Así, al construir $\mathrm{X}_{k+1}^\bullet$, podemos tomar la
resolución inyectiva de
$A_{k+1}/\mathrm{F}^{k+1}(A_{k+1})$ \emph{igual} a
$\mathrm{X}_k^\bullet$.
Se sigue entonces de (M, V, 2.3) que el morfismo de complejos
$\mathrm{X}_{k+1}^\bullet\longrightarrow\mathrm{X}_k^\bullet$ puede
construirse de modo que, al pasar a los cocientes, induzca un
isomorfismo
\[
  \mathrm{X}_{k+1}^\bullet/
  \mathrm{F}^{k+1}(\mathrm{X}_{k+1}^\bullet)
  \simto
  \mathrm{X}_k^\bullet
\]
que respeta las filtraciones; esta es precisamente la condición de
\sref{0.13.5.1}.
\end{env}

\begin{env}[13.7.3]
\label{0.13.7.3}
Por construcción, el sistema proyectivo
$\mathrm{T}(\bb{X}^\bullet)$ de complejos
$\mathrm{T}(\mathrm{X}_k^\bullet)$ satisface las hipótesis de
\sref{0.13.5.3}.
Los resultados de \sref{0.13.5.4}, \sref{0.13.5.5} y
\sref{0.13.5.6} se aplican, por tanto, a las sucesiones espectrales
\[
  \mathrm{E}\bigl(\mathrm{T}(\mathrm{X}_k^\bullet)\bigr)
  =
  \mathrm{E}(A_k).
\]
Para $1\leq r\leq+\infty$, escribimos
$\mathrm{E}_r^{pq}(\bb{A})$ en lugar de
$\mathrm{E}_r^{pq}(\mathrm{T}(\bb{X}^\bullet))$
(cf.\ \sref{0.13.5.7} cuando $r=+\infty$), y utilizamos la misma convención
para la notación análoga.
En particular,
\begin{equation}
\label{0.13.7.3.1}
\tag{13.7.3.1}
  \mathrm{E}_1^{pq}(\bb{A})
  =
  \RR^{p+q}\mathrm{T}\bigl(\mathrm{gr}^p(\bb{A})\bigr),
\end{equation}
por \sref{0.13.6.4.2} y el hecho de que el sistema proyectivo
$(\mathrm{gr}^p(A_k))_{k\in\bb{Z}}$ es esencialmente constante.
Estos resultados, junto con \sref{0.13.4.3}, dan la proposición
siguiente, demostrada por primera vez por Shih Weishu mediante un método
diferente (inédito).
\end{env}

\begin{proposition}[13.7.4]
\label{0.13.7.4}
\emph{(Shih).}
Sea $n$ un entero.
Las dos condiciones siguientes son equivalentes:
\begin{enumerate}
  \item[{\rm a)}] Para todo par $(p,q)$ tal que $p+q=n$, el
  sistema proyectivo
  $(\mathrm{E}_r^{pq}(\bb{A}))_{r\geq2}$ es esencialmente constante.
  \item[{\rm b)}] El sistema proyectivo
  \[
    \RR^n\mathrm{T}(\bb{A})
    =
    \bigl(\RR^n\mathrm{T}(A_k)\bigr)_{k\in\bb{Z}}
  \]
  satisface {\rm(ML)}.
\end{enumerate}
Además, cuando se cumplen estas condiciones, existe un
isomorfismo canónico
\begin{equation}
\label{0.13.7.4.1}
\tag{13.7.4.1}
  \mathrm{gr}^p\bigl(\RR^n\mathrm{T}(\bb{A})\bigr)
  \isoto
  \mathrm{E}_\infty^{p,n-p}(\bb{A})
  \qquad(p\in\bb{Z}).
\end{equation}
\end{proposition}

\begin{proof}
Por \sref{0.13.5.6}, la condición {\rm a)} dice precisamente que, para
$p+q=n$, el sistema proyectivo
$(\mathrm{E}_\infty^{pq}(A_k))_{k\in\bb{Z}}$ es esencialmente
constante.
Por otra parte,
$\mathrm{gr}^p(\RR^n\mathrm{T}(A_k))$ es canónicamente isomorfo a
$\mathrm{E}_\infty^{p,n-p}(A_k)$.
Se sigue de \sref{0.13.4.3} que {\rm a)} y {\rm b)} son
equivalentes, y \sref{0.13.7.4.1} no es sino
\sref{0.13.5.6.1} en el caso considerado.
\end{proof}

\begin{corollary}[13.7.5]
\label{0.13.7.5}
Sea $\sh{F}=(\sh{F}_k)_{k\in\bb{N}}$ un sistema proyectivo de
haces de grupos abelianos que satisface las condiciones {\rm(i)}, {\rm(ii)}
y {\rm(iii)} de \sref{0.13.3.1}, y pongamos
$\sh{F}=\varprojlim_k\sh{F}_k$.
Supongamos que, para el funtor
$\sh{G}\longmapsto\Gamma(X,\sh{G})$, el sistema proyectivo
$(\mathrm{E}_r^{pq}(\sh{F}))_r$ es esencialmente constante siempre que
$p+q=n$ o $p+q=n+1$.
Sobre $\HH^{n+1}(X,\sh{F})$, consideremos la filtración
\[
  \mathrm{F}^p\bigl(\HH^{n+1}(X,\sh{F})\bigr)
  =
  \Ker\bigl(
    \HH^{n+1}(X,\sh{F})
    \longrightarrow
    \HH^{n+1}(X,\sh{F}_{p-1})
  \bigr).
\]
Entonces existe un isomorfismo canónico
\begin{equation}
\label{0.13.7.5.1}
\tag{13.7.5.1}
  \mathrm{gr}^p\bigl(\HH^{n+1}(X,\sh{F})\bigr)
  \isoto
  \mathrm{E}_\infty^{p,n-p+1}(\sh{F})
  \qquad(p\in\bb{Z}).
\end{equation}
\end{corollary}

\begin{proof}
\oldpage[0\textsubscript{III}]{77}
Apliquemos \sref{0.13.7.4} tomando como $\cat{C}$ la categoría de haces de
grupos abelianos sobre $X$, como $\cat{C}'$ la categoría de grupos abelianos,
y $\mathrm{T}=\Gamma$.
Da, para todo $p\in\bb{Z}$, un isomorfismo canónico
\[
  \mathrm{gr}^p\bigl(\RR^{n+1}\Gamma(\sh{F})\bigr)
  \isoto
  \mathrm{E}_\infty^{p,n-p+1}(\sh{F}).
\]
Además, por \sref{0.13.7.4}, el sistema proyectivo
$(\HH^n(X,\sh{F}_k))_{k\in\bb{Z}}$ satisface {\rm(ML)}; por tanto
\sref{0.13.3.1} da un isomorfismo canónico
\[
\label{0.13.7.5.1b}
  \HH^{n+1}(X,\sh{F})
  \isoto
  \varprojlim_k\HH^{n+1}(X,\sh{F}_k).
\]

Como el sistema proyectivo
$\RR^{n+1}\Gamma(\sh{F})$ satisface {\rm(ML)} por
\sref{0.13.7.4}, existe un isomorfismo canónico
\[
  \mathrm{gr}^p\bigl(\RR^{n+1}\Gamma(\sh{F})\bigr)
  \isoto
  \varprojlim_k
  \mathrm{gr}^p\bigl(\HH^{n+1}(X,\sh{F}_k)\bigr)
\]
por \sref{0.13.4.3}, y un isomorfismo canónico
\[
  \varprojlim_k
  \mathrm{gr}^p\bigl(\HH^{n+1}(X,\sh{F}_k)\bigr)
  \isoto
  \mathrm{gr}^p\left(
    \varprojlim_k\HH^{n+1}(X,\sh{F}_k)
  \right)
\]
por \sref{0.13.4.5}.
Solo queda comprobar que el isomorfismo precedente es
compatible con las filtraciones de sus dos lados.
Esto se sigue inmediatamente de las definiciones y, para todo $p$,
de la conmutatividad de
\[
  \xymatrix{
    \HH^{n+1}(X,\sh{F})
      \ar[r]^-{\sim}\ar[dr]
    &
    \varprojlim_k\HH^{n+1}(X,\sh{F}_k)
      \ar[d]
    \\
    &
    \HH^{n+1}(X,\sh{F}_{p-1}).
  }
\]
\end{proof}

\begin{env}[13.7.6]
\label{0.13.7.6}
Sea $S$ un anillo dotado de una filtración
$(\mathrm{F}^i(S))_{i\in\bb{Z}}$ tal que
$\mathrm{F}^0(S)=S$ (y por tanto
$\mathrm{gr}^i(S)=0$ para $i<0$).
Supongamos que cada $A_k$, con la filtración definida en
\sref{0.13.7.1}, es un
$S$--$\cat{C}$-módulo filtrado en el sentido de \sref{0.13.6.6}, y que los
morfismos $A_h\longrightarrow A_k$ para $k\leq h$ son morfismos de
$S$--$\cat{C}$-módulos filtrados.
Para abreviar, decimos entonces que $\bb{A}$ es un
\emph{sistema proyectivo de $S$--$\cat{C}$-módulos filtrados}.

Es inmediato que, para $k\leq h$ y
$1\leq r\leq+\infty$, los morfismos
\[
  \mathrm{B}_r^{pq}(A_h)\longrightarrow\mathrm{B}_r^{pq}(A_k),
  \qquad
  \mathrm{Z}_r^{pq}(A_h)\longrightarrow\mathrm{Z}_r^{pq}(A_k)
\]
son morfismos de
$\mathrm{gr}^\bullet(S)$--$\cat{C}'$-módulos bigraduados
(\sref{0.13.6.5}), y que, para $r$ finito, las familias
\[
  \bigl(\mathrm{Z}_r^{pq}(\bb{A})\bigr),\qquad
  \bigl(\mathrm{B}_r^{pq}(\bb{A})\bigr),\qquad
  \bigl(\mathrm{E}_r^{pq}(\bb{A})\bigr)
\]
son $\mathrm{gr}^\bullet(S)$--$\cat{C}'$-módulos
bigraduados, siendo los dos primeros submódulos de
$(\mathrm{E}_1^{pq}(\bb{A}))$.
Escribimos de nuevo
$\mathrm{Z}_r^{(n)}(\bb{A})$,
$\mathrm{B}_r^{(n)}(\bb{A})$ y
$\mathrm{E}_r^{(n)}(\bb{A})$ para los respectivos subobjetos graduados
obtenidos conservando únicamente los términos con $p+q=n$; estos son
$\mathrm{gr}^\bullet(S)$--$\cat{C}'$-módulos graduados.

Cuando el sistema proyectivo
$(\mathrm{E}_r^{pq}(\bb{A}))_r$ es esencialmente constante,
$\mathrm{E}_\infty^{pq}(\bb{A})$ es asimismo un
$\mathrm{gr}^\bullet(S)$--$\cat{C}'$-módulo bigraduado, y cada
$\mathrm{E}_\infty^{(n)}(\bb{A})$ es un
$\mathrm{gr}^\bullet(S)$--$\cat{C}'$-módulo graduado.
Además, para cada $k$, los isomorfismos
\[
  \beta^{p,n-p}:
  \mathrm{E}_\infty^{p,n-p}(A_k)
  \isoto
  \mathrm{gr}^p\bigl(\RR^n\mathrm{T}(A_k)\bigr)
\]
forman un isomorfismo de
$\mathrm{gr}^\bullet(S)$--$\cat{C}'$-módulos graduados de
$\mathrm{E}_\infty^{(n)}(A_k)$ sobre
$\mathrm{gr}^\bullet(\RR^n\mathrm{T}(A_k))$.
Bajo las hipótesis precedentes, el isomorfismo
\[
  \beta^{p,n-p}:
  \mathrm{E}_\infty^{(n)}(\bb{A})
  \isoto
  \varprojlim_k
  \mathrm{gr}^\bullet\bigl(\RR^n\mathrm{T}(A_k)\bigr)
\]
respeta, por tanto, también estas estructuras.
Lo mismo vale manifiestamente para el isomorfismo canónico
\[
  \mathrm{gr}^\bullet\bigl(\RR^n\mathrm{T}(\bb{A})\bigr)
  \isoto
  \varprojlim_k
  \mathrm{gr}^\bullet\bigl(\RR^n\mathrm{T}(A_k)\bigr);
\]
por consiguiente, los isomorfismos \sref{0.13.7.4.1} forman un
isomorfismo de
$\mathrm{gr}^\bullet(S)$--$\cat{C}'$-módulos graduados.
\end{env}

\begin{proposition}[13.7.7]
\label{0.13.7.7}
Sea $S$ un anillo noetheriano $\mathfrak{J}$-ádico.
Supongamos que $\cat{C}$ es una categoría abeliana en la que todo objeto
es isomorfo a un subobjeto de un objeto inyectivo, y sea
$\mathrm{T}$ un funtor covariante aditivo de $\cat{C}$ en la
categoría de grupos abelianos.
Sea $\bb{A}=(A_k)_{k\in\bb{Z}}$ un
\oldpage[0\textsubscript{III}]{78}
sistema proyectivo estricto de $S$--$\cat{C}$-módulos filtrados, para la
filtración $\mathfrak{J}$-ádica sobre $S$, acotado inferiormente.
Para un entero fijo $n$, supongamos que se cumple la condición siguiente:
\begin{itemize}
  \item[$(\mathrm{F}_n)$]
  \emph{El $\mathrm{gr}^\bullet(S)$-módulo graduado}
  \[
    \mathrm{E}_1^{(m)}(\bb{A})
    =
    \bigl(
      \RR^m\mathrm{T}(\mathrm{gr}^p(\bb{A}))
    \bigr)_{p\in\bb{Z}}
  \]
  \emph{de \sref{0.13.7.3.1} es de tipo finito para
  $m=n$ y $m=n+1$.}
\end{itemize}
Entonces:
\begin{enumerate}
  \item[{\rm(i)}] Los sistemas proyectivos
  $(\RR^n\mathrm{T}(A_k))_{k\in\bb{Z}}$ y
  $(\RR^{n+1}\mathrm{T}(A_k))_{k\in\bb{Z}}$ satisfacen {\rm(ML)}.
  \item[{\rm(ii)}] Si
  \[
    \RR'^{\,n}\mathrm{T}(\bb{A})
    =
    \varprojlim_k\RR^n\mathrm{T}(A_k),
  \]
  entonces $\RR'^{\,n}\mathrm{T}(\bb{A})$ es un $S$-módulo de tipo
  finito.
  \item[{\rm(iii)}] La filtración sobre
  $\RR'^{\,n}\mathrm{T}(\bb{A})$ definida por
  \[
    \mathrm{F}^p\bigl(\RR'^{\,n}\mathrm{T}(\bb{A})\bigr)
    =
    \Ker\bigl(
      \RR'^{\,n}\mathrm{T}(\bb{A})
      \longrightarrow
      \RR^n\mathrm{T}(A_{p-1})
    \bigr)
    \qquad(p\in\bb{Z})
  \]
  es $\mathfrak{J}$-buena; es decir,
  \[
    \mathfrak{J}\,
    \mathrm{F}^p\bigl(\RR'^{\,n}\mathrm{T}(\bb{A})\bigr)
    \subset
    \mathrm{F}^{p+1}\bigl(\RR'^{\,n}\mathrm{T}(\bb{A})\bigr)
  \]
  para todo $p$, con igualdad para todo $p$ suficientemente grande.
  En particular, la topología sobre
  $\RR'^{\,n}\mathrm{T}(\bb{A})$ definida por esta filtración es la
  topología $\mathfrak{J}$-ádica.
  \item[{\rm(iv)}] El sistema proyectivo
  $(\mathrm{E}_r^{pq}(\bb{A}))_r$ es esencialmente constante siempre que
  $p+q=n$ o $p+q=n+1$.
  Así, $\mathrm{E}_\infty^{pq}(\bb{A})$ está definido
  (\sref{0.13.5.7}), y los isomorfismos
  \begin{equation}
  \label{0.13.7.7.1}
  \tag{13.7.7.1}
    \mathrm{gr}^p\bigl(\RR'^{\,n}\mathrm{T}(\bb{A})\bigr)
    \isoto
    \mathrm{E}_\infty^{p,n-p}(\bb{A})
    \qquad(p\in\bb{Z})
  \end{equation}
  forman un isomorfismo de
  $\mathrm{gr}^\bullet(S)$-módulos graduados.
\end{enumerate}
\end{proposition}

\begin{proof}
En vista de \sref{0.13.7.7.1}, y teniendo en cuenta los
isomorfismos \sref{0.13.7.4.1}, por abuso de notación escribiremos también
$\RR^n\mathrm{T}(\bb{A})$ para el límite proyectivo
$\RR'^{\,n}\mathrm{T}(\bb{A})$ del sistema proyectivo
$(\RR^n\mathrm{T}(A_k))$.

El anillo graduado $\mathrm{gr}^\bullet(S)$ es noetheriano
(Bourbaki, \emph{Alg.\ comm.}, cap.~III, \textsection2,
n\textsuperscript{o}~9, cor.~5 del teorema~2).
Por tanto, la sucesión creciente de
$\mathrm{gr}^\bullet(S)$-submódulos graduados
$\mathrm{B}_r^{(m)}(\bb{A})$ de
$\mathrm{E}_1^{(m)}(\bb{A})$ es estacionaria para $m=n$ y
$m=n+1$.
La condición {\rm(b)} de \sref{0.11.1.10} queda, por tanto, satisfecha.
Se sigue que la condición {\rm(a)} de \sref{0.13.7.4} se cumple para
$n$ y para $n+1$, lo que demuestra {\rm(i)}.

Además, los isomorfismos \sref{0.13.7.4.1}, junto con las
observaciones de \sref{0.13.7.6}, muestran que
$\mathrm{gr}^\bullet(\RR^n\mathrm{T}(\bb{A}))$ es un
$\mathrm{gr}^\bullet(S)$-módulo graduado isomorfo a
\[
  \mathrm{E}_\infty^{(n)}(\bb{A})
  =
  \mathrm{Z}_\infty^{(n)}(\bb{A})
  /
  \mathrm{B}_\infty^{(n)}(\bb{A}).
\]
Como $\mathrm{Z}_\infty^{(n)}(\bb{A})$ es un submódulo de
$\mathrm{E}_1^{(n)}(\bb{A})$, es de tipo finito, y también lo es
$\mathrm{E}_\infty^{(n)}(\bb{A})$.
Finalmente, para la filtración
$(\mathrm{F}^p(\RR'^{\,n}\mathrm{T}(\bb{A})))$, se sigue de
\sref{0.13.4.5} que
\[
  \mathrm{gr}^\bullet\bigl(\RR^n\mathrm{T}(\bb{A})\bigr)
  \quad\hbox{y}\quad
  \mathrm{gr}^\bullet\bigl(\RR'^{\,n}\mathrm{T}(\bb{A})\bigr)
\]
son $\mathrm{gr}^\bullet(S)$-módulos graduados isomorfos.
Esto demuestra {\rm(iv)}.
Las aserciones {\rm(ii)} y {\rm(iii)} se seguirán de estos resultados
y del lema siguiente.
\end{proof}

\begin{lemma}[13.7.7.2]
\label{0.13.7.7.2}
Sea $S$ un anillo noetheriano $\mathfrak{J}$-ádico, y sea $M$
un $S$-módulo dotado de una filtración codiscreta
$(\mathrm{F}^p(M))_{p\in\bb{Z}}$ tal que
\[
  \mathfrak{J}\mathrm{F}^p(M)
  \subset
  \mathrm{F}^{p+1}(M).
\]
Así, $M$ es un módulo filtrado sobre $S$ dotado de su
filtración $\mathfrak{J}$-ádica.
Supongamos además que $M$ es separado para la topología definida por
$(\mathrm{F}^p(M))$.
Entonces las condiciones siguientes son equivalentes:
\begin{enumerate}
  \item[{\rm a)}] $M$ es un $S$-módulo de tipo finito y
  $(\mathrm{F}^p(M))$ es una filtración $\mathfrak{J}$-buena.
  \item[{\rm b)}] $\mathrm{gr}^\bullet(M)$ es un
  $\mathrm{gr}^\bullet(S)$-módulo de tipo finito.
  \item[{\rm c)}] Los $\mathrm{gr}^p(M)$ son $S$-módulos de
  tipo finito y, para todo $p$ suficientemente grande, los homomorfismos
  canónicos
  \begin{equation}
  \label{0.13.7.7.3}
  \tag{13.7.7.3}
    \mathfrak{J}\otimes_S\mathrm{gr}^p(M)
    \longrightarrow
    \mathrm{gr}^{p+1}(M)
  \end{equation}
  \oldpage[0\textsubscript{III}]{79}
  son sobreyectivos.
  Están inducidos por
  $\mathfrak{J}\otimes_S\mathrm{F}^p(M)
  \longrightarrow\mathrm{F}^{p+1}(M)$, teniendo en cuenta el hecho
  de que la imagen del homomorfismo compuesto
  \[
    \mathfrak{J}\otimes_S\mathrm{F}^{p+1}(M)
    \longrightarrow
    \mathfrak{J}\otimes_S\mathrm{F}^p(M)
    \longrightarrow
    \mathrm{F}^{p+1}(M)
  \]
  es
  $\mathfrak{J}\mathrm{F}^{p+1}(M)
  \subset\mathrm{F}^{p+2}(M)$.
\end{enumerate}

Para la demostración, véase Bourbaki, \emph{Alg.\ comm.}, cap.~III,
\textsection3, n\textsuperscript{o}~1, prop.~3.
\end{lemma}

\begin{env}[13.7.7.4]
\label{0.13.7.7.4}
Para aplicar el lema \sref{0.13.7.7.2}, queda observar que la
topología sobre $\RR'^{\,n}\mathrm{T}(\bb{A})$ definida por la filtración
considerada lo convierte en un $S$-módulo separado y completo:
esta es la topología de límite proyectivo de los grupos discretos
$\RR^n\mathrm{T}(A_k)$.
Por otra parte, si $A_k=0$ para $k<k_0$, entonces también
$\RR^n\mathrm{T}(A_k)=0$ para $k<k_0$; por tanto
\[
  \mathrm{F}^{k_0}\bigl(\RR'^{\,n}\mathrm{T}(\bb{A})\bigr)
  =
  \RR'^{\,n}\mathrm{T}(\bb{A}).
\]
Así, todas las hipótesis del lema quedan efectivamente satisfechas.
\end{env}

\begin{corollary}[13.7.8]
\label{0.13.7.8}
Si se cumple la hipótesis $(\mathrm{F}_n)$, entonces para todo $f\in S$
existe un isomorfismo canónico
\begin{equation}
\label{0.13.7.8.1}
\tag{13.7.8.1}
  \varprojlim_k\bigl(\RR^n\mathrm{T}(A_k)\bigr)_f
  \isoto
  \RR^n\mathrm{T}(\bb{A})\otimes_S S_{\{f\}}.
\end{equation}
\end{corollary}

\begin{proof}
En efecto, $\RR^n\mathrm{T}(\bb{A})$ es un $S$-módulo de tipo finito,
y $S_{\{f\}}$ es una $S$-álgebra ádica noetheriana
(\sref[0\textsubscript{I}]{0.7.6.11}), a saber, la completación separada
de $S_f$ para la topología $\mathfrak{J}$-preádica
(\sref[0\textsubscript{I}]{0.7.6.2}).
Se sigue de \sref[0\textsubscript{I}]{0.7.7.8} y
\sref[0\textsubscript{I}]{0.7.7.1} que
\[
  \RR^n\mathrm{T}(\bb{A})\otimes_S S_{\{f\}}
\]
es isomorfo a la completación separada de
$\RR^n\mathrm{T}(\bb{A})\otimes_S S_f$ para la
topología $\mathfrak{J}$-preádica.
Un sistema fundamental de entornos de cero para esta topología
es
\[
  \bigl(
    \mathfrak{J}^p\RR^n\mathrm{T}(\bb{A})
  \bigr)\otimes_S S_f;
\]
los grupos
\[
  \mathrm{F}^p\bigl(\RR^n\mathrm{T}(\bb{A})\bigr)
  \otimes_S S_f
\]
forman, por tanto, también un sistema tal.
Este último grupo es el núcleo de la aplicación canónica
\[
  \bigl(\RR^n\mathrm{T}(\bb{A})\bigr)_f
  \longrightarrow
  \bigl(\RR^n\mathrm{T}(A_{p-1})\bigr)_f.
\]
Por consiguiente, el grupo separado asociado a
$\RR^n\mathrm{T}(\bb{A})\otimes_S S_f$ se identifica con un subgrupo
$G$ de
\[
  \varprojlim_k\bigl(\RR^n\mathrm{T}(A_k)\bigr)_f.
\]
Pero el sistema proyectivo
$((\RR^n\mathrm{T}(A_k))_f)_k$ satisface manifiestamente {\rm(ML)}, y la
imagen de $(\RR^n\mathrm{T}(\bb{A}))_f$ en cada
$(\RR^n\mathrm{T}(A_k))_f$ es igual a la imagen común de los
$(\RR^n\mathrm{T}(A_h))_f$ para todo $h\geq k$ suficientemente grande.
Se sigue inmediatamente que $G$ es \emph{denso en todas partes} en
$\varprojlim_k(\RR^n\mathrm{T}(A_k))_f$.
Como este último grupo es separado y completo, queda demostrado el
corolario.
\end{proof}


\clearpage
\phantomsection
\label{section:ega3}
\pdfbookmark[0]{Capítulo III}{chapter.III}
\addcontentsline{toc}{part}{Capítulo III: Estudio cohomológico de los haces coherentes}
\section*{CAPÍTULO III}
\setcounter{section}{0}
\setcounter{subsection}{0}
\setcounter{equation}{0}
\begin{center}
\Large\bfseries ESTUDIO COHOMOLÓGICO DE LOS HACES COHERENTES
\end{center}
\oldpage[III]{81}

\section*{Sumario}

\begin{longtable}{@{}p{0.08\textwidth}p{0.86\textwidth}@{}}
  \hyperref[section:III.1]{\textsection1}. & Cohomología de los esquemas afines.\\
  \hyperref[section:III.2]{\textsection2}. & Estudio cohomológico de los morfismos proyectivos.\\
  \hyperref[section:III.3]{\textsection3}. & Teorema de finitud para los morfismos propios.\\
  \hyperref[section:III.4]{\textsection4}. & Teorema fundamental de los morfismos propios. Aplicaciones.\\
  \hyperref[section:III.5]{\textsection5}. & Un teorema de existencia para los haces algebraicos coherentes.\\
  \textsection8. & Teorema de dualidad para los fibrados proyectivos.\\
  \textsection9. & Cohomología relativa y cohomología local; dualidad local.\\
  \textsection10. & Relaciones entre la cohomología proyectiva y la cohomología local. Técnica de compleción formal a lo largo de un divisor.\\
  \textsection11. & Grupos de Picard globales y locales.\footnotemark
\end{longtable}
\footnotetext{EGA IV no depende de los \textsection\textsection8--11 y probablemente se publicará antes que estos capítulos. \emph{[Trad.] Estos cuatro últimos capítulos nunca se publicaron.}}
\bigskip

Este capítulo presenta los teoremas fundamentales relativos a la cohomología de los haces algebraicos coherentes, con excepción de los teoremas que explican la teoría de residuos (teoremas de dualidad), que serán objeto de un capítulo posterior.
Estos resultados resultarán ser herramientas esenciales en todo lo que sigue, incluso para cuestiones cuya naturaleza no es propiamente cohomológica;
el lector verá los primeros ejemplos a partir del \hyperref[section:III.4]{\textsection4}.
Finalmente, en los \textsection\textsection8--11 desarrollaremos ciertos resultados, vinculados con la dualidad de los haces coherentes, que son particularmente importantes para las aplicaciones y pueden exponerse incluso antes de introducir la teoría general completa de residuos.

\oldpage[III]{82}
El contenido de los \hyperref[section:III.1]{\textsection1} y \hyperref[section:III.2]{\textsection2} se debe a J.-P.~Serre, y el lector observará que no hemos tenido más que seguir (FAC).
Los \textsection8 y 9 también se inspiran en (FAC) (aunque las modificaciones exigidas por los distintos contextos sean menos evidentes).
Finalmente, como dijimos en la Introducción, el \hyperref[section:III.4]{\textsection4} debe considerarse como la formalización, en lenguaje moderno, del ``teorema de invariancia'' fundamental de la ``teoría de funciones holomorfas'' de Zariski.

Llamamos la atención sobre el hecho de que los resultados del n.\textsuperscript{o}\hyperref[subsection:III.3.4]{3.4} (y las proposiciones preliminares de \sref[0]{0.13.4} a \sref[0]{0.13.7}) no se utilizarán en lo que sigue de este \hyperref[section:ega3]{Capítulo~III}, por lo que pueden omitirse en una primera lectura.
\bigskip

\section{Cohomología de los esquemas afines}
\label{section:III.1}

\subsection{Repaso del complejo del álgebra exterior}
\label{subsection:III.1.1}

\begin{env}[1.1.1]
\label{III.1.1.1}
Sea $A$ un anillo y sea $\mathbf{f}=(f_i)_{1\leq i\leq r}$ un sistema de $r$ elementos de $A$.
El \emph{complejo del álgebra exterior $K_\bullet(\mathbf{f})$} correspondiente a $\mathbf{f}$ es un complejo de cadenas (G, I, 2.2) definido del modo siguiente: el $A$-módulo graduado $K_\bullet(\mathbf{f})$ es igual al \emph{álgebra exterior $\wedge(A^r)$}, graduada de la manera usual, y la aplicación borde es la \emph{multiplicación interior $i_\mathbf{f}$} por $\mathbf{f}$, considerado como elemento del dual $\dual{(A^r)}$; recordemos que $i_\mathbf{f}$ es una \emph{antiderivación} de grado $-1$ de $\wedge(A^r)$, y que, si $(\mathbf{e}_i)_{1\leq i\leq r}$ es la base canónica de $A^r$, se tiene $i_\mathbf{f}(\mathbf{e}_i)=f_i$; la verificación de la condición $i_\mathbf{f}\circ i_\mathbf{f}=0$ es inmediata.

Una definición equivalente es la siguiente: para cada $i$, se considera un complejo de cadenas $K_\bullet(f_i)$ definido por $K_0(f_i)=K_1(f_i)=A$, $K_n(f_i)=0$ para $n\neq 0,1$; la aplicación borde queda definida por la condición de que $d_1:A\to A$ sea la \emph{multiplicación por $f_i$}.
Se toma entonces como $K_\bullet(\mathbf{f})$ el \emph{producto tensorial $K_\bullet(f_1)\otimes K_\bullet(f_2)\otimes\cdots\otimes K_\bullet(f_r)$} (G, I, 2.7), provisto de su grado total; la verificación del isomorfismo de este complejo con el definido anteriormente es inmediata.
\end{env}

\begin{env}[1.1.2]
\label{III.1.1.2}
Para todo $A$-módulo $M$, se define el \emph{complejo de cadenas}
\[
\label{III.1.1.2.1}
  K_\bullet(\mathbf{f},M)=K_\bullet(\mathbf{f})\otimes_A M
  \tag{1.1.2.1}
\]
y el \emph{complejo de cocadenas} (G, I, 2.2)
\[
\label{III.1.1.2.2}
  K^\bullet(\mathbf{f},M)=\Hom_A(K_\bullet(\mathbf{f}),M).
  \tag{1.1.2.2}
\]

Si $g$ es una $k$-cocadena de este último complejo y se pone
\[
  g(i_1,\dots,i_k)=g(\mathbf{e}_{i_1}\wedge\cdots\wedge\mathbf{e}_{i_k}),
\]
entonces $g$ se identifica con una aplicación \emph{alternada} de $[1,r]^k$ en $M$, y de las definiciones anteriores se deduce que
\[
\label{III.1.1.2.3}
  d^k g(i_1,i_2,\dots,i_{k+1})=\sum_{h=1}^{k+1}(-1)^{h-1}f_{i_h}g(i_1,\dots,\widehat{i_h},\dots,i_{k+1}).
  \tag{1.1.2.3}
\]
\end{env}

\begin{env}[1.1.3]
\label{III.1.1.3}
\oldpage[III]{83}
De los complejos anteriores se deducen, como de costumbre, los \emph{$A$-módulos de homología y de cohomología} (G, I, 2.2)
\[
  \HH_\bullet(\mathbf{f},M)=\HH_\bullet(K_\bullet(\mathbf{f},M)),
  \tag{1.1.3.1}
\]
\[
  \HH^\bullet(\mathbf{f},M)=\HH^\bullet(K^\bullet(\mathbf{f},M)).
  \tag{1.1.3.2}
\]

Se define un \emph{$A$-isomorfismo $K_i(\mathbf{f},M)\isoto K^{r-i}(\mathbf{f},M)$} que envía cada cadena $z=\sum(\mathbf{e}_{i_1}\wedge\cdots\wedge\mathbf{e}_{i_k})\otimes z_{i_1,\dots,i_k}$ a la cocadena $g_z$ tal que $g_z(j_1,\dots,j_{r-k})=\varepsilon z_{i_1,\dots,i_k}$, donde $(j_h)_{1\leq h\leq r-k}$ es la sucesión estrictamente creciente complementaria a la sucesión estrictamente creciente $(i_h)_{1\leq h\leq k}$ en $[1,r]$ y $\varepsilon=(-1)^\nu$, siendo $\nu$ el número de inversiones de la permutación $i_1,\dots,i_k,j_1,\dots,j_{r-k}$ de $[1,r]$.
Se comprueba que $g_{dz}=d(g_z)$, lo que da un isomorfismo
\[
  \HH^i(\mathbf{f},M)\isoto\HH_{r-i}(\mathbf{f},M)\text{ para }0\leq i\leq r.
  \tag{1.1.3.3}\label{III.1.1.3.3}
\]

En este capítulo se considerarán especialmente los módulos de cohomología $\HH^\bullet(\mathbf{f},M)$.

Para un $\mathbf{f}$ dado, es inmediato (G, I, 2.1) que $M\mapsto\HH^\bullet(\mathbf{f},M)$ es un \emph{funtor cohomológico} (T, II, 2.1) de la categoría de $A$-módulos en la categoría de $A$-módulos graduados, nulo en los grados $<0$ y $>r$.
Además, se tiene
\[
  \HH^0(\mathbf{f},M)=\Hom_A(A/(\mathbf{f}),M),
  \tag{1.1.3.4}
\]
denotando por $(\mathbf{f})$ el ideal de $A$ generado por $f_1,\dots,f_r$; esto se deduce inmediatamente de (\hyperref[III.1.1.2.3]{1.1.2.3}), y es claro que $\HH^0(\mathbf{f},M)$ se identifica con el submódulo de $M$ \emph{anulado por $(\mathbf{f})$}.
De manera análoga, por (\hyperref[III.1.1.2.3]{1.1.2.3}) se tiene
\[
  \HH^r(\mathbf{f},M)=M/\bigg(\sum_{i=1}^r f_i M\bigg)=(A/(\mathbf{f}))\otimes_A M.
  \tag{1.1.3.5}\label{III.1.1.3.5}
\]

Se utilizará el siguiente resultado conocido, del que se recordará una demostración para ser completos:
\end{env}

\begin{proposition}[1.1.4]
\label{III.1.1.4}
Sean $A$ un anillo, $\mathbf{f}=(f_i)_{1\leq i\leq r}$ una familia finita de elementos de $A$ y $M$ un $A$-módulo.
Si, para $1\leq i\leq r$, la aplicación $z\mapsto f_i\cdot z$ sobre $M_{i-1}=M/(f_1 M+\cdots+f_{i-1}M)$ es inyectiva, entonces $\HH^i(\mathbf{f},M)=0$ para $i\neq r$.
\end{proposition}

\begin{proof}
Basta probar que $\HH_i(\mathbf{f},M)=0$ para todo $i>0$, en virtud de (\hyperref[III.1.1.3.3]{1.1.3.3}).
Se procede por inducción sobre $r$, siendo trivial el caso $r=0$.
Póngase $\mathbf{f}'=(f_i)_{1\leq i\leq r-1}$; esta familia satisface las condiciones del enunciado, de modo que, si se pone $L_\bullet=K_\bullet(\mathbf{f}',M)$, por hipótesis se tiene $\HH_i(L_\bullet)=0$ para $i>0$, y $\HH_0(L_\bullet)=M_{r-1}$ en virtud de (\hyperref[III.1.1.3.3]{1.1.3.3}) y (\hyperref[III.1.1.3.5]{1.1.3.5}).
Para abreviar, póngase $K_\bullet=K_\bullet(f_r)=K_0\oplus K_1$, con $K_0=K_1=A$ y $d_1:K_1\to K_0$ igual a la multiplicación por $f_r$; por definición \sref{III.1.1.1}, se tiene $K_\bullet(\mathbf{f},M)=K_\bullet\otimes_A L_\bullet$.
Se dispone del siguiente lema:

\begin{lemma}[1.1.4.1]
\label{III.1.1.4.1}
Sea $K_\bullet$ un complejo de cadenas de $A$-módulos libres, nulo salvo en las dimensiones $0$ y $1$.
Para todo complejo de cadenas $L_\bullet$ de $A$-módulos se tiene una sucesión exacta
\[
  0\to\HH_0(K_\bullet\otimes\HH_p(L_\bullet))\to\HH_p(K_\bullet\otimes L_\bullet)\to\HH_1(K_\bullet\otimes\HH_{p-1}(L_\bullet))\to 0
\]
para todo índice $p$.
\end{lemma}

\oldpage[III]{84}
Este es un caso particular de una sucesión exacta de términos de grado bajo de la sucesión espectral de K\"unneth (M, XVII, 5.2 (a) y G, I, 5.5.2); puede demostrarse directamente como sigue.
Considérense $K_0$ y $K_1$ como complejos de cadenas (nulos en las dimensiones $\neq 0$ y $\neq 1$, respectivamente); se tiene entonces una sucesión exacta de complejos
\[
  0\to K_0\otimes L_\bullet\to K_\bullet\otimes L_\bullet\to K_1\otimes L_\bullet\to 0,
\]
a la que puede aplicarse la sucesión exacta de homología
\[
  \begin{aligned}
  \cdots&\to\HH_{p+1}(K_1\otimes L_\bullet)
  \xrightarrow{\partial}\HH_p(K_0\otimes L_\bullet)
  \to\HH_p(K_\bullet\otimes L_\bullet)\\
  &\to\HH_p(K_1\otimes L_\bullet)
  \xrightarrow{\partial}\HH_{p-1}(K_0\otimes L_\bullet)\to\cdots.
  \end{aligned}
\]
Pero es evidente que $\HH_p(K_0\otimes L_\bullet)=K_0\otimes\HH_p(L_\bullet)$ y $\HH_p(K_1\otimes L_\bullet)=K_1\otimes\HH_{p-1}(L_\bullet)$ para todo $p$; además, se comprueba inmediatamente que el operador $\partial:K_1\otimes\HH_p(L_\bullet)\to K_0\otimes\HH_p(L_\bullet)$ no es otro que $d_1\otimes 1$; el lema se deduce así de la sucesión exacta anterior y de la definición de $\HH_0(K_\bullet\otimes\HH_p(L_\bullet))$ y $\HH_1(K_\bullet\otimes\HH_{p-1}(L_\bullet))$.

Establecido el lema, el final de la demostración de la Proposición \sref{III.1.1.4} es inmediato: la hipótesis de inducción y el Lema \sref{III.1.1.4.1} dan $\HH_p(K_\bullet\otimes L_\bullet)=0$ para $p\geq 2$; además, si se prueba que $\HH_1(K_\bullet,\HH_0(L_\bullet))=0$, se deduce también del Lema \sref{III.1.1.4.1} que $\HH_1(K_\bullet\otimes L_\bullet)=0$; pero, por definición, $\HH_1(K_\bullet,\HH_0(L_\bullet))$ no es otro que el núcleo de la aplicación $z\mapsto f_r\cdot z$ sobre $M_{r-1}$, y, como por hipótesis este núcleo es nulo, queda terminada la demostración.
\end{proof}

\begin{env}[1.1.5]
\label{III.1.1.5}
Sea $\mathbf{g}=(g_i)_{1\leq i\leq r}$ una segunda sucesión de $r$ elementos de $A$, y póngase $\mathbf{f}\mathbf{g}=(f_i g_i)_{1\leq i\leq r}$.
Puede definirse un homomorfismo canónico de complejos
\[
  \vphi_\mathbf{g}:K_\bullet(\mathbf{f}\mathbf{g})\to K_\bullet(\mathbf{f})
  \tag{1.1.5.1}
\]
como la extensión canónica al álgebra exterior $\wedge(A^r)$ de la aplicación $A$-lineal $(x_1,\dots,x_r)\mapsto(g_1 x_1,\dots,g_r x_r)$ de $A^r$ en sí mismo.
Para ver que se trata de un homomorfismo de complejos, basta observar, en general, que, si $u:E\to F$ es una aplicación $A$-lineal, $\mathbf{x}\in\dual{F}$ y $\mathbf{y}={}^t u(\mathbf{x})\in\dual{E}$, se tiene la fórmula
\[
  (\wedge u)\circ i_\mathbf{y}=i_\mathbf{x}\circ(\wedge u);
  \tag{1.1.5.2}
\]
en efecto, ambos miembros son antiderivaciones de $\wedge F$, y basta comprobar que coinciden sobre $F$, lo cual se deduce inmediatamente de las definiciones.

Cuando se identifica $K_\bullet(\mathbf{f})$ con el producto tensorial de los $K_\bullet(f_i)$ \sref{III.1.1.1}, $\vphi_\mathbf{g}$ es el producto tensorial de los $\vphi_{g_i}$, donde $\vphi_{g_i}$ es la identidad en grado $0$ y la multiplicación por $g_i$ en grado $1$.
\end{env}

\begin{env}[1.1.6]
\label{III.1.1.6}
En particular, para todo par de enteros $m$ y $n$ tales que $0\leq n\leq m$, se tienen homomorfismos de complejos
\[
  \vphi_{\mathbf{f}^{m-n}}:K_\bullet(\mathbf{f}^m)\to K_\bullet(\mathbf{f}^n)
  \tag{1.1.6.1}
\]
y, por consiguiente, homomorfismos
\[
  \vphi_{\mathbf{f}^{m-n}}:K^\bullet(\mathbf{f}^n,M)\to K^\bullet(\mathbf{f}^m,M),
  \tag{1.1.6.2}
\]
\[
  \vphi_{\mathbf{f}^{m-n}}:\HH^\bullet(\mathbf{f}^n,M)\to\HH^\bullet(\mathbf{f}^m,M).
  \tag{1.1.6.3}
\]

\oldpage[III]{85}
Estos últimos homomorfismos satisfacen evidentemente la condición de transitividad $\vphi_{\mathbf{f}^{m-p}}=\vphi_{\mathbf{f}^{m-n}}\circ\vphi_{\mathbf{f}^{n-p}}$ para $p\leq n\leq m$; definen, por tanto, dos \emph{sistemas inductivos} de $A$-módulos; se pone
\[
  C^\bullet((\mathbf{f}),M)=\varinjlim_n K^\bullet(\mathbf{f}^n,M),
  \tag{1.1.6.4}
\]
\[
  \HH^\bullet((\mathbf{f}),M)=\HH^\bullet(C^\bullet((\mathbf{f}),M))=\varinjlim_n\HH^\bullet(\mathbf{f}^n,M),
  \tag{1.1.6.5}
\]
deduciéndose la última igualdad del hecho de que el paso al límite inductivo conmuta con el funtor $\HH^\bullet$ (G, I, 2.1).
Se verá más adelante \sref{III.1.4.3} que $\HH^\bullet((\mathbf{f}),M)$ depende en realidad solo del \emph{ideal $(\mathbf{f})$} de $A$ (e incluso solo de la topología $(\mathbf{f})$-preádica sobre $A$), lo que justifica las notaciones.

Es claro que $M\mapsto C^\bullet((\mathbf{f}),M)$ es un funtor exacto $A$-lineal y que $M\mapsto\HH^\bullet((\mathbf{f}),M)$ es un funtor cohomológico.
\end{env}

\begin{env}[1.1.7]
\label{III.1.1.7}
Pónganse $\mathbf{f}=(f_i)\in A^r$ y $\mathbf{g}=(g_i)\in A^r$; denótese por $e_\mathbf{g}$ la multiplicación por la izquierda por el vector $\mathbf{g}\in A^r$ sobre el álgebra exterior $\wedge(A^r)$; se sabe que se tiene la \emph{fórmula de homotopía}
\[
  i_\mathbf{f}e_\mathbf{g}+e_\mathbf{g}i_\mathbf{f}=\langle\mathbf{g},\mathbf{f}\rangle 1
  \tag{1.1.7.1}
\]
en el $A$-módulo $\wedge(A^r)$ ($1$ denota el automorfismo identidad de $\wedge(A^r)$); esta relación implica también que \emph{en el complejo $K_\bullet(\mathbf{f})$} se tiene
\[
  de_\mathbf{g}+e_\mathbf{g}d=\langle\mathbf{g},\mathbf{f}\rangle 1.
  \tag{1.1.7.2}
\]

Si el ideal $(\mathbf{f})$ es igual a $A$, existe $\mathbf{g}\in A^r$ tal que $\langle\mathbf{g},\mathbf{f}\rangle=\sum_{i=1}^r g_i f_i=1$.
Por consiguiente (G, I, 2.4):
\end{env}

\begin{proposition}[1.1.8]
\label{III.1.1.8}
Supóngase que el ideal $(\mathbf{f})$ generado por los $f_i$ es igual a $A$.
Entonces el complejo $K_\bullet(\mathbf{f})$ es homotópicamente trivial, al igual que los complejos $K_\bullet(\mathbf{f},M)$ y $K^\bullet(\mathbf{f},M)$ para todo $A$-módulo $M$.
\end{proposition}

\begin{corollary}[1.1.9]
\label{III.1.1.9}
Si $(\mathbf{f})=A$, entonces $\HH^\bullet(\mathbf{f},M)=0$ y $\HH^\bullet((\mathbf{f}),M)=0$ para todo $A$-módulo $M$.
\end{corollary}

\begin{proof}
En efecto, se tiene entonces $(\mathbf{f}^n)=A$ para todo $n$.
\end{proof}

\begin{remark}[1.1.10]
\label{III.1.1.10}
Con las mismas notaciones anteriores, póngase $X=\Spec(A)$ y sea $Y$ el subpreesquema cerrado de $X$ definido por el ideal $(\mathbf{f})$.
Se probará en \textsection9 que $\HH^\bullet((\mathbf{f}),M)$ es isomorfo a la cohomología $\HH_Y^\bullet(X,\widetilde{M})$ correspondiente al antifiltro $\Phi$ de subconjuntos cerrados de $Y$ (T, 3.2).
Se mostrará también que la Proposición \sref{III.1.2.3}, aplicada a $X$ y a $\sh{F}=\widetilde{M}$, es un caso particular de una sucesión exacta de cohomología
\[
  \cdots\to\HH_Y^p(X,\sh{F})\to\HH^p(X,\sh{F})\to\HH^p(X\setmin Y,\sh{F})\to\HH_Y^{p+1}(X,\sh{F})\to\cdots.
\]
\end{remark}

\subsection{Cohomología de \v Cech de un recubrimiento abierto}
\label{subsection:III.1.2}

\begin{notation}[1.2.1]
\label{III.1.2.1}
En esta sección se denota por:
\begin{enumerate}
  \item $X$ un preesquema;
  \item $\sh{F}$ un $\sh{O}_X$-módulo cuasi-coherente;
\oldpage[III]{86}
  \item $A=\Gamma(X,\sh{O}_X)$, $M=\Gamma(X,\sh{F})$;
  \item $\mathbf{f}=(f_i)_{1\leq i\leq r}$ un sistema finito de elementos de $A$;
  \item $U_i=X_{f_i}$, el abierto \sref[0]{0.5.5.2} de los $x\in X$ tales que $f_i(x)\neq 0$;
  \item $U=\bigcup_{i=1}^r U_i$;
  \item $\mathfrak{U}$ el recubrimiento $(U_i)_{1\leq i\leq r}$ de $U$.
\end{enumerate}
\end{notation}

\begin{env}[1.2.2]
\label{III.1.2.2}
Supóngase que $X$ es un preesquema cuyo espacio subyacente es \emph{noetheriano}, o bien un \emph{esquema} cuyo espacio subyacente es \emph{cuasicompacto}.
Se sabe entonces \sref[I]{I.9.3.3} que $\Gamma(U_i,\sh{F})=M_{f_i}$.
Póngase
\[
  U_{i_0 i_1\cdots i_p}=\bigcap_{k=0}^p U_{i_k}=X_{f_{i_0}f_{i_1}\cdots f_{i_p}}
\]
\sref[0]{0.5.5.3}; se tiene, por tanto, también
\[
  \Gamma(U_{i_0 i_1\cdots i_p},\sh{F})=M_{f_{i_0}f_{i_1}\cdots f_{i_p}}.
  \tag{1.2.2.1}
\]

Por \sref[0]{0.1.6.1}, $M_{f_{i_0}f_{i_1}\cdots f_{i_p}}$ se identifica con el límite inductivo $\varinjlim_n M_{i_0 i_1\cdots i_p}^{(n)}$, donde el sistema inductivo está formado por los $M_{i_0 i_1\cdots i_p}^{(n)}=M$, siendo los homomorfismos $\vphi_{nm}:M_{i_0 i_1\cdots i_p}^{(m)}\to M_{i_0 i_1\cdots i_p}^{(n)}$ la multiplicación por $(f_{i_0}f_{i_1}\cdots f_{i_p})^{n-m}$ para $m\leq n$.
Se denota por $C_n^p(M)$ el conjunto de aplicaciones \emph{alternadas} de $[1,r]^{p+1}$ en $M$ (para todo $n$); estos $A$-módulos forman también un sistema inductivo respecto de los $\vphi_{nm}$.
Si $C^p(\mathfrak{U},\sh{F})$ es el grupo de las $p$-cocadenas \emph{alternadas} de \v Cech relativas al recubrimiento $\mathfrak{U}$, con coeficientes en $\sh{F}$ (G, II, 5.1), de lo anterior se deduce que puede escribirse
\[
  C^p(\mathfrak{U},\sh{F})=\varinjlim_n C_n^p(M).
  \tag{1.2.2.2}
\]

Con las notaciones de \sref{III.1.1.2}, $C_n^p(M)$ se identifica con $K^{p+1}(\mathbf{f}^n,M)$, y la aplicación $\vphi_{nm}$ se identifica con la aplicación $\vphi_{\mathbf{f}^{n-m}}$ definida en \sref{III.1.1.6}.
Se tiene así, para todo $p\geq 0$, un isomorfismo canónico funtorial
\[
  C^p(\mathfrak{U},\sh{F})\isoto C^{p+1}((\mathbf{f}),M).
  \tag{1.2.2.3}\label{III.1.2.2.3}
\]

Además, la fórmula (\hyperref[III.1.1.2.3]{1.1.2.3}) y la definición de la cohomología de un recubrimiento (G, II, 5.1) muestran que los isomorfismos (\hyperref[III.1.2.2.3]{1.2.2.3}) son compatibles con las aplicaciones coborde.
\end{env}

\begin{proposition}[1.2.3]
\label{III.1.2.3}
Si $X$ es un preesquema cuyo espacio subyacente es noetheriano, o un esquema cuyo espacio subyacente es cuasicompacto, existe un isomorfismo canónico funtorial en $\sh{F}$
\[
  \HH^p(\mathfrak{U},\sh{F})\isoto\HH^{p+1}((\mathbf{f}),M)\text{ para }p\geq 1.
  \tag{1.2.3.1}\label{III.1.2.3.1}
\]

Además, se tiene una sucesión exacta funtorial en $\sh{F}$
\[
  0\to\HH^0((\mathbf{f}),M)\to M\to\HH^0(\mathfrak{U},\sh{F})\to\HH^1((\mathbf{f}),M)\to 0.
  \tag{1.2.3.2}\label{III.1.2.3.2}
\]
\end{proposition}

\oldpage[III]{87}
\begin{proof}
Los isomorfismos (\hyperref[III.1.2.3.1]{1.2.3.1}) son consecuencias inmediatas de lo visto en \sref{III.1.2.2}.
Por otra parte, se tiene $C^0(\mathfrak{U},\sh{F})=C^1((\mathbf{f}),M)$; por consiguiente, $\HH^0(\mathfrak{U},\sh{F})$ se identifica con el subgrupo de los $1$-cociclos de $C^1((\mathbf{f}),M)$; como $M=C^0((\mathbf{f}),M)$, la sucesión exacta (\hyperref[III.1.2.3.2]{1.2.3.2}) no es otra que la dada por la definición de los grupos de cohomología $\HH^0((\mathbf{f}),M)$ y $\HH^1((\mathbf{f}),M)$.
\end{proof}

\begin{corollary}[1.2.4]
\label{III.1.2.4}
Supóngase que los $X_{f_i}$ son cuasicompactos y que existen $g_i\in\Gamma(U,\sh{O}_X)$ tales que $\sum_i g_i(f_i|U)=1|U$.
Entonces, para todo $(\sh{O}_X|U)$-módulo cuasi-coherente $\sh{G}$, se tiene $\HH^p(\mathfrak{U},\sh{G})=0$ para $p>0$; si además $U=X$, el homomorfismo canónico (\hyperref[III.1.2.3.2]{1.2.3.2}) $\Gamma(X,\sh{G})\to\HH^0(\mathfrak{U},\sh{G})$ es biyectivo.
\end{corollary}

\begin{proof}
Como, por hipótesis, los $U_i=X_{f_i}$ son cuasicompactos, también lo es $U$, y puede reducirse al caso $U=X$; poniendo $N=\Gamma(X,\sh{G})$, la hipótesis implica entonces que $\HH^p((\mathbf{f}),N)=0$ para todo $p\geq 0$ \sref{III.1.1.9}.
El corolario se deduce inmediatamente de (\hyperref[III.1.2.3.1]{1.2.3.1}) y (\hyperref[III.1.2.3.2]{1.2.3.2}).
\end{proof}

Obsérvese que, puesto que $\HH^0(\mathfrak{U},\sh{F})=\HH^0(U,\sh{F})$ (G, II, 5.2.2), se ha demostrado de nuevo \sref[I]{I.1.3.7} como caso particular.

\begin{remark}[1.2.5]
\label{III.1.2.5}
Supóngase que $X$ es un \emph{esquema afín}; entonces los $U_i=X_{f_i}=D(f_i)$ son abiertos afines, al igual que los $U_{i_0 i_1\cdots i_p}$ (pero $U$ no es necesariamente afín).
En este caso, los funtores $\Gamma(X,\sh{F})$ y $\Gamma(U_{i_0 i_1\cdots i_p},\sh{F})$ son exactos en $\sh{F}$ \sref[I]{I.1.3.11}.
Si se tiene una sucesión exacta $0\to\sh{F}'\to\sh{F}\to\sh{F}''\to 0$ de $\sh{O}_X$-módulos cuasi-coherentes, la sucesión de complejos
\[
  0\to C^\bullet(\mathfrak{U},\sh{F}')\to C^\bullet(\mathfrak{U},\sh{F})\to C^\bullet(\mathfrak{U},\sh{F}'')\to 0
\]
es exacta y da, por tanto, una sucesión exacta de cohomología
\[
  \cdots\to\HH^p(\mathfrak{U},\sh{F}')\to\HH^p(\mathfrak{U},\sh{F})\to\HH^p(\mathfrak{U},\sh{F}'')\xrightarrow{\partial}\HH^{p+1}(\mathfrak{U},\sh{F}')\to\cdots.
\]
Por otra parte, si se pone $M'=\Gamma(X,\sh{F}')$ y $M''=\Gamma(X,\sh{F}'')$, la sucesión $0\to M'\to M\to M''\to 0$ es exacta; como $C^\bullet((\mathbf{f}),M)$ es un funtor exacto en $M$, se tiene asimismo la sucesión exacta de cohomología
\[
  \cdots\to\HH^p((\mathbf{f}),M')\to\HH^p((\mathbf{f}),M)\to\HH^p((\mathbf{f}),M'')\xrightarrow{\partial}\HH^{p+1}((\mathbf{f}),M')\to\cdots.
\]

Así pues, como el diagrama
\[
  \xymatrix{
    0\ar[r] &
    C^\bullet(\mathfrak{U},\sh{F}')\ar[r]\ar[d] &
    C^\bullet(\mathfrak{U},\sh{F})\ar[r]\ar[d] &
    C^\bullet(\mathfrak{U},\sh{F}'')\ar[r]\ar[d] &
    0\\
    0\ar[r] &
    C^\bullet((\mathbf{f}),M')\ar[r] &
    C^\bullet((\mathbf{f}),M)\ar[r] &
    C^\bullet((\mathbf{f}),M'')\ar[r] &
    0
  }
\]
\oldpage[III]{88}
es conmutativo, se concluye que los diagramas
\[
  \xymatrix{
    \HH^p(\mathfrak{U},\sh{F}'')\ar[r]^\partial\ar[d] &
    \HH^{p+1}(\mathfrak{U},\sh{F}')\ar[d]\\
    \HH^{p+1}((\mathbf{f}),M'')\ar[r]^\partial &
    \HH^{p+2}((\mathbf{f}),M')
  }
  \tag{1.2.5.1}
\]
son conmutativos para todo $p$~(G, I, 2.1.1).
\end{remark}

\subsection{Cohomología de un esquema afín}
\label{subsection:III.1.3}

\begin{theorem}[1.3.1]
\label{III.1.3.1}
Sea $X$ un esquema afín.
Para todo $\sh{O}_X$-módulo cuasi-coherente $\sh{F}$, se tiene $\HH^p(X,\sh{F})=0$ para todo $p>0$.
\end{theorem}

\begin{proof}
Sea $\mathfrak{U}$ un recubrimiento finito de $X$ por los abiertos afines $X_{f_i}=D(f_i)$ ($1\leq i\leq r$); se sabe entonces que el ideal de $A=\Gamma(X,\sh{O}_X)$ generado por los $f_i$ es igual a $A$.
Del Corolario~\sref{III.1.2.4} se concluye, pues, que $\HH^p(\mathfrak{U},\sh{F})=0$ para $p>0$.
Como existen recubrimientos finitos de $X$ por abiertos afines arbitrariamente finos \sref[I]{I.1.1.10}, la definición de la cohomología de \v Cech (G, II, 5.8) muestra que también $\CHH^p(X,\sh{F})=0$ para $p>0$.
Pero esto se aplica asimismo a todo preesquema $X_f$, para $f\in A$ \sref[I]{I.1.3.6}; por tanto, $\CHH^p(X_f,\sh{F})=0$ para $p>0$.
Como $X_f\cap X_g=X_{fg}$, se deduce también que $\HH^p(X,\sh{F})=0$ para todo $p>0$, en virtud de (G, II, 5.9.2).
\end{proof}

\begin{corollary}[1.3.2]
\label{III.1.3.2}
Sea $Y$ un preesquema y sea $f:X\to Y$ un morfismo afín \sref[II]{II.1.6.1}.
Para todo $\sh{O}_X$-módulo cuasi-coherente $\sh{F}$, se tiene $\RR^q f_*(\sh{F})=0$ para $q>0$.
\end{corollary}

\begin{proof}
Por definición, $\RR^q f_*(\sh{F})$ es el $\sh{O}_Y$-módulo asociado al prehaz $U\mapsto\HH^q(f^{-1}(U),\sh{F})$, donde $U$ recorre los abiertos de $Y$.
Pero los abiertos afines forman una base de $Y$, y, para un abierto tal $U$, $f^{-1}(U)$ es afín \sref[II]{II.1.3.2}; por tanto, $\HH^q(f^{-1}(U),\sh{F})=0$ por el Teorema~\sref{III.1.3.1}, lo que prueba el corolario.
\end{proof}

\begin{corollary}[1.3.3]
\label{III.1.3.3}
Sea $Y$ un preesquema y sea $f:X\to Y$ un morfismo afín.
Para todo $\sh{O}_X$-módulo cuasi-coherente $\sh{F}$, el homomorfismo canónico $\HH^p(Y,f_*(\sh{F}))\to\HH^p(X,\sh{F})$ (\textbf{0},~12.1.3.1) es biyectivo para todo $p$.
\end{corollary}

\begin{proof}
Basta (por \sref[0]{0.12.1.7}) mostrar que los homomorfismos de borde $''E_2^{p0}=\HH^p(Y,f_*(\sh{F}))\to\HH^p(X,\sh{F})$ de la segunda sucesión espectral del funtor compuesto $\Gamma f_*$ son biyectivos.
Pero el término $E_2$ de esta sucesión espectral viene dado por $''E_2^{pq}=\HH^p(Y,\RR^q f_*(\sh{F}))$ (G, II, 4.17.1); por el Corolario~\sref{III.1.3.2}, $''E_2^{pq}=0$ para $q>0$, y la sucesión espectral degenera; de ahí la afirmación \sref[0]{0.11.1.6}.
\end{proof}

\begin{corollary}[1.3.4]
\label{III.1.3.4}
Sean $f:X\to Y$ un morfismo afín y $g:Y\to Z$ un morfismo.
\oldpage[III]{89}
Para todo $\sh{O}_X$-módulo cuasi-coherente $\sh{F}$, el homomorfismo canónico $\RR^p g_*(f_*(\sh{F}))\to\RR^p(g\circ f)_*(\sh{F})$ (\textbf{0},~12.2.5.1) es biyectivo para todo $p$.
\end{corollary}

\begin{proof}
Basta observar que, según el Corolario~\sref{III.1.3.3}, para todo abierto afín $W$ de $Z$, el homomorfismo canónico $\HH^p(g^{-1}(W),f_*(\sh{F}))\to\HH^p(f^{-1}(g^{-1}(W)),\sh{F})$ es biyectivo; esto prueba que el homomorfismo de prehaces que define el homomorfismo canónico $\RR^p g_*(f_*(\sh{F}))\to\RR^p(g\circ f)_*(\sh{F})$ es biyectivo \sref[0]{0.12.2.5}.
\end{proof}

\subsection{Aplicación a la cohomología de preesquemas arbitrarios}
\label{subsection:III.1.4}

\begin{proposition}[1.4.1]
\label{III.1.4.1}
Sea $X$ un esquema y sea $\mathfrak{U}=(U_\alpha)$ un recubrimiento de $X$ por abiertos afines.
Para todo $\sh{O}_X$-módulo cuasi-coherente $\sh{F}$, los módulos de cohomología $\HH^\bullet(X,\sh{F})$ y $\HH^\bullet(\mathfrak{U},\sh{F})$ (sobre $\Gamma(X,\sh{O}_X)$) son canónicamente isomorfos.
\end{proposition}

\begin{proof}
Como $X$ es un esquema, toda intersección finita $V$ de abiertos del recubrimiento $\mathfrak{U}$ es afín \sref[I]{I.5.5.6}; por tanto, $\HH^q(V,\sh{F})=0$ para $q\geq 1$ por el Teorema~\sref{III.1.3.1}.
La proposición se deduce entonces de un teorema de Leray (G, II, 5.4.1).
\end{proof}

\begin{remark}[1.4.2]
\label{III.1.4.2}
Obsérvese que el resultado de la Proposición~\sref{III.1.4.1} sigue siendo válido cuando las intersecciones finitas de los conjuntos $U_\alpha$ son afines, incluso si no se supone necesariamente que $X$ sea un esquema.
\end{remark}

\begin{corollary}[1.4.3]
\label{III.1.4.3}
Sea $X$ un esquema de espacio subyacente cuasicompacto, sea $A=\Gamma(X,\sh{O}_X)$ y sea $\mathbf{f}=(f_i)_{1\leq i\leq r}$ una sucesión finita de elementos de $A$ tal que los $X_{f_i}$ (notación de \sref{III.1.2.1}) sean afines.
Entonces (con las notaciones de \sref{III.1.2.1}), para todo $\sh{O}_X$-módulo cuasi-coherente $\sh{F}$, se tiene un isomorfismo canónico funtorial en $\sh{F}$
\[
  \HH^q(U,\sh{F})\isoto\HH^{q+1}((\mathbf{f}),M)\text{ para }q\geq 1,
  \tag{1.4.3.1}\label{III.1.4.3.1}
\]
y una sucesión exacta funtorial en $\sh{F}$
\[
  0\to\HH^0((\mathbf{f}),M)\to M\to\HH^0(U,\sh{F})\to\HH^1((\mathbf{f}),M)\to 0.
  \tag{1.4.3.2}
\]
\end{corollary}

\begin{proof}
Esto se deduce inmediatamente de las Proposiciones~\sref{III.1.4.1} y \sref{III.1.2.3}.
\end{proof}

\begin{env}[1.4.4]
\label{III.1.4.4}
Si $X$ es un \emph{esquema afín}, de la Observación~\sref{III.1.2.5} y de la Proposición~\sref{III.1.4.1} se deduce que, para todo $q\geq 0$, los diagramas
\[
  \xymatrix{
    \HH^q(U,\sh{F}'')\ar[r]^\partial\ar[d] &
    \HH^{q+1}(U,\sh{F}')\ar[d]\\
    \HH^{q+1}((\mathbf{f}),M'')\ar[r]^\partial &
    \HH^{q+2}((\mathbf{f}),M')
  }
  \tag{1.4.4.1}
\]
correspondientes a una sucesión exacta $0\to\sh{F}'\to\sh{F}\to\sh{F}''\to 0$ de $\sh{O}_X$-módulos cuasi-coherentes (con las notaciones de la Observación~\sref{III.1.2.5}) son conmutativos.
\end{env}

\begin{proposition}[1.4.5]
\label{III.1.4.5}
\oldpage[III]{90}
Sea $X$ un esquema cuasicompacto, sea $\sh{L}$ un $\sh{O}_X$-módulo invertible, y considérese el anillo graduado $A_\bullet=\Gamma_\bullet(\sh{L})$ \sref[0]{0.5.4.6}; entonces $\HH^\bullet(\sh{F},\sh{L})=\bigoplus_{n\in\bb{Z}}\HH^\bullet(X,\sh{F}\otimes\sh{L}^{\otimes n})$ es un $A_\bullet$-módulo graduado y, para todo $f\in A_n$, se tiene un isomorfismo canónico
\[
  \HH^\bullet(X_f,\sh{F})\isoto(\HH^\bullet(\sh{F},\sh{L}))_{(f)}
  \tag{1.4.5.1}\label{III.1.4.5.1}
\]
de $(A_\bullet)_{(f)}$-módulos.
\end{proposition}

\begin{proof}
Como $X$ es un esquema cuasicompacto, puede calcularse la cohomología de todos los $\sh{O}_X$-módulos $\sh{F}\otimes\sh{L}^{\otimes n}$ utilizando un mismo recubrimiento finito $\mathfrak{U}=(U_i)$, formado por abiertos afines tales que la restricción $\sh{L}|U_i$ sea isomorfa a $\sh{O}_X|U_i$ para cada $i$ \sref{III.1.4.1}.
Es entonces inmediato que los $U_i\cap X_f$ son abiertos afines \sref[I]{I.1.3.6}, y puede calcularse así la cohomología $\HH^\bullet(X_f,\sh{F}\otimes\sh{L}^{\otimes n})$ utilizando el recubrimiento $\mathfrak{U}|X_f=(U_i\cap X_f)$ \sref{III.1.4.1}.
Es inmediato que, para todo $f\in A_n$, la multiplicación por $f$ define un homomorfismo $C^\bullet(\mathfrak{U},\sh{F}\otimes\sh{L}^m)\to C^\bullet(\mathfrak{U},\sh{F}\otimes\sh{L}^{\otimes(m+n)})$ y, por tanto, un homomorfismo $\HH^\bullet(\mathfrak{U},\sh{F}\otimes\sh{L}^{\otimes m})\to\HH^\bullet(\mathfrak{U},\sh{F}\otimes\sh{L}^{\otimes(m+n)})$, lo que establece la primera afirmación.
Por otra parte, para un $f\in A_n$ dado, de \sref[I]{I.9.3.2} se deduce que se tiene un isomorfismo de complejos de $(A_\bullet)_{(f)}$-módulos
\[
  C^\bullet(\mathfrak{U}|X_f,\sh{F})\isoto\bigg(\!\!C^\bullet\bigg(\mathfrak{U},\bigoplus_{n\in\bb{Z}}\sh{F}\otimes\sh{L}^{\otimes n}\bigg)\!\!\bigg)_{(f)},
\]
teniendo en cuenta \sref[I]{I.1.3.9}[ii].
Al pasar a la cohomología de estos complejos se obtiene el isomorfismo (\hyperref[III.1.4.5.1]{1.4.5.1}), recordando que el funtor $M\mapsto M_{(f)}$ es exacto sobre la categoría de $A_\bullet$-módulos graduados.
\end{proof}

\begin{corollary}[1.4.6]
\label{III.1.4.6}
Supóngase que se satisfacen las hipótesis de la Proposición~\sref{III.1.4.5} y, además, que $\sh{L}=\sh{O}_X$.
Si se pone $A=\Gamma(X,\sh{O}_X)$, entonces, para todo $f\in A$, se tiene un isomorfismo canónico $\HH^\bullet(X_f,\sh{F})\isoto(\HH^\bullet(X,\sh{F}))_f$ de $A_f$-módulos.
\end{corollary}

\begin{corollary}[1.4.7]
\label{III.1.4.7}
Sea $X$ un esquema cuasicompacto y sea $f$ un elemento de $\Gamma(X,\sh{O}_X)$.
\begin{enumerate}
  \item[{\rm(i)}] Supóngase que el abierto $X_f$ es afín.
    Entonces, para todo $\sh{O}_X$-módulo cuasi-coherente $\sh{F}$, todo $i>0$ y todo $\xi\in\HH^i(X,\sh{F})$, existe un entero $n>0$ tal que $f^n\xi=0$.
  \item[{\rm(ii)}] Recíprocamente, supóngase que $X_f$ es cuasicompacto y que, para todo haz cuasi-coherente de ideales $\sh{J}$ de $\sh{O}_X$ y todo $\zeta\in\HH^1(X,\sh{J})$, existe un $n>0$ tal que $f^n\zeta=0$.
    Entonces $X_f$ es afín.
\end{enumerate}
\end{corollary}

\begin{proof}
\medskip\noindent
\begin{enumerate}
  \item[(i)] Si $X_f$ es afín, se tiene $\HH^i(X_f,\sh{F})=0$ para todo $i>0$ \sref{III.1.3.1}, de modo que la afirmación se deduce directamente del Corolario~\sref{III.1.4.6}.
  \item[(ii)] En virtud del criterio de Serre \sref[II]{II.5.2.1}, basta probar que, para todo haz cuasi-coherente de ideales $\sh{K}$ de $\sh{O}_X|X_f$, se tiene $\HH^1(X_f,\sh{K})=0$.
    Como $X_f$ es un abierto cuasicompacto de un esquema cuasicompacto $X$, existe un haz cuasi-coherente de ideales $\sh{J}$ de $\sh{O}_X$ tal que $\sh{K}=\sh{J}|X_f$ \sref[I]{I.9.4.2}.
    Según el Corolario~\sref{III.1.4.6}, se tiene $\HH^1(X_f,\sh{K})=(\HH^1(X,\sh{J}))_f$, y la hipótesis implica que el miembro de la derecha es nulo, lo que da la afirmación.
\end{enumerate}
\end{proof}

\begin{remark}[1.4.8]
\label{III.1.4.8}
Obsérvese que el Corolario~\sref{III.1.4.7}[i] da una demostración más sencilla de la relación (\textbf{II},~4.5.13.2).
\end{remark}

\begin{lemma}[1.4.9]
\label{III.1.4.9}
\oldpage[III]{91}
Sea $X$ un esquema cuasicompacto, sea $\mathfrak{U}=(U_i)_{1\leq i\leq n}$ un recubrimiento finito de $X$ por abiertos afines y sea $\sh{F}$ un $\sh{O}_X$-módulo cuasi-coherente.
El complejo de haces $\sh{C}^\bullet(\mathfrak{U},\sh{F})$ definido por el recubrimiento $\mathfrak{U}$ (G, II, 5.2) es entonces un $\sh{O}_X$-módulo cuasi-coherente.
\end{lemma}

\begin{proof}
De las definiciones (G, II, 5.2) se deduce que $\sh{C}^p(\mathfrak{U},\sh{F})$ es la suma directa de los haces imagen directa de los $\sh{F}|U_{i_0\cdots i_p}$ por la inyección canónica $U_{i_0\cdots i_p}\to X$.
La hipótesis de que $X$ es un esquema implica que estas inyecciones son morfismos afines \sref[I]{I.5.5.6}; por tanto, los $\sh{C}^p(\mathfrak{U},\sh{F})$ son cuasi-coherentes \sref[II]{II.1.2.6}.
\end{proof}

\begin{proposition}[1.4.10]
\label{III.1.4.10}
Sea $u:X\to Y$ un morfismo separado y cuasicompacto.
Para todo $\sh{O}_X$-módulo cuasi-coherente $\sh{F}$, los $\RR^q u_*(\sh{F})$ son $\sh{O}_Y$-módulos cuasi-coherentes.
\end{proposition}

\begin{proof}
La cuestión es local sobre $Y$, de modo que puede suponerse que $Y$ es afín.
Entonces $X$ es una unión finita de abiertos afines $U_i$ ($1\leq i\leq n$); sea $\mathfrak{U}$ el recubrimiento $(U_i)$.
Además, como $Y$ es un esquema, de \sref[I]{I.5.5.10} se deduce que, para todo abierto afín $V\subset Y$, la inyección canónica $u^{-1}(V)\to X$ es un morfismo afín; se concluye (Proposición~\sref{III.1.4.1} y (G, II, 5.2)) que se tiene un isomorfismo canónico
\[
  \HH^\bullet(u^{-1}(V),\sh{F})\isoto\HH^\bullet(\Gamma(V,\sh{K}^\bullet)),
  \tag{1.4.10.1}\label{III.1.4.10.1}
\]
donde se pone $\sh{K}^\bullet=u_*(\sh{C}^\bullet(\mathfrak{U},\sh{F}))$.
Según el Lema~\sref{III.1.4.9} y \sref[I]{I.9.2.2}, $\sh{K}^\bullet$ es un $\sh{O}_Y$-módulo cuasi-coherente; además, constituye un \emph{complejo de haces}, puesto que también lo es $\sh{C}^\bullet(\mathfrak{U},\sh{F})$.
De la definición de la cohomología $\sh{H}^\bullet(\sh{K}^\bullet)$ (G, II, 4.1) se deduce entonces que esta última está formada por $\sh{O}_Y$-módulos cuasi-coherentes \sref[I]{I.4.1.1}.
Como (para $V$ afín en $Y$) el funtor $\Gamma(V,\sh{G})$ es exacto en $\sh{G}$ sobre la categoría de $\sh{O}_Y$-módulos cuasi-coherentes, se tiene (G, II, 4.1)
\[
  \HH^\bullet(\Gamma(V,\sh{K}^\bullet))=\Gamma(V,\sh{H}^\bullet(\sh{K}^\bullet)).
  \tag{1.4.10.2}\label{III.1.4.10.2}
\]

Finalmente, obsérvese que de la definición del homomorfismo canónico
\[
  \HH^\bullet(\mathfrak{U},\sh{F})\to\HH^\bullet(X,\sh{F}),
\]
dado en (G, II, 5.2), se deduce que, si $V'\subset V$ es otro abierto afín de $Y$, el
diagrama
\[
  \xymatrix{
    \HH^\bullet(u^{-1}(V),\sh{F})\ar[r]^\sim\ar[d] &
    \HH^\bullet(\Gamma(V,\sh{K}^\bullet))\ar[d]\\
    \HH^\bullet(u^{-1}(V'),\sh{F})\ar[r]^\sim &
    \HH^\bullet(\Gamma(V',\sh{K}^\bullet))
  }
\]
es conmutativo.
De lo anterior se concluye que los isomorfismos (\hyperref[III.1.4.10.1]{1.4.10.1}) definen un isomorfismo de $\sh{O}_Y$-módulos
\[
  \RR^\bullet u_*(\sh{F})\isoto\sh{H}^\bullet(\sh{K}^\bullet),
  \tag{1.4.10.3}\label{III.1.4.10.3}
\]
\oldpage[III]{92}
y, por consiguiente, $\RR^\bullet u_*(\sh{F})$ es cuasi-coherente.
\end{proof}

Además, de (\hyperref[III.1.4.10.3]{1.4.10.3}), (\hyperref[III.1.4.10.2]{1.4.10.2}) y (\hyperref[III.1.4.10.1]{1.4.10.1}) se deduce que:
\begin{corollary}[1.4.11]
\label{III.1.4.11}
Bajo las hipótesis de la Proposición~\sref{III.1.4.10}, para todo abierto afín $V$ de $Y$, el homomorfismo canónico
\[
  \HH^q(u^{-1}(V),\sh{F})\to\Gamma(V,\RR^q u_*(\sh{F}))
  \tag{1.4.11.1}\label{III.1.4.11.1}
\]
es un isomorfismo para todo $q\geq 0$.
\end{corollary}

\begin{corollary}[1.4.12]
\label{III.1.4.12}
Supóngase que se satisfacen las hipótesis de la Proposición~\sref{III.1.4.10} y, además, que $Y$ es cuasicompacto.
Entonces existe un entero $r>0$ tal que, para todo $\sh{O}_X$-módulo cuasi-coherente $\sh{F}$ y todo entero $q>r$, se tiene $\RR^q u_*(\sh{F})=0$.
Si $Y$ es afín, puede tomarse como $r$ un entero tal que exista un recubrimiento de $X$ formado por $r$ abiertos afines.
\end{corollary}

\begin{proof}
Como puede recubrirse $Y$ por un número finito de abiertos afines, en virtud del Corolario~\sref{III.1.4.11} puede reducirse la demostración a la segunda afirmación.
Si $\mathfrak{U}$ es un recubrimiento de $X$ por $r$ abiertos afines, se tiene $\HH^q(\mathfrak{U},\sh{F})=0$ para $q>r$, puesto que las cocadenas de $C^q(\mathfrak{U},\sh{F})$ son alternadas; la afirmación se deduce entonces de la Proposición~\sref{III.1.4.1}.
\end{proof}

\begin{corollary}[1.4.13]
\label{III.1.4.13}
Supóngase que se satisfacen las hipótesis de la Proposición~\sref{III.1.4.10} y, además, que $Y=\Spec(A)$ es afín.
Entonces, para todo $\sh{O}_X$-módulo cuasi-coherente $\sh{F}$ y todo $f\in A$, se tiene
\[
  \Gamma(Y_f,\RR^q u_*(\sh{F}))=(\Gamma(Y,\RR^q u_*(\sh{F})))_f
\]
salvo isomorfismo canónico.
\end{corollary}

\begin{proof}
Esto se deduce del hecho de que $\RR^q u_*(\sh{F})$ es un $\sh{O}_Y$-módulo cuasi-coherente \sref[I]{I.1.3.7}.
\end{proof}

\begin{proposition}[1.4.14]
\label{III.1.4.14}
Sean $f:X\to Y$ un morfismo separado y cuasicompacto y $g:Y\to Z$ un morfismo afín.
Para todo $\sh{O}_X$-módulo cuasi-coherente $\sh{F}$, el homomorfismo canónico $\RR^p(g\circ f)_*(\sh{F})\to g_*(\RR^p f_*(\sh{F}))$ (\textbf{0},~12.2.5.2) es biyectivo para todo $p$.
\end{proposition}

\begin{proof}
Para todo abierto afín $W$ de $Z$, $g^{-1}(W)$ es un abierto afín de $Y$.
El homomorfismo de prehaces que define el homomorfismo canónico
\[
  \RR^p(g\circ f)_*(\sh{F})\to g_*(\RR^p f_*(\sh{F}))
\]
\sref[0]{0.12.2.5} es, por tanto, biyectivo en virtud del Corolario~\sref{III.1.4.11}.
\end{proof}

\begin{proposition}[1.4.15]
\label{III.1.4.15}
Sean $u:X\to Y$ un morfismo separado cuasicompacto y $v:Y'\to Y$ un morfismo plano de preesquemas \sref[0]{0.6.7.1}; sea $u'=u_{(Y')}$, de modo que se tiene el diagrama conmutativo
\[
  \xymatrix{
    X\ar[d]_u &
    X'=X_{(Y')}\ar[l]_{v'}\ar[d]^{u'}\\
    Y &
    Y'\ar[l]_{v}.
  }
  \tag{1.4.15.1}
  \label{III.1.4.15.1}
\]
Entonces, para todo $\sh{O}_X$-módulo cuasi-coherente $\sh{F}$, $\RR^q u_*'(\sh{F}')$ es canónicamente isomorfo a $\RR^q u_*(\sh{F})\otimes_{\sh{O}_Y}\sh{O}_{Y'}=v^*(\RR^q u_*(\sh{F}))$ para todo $q\geq 0$, donde $\sh{F}'={v'}^*(\sh{F})=\sh{F}\otimes_{\sh{O}_Y}\sh{O}_{Y'}$.
\end{proposition}

\oldpage[III]{93}
\begin{proof}
El homomorfismo canónico $\rho:\sh{F}\to v_*'({v'}^*(\sh{F}))$ (\textbf{0},~4.4.3.2) define por funtorialidad un homomorfismo
\[
  \RR^q u_*(\sh{F})\to\RR^q u_*(v_*'(\sh{F}')).
  \tag{1.4.15.2}\label{III.1.4.15.2}
\]

Por otra parte, poniendo $w=u\circ v'=v\circ u'$, se tienen los homomorfismos canónicos (\textbf{0}~,12.2.5.1 y 12.2.5.2)
\[
  \RR^q u_*(v_*'(\sh{F}'))\to\RR^q w_*(\sh{F}')\to v_*(\RR^q u_*'(\sh{F}')).
  \tag{1.4.15.3}\label{III.1.4.15.3}
\]

Componiendo (\hyperref[III.1.4.15.3]{1.4.15.3}) y (\hyperref[III.1.4.15.2]{1.4.15.2}), se obtiene un homomorfismo
\[
  \psi:\RR^q u_*(\sh{F})\to v_*(\RR^q u_*'(\sh{F}')),
\]
y se obtiene finalmente un homomorfismo canónico (cuya definición no impone \emph{hipótesis alguna} sobre $v$)
\[
  \psi^\sharp:v^*(\RR^q u_*(\sh{F}))\to\RR^q u_*'(\sh{F}'),
  \tag{1.4.15.4}
\]
y hay que probar que es un isomorfismo cuando $v$ es \emph{plano}.
Es claro que la cuestión es local sobre $Y$ y $Y'$; puede suponerse, por tanto, que $Y=\Spec(A)$ y $Y'=\Spec(B)$; se utilizará además el siguiente lema:
\begin{lemma}[1.4.15.5]
\label{III.1.4.15.5}
Sean $\vphi:A\to B$ un homomorfismo de anillos, $Y=\Spec(A)$, $X=\Spec(B)$, $f:X\to Y$ el morfismo correspondiente a $\vphi$, y $M$ un $B$-módulo.
Para que el $\sh{O}_X$-módulo $\widetilde{M}$ sea $f$-plano \sref[0]{0.6.7.1}, es necesario y suficiente que $M$ sea un $A$-módulo plano.
En particular, para que el morfismo $f$ sea plano, es necesario y suficiente que $B$ sea un $A$-módulo plano.
\end{lemma}

Esto se deduce de la definición \sref[0]{0.6.7.1} y de \sref[0]{0.6.3.3}, teniendo en cuenta \sref[I]{I.1.3.4}.

Así pues, de (\hyperref[III.1.4.11.1]{1.4.11.1}) y de las definiciones de los homomorfismos (\hyperref[III.1.4.15.3]{1.4.15.3}) (cf.~\sref[0]{0.12.2.5}) se deduce que $\psi$ corresponde entonces al morfismo compuesto
\[
  \HH^q(X,\sh{F})\xrightarrow{\rho_q}\HH^q(X,v_*'({v'}^*(\sh{F})))\xrightarrow{\theta_q}\HH^q(X',{v'}^*(v_*'({v'}^*(\sh{F}))))\xrightarrow{\sigma_q}\HH^q(X',{v'}^*(\sh{F})),
\]
donde $\rho_q$ y $\sigma_q$ son los homomorfismos de cohomología correspondientes a los morfismos canónicos $\rho$ y $\sigma:{v'}^*(v_*'(\sh{G}'))\to\sh{G}'$, y $\theta_q$ es el $\vphi$-morfismo (\textbf{0},~12.1.3.1) relativo al $\sh{O}_X$-módulo $v_*'({v'}^*(\sh{F}))$.
Pero, por la funtorialidad de $\theta_q$, se tiene el diagrama conmutativo
\[
  \xymatrix{
    \HH^q(X,\sh{F})\ar[rr]^{\rho_q}\ar[dd]_{\theta_q} & &
    \HH^q(X,v_*'({v'}^*(\sh{F})))\ar[dd]^{\theta_q}\\\\
    \HH^q(X',{v'}^*(\sh{F}))\ar[rr]^{{v'}^*(\rho_q)} & &
    \HH^q(X',{v'}^*(v_*'({v'}^*(\sh{F})))),
  }
\]
\oldpage[III]{94}
y, como por definición \sref[0]{0.4.4.3} ${v'}^*(\rho)$ es el inverso de $\sigma$, se ve que el morfismo compuesto considerado anteriormente no es finalmente otro que $\theta_q$; por consiguiente, $\psi^\sharp$ es el $B$-homomorfismo asociado $\HH^q(X,\sh{F})\otimes_A B\to\HH^q(X',\sh{F}')$.
Como $u$ es cuasicompacto, $X$ es una unión finita de abiertos afines $U_i$ ($1\leq i\leq r$); sea $\mathfrak{U}$ el recubrimiento $(U_i)$.
Como $v$ es un morfismo afín, también lo es $v'$ \sref[II]{II.1.6.2}[iii]; por consiguiente, los $U_i'={v'}^{-1}(U_i)$ forman un recubrimiento abierto afín $\mathfrak{U}'$ de $X'$.
Se sabe entonces \sref[0]{0.12.1.4.2} que el diagrama
\[
  \xymatrix{
    \HH^q(\mathfrak{U},\sh{F})\ar[r]^{\theta_q}\ar[d] &
    \HH^q(\mathfrak{U}',\sh{F}')\ar[d]\\
    \HH^q(X,\sh{F})\ar[r]^{\theta_q} &
    \HH^q(X',\sh{F}')
  }
\]
es conmutativo, y las flechas verticales son isomorfismos puesto que $X$ y $X'$ son esquemas \sref{III.1.4.1}.
Por consiguiente, basta probar que el $\vphi$-morfismo canónico $\theta_q:\HH^q(\mathfrak{U},\sh{F})\to\HH^q(\mathfrak{U}',\sh{F}')$ es tal que el $B$-homomorfismo asociado
\[
  \HH^q(\mathfrak{U},\sh{F})\otimes_A B\to\HH^q(\mathfrak{U}',\sh{F}')
\]
sea un isomorfismo.
Para toda sucesión $\mathbf{s}=(i_k)_{0\leq k\leq p}$ de $p+1$ índices de $[1,r]$, pónganse $U_\mathbf{s}=\bigcap_{k=0}^p U_{i_k}$, $U_\mathbf{s}'=\bigcap_{k=0}^p U_{i_k}'={v'}^{-1}(U_\mathbf{s})$, $M_\mathbf{s}=\Gamma(U_\mathbf{s},\sh{F})$ y $M_\mathbf{s}'=\Gamma(U_\mathbf{s}',\sh{F}')$.
La aplicación canónica $M_\mathbf{s}\otimes_A B\to M_\mathbf{s}'$ es un isomorfismo \sref[I]{I.1.6.5}; por tanto, la aplicación canónica $C^p(\mathfrak{U},\sh{F})\otimes_A B\to C^p(\mathfrak{U}',\sh{F}')$ es un isomorfismo, mediante el cual $d\otimes 1$ se identifica con la aplicación coborde $C^p(\mathfrak{U}',\sh{F}')\to C^{p+1}(\mathfrak{U}',\sh{F}')$.
Como $B$ es un $A$-módulo \emph{plano}, de la definición de los módulos de cohomología se deduce que la aplicación canónica $\HH^q(\mathfrak{U},\sh{F})\otimes_A B\to\HH^q(\mathfrak{U}',\sh{F}')$ es un isomorfismo \sref[0]{0.6.1.1}.
Este resultado se generalizará más adelante en \EGAUnavailableHyperref{section:III.6}{\textsection6}.
\end{proof}

\begin{corollary}[1.4.16]
\label{III.1.4.16}
Sean $A$ un anillo, $X$ un $A$-esquema de tipo finito y $B$ una $A$-álgebra fielmente plana sobre $A$.
Para que $X$ sea afín, es necesario y suficiente que $X\otimes_A B$ lo sea.
\end{corollary}

\begin{proof}
La condición es evidentemente necesaria \sref[I]{I.3.2.2}; se probará que es suficiente.
Como $X$ es separado sobre $A$ y el morfismo $\Spec(B)\to\Spec(A)$ es plano, de la Proposición~\sref{III.1.4.15} se deduce que
\[
\label{III.1.4.16.1}
  \HH^i(X\otimes_A B,\sh{F}\otimes_A B)=\HH^i(X,\sh{F})\otimes_A B
  \tag{1.4.16.1}
\]
para todo $i\geq 0$ y todo $\sh{O}_X$-módulo cuasi-coherente $\sh{F}$.
Si $X\otimes_A B$ es afín, el miembro izquierdo de (\hyperref[III.1.4.16.1]{1.4.16.1}) es nulo para $i=1$; por tanto, también lo es $\HH^1(X,\sh{F})$, puesto que $B$ es un $A$-módulo fielmente plano.
Como $X$ es un esquema cuasicompacto, la demostración termina por el criterio de Serre \sref[II]{II.5.2.1}.
\end{proof}

\begin{proposition}[1.4.17]
\label{III.1.4.17}
Sea $X$ un preesquema y sea $0\to\sh{F}\xrightarrow{u}\sh{G}\xrightarrow{v}\sh{H}\to 0$ una sucesión exacta de $\sh{O}_X$-módulos.
Si $\sh{F}$ y $\sh{H}$ son cuasi-coherentes, también lo es $\sh{G}$.
\end{proposition}

\oldpage[III]{95}
\begin{proof}
La cuestión es local sobre $X$, de modo que puede suponerse que $X=\Spec(A)$ es afín; basta entonces probar que $\sh{G}$ satisface las condiciones (d1) y (d2) de \sref[I]{I.1.4.1} (con $V=X$).
La verificación de (d2) es inmediata, pues, si $t\in\Gamma(X,\sh{G})$ se anula al restringirlo a $D(f)$, lo mismo ocurre con su imagen $v(t)\in\Gamma(X,\sh{H})$; por tanto, existe un $m>0$ tal que $f^m v(t)=v(f^m t)=0$ \sref[I]{I.1.4.1}, y, como $\Gamma$ es exacto por la izquierda, $f^m t=u(s)$, donde $s\in\Gamma(X,\sh{F})$; como $u$ es inyectivo, la restricción de $s$ a $D(f)$ es nula, luego por \sref[I]{I.1.4.1} existe un entero $n>0$ tal que $f^n s=0$; se deduce finalmente que $f^{m+n}t=u(f^n s)=0$.

Compruébese ahora (d1); sea $t'\in\Gamma(D(f),\sh{G})$; como $\sh{H}$ es cuasi-coherente, existe un entero $m$ tal que $f^m v(t')=v(f^m t')$ se prolonga a una sección $z\in\Gamma(X,\sh{H})$ \sref[I]{I.1.4.1}.
Pero, en virtud del Teorema~\sref{III.1.3.1} (o de (\textbf{I},~5.1.9.2)) aplicado al $\sh{O}_X$-módulo cuasi-coherente $\sh{F}$, la sucesión $\Gamma(X,\sh{G})\to\Gamma(X,\sh{H})\to 0$ es exacta; existe, pues, $t\in\Gamma(X,\sh{G})$ tal que $z=v(t)$; se ve así que $v(f^m t'-t'')=0$, denotando por $t''$ la restricción de $t$ a $D(f)$; por tanto, $f^m t'-t''=u(s')$, donde $s'\in\Gamma(D(f),\sh{F})$.
Pero, como $\sh{F}$ es cuasi-coherente, existe un entero $n>0$ tal que $f^n s'$ se prolonga a una sección $s\in\Gamma(X,\sh{F})$; como $f^{m+n}t'-f^n t''=u(f^n s')$, se ve que $f^{m+n}t'$ es la restricción a $D(f)$ de una sección $f^n t+u(s)\in\Gamma(X,\sh{G})$, lo que termina la demostración.
\end{proof}

\section{Estudio cohomológico de los morfismos proyectivos}
\label{section:III.2}


\subsection{Cálculos explícitos de ciertos grupos de cohomología}
\label{subsection:III.2.1}

\begin{env}[2.1.1]
\label{III.2.1.1}
\oldpage[III]{95}
Sea $X$ un preesquema y sea $\sh{L}$ un $\sh{O}_X$-módulo invertible; consideremos el anillo graduado \sref[0]{0.5.4.6}
\[
\label{III.2.1.1.1}
  S = \Gamma_*(X,\sh{L}) = \bigoplus_{n\in\bb{Z}} \Gamma(X,\sh{L}^{\otimes n}).
  \tag{2.1.1.1}
\]

Sea $(f_i)_{1\leq i\leq r}$ una familia finita de elementos \emph{homogéneos} de $S$, digamos $f_i\in S_{d_i}$; pongamos $U_i = X_{f_i}$, $U = \bigcup_i U_i$, y denotemos por $\mathfrak{U}$ el recubrimiento $(U_i)$ de $U$.
Para todo $\sh{O}_X$-módulo cuasicoherente $\sh{F}$, ponemos
\[
\label{III.2.1.1.2}
  \HH^\bullet(U,\sh{F}(*)) = \bigoplus_{n\in\bb{Z}} \HH^\bullet(U,\sh{F}\otimes\sh{L}^{\otimes n})
  \tag{2.1.1.2}
\]
\[
\label{III.2.1.1.3}
  \HH^\bullet(\mathfrak{U},\sh{F}(*)) = \bigoplus_{n\in\bb{Z}} \HH^\bullet(\mathfrak{U},\sh{F}\otimes\sh{L}^{\otimes n}).
  \tag{2.1.1.3}
\]

Los grupos abelianos (\hyperref[III.2.1.1.2]{2.1.1.2}) y (\hyperref[III.2.1.1.3]{2.1.1.3}) son \emph{bigraduados}, tomando
\[
  (\HH^\bullet(U,\sh{F}(*)))_{mn} = \HH^m(U,\sh{F}\otimes\sh{L}^{\otimes n})
\]
y una definición análoga para (\hyperref[III.2.1.1.3]{2.1.1.3}).
Para la segunda graduación, está claro que estos grupos son $S$-módulos graduados, como se deduce, por ejemplo, del hecho de que $\sh{F}\mapsto\HH^m(U,\sh{F})$ y $\sh{F}\mapsto\HH^m(\mathfrak{U},\sh{F})$ son funtores.
\end{env}

\begin{env}[2.1.2]
\label{III.2.1.2}
Consideremos ahora el $S$-módulo graduado \sref[0]{0.5.4.6}
\[
\label{III.2.1.2.1}
  M = \Gamma_*(\sh{L},\sh{F}) = \HH^0(X,\sh{F}(*)) = \bigoplus_{n\in\bb{Z}} \Gamma(X,\sh{F}\otimes\sh{L}^{\otimes n}).
  \tag{2.1.2.1}
\]

\oldpage[III]{96}
Si $X$ es un preesquema cuyo espacio subyacente es noetheriano, o un esquema cuasicompacto, se deduce de \sref[I]{I.9.3.1} que, poniendo como de costumbre $U_{i_0 i_1\cdots i_p} = \bigcap_{k=0}^p U_{i_k}$, se tiene, salvo isomorfismo canónico,
\[
  \Gamma(U_{i_0 i_1\cdots i_p},\sh{F}(*)) = \HH^0(U_{i_0 i_1\cdots i_p},\sh{F}(*)) = M_{f_{i_0}f_{i_1}\cdots f_{i_p}}.
\]

Además, con la notación de (\hyperref[III.1.2.2]{1.2.2}), podemos identificar $M_{f_{i_0}f_{i_1}\cdots f_{i_p}}$ con $\varinjlim_n M_{i_0 i_1\cdots i_p}^{(n)}$.
Esta identificación es un isomorfismo de $S$-módulos \emph{graduados} si definimos el grado de un elemento homogéneo $z\in\varinjlim_n M_{i_0 i_1\cdots i_p}^{(n)}$ como sigue: $z$ es la imagen canónica de un elemento homogéneo $x\in M_{i_0 i_1\cdots i_p}^{(n)} = M$ de grado $m$, y tomamos como grado de $z$ el número $m-n(d_{i_0}+d_{i_1}+\cdots+d_{i_p})$.
Teniendo en cuenta la definición de los homomorfismos $\vphi_{kh}:M_{i_0 i_1\cdots i_p}^{(h)}\to M_{i_0 i_1\cdots i_p}^{(k)}$ (\hyperref[III.1.2.2]{1.2.2}), se ve inmediatamente que esta definición no depende del ``representante'' $x$ de $z$ elegido.
Denotando, como en (\hyperref[III.1.2.2]{1.2.2}), por $C_n^p(M)$ el conjunto de aplicaciones alternadas de $[1,r]^{p+1}$ en $M$ (para todo $n$), definimos de manera análoga una estructura de $S$-módulo graduado sobre $\varinjlim_n C_n^p(M)$; como en (\hyperref[III.1.2.2]{1.2.2}), se tiene entonces
\[
\label{III.2.1.2.2}
  C^p(\mathfrak{U},\sh{F}(*)) = \varinjlim_n C_n^p(M)
  \tag{2.1.2.2}
\]
donde el isomorfismo entre ambos miembros \emph{respeta los grados}.
Se tiene también, como en (\hyperref[III.1.2.2]{1.2.2}),
\[
\label{III.2.1.2.3}
  C^p(\mathfrak{U},\sh{F}(*)) = C^{p+1}((\mathbf{f}),M) = \varinjlim_n K^{p+1}(\mathbf{f}^n,M)
  \tag{2.1.2.3}
\]
donde el isomorfismo \emph{conserva los grados}: el grado de un elemento de $\varinjlim_n K^{p+1}(\mathbf{f}^n,M)$ que sea la imagen canónica de una cocadena $g\in K^{p+1}(\mathbf{f}^n,M)$ cuyos valores $g(i_0,\dots,i_p)$ pertenezcan a una misma componente homogénea $M_m$ de $M$ es $m-n(d_{i_0}+\cdots+d_{i_p})$, independientemente de la elección de esta cocadena como representante del elemento considerado.

Como los isomorfismos anteriores son compatibles con los operadores de coborde, se concluye, como en (\hyperref[III.1.2.2]{1.2.2}), que:
\end{env}

\begin{proposition}[2.1.3]
\label{III.2.1.3}
Sea $X$ un preesquema cuyo espacio subyacente es noetheriano, o un esquema cuasicompacto.
Existe un isomorfismo canónico, funtorial en $\sh{F}$,
\[
\label{III.2.1.3.1}
  \HH^p(\mathfrak{U},\sh{F}(*)) \simto \HH^{p+1}((\mathbf{f}),M)
  \qquad\text{para todo } p\geq 1.
  \tag{2.1.3.1}
\]
Además, existe una sucesión exacta, funtorial en $\sh{F}$,
\[
\label{III.2.1.3.2}
  0 \to \HH^0((\mathbf{f}),M) \to M \to \HH^0(\mathfrak{U},\sh{F}(*)) \to \HH^1((\mathbf{f}),M) \to 0.
  \tag{2.1.3.2}
\]
Todos los homomorfismos introducidos son de grado $0$ para las estructuras de $S$-módulo graduado, donde $S$ es el anillo (\hyperref[III.2.1.1.1]{2.1.1.1}).
\end{proposition}

\begin{corollary}[2.1.4]
\label{III.2.1.4}
\oldpage[III]{97}
Si $X$ es un esquema cuasicompacto y los $U_i=X_{f_i}$ son afines, existe un isomorfismo canónico, funtorial en $\sh{F}$ y de grado $0$,
\[
\label{III.2.1.4.1}
  \HH^p(U,\sh{F}(*)) \simto \HH^{p+1}((\mathbf{f}),M)
  \qquad\text{para todo } p\geq 1,
  \tag{2.1.4.1}
\]
y una sucesión exacta, funtorial en $\sh{F}$,
\[
\label{III.2.1.4.2}
  0 \to \HH^0((\mathbf{f}),M) \to M \to \HH^0(U,\sh{F}(*)) \to \HH^1((\mathbf{f}),M) \to 0,
  \tag{2.1.4.2}
\]
en la que todos los homomorfismos son de grado $0$.
\end{corollary}

\begin{proof}
Basta aplicar \sref{III.1.4.1} al resultado de \sref{III.2.1.3}.
\end{proof}

La proposición ``local'' análoga a \sref{III.2.1.3} es la siguiente.

\begin{proposition}[2.1.5]
\label{III.2.1.5}
Sea $S$ un anillo graduado positivamente, sean los $f_i$ ($1\leq i\leq r$) elementos homogéneos de $S_+$ de grados $d_i$, y sea $M$ un $S$-módulo graduado.
Sea $X=\Proj(S)$ el espectro primo homogéneo de $S$, y pongamos $U_i=\mathrm{D}_+(f_i)$, $U=\bigcup_i U_i$ y $\HH^\bullet(U,\widetilde{M}(*))=\bigoplus_{n\in\bb{Z}}\HH^\bullet(U,(M(n))^\sim)$.
Existen isomorfismos canónicos, funtoriales en $M$ y de grado $0$ para las estructuras de $S$-módulo graduado,
\[
\label{III.2.1.5.1}
  \HH^p(U,\widetilde{M}(*)) \simto \HH^{p+1}((\mathbf{f}),M)
  \qquad\text{para todo } p\geq 1,
  \tag{2.1.5.1}
\]
y una sucesión exacta, funtorial en $M$,
\[
\label{III.2.1.5.2}
  0 \to \HH^0((\mathbf{f}),M) \to M \to \HH^0(U,\widetilde{M}(*)) \to \HH^1((\mathbf{f}),M) \to 0,
  \tag{2.1.5.2}
\]
en la que todos los homomorfismos son de grado $0$.
\end{proposition}

\begin{proof}
En efecto, por definición \sref[II]{II.2.5.2},
\[
  \Gamma(U_{i_0i_1\cdots i_p},(M(n))^\sim)
  =(M_{f_{i_0}\cdots f_{i_p}})_n,
\]
y, por tanto, $\Gamma(U_{i_0i_1\cdots i_p},\widetilde{M}(*))=M_{f_{i_0}\cdots f_{i_p}}$.
El resto del razonamiento es el mismo que en la demostración de \sref{III.2.1.4}, teniendo en cuenta que $X$ es un esquema.
\end{proof}

\begin{env}[2.1.6]
\label{III.2.1.6}
\emph{Observaciones.} ---
(i) Bajo las hipótesis de \sref{III.2.1.5}, los funtores $\Gamma(U_{i_0i_1\cdots i_p},\widetilde{M}(*))$ son exactos en $M$, por \sref[0]{0.1.3.2}.
El mismo razonamiento que en (\hyperref[III.1.2.5]{1.2.5}) muestra, por tanto, que, si $0\to M'\to M\to M''\to0$ es una sucesión exacta de $S$-módulos graduados cuyos homomorfismos tienen grado $0$, entonces para todo $p\geq0$ existen diagramas conmutativos
\[
\label{III.2.1.6.1}
  \xymatrix{
    \HH^p(U,\widetilde{M''}(*)) \ar[r]^-\partial \ar[d] &
    \HH^{p+1}(U,\widetilde{M'}(*)) \ar[d] \\
    \HH^{p+1}((\mathbf{f}),M'') \ar[r]^-\partial &
    \HH^{p+2}((\mathbf{f}),M')
  }
  \tag{2.1.6.1}
\]

(ii) La Proposición~\sref{III.2.1.5} es especialmente útil cuando $S$ es una $A$-álgebra engendrada por un número finito de elementos de grado $1$, con $A$ noetheriano; en ese caso, todo $\sh{O}_X$-módulo cuasicoherente es de la forma $\widetilde{M}$ \sref[II]{II.2.7.7}.
\end{env}

\begin{env}[2.1.7]
\label{III.2.1.7}
\oldpage[III]{98}
Aplicamos ahora \sref{III.2.1.5} al caso $S=A[\mathrm{T}_0,\dots,\mathrm{T}_r]$, donde $A$ es un anillo arbitrario y las $\mathrm{T}_i$ son indeterminadas, tomando $M=S$ y $f_i=\mathrm{T}_i$.
Quedamos así reducidos esencialmente a calcular $\HH^\bullet((\mathbf{T}),S)$, donde $\mathbf{T}=(\mathrm{T}_i)_{0\leq i\leq r}$.
\end{env}

\begin{lemma}[2.1.8]
\label{III.2.1.8}
Si $S=A[\mathrm{T}_0,\dots,\mathrm{T}_r]$, entonces, con $\mathbf{T}=(\mathrm{T}_i)_{0\leq i\leq r}$,
\[
\label{III.2.1.8.1}
  \HH^i(\mathbf{T}^n,S)=0
  \qquad\text{si }i\neq r+1,
  \tag{2.1.8.1}
\]
\[
\label{III.2.1.8.2}
  \HH^{r+1}(\mathbf{T}^n,S)=S/(\mathbf{T}^n).
  \tag{2.1.8.2}
\]
Por tanto, el $A$-módulo $\HH^{r+1}(\mathbf{T}^n,S)$ tiene una $A$-base formada por las clases módulo $(\mathbf{T}^n)$ de los monomios
\[
  \mathrm{T}_{\mathbf{p}}=\mathrm{T}_0^{p_0}\cdots\mathrm{T}_r^{p_r},
  \qquad \mathbf{p}=(p_0,\dots,p_r),\quad 0\leq p_i<n\ \text{para todo }i.
\]
\end{lemma}

\begin{proof}
Esto se deduce inmediatamente de (\hyperref[III.1.1.3.5]{1.1.3.5}) y de la Proposición~\sref{III.1.1.4}, cuyas hipótesis se satisfacen trivialmente.
\end{proof}

\begin{env}[2.1.9]
\label{III.2.1.9}
Pasando al límite inductivo respecto de $n$, las relaciones (\hyperref[III.2.1.8.1]{2.1.8.1}) dan $\HH^i((\mathbf{T}),S)=0$ para $i\neq r+1$.
Para $i=r+1$, el sistema inductivo consta de los módulos $S/(\mathbf{T}^n)$, siendo el homomorfismo
\[
  \vphi_{nm}:S/(\mathbf{T}^n)\longrightarrow S/(\mathbf{T}^m)
  \qquad (0\leq n\leq m)
\]
la multiplicación por $(\mathrm{T}_0\cdots\mathrm{T}_r)^{m-n}$.
Para $n\geq\sup(p_i)_{0\leq i\leq r}$, denotemos por $\xi_{\mathbf{p}}^{(n)}=\xi_{p_0,\dots,p_r}^{(n)}$ la clase de $\mathrm{T}_0^{n-p_0}\cdots\mathrm{T}_r^{n-p_r}$ módulo $(\mathbf{T}^n)$.
Entonces $\vphi_{nm}(\xi_{\mathbf{p}}^{(n)})=\xi_{\mathbf{p}}^{(m)}$, de modo que estos elementos tienen la misma imagen canónica $\xi_{\mathbf{p}}=\xi_{p_0,\dots,p_r}$ en el límite inductivo $\HH^{r+1}((\mathbf{T}),S)$.
Por la definición del grado en (\hyperref[III.2.1.2]{2.1.2}), el grado de $\xi_{\mathbf{p}}$ es $-|\mathbf{p}|=-(p_0+p_1+\cdots+p_r)$.
Está claro que los $\xi_{\mathbf{p}}^{(n)}$ con $0<p_i\leq n$ para $0\leq i\leq r$ forman una base de $S/(\mathbf{T}^n)$.
Se deduce, pues, inmediatamente de \sref{III.2.1.8}:
\end{env}

\begin{corollary}[2.1.10]
\label{III.2.1.10}
Con la notación de \sref{III.2.1.8},
\[
\label{III.2.1.10.1}
  \HH^i((\mathbf{T}),S)=0
  \qquad\text{para }i\neq r+1,
  \tag{2.1.10.1}
\]
y $\HH^{r+1}((\mathbf{T}),S)$ es un $A$-módulo libre con base los elementos $\xi_{p_0,\dots,p_r}$ tales que $p_i>0$ para $0\leq i\leq r$.
\end{corollary}

\begin{env}[2.1.11]
\label{III.2.1.11}
\emph{Observación.} ---
Sea $N$ un $A$-módulo cualquiera y sea $M=S\otimes_A N$.
El razonamiento de \sref{III.2.1.8} muestra más generalmente que
\[
\label{III.2.1.11.1}
  \HH^i(\mathbf{T}^n,M)=0
  \qquad\text{si }i\neq r+1,
  \tag{2.1.11.1}
\]
\[
\label{III.2.1.11.2}
  \HH^{r+1}(\mathbf{T}^n,M)=(S/(\mathbf{T}^n))\otimes_A N.
  \tag{2.1.11.2}
\]
En efecto, esta última fórmula se deduce directamente de (\hyperref[III.1.1.3.5]{1.1.3.5}).
Además, está claro que
\[
  M/(\mathrm{T}_0^nM+\cdots+\mathrm{T}_{i-1}^nM)
  =\bigl(S/(\mathrm{T}_0^nS+\cdots+\mathrm{T}_{i-1}^nS)\bigr)\otimes_A N,
\]
puesto que el ideal $\mathrm{T}_0^nS+\cdots+\mathrm{T}_{i-1}^nS$ es un sumando directo del $A$-módulo $S$.
Esto permite aplicar \sref{III.1.1.4} a $M$ y da (\hyperref[III.2.1.11.1]{2.1.11.1}).
\end{env}

Combinando \sref{III.2.1.10} y \sref{III.2.1.5}, se obtiene:

\begin{proposition}[2.1.12]
\label{III.2.1.12}
Sea $A$ un anillo, sea $r$ un entero $>0$ y sea $X=\bb{P}_A^r$ \sref[II]{II.4.1.1}.
Entonces:
\begin{enumerate}
  \item[{\rm(i)}] $\HH^i(X,\sh{O}_X(*))=0$ para $i\neq0,r$.
  \item[{\rm(ii)}] El homomorfismo canónico $\alpha:S\to\HH^0(X,\sh{O}_X(*))$ \sref[II]{II.2.6.2} es biyectivo.
\oldpage[III]{99}
  \item[{\rm(iii)}] $\HH^r(X,\sh{O}_X(*))$ es un $A$-módulo libre con base los elementos $\xi_{p_0,\dots,p_r}$, donde $p_i>0$ para $0\leq i\leq r$; el elemento $\xi_{p_0,\dots,p_r}$ tiene grado $-|\mathbf{p}|=-(p_0+\cdots+p_r)$, y
  \[
    \mathrm{T}_i\xi_{p_0,\dots,p_r}
    =\xi_{p_0,\dots,p_{i-1},p_i-1,p_{i+1},\dots,p_r}.
  \]
\end{enumerate}
\end{proposition}

\begin{proof}
En la sucesión exacta (\hyperref[III.2.1.5.2]{2.1.5.2}) aplicada a $M=S=A[\mathrm{T}_0,\dots,\mathrm{T}_r]$, se tiene $\HH^0((\mathbf{T}),S)=0$ y $\HH^1((\mathbf{T}),S)=0$ por \sref{III.2.1.10.1}; además, la Proposición~\sref{III.2.1.5} se aplica con $U=X$, puesto que $X$ es la unión de los $\mathrm{D}_+(\mathrm{T}_i)$ \sref[II]{II.2.3.14}.
Queda identificar la aplicación $S\to\HH^0(X,\sh{O}_X(*))$ de (\hyperref[III.2.1.5.2]{2.1.5.2}) con la aplicación canónica $\alpha$.
Esto se deduce de la identificación canónica de $\HH^0(U,\sh{O}_X(*))$ y $\HH^0(\mathfrak{U},\sh{O}_X(*))$.
\end{proof}

\begin{corollary}[2.1.13]
\label{III.2.1.13}
Los únicos pares $(i,n)$ para los cuales puede ocurrir que $\HH^i(X,\sh{O}_X(n))\neq0$ son $i=0$, $n\geq0$, e $i=r$, $n\leq-(r+1)$.
\end{corollary}

Si $A\neq0$, entonces se tiene efectivamente $\HH^i(X,\sh{O}_X(n))\neq0$ para todos los pares enumerados en \sref{III.2.1.13}.
Esto se deduce de \sref{III.2.1.12}, puesto que $S_n\neq0$ para todo grado $n\geq0$.

En las aplicaciones de este capítulo, utilizaremos principalmente el siguiente resultado menos preciso.

\begin{corollary}[2.1.14]
\label{III.2.1.14}
Los $A$-módulos $\HH^i(X,\sh{O}_X(n))$ son libres de tipo finito; si $i>0$, se anulan para $n>0$.
\end{corollary}

\begin{proposition}[2.1.15]
\label{III.2.1.15}
Sea $Y$ un preesquema, sea $\sh{E}$ un $\sh{O}_Y$-módulo localmente libre de rango $r+1$, sea $X=\bb{P}(\sh{E})$ el fibrado proyectivo definido por $\sh{E}$, y sea $f:X\to Y$ el morfismo estructural.
Los únicos valores de $i$ y $n$ para los cuales puede ocurrir que $\RR^if_*(\sh{O}_X(n))\neq0$ son $i=0$, $n\geq0$, e $i=r$, $n\leq-(r+1)$.
Además, el homomorfismo canónico \sref[II]{II.3.3.2}
\[
  \alpha:\bb{S}_{\sh{O}_Y}(\sh{E})
  \longrightarrow \Gamma_*(\sh{O}_X)
  =\RR^0f_*(\sh{O}_X(*))
  =\bigoplus_{n\in\bb{Z}}f_*(\sh{O}_X(n))
\]
es un isomorfismo.
\end{proposition}

\begin{proof}
La cuestión es local sobre $Y$, por lo que podemos suponer que $Y$ es afín de anillo $A$ y que $\sh{E}=\widetilde{E}$, donde $E=A^{r+1}$.
Quedamos entonces reducidos inmediatamente a \sref{III.2.1.12}, teniendo en cuenta \sref{III.1.4.11}.
\end{proof}

\begin{env}[2.1.16]
\label{III.2.1.16}
\emph{Observación.} ---
Más adelante completaremos los resultados de \sref{III.2.1.15} demostrando las proposiciones siguientes.
Pongamos
\[
  \boldsymbol{\omega}=f^*\bigl(\bigwedge^{r+1}\sh{E}\bigr)(-r-1),
\]
que es un $\sh{O}_X$-módulo invertible.
Entonces:
\begin{enumerate}
  \item[{\rm(i)}] Existe un isomorfismo canónico
  \[
  \label{III.2.1.16.1}
    \rho:\sh{O}_Y \simto \RR^rf_*(\boldsymbol{\omega}).
    \tag{2.1.16.1}
  \]
  \item[{\rm(ii)}] El apareamiento dado por el producto cup \sref[0]{0.12.2.2}
  \[
  \label{III.2.1.16.2}
    \RR^rf_*(\sh{O}_X(n))
    \times \RR^0f_*(\boldsymbol{\omega}(-n))
    \longrightarrow \RR^rf_*(\boldsymbol{\omega}),
    \tag{2.1.16.2}
  \]
  compuesto con $\rho^{-1}$, define un \emph{isomorfismo} de $\RR^rf_*(\sh{O}_X(n))$ sobre el \emph{dual} del $\sh{O}_Y$-módulo localmente libre
  \[
    \RR^0f_*(\boldsymbol{\omega}(-n))
    =\bigl(\bigwedge^{r+1}\sh{E}\bigr)
      \otimes_{\sh{O}_Y}
      \bigl(\bb{S}_{\sh{O}_Y}(\sh{E})\bigr)_{-n}.
  \]
\end{enumerate}
\end{env}

\subsection{El teorema fundamental de los morfismos proyectivos}
\label{subsection:III.2.2}
\oldpage[III]{100}

\begin{theorem}[2.2.1]
\label{III.2.2.1}
(Serre).
Sea $Y$ un preesquema noetheriano, sea $f:X\to Y$ un morfismo propio y sea $\sh{L}$ un $\sh{O}_X$-módulo invertible amplio respecto de $f$.
Para todo $\sh{O}_X$-módulo $\sh{F}$, pongamos
\[
  \sh{F}(n)=\sh{F}\otimes_{\sh{O}_X}\sh{L}^{\otimes n}
  \qquad(n\in\bb{Z}).
\]
Entonces, para todo $\sh{O}_X$-módulo coherente $\sh{F}$:
\begin{enumerate}
  \item[{\rm(i)}] los $\sh{O}_Y$-módulos $\RR^qf_*(\sh{F})$ son coherentes;
  \item[{\rm(ii)}] existe un entero $N$ tal que, para $n\geq N$,
  \[
    \RR^qf_*(\sh{F}(n))=0
    \qquad\text{para todo }q>0;
  \]
  \item[{\rm(iii)}] existe un entero $N$ tal que, para $n\geq N$, el homomorfismo canónico
  \[
    f^*\bigl(f_*(\sh{F}(n))\bigr)\longrightarrow\sh{F}(n)
  \]
  es sobreyectivo.
\end{enumerate}
\end{theorem}

\begin{proof}
Observemos primero que, si el teorema es cierto tras sustituir $\sh{L}$ por $\sh{L}^{\otimes d}$, donde $d>0$, entonces es cierto en su forma original.
En efecto, escribiendo
\[
  \sh{F}(n)
  =\bigl(\sh{F}\otimes\sh{L}^{\otimes r}\bigr)
   \otimes\sh{L}^{\otimes hd},
   \qquad h>0,\quad 0\leq r<d,
\]
la hipótesis da, para cada $r$, un entero $N_r$ tal que (ii) y (iii) se cumplen para $\sh{F}\otimes\sh{L}^{\otimes r}$ siempre que $h\geq N_r$.
Tomando para $N$ el mayor de los enteros $dN_r$, las afirmaciones (ii) y (iii) se cumplen, por tanto, para $n\geq N$.

Podemos, por consiguiente, suponer que $\sh{L}$ es muy amplio respecto de $f$ \sref[II]{II.4.6.11}.
Existe entonces una inmersión abierta dominante sobre $Y$
\[
  i:X\longrightarrow P,
  \qquad P=\Proj(\sh{S}),
\]
donde $\sh{S}$ es una $\sh{O}_Y$-álgebra cuasicoherente graduada en grados positivos, $\sh{S}_1$ es de tipo finito y engendra $\sh{S}$, y $\sh{L}$ es isomorfo a $i^*(\sh{O}_P(1))$ \sref[II]{II.4.4.7}.
Como $f$ es propio, $i$ también es propio \sref[II]{II.5.4.4} y, por tanto, es un isomorfismo $X\simto P$.
Quedamos así reducidos al caso $X=\Proj(\sh{S})$ y $\sh{L}=\sh{O}_X(1)$.
El teorema es entonces consecuencia de la proposición siguiente.
\end{proof}

\begin{proposition}[2.2.2]
\label{III.2.2.2}
Sea $A$ un anillo noetheriano y sea $S$ una $A$-álgebra graduada en grados positivos tal que el $A$-módulo $S_1$ posee un sistema de $r+1$ generadores y engendra el álgebra $S$.
Pongamos $X=\Proj(S)$.
Para todo $\sh{O}_X$-módulo coherente $\sh{F}$:
\begin{enumerate}
  \item[{\rm(i)}] los $A$-módulos $\HH^q(X,\sh{F})$ son de tipo finito;
  \item[{\rm(ii)}] $\HH^q(X,\sh{F})=0$ para $q>r$;
  \item[{\rm(iii)}] existe un entero $N$ tal que, para $n\geq N$,
  \[
    \HH^q(X,\sh{F}(n))=0
    \qquad\text{para todo }q>0;
  \]
  \item[{\rm(iv)}] existe un entero $N$ tal que, para $n\geq N$, $\sh{F}(n)$ está engendrado por sus secciones sobre $X$.
\end{enumerate}
\end{proposition}

\begin{proof}
Mostremos primero cómo \sref{III.2.2.2} implica \sref{III.2.2.1}.
En el caso particular $X=\Proj(\sh{S})$ al que se redujo anteriormente \sref{III.2.2.1}, $Y$ es cuasicompacto.
Puede, por tanto, recubrirse por un número finito de abiertos afines $U_\alpha$, de anillos noetherianos, tales que la restricción de $\sh{S}_1$ a cada $U_\alpha$ esté engendrada por un número finito de secciones sobre $U_\alpha$.
Una vez conocida \sref{III.2.2.2}, basta tomar en las partes (ii) y (iii) de \sref{III.2.2.1} el mayor de los enteros correspondientes a los $U_\alpha$, teniendo en cuenta \sref{III.1.4.11} y \sref[II]{II.3.4.7}.

Para demostrar \sref{III.2.2.2}, observemos que $X$ se identifica con un subpreesquema cerrado de $P=\bb{P}_A^r$ \sref[II]{II.3.6.2}.
Además, si $j:X\to P$ es la inyección canónica, entonces $j_*(\sh{F})$ es un $\sh{O}_P$-módulo coherente y
\[
  j_*(\sh{F}(n))=(j_*\sh{F})(n)
\]
por \sref[II]{II.3.4.5} y \sref[II]{II.3.5.2}.
En virtud de (G, II, Corolario del Teorema~4.9.1), quedamos, por tanto, reducidos a demostrar \sref{III.2.2.2} en el caso particular
\[
  X=\bb{P}_A^r,
  \qquad S=A[T_0,\ldots,T_r].
\]

\oldpage[III]{101}

Como $X$ está recubierto por los $r+1$ abiertos afines $D_+(T_i)$, la afirmación (ii) se deduce de \sref{III.1.4.12}.
La afirmación (iv) ya fue demostrada en \sref[II]{II.2.7.9}.

Demostremos ahora simultáneamente (i) y (iii).
Se cumplen para $\sh{F}=\sh{O}_X(m)$ por \sref{III.2.1.13} y, por tanto, también cuando $\sh{F}$ es suma directa de un número finito de $\sh{O}_X$-módulos de la forma $\sh{O}_X(m_j)$.
Son trivialmente ciertas para $q>r$ por (ii).
Procedemos por inducción descendente sobre $q$.

El haz $\sh{F}$ es isomorfo a un cociente de una suma directa $\sh{E}$ de un número finito de haces $\sh{O}_X(m_j)$ \sref[II]{II.2.7.10}.
Existe, pues, una sucesión exacta
\[
  0\longrightarrow\sh{R}\longrightarrow\sh{E}
   \longrightarrow\sh{F}\longrightarrow0,
\]
donde $\sh{R}$ es coherente \sref[0]{0.5.3.3} y $\sh{E}$ satisface (i) y (iii).
Como el cambio de grado es exacto, para todo $n\in\bb{Z}$ se tiene también una sucesión exacta
\[
  0\longrightarrow\sh{R}(n)\longrightarrow\sh{E}(n)
   \longrightarrow\sh{F}(n)\longrightarrow0,
\]
y, por tanto, una sucesión exacta de cohomología
\[
  \HH^{q-1}(X,\sh{E}(n))
  \longrightarrow\HH^{q-1}(X,\sh{F}(n))
  \longrightarrow\HH^q(X,\sh{R}(n)).
\]
Como $\sh{E}(n)$ es una suma directa de haces $\sh{O}_X(n+m_j)$ \sref[II]{II.2.5.14}, el módulo $\HH^{q-1}(X,\sh{E}(n))$ es de tipo finito.
Análogamente, el módulo $\HH^q(X,\sh{R}(n))$ es de tipo finito por la hipótesis de inducción.
Como $A$ es noetheriano, se deduce que $\HH^{q-1}(X,\sh{F}(n))$ es de tipo finito, lo que demuestra (i).

Por la hipótesis de inducción podemos elegir $N$ de modo que $\HH^q(X,\sh{R}(n))=0$ para $n\geq N$.
Podemos elegir asimismo $N$ de modo que $\HH^{q-1}(X,\sh{E}(n))=0$ para $n\geq N$, puesto que $\sh{E}$ satisface (iii).
La sucesión exacta de cohomología da entonces
\[
  \HH^{q-1}(X,\sh{F}(n))=0
  \qquad(n\geq N),
\]
lo que termina la demostración.
\end{proof}

\begin{corollary}[2.2.3]
\label{III.2.2.3}
Bajo las hipótesis de \sref{III.2.2.1}, sea
\[
  \sh{F}\longrightarrow\sh{G}\longrightarrow\sh{H}
\]
una sucesión exacta de $\sh{O}_X$-módulos coherentes.
Existe un entero $N$ tal que, para $n\geq N$, la sucesión
\[
  f_*(\sh{F}(n))
  \longrightarrow f_*(\sh{G}(n))
  \longrightarrow f_*(\sh{H}(n))
\]
es exacta.
\end{corollary}

\begin{proof}
Sean $\sh{F}'$, $\sh{G}'$ y $\sh{G}''$, respectivamente, el núcleo, la imagen y el conúcleo de $\sh{F}\to\sh{G}$.
Entonces $\sh{G}'$ es el núcleo y $\sh{G}''$ la imagen de $\sh{G}\to\sh{H}$; sea $\sh{H}''$ el conúcleo de este último homomorfismo.
Todos estos $\sh{O}_X$-módulos son coherentes \sref[0]{0.5.3.4}.
Como el cambio de grado es exacto, basta demostrar que, para $n$ suficientemente grande, cada una de las sucesiones
\begin{align*}
  0&\longrightarrow f_*(\sh{F}'(n))
    \longrightarrow f_*(\sh{F}(n))
    \longrightarrow f_*(\sh{G}'(n))
    \longrightarrow0,\\
  0&\longrightarrow f_*(\sh{G}'(n))
    \longrightarrow f_*(\sh{G}(n))
    \longrightarrow f_*(\sh{G}''(n))
    \longrightarrow0,\\
  0&\longrightarrow f_*(\sh{G}''(n))
    \longrightarrow f_*(\sh{H}(n))
    \longrightarrow f_*(\sh{H}''(n))
    \longrightarrow0
\end{align*}
es exacta.
Por consiguiente, podemos suponer que
\[
  0\longrightarrow\sh{F}\longrightarrow\sh{G}
   \longrightarrow\sh{H}\longrightarrow0
\]
es exacta.
La sucesión de cohomología es entonces
\[
  0\longrightarrow f_*(\sh{F}(n))
  \longrightarrow f_*(\sh{G}(n))
  \longrightarrow f_*(\sh{H}(n))
  \longrightarrow \RR^1f_*(\sh{F}(n))
  \longrightarrow\cdots,
\]
y la conclusión se deduce de \sref{III.2.2.1}, (ii).
\end{proof}

\oldpage[III]{102}

\begin{corollary}[2.2.4]
\label{III.2.2.4}
Sea $Y$ un preesquema noetheriano, sea $f:X\to Y$ un morfismo de tipo finito y sea $\sh{L}$ un $\sh{O}_X$-módulo invertible amplio respecto de $f$.
Para todo $\sh{O}_X$-módulo $\sh{F}$, pongamos
\[
  \sh{F}(n)=\sh{F}\otimes_{\sh{O}_X}\sh{L}^{\otimes n}
  \qquad(n\in\bb{Z}).
\]
Sea
\[
  \sh{F}\longrightarrow\sh{G}\longrightarrow\sh{H}
\]
una sucesión exacta de $\sh{O}_X$-módulos coherentes tal que el soporte de $\operatorname{Im}(\sh{F}\to\sh{G})$ sea propio sobre $Y$ \sref[II]{II.5.4.10}.
Entonces existe un entero $N$ tal que, para $n\geq N$, la sucesión
\[
  f_*(\sh{F}(n))
  \longrightarrow f_*(\sh{G}(n))
  \longrightarrow f_*(\sh{H}(n))
\]
es exacta.
\end{corollary}

\begin{proof}
El razonamiento del comienzo de \sref{III.2.2.1} muestra que, si el corolario es cierto para $\sh{L}^{\otimes d}$, donde $d>0$, entonces es cierto para $\sh{L}$.
Podemos, por tanto, suponer que $\sh{L}$ es muy amplio respecto de $f$ \sref[II]{II.4.6.11}.
Podemos entonces identificar $X$ con un subconjunto abierto de un $Y$-preesquema
\[
  Z=\Proj(\sh{S}),
\]
donde $\sh{S}$ es una $\sh{O}_Y$-álgebra cuasicoherente graduada en grados positivos, $\sh{S}_1$ es de tipo finito y engendra $\sh{S}$, y
\[
  \sh{L}=i^*(\sh{O}_Z(1)),
\]
siendo $i:X\to Z$ la inmersión canónica \sref[II]{II.4.4.7}.

Pongamos
\[
  \sh{G}_1=\Im(\sh{F}\to\sh{G})
  =\Ker(\sh{G}\to\sh{H}),
\]
que es coherente. Como se tiene la sucesión exacta
\[
  0\to\sh{G}_1(n)\to\sh{G}(n)\to\sh{H}(n)
\]
y el funtor $f_*$ es exacto por la izquierda, la sucesión
\[
  0\to f_*(\sh{G}_1(n))\to f_*(\sh{G}(n))
  \to f_*(\sh{H}(n))
\]
es exacta para todo $n$; basta, pues, demostrar que, para $n$ suficientemente grande, la sucesión
\[
  f_*(\sh{F}(n))\to f_*(\sh{G}_1(n))\to0
\]
es exacta. En otras palabras, podemos reducirnos al caso en que $\sh{H}=0$ y el soporte de $\sh{G}$ es propio sobre $Y$, y por tanto \emph{cerrado en $Z$} (II, 5.4.10). En consecuencia, $\sh{G}'=i_*(\sh{G})$ es un $\sh{O}_Z$-módulo coherente tal que $\sh{G}'|X=\sh{G}$. Sabemos por (I, 9.4.3) que existe un $\sh{O}_Z$-módulo coherente $\sh{F}'$ tal que $\sh{F}'|X=\sh{F}$. Además, como $X$ es abierto en $Z$ y $\Supp(\sh{G})\subset X$ es cerrado en $Z$, se define inmediatamente un homomorfismo sobreyectivo $u':\sh{F}'\to\sh{G}'$ cuya restricción a $X$ es el homomorfismo sobreyectivo dado $u:\sh{F}\to\sh{G}$, tomando $u'|U=0$ para todo abierto $U$ de $Z$ que no interseque $\Supp(\sh{G})$, y $u'|U=u|U$ para todo abierto $U\subset X$.

Como al comienzo de la demostración de (2.2.2), podemos reducirnos al caso en que $Y$ es afín; se trata entonces de demostrar que, para $n$ suficientemente grande,
\[
  \Gamma(u):\Gamma(X,\sh{F}(n))
  \to\Gamma(X,\sh{G}(n))
\]
es sobreyectivo. Se tiene el diagrama conmutativo
\[
  \xymatrix@C=35mm@R=12mm{
    \Gamma(Z,\sh{F}'(n)) \ar[r]^{\Gamma(u')} \ar[d]_v &
    \Gamma(Z,\sh{G}'(n)) \ar[d]^w \\
    \Gamma(X,\sh{F}(n)) \ar[r]_{\Gamma(u)} &
    \Gamma(X,\sh{G}(n)).
  }
\]
La definición de $\sh{G}'$ muestra que $w$ es \emph{biyectivo}; por otra parte, en virtud de (2.2.3), $\Gamma(u')$ es \emph{sobreyectivo} para $n$ suficientemente grande, y lo mismo ocurre, por tanto, con $\Gamma(u)$.
\end{proof}

\begin{env}[2.2.5]
\label{III.2.2.5}
\emph{Observaciones.}

(i) La afirmación (i) de \sref{III.2.2.1} sigue siendo cierta si solo se supone que $Y$ es \emph{localmente noetheriano}.
En efecto, la propiedad es manifiestamente local sobre $Y$.
Además, las hipótesis de \sref{III.2.2.1} implican que, para todo subconjunto abierto $U\subset Y$, la restricción de $f$ a $f^{-1}(U)$ es un morfismo proyectivo
\[
  f^{-1}(U)\longrightarrow U
\]
\sref[II]{II.5.5.5}[iii], y $\sh{L}|_{f^{-1}(U)}$ es amplio para este morfismo \sref[II]{II.4.6.4}.

(ii) La afirmación (iii) de \sref{III.2.2.1} sigue siendo válida, como hemos visto, si solo se supone que $X$ es un esquema cuasicompacto, o un preesquema cuyo espacio subyacente es noetheriano, y que $f:X\to Y$ es cuasicompacto \sref[II]{II.4.6.8}.
Conviene observar, sin embargo, que incluso cuando $Y$ es el \emph{espectro de un cuerpo} $K$ y $f$ es \emph{cuasiproyectivo}, la afirmación (ii) de \sref{III.2.2.1} puede dejar de cumplirse.

Por ejemplo, sea
\[
  X'=\Spec(K[T_0,\ldots,T_r])
\]
y sea $X$ la unión de los subconjuntos abiertos afines $D(T_i)$ de $X'$, para $0\leq i\leq r$.
Como la inmersión $X\to X'$ es cuasicompacta, el morfismo estructural $f:X\to Y$ es cuasiafín \sref[II]{II.5.1.10} y, por tanto, $\sh{O}_X$ es muy amplio respecto de $f$ \sref[II]{II.5.1.6}.
Pero
\[
  \Gamma(X,\sh{O}_X)
  =\bigcap_{i=0}^r(K[T_0,\ldots,T_r])_{T_i}
  =K[T_0,\ldots,T_r]
\]
por \sref[I]{I.8.2.1.1}.
Se deduce de \sref{III.1.4.3.1} y \sref{III.1.1.3.5} que
\[
  \HH^r(X,\sh{O}_X^{\otimes n})
  =\HH^r(X,\sh{O}_X)
  \neq0
\]
para todo $n$.

(iii) Bajo las hipótesis de (2.2.4) relativas a $X$, $Y$, $f$ y $\sh{L}$, es inmediato que, si $\sh{H}$ es un $\sh{O}_X$-módulo coherente cuyo soporte es \emph{propio sobre $Y$}, un razonamiento análogo al de (2.2.4) muestra que existe un entero $N$ tal que, para $n\geq N$,
\[
  \RR^i f_*(\sh{H}(n))=0
\]
para todo $i>0$. Se deduce de la sucesión exacta de cohomología que, si $u:\sh{F}\to\sh{G}$ es un homomorfismo de $\sh{O}_X$-módulos coherentes tal que $\Ker(u)$ y $\Coker(u)$ tienen soportes \emph{propios sobre $Y$}, existe un entero $N$ tal que, para $n\geq N$, el homomorfismo correspondiente
\[
  \RR^i f_*(\sh{F}(n))\to\RR^i f_*(\sh{G}(n))
\]
es \emph{biyectivo} para todo $i>0$.
\end{env}


\subsection{Aplicación a los haces graduados de álgebras y de módulos}
\label{subsection:III.2.3}

\begin{theorem}[2.3.1]
\label{III.2.3.1}
Sea $Y$ un preesquema noetheriano, sea $\sh{S}$ una $\sh{O}_Y$-álgebra cuasicoherente de tipo finito, graduada en grados positivos, sea
\[
  X=\Proj(\sh{S}),
\]
y sea $q:X\to Y$ el morfismo estructural.
Sea $\sh{M}$ un $\sh{S}$-módulo graduado cuasicoherente que satisface la condición \textbf{(TF)}.
Existe entonces un entero $N$ tal que, para $n\geq N$, el homomorfismo canónico \sref[II]{II.8.14.5.1}
\[
  \alpha_n:\sh{M}_n
  \longrightarrow q_*\bigl(\shProj_0(\sh{M}(n))\bigr)
  =q_*\bigl((\shProj(\sh{M}))_n\bigr)
\]
\oldpage[III]{103}
es biyectivo.
En otras palabras, el homomorfismo canónico
\[
  \alpha:\sh{M}\longrightarrow\bbGamma_*(\shProj(\sh{M}))
\]
es un \textbf{(TN)}-isomorfismo.
\end{theorem}

\begin{proof}
Podemos restringirnos al caso en que $\sh{M}$ es un $\sh{S}$-módulo de tipo finito \sref[II]{II.3.4.2}.
Como $Y$ es cuasicompacto, existe un entero $d>0$ tal que $\sh{S}^{(d)}$ está engendrada por el $\sh{O}_Y$-módulo cuasicoherente $\sh{S}_d$, que es de tipo finito \sref[II]{II.3.1.10} y, por tanto, coherente porque $Y$ es noetheriano.
Ahora bien, $\sh{M}$ es la suma directa de los $\sh{M}^{(d,k)}$, para $0\leq k<d$, y cada $\sh{M}^{(d,k)}$ es un $\sh{S}^{(d)}$-módulo cuasicoherente de tipo finito, por \sref[II]{II.2.1.6}[iii].
Como la cuestión es local sobre $Y$, basta manifiestamente demostrar que cada homomorfismo canónico
\[
  \alpha:\sh{M}^{(d,k)}
  \longrightarrow
  \bbGamma_*\bigl((\shProj(\sh{M}))^{(d,k)}\bigr)
\]
es un \textbf{(TN)}-isomorfismo.
Teniendo en cuenta \sref[II]{II.8.14.13}, y en particular el diagrama \sref[II]{II.8.14.13.4}, quedamos así reducidos a demostrar el teorema cuando $\sh{S}$ está engendrada por $\sh{S}_1$ y $\sh{S}_1$ es un $\sh{O}_Y$-módulo coherente.

Como $Y$ es noetheriano, el razonamiento utilizado al comienzo de la demostración de \sref{III.2.2.2} muestra que podemos reducirnos al caso
\[
  Y=\Spec(A),\qquad
  \sh{S}=\widetilde{S},\qquad
  \sh{M}=\widetilde{M},
\]
donde $A$ es un anillo noetheriano, $S_1$ es un $A$-módulo de tipo finito y $M$ es un $S$-módulo graduado de tipo finito.
Mostraremos que basta entonces demostrar el teorema para $M=S$.
En efecto, en el caso general existe una sucesión exacta
\[
  L'\longrightarrow L\longrightarrow M\longrightarrow0,
\]
donde $L$ y $L'$ son sumas directas de módulos graduados de la forma $S(m)$.
Si el resultado se cumple para $M=S$, se cumple también para $M=S(m)$ y, por tanto, para $L$ y $L'$.
Consideremos el diagrama conmutativo
\[
  \xymatrix{
    \widetilde{L'_n}\ar[r]\ar[d]_{\alpha_n} &
    \widetilde{L_n}\ar[r]\ar[d]^{\alpha_n} &
    \widetilde{M_n}\ar[r]\ar[d]^{\alpha_n} &
    0\\
    q_*(\widetilde{L'}(n))\ar[r] &
    q_*(\widetilde{L}(n))\ar[r] &
    q_*(\widetilde{M}(n))\ar[r] &
    0 .
  }
\]
\oldpage[III]{104}
La segunda fila es exacta por \sref{III.2.2.3} cuando $n$ es suficientemente grande.
Lo mismo ocurre con la primera fila y, como las dos flechas verticales de la izquierda son isomorfismos, también lo es la tercera.

Queda, pues, demostrar el teorema para $M=S$.
Supongamos primero que
\[
  S=A[T_0,\ldots,T_r],
\]
donde los $T_i$ son indeterminadas.
En este caso, la afirmación es precisamente \sref{III.2.1.11}[ii].
En el caso general, $S$ se identifica con un cociente de un anillo
\[
  S'=A[T_0,\ldots,T_r]
\]
por un ideal graduado, de modo que $X$ es un subesquema cerrado de
\[
  X'=\bb{P}_A^r
\]
\sref[II]{II.2.9.2}.
Si $j:X\to X'$ es la inyección canónica, entonces $j_*(\widetilde{S}(n))$ es el $\sh{O}_{X'}$-módulo
\[
  (\shProj(\widetilde{S}))(n),
\]
donde $S$ se considera como $S'$-módulo graduado; esto se deduce inmediatamente de \sref[II]{II.2.8.7}.
Como $j_*(\widetilde{S}(n))$ es un $\sh{O}_{X'}$-módulo que satisface \textbf{(TF)}, el homomorfismo canónico
\[
  \alpha_n:S_n\longrightarrow
  \Gamma\bigl(X',j_*(\widetilde{S}(n))\bigr)
\]
es biyectivo para todo $n$ suficientemente grande, por el caso polinómico que acabamos de demostrar.
Esto termina la demostración, puesto que
\[
  \Gamma\bigl(X',j_*(\widetilde{S}(n))\bigr)
  =\Gamma(X,\widetilde{S}(n)).
\]
\end{proof}

\begin{corollary}[2.3.2]
\label{III.2.3.2}
Bajo las hipótesis de \sref{III.2.3.1}, sea
\[
  \sh{S}_X=\bigoplus_{n\in\bb{Z}}\sh{O}_X(n),
\]
y sea $\sh{F}$ un $\sh{S}_X$-módulo graduado cuasicoherente de tipo finito.
Entonces $\bbGamma_*(\sh{F})$ satisface la condición \textbf{(TF)}.
\end{corollary}

\begin{proof}
En la demostración de \sref{III.2.3.1} vimos que $X$, que es isomorfo a $\Proj(\sh{S}^{(d)})$ \sref[II]{II.3.1.8}, es de tipo finito sobre $Y$ \sref[II]{II.3.4.1}.
Se deduce de \sref[II]{II.8.14.9} que $\sh{F}$ es isomorfo a un $\sh{S}_X$-módulo graduado de la forma $\shProj(\sh{M})$, donde $\sh{M}$ es un $\sh{S}$-módulo graduado cuasicoherente de tipo finito.
Por \sref{III.2.3.1}, $\bbGamma_*(\sh{F})$ es \textbf{(TN)}-isomorfo a $\sh{M}$ y, por consiguiente, satisface \textbf{(TF)}.
\end{proof}

\begin{env}[2.3.3]
\label{III.2.3.3}
\emph{Escolio.}
Sea $Y$ un preesquema noetheriano, sea $\sh{S}$ una $\sh{O}_Y$-álgebra graduada que satisface las hipótesis de \sref{III.2.3.1}, y sea $X=\Proj(\sh{S})$.
Sea $\cat{K}_{\sh{S}}$ la categoría abeliana de los $\sh{S}$-módulos graduados cuasicoherentes que satisfacen \textbf{(TF)}, y sea $\cat{K}'_{\sh{S}}$ la subcategoría de $\cat{K}_{\sh{S}}$ formada por los $\sh{S}$-módulos que satisfacen \textbf{(TN)}.
Por último, sea $\cat{K}_X$ la categoría de los $\sh{S}_X$-módulos graduados cuasicoherentes $\sh{F}$ de tipo finito; como $\sh{S}_X$ es periódica \sref[II]{II.8.14.4} y \sref[II]{II.8.14.12}, esto equivale a exigir que los $\sh{F}_i$ sean $\sh{O}_X$-módulos coherentes.

Por \sref[II]{II.8.14.8}, \sref[II]{II.8.14.10} y \sref{III.2.3.2}, los funtores
\[
  \sh{M}\longmapsto\shProj(\sh{M}),
  \qquad
  \sh{F}\longmapsto\bbGamma_*(\sh{F})
\]
definen una equivalencia (T, I, 1.2) entre la categoría cociente
\[
  \cat{K}_{\sh{S}}/\cat{K}'_{\sh{S}}
\]
(T, I, 1.11) y la categoría $\cat{K}_X$.
Cuando $\sh{S}$ está engendrada por $\sh{S}_1$, la categoría $\cat{K}_X$ puede sustituirse por la categoría de los $\sh{O}_X$-módulos coherentes \sref[II]{II.8.14.12}.
\end{env}

\begin{proposition}[2.3.4]
\label{III.2.3.4}
Sea $Y$ un preesquema noetheriano.

(i) Sea $\sh{S}$ una $\sh{O}_Y$-álgebra cuasicoherente de tipo finito, graduada en grados positivos.
Pongamos
\[
  X=\Proj(\sh{S}),\qquad
  \sh{S}_X=\shProj(\sh{S})
  =\bigoplus_{n\in\bb{Z}}\sh{O}_X(n).
\]
Entonces $\sh{S}_X$ es una $\sh{O}_X$-álgebra graduada cuasicoherente periódica \sref[II]{II.8.14.12}, sus componentes homogéneas
\[
  (\sh{S}_X)_n=\sh{O}_X(n)
\]
son $\sh{O}_X$-módulos coherentes y, si $d>0$ es un período de $\sh{S}_X$, entonces
\[
  (\sh{S}_X)_d=\sh{O}_X(d)
\]
es $Y$-amplio como $\sh{O}_X$-módulo invertible.
Además, el homomorfismo canónico
\[
  \alpha:\sh{S}\longrightarrow\bbGamma_*(\sh{S}_X)
\]
es un \textbf{(TN)}-isomorfismo.

(ii) Recíprocamente, sea $q:X\to Y$ un morfismo proyectivo y sea $\sh{S}'$ una $\sh{O}_X$-álgebra graduada cuyas componentes homogéneas $\sh{S}'_n$, para $n\in\bb{Z}$, son $\sh{O}_X$-módulos coherentes.
Supongamos que $\sh{S}'$ tiene un período $d>0$ tal que $\sh{S}'_d$ es un $\sh{O}_X$-módulo invertible amplio para $q$.
Entonces
\[
  \sh{S}=\bigoplus_{n\geq0}q_*(\sh{S}'_n)
\]
es una $\sh{O}_Y$-álgebra cuasicoherente de tipo finito, graduada en grados positivos, y existe un $Y$-isomorfismo
\[
  r:X\simto\Proj(\sh{S})
\]
tal que $r^*(\shProj(\sh{S}))$ es isomorfo a $\sh{S}'$ como $\sh{O}_X$-álgebra graduada.
\end{proposition}

\begin{proof}
(i) Casi todas las afirmaciones ya han sido demostradas; la última es un caso particular de \sref{III.2.3.2}.
La periodicidad de $\sh{S}_X$ fue demostrada en \sref[II]{II.8.14.14}, y la afirmación de que existe un período $d>0$ para el cual $\sh{O}_X(d)$ es invertible y $Y$-amplio es precisamente \sref[II]{II.4.6.18}.
Finalmente, para $0\leq k<d$, $(\sh{S}_X)^{(d,k)}$ es un $(\sh{S}_X)^{(d)}$-módulo de tipo finito \sref[II]{II.8.14.14}.
Se deduce que cada $(\sh{S}_X)_n$ es un $\sh{O}_X$-módulo cuasicoherente de tipo finito \sref[II]{II.2.1.6}[ii]; como $\sh{O}_X$ es coherente, también lo es cada $(\sh{S}_X)_n$.

(ii) Tras sustituir el período $d$ por uno de sus múltiplos, podemos suponer que
\[
  \sh{L}=\sh{S}'_d
\]
es muy amplio respecto de $q$ \sref[II]{II.4.6.11}.
Por hipótesis,
\[
  \sh{S}'^{(d)}
  =\bigoplus_{n\in\bb{Z}}\sh{L}^{\otimes n},
\]
y, por tanto,
\[
  \sh{S}^{(d)}
  =\bigoplus_{n\geq0}q_*(\sh{L}^{\otimes n}).
\]
Por \sref[II]{II.3.1.8} y \sref[II]{II.3.2.9}, existe un $Y$-isomorfismo
\[
  s:X'=\Proj(\sh{S})\simto
  X''=\Proj(\sh{S}^{(d)})
\]
tal que
\oldpage[III]{105}
\[
  s^*(\sh{O}_{X''}(n))=\sh{O}_{X'}(nd)
\]
para todo $n$ \sref[II]{II.3.5.2}.
Estableceremos, por tanto, la existencia de un $Y$-isomorfismo $X\simto X'$ demostrando la afirmación siguiente.

\begin{env}[2.3.4.1]
\label{III.2.3.4.1}
Sea $Y$ un preesquema noetheriano, sea $q:X\to Y$ un morfismo proyectivo y sea $\sh{L}$ un $\sh{O}_X$-módulo invertible muy amplio para $q$.
Entonces
\[
  \sh{S}=\bigoplus_{n\geq0}q_*(\sh{L}^{\otimes n})
\]
es una $\sh{O}_Y$-álgebra graduada cuasicoherente de tipo finito tal que
\[
  \sh{S}_n=\sh{S}_1^n
\]
para todo $n$ suficientemente grande, y existe un $Y$-isomorfismo
\[
  r:X\simto P=\Proj(\sh{S})
\]
tal que
\[
  \sh{L}=r^*(\sh{O}_P(1)).
\]
\end{env}

Como $q$ es proyectivo, se deduce de \sref[II]{II.5.4.4} y \sref[II]{II.4.4.7} que existe un $Y$-isomorfismo
\[
  r':X\simto P'=\Proj(\sh{T}),
\]
donde $\sh{T}$ es una $\sh{O}_Y$-álgebra cuasicoherente, $\sh{T}_1$ es de tipo finito y engendra $\sh{T}$, y
\[
  \sh{L}=r'^*(\sh{O}_{P'}(1)).
\]
Si $q':P'\to Y$ es el morfismo estructural, entonces
\[
  \sh{S}=\bigoplus_{n\geq0}q'_*(\sh{O}_{P'}(n)).
\]
Por \sref{III.2.3.1}, el homomorfismo canónico
\[
  \alpha_n:\sh{T}_n\longrightarrow
  \sh{S}_n=q'_*(\sh{O}_{P'}(n))
\]
es biyectivo para todo $n$ suficientemente grande.
Como $\sh{T}_n=\sh{T}_1^n$, se deduce \emph{a fortiori} que $\sh{S}_n=\sh{S}_1^n$ para todo $n$ suficientemente grande.

El homomorfismo canónico $\alpha:\sh{T}\to\sh{S}$ de $\sh{O}_Y$-álgebras graduadas es un \textbf{(TN)}-isomorfismo.
Por consiguiente,
\phantomsection\label{III.2.3.5.1}
\[
  \Phi=\Proj(\alpha):\Proj(\sh{S})\longrightarrow\Proj(\sh{T})
\]
es un isomorfismo \sref[II]{II.3.6.1}.
Además, \sref[II]{II.3.5.2} y el hecho de que los $\sh{T}$-módulos graduados obtenidos de $\sh{S}(n)$ por restricción de escalares son \textbf{(TN)}-isomorfos a $\sh{T}(n)$ dan
\[
  \Phi_*(\sh{O}_P(n))=\sh{O}_{P'}(n)
\]
para todo $n$ \sref[II]{II.3.4.2}.
Solo queda mostrar que $\sh{S}$ es de tipo finito sobre $\sh{O}_Y$.
Los módulos
\[
  \sh{S}_n=q'_*(\sh{O}_{P'}(n))
\]
son coherentes por \sref{III.2.2.1}; y si $\sh{S}_n=\sh{S}_1^n$ para $n\geq n_0$, entonces $\sh{S}$ está engendrada por el módulo coherente
\[
  \bigoplus_{i\leq n_0}\sh{S}_i.
\]
Esto demuestra la afirmación \sref[I]{I.9.6.2}.

Volvamos a la demostración de \sref{III.2.3.4}, conservando su notación.
Hemos demostrado que existe un $Y$-isomorfismo
\[
  r'':X\simto X''
\]
tal que
\[
  r''_*(\sh{L}^{\otimes n})=\sh{O}_{X''}(n)
\]
para todo $n\in\bb{Z}$.
Sea $q'':X''\to Y$ el morfismo estructural.
El álgebra $\sh{S}'$ es la suma directa de los $\sh{S}'^{(d)}$-módulos graduados $\sh{S}'^{(d,k)}$.
Cada uno de ellos es un $\sh{S}'^{(d)}$-módulo cuasicoherente de tipo finito, por la periodicidad de $\sh{S}'$ y la hipótesis de que los $\sh{S}'_n$ son $\sh{O}_X$-módulos de tipo finito \sref[II]{II.8.14.12}.
Pongamos
\[
  \sh{F}^{(k)}=r''_*(\sh{S}'^{(d,k)}),
\]
de modo que los $\sh{F}^{(k)}$ son $\sh{S}_{X''}$-módulos graduados cuasicoherentes de tipo finito.
Por \sref[II]{II.8.14.8}, el homomorfismo canónico
\[
  \beta:\shProj\bigl(\bbGamma_*(\sh{F}^{(k)})\bigr)
  \longrightarrow\sh{F}^{(k)}
\]
es un isomorfismo de $\sh{S}_{X''}$-módulos.
Por otra parte,
\[
  q''_*\bigl((\sh{F}^{(k)})_n\bigr)
  =q_*\bigl((\sh{S}'^{(d,k)})_n\bigr),
\]
y para $n\geq0$ el último $\sh{O}_Y$-módulo es, por definición, igual a $(\sh{S}^{(d,k)})_n$.
En otras palabras, la inyección canónica
\[
  \sh{S}^{(d,k)}
  \longrightarrow\bbGamma_*(\sh{F}^{(k)})
\]
es un \textbf{(TN)}-isomorfismo.
Por tanto,
\[
  \shProj(\sh{S}^{(d,k)})
  =\shProj\bigl(\bbGamma_*(\sh{F}^{(k)})\bigr),
\]
y, en consecuencia,
\[
  r''{}^*\bigl(\shProj(\sh{S}^{(d,k)})\bigr)
  \simeq\sh{S}'^{(d,k)}.
\]
Queda observar que, salvo isomorfismo canónico,
\[
  \shProj(\sh{S}^{(d,k)})
  =s_*\bigl((\shProj(\sh{S}))^{(d,k)}\bigr)
\]
\sref[II]{II.8.14.13.1}; esto demuestra que $r^*(\shProj(\sh{S}))$ es isomorfo a $\sh{S}'$.

Finalmente, por \sref{III.2.3.2}, cada $\bbGamma_*(\sh{F}^{(k)})$ satisface la condición \textbf{(TF)} y, por tanto, también la satisface cada $\sh{S}^{(d,k)}$.
Además, los $\sh{S}_n=q_*(\sh{S}'_n)$ son coherentes por \sref{III.2.2.1}, porque los $\sh{S}'_n$ son coherentes.
Se deduce inmediatamente que los $\sh{S}^{(d,k)}$ son $\sh{S}^{(d)}$-módulos de tipo finito.
Como \sref{III.2.3.4.1} muestra que $\sh{S}^{(d)}$ es una $\sh{O}_Y$-álgebra de tipo finito, $\sh{S}$ es asimismo una $\sh{O}_Y$-álgebra de tipo finito.
\end{proof}

\oldpage[III]{106}
\begin{proposition}[2.3.5]
\label{III.2.3.5}
Sea $Y$ un preesquema noetheriano integral, sea $X$ un preesquema integral y sea
\[
  f:X\longrightarrow Y
\]
un morfismo proyectivo birracional.
Existe entonces un ideal fraccionario coherente
\[
  \sh{J}\subset\sh{R}(Y)
\]
\sref[II]{II.8.1.2} tal que $X$ es $Y$-isomorfo al preesquema obtenido por explosión de $\sh{J}$ \sref[II]{II.8.1.3}.
Además, existe un subconjunto abierto $U$ de $Y$ tal que la restricción de $f$ a $f^{-1}(U)$ es un isomorfismo de $f^{-1}(U)$ sobre $U$ (cf.\ \textbf{I}, 6.5.5), y tal que $\sh{J}|U$ es invertible.
\end{proposition}

\begin{proof}
Existe un $\sh{O}_X$-módulo invertible $\sh{L}$ muy amplio para $f$ \sref[II]{II.4.4.2} y \sref[II]{II.5.3.2}.
Aplicando \sref{III.2.3.4.1}, identificamos $X$ con $\Proj(\sh{S})$, donde
\[
  \sh{S}=\bigoplus_{n\geq0}f_*(\sh{L}^{\otimes n}).
\]
Los módulos $f_*(\sh{L}^{\otimes n})$ son sin torsión sobre $\sh{O}_Y$ \sref[I]{I.7.4.5}.
Por tanto, $\sh{S}$ se identifica canónicamente con un sub-$\sh{O}_Y$-módulo de
\[
  \sh{S}\otimes_{\sh{O}_Y}\sh{R}(Y)
\]
\sref[I]{I.7.4.1}.
Este último es un haz simple \sref[I]{I.7.3.6} y está, por tanto, determinado por su restricción a un subconjunto abierto no vacío.

Elijamos un subconjunto abierto no vacío $U'\subset U$ sobre el cual $\sh{L}|f^{-1}(U')$ sea isomorfo a $\sh{O}_X|f^{-1}(U')$ y, por tanto, a $\sh{O}_Y|U'$.
Las restricciones $f_*(\sh{L}^{\otimes n})|U'$ son entonces isomorfas a $\sh{O}_Y|U'$.
Se sigue que
\[
  \sh{S}\otimes_{\sh{O}_Y}\sh{R}(Y)
  \simeq\sh{R}(Y)[T],
\]
donde $T$ es una indeterminada.
Por \sref{III.2.3.4.1}, $\sh{S}$ es \textbf{(TN)}-isomorfa a la sub-$\sh{O}_Y$-álgebra engendrada por la imagen canónica de $f_*(\sh{L})$ en esta álgebra de polinomios.
Bajo la identificación mostrada, dicha imagen es $\sh{J}\cdot T$, donde $\sh{J}$ es un sub-$\sh{O}_Y$-módulo coherente de $\sh{R}(Y)$ \sref{III.2.2.1}.
Su restricción a $U'$ es isomorfa a $\sh{O}_Y|U'$ y, reduciendo $U$ si es necesario, $\sh{J}|U$ es invertible.
Por consiguiente,
\[
  \sh{S}\quad\text{es \textbf{(TN)}-isomorfa a}\quad
  \bigoplus_{n\geq0}\sh{J}^n,
\]
lo que termina la demostración.
\end{proof}

\begin{corollary}[2.3.6]
\label{III.2.3.6}
Bajo las hipótesis de \sref{III.2.3.5}, supongamos además que para todo sub-$\sh{O}_Y$-módulo coherente no nulo
\[
  \sh{J}\subset\sh{R}(Y)
\]
existe un $\sh{O}_Y$-módulo invertible $\sh{L}$ tal que
\[
  \Gamma\bigl(Y,\sh{L}\otimes_{\sh{O}_Y}
  \shHom(\sh{J},\sh{O}_Y)\bigr)\neq0.
\]
Entonces, en el enunciado de \sref{III.2.3.5}, puede tomarse $\sh{J}$ como un ideal de $\sh{O}_Y$.
Esta condición adicional se satisface siempre si existe un $\sh{O}_Y$-módulo amplio.
\end{corollary}

\begin{proof}
En efecto, por \sref[0]{0.5.4.2},
\[
  \sh{L}\otimes\shHom(\sh{J},\sh{O}_Y)
  =\shHom\bigl(\sh{L}^{-1},\shHom(\sh{J},\sh{O}_Y)\bigr)
  =\shHom(\sh{J}\otimes\sh{L}^{-1},\sh{O}_Y).
\]
La hipótesis dice, por tanto, que existe un homomorfismo no nulo
\[
  u:\sh{J}\otimes\sh{L}^{-1}\longrightarrow\sh{O}_Y.
\]
Para todo $y\in Y$, el módulo $(\sh{J}\otimes\sh{L}^{-1})_y$ se identifica con un sub-$\sh{O}_y$-módulo del cuerpo de fracciones $\sh{R}(Y)_y$ de $\sh{O}_y$ \sref[I]{I.7.1.5}.
Así, $u_y$ es necesariamente inyectivo, y $u$ identifica $\sh{J}\otimes\sh{L}^{-1}$ con un ideal $\sh{J}'$ de $\sh{O}_Y$.
Pero
\[
  \Proj\left(\bigoplus_{n\geq0}\sh{J}^n\right)
  \quad\text{y}\quad
  \Proj\left(\bigoplus_{n\geq0}
  (\sh{J}\otimes\sh{L}^{-1})^n\right)
\]
son $Y$-isomorfos \sref[II]{II.3.1.8}; esto demuestra la primera aserción.

Para la segunda, observemos que
\[
  \sh{F}=\shHom_{\sh{O}_Y}(\sh{J},\sh{O}_Y)
\]
es coherente y no nulo, pues existe un subconjunto abierto $U$ de $Y$ sobre el cual $\sh{J}|U$ es invertible.
Si $\sh{L}$ es un $\sh{O}_Y$-módulo amplio, existe un entero $n$ tal que
\[
  \sh{F}(n)=\sh{F}\otimes\sh{L}^{\otimes n}
\]
está engendrado por sus secciones sobre $Y$ \sref[II]{II.4.5.5}.
\emph{A fortiori},
\[
  \Gamma(Y,\sh{F}(n))\neq0,
\]
lo que demuestra la aserción.
\end{proof}

\begin{corollary}[2.3.7]
\label{III.2.3.7}
Sean $X$ e $Y$ esquemas integrales proyectivos sobre un cuerpo $k$, y sea
\[
  f:X\longrightarrow Y
\]
un $k$-morfismo birracional.
Entonces $X$ es $k$-isomorfo a un $Y$-esquema obtenido por explosión de un subesquema cerrado $Y'$ de $Y$, no necesariamente reducido.
\end{corollary}

\begin{proof}
\oldpage[III]{107}
En efecto, $f$ es proyectivo \sref[II]{II.5.5.5}[v], y como $Y$ es proyectivo sobre $k$, se satisface la condición adicional de \sref{III.2.3.6}.
Basta, por tanto, tomar el subesquema cerrado $Y'$ de $Y$ definido por el ideal coherente $\sh{J}$ de \sref{III.2.3.6}.
\end{proof}

\begin{env}[2.3.8]
\label{III.2.3.8}
\emph{Observación.}
En el Capítulo~IV, al estudiar la noción de divisor, veremos que si, en el enunciado de \sref{III.2.3.5}, se supone que los anillos locales $\sh{O}_y$, para $y\in Y$, son factoriales---como ocurre, por ejemplo, cuando $Y$ es no singular---, entonces $X$ puede obtenerse de $Y$ por explosión de un subpreesquema cerrado $Y'$ de $Y$ cuyo espacio subyacente está contenido en $Y-U$.
\end{env}


\subsection{Una generalización del teorema fundamental}
\label{subsection:III.2.4}

\begin{theorem}[2.4.1]
\label{III.2.4.1}
Sea $Y$ un preesquema noetheriano, sea $\sh{S}$ una $\sh{O}_Y$-álgebra cuasicoherente de tipo finito, sea
\[
  f:X\longrightarrow Y
\]
un morfismo proyectivo, pongamos $\sh{S}'=f^*(\sh{S})$, y sea $\sh{M}$ un $\sh{S}'$-módulo cuasicoherente de tipo finito.
Entonces:
\begin{enumerate}
  \item[(i)] Para todo $p\in\bb Z$, $\RR^p f_*(\sh{M})$ es un $\sh{S}$-módulo de tipo finito.

  \item[(ii)] Supongamos además que $\sh{L}$ es un $\sh{O}_X$-módulo invertible amplio para $f$ y, para todo $n\in\bb Z$, pongamos
  \[
    \sh{M}(n)=\sh{M}\otimes\sh{L}^{\otimes n}.
  \]
  Existe un entero $N$ tal que, para $n\geq N$,
  \[
  \label{III.2.4.1.1}
    \RR^p f_*(\sh{M}(n))=0
    \tag{2.4.1.1}
  \]
  para todo $p>0$, y el homomorfismo canónico
  \[
    f^*\bigl(f_*(\sh{M}(n))\bigr)\longrightarrow\sh{M}(n)
  \]
  \sref[0]{0.4.4.3} es sobreyectivo.
\end{enumerate}
\end{theorem}

\begin{proof}
Pongamos
\[
  Y'=\Spec(\sh{S}),\qquad X'=\Spec(\sh{S}'),
\]
de modo que $X'=X\times_Y Y'$ \sref[II]{II.1.5.5}.
Sean
\[
  g:Y'\longrightarrow Y,\qquad g':X'\longrightarrow X
\]
los morfismos estructurales, que son afines por definición, y sea
\[
  f'=f_{(Y')}:X'\longrightarrow Y'.
\]
Tenemos así un diagrama conmutativo
\[
  \xymatrix{
    X\ar[d]_f &
    X'\ar[l]_{g'}\ar[d]^{f'}\ar[dl]^{h}\\
    Y &
    Y'\ar[l]^{g},
  }
\]
donde $f'$ es proyectivo \sref[II]{II.5.5.5}[iii] y
\[
  h=f\circ g'=g\circ f'.
\]

\emph{(i)} Sea $\widetilde{\sh{M}}$ el $\sh{O}_{X'}$-módulo asociado al $\sh{S}'$-módulo cuasicoherente $\sh{M}$ cuando se considera $X'$ como un $X$-esquema afín \sref[II]{II.1.4.3}; así,
\[
  \sh{M}=g'_*(\widetilde{\sh{M}}).
\]
Como $\sh{M}$ es un $\sh{S}'$-módulo de tipo finito, $\widetilde{\sh{M}}$ es un $\sh{O}_{X'}$-módulo de tipo finito \sref[II]{II.1.4.5}.
Además, $h$ es de tipo finito, puesto que $g$ y $f'$ son de tipo finito \sref[II]{II.1.3.7} y \sref[I]{I.6.3.4}[ii].
Por tanto, $X'$ es noetheriano \sref[I]{I.6.3.7}, y $\widetilde{\sh{M}}$ es, en consecuencia, coherente.

Como $g'$ es afín, el homomorfismo canónico
\[
  \RR^p f_*(\sh{M})
  \longrightarrow
  \RR^p h_*(\widetilde{\sh{M}})
\]
es biyectivo \sref{III.1.3.4}.
Es, además, un homomorfismo de $\sh{S}$-módulos.
En efecto, el homomorfismo canónico
\[
\label{III.2.4.1.2}
  g'_*(\sh{O}_{X'})
  \otimes_{\sh{O}_X}
  g'_*(\widetilde{\sh{M}})
  \longrightarrow
  g'_*(\widetilde{\sh{M}})
  \tag{2.4.1.2}
\]
define la estructura de $\sh{S}'$-módulo sobre $\sh{M}$, puesto que
\[
  \sh{S}'=g'_*(\sh{O}_{X'}).
\]
Da lugar canónicamente a un homomorfismo
\[
  f_*\bigl(g'_*(\sh{O}_{X'})\bigr)
  \otimes
  \RR^p f_*\bigl(g'_*(\widetilde{\sh{M}})\bigr)
  \longrightarrow
  \RR^p f_*\bigl(g'_*(\widetilde{\sh{M}})\bigr).
\]
\oldpage[III]{108}
Por \sref[0]{0.12.2.2}, y dado que \eqref{III.2.4.1.2} se obtiene a su vez, aplicando \sref[0]{0.4.2.2.1}, del homomorfismo
\[
  \sh{O}_{X'}\otimes\widetilde{\sh{M}}
  \longrightarrow\widetilde{\sh{M}}
\]
que define la estructura de $\sh{O}_{X'}$-módulo sobre $\widetilde{\sh{M}}$, el diagrama
\[
  \xymatrix{
    f_*\bigl(g'_*(\sh{O}_{X'})\bigr)
    \otimes\RR^p f_*\bigl(g'_*(\widetilde{\sh{M}})\bigr)
    \ar[r]\ar[d] &
    \RR^p f_*\bigl(g'_*(\widetilde{\sh{M}})\bigr)\ar[d]\\
    h_*(\sh{O}_{X'})
    \otimes\RR^p h_*(\widetilde{\sh{M}})
    \ar[r] &
    \RR^p h_*(\widetilde{\sh{M}})
  }
\]
es conmutativo \sref[0]{0.12.2.6}.
Componiendo las flechas horizontales con el homomorfismo inducido por el homomorfismo canónico
\[
  \sh{S}
  \longrightarrow f_*(f^*(\sh{S}))
  =f_*(\sh{S}')
  =f_*\bigl(g'_*(\sh{O}_{X'})\bigr)
  =h_*(\sh{O}_{X'}),
\]
obtenemos la aserción.

Por otra parte, como $g$ es afín y $f'$ es separado y cuasicompacto, el homomorfismo canónico
\[
  \RR^p h_*(\widetilde{\sh{M}})
  \longrightarrow
  g_*\bigl(\RR^p f'_*(\widetilde{\sh{M}})\bigr)
\]
es biyectivo \sref{III.1.4.14}.
Como antes, es un isomorfismo de $\sh{S}$-módulos, esta vez utilizando la conmutatividad de \sref[0]{0.12.2.6.2}.
Ahora bien, $f'$ es proyectivo y $\widetilde{\sh{M}}$ es coherente, por lo que $\RR^p f'_*(\widetilde{\sh{M}})$ es un $\sh{O}_{Y'}$-módulo coherente por \sref{III.2.2.1}.
Se sigue que
\[
  g_*\bigl(\RR^p f'_*(\widetilde{\sh{M}})\bigr)
\]
es un $\sh{S}$-módulo de tipo finito \sref[II]{II.1.4.5}.

\emph{(ii)} Pongamos
\[
  \sh{L}'=g'^*(\sh{L}),
\]
un $\sh{O}_{X'}$-módulo invertible.
Para todo $n\in\bb Z$, salvo isomorfismo canónico,
\[
  g'_*\bigl(\widetilde{\sh{M}}\otimes\sh{L}'^{\otimes n}\bigr)
  =
  g'_*(\widetilde{\sh{M}})\otimes\sh{L}^{\otimes n}
  =
  \sh{M}(n)
\]
\sref[0]{0.5.4.10}.
Aplicando el razonamiento de \emph{(i)} a
\[
  \widetilde{\sh{M}}\otimes\sh{L}'^{\otimes n}
\]
muestra que
\[
  \RR^p f_*
  \bigl(g'_*(\widetilde{\sh{M}}\otimes\sh{L}'^{\otimes n})\bigr)
  \simeq
  g_*\bigl(\RR^p f'_*
  (\widetilde{\sh{M}}\otimes\sh{L}'^{\otimes n})\bigr).
\]
Ahora bien, $\sh{L}'$ es amplio para $f'$ \sref[II]{II.4.6.13}[iii].
Se deduce de \sref{III.2.2.1} que existe un entero $N$ tal que
\[
  \RR^p f'_*
  \bigl(\widetilde{\sh{M}}\otimes\sh{L}'^{\otimes n}\bigr)=0
\]
para todo $p>0$ y todo $n\geq N$.
Esto demuestra \eqref{III.2.4.1.1}.

Se deduce una vez más de \sref{III.2.2.1} que puede elegirse $N$ de modo que, para $n\geq N$, el homomorfismo canónico
\[
  f'^*\bigl(f'_*(\widetilde{\sh{M}}\otimes\sh{L}'^{\otimes n})\bigr)
  \longrightarrow
  \widetilde{\sh{M}}\otimes\sh{L}'^{\otimes n}
\]
sea sobreyectivo.
Como $g'_*$ es exacto \sref[II]{II.1.4.4}, el homomorfismo correspondiente
\[
  g'_*\Bigl(
    f'^*\bigl(f'_*(\widetilde{\sh{M}}\otimes\sh{L}'^{\otimes n})\bigr)
  \Bigr)
  \longrightarrow
  g'_*(\widetilde{\sh{M}}\otimes\sh{L}'^{\otimes n})
  =\sh{M}(n)
\]
es sobreyectivo.
Pero \sref[II]{II.1.5.2} da
\[
  g'_*\Bigl(
    f'^*\bigl(f'_*(\widetilde{\sh{M}}\otimes\sh{L}'^{\otimes n})\bigr)
  \Bigr)
  =
  f^*\Bigl(
    g_*\bigl(f'_*(\widetilde{\sh{M}}\otimes\sh{L}'^{\otimes n})\bigr)
  \Bigr),
\]
y $g_*\circ f'_*=f_*\circ g'_*$.
Por consiguiente,
\[
  g'_*\Bigl(
    f'^*\bigl(f'_*(\widetilde{\sh{M}}\otimes\sh{L}'^{\otimes n})\bigr)
  \Bigr)
  =
  f^*\bigl(f_*(\sh{M}(n))\bigr),
\]
lo que termina la demostración.
\end{proof}

\begin{env}[2.4.2]
\label{III.2.4.2}
Aplicaremos \sref{III.2.4.1}, en particular, cuando $\sh{S}$ sea una $\sh{O}_Y$-álgebra graduada en grados positivos y
\[
  \sh{M}=\bigoplus_{k\in\bb Z}\sh{M}_k
\]
sea un $\sh{S}'$-módulo graduado.
\oldpage[III]{109}
Bajo las mismas hipótesis de finitud sobre $\sh{S}$ y $\sh{M}$, se deduce de \sref{III.2.4.1} que
\[
  \bigoplus_{k\in\bb Z}\RR^p f_*(\sh{M}_k)
\]
es un $\sh{S}$-módulo de tipo finito para todo $p$.
Bajo las hipótesis adicionales de \sref{III.2.4.1}[ii], existe además un entero $N$ tal que, para $n\geq N$,
\[
  \RR^p f_*(\sh{M}_k(n))=0
\]
para todo $p>0$ y todo $k\in\bb Z$, y el homomorfismo canónico
\[
  f^*\bigl(f_*(\sh{M}_k(n))\bigr)
  \longrightarrow
  \sh{M}_k(n)
\]
es sobreyectivo para todo $k\in\bb Z$.
\end{env}


\subsection{Característica de Euler--Poincar\'e y polinomio de Hilbert}
\label{subsection:III.2.5}

\begin{env}[2.5.1]
\label{III.2.5.1}
Sea $A$ un anillo artiniano y sea $X$ un $A$-esquema proyectivo sobre
\[
  Y=\Spec(A).
\]
Para todo $\sh{O}_X$-módulo coherente $\sh{F}$, los $A$-módulos
\[
  \HH^i(X,\sh{F})\qquad(i\geq0)
\]
son de tipo finito \sref{III.2.2.1} y, por tanto, tienen longitud finita, puesto que $A$ es artiniano.
Además, \sref{III.2.2.1} muestra que $\HH^i(X,\sh{F})=0$ salvo para un número finito de $i\geq0$.
Así, el entero
\[
\label{III.2.5.1.1}
  \chi_A(\sh{F})
  =
  \sum_{i=0}^{\infty}
  (-1)^i\operatorname{length}_A\bigl(\HH^i(X,\sh{F})\bigr)
  \tag{2.5.1.1}
\]
está definido para todo $\sh{O}_X$-módulo coherente $\sh{F}$.
Cuando $A$ es un anillo local artiniano, $\chi_A(\sh{F})$ se llama la \emph{característica de Euler--Poincar\'e} de $\sh{F}$ con respecto a $A$.
Para $\sh{F}=\sh{O}_X$, el entero $\chi_A(\sh{O}_X)$ se llama el \emph{género aritmético} de $X$ con respecto a $A$.
\end{env}

\begin{proposition}[2.5.2]
\label{III.2.5.2}
Sea
\[
  0\longrightarrow\sh{F}'
  \longrightarrow\sh{F}
  \longrightarrow\sh{F}''
  \longrightarrow0
\]
una sucesión exacta de $\sh{O}_X$-módulos coherentes.
Entonces
\[
\label{III.2.5.2.1}
  \chi_A(\sh{F})
  =
  \chi_A(\sh{F}')+\chi_A(\sh{F}'').
  \tag{2.5.2.1}
\]
\end{proposition}

\begin{proof}
Los módulos de cohomología de $\sh{F}$, $\sh{F}'$ y $\sh{F}''$ se anulan en todos los grados suficientemente grandes.
Existe, por tanto, un entero $r>0$ para el cual la sucesión exacta larga de cohomología tiene la forma
\[
  0\longrightarrow\HH^0(X,\sh{F}')
  \longrightarrow\HH^0(X,\sh{F})
  \longrightarrow\HH^0(X,\sh{F}'')
  \longrightarrow\HH^1(X,\sh{F}')
  \longrightarrow\cdots
\]
\[
  \cdots\longrightarrow\HH^r(X,\sh{F}')
  \longrightarrow\HH^r(X,\sh{F})
  \longrightarrow\HH^r(X,\sh{F}'')
  \longrightarrow0.
\]
En una sucesión exacta de $A$-módulos de longitud finita con términos nulos en ambos extremos, la suma alternada de las longitudes es nula \sref[0]{0.11.10.1}.
La aplicación de este hecho da inmediatamente \eqref{III.2.5.2.1}.
\end{proof}

El resultado de \sref{III.2.5.2} se aplica siempre que $X$ sea un $A$-esquema cuasicompacto y que los $A$-módulos $\HH^i(X,\sh{F})$ sean de tipo finito para todo $\sh{O}_X$-módulo coherente $\sh{F}$ \sref{III.1.4.12}.

\begin{theorem}[2.5.3]
\label{III.2.5.3}
Sea $A$ un anillo local artiniano, sea $X$ un esquema proyectivo sobre
\[
  Y=\Spec(A),
\]
sea $\sh{L}$ un $\sh{O}_X$-módulo invertible muy amplio relativo a $Y$, y sea $\sh{F}$ un $\sh{O}_X$-módulo coherente.
Para todo $n\in\bb Z$, pongamos
\[
  \sh{F}(n)=\sh{F}\otimes_{\sh{O}_X}\sh{L}^{\otimes n}.
\]
\begin{enumerate}
  \item[(i)] Existe un único polinomio $P\in\bb Q[T]$ tal que
  \[
    \chi_A(\sh{F}(n))=P(n)
  \]
  para todo $n\in\bb Z$.
  El polinomio $P$ se llama el \emph{polinomio de Hilbert} de $\sh{F}$ con respecto a $A$.

  \item[(ii)] Para todo $n$ suficientemente grande,
  \[
    \chi_A(\sh{F}(n))
    =
    \operatorname{length}_A\Gamma(X,\sh{F}(n)).
  \]

  \item[(iii)] El coeficiente del término de mayor grado de $\chi_A(\sh{F}(n))$ es no negativo.
\end{enumerate}
\end{theorem}

Además, demostraremos en el Capítulo~IV, en la sección dedicada a la dimensión, que \emph{el grado de $\chi_A(\sh{F}(n))$ es igual a la dimensión del soporte de $\sh{F}$}.

\begin{proof}
Para todo $n$ suficientemente grande, se tiene
\[
  \HH^i(X,\sh{F}(n))=0
\]
para todo $i>0$ \sref{III.2.2.1}.
\oldpage[III]{110}
Por consiguiente,
\[
  \chi_A(\sh{F}(n))
  =
  \operatorname{length}_A\HH^0(X,\sh{F}(n))
  =
  \operatorname{length}_A\Gamma(X,\sh{F}(n))
\]
para todo $n$ suficientemente grande, lo que demuestra \emph{(ii)}.
Se sigue que $\chi_A(\sh{F}(n))\geq0$ para todo $n$ suficientemente grande y, por tanto, \emph{(iii)} se deduce de \emph{(i)}.
La aserción de unicidad de \emph{(i)} es inmediata, por lo que queda demostrar la existencia de $P$.

Mostremos primero que podemos suponer
\[
  \mathfrak{m}\sh{F}=0,
\]
donde $\mathfrak{m}$ es el ideal maximal de $A$.
En efecto, existe un entero $s>0$ tal que $\mathfrak{m}^s=0$ y, por tanto, $\sh{F}(n)$ posee la filtración finita
\[
  \sh{F}(n)
  \supset\mathfrak{m}\sh{F}(n)
  \supset\cdots
  \supset\mathfrak{m}^{s-1}\sh{F}(n)
  \supset0.
\]
Una inducción que utiliza \eqref{III.2.5.2.1} da
\[
  \chi_A(\sh{F}(n))
  =
  \sum_{k=1}^{s}
  \chi_A\left(
    \mathfrak{m}^{k-1}\sh{F}(n)/
    \mathfrak{m}^{k}\sh{F}(n)
  \right).
\]
Pero
\[
  \mathfrak{m}^{k-1}\sh{F}(n)/
  \mathfrak{m}^{k}\sh{F}(n)
  =
  \sh{F}'_k(n),
  \qquad
  \sh{F}'_k=
  \mathfrak{m}^{k-1}\sh{F}/\mathfrak{m}^{k}\sh{F},
\]
lo que demuestra la reducción.

Supongamos, pues, que $\mathfrak{m}\sh{F}=0$.
Sea $X'$ el subesquema cerrado de $X$ obtenido como imagen inversa, por el morfismo estructural
\[
  X\longrightarrow\Spec(A),
\]
del único punto cerrado de $\Spec(A)$, y sea $j:X'\to X$ la inyección canónica.
Entonces
\[
  \sh{F}=j_*(\sh{F}')
\]
para un $\sh{O}_{X'}$-módulo coherente $\sh{F}'$, y $X'$ es proyectivo sobre $\Spec(K)$, donde
\[
  K=A/\mathfrak{m}.
\]
Si $\sh{L}'=j^*(\sh{L})$, entonces $\sh{L}'$ es muy amplio relativo a $\Spec(K)$ \sref[II]{II.4.4.10}, y
\[
  \sh{F}(n)=j_*(\sh{F}'(n)),
  \qquad
  \sh{F}'(n)
  =
  \sh{F}'\otimes_{\sh{O}_{X'}}\sh{L}'^{\otimes n}
\]
\sref[0]{0.5.4.10}.
Se sigue que
\[
  \chi_A(\sh{F}(n))
  =
  \chi_K(\sh{F}'(n))
\]
(G, II, 4.9.1), y quedamos reducidos al caso en que $A$ es un cuerpo.

Ahora puede considerarse $X$ como un subesquema cerrado de
\[
  P=\mathbf{P}_A^r
\]
para un $r$ conveniente \sref[II]{II.5.5.4}[ii].
Si $i:X\to P$ es la inyección canónica, el mismo razonamiento da
\[
  \chi_A(\sh{F}(n))
  =
  \chi_A\bigl(i_*(\sh{F})(n)\bigr).
\]
Podemos, por tanto, limitarnos al caso
\[
  X=\mathbf{P}_A^r=\Proj(S),
  \qquad
  S=A[T_0,\ldots,T_r],
\]
donde $A$ es un cuerpo.

En este caso, $\sh{F}=\widetilde{M}$ para un $S$-módulo graduado $M$ de tipo finito \sref[II]{II.2.7.8}.
Por el teorema de las sicigias de Hilbert (M, VIII, 6.5), $M$ posee una resolución finita por $S$-módulos libres graduados de tipo finito:
\[
  0\longrightarrow L_q
  \longrightarrow L_{q-1}
  \longrightarrow\cdots
  \longrightarrow L_1
  \longrightarrow M
  \longrightarrow0.
\]
Como el funtor $M\mapsto\widetilde{M}$ es exacto \sref[II]{II.2.5.4}, obtenemos una sucesión exacta
\[
  0\longrightarrow\widetilde{L}_q
  \longrightarrow\widetilde{L}_{q-1}
  \longrightarrow\cdots
  \longrightarrow\widetilde{L}_1
  \longrightarrow\widetilde{M}
  \longrightarrow0,
\]
y, por tanto, para todo $n\in\bb Z$, una sucesión exacta
\[
  0\longrightarrow\widetilde{L}_q(n)
  \longrightarrow\widetilde{L}_{q-1}(n)
  \longrightarrow\cdots
  \longrightarrow\widetilde{L}_1(n)
  \longrightarrow\widetilde{M}(n)
  \longrightarrow0.
\]
Una inducción sobre $q$, utilizando \sref{III.2.5.2}, da
\[
  \chi_A(\widetilde{M}(n))
  =
  \sum_{j=1}^{q}
  (-1)^{j+1}\chi_A(\widetilde{L}_j(n)).
\]
Para demostrar \emph{(i)}, quedamos, por tanto, reducidos al caso en que $M$ es graduado, libre y de tipo finito y, por consiguiente, al caso
\[
  M=S(h)
\]
para algún $h\in\bb Z$.
Entonces
\[
  \widetilde{M}(n)
  =
  \widetilde{M(n)}
  =
  \widetilde{S(n+h)}
\]
\sref[II]{II.2.5.15}, de modo que el teorema se deduce del lema siguiente.
\end{proof}

\oldpage[III]{111}
\begin{lemma}[2.5.3.1]
\label{III.2.5.3.1}
Sea $A$ un cuerpo, sea $r>0$ un entero y sea
\[
  X=\mathbf{P}_A^r.
\]
Entonces
\[
  \chi_A(\sh{O}_X(n))
  =
  \binom{n+r}{r}
\]
para todo $n\in\bb Z$.
\end{lemma}

\begin{proof}
Para $n\geq0$,
\[
  \chi_A(\sh{O}_X(n))
  =
  \operatorname{length}_A\HH^0(X,\sh{O}_X(n)),
\]
que es el número de monomios en los $T_i$ de grado total $n$, a saber,
\[
  \binom{n+r}{r}
\]
\sref{III.2.1.12}.

Para $n\leq-r-1$, se tiene análogamente
\[
  \chi_A(\sh{O}_X(n))
  =
  (-1)^r\operatorname{length}_A\HH^r(X,\sh{O}_X(n)).
\]
Escribiendo $n=-r-h$, la dimensión de $\HH^r(X,\sh{O}_X(n))$ sobre $A$ es el número de sucesiones
\[
  (p_i)_{0\leq i\leq r}
\]
de enteros positivos que satisfacen
\[
  \sum_{i=0}^{r}p_i=r+h
\]
\sref{III.2.1.12}.
Equivalentemente, es el número de sucesiones de enteros no negativos
\[
  (q_i)_{0\leq i\leq r}
\]
tales que
\[
  \sum_{i=0}^{r}q_i=h-1.
\]
Este número es
\[
  \binom{h+r-1}{r}
  =
  (-1)^r\binom{n+r}{r}.
\]
Finalmente, si $-r\leq n<0$, entonces
\[
  \binom{n+r}{r}=0,
\]
y $\HH^i(X,\sh{O}_X(n))=0$ para todo $i\geq0$ \sref{III.2.1.12}.
Esto demuestra el lema.
\end{proof}

\begin{corollary}[2.5.4]
\label{III.2.5.4}
Sea $A$ un anillo local artiniano, sea $S$ una $A$-álgebra graduada de tipo finito engendrada por $S_1$, sea $M$ un $S$-módulo graduado de tipo finito y sea
\[
  X=\Proj(S).
\]
Entonces, para todo $n$ suficientemente grande,
\[
  \chi_A(\widetilde{M}(n))
  =
  \operatorname{length}_A(M_n).
\]
\end{corollary}

\begin{proof}
Para todo $n$ suficientemente grande, los módulos
\[
  M_n
  \quad\text{y}\quad
  \Gamma(X,\widetilde{M}(n))
\]
son isomorfos \sref{III.2.3.1}.
\end{proof}


\subsection{Aplicación: criterios de amplitud}
\label{subsection:III.2.6}

\begin{proposition}[2.6.1]
\label{III.2.6.1}
Sea $Y$ un preesquema noetheriano, sea $f:X\to Y$ un morfismo propio,
y sea $\sh{L}$ un $\sh{O}_X$-módulo invertible.
Las condiciones siguientes son equivalentes:
\begin{enumerate}
  \item[\emph{a)}] $\sh{L}$ es amplio para $f$.
  \item[\emph{b)}] Para todo $\sh{O}_X$-módulo coherente $\sh{F}$, existe
  un entero $N$ tal que, para $n\geq N$,
  \[
    \RR^qf_*(\sh{F}\otimes\sh{L}^{\otimes n})=0
  \]
  para todo $q>0$.
  \item[\emph{c)}] Para todo ideal coherente $\sh{J}$ de $\sh{O}_X$, existe
  un entero $N$ tal que, para $n\geq N$,
  \[
    \RR^1f_*(\sh{J}\otimes\sh{L}^{\otimes n})=0.
  \]
\end{enumerate}
\end{proposition}

\begin{proof}
Hemos visto que \emph{a)} implica \emph{b)} \sref{III.2.2.1},~(ii).
Es inmediato que \emph{b)} implica \emph{c)}, y queda demostrar
que \emph{c)} implica \emph{a)}.
Podemos limitarnos al caso en que $Y$ es afín \sref[II]{II.4.6.4}.
Basta mostrar que, cuando $h$ recorre las secciones de
$\sh{L}^{\otimes n}$ sobre $X$, con $n>0$, los subconjuntos abiertos afines $X_h$
forman un recubrimiento de $X$ \sref[II]{II.4.5.2}.

Para ello, sea $x$ un punto cerrado de $X$ y sea $U$ un entorno
abierto afín de $x$.
Encontraremos un entero $n$ y una sección
\[
  h\in\Gamma(X,\sh{L}^{\otimes n})
\]
tal que $x\in X_h\subset U$.
El subconjunto abierto $X_h$ será entonces afín \sref[I]{I.1.3.6}.
La unión de todos esos $X_h$ es un subconjunto abierto de $X$ que contiene todo
punto cerrado y, por tanto, es todo $X$, puesto que $X$ es noetheriano
\sref[I]{I.6.3.7} y \sref[0]{0.2.1.3}.

Sean $\sh{J}$ y $\sh{J}'$, respectivamente, los ideales cuasicoherentes de
$\sh{O}_X$ que definen los subpreesquemas cerrados reducidos cuyos espacios
subyacentes son $X-U$ y $(X-U)\cup\{x\}$, respectivamente \sref[I]{I.5.2.1}.
Los ideales $\sh{J}$ y $\sh{J}'$ son coherentes \sref[I]{I.6.1.1},
$\sh{J}'\subset\sh{J}$, y
\[
  \sh{J}''=\sh{J}/\sh{J}'
\]
es un $\sh{O}_X$-módulo coherente \sref[0]{0.5.3.3}, con soporte
$\{x\}$, y $\sh{J}''_x=\kres(x)$.
Como $\sh{L}$ es localmente libre, la sucesión
\[
  0\longrightarrow
  \sh{J}'\otimes\sh{L}^{\otimes n}
  \longrightarrow
  \sh{J}\otimes\sh{L}^{\otimes n}
  \longrightarrow
  \sh{J}''\otimes\sh{L}^{\otimes n}
  \longrightarrow0
\]
es exacta para todo $n$.
Por hipótesis, para todo $n$ suficientemente grande,
\[
  \HH^1(X,\sh{J}'\otimes\sh{L}^{\otimes n})=0.
\]

\oldpage[III]{112}

La sucesión exacta de cohomología muestra, por tanto, que
\[
  \Gamma(X,\sh{J}\otimes\sh{L}^{\otimes n})
  \longrightarrow
  \Gamma(X,\sh{J}''\otimes\sh{L}^{\otimes n})
\]
es sobreyectivo.
Una sección $g$ de $\sh{J}''\otimes\sh{L}^{\otimes n}$ sobre $X$ tal que
$g(x)\neq0$ es, por consiguiente, la imagen de una sección
\[
  h\in\Gamma(X,\sh{J}\otimes\sh{L}^{\otimes n})
  \subset\Gamma(X,\sh{L}^{\otimes n});
\]
la inclusión se deduce de \sref[0]{0.5.4.1}.
Por construcción, $h(x)\neq0$, mientras que $h(z)=0$ para todo $z\notin U$.
Así, $x\in X_h\subset U$, lo que demuestra la aserción.
\end{proof}

\begin{proposition}[2.6.2]
\label{III.2.6.2}
Sea $Y$ un preesquema noetheriano, sea $f:X\to Y$ un morfismo de
tipo finito, sea $g:X'\to X$ un morfismo finito sobreyectivo, sea
$\sh{L}$ un $\sh{O}_X$-módulo invertible, y pongamos
\[
  \sh{L}'=g^*(\sh{L}).
\]
Supongamos que existe un subconjunto $Z$ de $X$, propio sobre $Y$
\sref[II]{II.5.4.10}, tal que, para todo $x\in X-Z$, o bien $X$ es
normal en $x$, o bien
\[
  \bigl(g_*(\sh{O}_{X'})\bigr)_x
\]
es un $\sh{O}_{X,x}$-módulo libre.
Entonces $\sh{L}$ es amplio para $f$ si y solo si $\sh{L}'$ es amplio para
$f\circ g$.
\end{proposition}

\begin{proof}
\begin{env}[2.6.2.1]
\label{III.2.6.2.1}
Como $g$ es afín, la condición es necesaria \sref[II]{II.5.1.12}.
Para demostrar que es suficiente, podemos suponer que $Y$ es afín
\sref[II]{II.4.6.4}.
Mostremos primero que también podemos limitarnos al caso en que $X$ es
reducido.

Sea $j:X_{\mathrm{red}}\to X$ la inmersión canónica, y pongamos
\[
  X_1=X_{\mathrm{red}},
  \qquad
  X'_1=X'\times_X X_1.
\]
Tenemos entonces el diagrama conmutativo
\[
  \xymatrix{
    X'\ar[d]_g &
    X'_1\ar[l]_{j'}\ar[d]^{g_1}\\
    X &
    X_1\ar[l]^{j}
  }
  \label{III.2.6.2.2}\tag{2.6.2.2}
\]
El morfismo $f\circ j$ es de tipo finito \sref[I]{I.6.3.4}, y
$g_1$ es finito \sref[II]{II.6.1.5},~(iii).
Si $\sh{L}'$ es amplio para $f\circ g$, entonces
$j'^*(\sh{L}')$ es amplio para $f\circ g\circ j'$, porque $j'$ es una
inmersión cerrada \sref[II]{II.5.1.12} y \sref[I]{I.4.3.2}.
Si $Z_1=j^{-1}(Z)$, entonces $Z_1$ es propio sobre $Y$
\sref[II]{II.5.4.10}.
Además, si $X$ es normal en un punto $x$, también lo es $X_{\mathrm{red}}$
en ese punto; y si
$\bigl(g_*(\sh{O}_{X'})\bigr)_x$ es un $\sh{O}_{X,x}$-módulo libre,
entonces \sref[II]{II.1.5.2} muestra que
\[
  \bigl((g_1)_*(\sh{O}_{X'_1})\bigr)_x
\]
es un $\sh{O}_{X_1,x}$-módulo libre.
Finalmente, como $X$ es noetheriano \sref[I]{I.6.3.7}, la amplitud de
$j^*(\sh{L})$ implica la de $\sh{L}$ \sref[II]{II.4.5.14}.
Como
\[
  j'^*(\sh{L}')=g_1^*(j^*(\sh{L})),
\]
esto demuestra la reducción anunciada.
En lo sucesivo suponemos que $Y$ es afín y que $X$ es reducido.
\end{env}

Se satisfacen ahora las hipótesis de \sref[II]{II.6.6.11}.
Existen un $Y$-preesquema reducido $X_2$ y un
$Y$-morfismo finito birracional
\[
  h:X_2\to X
\]
cuya restricción a $h^{-1}(X-Z)$ es un isomorfismo sobre $X-Z$, y
tal que $h^*(\sh{L})$ es amplio.
Reemplazando $X'$ por $X_2$, quedamos reducidos a demostrar la proposición bajo
la hipótesis adicional de que $g$ posee las propiedades que acabamos de enumerar para
$h$.
Seguimos denotando por $Z$ un subpreesquema cerrado de $X$, propio sobre
$Y$ \sref[II]{II.5.4.10}, cuyo espacio subyacente es el subconjunto dado
$Z$.

\begin{env}[2.6.2.3]
\label{III.2.6.2.3}
Sea $X_1$ un subpreesquema cerrado de $X$, sea $j:X_1\to X$ la
inmersión canónica, y sea
\[
  X'_1=g^{-1}(X_1)=X'\times_X X_1,
\]
con inmersión canónica $j':X'_1\to X'$.
El cuadrado correspondiente es de nuevo \eqref{III.2.6.2.2}.
Pongamos
\[
  \sh{L}_1=j^*(\sh{L}),
  \qquad
  \sh{L}'_1=j'^*(\sh{L}')=g_1^*(\sh{L}_1).
\]
Entonces $\sh{L}'_1$ es amplio para $f\circ g\circ j'$
\sref[II]{II.5.1.12}.
Si $Z_1=j^{-1}(Z)$, el subpreesquema cerrado $Z_1$ de $X_1$ es propio
sobre $Y$ \sref[II]{II.5.4.2},~(ii).
En otras palabras, las hipótesis de \sref{III.2.6.2} se cumplen para
$X_1$, $\sh{L}_1$, $g_1$ y $Z_1$.
\end{env}

\oldpage[III]{113}

Esto nos permite demostrar \sref{III.2.6.2} por inducción noetheriana
\sref[0]{0.2.2.2} cuando la restricción de $g$ a $g^{-1}(X-Z)$ es un
isomorfismo sobre $X-Z$; como se vio en \sref{III.2.6.2.1}, ese caso es
suficiente.
Basta mostrar que, si la conclusión de \sref{III.2.6.2} se cumple
para $\sh{L}_1$ sobre todo subpreesquema cerrado $X_1$ de $X$ cuyo
espacio subyacente no es $X$, entonces también se cumple para $\sh{L}$.

\begin{env}[2.6.2.4]
\label{III.2.6.2.4}
Pongamos
\[
  \sh{A}=\sh{O}_X,
  \qquad
  \sh{B}=g_*(\sh{O}_{X'}).
\]
Como $g$ es birracional y $X$ es reducido, $\sh{B}$ es una
sub-$\sh{A}$-álgebra del haz $\sh{R}(X)$ de anillos totales de fracciones,
y es un $\sh{A}$-módulo coherente.
Además,
\[
  \sh{B}|_{X-Z}=\sh{A}|_{X-Z}.
\]
Sea $\sh{K}$ el conductor de $\sh{B}$ sobre $\sh{A}$, es decir, el mayor
sub-$\sh{A}$-módulo cuasicoherente de $\sh{A}$ tal que
\[
  \sh{B}\sh{K}\subset\sh{A}.
\]
Equivalentemente, $\sh{K}$ es el anulador del
$\sh{A}$-módulo $\sh{B}/\sh{A}$ \sref[0]{0.5.3.7}.
Se sigue que
\[
  \sh{B}\sh{K}=\sh{K}.
\]
Está claro que $\sh{K}_x=\sh{A}_x$ en todo punto que admite un entorno
$W_x$ sobre el cual $g:g^{-1}(W_x)\to W_x$ es un isomorfismo.
Esto ocurre, en particular, en todo punto de $X-Z$, y en un
entorno de todo punto genérico de una componente irreducible de
$X$.

Sea
\[
  Z_1=\Spec(\sh{A}/\sh{K})
\]
el subpreesquema cerrado de $X$ definido por $\sh{K}$.
Es propio sobre $Y$, porque su espacio subyacente es cerrado en $Z$
\sref[II]{II.5.4.10}.
La definición de $\sh{K}$ da también
\[
  \sh{B}|_{X-Z_1}=\sh{A}|_{X-Z_1}.
\]
Podemos, por tanto, reducirnos al caso
\[
  Z=\Spec(\sh{A}/\sh{K}).
\]
Como $X-Z_1$ es un subconjunto abierto no vacío de $X$, podemos además
suponer que el espacio subyacente de $Z$ es distinto de $X$.
\end{env}

\begin{env}[2.6.2.5]
\label{III.2.6.2.5}
Consideremos $X'$ como $\Spec(\sh{B})$, puesto que $g$ es afín, y sea
\[
  \sh{K}'=\widetilde{\sh{K}}.
\]
Este es un ideal coherente de $\sh{O}_{X'}$ que satisface
$g_*(\sh{K}')=\sh{K}$ \sref[II]{II.1.4.1}.
El subpreesquema cerrado
\[
  Z'=g^{-1}(Z)=Z\times_X X'
\]
de $X'$ está definido por $\sh{K}'$ y es igual a
\[
  \Spec(\sh{B}/\sh{K})
\]
\sref[II]{II.1.4.10}.
El morfismo $h:Z'\to Z$ es finito \sref[II]{II.6.1.5},~(iii);
por consiguiente, $Z'$ es propio sobre $Y$ \sref[II]{II.6.1.11} y
\sref[II]{II.5.4.2},~(ii).

Debemos demostrar que, para todo $x\in X$ y todo entorno abierto $U$
de $x$, existen un entero $n>0$ y una sección $s$ de
$\sh{L}^{\otimes n}$ sobre $X$ tales que
\[
  x\in X_s\subset U
\]
\sref[II]{II.4.5.2}.
Distinguimos dos casos.

\smallskip
\noindent
$1^\circ$ Supongamos primero que $x\in X-Z$.
Podemos suponer evidentemente que $U\subset X-Z$, de modo que
$U'=g^{-1}(U)$ no corta a $Z'$.
Como $\sh{L}'$ es amplio, existen $n>0$ y una sección $s'$ de
$\sh{L}'^{\otimes n}$ sobre $X'$ tales que
\[
  x'=g^{-1}(x)\in X'_{s'}\subset g^{-1}(U)
\]
\sref[II]{II.4.5.2}.
También podemos suponer que
$\sh{K}'\otimes\sh{L}'^{\otimes n}$ está engendrado por sus secciones globales
\sref[II]{II.4.5.5}.
Como $\sh{K}'_{x'}=\sh{O}_{X',x'}$, una de estas secciones, digamos $s''$, es
no nula en $x'$.
Multiplicándola por $s'$ y reemplazando así $n$ por $2n$, podemos conseguir
que
\[
  x'\in X'_{s''}\subset g^{-1}(U).
\]
Por \sref[0]{0.5.4.10}, existe un isomorfismo canónico
\[
  \Gamma(X,\sh{K}\otimes\sh{L}^{\otimes n})
  \xrightarrow{\sim}
  \Gamma(X',\sh{K}'\otimes\sh{L}'^{\otimes n}).
\]
La sección $s$ de $\sh{K}\otimes\sh{L}^{\otimes n}$ correspondiente a
$s''$ posee manifiestamente las propiedades requeridas.

\smallskip
\noindent
$2^\circ$ Supongamos ahora que $x\in Z$.
Sea $\sh{J}$ el ideal coherente de $\sh{O}_X$ que define el
subpreesquema cerrado reducido cuyo espacio subyacente es $X-U$, y consideremos en
$\sh{B}$ los ideales coherentes

\oldpage[III]{114}

\[
  \sh{J}\sh{B}
  \qquad\text{y}\qquad
  \sh{J}\sh{K}
  =
  \sh{J}(\sh{K}\sh{B})
  =
  \sh{K}(\sh{J}\sh{B}).
\]
Tenemos el diagrama de inclusiones
\[
  \xymatrix@R=1.5em{
    \sh{J}\sh{B}\ar@{^{(}->}[r] &
    \sh{B}\\
    \sh{J}\ar@{^{(}->}[r]\ar@{^{(}->}[u] &
    \sh{A}\ar@{^{(}->}[u]\\
    \sh{J}\sh{K}\sh{B}=\sh{J}\sh{K}
      \ar@{^{(}->}[r]\ar@{^{(}->}[u] &
    \sh{K}\ar@{^{(}->}[u]
  }
  \label{III.2.6.2.6}\tag{2.6.2.6}
\]
Sea $\sh{J}'=(\sh{J}\sh{B})^\sim$ el ideal coherente de
$\sh{O}_{X'}$ tal que
\[
  g_*(\sh{J}')=\sh{J}\sh{B}.
\]
Entonces
\[
  \sh{J}'\sh{K}'=(\sh{J}\sh{K}\sh{B})^\sim,
  \qquad
  \sh{J}'/\sh{J}'\sh{K}'
  =
  (\sh{J}\sh{B}/\sh{J}\sh{K}\sh{B})^\sim
\]
\sref[II]{II.1.4.4}.
Como $\sh{J}|_V=\sh{J}\sh{K}|_V$ sobre todo subconjunto abierto $V$ disjunto
de $Z$, el soporte de $\sh{J}'/\sh{J}'\sh{K}'$ está contenido en
$Z'$.
Como $Z'$ es propio sobre $Y$, \sref{III.2.2.4} muestra que, para todo
$n$ suficientemente grande, la aplicación canónica
\[
  \Gamma(X',\sh{J}'\otimes\sh{L}'^{\otimes n})
  \longrightarrow
  \Gamma\bigl(
    X',
    (\sh{J}'/\sh{J}'\sh{K}')\otimes\sh{L}'^{\otimes n}
  \bigr)
\]
es sobreyectiva.
Por \sref[0]{0.5.4.10}, se sigue que la aplicación canónica
\[
  \Gamma(X,\sh{J}\sh{B}\otimes\sh{L}^{\otimes n})
  \longrightarrow
  \Gamma\bigl(
    X,
    (\sh{J}\sh{B}/\sh{J}\sh{K}\sh{B})
      \otimes\sh{L}^{\otimes n}
  \bigr)
\]
es sobreyectiva.

Sean $i:Z\to X$ e $i':Z'\to X'$ las inmersiones canónicas.
Forman parte del diagrama conmutativo
\[
  \xymatrix{
    X'\ar[d]_g &
    Z'\ar[l]_{i'}\ar[d]^h\\
    X &
    Z\ar[l]^{i}
  }
\]
Pongamos
\[
  \sh{M}=i^*(\sh{L}),
  \qquad
  \sh{M}'=i'^*(\sh{L}').
\]
Como $\sh{L}'$ es amplio, $\sh{M}'$ es amplio
\sref[II]{II.5.1.12}, y
\[
  \sh{M}'=h^*(\sh{M}).
\]
La hipótesis de inducción noetheriana, aplicada puesto que $Z\neq X$, muestra ahora
que $\sh{M}$ es amplio.
Por consiguiente, para todo $n$ suficientemente grande,
\[
  i^*(\sh{J}/\sh{J}\sh{K})\otimes\sh{M}^{\otimes n}
\]
está engendrado por sus secciones sobre $Z$ \sref[II]{II.4.5.5}.
Como
\[
  \sh{J}/\sh{J}\sh{K}
  =
  i_*\bigl(i^*(\sh{J}/\sh{J}\sh{K})\bigr),
\]
\sref[0]{0.5.4.10} da una sección $s$ de
\[
  (\sh{J}/\sh{J}\sh{K})\otimes\sh{L}^{\otimes n}
\]
sobre $X$, para algún $n$ suficientemente grande, tal que $s(x)\neq0$.
En efecto, $\sh{J}_x=\sh{A}_x$ por la definición de $\sh{J}$, mientras que
$\sh{K}_x\neq\sh{A}_x$ por hipótesis.

El diagrama \eqref{III.2.6.2.6} muestra que $s$ es también una sección de
\[
  (\sh{J}\sh{B}/\sh{J}\sh{K}\sh{B})
  \otimes\sh{L}^{\otimes n}
\]
sobre $X$.
Es, por tanto, la imagen canónica de una sección
\[
  t\in\Gamma(X,\sh{J}\sh{B}\otimes\sh{L}^{\otimes n}).
\]
Por definición, la imagen de $t$ módulo
$\sh{J}\sh{K}\sh{B}\otimes\sh{L}^{\otimes n}$ pertenece a
\[
  \frac{\sh{J}\otimes\sh{L}^{\otimes n}}
       {\sh{J}\sh{K}\sh{B}\otimes\sh{L}^{\otimes n}}.
\]
El diagrama \eqref{III.2.6.2.6} implica, por tanto, que $t$ es una sección
de $\sh{J}\otimes\sh{L}^{\otimes n}$ y, en particular, una sección de
$\sh{L}^{\otimes n}$.
Se tiene $t(x)\neq0$, luego $x\in X_t$.
Por la definición de $\sh{J}$, la sección $t$ se anula sobre $X-U$, el
soporte de $\sh{O}_X/\sh{J}$.
Así, $X_t\subset U$, y la demostración está terminada.
\end{env}
\end{proof}

\oldpage[III]{115}

\begin{env}[2.6.3]
\label{III.2.6.3}
\emph{Observación.}
Cuando $X$ es propio sobre $Y$, la Proposición~\sref{III.2.6.2} puede demostrarse
más sencillamente razonando como en el teorema de Chevalley
\sref[II]{II.6.7.1}, con ayuda de la Proposición~\sref{III.2.6.1} y
del lema \sref[II]{II.6.7.1.1}.
\end{env}

\section{Teorema de finitud para los morfismos propios}
\label{section:III.3}

\subsection{El lema de d\'evissage}
\label{subsection:III.3.1}

\begin{definition}[3.1.1]
\label{III.3.1.1}
Sea $\C$ una categoría abeliana.
Diremos que un subconjunto $\C'$ del conjunto de objetos de $\C$ es \emph{exacto} si $0\in\C'$ y si, para toda sucesión exacta $0\to A'\to A\to A''\to 0$ en $\C$ tal que dos de los objetos $A$, $A'$, $A''$ pertenecen a $\C'$, el tercero pertenece también a $\C'$.
\end{definition}

\begin{theorem}[3.1.2]
\label{III.3.1.2}
Sea $X$ un preesquema noetheriano; denotemos por $\C$ la categoría abeliana de los $\sh{O}_X$-módulos coherentes.
Sea $\C'$ un subconjunto exacto de $\C$, y sea $X'$ un subconjunto cerrado del espacio subyacente a $X$.
Supongamos que, para todo subconjunto cerrado irreducible $Y$ de $X'$, de punto genérico $y$, existe un $\sh{O}_X$-módulo $\sh{G}\in\C'$ de soporte $Y$ tal que $\sh{G}_y$ es un espacio vectorial sobre $\kres(y)$ de dimensión~$1$.
Entonces todo $\sh{O}_X$-módulo coherente cuyo soporte esté contenido en $X'$ pertenece a $\C'$ (y, en particular, si $X'=X$, se tiene $\C'=\C$).
\end{theorem}

\begin{proof}
Consideremos la siguiente propiedad $\textbf{P}(Y)$ de un subconjunto cerrado $Y$ de $X'$: todo $\sh{O}_X$-módulo coherente cuyo soporte esté contenido en $Y$ pertenece a $\C'$.
En virtud del principio de inducción noetheriana \sref[0]{0.2.2.2}, vemos que basta demostrar que \emph{si $Y$ es un subconjunto cerrado de $X'$ tal que la propiedad $\textbf{P}(Y')$ es verdadera para todo subconjunto cerrado $Y'$ de $Y$, distinto de $Y$, entonces $\textbf{P}(Y)$ es verdadera}.

Sea, pues, $\sh{F}\in\C$ de soporte contenido en $Y$, y demostremos que $\sh{F}\in\C'$.
Denotemos también por $Y$ el subpreesquema cerrado reducido de $X$ cuyo espacio subyacente es $Y$ \sref[I]{I.5.2.1}; está definido por un haz coherente de ideales $\sh{J}$ de $\sh{O}_X$.
Sabemos \sref[I]{I.9.3.4} que existe un entero $n>0$ tal que $\sh{J}^n\sh{F}=0$; para $1\leq k\leq n$ se tiene, por tanto, una sucesión exacta
\[
  0\to\sh{J}^{k-1}\sh{F}/\sh{J}^k\sh{F}\to\sh{F}/\sh{J}^k\sh{F}\to\sh{F}/\sh{J}^{k-1}\sh{F}\to 0
\]
 de $\sh{O}_X$-módulos coherentes (\sref[0]{0.5.3.6} y \sref[0]{0.5.3.3}); como $\C'$ es exacto, una inducción sobre $k$ muestra que basta probar que cada $\sh{F}_k=\sh{J}^{k-1}\sh{F}/\sh{J}^k\sh{F}$ pertenece a $\C'$.
Nos reducimos así a probar que $\sh{F}\in\C'$ bajo la hipótesis adicional $\sh{J}\sh{F}=0$; esto equivale a decir que $\sh{F}=j_*(j^*(\sh{F}))$, donde $j$ es la inmersión canónica $Y\to X$.
Consideremos ahora dos casos:
\begin{enumerate}
  \item[(a)] $Y$ es \emph{reducible}.
    Escribamos $Y=Y'\cup Y''$, donde $Y'$ e $Y''$ son subconjuntos cerrados de $Y$, distintos de $Y$; denotemos también por $Y'$ e $Y''$ los subpreesquemas cerrados reducidos de $X$ cuyos espacios subyacentes son, respectivamente, $Y'$ e $Y''$, definidos respectivamente por los haces de ideales $\sh{J}'$ y $\sh{J}''$ de $\sh{O}_X$.
    Pongamos $\sh{F}'=\sh{F}\otimes_{\sh{O}_X}(\sh{O}_X/\sh{J}')$ y $\sh{F}''=\sh{F}\otimes_{\sh{O}_X}(\sh{O}_X/\sh{J}'')$.
    Los homomorfismos canónicos $\sh{F}\to\sh{F}'$ y $\sh{F}\to\sh{F}''$ definen así un homomorfismo $u:\sh{F}\to\sh{F}'\oplus\sh{F}''$.
    Demostremos que, para todo $z\not\in Y'\cap Y''$, el homomorfismo $u_z:\sh{F}_z\to\sh{F}_z'\oplus\sh{F}_z''$ es \emph{biyectivo}.
    En efecto, se tiene $\sh{J}'\cap\sh{J}''=\sh{J}$, pues la cuestión es local y
\oldpage[III]{116}
    la igualdad anterior se deduce de (\sref[I]{I.5.2.1} y \sref[I]{I.1.1.5}); si $z\not\in Y''$, se tiene $\sh{J}_z'=\sh{J}_z$, de donde $\sh{F}_z'=\sh{F}_z$ y $\sh{F}_z''=0$, lo que prueba nuestra afirmación en este caso; se razona del mismo modo si $z\not\in Y'$.
    Por consiguiente, el núcleo y el conúcleo de $u$, que pertenecen a $\C$ \sref[0]{0.5.3.4}, tienen su soporte en $Y'\cap Y''$ y, por hipótesis, pertenecen a $\C'$; por la misma razón, $\sh{F}'$ y $\sh{F}''$ pertenecen a $\C'$, y por tanto también $\sh{F}'\oplus\sh{F}''$, pues $\C'$ es exacto.
    La conclusión se obtiene entonces considerando las dos sucesiones exactas
    \[
      0\to\Im u\to\sh{F}'\oplus\sh{F}''\to\Coker u\to 0,
    \]
    \[
      0\to\Ker u\to\sh{F}\to\Im u\to 0,
    \]
    y la hipótesis de que $\C'$ es exacto.
  \item[(b)] $Y$ es irreducible y, por consiguiente, el subpreesquema $Y$ de $X$ es \emph{íntegro}.
    Si $y$ es su punto genérico, se tiene $(\sh{O}_Y)_y=\kres(y)$ y, como $j^*(\sh{F})$ es un $\sh{O}_Y$-módulo coherente, $\sh{F}_y=(j^*(\sh{F}))_y$ es un espacio vectorial sobre $\kres(y)$ de dimensión finita~$m$.
    Por hipótesis existe un $\sh{O}_X$-módulo coherente $\sh{G}\in\C'$ (necesariamente de soporte $Y$) tal que $\sh{G}_y$ es un espacio vectorial sobre $\kres(y)$ de dimensión~$1$.
    Existe, por tanto, un $\kres(y)$-isomorfismo $(\sh{G}_y)^m\isoto\sh{F}_y$, que es también un $\sh{O}_Y$-isomorfismo; como $\sh{G}^m$ y $\sh{F}$ son coherentes, existe un entorno abierto $W$ de $y$ en $X$ y un isomorfismo $\sh{G}^m|W\isoto\sh{F}|W$ \sref[0]{0.5.2.7}.
    Sea $\sh{H}$ el grafo de este isomorfismo, que es un $(\sh{O}_X|W)$-submódulo coherente de $(\sh{G}^m\oplus\sh{F})|W$, canónicamente isomorfo a $\sh{G}^m|W$ y a $\sh{F}|W$; existe entonces un $\sh{O}_X$-submódulo coherente $\sh{H}_0$ de $\sh{G}^m\oplus\sh{F}$ que induce $\sh{H}$ sobre $W$ y $0$ sobre $X\setmin Y$, puesto que $\sh{G}^m$ y $\sh{F}$ tienen soporte $Y$~\sref[I]{I.9.4.7}.
    Las restricciones $v:\sh{H}_0\to\sh{G}^m$ y $w:\sh{H}_0\to\sh{F}$ de las proyecciones canónicas de $\sh{G}^m\oplus\sh{F}$ son entonces homomorfismos de $\sh{O}_X$-módulos coherentes que, sobre $W$ y sobre $X\setmin Y$, se reducen a isomorfismos; en otras palabras, los núcleos y conúcleos de $v$ y $w$ tienen su soporte en el cerrado $Y\setmin(Y\cap W)$, distinto de $Y$.
    Pertenecen a $\C'$; por otra parte, $\sh{G}^m\in\C'$, pues $\sh{G}\in\C'$ y $\C'$ es exacto.
    De la exactitud de $\C'$ concluimos sucesivamente que $\sh{H}_0\in\C'$ y luego que $\sh{F}\in\C'$.
C.Q.D.
\end{enumerate}
\end{proof}

\begin{corollary}[3.1.3]
\label{III.3.1.3}
Supongamos además que el subconjunto exacto $\C'$ de $\C$ tiene la propiedad de que todo factor directo coherente de un $\sh{O}_X$-módulo coherente $\sh{M}\in\C'$ pertenece también a $\C'$.
En este caso, la conclusión del Teorema~\sref{III.3.1.2} sigue siendo válida cuando la condición «$\sh{G}_y$ es un espacio vectorial sobre $\kres(y)$ de dimensión~$1$» se sustituye por $\sh{G}_y\neq 0$ (lo que equivale a $\Supp(\sh{G})=Y$).
\end{corollary}

\begin{proof}
El razonamiento del Teorema~\sref{III.3.1.2} solo debe modificarse en el caso~(b); ahora $\sh{G}_y$ es un espacio vectorial sobre $\kres(y)$ de dimensión $q>0$ y, por tanto, se tiene un $\sh{O}_Y$-isomorfismo $(\sh{G}_y)^m\isoto(\sh{F}_y)^q$; el final del razonamiento del Teorema~\sref{III.3.1.2} prueba entonces que $\sh{F}^q\in\C'$, y la hipótesis adicional sobre $\C'$ implica que $\sh{F}\in\C'$.
\end{proof}

\subsection{El teorema de finitud: caso de los esquemas ordinarios}
\label{subsection:III.3.2}
\label{III.3.2}

\begin{theorem}[3.2.1]
\label{III.3.2.1}
Sea $Y$ un preesquema localmente noetheriano y sea $f:X\to Y$ un morfismo propio.
Para todo $\sh{O}_X$-módulo coherente $\sh{F}$, los $\sh{O}_Y$-módulos $\RR^q f_*(\sh{F})$ son coherentes para $q\geq 0$.
\end{theorem}

\begin{proof}
Como la cuestión es local sobre $Y$, podemos suponer que $Y$ es noetheriano, y por tanto que $X$ es noetheriano~\sref[I]{I.6.3.7}.
Los $\sh{O}_X$-módulos coherentes $\sh{F}$ para los cuales es válida la conclusión del Teorema~\sref{III.3.2.1} forman un subconjunto \emph{exacto} $\C'$ de la categoría $\C$ de los $\sh{O}_X$-módulos coherentes.
\oldpage[III]{117}
En efecto, sea $0\to\sh{F}'\to\sh{F}\to\sh{F}''\to 0$ una sucesión exacta de $\sh{O}_X$-módulos coherentes; supongamos, por ejemplo, que $\sh{F}'$ y $\sh{F}''$ pertenecen a $\C'$; se tiene la sucesión exacta larga de cohomología
\[
  \RR^{q-1}f_*(\sh{F}'')\xrightarrow{\partial}\RR^q f_*(\sh{F}')\to\RR^q f_*(\sh{F})\to\RR^q f_*(\sh{F}'')\xrightarrow{\partial}\RR^{q+1}f_*(\sh{F}'),
\]
en la que, por hipótesis, los cuatro términos exteriores son coherentes; lo mismo ocurre con el término central $\RR^q f_*(\sh{F})$ por (\sref[0]{0.5.3.4}~y~\sref[0]{0.5.3.3}).
Se demuestra del mismo modo que, cuando $\sh{F}$ y $\sh{F}'$ (resp.~$\sh{F}$ y $\sh{F}''$) pertenecen a $\C'$, también pertenece $\sh{F}''$ (resp.~$\sh{F}'$).
Además, todo \emph{factor directo} coherente $\sh{F}'$ de un $\sh{O}_X$-módulo $\sh{F}\in\C'$ pertenece a $\C'$: en efecto, $\RR^q f_*(\sh{F}')$ es entonces un factor directo de $\RR^q f_*(\sh{F})$~(G,~II,~4.4.4), por lo que es de tipo finito; como es cuasi-coherente~\sref{III.1.4.10}, es coherente, pues $Y$ es noetheriano.
En virtud del Corolario~\sref{III.3.1.3}, nos reducimos a probar que, cuando $X$ es \emph{irreducible} de punto genérico $x$, existe \emph{un} $\sh{O}_X$-módulo coherente $\sh{F}$ perteneciente a $\C'$ tal que $\sh{F}_x\neq 0$: en efecto, una vez establecido este punto, puede aplicarse a cualquier subpreesquema cerrado irreducible $Y$ de $X$, pues si $j:Y\to X$ es la inmersión canónica, $f\circ j$ es propio~\sref[II]{II.5.4.2}, y si $\sh{G}$ es un $\sh{O}_Y$-módulo coherente de soporte $Y$, entonces $j_*(\sh{G})$ es un $\sh{O}_X$-módulo coherente tal que $\RR^q(f\circ j)_*(\sh{G})=\RR^q f_*(j_*(\sh{G}))$~(G,~II,~4.9.1); podemos, por tanto, aplicar el Corolario~\sref{III.3.1.3}.

En virtud del lema de Chow~\sref[II]{II.5.6.2}, existen un preesquema irreducible $X'$ y un morfismo \emph{proyectivo} y sobreyectivo $g:X'\to X$ tales que $f\circ g:X'\to Y$ es \emph{proyectivo}.
Existe un $\sh{O}_{X'}$-módulo $\sh{L}$ amplio para $g$~\sref[II]{II.5.3.1}; aplicamos el teorema fundamental de los morfismos proyectivos~\sref{III.2.2.1} a $g:X'\to X$ y a $\sh{L}$: existe entonces un entero $n$ tal que $\sh{F}=g_*(\sh{O}_{X'}(n))$ es un $\sh{O}_X$-módulo coherente y $\RR^q g_*(\sh{O}_{X'}(n))=0$ para todo $q>0$; además, como $g^*(g_*(\sh{O}_{X'}(n)))\to\sh{O}_{X'}(n)$ es sobreyectivo para $n$ suficientemente grande~\sref{III.2.2.1}, vemos que podemos suponer, en el punto genérico $x$ de $X$, que $\sh{F}_x\neq 0$~\sref[II]{II.3.4.7}.
Por otra parte, como $f\circ g$ es proyectivo e $Y$ es noetheriano, los $\RR^q(f\circ g)_*(\sh{O}_{X'}(n))$ son \emph{coherentes}~\sref{III.2.2.1}.
Ahora bien, $\RR^\bullet(f\circ g)_*(\sh{O}_{X'}(n))$ es la convergencia de una sucesión espectral de Leray cuyo término $E_2$ viene dado por $E_2^{pq}=\RR^p f_*(\RR^q g_*(\sh{O}_{X'}(n)))$; lo anterior muestra que esta sucesión espectral degenera, y sabemos entonces~\sref[0]{0.11.1.6} que $E_2^{p0}=\RR^p f_*(\sh{F})$ es isomorfo a $\RR^p(f\circ g)_*(\sh{O}_{X'}(n))$, lo que termina la demostración.
\end{proof}

\begin{corollary}[3.2.2]
\label{III.3.2.2}
Sea $Y$ un preesquema localmente noetheriano.
Para todo morfismo propio $f:X\to Y$, la imagen directa por $f$ de cualquier $\sh{O}_X$-módulo coherente es un $\sh{O}_Y$-módulo coherente.
\end{corollary}

\begin{corollary}[3.2.3]
\label{III.3.2.3}
Sea $A$ un anillo noetheriano y sea $X$ un esquema propio sobre $A$; para todo $\sh{O}_X$-módulo coherente $\sh{F}$, los $\HH^p(X,\sh{F})$ son $A$-módulos de tipo finito, y existe un entero $r>0$ tal que, para todo $\sh{O}_X$-módulo coherente $\sh{F}$ y todo $p>r$, $\HH^p(X,\sh{F})=0$.
\end{corollary}

\begin{proof}
La segunda afirmación ya se demostró en~\sref{III.1.4.12}; la primera se deduce del teorema de finitud~\sref{III.3.2.1}, teniendo en cuenta el Corolario~\sref{III.1.4.11}.
\end{proof}

En particular, si $X$ es un \emph{esquema algebraico propio} sobre un cuerpo $k$, entonces, para todo $\sh{O}_X$-módulo coherente $\sh{F}$, los $\HH^p(X,\sh{F})$ son espacios vectoriales sobre $k$ de \emph{dimensión finita}.

\begin{corollary}[3.2.4]
\label{III.3.2.4}
Sea $Y$ un preesquema localmente noetheriano y sea $f:X\to Y$ un morfismo de tipo finito.
Para todo $\sh{O}_X$-módulo coherente $\sh{F}$ cuyo soporte sea propio sobre $Y$~\sref[II]{II.5.4.10}, los $\sh{O}_Y$-módulos $\RR^q f_*(\sh{F})$ son coherentes.
\end{corollary}

\oldpage[III]{118}
\begin{proof}
Como la cuestión es local sobre $Y$, podemos suponer que $Y$ es noetheriano, y lo mismo ocurre entonces con $X$~\sref[I]{I.6.3.7}.
Por hipótesis, todo subpreesquema cerrado $Z$ de $X$ cuyo espacio subyacente sea $\Supp(\sh{F})$ es propio sobre $Y$; en otras palabras, si $j:Z\to X$ es la inmersión canónica, entonces $f\circ j:Z\to Y$ es propio.
Podemos suponer que $Z$ es tal que $\sh{F}=j_*(\sh{G})$, donde $\sh{G}=j^*(\sh{F})$ es un $\sh{O}_Z$-módulo coherente~\sref[I]{I.9.3.5}; como $\RR^q f_*(\sh{F})=\RR^q(f\circ j)_*(\sh{G})$ por el Corolario~\sref{III.1.3.4}, la conclusión se deduce inmediatamente del Teorema~\sref{III.3.2.1}.
\end{proof}

\subsection{Generalización del teorema de finitud (esquemas ordinarios)}
\label{subsection:III.3.3}

\begin{proposition}[3.3.1]
\label{III.3.3.1}
Sea $Y$ un preesquema noetheriano, sea $\sh{S}$ una $\sh{O}_Y$-álgebra cuasi-coherente de tipo finito, graduada en grados positivos, sea $Y'=\Proj(\sh{S})$, y sea $g:Y'\to Y$ el morfismo estructural.
Sea $f:X\to Y$ un morfismo propio, sea $\sh{S}'=f^*(\sh{S})$, y sea $\sh{M}=\bigoplus_{k\in\bb{Z}}\sh{M}_k$ un $\sh{S}'$-módulo graduado cuasi-coherente de tipo finito.
Entonces los $\RR^p f_*(\sh{M})=\bigoplus_{k\in\bb{Z}}\RR^p f_*(\sh{M}_k)$ son $\sh{S}$-módulos graduados de tipo finito para todo $p$.
Supongamos además que $\sh{S}$ está generada por $\sh{S}_1$; entonces, para todo $p\in\bb{Z}$, existe un entero $k_p$ tal que, para todo $k\geq k_p$ y todo $r\geq 0$, se tiene
\[
\label{III.3.3.1.1}
  \RR^p f_*(\sh{M}_{k+r})=\sh{S}_r\RR^p f_*(\sh{M}_k).
  \tag{3.3.1.1}
\] 
\end{proposition}

\begin{proof}
La primera afirmación es idéntica al enunciado del Teorema~\sref{III.2.4.1}[i], en el que simplemente se sustituye «morfismo proyectivo» por «morfismo propio».
En la demostración del Teorema~\sref{III.2.4.1}[i], la hipótesis sobre $f$ solo se utilizó para demostrar (con la notación de aquella demostración) que $\RR^p f_*'(\widetilde{\sh{M}})$ es un $\sh{O}_{Y'}$-módulo coherente.
Bajo las hipótesis de la Proposición~\sref{III.3.3.1}, $f'$ es propio~\sref[II]{II.5.4.2}[iii], de modo que podemos retomar sin modificación la demostración del Teorema~\sref{III.2.4.1}[i], gracias al teorema de finitud~\sref{III.3.2.1}.

Para la segunda afirmación, basta observar que existe un recubrimiento abierto afín finito $(U_i)$ de $Y$ tal que las restricciones a los $U_i$ de los dos miembros de~(\hyperref[III.3.3.1.1]{3.3.1.1}) son iguales para todo $k\geq k_{p,i}$~\sref[II]{II.2.1.6}[ii]; basta tomar como $k_p$ el mayor de los $k_{p,i}$.
\end{proof}

\begin{corollary}[3.3.2]
\label{III.3.3.2}
Sea $A$ un anillo noetheriano, sea $\mathfrak{m}$ un ideal de $A$, sea $X$ un $A$-esquema propio y sea $\sh{F}$ un $\sh{O}_X$-módulo coherente.
Entonces, para todo $p\geq 0$, la suma directa $\bigoplus_{k\geq 0}\HH^p(X,\mathfrak{m}^k\sh{F})$ es un módulo de tipo finito sobre el anillo $S=\bigoplus_{k\geq 0}\mathfrak{m}^k$; en particular, existe un entero $k_p\geq 0$ tal que, para todo $k\geq k_p$ y todo $r\geq 0$, se tiene
\[
\label{III.3.3.2.1}
  \HH^p(X,\mathfrak{m}^{k+r}\sh{F})=\mathfrak{m}^r\HH^p(X,\mathfrak{m}^k\sh{F}).
  \tag{3.3.2.1}
\]
\end{corollary}

\begin{proof}
Basta aplicar la Proposición~\sref{III.3.3.1} con $Y=\Spec(A)$, $\sh{S}=\widetilde{S}$ y $\sh{M}_k=\mathfrak{m}^k\sh{F}$, teniendo en cuenta el Corolario~\sref{III.1.4.11}.
\end{proof}

Recordemos que la estructura de $S$-módulo sobre $\bigoplus_{k\geq 0}\HH^p(X,\mathfrak{m}^k\sh{F})$ se obtiene considerando, para todo $a\in\mathfrak{m}^r$, la aplicación $\HH^p(X,\mathfrak{m}^k\sh{F})\to\HH^p(X,\mathfrak{m}^{k+r}\sh{F})$ que procede del paso a cohomología de la aplicación de multiplicación $\mathfrak{m}^k\sh{F}\to\mathfrak{m}^{k+r}\sh{F}$ definida por $a$~\sref{III.2.4.1}.

\subsection{Teorema de finitud: caso de los esquemas formales}
\label{subsection:III.3.4}

\oldpage[III]{119}
Los resultados de esta sección (salvo la definición~\sref{III.3.4.1}) no se utilizarán en el resto de este capítulo.

\begin{env}[3.4.1]
\label{III.3.4.1}
Sean $\mathfrak{X}$ y $\mathfrak{S}$ dos preesquemas formales localmente noetherianos~\sref[I]{I.10.4.2}, y sea $f:\mathfrak{X}\to\mathfrak{S}$ un morfismo de preesquemas formales.
Diremos que $f$ es un morfismo \emph{propio} si satisface las condiciones siguientes:
\begin{enumerate}
  \item[1.º] \emph{$f$ es un morfismo de tipo finito~\sref[I]{I.10.13.3}}.
  \item[2.º] \emph{Si $\sh{K}$ es un haz de ideales de definición de $\mathfrak{S}$ y se pone $\sh{J}=f^*(\sh{K})\sh{O}_\mathfrak{X}$, $X_0=(\mathfrak{X},\sh{O}_\mathfrak{X}/\sh{J})$, $S_0=(\mathfrak{S},\sh{O}_\mathfrak{S}/\sh{K})$, entonces el morfismo $f_0:X_0\to S_0$ inducido por $f$~\sref[I]{I.10.5.6} es propio}.
\end{enumerate}
Es inmediato que esta definición no depende del haz de ideales de definición $\sh{K}$ de $\mathfrak{S}$ considerado; en efecto, si $\sh{K}'$ es un segundo haz de ideales de definición tal que $\sh{K}'\subset\sh{K}$, y si se pone $\sh{J}'=f^*(\sh{K}')\sh{O}_\mathfrak{X}$, $X_0'=(\mathfrak{X},\sh{O}_\mathfrak{X}/\sh{J}')$, $S_0'=(\mathfrak{S},\sh{O}_\mathfrak{S}/\sh{K}')$, entonces el morfismo $f_0':X_0'\to S_0'$ inducido por $f$ es tal que el diagrama
\[
  \xymatrix{
    X_0\ar[r]^{f_0}\ar[d]_i &
    S_0\ar[d]^j\\
    X_0'\ar[r]_{f_0'} &
    S_0'
  }
\]
es conmutativo, siendo $i$ y $j$ inmersiones sobreyectivas; en virtud de~\sref[II]{II.5.4.5}, es equivalente decir que $f_0$ o $f_0'$ es propio.

Observemos que, para todo $n\geq 0$, si se pone $X_n=(\mathfrak{X},\sh{O}_\mathfrak{X}/\sh{J}^{n+1})$, $S_n=(\mathfrak{S},\sh{O}_\mathfrak{S}/\sh{K}^{n+1})$, entonces el morfismo $f_n:X_n\to S_n$ inducido por $f$~\sref[I]{I.10.5.6} es propio para todo $n$ en cuanto lo es para $n=0$~\sref[II]{II.5.4.6}.

Si $g:Y\to Z$ es un morfismo propio de preesquemas ordinarios localmente noetherianos, $Z'$ es un subconjunto cerrado de $Z$, e $Y'$ es un subconjunto cerrado de $Y$ tal que $g(Y')\subset Z'$, entonces la extensión $\widehat{g}:Y_{/Y'}\to Z_{/Z'}$ de $g$ a las compleciones~\sref[I]{I.10.9.1} es un morfismo propio de preesquemas formales, como se deduce de la definición y de~\sref[II]{II.5.4.5}.

Sean $\mathfrak{X}$ y $\mathfrak{S}$ dos preesquemas formales localmente noetherianos y sea $f:\mathfrak{X}\to\mathfrak{S}$ un morfismo \emph{de tipo finito}~\sref[I]{I.10.13.3}; conservando la notación anterior, diremos que un subconjunto $Z$ del espacio subyacente a $\mathfrak{X}$ es \emph{propio} sobre $\mathfrak{S}$ (o propio para $f$) si, considerado como subconjunto de $X_0$, $Z$ es \emph{propio sobre $S_0$}~\sref[II]{II.5.4.10}.
Todas las propiedades de los subconjuntos propios de preesquemas ordinarios enunciadas en~\sref[II]{II.5.4.10} siguen siendo válidas para los subconjuntos propios de preesquemas formales, como se deduce inmediatamente de las definiciones.
\end{env}

\begin{theorem}[3.4.2]
\label{III.3.4.2}
Sean $\mathfrak{X}$ y $\mathfrak{Y}$ preesquemas formales localmente noetherianos, y sea $f:\mathfrak{X}\to\mathfrak{Y}$ un morfismo propio.
Para todo $\sh{O}_\mathfrak{X}$-módulo coherente $\sh{F}$, los $\sh{O}_\mathfrak{Y}$-módulos $\RR^q f_*(\sh{F})$ son coherentes para todo $q\geq 0$.
\end{theorem}

Sea $\sh{J}$ un haz de ideales de definición de $\mathfrak{Y}$, sea $\sh{K}=f^*(\sh{J})\sh{O}_\mathfrak{X}$, y consideremos los $\sh{O}_\mathfrak{X}$-módulos
\[
\label{III.3.4.2.1}
  \sh{F}_k=\sh{F}\otimes_{\sh{O}_\mathfrak{Y}}(\sh{O}_\mathfrak{Y}/\sh{J}^{k+1})=\sh{F}/\sh{K}^{k+1}\sh{F}\quad(k\geq 0)
  \tag{3.4.2.1}
\]
que forman evidentemente un \emph{sistema proyectivo} de $\sh{O}_\mathfrak{X}$-módulos topológicos, tal que
\oldpage[III]{120}
\[
  \sh{F}=\varprojlim_k\sh{F}_k
\]
por \sref[I]{I.10.11.3}.
Por otra parte, del Teorema~\sref{III.3.4.2} se deduce que cada $\RR^q f_*(\sh{F})$, por ser coherente, está dotado naturalmente de una estructura de $\sh{O}_\mathfrak{Y}$-módulo topológico~\sref[I]{I.10.11.6}, y lo mismo ocurre con los $\RR^q f_*(\sh{F}_k)$.
A los homomorfismos canónicos $\sh{F}\to\sh{F}_k=\sh{F}/\sh{K}^{k+1}\sh{F}$ corresponden canónicamente homomorfismos
\[
  \RR^q f_*(\sh{F})\to\RR^q f_*(\sh{F}_k),
\]
que son necesariamente continuos para las estructuras de $\sh{O}_\mathfrak{Y}$-módulo topológico anteriores~\sref[I]{I.10.11.6} y forman un sistema proyectivo, dando en el límite un homomorfismo funtorial canónico
\[
\label{III.3.4.2.2}
  \RR^q f_*(\sh{F})\to\varprojlim_k\RR^q f_*(\sh{F}_k),
  \tag{3.4.2.2}
\]
que es un homomorfismo continuo de $\sh{O}_\mathfrak{Y}$-módulos topológicos.
Junto con el Teorema~\sref{III.3.4.2}, demostraremos el
\begin{corollary}[3.4.3]
\label{III.3.4.3}
Cada uno de los homomorfismos~(\hyperref[III.3.4.2.2]{3.4.2.2}) es un isomorfismo topológico.
Además, si $\mathfrak{Y}$ es noetheriano, el sistema proyectivo $(\RR^q f_*(\sh{F}/\sh{K}^{k+1}\sh{F}))_{k\geq 0}$ satisface la condición \emph{(ML)}~\sref[0]{0.13.1.1}.
\end{corollary}
Comenzaremos estableciendo el Teorema~\sref{III.3.4.2} y el Corolario~\sref{III.3.4.3} cuando $\mathfrak{Y}$ es un esquema formal afín noetheriano~\sref[I]{I.10.4.1}:
\begin{corollary}[3.4.4]
\label{III.3.4.4}
Bajo las hipótesis del Teorema~\sref{III.3.4.2}, supongamos además que $\mathfrak{Y}=\Spf(A)$, donde $A$ es un anillo ádico noetheriano.
Sea $\mathfrak{J}$ un ideal de definición de $A$, y pongamos $\sh{F}_k=\sh{F}/\mathfrak{J}^{k+1}\sh{F}$ para $k\geq 0$.
Entonces los $\HH^n(\mathfrak{X},\sh{F})$ son $A$-módulos de tipo finito; el sistema proyectivo $(\HH^n(\mathfrak{X},\sh{F}_k))_{k\geq 0}$ satisface la condición \emph{(ML)} para todo $n$; si se pone
\[
\label{III.3.4.4.1}
  N_{n,k}=\Ker\big(\HH^n(\mathfrak{X},\sh{F})\to\HH^n(\mathfrak{X},\sh{F}_k)\big)
  \tag{3.4.4.1}
\]
(igual también a $\Im(\HH^n(\mathfrak{X},\mathfrak{J}^{k+1}\sh{F})\to\HH^n(\mathfrak{X},\sh{F}))$ por la sucesión exacta de cohomología), entonces los $N_{n,k}$ definen sobre $\HH^n(\mathfrak{X},\sh{F})$ una filtración $\mathfrak{J}$-buena~\sref[0]{0.13.7.7}; finalmente, el homomorfismo canónico
\[
\label{III.3.4.4.2}
  \HH^n(\mathfrak{X},\sh{F})\to\varprojlim_k\HH^n(\mathfrak{X},\sh{F}_k)
  \tag{3.4.4.2}
\]
es un isomorfismo topológico para todo $n$ (el miembro izquierdo está dotado de la topología $\mathfrak{J}$-ádica y los $\HH^n(\mathfrak{X},\sh{F}_k)$ de la topología discreta).
\end{corollary}

Pongamos
\[
\label{III.3.4.4.3}
  S=\gr(A)=\bigoplus_{k\geq 0}\mathfrak{J}^k/\mathfrak{J}^{k+1},\ \sh{M}=\gr(\sh{F})=\bigoplus_{k\geq 0}\mathfrak{J}^k\sh{F}/\mathfrak{J}^{k+1}\sh{F}. \tag{3.4.4.3}
\]
Sabemos que $\mathfrak{J}^\Delta$ es un haz de ideales de definición de $\mathfrak{Y}$~\sref[I]{I.10.3.1}; sean $\sh{K}=f^*(\mathfrak{J}^\Delta)\sh{O}_\mathfrak{X}$, $X_0=(\mathfrak{X},\sh{O}_\mathfrak{X}/\sh{K})$, $Y_0=(\mathfrak{Y},\sh{O}_\mathfrak{Y}/\mathfrak{J}^\Delta)=\Spec(A_0)$, con $A_0=A/\mathfrak{J}$.
Está claro que los $\sh{M}_k=\mathfrak{J}^k\sh{F}/\mathfrak{J}^{k+1}\sh{F}$ son $\sh{O}_{X_0}$-módulos coherentes~\sref[I]{I.10.11.3}.
Consideremos, por otra parte, la $\sh{O}_{X_0}$-álgebra graduada cuasi-coherente
\[
\label{III.3.4.4.4}
  \sh{S}=\sh{O}_{X_0}\otimes_{A_0}S.
  \tag{3.4.4.4}
\]
La multiplicación define un homomorfismo sobreyectivo de álgebras graduadas
\[
  \sh{S}\longrightarrow\gr(\sh{O}_\mathfrak{X})
  =\bigoplus_{k\geq 0}\sh{K}^k/\sh{K}^{k+1},
\]
por medio del cual $\sh{M}$ es un $\sh{S}$-módulo graduado.

\begingroup
\microtypesetup{protrusion=false,expansion=false}
La hipótesis de que $\sh{F}$ es un $\sh{O}_\mathfrak{X}$-módulo de tipo finito implica primero que $\sh{M}$ es
\oldpage[III]{121}
un $\sh{S}$-módulo graduado \emph{de tipo finito}.
En efecto, la cuestión es local sobre $\mathfrak{X}$, y podemos suponer que $\mathfrak{X}=\Spf(B)$, donde $B$ es un anillo ádico noetheriano, y que $\sh{F}=N^\Delta$, donde $N$ es un $B$-módulo de tipo finito~\sref[I]{I.10.10.5}.
Pongamos $B_0=B/\mathfrak{J}B$ y $N_0=N/\mathfrak{J}N$.
Entonces $X_0=\Spec(B_0)$, y los $\sh{O}_{X_0}$-módulos cuasi-coherentes $\sh{S}$ y $\sh{M}$ están asociados respectivamente a
\[
  S'=\bigoplus_{k\geq 0}((\mathfrak{J}^k/\mathfrak{J}^{k+1})\otimes_{A_0}B_0),\qquad
  M'=\bigoplus_{k\geq 0}\mathfrak{J}^kN/\mathfrak{J}^{k+1}N.
\]
Para todo $k\geq 0$, la multiplicación induce una aplicación sobreyectiva
\[
  (\mathfrak{J}^k/\mathfrak{J}^{k+1})\otimes_{A_0}N_0
  \longrightarrow\mathfrak{J}^kN/\mathfrak{J}^{k+1}N,
  \qquad \overline a\otimes\overline n\longmapsto\overline{an}.
\]
Estas aplicaciones definen un homomorfismo sobreyectivo de $S'$-módulos graduados $S'\otimes_{B_0}N_0\longrightarrow M'$.
Como $N_0$ es un $B_0$-módulo de tipo finito, $M'$ es, por tanto, un $S'$-módulo de tipo finito, lo que prueba nuestra afirmación~\sref[I]{I.1.3.13}.
\par\endgroup

Como el morfismo $f_0:X_0\to Y_0$ es \emph{propio} por hipótesis, podemos aplicar la Proposición~\sref{III.3.3.1} a $\sh{S}$, $\sh{M}$ y al morfismo $f_0$: teniendo en cuenta el Corolario~\sref{III.1.4.11}, concluimos que, para \emph{todo $n\geq 0$}, $\bigoplus_{k\geq 0}\HH^n(X_0,\sh{M}_k)$ es un $S$-módulo graduado \emph{de tipo finito}.
Esto prueba que la condición ($\text{F}_n$) de~\sref[0]{0.13.7.7} se satisface para \emph{todo $n\geq 0$}, cuando se considera el sistema proyectivo estricto $(\sh{F}/\mathfrak{J}^k\sh{F})_{k\geq 0}$ de haces de grupos abelianos sobre $X_0$, cada uno dotado de su estructura natural de «$A$-módulo filtrado».
Podemos, pues, aplicar~\sref[0]{0.13.7.7}, que demuestra que:
\begin{enumerate}
  \item[1.º] El sistema proyectivo $(\HH^n(\mathfrak{X},\sh{F}_k))_{k\geq 0}$ satisface la condición (ML).
  \item[2.º] Si $\HH^{\prime n}=\varprojlim_k\HH^n(\mathfrak{X},\sh{F}_k)$, entonces $\HH^{\prime n}$ es un $A$-módulo de tipo finito.
  \item[3.º] La filtración definida sobre $\HH^{\prime n}$ por los núcleos de los homomorfismos canónicos $\HH^{\prime n}\to\HH^n(\mathfrak{X},\sh{F}_k)$ es $\mathfrak{J}$-buena.
\end{enumerate}

Observemos, por otra parte, que si se pone $X_k=(\mathfrak{X},\sh{O}_\mathfrak{X}/\sh{K}^{k+1})$, entonces $\sh{F}_k$ es un $\sh{O}_{X_k}$-módulo coherente~\sref[I]{I.10.11.3}; y si $U$ es un abierto afín de $X_0$, entonces $U$ es también un abierto afín de cada $X_k$~\sref[I]{I.5.1.9}, de modo que $\HH^n(U,\sh{F}_k)=0$ para todo $n>0$ y todo $k$~\sref{III.1.3.1}, y $\HH^0(U,\sh{F}_k)\to\HH^0(U,\sh{F}_h)$ es sobreyectivo para $h\leq k$~\sref[I]{I.1.3.9}.
Estamos, pues, en las condiciones de~\sref[0]{0.13.3.2}, y la aplicación de~\sref[0]{0.13.3.1} demuestra que $\HH^{\prime n}$ se identifica canónicamente con $\HH^n(\mathfrak{X},\varprojlim_k\sh{F}_k)=\HH^n(\mathfrak{X},\sh{F})$; esto termina la demostración del Corolario~\sref{III.3.4.4}.

\begin{env}[3.4.5]
\label{III.3.4.5}
Volvamos a la demostración de~\sref{III.3.4.2} y~\sref{III.3.4.3}.
Demostremos primero las proposiciones en el caso $\mathfrak{Y}=\Spf(A)$ considerado en~\sref{III.3.4.4}; para ello, para todo $g\in A$, apliquemos~\sref{III.3.4.4} al esquema formal afín noetheriano inducido sobre el abierto $\mathfrak{Y}_g=\mathfrak{D}(g)$ de $\mathfrak{Y}$, que es igual a $\Spf(A_{\{g\}})$, y al preesquema formal inducido por $\mathfrak{X}$ sobre $f^{-1}(\mathfrak{Y}_g)$; observemos que $\mathfrak{Y}_g$ es también un abierto afín del preesquema $Y_k=(\mathfrak{Y},\sh{O}_\mathfrak{Y}/(\mathfrak{J}^\Delta)^{k+1})$, y como $\sh{F}_k$ es un $\sh{O}_{X_k}$-módulo coherente, se tiene
\[
  \HH^n(f^{-1}(\mathfrak{Y}_g),\sh{F}_k)=\Gamma(\mathfrak{Y}_g,\RR^n f_*(\sh{F}_k))
\]
para todo $k\geq 0$, en virtud del Corolario~\sref{III.1.4.11}.
El homomorfismo canónico
\[
  \HH^n(f^{-1}(\mathfrak{Y}_g),\sh{F})\to\varprojlim_k\Gamma(\mathfrak{Y}_g,\RR^n f_*(\sh{F}_k))
\]
es un isomorfismo; pero se tiene~\sref[0]{0.3.2.6}
\[
  \varprojlim_k\Gamma(\mathfrak{Y}_g,\RR^n f_*(\sh{F}_k))=\Gamma(\mathfrak{Y}_g,\varprojlim_k\RR^n f_*(\sh{F}_k)),
\]
\oldpage[III]{122}
y como el haz $\RR^n f_*(\sh{F})$ es el haz asociado al prehaz $\mathfrak{Y}_g\mapsto\HH^n(f^{-1}(\mathfrak{Y}_g),\sh{F})$ sobre los $\mathfrak{Y}_g$~\sref[0]{0.3.2.1}, hemos demostrado que el homomorfismo~(\hyperref[III.3.4.2.2]{3.4.2.2}) es \emph{biyectivo}.
Demostremos ahora que $\RR^n f_*(\sh{F})$ es un $\sh{O}_\mathfrak{Y}$-módulo coherente y, más precisamente, que se tiene
\[
\label{III.3.4.5.1}
  \RR^n f_*(\sh{F})=\big(\HH^n(\mathfrak{X},\sh{F})\big)^\Delta.
  \tag{3.4.5.1}
\]

Con la notación anterior, como $\sh{F}_k$ es un $\sh{O}_{X_k}$-módulo coherente~\sref{III.1.4.13}, se tiene
\[
  \Gamma(\mathfrak{Y}_g,\RR^n f_*(\sh{F}_k))=(\Gamma(\mathfrak{Y},\RR^n f_*(\sh{F}_k)))_g=(\HH^n(\mathfrak{X},\sh{F}_k))_g.
\]

Ahora bien, los $\HH^n(\mathfrak{X},\sh{F}_k)$ forman un sistema proyectivo que satisface (ML), y su límite proyectivo $\HH^n(\mathfrak{X},\sh{F})$ es un $A$-módulo de tipo finito.
Concluimos~\sref[0]{0.13.7.8} que se tiene
\[
  \varprojlim_k\big((\HH^n(\mathfrak{X},\sh{F}_k))_g\big)=\HH^n(\mathfrak{X},\sh{F})\otimes_A A_{\{g\}}=\Gamma(\mathfrak{Y}_g,(\HH^n(\mathfrak{X},\sh{F}))^\Delta),
\]
teniendo en cuenta~\sref[I]{I.10.10.8} aplicado a $A$ y $A_{\{g\}}$; esto demuestra~(\hyperref[III.3.4.5.1]{3.4.5.1}), pues $\Gamma(\mathfrak{Y}_g,\RR^n f_*(\sh{F}))=\varprojlim_k\Gamma(\mathfrak{Y}_g,\RR^n f_*(\sh{F}_k))$.

Como (\hyperref[III.3.4.2.2]{3.4.2.2}) es entonces un isomorfismo de $\sh{O}_\mathfrak{Y}$-módulos coherentes, es necesariamente un isomorfismo \emph{topológico}~\sref[I]{I.10.11.6}.
Finalmente, de las relaciones $\RR^n f_*(\sh{F}_k)=(\HH^n(\mathfrak{X},\sh{F}_k))^\Delta$ se deduce que el sistema proyectivo $(\RR^n f_*(\sh{F}_k))_{k\geq 0}$ satisface (ML)~\sref[I]{I.10.10.2}.

Una vez demostrados \sref{III.3.4.2} y \sref{III.3.4.3} en el caso en que el preesquema formal $\mathfrak{Y}$ es afín noetheriano, el paso al caso general es inmediato para~\sref{III.3.4.2} y para la primera afirmación de~\sref{III.3.4.3}, que son locales sobre $\mathfrak{Y}$.
En cuanto a la segunda afirmación de~\sref{III.3.4.3}, como $\mathfrak{Y}$ es noetheriano, basta recubrirlo por un número finito de abiertos afines noetherianos $U_i$ y observar que las restricciones del sistema proyectivo $(\RR^q f_*(\sh{F}_k))$ a cada uno de los $U_i$ satisfacen (ML).
\end{env}

En el curso de la demostración hemos probado además:
\begin{corollary}[3.4.6]
\label{III.3.4.6}
Bajo las hipótesis del Corolario~\sref{III.3.4.4}, el homomorfismo canónico
\[
\label{III.3.4.6.1}
  \HH^q(\mathfrak{X},\sh{F})\to\Gamma(\mathfrak{Y},\RR^q f_*(\sh{F}))
  \tag{3.4.6.1}
\]
es biyectivo.
\end{corollary}

\section{El teorema fundamental de los morfismos propios. Aplicaciones}
\label{section:III.4}


\subsection{El teorema fundamental}
\label{subsection:III.4.1}

\begin{env}[4.1.1]
\label{III.4.1.1}
Sean $X$ e $Y$ esquemas ordinarios noetherianos, sea $f:X\to Y$ un
morfismo propio, sea $Y'$ un subconjunto cerrado de $Y$, y pongamos
\[
  X'=f^{-1}(Y').
\]
Denotamos por $\widehat{X}$ y $\widehat{Y}$ los esquemas formales
$X_{/X'}$ y $Y_{/Y'}$, obtenidos completando $X$ e $Y$ a lo largo de
$X'$ e $Y'$, respectivamente \sref[I]{I.10.8.5}.
Denotamos por $\widehat{f}:\widehat{X}\to\widehat{Y}$ la extensión de
$f$ a estas completaciones \sref[I]{I.10.9.1}.
Para todo $\sh{O}_X$-módulo coherente $\sh{F}$, denotamos por

\oldpage[III]{123}

\[
  \widehat{\sh{F}}=\sh{F}_{/X'}
\]
su completación a lo largo de $X'$ \sref[I]{I.10.8.4}; este es un
$\sh{O}_{\widehat{X}}$-módulo coherente \sref[I]{I.10.8.8}.
\end{env}

\begin{env}[4.1.2]
\label{III.4.1.2}
Sea $\sh{J}$ un ideal coherente de $\sh{O}_Y$ tal que
\[
  \Supp(\sh{O}_Y/\sh{J})=Y'
\]
\sref[I]{I.5.2.1}.
Sabemos \sref[I]{I.4.4.5} que
\[
  \sh{K}=f^*(\sh{J})\sh{O}_X
\]
es un ideal coherente de $\sh{O}_X$ que satisface
\[
  \Supp(\sh{O}_X/\sh{K})=X'.
\]
Para todo $k\geq0$, consideremos el $\sh{O}_X$-módulo coherente
\[
  \sh{F}_k
  =
  \sh{F}\otimes_{\sh{O}_Y}(\sh{O}_Y/\sh{J}^{k+1})
  =
  \sh{F}/\sh{K}^{k+1}\sh{F}.
\]
Los $\sh{O}_Y$-módulos
\[
  \RR^n f_*(\sh{F})
  \quad\text{y}\quad
  \RR^n f_*(\sh{F}_k)
\]
son coherentes para todo $n$ \sref{III.3.2.1}.
Para todo $k\geq0$ y todo $n$, el homomorfismo canónico
$\sh{F}\to\sh{F}_k$ induce, por funtorialidad, un homomorfismo
\[
  \RR^n f_*(\sh{F})
  \longrightarrow
  \RR^n f_*(\sh{F}_k).
  \tag{4.1.2.1}
  \label{III.4.1.2.1}
\]
Como $\sh{F}_k$ es un $\sh{O}_X/\sh{K}^{k+1}$-módulo,
$\RR^n f_*(\sh{F}_k)$ es un
$\sh{O}_Y/\sh{J}^{k+1}$-módulo \sref[0]{0.12.2.1}.
Deducimos, pues, de \eqref{III.4.1.2.1} un homomorfismo
\[
  \RR^n f_*(\sh{F})
  \otimes_{\sh{O}_Y}(\sh{O}_Y/\sh{J}^{k+1})
  \longrightarrow
  \RR^n f_*(\sh{F}_k).
  \tag{4.1.2.2}
  \label{III.4.1.2.2}
\]
Los dos miembros de \eqref{III.4.1.2.2} forman sistemas proyectivos.
El límite proyectivo del miembro izquierdo es la completación
\[
  \bigl(\RR^n f_*(\sh{F})\bigr)_{/Y'},
\]
que denotaremos por
\[
  \widehat{\RR^n f_*(\sh{F})}.
\]
Los homomorfismos \eqref{III.4.1.2.2} son compatibles con los sistemas
proyectivos y definen, por tanto, al pasar al límite, un homomorfismo
canónico
\[
  \vphi_n:
  \widehat{\RR^n f_*(\sh{F})}
  \longrightarrow
  \varprojlim_k\RR^n f_*(\sh{F}_k).
  \tag{4.1.2.3}
  \label{III.4.1.2.3}
\]
Además, \eqref{III.4.1.2.2} es un homomorfismo de
$\sh{O}_Y/\sh{J}^{k+1}$-módulos.
Por \sref[I]{I.10.8.3}, puede considerarse, por tanto, como un
homomorfismo continuo de $\sh{O}_{\widehat{Y}}$-módulos topológicos
provistos de sus topologías seudodiscretas \sref[0]{0.3.8.1}.
En consecuencia, $\vphi_n$ es un homomorfismo continuo de
$\sh{O}_{\widehat{Y}}$-módulos topológicos.
\end{env}

\begin{env}[4.1.3]
\label{III.4.1.3}
Sea $i:\widehat{X}\to X$ el morfismo canónico de espacios anillados
definido en \sref[I]{I.10.8.7}.
Tenemos el diagrama conmutativo
\[
  \xymatrix{
    X_k\ar[r]^{h_k}\ar[dr]_{i_k} &
    \widehat{X}\ar[d]^{i}\\
    & X
  }
  \tag{4.1.3.1}
  \label{III.4.1.3.1}
\]
donde $X_k$ es el subesquema cerrado de $X$ definido por
$\sh{K}^{k+1}$, $i_k$ es la inmersión canónica, y $h_k$ es el
morfismo de espacios anillados que corresponde a la identidad sobre
los espacios subyacentes y al homomorfismo canónico
\[
  \sh{O}_{\widehat{X}}
  \longrightarrow
  \sh{O}_X/\sh{K}^{k+1}
\]
\sref[I]{I.10.5.2}.
Además,
\[
  \widehat{\sh{F}}=i^*(\sh{F})
\]
salvo isomorfismo canónico \sref[I]{I.10.8.8}.
Tenemos
\[
  \HH^n(X_k,i_k^*(\sh{F}_k))
  =
  \HH^n(X,\sh{F}_k)
  \tag{4.1.3.2}
  \label{III.4.1.3.2}
\]
salvo isomorfismo canónico.

\oldpage[III]{124}

En efecto,
\[
  \sh{F}_k=(i_k)_*(i_k^*(\sh{F}_k))
\]
(G,~II,~4.9.1).
El homomorfismo canónico
\[
  \HH^n(\widehat{X},\widehat{\sh{F}})
  \longrightarrow
  \HH^n(X_k,h_k^*(\widehat{\sh{F}}))
\]
\sref[0]{0.12.1.3.5} puede escribirse también
\[
  \HH^n(\widehat{X},\widehat{\sh{F}})
  \longrightarrow
  \HH^n(X,\sh{F}_k).
  \tag{4.1.3.3}
  \label{III.4.1.3.3}
\]
Estos homomorfismos forman un sistema proyectivo y definen, por tanto,
un homomorfismo canónico
\[
  \psi_{n,X}:
  \HH^n(\widehat{X},\widehat{\sh{F}})
  \longrightarrow
  \varprojlim_k\HH^n(X,\sh{F}_k).
  \tag{4.1.3.4}
  \label{III.4.1.3.4}
\]
Reemplazando $X$ por un subconjunto abierto de la forma $f^{-1}(V)$,
donde $V$ es un abierto afín de $Y$, y teniendo en cuenta
\sref{III.1.4.11}, obtenemos homomorfismos
\[
  \psi_{n,V}:
  \HH^n(\widehat{X}\cap f^{-1}(V),\widehat{\sh{F}})
  \longrightarrow
  \varprojlim_k\Gamma(V,\RR^n f_*(\sh{F}_k)).
  \tag{4.1.3.5}
  \label{III.4.1.3.5}
\]
Estos homomorfismos conmutan manifiestamente con la restricción de $V$
a un subconjunto abierto afín más pequeño.
Definen, por tanto, un homomorfismo canónico de haces
\[
  \psi_n:
  \RR^n\widehat{f}_*(\widehat{\sh{F}})
  \longrightarrow
  \varprojlim_k\RR^n f_*(\sh{F}_k).
  \tag{4.1.3.6}
  \label{III.4.1.3.6}
\]
\end{env}

\begin{env}[4.1.4]
\label{III.4.1.4}
Por último, sea $j:\widehat{Y}\to Y$ el morfismo canónico de espacios
anillados \sref[I]{I.10.8.7}.
Como $\RR^n f_*(\sh{F})$ es un $\sh{O}_Y$-módulo coherente
\sref{III.3.2.1}, tenemos, salvo isomorfismo canónico,
\[
  j^*(\RR^n f_*(\sh{F}))
  =
  \widehat{\RR^n f_*(\sh{F})}
\]
\sref[I]{I.10.8.8}.
Existe, en consecuencia, un homomorfismo canónico
\[
  \rho_n:
  \widehat{\RR^n f_*(\sh{F})}
  =
  j^*(\RR^n f_*(\sh{F}))
  \longrightarrow
  \RR^n\widehat{f}_*(i^*(\sh{F}))
  =
  \RR^n\widehat{f}_*(\widehat{\sh{F}})
  \tag{4.1.4.1}
  \label{III.4.1.4.1}
\]
definido para los espacios anillados en general; véase la demostración
de \sref{III.1.4.15}.
Demostraremos que el diagrama
\[
  \xymatrix{
    \widehat{\RR^n f_*(\sh{F})}
      \ar[rr]^{\rho_n}\ar[dr]_{\vphi_n}
    &&
    \RR^n\widehat{f}_*(\widehat{\sh{F}})
      \ar[dl]^{\psi_n}\\
    &
    \displaystyle\varprojlim_k\RR^n f_*(\sh{F}_k)
    &
  }
  \tag{4.1.4.2}
  \label{III.4.1.4.2}
\]
es conmutativo.
Basta demostrar la conmutatividad del diagrama formado por los
homomorfismos correspondientes de prehaces.
Podemos, por tanto, restringirnos al caso en que $Y$ es afín, y basta
demostrar que
\[
  \xymatrix{
    \widehat{\HH^n(X,\sh{F})}
      \ar[rr]^{\rho_n}\ar[dr]_{\vphi_n}
    &&
    \HH^n(\widehat{X},\widehat{\sh{F}})
      \ar[dl]^{\psi_{n,X}}\\
    &
    \displaystyle\varprojlim_k\HH^n(X,\sh{F}_k)
    &
  }
  \tag{4.1.4.3}
  \label{III.4.1.4.3}
\]
es conmutativo.
La conmutatividad de \eqref{III.4.1.3.1}, junto con las relaciones
entre los grupos de cohomología establecidas en \sref{III.4.1.3}, da
el diagrama conmutativo

\oldpage[III]{125}

\[
  \xymatrix{
    \HH^n(X,\sh{F})
      \ar[r]\ar[dr]
    &
    \HH^n(\widehat{X},\widehat{\sh{F}})
      =
    \HH^n(\widehat{X},i^*(\sh{F}))
      \ar[d]\\
    &
    \HH^n(X,\sh{F}_k)
      =
    \HH^n(X_k,i_k^*(\sh{F})).
  }
\]
La conmutatividad de \eqref{III.4.1.4.3} se sigue inmediatamente.
\end{env}

\begin{theorem}[4.1.5]
\label{III.4.1.5}
Sea $f:X\to Y$ un morfismo propio de esquemas noetherianos, sea
$Y'$ un subconjunto cerrado de $Y$, y pongamos $X'=f^{-1}(Y')$.
Entonces, para todo $\sh{O}_X$-módulo coherente $\sh{F}$,
\[
  \RR^n\widehat{f}_*(\widehat{\sh{F}})
\]
es un $\sh{O}_{\widehat{Y}}$-módulo coherente, y los homomorfismos
$\vphi_n$, $\psi_n$ y $\rho_n$ del diagrama
\eqref{III.4.1.4.2} son isomorfismos topológicos.
\end{theorem}

\begin{proof}
Basta demostrar que $\vphi_n$ y $\psi_n$ son isomorfismos.
En efecto, $\RR^n f_*(\sh{F})$ es coherente \sref{III.3.2.1}, luego
su completación $\widehat{\RR^n f_*(\sh{F})}$ es coherente
\sref[I]{I.10.8.8}; la bicontinuidad de $\vphi_n$, $\psi_n$ y
$\rho_n$ es entonces automática \sref[I]{I.10.11.6}.
\end{proof}

\begin{env}[4.1.6]
\label{III.4.1.6}
\emph{Observaciones.}
\begin{enumerate}
  \item[(i)] Pongamos
  \[
    \widehat{\sh{F}}_k
    =
    \widehat{\sh{F}}/
    \widehat{\sh{K}}^{\,k+1}\widehat{\sh{F}}.
  \]
  Entonces, inmediatamente,
  \[
    \sh{F}_k=i_*(\widehat{\sh{F}}_k),
  \]
  y el homomorfismo canónico \eqref{III.4.1.3.6} es precisamente el
  homomorfismo ya definido en (\hyperref[III.3.4.2.2]{3.4.2.2}):
  \[
    \RR^n\widehat{f}_*(\widehat{\sh{F}})
    \longrightarrow
    \varprojlim_k
    \RR^n\widehat{f}_*(\widehat{\sh{F}}_k).
    \tag{4.1.6.1}
    \label{III.4.1.6.1}
  \]
  Así, el hecho de que $\psi_n$ sea un isomorfismo es un caso particular
  de \sref{III.3.4.3}.
  Daremos, no obstante, más adelante una demostración directa, evitando
  las delicadas consideraciones relativas a los límites proyectivos de
  sucesiones espectrales \sref[0]{0.13.7} en las que se apoya el teorema
  general \sref{III.3.4.3}.

  \item[(ii)] Como los $\psi_n$ son isomorfismos, decir que los
  $\vphi_n$ son isomorfismos equivale a decir que
  \[
    \rho_n=\psi_n^{-1}\circ\vphi_n
  \]
  son isomorfismos.
  El Teorema~\sref{III.4.1.5} afirma, en particular, que la formación de
  $\RR^n f_*$ conmuta con la completación; puede llamarse, por tanto, el
  primer teorema de comparación entre las teorías ``algebraica'' y
  ``formal''.
\end{enumerate}
Comenzaremos demostrando la forma afín de \sref{III.4.1.5}.
\end{env}

\begin{corollary}[4.1.7]
\label{III.4.1.7}
Bajo las hipótesis de \sref{III.4.1.5}, supongamos además que
$Y=\Spec(A)$, donde $A$ es noetheriano, y que
$\sh{J}=\widetilde{\mathfrak{J}}$, donde $\mathfrak{J}$ es un ideal
de $A$, de modo que
\[
  \sh{F}_k=\sh{F}/\mathfrak{J}^{k+1}\sh{F}.
\]
El homomorfismo canónico
\[
  \vphi_n:
  \widehat{\HH^n(X,\sh{F})}
  \longrightarrow
  \varprojlim_k\HH^n(X,\sh{F}_k)
  \tag{4.1.7.1}
  \label{III.4.1.7.1}
\]
donde el primer término es la completación separada de
$\HH^n(X,\sh{F})$ para la topología $\mathfrak{J}$-preádica, es un
isomorfismo.
Para todo $n$, el sistema proyectivo
\[
  \bigl(\HH^n(X,\sh{F}_k)\bigr)_{k\geq0}
\]
satisface la condición \emph{(ML)}, y el homomorfismo canónico
\[
  \psi_n:
  \HH^n(\widehat{X},\widehat{\sh{F}})
  \longrightarrow
  \varprojlim_k\HH^n(X,\sh{F}_k)
  \tag{4.1.7.2}
  \label{III.4.1.7.2}
\]

\oldpage[III]{126}

es un isomorfismo.
Por último, la filtración sobre $\HH^n(X,\sh{F})$ definida por los
núcleos de los homomorfismos canónicos
\[
  \HH^n(X,\sh{F})\longrightarrow\HH^n(X,\sh{F}_k)
\]
es $\mathfrak{J}$-buena \sref[0]{0.13.7.7}, y $\vphi_n$ es un
isomorfismo topológico.\footnote{La demostración siguiente, más sencilla
que la demostración original, y el complemento relativo a la filtración
de $\HH^n(X,\sh{F})$, nos fueron comunicados por J.-P. Serre.}
\end{corollary}

\begin{proof}
Fijemos $n\geq0$ durante toda esta demostración y, para simplificar,
pongamos
\[
  \mathrm{H}=\HH^n(X,\sh{F}),
  \qquad
  \mathrm{H}_k=\HH^n(X,\sh{F}_k).
  \tag{4.1.7.3}
  \label{III.4.1.7.3}
\]
Pongamos también
\[
  \mathrm{R}_k
  =
  \Ker(\mathrm{H}\longrightarrow\mathrm{H}_k),
  \qquad
  \text{un sub-$A$-módulo de $\mathrm{H}$.}
  \tag{4.1.7.4}
  \label{III.4.1.7.4}
\]
La sucesión exacta de cohomología
\[
  \HH^n(X,\mathfrak{J}^{k+1}\sh{F})
  \longrightarrow
  \HH^n(X,\sh{F})
  \longrightarrow
  \HH^n(X,\sh{F}_k)
  \xrightarrow{\partial}
  \HH^{n+1}(X,\mathfrak{J}^{k+1}\sh{F})
  \longrightarrow
  \HH^{n+1}(X,\sh{F})
\]
muestra que
\[
  \mathrm{R}_k
  =
  \Im\bigl(
    \HH^n(X,\mathfrak{J}^{k+1}\sh{F})
    \longrightarrow
    \HH^n(X,\sh{F})
  \bigr).
\]
Ponemos
\[
  \begin{split}
  \mathrm{Q}_k
  &=
  \Ker\bigl(
    \HH^{n+1}(X,\mathfrak{J}^{k+1}\sh{F})
    \longrightarrow
    \HH^{n+1}(X,\sh{F})
  \bigr)
  \\
  &=
  \Im\bigl(
    \HH^n(X,\sh{F}_k)
    \longrightarrow
    \HH^{n+1}(X,\mathfrak{J}^{k+1}\sh{F})
  \bigr).
  \end{split}
  \tag{4.1.7.5}
  \label{III.4.1.7.5}
\]
Tenemos, por tanto, la sucesión exacta
\[
  0\longrightarrow\mathrm{R}_k
  \longrightarrow\mathrm{H}
  \longrightarrow\mathrm{H}_k
  \longrightarrow\mathrm{Q}_k
  \longrightarrow0.
  \tag{4.1.7.6}
  \label{III.4.1.7.6}
\]

\begin{env}[4.1.7.7]
\label{III.4.1.7.7}
Sea $x\in\mathfrak{J}^m$, donde $m\geq0$.
La multiplicación por $x$ sobre $\mathfrak{J}^k\sh{F}$ es un
homomorfismo
\[
  \mathfrak{J}^k\sh{F}
  \longrightarrow
  \mathfrak{J}^{k+m}\sh{F},
\]
y da, por tanto, un homomorfismo
\[
  \mu_{x,m}:
  \HH^n(X,\mathfrak{J}^k\sh{F})
  \longrightarrow
  \HH^n(X,\mathfrak{J}^{k+m}\sh{F}).
  \tag{4.1.7.8}
  \label{III.4.1.7.8}
\]
Sea
\[
  S=\bigoplus_{k\geq0}\mathfrak{J}^k
\]
la $A$-álgebra graduada.
Sabemos que las multiplicaciones $\mu_{x,m}$ dotan a
\[
  \mathrm{E}
  =
  \bigoplus_{k\geq0}
  \HH^n(X,\mathfrak{J}^k\sh{F})
\]
de una estructura de $S$-módulo graduado de tipo finito
\sref{III.3.3.2}; el anillo graduado $S$ es noetheriano
\sref[II]{II.2.1.5}.
\end{env}

\begin{lemma}[4.1.7.9]
\label{III.4.1.7.9}
Los submódulos $\mathrm{R}_k$ de $\mathrm{H}$ definen una
filtración $\mathfrak{J}$-buena sobre $\mathrm{H}$.
\end{lemma}

\begin{proof}
Demostremos primero que
\[
  \mathfrak{J}^m\mathrm{R}_k
  \subset
  \mathrm{R}_{k+m},
  \tag{4.1.7.10}
  \label{III.4.1.7.10}
\]
donde la multiplicación sobre
$\mathrm{H}=\HH^n(X,\sh{F})$ por un elemento
$x\in\mathfrak{J}^m$ es la aplicación $\mu_{x,0}$.
Para todo $x\in\mathfrak{J}^m$, el diagrama
\[
  \xymatrix{
    \mathfrak{J}^{k+1}\sh{F}
      \ar[r]^{x}\ar[d]
    &
    \mathfrak{J}^{k+m+1}\sh{F}
      \ar[d]
    \\
    \sh{F}
      \ar[r]_{x}
    &
    \sh{F}
  }
  \tag{4.1.7.11}
  \label{III.4.1.7.11}
\]

\oldpage[III]{127}

es conmutativo; sus flechas horizontales son las multiplicaciones por
$x$, y sus flechas verticales son las inyecciones canónicas.
En consecuencia, el diagrama correspondiente
\[
  \xymatrix{
    \HH^n(X,\mathfrak{J}^{k+1}\sh{F})
      \ar[r]^{\mu_{x,m}}\ar[d]
    &
    \HH^n(X,\mathfrak{J}^{k+m+1}\sh{F})
      \ar[d]
    \\
    \HH^n(X,\sh{F})
      \ar[r]_{\mu_{x,0}}
    &
    \HH^n(X,\sh{F})
  }
\]
es conmutativo.
La interpretación de $\mathrm{R}_k$ como la imagen de
\[
  \HH^n(X,\mathfrak{J}^{k+1}\sh{F})
  \longrightarrow
  \HH^n(X,\sh{F})
\]
demuestra, por tanto, \eqref{III.4.1.7.10}.
Muestra también que el $S$-módulo graduado
\[
  \mathrm{R}=\bigoplus_{k\geq0}\mathrm{R}_k
\]
es un cociente del sub-$S$-módulo
\[
  \mathrm{M}
  =
  \bigoplus_{k\geq0}
  \HH^n(X,\mathfrak{J}^{k+1}\sh{F})
\]
de $\mathrm{E}$.
La observación anterior muestra, pues, que $\mathrm{R}$ es un
$S$-módulo de tipo finito.
Esto equivale a la afirmación del lema
(Bourbaki, \emph{Alg.\ comm.}, cap.~III, \S3, n.º~1, teorema~1).
\end{proof}

\begin{env}[4.1.7.12]
\label{III.4.1.7.12}
Consideremos ahora el $S$-módulo graduado
\[
  \mathrm{N}
  =
  \bigoplus_{k\geq0}
  \HH^{n+1}(X,\mathfrak{J}^{k+1}\sh{F}),
\]
definido como en \eqref{III.4.1.7.8}.
También es un $S$-módulo de tipo finito por \sref{III.3.3.2}.
Por \eqref{III.4.1.7.5}, tenemos
$\mathrm{Q}_k\subset\mathrm{N}_k$ para todo $k$, y el diagrama
\eqref{III.4.1.7.11}, sustituyendo $n$ por $n+1$, muestra que
\[
  S_m\mathrm{Q}_k
  =
  \mathfrak{J}^m\mathrm{Q}_k
  \subset
  \mathrm{Q}_{k+m}.
\]
En otros términos, $\mathrm{Q}$ es un sub-$S$-módulo graduado de
$\mathrm{N}$ y es, por tanto, de tipo finito.
\end{env}

\begin{env}[4.1.7.13]
\label{III.4.1.7.13}
Denotemos por $\alpha_m$ la inyección canónica
\[
  \mathfrak{J}^m\longrightarrow A,
\]
que también puede escribirse $S_m\to S_0$.
Como $\mathfrak{J}^{k+1}\sh{F}_k=0$, el $A$-módulo
$\HH^n(X,\sh{F}_k)$ es anulado por $\mathfrak{J}^{k+1}$.
Como $\mathrm{Q}_k$ es la imagen del $A$-homomorfismo
\[
  \HH^n(X,\sh{F}_k)
  \longrightarrow
  \HH^{n+1}(X,\mathfrak{J}^{k+1}\sh{F}),
\]
$\mathrm{Q}_k$ es igualmente anulado por $\mathfrak{J}^{k+1}$.
Equivalentemente, en el $S$-módulo $\mathrm{Q}$,
\[
  \alpha_{k+1}(S_{k+1})\mathrm{Q}_k=0.
  \tag{4.1.7.14}
  \label{III.4.1.7.14}
\]
Como $\mathrm{Q}$ es un $S$-módulo de tipo finito, existen enteros
$k_0$ y $h$ tales que
\[
  \mathrm{Q}_{k+h}=S_h\mathrm{Q}_k
  \qquad(k\geq k_0)
\]
\sref[II]{II.2.1.6}[ii].
Esta relación y \eqref{III.4.1.7.14} implican que existe un entero
$r>0$ tal que
\[
  \alpha_r(S_r)\mathrm{Q}=0.
  \tag{4.1.7.15}
  \label{III.4.1.7.15}
\]
\end{env}

\begin{env}[4.1.7.16]
\label{III.4.1.7.16}
La inyección canónica
\[
  \mathfrak{J}^{k+m}\sh{F}
  \longrightarrow
  \mathfrak{J}^k\sh{F}
\]
induce en cohomología un $A$-homomorfismo
\[
  \nu_m:
  \HH^{n+1}(X,\mathfrak{J}^{k+m}\sh{F})
  \longrightarrow
  \HH^{n+1}(X,\mathfrak{J}^k\sh{F}).
  \tag{4.1.7.17}
  \label{III.4.1.7.17}
\]
Para todo $x\in\mathfrak{J}^m$, tenemos manifiestamente la
factorización
\[
  \xymatrix@C=4.5em{
    \HH^{n+1}(X,\mathfrak{J}^k\sh{F})
      \ar[r]^{\mu_{x,m}}
    &
    \HH^{n+1}(X,\mathfrak{J}^{k+m}\sh{F})
      \ar[r]^{\nu_m}
    &
    \HH^{n+1}(X,\mathfrak{J}^k\sh{F}),
  }
  \tag{4.1.7.18}
  \label{III.4.1.7.18}
\]
cuya compuesta es $\mu_{x,0}$.
\end{env}

\oldpage[III]{128}

Se sigue que, para todo sub-$A$-módulo $\mathrm{P}$ de
$\HH^{n+1}(X,\mathfrak{J}^k\sh{F})$, tenemos en el $S$-módulo
$\mathrm{N}$
\[
  \nu_m(S_m\mathrm{P})
  =
  \alpha_m(S_m)\mathrm{P}.
  \tag{4.1.7.19}
  \label{III.4.1.7.19}
\]

\begin{lemma}[4.1.7.20]
\label{III.4.1.7.20}
Existe un entero $m>0$ tal que
\[
  \nu_m(\mathrm{Q}_{k+m})=0
\]
para todo $k\geq k_0$.
\end{lemma}

\begin{proof}
Tomemos $m$ múltiplo de $h$ con $m\geq r$.
Como
\[
  \mathrm{Q}_{k+m}=S_m\mathrm{Q}_k
  \qquad(k\geq k_0),
\]
las relaciones \eqref{III.4.1.7.19} y \eqref{III.4.1.7.15} dan
\[
  \nu_m(\mathrm{Q}_{k+m})
  =
  \alpha_m(S_m)\mathrm{Q}_k
  \subset
  \alpha_r(S_r)\mathrm{Q}_k
  =
  0.
\]
\end{proof}

\begin{env}[4.1.7.21]
\label{III.4.1.7.21}
El diagrama conmutativo
\[
  \xymatrix@C=2.5em{
    \HH^n(X,\sh{F})
      \ar[r]
    &
    \HH^n(X,\sh{F}_k)
      \ar[r]^-{\partial}
    &
    \HH^{n+1}(X,\mathfrak{J}^{k+1}\sh{F})
      \ar[r]
    &
    \HH^{n+1}(X,\sh{F})
    \\
    \HH^n(X,\sh{F})
      \ar[r]\ar[u]
    &
    \HH^n(X,\sh{F}_{k+m})
      \ar[r]^-{\partial}\ar[u]
    &
    \HH^{n+1}(X,\mathfrak{J}^{k+m+1}\sh{F})
      \ar[r]\ar[u]
    &
    \HH^{n+1}(X,\sh{F})
      \ar[u]
  }
\]
procede a su vez del diagrama conmutativo
\[
  \xymatrix{
    0\ar[r]
    &
    \mathfrak{J}^{k+1}\sh{F}\ar[r]
    &
    \sh{F}\ar[r]
    &
    \sh{F}_k\ar[r]
    &
    0
    \\
    0\ar[r]
    &
    \mathfrak{J}^{k+m+1}\sh{F}\ar[r]\ar[u]
    &
    \sh{F}\ar[r]\ar[u]
    &
    \sh{F}_{k+m}\ar[r]\ar[u]
    &
    0.
  }
\]
Aquí las flechas verticales son las aplicaciones canónicas.
Obtenemos, por tanto, el diagrama conmutativo
\[
  \xymatrix@C=2.2em{
    0\ar[r]
    &
    \mathrm{R}_k\ar[r]
    &
    \mathrm{H}\ar[r]
    &
    \mathrm{H}_k\ar[r]
    &
    \mathrm{Q}_k\ar[r]
    &
    0
    \\
    0\ar[r]
    &
    \mathrm{R}_{k+m}\ar[r]\ar[u]
    &
    \mathrm{H}\ar[r]\ar@{=}[u]
    &
    \mathrm{H}_{k+m}\ar[r]\ar[u]
    &
    \mathrm{Q}_{k+m}\ar[r]\ar[u]^{\nu_m}
    &
    0,
  }
\]
de filas exactas.
Para $k\geq k_0$, la última flecha vertical es nula por
\sref{III.4.1.7.20}.
Así, la imagen de $\mathrm{H}_{k+m}$ en $\mathrm{H}_k$ está contenida
en
\[
  \Ker(\mathrm{H}_k\longrightarrow\mathrm{Q}_k)
  =
  \Im(\mathrm{H}\longrightarrow\mathrm{H}_k).
\]
Recíprocamente, la conmutatividad muestra que contiene
$\Im(\mathrm{H}\to\mathrm{H}_k)$ y, por tanto, que es igual a esta
imagen.
Lo mismo ocurre, pues, con las imágenes en $\mathrm{H}_k$ de
$\mathrm{H}_{k'}$ para $k'\geq k+m$.
Esto demuestra la condición \emph{(ML)} para el sistema proyectivo
$(\mathrm{H}_k)_{k\geq0}$.

Además, para todo subconjunto abierto afín $U$ de $X$,
\[
  \HH^i(U,\sh{F}_k)=0
  \qquad(i>0)
\]
por \sref{III.1.3.1}, y, para $m>0$, la aplicación
\[
  \HH^0(U,\sh{F}_{k+m})
  \longrightarrow
  \HH^0(U,\sh{F}_k)
\]
es sobreyectiva \sref[I]{I.1.3.9}.
Podemos, por tanto,

\oldpage[III]{129}

aplicar \sref[0]{0.13.3.1}; el homomorfismo canónico
\[
  \HH^n(\widehat{X},\widehat{\sh{F}})
  \longrightarrow
  \varprojlim_k\HH^n(X,\sh{F}_k)
\]
es biyectivo para todo $n\geq0$.

Como el sistema proyectivo
$(\mathrm{H}/\mathrm{R}_k)_{k\geq0}$ es estricto, podemos pasar al
límite proyectivo en las sucesiones exactas
\[
  0
  \longrightarrow
  \mathrm{H}/\mathrm{R}_k
  \longrightarrow
  \mathrm{H}_k
  \longrightarrow
  \mathrm{Q}_k
  \longrightarrow
  0
  \tag{4.1.7.22}
  \label{III.4.1.7.22}
\]
por \sref[0]{0.13.2.2}.
Como $\nu_m(\mathrm{Q}_{k+m})=0$, tenemos
\[
  \varprojlim_k\mathrm{Q}_k=0,
\]
y, por tanto, un isomorfismo topológico
\[
  \varprojlim_k(\mathrm{H}/\mathrm{R}_k)
  \isoto
  \varprojlim_k\mathrm{H}_k.
\]
Pero la filtración $(\mathrm{R}_k)$ de $\mathrm{H}$ es
$\mathfrak{J}$-buena y define, por tanto, sobre $\mathrm{H}$ la
topología $\mathfrak{J}$-preádica.
Así, $\varprojlim_k(\mathrm{H}/\mathrm{R}_k)$ es la completación
separada de $\mathrm{H}$ para esta topología.
Esto termina la demostración de \sref{III.4.1.7}.
\end{env}
\end{proof}

\begin{env}[4.1.8]
\label{III.4.1.8}
Podemos terminar ahora la demostración de \sref{III.4.1.5}.
Para todo subconjunto abierto afín $V$ de $Y$,
\[
  \Gamma\bigl(
    V,\widehat{\RR^n f_*(\sh{F})}
  \bigr)
\]
es la completación separada de
$\Gamma(V,\RR^n f_*(\sh{F}))$ para la topología
$\mathfrak{J}$-preádica, si
$\sh{J}|V=\widetilde{\mathfrak{J}}$, porque
$\RR^n f_*(\sh{F})$ es un $\sh{O}_Y$-módulo coherente
\sref[I]{I.10.8.4}.
Además,
\[
  \Gamma\left(
    V,\varprojlim_k\RR^n f_*(\sh{F}_k)
  \right)
  =
  \varprojlim_k
  \Gamma\bigl(V,\RR^n f_*(\sh{F}_k)\bigr)
\]
por \sref[0]{0.3.2.6}.
De \sref{III.4.1.7} y \sref{III.1.4.11} se sigue que
$\vphi_n$ es un isomorfismo topológico.
Del mismo modo, de nuevo por \sref{III.1.4.11}, \sref{III.4.1.7}
muestra que el homomorfismo $\psi_{n,V}$ de \eqref{III.4.1.3.5} es un
isomorfismo.
Por la definición de
$\RR^n\widehat{f}_*(\widehat{\sh{F}})$, se sigue que $\psi_n$ es un
isomorfismo.
\end{env}

\begin{corollary}[4.1.9]
\label{III.4.1.9}
Bajo las hipótesis de \sref{III.4.1.4}, para todo subconjunto abierto
afín $V$ de $Y$, el homomorfismo canónico
\[
  \HH^n\bigl(
    \widehat{X}\cap f^{-1}(V),
    \widehat{\sh{F}}
  \bigr)
  \longrightarrow
  \Gamma\bigl(
    \widehat{Y}\cap V,
    \RR^n\widehat{f}_*(\widehat{\sh{F}})
  \bigr)
\]
es biyectivo.
\end{corollary}

\begin{env}[4.1.10]
\label{III.4.1.10}
\emph{Observación.}
Sea $f:X\to Y$ un morfismo de tipo finito entre esquemas ordinarios
noetherianos, y sea $\sh{F}$ un $\sh{O}_X$-módulo coherente cuyo
soporte es propio sobre $Y$ \sref[II]{II.5.4.10}.
Sabemos entonces \sref{III.3.2.4} que $\RR^n f_*(\sh{F})$ es un
$\sh{O}_Y$-módulo coherente para todo $n\geq0$.
Además, siempre podemos suponer que
\[
  \sh{F}=u_*(\sh{G}),
  \qquad
  \sh{G}=u^*(\sh{F}),
\]
donde $\sh{G}$ es un $\sh{O}_Z$-módulo coherente, $Z$ es un
subesquema cerrado adecuado de $X$ cuyo espacio subyacente es
$\Supp(\sh{F})$, y $u:Z\to X$ es la inyección canónica
\sref[I]{I.9.3.5}.
Si
\[
  \sh{G}_k
  =
  \sh{G}\otimes_{\sh{O}_Y}
  (\sh{O}_Y/\sh{J}^{k+1}),
\]
entonces
\[
  \sh{G}_k=u^*(\sh{F}_k),
  \qquad
  \RR^n f_*(\sh{F}_k)
  =
  \RR^n(f\circ u)_*(\sh{G}_k)
\]
y
\[
  \RR^n f_*(\sh{F})
  =
  \RR^n(f\circ u)_*(\sh{G})
\]
por \sref{III.1.3.4}; y, teniendo en cuenta \sref[I]{I.10.9.5},
\[
  \RR^n\widehat{f}_*(\widehat{\sh{F}})
  =
  \RR^n\widehat{(f\circ u)}_*(\widehat{\sh{G}}).
\]
Podemos, por tanto, aplicar \sref{III.4.1.5} a $\sh{G}$ y al morfismo
propio $f\circ u$.
Se sigue que, bajo estas hipótesis, los resultados de
\sref{III.4.1.5} siguen siendo válidos para $\sh{F}$ y $f$.
\end{env}


\subsection{Casos particulares y variantes}
\label{subsection:III.4.2}

La forma más útil del teorema de comparación \sref{III.4.1.5} es la
siguiente.

\begin{proposition}[4.2.1]
\label{III.4.2.1}
Sea $Y$ un esquema localmente noetheriano, sea $f:X\to Y$ un
morfismo propio, y sea $\sh{F}$ un $\sh{O}_X$-módulo coherente.
Entonces, para todo $y\in Y$ y todo $p$,
\[
  \bigl(\RR^p f_*(\sh{F})\bigr)_y
\]

\oldpage[III]{130}

es un $\sh{O}_y$-módulo de tipo finito y está, por tanto, separado
para la topología $\mathfrak{m}_y$-preádica.
Además, existe un isomorfismo topológico canónico
\[
  \widehat{\bigl(\RR^p f_*(\sh{F})\bigr)_y}
  \isoto
  \varprojlim_k
  \HH^p\left(
    f^{-1}(y),
    \sh{F}\otimes_{\sh{O}_Y}
    (\sh{O}_y/\mathfrak{m}_y^k)
  \right).
  \tag{4.2.1.1}
  \label{III.4.2.1.1}
\]
Aquí el primer término es la completación de
$\bigl(\RR^p f_*(\sh{F})\bigr)_y$ para la topología
$\mathfrak{m}_y$-preádica.
En el segundo término, para todo $k\geq0$, $f^{-1}(y)$ se considera
como el espacio subyacente del esquema
\[
  X\times_Y\Spec(\sh{O}_y/\mathfrak{m}_y^k)
\]
\sref[I]{I.3.6.1}.
\end{proposition}

\begin{proof}
Como $\sh{O}_y$ es un anillo local noetheriano y
$\bigl(\RR^p f_*(\sh{F})\bigr)_y$ es un $\sh{O}_y$-módulo de tipo
finito \sref{III.3.2.1}, la topología $\mathfrak{m}_y$-preádica sobre
este módulo está separada \sref[0]{0.7.3.5}.

Cuando $Y$ es noetheriano e $y$ es cerrado, las demás afirmaciones se
siguen de \sref{III.4.1.7}, reemplazando $Y$ por un entorno afín de
$y$ y tomando $Y'=\{y\}$; aquí utilizamos (G,~II,~4.9.1).

En general, pongamos
\[
  Y_1=\Spec(\sh{O}_y),
  \qquad
  X_1=X\times_Y Y_1,
  \qquad
  \sh{F}_1
  =
  \sh{F}\otimes_{\sh{O}_Y}\sh{O}_{Y_1},
\]
y sea
\[
  f_1=f\times 1_{Y_1}:X_1\longrightarrow Y_1.
\]
El esquema $Y_1$ es noetheriano, el morfismo $f_1$ es propio
\sref[II]{II.5.4.2}[iii], y $\sh{F}_1$ es coherente
\sref[0]{0.5.3.11}.
Sea $y_1$ el único punto cerrado de $Y_1$.
La proposición ya demostrada se aplica a $f_1$, $\sh{F}_1$ y
$y_1$.
Tenemos
\[
  \sh{O}_{y_1}=\sh{O}_y,
  \qquad
  f_1^{-1}(y_1)=f^{-1}(y)
\]
\sref[I]{I.3.6.5}.
Además, los esquemas
\[
  X\times_Y\Spec(\sh{O}_y/\mathfrak{m}_y^k)
  \quad\text{y}\quad
  X_1\times_{Y_1}\Spec(\sh{O}_y/\mathfrak{m}_y^k)
\]
se identifican canónicamente \sref[I]{I.3.3.9}, al igual que
\[
  \sh{F}_1
  \otimes_{\sh{O}_{Y_1}}
  (\sh{O}_y/\mathfrak{m}_y^k)
  \quad\text{y}\quad
  \sh{F}
  \otimes_{\sh{O}_Y}
  (\sh{O}_y/\mathfrak{m}_y^k)
\]
\sref[I]{I.9.1.6}.
Solo queda observar que
\[
  \RR^p f_{1*}(\sh{F}_1)
  \simeq
  \RR^p f_*(\sh{F})
  \otimes_{\sh{O}_Y}
  \sh{O}_{Y_1}.
\]
Esto se sigue de \sref{III.1.4.15}, pues el morfismo local
$\Spec(\sh{O}_y)\to Y$ es plano
\sref[0]{0.6.7.1}, \sref[I]{I.2.4.2}.
\end{proof}

El corolario siguiente utiliza la terminología de la teoría de la
dimensión (capítulo~IV) y no se aplicará antes de dicho capítulo.

\begin{corollary}[4.2.2]
\label{III.4.2.2}
Sea $Y$ un esquema localmente noetheriano, sea $f:X\to Y$ un
morfismo propio, sea $y\in Y$, y sea $r$ la dimensión de
$f^{-1}(y)$.
Entonces, para todo $\sh{O}_X$-módulo coherente $\sh{F}$, los haces
$\RR^p f_*(\sh{F})$ se anulan en un entorno de $y$ para todo
$p>r$.
\end{corollary}

\begin{proof}
Para todo $k$, tenemos
\[
  \HH^p\left(
    f^{-1}(y),
    \sh{F}\otimes_{\sh{O}_Y}
    (\sh{O}_y/\mathfrak{m}_y^k)
  \right)
  =
  0
\]
por (G,~II,~4.15.2).
De \sref{III.4.2.1} se sigue que la completación separada de
$\bigl(\RR^p f_*(\sh{F})\bigr)_y$ para la topología
$\mathfrak{m}_y$-preádica es nula.
Como esta topología está separada, tenemos también
\[
  \bigl(\RR^p f_*(\sh{F})\bigr)_y=0.
\]
La conclusión se sigue porque $\RR^p f_*(\sh{F})$ es coherente
\sref[0]{0.5.2.2}.
\end{proof}

\begin{env}[4.2.3]
\label{III.4.2.3}
El resultado \sref{III.4.2.1} se utiliza especialmente cuando $p=0$.
Obtenemos el corolario siguiente.
\end{env}

\begin{corollary}[4.2.4]
\label{III.4.2.4}
Bajo las hipótesis de \sref{III.4.2.1}, existe un isomorfismo
topológico canónico
\[
  \widehat{\bigl(f_*(\sh{F})\bigr)_y}
  \isoto
  \varprojlim_k
  \Gamma\left(
    f^{-1}(y),
    \sh{F}_y/\mathfrak{m}_y^k\sh{F}_y
  \right).
\]
\end{corollary}


\subsection{Teorema de conexidad de Zariski}
\label{subsection:III.4.3}
\label{III.4.3}

Los resultados de esta sección y de la siguiente generalizan teoremas
bien conocidos de Zariski, y todos pueden deducirse de
\sref{III.4.2.4}. Son consecuencias del teorema siguiente.

\begin{theorem}[4.3.1]
\label{III.4.3.1}
\textnormal{(Teorema de conexidad).}
Sea $Y$ un esquema localmente noetheriano y sea $f:X\to Y$ un
morfismo propio. Entonces
\[
  \sh{A}(X)=f_*(\sh{O}_X)
\]
es una $\sh{O}_Y$-álgebra coherente.

\oldpage[III]{131}

Sea $Y'$ el $Y$-esquema finito tal que
$\sh{A}(Y')=\sh{A}(X)$; está determinado salvo $Y$-isomorfismo
\sref[II]{II.1.3.1}, \sref[II]{II.6.1.3}.
Si $f'=\sh{A}(e)$ es el $Y$-morfismo $X\to Y'$ deducido del
isomorfismo identidad
\[
  e:\sh{A}(Y')\isoto\sh{A}(X)
\]
\sref[II]{II.1.2.7}, entonces $f'$ es propio,
$f'_*(\sh{O}_X)$ es isomorfo a $\sh{O}_{Y'}$, y toda fibra
$f'^{-1}(y')$, $y'\in Y'$, es conexa y no vacía.
\end{theorem}

\begin{proof}
Sea $g:Y'\to Y$ el morfismo estructural. Para demostrar que el
homomorfismo
\[
  \theta:\sh{O}_{Y'}\longrightarrow f'_*(\sh{O}_X)
\]
que interviene en la definición de $f'$ es biyectivo, basta, como
$Y'$ es afín sobre $Y$, demostrar que
\[
  g_*(\theta):
  g_*(\sh{O}_{Y'})
  \longrightarrow
  g_*\bigl(f'_*(\sh{O}_X)\bigr)
  =
  f_*(\sh{O}_X)
\]
es la identidad \sref[II]{II.1.4.2}. Esto se sigue de las
definiciones, pues
$g_*(\sh{O}_{Y'})=\sh{A}(Y')$ y
$f_*(\sh{O}_X)=\sh{A}(X)$.
La coherencia de $\sh{A}(X)$ es un caso particular del teorema de
finitud \sref{III.3.2.1}. Como $f$ es propio y $g$ es separado,
$f'$ es propio \sref[II]{II.5.4.3}[i]. Solo queda demostrar
\sref{III.4.3.2}.
\end{proof}

\begin{corollary}[4.3.2]
\label{III.4.3.2}
Bajo las hipótesis de \sref{III.4.3.1}, supongamos además que
$f_*(\sh{O}_X)$ es isomorfo a $\sh{O}_Y$. Entonces toda fibra
$f^{-1}(y)$ de $f$ es conexa y no vacía.
\end{corollary}

\begin{proof}
La hipótesis de que $f_*(\sh{O}_X)$ es isomorfo a $\sh{O}_Y$ implica
ya que $f$ es dominante, luego sobreyectivo porque $f$ es una
aplicación cerrada. Como en \sref{III.4.2.1}, podemos reducirnos al
caso en que $y$ es cerrado en $Y$. Como $f^{-1}(y)$ es un espacio
noetheriano, tiene un número finito de componentes conexas; estas son
también las componentes conexas del espacio subyacente de la
completación $\widehat{X}$ a lo largo de $f^{-1}(y)$.

Si $Z_i$, $1\leq i\leq n$, son estas componentes conexas, entonces
\[
  \Gamma(\widehat{X},\sh{O}_{\widehat{X}})
  =
  \prod_{i=1}^{n}\Gamma(Z_i,\sh{O}_{\widehat{X}}),
\]
y ninguno de los anillos del miembro derecho es nulo, pues la sección
unidad no es nula en ningún punto de $\widehat{X}$. Aplicando
\sref{III.4.1.5} a $\sh{F}=\sh{O}_X$, cuya completación a lo largo de
$f^{-1}(y)$ es $\sh{O}_{\widehat{X}}$, se ve que
$\Gamma(\widehat{X},\sh{O}_{\widehat{X}})$ es isomorfo a la
completación $\mathfrak{m}_y$-ádica separada
$\widehat{\sh{O}}_y$ del anillo local $\sh{O}_y$. Es, por tanto, un
anillo local y no puede ser un producto directo de varios anillos no
nulos. Así, $n=1$.
\end{proof}

\begin{corollary}[4.3.3]
\label{III.4.3.3}
Bajo las hipótesis de \sref{III.4.3.1}, para todo $y\in Y$ el conjunto
de componentes conexas de la fibra $f^{-1}(y)$ está en correspondencia
biyectiva con el conjunto finito de puntos de la fibra
$g^{-1}(y)$, donde $g:Y'\to Y$ es el morfismo estructural; en otros
términos, corresponde al conjunto de ideales maximales de
$\bigl(f_*(\sh{O}_X)\bigr)_y$.
\end{corollary}

\begin{proof}
Como $Y'$ es finito sobre $Y$, el espacio $g^{-1}(y)$ es finito y
discreto \sref[II]{II.6.1.7}. Como
\[
  f^{-1}(y)=f'^{-1}\bigl(g^{-1}(y)\bigr),
\]
la afirmación se sigue de esta observación y de
\sref{III.4.3.1}.
\end{proof}

Esto da una interpretación notable del $Y$-esquema $Y'$ definido en
\sref{III.4.3.1}. La factorización
$f=g\circ f'$ del morfismo propio $f$ es análoga a la factorización
obtenida por K.~Stein para aplicaciones holomorfas de espacios
analíticos, y se llamará en adelante la \emph{factorización de Stein}
de $f$.

\begin{remark}[4.3.4]
\label{III.4.3.4}
Sea $k$ una extensión del cuerpo $\kres(y)$. Si
\[
  f^{-1}(y)\otimes_{\kres(y)}k
  =
  X\times_Y\Spec(k)
\]
es conexo, entonces $f^{-1}(y)$ también es conexo, pues es la imagen
del primero por un morfismo de proyección
\sref[I]{I.3.4.7}. Para un morfismo $f:X\to Y$ de esquemas y un punto
$y\in Y$, diremos que la fibra $f^{-1}(y)$ es
\emph{geométricamente conexa} si, para toda extensión $k$ de
$\kres(y)$, el esquema
$f^{-1}(y)\otimes_{\kres(y)}k=X\times_Y\Spec(k)$ es conexo.

\oldpage[III]{132}

Bajo las hipótesis de \sref{III.4.3.2}, la conclusión puede
reforzarse: las fibras $f^{-1}(y)$ son geométricamente conexas.
En efecto, para toda extensión $k$ de $\kres(y)$ existen un anillo
local noetheriano $A$ y un homomorfismo local
\[
  \varphi:\sh{O}_y\longrightarrow A
\]
tal que $A$ es un $\sh{O}_y$-módulo plano y el cuerpo residual de
$A$ es $\kres(y)$-isomorfo a $k$ \sref[0]{0.10.3.1}.
Sea $Y_1=\Spec(A)$, y sea $h:Y_1\to Y$ el morfismo local
correspondiente a $\varphi$, que lleva el único punto cerrado $y_1$ de
$Y_1$ a $y$ \sref[I]{I.2.4.1}. Pongamos
\[
  X_1=X\times_Y Y_1,
  \qquad
  f_1=f\times 1_{Y_1}.
\]
El esquema $Y_1$ es noetheriano, $f_1$ es propio
\sref[II]{II.5.4.2}, y $f_1^{-1}(y_1)$ es un
$\kres(y_1)$-esquema isomorfo a $X\times_Y\Spec(k)$. Basta, por
tanto, demostrar que
\[
  f_{1*}(\sh{O}_{X_1})=\sh{O}_{Y_1},
\]
de modo que \sref{III.4.3.2} pueda aplicarse a $f_1$. Ahora $h$ es
plano por \sref[I]{I.2.4.2}; por tanto, por \sref{III.1.4.15} con
$q=0$,
\[
  f_{1*}(\sh{O}_{X_1})
  =
  h^*\bigl(f_*(\sh{O}_X)\bigr)
  =
  h^*(\sh{O}_Y)
  =
  \sh{O}_{Y_1}.
\]

En el caso general de \sref{III.4.3.1}, el mismo argumento muestra,
con la notación de dicho teorema, que
\[
  f_{1*}(\sh{O}_{X_1})
  =
  h^*\bigl(g_*(\sh{O}_{Y'})\bigr),
\]
y que la factorización de Stein $f_1=g_1\circ f'_1$ de $f_1$
satisface
\[
  g_1=g\times 1_{Y_1}
\]
\sref[II]{II.1.5.2}, siendo el $Y_1$-esquema finito correspondiente
$Y'_1=Y'\times_Y Y_1$. Teniendo en cuenta la transitividad de las
fibras \sref[I]{I.3.6.4}, se sigue de
\sref{III.4.3.3} que el número de componentes conexas de
$f_1^{-1}(y_1)$ es el número de elementos de
\[
  g_1^{-1}(y_1)
  =
  g^{-1}(y)\otimes_{\kres(y)}k.
\]
Si $k$ se elige algebraicamente cerrado, este número es independiente
de la extensión algebraicamente cerrada elegida y es igual al
\emph{número geométrico de puntos} de $g^{-1}(y)$
\sref[I]{I.6.4.7}, o equivalentemente a
\[
  \sum_i[\kres(y'_i):\kres(y)]_s,
\]
donde los $y'_i$ recorren el conjunto finito $g^{-1}(y)$. Este número
se llama \emph{número geométrico de componentes conexas} de
$f^{-1}(y)$. Los cuerpos $\kres(y'_i)$ son precisamente los cuerpos
residuales del anillo semilocal $\bigl(f_*(\sh{O}_X)\bigr)_y$.
\end{remark}

\begin{proposition}[4.3.5]
\label{III.4.3.5}
Sean $X$ e $Y$ esquemas íntegros localmente noetherianos, y sea
$f:X\to Y$ un morfismo propio dominante. Para todo $y\in Y$, el
número de componentes conexas de $f^{-1}(y)$ es a lo sumo el número
de ideales maximales de la clausura entera $\sh{O}'_y$ de
$\sh{O}_y$ en el cuerpo de funciones racionales $\mathrm{R}(X)$.
\end{proposition}

\begin{proof}
Para todo abierto $U$ de $Y$,
\[
  \Gamma\bigl(U,f_*(\sh{O}_X)\bigr)
  =
  \Gamma\bigl(f^{-1}(U),\sh{O}_X\bigr)
\]
es la intersección de los anillos locales $\sh{O}_x$ para
$x\in f^{-1}(U)$ \sref[I]{I.8.2.1.1}. Se sigue inmediatamente que
$\bigl(f_*(\sh{O}_X)\bigr)_y$ es un subanillo de $\mathrm{R}(X)$ que
contiene a $\sh{O}_y$. Como $f_*(\sh{O}_X)$ es un
$\sh{O}_Y$-módulo coherente, $\bigl(f_*(\sh{O}_X)\bigr)_y$ es un
$\sh{O}_y$-módulo finito y, por tanto, está contenido en $\sh{O}'_y$.
Todo ideal maximal de tal anillo es la intersección de este anillo con
un ideal maximal de $\sh{O}'_y$
\cite[vol.~I, pp.~257 and 259]{I-13}, lo que demuestra la proposición.
\end{proof}

\begin{definition}[4.3.6]
\label{III.4.3.6}
Un anillo local íntegro se dice \emph{unirrama} si su clausura entera
es de nuevo un anillo local. Un punto $y$ de un esquema íntegro
$Y$ se dice \emph{unirrama} si el anillo local $\sh{O}_y$ es
unirrama; este es, en particular, el caso cuando $Y$ es normal en
$y$.
\end{definition}

Sea $A$ un anillo local íntegro y sea $K$ su cuerpo de fracciones. El
anillo $A$ es unirrama si y solo si todo subanillo $A_1$ de $K$ que
contenga a $A$ y sea una $A$-álgebra finita es local. En efecto, sea
$A'$ la clausura entera de $A$. Por el primer teorema de
Cohen--Seidenberg (Bourbaki, \emph{Alg\`ebre commutative}, capítulo~V,
\S2, n.º~1, teorema~1), todo ideal maximal de $A_1$ es la contracción
a $A_1$ de un ideal maximal de $A'$; así, si $A'$ es local, también
lo es $A_1$.

\oldpage[III]{133}

Recíprocamente, $A'$ es el límite inductivo de la familia dirigida de
$A$-subálgebras finitas $A_\alpha$ de $A'$. Si cada $A_\alpha$ es
local, entonces, para $A_\alpha\subset A_\beta$, el ideal maximal de
$A_\alpha$ es la contracción del ideal maximal de $A_\beta$, por el
mismo argumento; por tanto, $A'$ es local
\sref[0]{0.10.3.1.3}.

Observemos también que, si la completación de un anillo local
noetheriano $A$ es íntegra ---propiedad que se expresa diciendo que
$A$ es \emph{analíticamente íntegro}---, entonces $A$ es unirrama.
Sea $\mathfrak{m}$ el ideal maximal de $A$, sea $K$ su cuerpo de
fracciones, y sea $K'$ el cuerpo de fracciones de $\widehat{A}$;
entonces
\[
  K'=K\otimes_A\widehat{A}.
\]
Sea $A'_F$ una $A$-subálgebra finita de $K$. El subanillo $B_F$ de
$K'$ generado por $\widehat{A}$ y $A'_F$ es isomorfo a
\[
  A'_F\otimes_A\widehat{A};
\]
es un $\widehat{A}$-módulo finito y es la completación de $A'_F$
para la topología $\mathfrak{m}$-ádica
\sref[0]{0.7.3.3}, \sref[0]{0.7.3.6}. Como $A'_F$ es semilocal
(Bourbaki, \emph{Alg\`ebre commutative}, capítulo~IV, \S2, n.º~5,
corolario~3 de la proposición~9), mientras que su completación es
íntegra, $A'_F$ solo puede tener un ideal maximal $\mathfrak{m}'_F$, y
$\mathfrak{m}'_F\cap A=\mathfrak{m}$. Esto demuestra la afirmación.

\begin{corollary}[4.3.7]
\label{III.4.3.7}
Bajo las hipótesis de \sref{III.4.3.5}, supongamos que la clausura
algebraica de $\mathrm{R}(Y)$ en $\mathrm{R}(X)$ tiene grado separable
$n$ y que $y\in Y$ es unirrama. Entonces la fibra $f^{-1}(y)$ tiene a
lo sumo $n$ componentes conexas. En particular, si la clausura
algebraica de $\mathrm{R}(Y)$ en $\mathrm{R}(X)$ es radicial sobre
$\mathrm{R}(Y)$, entonces $f^{-1}(y)$ es conexa.
\end{corollary}

\begin{proof}
Sea $\sh{O}'_y$ la clausura entera de $\sh{O}_y$, y sea
$\sh{O}''_y$ la clausura entera de $\sh{O}_y$ en $\mathrm{R}(X)$.
Esta última es también la clausura entera de $\sh{O}'_y$ en
$\mathrm{R}(X)$. Si $\sh{O}'_y$ es local, entonces $\sh{O}''_y$ es
un anillo semilocal que tiene a lo sumo $n$ ideales maximales
\cite[vol.~I, p.~289, Theorem~22]{I-13}.

Este corolario es esencialmente la forma en la que Zariski enuncia su
``teorema de conexidad'' para esquemas algebraicos.
\end{proof}

\begin{remark}[4.3.8]
\label{III.4.3.8}
Si, además de las hipótesis de \sref{III.4.3.7}, $Y$ es normal en
$y$, entonces la fibra $f^{-1}(y)$ es geométricamente conexa: con la
notación de \sref{III.4.3.4}, $g^{-1}(y)$ es un solo punto $y'$, y
$\kres(y')$ es radicial sobre $\kres(y)$.
\end{remark}

\begin{definition}[4.3.9]
\label{III.4.3.9}
Sea $Y$ un esquema localmente noetheriano. Un morfismo de tipo finito
$f:X\to Y$ se dice \emph{universalmente abierto} si, para todo esquema
irreducible localmente noetheriano $Y'$ y todo morfismo dominante
$g:Y'\to Y$, toda componente irreducible de
$X'=X\times_Y Y'$ domina a $Y'$.
\end{definition}

Si $Y$ es irreducible, esto significa que, si $\eta$ y $\eta'$ son
los puntos genéricos de $Y$ e $Y'$, respectivamente, y si
$f'=f\times 1_{Y'}$, toda componente irreducible de $X'$ tiene su
punto genérico en $f'^{-1}(\eta')$. En consecuencia, para todo
abierto $U$ de $Y$, el morfismo restringido $f^{-1}(U)\to U$ es
universalmente abierto.

\begin{corollary}[4.3.10]
\label{III.4.3.10}
Sean $X$ e $Y$ esquemas íntegros localmente noetherianos, y sea
$f:X\to Y$ propio, dominante y universalmente abierto. Si la clausura
algebraica de $\mathrm{R}(Y)$ en $\mathrm{R}(X)$ es radicial sobre
$\mathrm{R}(Y)$, entonces toda fibra $f^{-1}(y)$, $y\in Y$, es
geométricamente conexa.
\end{corollary}

\begin{proof}
Podemos reducirnos al caso $Y=\Spec(B)$, donde $B$ es un anillo
íntegro noetheriano. Por \sref[II]{II.7.1.7}, existe un anillo local
noetheriano íntegramente cerrado $A$ que domina a $\sh{O}_y$ y tiene a
$\mathrm{R}(Y)$ como cuerpo de fracciones. Pongamos $Y'=\Spec(A)$, y
sea $h:Y'\to Y$ el morfismo correspondiente a la inyección canónica
$B\to A$. Es birracional, luego dominante; además, si $y'$ es el
único punto cerrado de $Y'$, entonces $h(y')=y$.

Pongamos
\[
  X'=X\times_Y Y',
  \qquad
  f'=f\times 1_{Y'}.
\]
Sean $\eta$, $\eta'$ y $\xi$ los puntos genéricos de $Y$, $Y'$ y
$X$, respectivamente. Entonces $f(\xi)=\eta$ y
$h^{-1}(\eta)=\{\eta'\}$; además,
$\kres(\eta)=\kres(\eta')=\mathrm{R}(Y)$. Así, $f'^{-1}(\eta')$ es
isomorfo a $f^{-1}(\eta)$ \sref[I]{I.3.6.4}, y tiene un único punto
genérico, pues $\xi$ es el punto genérico de
$f^{-1}(\eta)$ \sref[0]{0.2.1.8}. Por la apertura universal, toda
componente irreducible de $X'$ tiene su punto genérico en
$f'^{-1}(\eta')$. Por tanto, $X'$ es irreducible; su punto genérico
$\xi'$ es el de $f'^{-1}(\eta')$, y
$\kres(\xi')=\kres(\xi)$.

\oldpage[III]{134}

Pongamos $X''=X'_{\mathrm{red}}$ y $f''=f'_{\mathrm{red}}$. Entonces
$X''$ es íntegro y noetheriano, $f''$ es propio
\sref[II]{II.5.4.6}, y los espacios subyacentes de las fibras
$f'^{-1}(y')$ y $f''^{-1}(y')$ son iguales. Además,
\[
  \mathrm{R}(X'')=\kres(\xi')=\mathrm{R}(X).
\]
Así, $f''$ satisface las hipótesis de \sref{III.4.3.8}, y
$f''^{-1}(y')$ es geométricamente conexa.

Sea ahora $k$ una extensión cualquiera de $\kres(y)$. Existe una
extensión $k_1$ de $\kres(y)$ en la que tanto $\kres(y')$ como $k$
pueden considerarse subextensiones (Bourbaki, \emph{Alg\`ebre},
capítulo~V, \S4, proposición~2). Por el párrafo anterior,
\[
  f''^{-1}(y')\times_{Y'}\Spec(k_1)
\]
es conexo. Tiene el mismo esquema reducido asociado que
$f'^{-1}(y')\times_{Y'}\Spec(k_1)$
\sref[I]{I.5.1.8}, luego este último es conexo; es isomorfo a
$f^{-1}(y)\times_Y\Spec(k_1)$ \sref[I]{I.3.6.4}. Por tanto,
$f^{-1}(y)\times_Y\Spec(k)$ es conexo, y la observación al comienzo
de \sref{III.4.3.4} termina la demostración.
\end{proof}

\begin{remark}[4.3.11]
\label{III.4.3.11}
\emph{(i)}
El argumento anterior se debe, en lo esencial, a Zariski
\cite{I-20}, salvo que, en su marco, puede tomarse como $A$ la
clausura entera de $\sh{O}_y$, pues esta es un anillo noetheriano para
los anillos locales de la geometría algebraica clásica. Zariski
demuestra también que, si $Y$ es la variedad de Chow de un espacio
proyectivo $\mathbf{P}_k^r$ sobre un cuerpo $k$, y $X$ es el
subconjunto cerrado de $\mathbf{P}_k^r\times_kY$ que define la
correspondencia de Chow entre $\mathbf{P}_k^r$ e $Y$, entonces la
proyección $X\to Y$ es universalmente abierta (\emph{loc.\ cit.},
lema de la p.~82).
Parece que esta es la única propiedad formal de las ``coordenadas de
Chow'' utilizada en ciertas aplicaciones. En tal situación es, por
tanto, útil reemplazar el lenguaje de la especialización de ciclos en
el espacio proyectivo por el de las fibras de un morfismo propio,
supuesto eventualmente universalmente abierto o sometido a condiciones
análogas de regularidad local.

\emph{(ii)}
En el capítulo~IV veremos que un morfismo universalmente abierto
$f:X\to Y$ puede definirse equivalentemente exigiendo que, para todo
morfismo $Z\to Y$, el morfismo obtenido por cambio de base
$f_{(Z)}:X_{(Z)}\to Z$ sea abierto. También puede demostrarse que, si
$f$ satisface las hipótesis de \sref{III.4.3.10}, y si $y$ es una
especialización de $y'$, entonces el número geométrico de componentes
conexas de $f^{-1}(y)$ es a lo sumo el de $f^{-1}(y')$.
\end{remark}

\begin{corollary}[4.3.12]
\label{III.4.3.12}
Bajo las hipótesis de \sref{III.4.3.5}, supongamos además que
$\mathrm{R}(Y)$ es algebraicamente cerrado en $\mathrm{R}(X)$, y sea
$y$ un punto normal de $Y$. Entonces $f^{-1}(y)$ es geométricamente
conexa, y existe un entorno abierto $U$ de $y$ en $Y$ tal que
\[
  f_*\bigl(\sh{O}_X|_{f^{-1}(U)}\bigr)
  \simeq
  \sh{O}_Y|_U.
\]
En particular, si $Y$ es normal y $\mathrm{R}(Y)$ es algebraicamente
cerrado en $\mathrm{R}(X)$, entonces
$f_*(\sh{O}_X)\simeq\sh{O}_Y$.
\end{corollary}

\begin{proof}
La afirmación relativa a $f^{-1}(y)$ es un caso particular de
\sref{III.4.3.8}.

\oldpage[III]{135}

Si
\[
  X\xrightarrow{f'}Y'\xrightarrow{g}Y
\]
es la factorización de Stein de $f$ \sref{III.4.3.3}, se sigue que
$g^{-1}(y)$ consta de un solo punto $y'$. Además,
\[
  \sh{O}_y
  \subset
  \sh{O}_{y'}
  =
  \bigl(f_*(\sh{O}_X)\bigr)_y
  \subset
  \mathrm{R}(X).
\]
Como $\sh{O}_{y'}$ es finito sobre $\sh{O}_y$, es algebraico sobre $\mathrm{R}(Y)$ y, por consiguiente, está contenido en $\mathrm{R}(Y)$ en virtud de la hipótesis. Como $y$ es normal, necesariamente
$\sh{O}_{y'}=\sh{O}_y$. Así, $g$ es un isomorfismo local en $y'$
\sref[I]{I.6.5.4}, lo que demuestra la primera parte. La segunda
afirmación se sigue de la primera: la hipótesis adicional hace que $g$
sea biyectivo y un isomorfismo local en un entorno de cada punto de
$Y'$, luego un isomorfismo.
\end{proof}

El hecho de que \sref{III.4.3.7} esté disponible para esquemas da, por
ejemplo, la aplicación siguiente.

\begin{proposition}[4.3.13]
\label{III.4.3.13}
Sea $A$ un anillo local noetheriano unirrama, sea $\mathfrak{a}$ un
ideal de definición de $A$, y sean
\[
  A_0=A/\mathfrak{a},
  \qquad
  S=\operatorname{gr}_{\mathfrak{a}}(A)
\]
el anillo graduado asociado a la filtración $\mathfrak{a}$-preádica.
Supongamos que $S$ es una $A_0$-álgebra graduada generada por $S_1$,
donde $S_1$ es un $A_0$-módulo finito. Entonces $\Proj(S)$ es un
$A_0$-esquema conexo.
\end{proposition}

\begin{proof}
Sea $\mathfrak{m}$ el ideal maximal de $A$. El esquema
$Y=\Spec(A)$ es íntegro, y el punto $y$ correspondiente a
$\mathfrak{m}$ es su único punto cerrado. Para cierto entero $p$,
\[
  \mathfrak{m}^p\subset\mathfrak{a}\subset\mathfrak{m},
\]
de modo que $V(\mathfrak{a})=\{\mathfrak{m}\}$. Pongamos
\[
  S'=\bigoplus_{n\geq0}\mathfrak{a}^n
  \qquad\text{y}\qquad
  X=\Proj(S').
\]
Este es el $Y$-esquema obtenido por la explosión del ideal
$\mathfrak{a}$. Es íntegro, y el morfismo estructural
$f:X\to Y$ es birracional \sref[II]{II.8.1.4} y proyectivo.
Por tanto, \sref{III.4.3.7} se aplica y muestra que $f^{-1}(y)$ es
conexo. Pero el espacio subyacente de $f^{-1}(y)$ es el de
\[
  \Proj(S'\otimes_A A_0)
\]
\sref[I]{I.3.6.1}, \sref[II]{II.2.8.10}; como
$S'\otimes_AA_0=S$ por definición, se sigue la proposición.
\end{proof}


\subsection{El ``teorema principal'' de Zariski}
\label{subsection:III.4.4}

\begin{proposition}[4.4.1]
\label{III.4.4.1}
Sea $Y$ un esquema localmente noetheriano y sea $f:X\to Y$ un
morfismo propio. Sea $X'$ el conjunto de puntos $x\in X$ que están
aislados en su fibra $f^{-1}(f(x))$. Entonces $X'$ es abierto en $X$.
Si
\[
  f=g\circ f'
\]
es la factorización de Stein de $f$ \sref{III.4.3.3}, la restricción
de $f'$ a $X'$ es un isomorfismo de $X'$ sobre el subesquema inducido
en un subconjunto abierto $U$ de $Y'$, y
$X'=f'^{-1}(U)$.
\end{proposition}

\begin{proof}
Como $g^{-1}(f(x))$ es finito y discreto
\sref{III.4.3.3}, \sref[II]{II.6.1.7}, el punto $x$ está aislado en
$f^{-1}(f(x))$ si y solo si está aislado en
$f'^{-1}(f'(x))$. Podemos, por tanto, reducirnos al caso $f'=f$, de
modo que $f_*(\sh{O}_X)=\sh{O}_Y$. Si $x\in X'$, la fibra
$f^{-1}(f(x))$, que es conexa por \sref{III.4.3.2}, debe constar
entonces del único punto $x$.

Como $f$ es cerrado, para todo entorno abierto $V$ de $x$ en $X$, el
conjunto cerrado $f(X-V)$ no contiene a $y=f(x)$. Si $U$ es su
complementario en $Y$, entonces $f^{-1}(U)\subset V$. Así, las
imágenes inversas por $f$ de un sistema fundamental de entornos
abiertos de $y$ forman un sistema fundamental de entornos abiertos de
$x$. La igualdad $f_*(\sh{O}_X)=\sh{O}_Y$ y la definición de la
imagen directa de un haz \sref[0]{0.3.4.1}, \sref{III.4.2.1} implican
ahora que, si $f=(\psi,\theta)$, el homomorfismo
\[
  \theta_y^\#:\sh{O}_y\longrightarrow\sh{O}_x
\]
es un isomorfismo. Por tanto, existen entornos abiertos $V$ de $x$ y
$U$ de $y$ tales que la restricción de $f$ es un isomorfismo
$V\isoto U$ \sref[I]{I.6.5.4}.

\oldpage[III]{136}

Por el párrafo anterior podemos disponer además que
$f^{-1}(U)=V$. Se sigue inmediatamente que $V\subset X'$, lo que
termina la demostración.
\end{proof}

La proposición siguiente fue demostrada por Chevalley para esquemas
algebraicos.

\begin{proposition}[4.4.2]
\label{III.4.4.2}
Sea $Y$ un esquema localmente noetheriano y sea $f:X\to Y$ un
morfismo. Las condiciones siguientes son equivalentes:
\begin{enumerate}
  \item[{\rm a)}] $f$ es finito;
  \item[{\rm b)}] $f$ es afín y propio;
  \item[{\rm c)}] $f$ es propio y, para todo $y\in Y$, el conjunto
    $f^{-1}(y)$ es finito.
\end{enumerate}
\end{proposition}

\begin{proof}
La condición~a) implica b) por \sref[II]{II.6.1.2} y
\sref[II]{II.6.1.11}. Si $f$ es propio y afín, también lo es
\[
  f^{-1}(y)\longrightarrow\Spec(\kres(y))
\]
\sref[II]{II.1.6.2}, \sref[II]{II.5.4.2}. El teorema de finitud
\sref{III.3.2.1}, aplicado al haz estructural de $f^{-1}(y)$, muestra
que $f^{-1}(y)=\Spec(A)$ para una $\kres(y)$-álgebra finita $A$. Así,
$f^{-1}(y)$ es un conjunto finito \sref[II]{II.6.1.7}, y b) implica
c).

Por último, si se cumple c), la fibra $f^{-1}(y)$ es un esquema
algebraico sobre $\kres(y)$ cuyo conjunto subyacente es finito y es,
por tanto, discreto \sref[I]{I.6.4.4}. Con la notación de
\sref{III.4.4.1}, tenemos $X'=X$, luego
$f':X\to Y'$ es un isomorfismo. Como $g$ es finito, c) implica a).
\end{proof}

\begin{theorem}[4.4.3]
\label{III.4.4.3}
\textnormal{(``Teorema principal'' de Zariski).}
Sea $Y$ un esquema noetheriano, sea $f:X\to Y$ un morfismo
cuasiproyectivo, y sea $X'$ el conjunto de puntos $x\in X$ que están
aislados en su fibra $f^{-1}(f(x))$. Entonces $X'$ es abierto en $X$,
y el subesquema inducido por $X$ sobre $X'$ es isomorfo a un
subesquema inducido sobre un subconjunto abierto de un $Y$-esquema
$Y'$ finito sobre $Y$.
\end{theorem}

\begin{proof}
Existe un $Y$-esquema proyectivo $Z$ tal que $X$ es
$Y$-isomorfo al subesquema inducido sobre un subconjunto abierto de
$Z$ \sref[II]{II.5.3.2}, \sref[II]{II.5.5.1}. Nos reducimos así al
caso en que $f$ es proyectivo, luego propio \sref[II]{II.5.5.3}, y el
teorema se sigue inmediatamente de
\sref{III.4.4.1}.
\end{proof}

\begin{remark}[4.4.4]
\label{III.4.4.4}
Si $X$ es reducido (respectivamente, si $X$ es irreducible y $X'$ no
es vacío), puede exigirse en \sref{III.4.4.3} que $Y'$ sea reducido
(respectivamente, irreducible). En efecto, $Y'$ puede reemplazarse por
la clausura esquemática $\overline{X'}$ de $X'$ en $Y'$
\sref[I]{I.9.5.11}, \sref[II]{II.6.1.5}. Si $X'$ es reducido, también
lo es $\overline{X'}$ \sref[I]{I.9.5.9}; si $X'$ no es vacío y es
irreducible, entonces $\overline{X'}$ también es irreducible.
\end{remark}

\begin{corollary}[4.4.5]
\label{III.4.4.5}
Sea $Y$ un esquema localmente noetheriano, sea $f:X\to Y$ un
morfismo de tipo finito, y sea $x$ un punto de $X$ aislado en su fibra
$f^{-1}(f(x))$. Entonces existe un entorno abierto de $x$ en $X$
isomorfo a un subconjunto abierto de un $Y$-esquema finito sobre $Y$.
\end{corollary}

\begin{proof}
Pongamos $y=f(x)$, sea $U$ un entorno abierto afín de $y$ en $Y$, y
sea $V$ un entorno abierto afín de $x$ en $X$, contenido en
$f^{-1}(U)$. Como $Y$ es separado, la inclusión $U\to Y$ es afín
\sref[II]{II.1.6.3}; y como $V$ es afín sobre $U$ (\emph{ibid.}), la
restricción de $f$ a $V$ es un morfismo afín $V\to Y$
\sref[II]{II.1.6.2}. A fortiori, esta restricción es
cuasiproyectiva, pues es de tipo finito
\sref[I]{I.6.3.5}, \sref[II]{II.5.3.4}. Solo queda aplicar
\sref{III.4.4.3} a esta restricción.
\end{proof}

El Corolario \sref{III.4.4.5} puede enunciarse en el lenguaje del
álgebra conmutativa como sigue.

\begin{corollary}[4.4.6]
\label{III.4.4.6}
Sea $A$ un anillo noetheriano, sea $B$ una $A$-álgebra de tipo
finito, sea $\mathfrak q$ un ideal primo de $B$, y sea $\mathfrak p$
su imagen inversa en $A$. Supongamos que $\mathfrak q$ es a la vez
maximal y minimal en el conjunto de ideales primos de $B$ cuya imagen
inversa es $\mathfrak p$. Entonces existen un elemento
$g\in B-\mathfrak q$, una $A$-álgebra finita $A'$, y un elemento
$f'\in A'$ tales que las $A$-álgebras $B_g$ y $A'_{f'}$ son
isomorfas.
\end{corollary}

\begin{proof}
Apliquemos \sref{III.4.4.5} a $Y=\Spec(A)$ y $X=\Spec(B)$. La
hipótesis sobre $\mathfrak q$ dice precisamente que $\mathfrak q$
está aislado en su fibra \sref[I]{I.1.1.7}.
\end{proof}

La consecuencia siguiente parece menos general.

\oldpage[III]{137}

\begin{corollary}[4.4.7]
\label{III.4.4.7}
Sea $A$ un anillo local noetheriano, sea $B$ una $A$-álgebra de tipo
finito, y sea $\mathfrak n$ un ideal primo de $B$ cuya imagen inversa
en $A$ es el ideal maximal $\mathfrak m$. Supongamos que $\mathfrak n$
es maximal en $B$ y minimal en el conjunto de ideales primos de $B$
cuya imagen inversa es $\mathfrak m$ (lo que significa también que
$\mathfrak mB$ es $\mathfrak n$-primario). Entonces existen una
$A$-álgebra finita $A'$ y un ideal maximal $\mathfrak m'$ de $A'$
(cuya imagen inversa en $A$ es $\mathfrak m$) tales que $B_{\mathfrak
n}$ es isomorfo, como $A$-álgebra, a $A'_{\mathfrak m'}$.
\end{corollary}

El caso particular siguiente de \sref{III.4.4.7} se llama también a
veces el ``teorema principal''.

\begin{corollary}[4.4.8]
\label{III.4.4.8}
Bajo las hipótesis de \sref{III.4.4.7}, supongamos además que $A$ y
$B$ son dominios íntegros con el mismo cuerpo de fracciones $K$.
Si $A$ es íntegramente cerrado, entonces $A=B$.
\end{corollary}

\begin{proof}
La Observación \sref{III.4.4.4} muestra que, al aplicar
\sref{III.4.4.7}, podemos suponer que $A'$ es un dominio íntegro con
cuerpo de fracciones $K$. La hipótesis sobre $A$ da entonces
$A'=A$, luego $B_{\mathfrak n}=A$. Como
$A\subset B\subset B_{\mathfrak n}$, se sigue que $A=B$.
\end{proof}

El enunciado \sref{III.4.4.8} es la forma en la que Zariski dio su
``teorema principal'' (extendido aquí a anillos locales íntegros
noetherianos arbitrarios).

Los corolarios anteriores eran variantes locales de
\sref{III.4.4.3}, que es un resultado global. He aquí otra
consecuencia de naturaleza global.

\begin{corollary}[4.4.9]
\label{III.4.4.9}
Sean $X$ e $Y$ dos preesquemas íntegros localmente noetherianos, y sea
$f:X\to Y$ un morfismo separado, localmente de tipo finito y birracional.
Supongamos además que $Y$ es normal y que toda fibra $f^{-1}(y)$,
para $y\in Y$, es finita. Entonces $f$ es una inmersión abierta. Si,
además, $f$ es cerrado (en particular, si $f$ es propio), entonces
$f$ es un isomorfismo.
\end{corollary}

\begin{proof}
Sea $x\in X$ y pongamos $y=f(x)$. Como $f^{-1}(y)$ es un esquema
localmente de tipo finito sobre $\kres(y)$, la hipótesis de que es finito implica que
es discreto \sref[I]{I.6.4.4}. Además, $\sh{O}_y$ es íntegramente
cerrado, y $\sh{O}_x$ y $\sh{O}_y$ tienen el mismo cuerpo de fracciones
\sref[I]{I.7.1.5}. Podemos, por tanto, aplicar
\sref{III.4.4.8}; si $f=(\psi,\theta)$, el homomorfismo
\[
  \theta_y^\#:\sh{O}_y\longrightarrow\sh{O}_x
\]
es biyectivo. De \sref[I]{I.6.5.4} se sigue que $f$ es un isomorfismo
local. Como $f$ es separado y $X$ es íntegro, $f$ es una inmersión
abierta \sref[I]{I.8.2.8}. La última afirmación se sigue del hecho de
que $f$ es dominante.
\end{proof}

\begin{proposition}[4.4.10]
\label{III.4.4.10}
Sea $Y$ un esquema localmente noetheriano, y sea $f:X\to Y$ un
morfismo localmente de tipo finito. El conjunto $X'$ de puntos
$x\in X$ que están aislados en su fibra $f^{-1}(f(x))$ es abierto en
$X$.
\end{proposition}

\begin{proof}
La cuestión es local sobre $X$ e $Y$, por lo que podemos suponer que
$X$ e $Y$ son afines noetherianos y que $f$ es de tipo finito. El
morfismo $f$ es entonces afín y de tipo finito, luego
cuasiproyectivo \sref[II]{II.5.3.4}; basta aplicar
\sref{III.4.4.3}.
\end{proof}

\oldpage[III]{138}

\begin{corollary}[4.4.11]
\label{III.4.4.11}
Sea $Y$ un esquema localmente noetheriano, y sea $f:X\to Y$ un
morfismo propio. El conjunto $U$ de puntos $y\in Y$ para los que
$f^{-1}(y)$ es discreto es abierto en $Y$, y el morfismo restringido
$f^{-1}(U)\to U$ es finito. En particular, todo morfismo propio
cuasifinito $X\to Y$ es finito.
\end{corollary}

\begin{proof}
El complementario de $U$ en $Y$ es la imagen por $f$ de $X-X'$, que
es cerrado en $X$ por \sref{III.4.4.10}; como $f$ es una aplicación
cerrada, $U$ es abierto. Además, \sref[II]{II.6.2.2} muestra que
$f^{-1}(y)$ es finita para todo $y\in U$. Como el morfismo restringido
$f^{-1}(U)\to U$ es propio \sref[II]{II.5.4.1}, es finito por
\sref{III.4.4.2}.
\end{proof}

\begin{remark}[4.4.12]
\label{III.4.4.12}
\emph{(i)}
Como se anunció en \sref[II]{II.6.2.7}, demostraremos en el capítulo~V
que, si $Y$ es localmente noetheriano, todo morfismo separado
cuasifinito $f:X\to Y$ es cuasiafín, luego cuasiproyectivo. Se seguirá
que la conclusión del teorema principal \sref{III.4.4.3} sigue siendo
válida cuando solo se supone que $f$ es separado y de tipo finito.
En efecto, \sref{III.4.4.10} muestra que $X'$ es abierto en $X$; y
como $X$ es localmente noetheriano, la restricción de $f$ a $X'$ sigue
siendo de tipo finito \sref[I]{I.6.3.5}, luego cuasifinita por la
definición de $X'$, y claramente separada. Podemos, por tanto, aplicar
\sref{III.4.4.3} a esta restricción.

\emph{(ii)}
En el capítulo~IV daremos una demostración más elemental de
\sref{III.4.4.10}, utilizando la teoría de la dimensión.
\end{remark}


\subsection{Completaciones de módulos de homomorfismos}
\label{subsection:III.4.5}

\begin{proposition}[4.5.1]
\label{III.4.5.1}
Sea $A$ un anillo noetheriano, sea $\mathfrak J$ un ideal de $A$, sea
$X$ un $A$-esquema de tipo finito, y sean $\sh{F}$ y $\sh{G}$ dos
$\sh{O}_X$-módulos coherentes cuyos soportes tienen intersección
propia sobre $Y=\Spec(A)$ \sref[II]{II.5.4.10}. Entonces, para todo
entero $n\geq0$,
\[
  \Ext^n_{\sh{O}_X}(X;\sh{F},\sh{G})
\]
es un $A$-módulo finito, y su completación separada para la topología
$\mathfrak J$-preádica se identifica canónicamente, con la notación de
\sref{III.4.1.7}, con
\[
  \Ext^n_{\sh{O}_{\hat{X}}}
    (\hat{X};\hat{\sh{F}},\hat{\sh{G}}).
\]
\end{proposition}

\begin{proof}
Sabemos por (T, 4.2) que existe una sucesión espectral birregular
$\mathrm{E}(\sh{F},\sh{G})$ cuyo término límite es
\[
  \Ext_{\sh{O}_X}(X;\sh{F},\sh{G})
\]
y cuyos términos $\mathrm{E}_2$ son
\[
  \mathrm{E}_2^{pq}
  =\HH^p\bigl(X,\shExt^q_{\sh{O}_X}(\sh{F},\sh{G})\bigr).
\]
El $\sh{O}_X$-módulo
$\shExt^q_{\sh{O}_X}(\sh{F},\sh{G})$ es coherente
\sref[0]{0.12.3.3}; su soporte está contenido en la intersección de
los soportes de $\sh{F}$ y $\sh{G}$ (T, 4.2.2) y, en consecuencia, es
propio sobre $Y$ \sref[II]{II.5.4.10}. De \sref{III.3.2.4} se sigue
que los $\mathrm{E}_2^{pq}$ son $A$-módulos finitos, y después, de
\sref[0]{0.11.1.8}, que lo mismo ocurre con todo término
$\mathrm{E}_r^{pq}$ de la sucesión espectral y con su término límite.

Por otra parte, si $i:\hat{X}\to X$ es el morfismo canónico, entonces
$\hat{\sh{F}}$ y $\hat{\sh{G}}$ se identifican canónicamente con
$i^*(\sh{F})$ e $i^*(\sh{G})$, y $i$ es plano
\sref[I]{I.10.8.8}, \sref[I]{I.10.8.9}. Para todo $q\geq0$ existe,
por tanto, un $i$-morfismo canónico
\[
  u_q:\shExt^q_{\sh{O}_X}(\sh{F},\sh{G})
  \longrightarrow
  \shExt^q_{\sh{O}_{\hat{X}}}
    (\hat{\sh{F}},\hat{\sh{G}})
\]
\sref[0]{0.12.3.4}, y el $\sh{O}_{\hat{X}}$-homomorfismo
correspondiente
\[
  u_q^\#:
  i^*\bigl(\shExt^q_{\sh{O}_X}(\sh{F},\sh{G})\bigr)
  \longrightarrow
  \shExt^q_{\sh{O}_{\hat{X}}}
    (\hat{\sh{F}},\hat{\sh{G}})
\]
es un isomorfismo \sref[0]{0.12.3.5}. En otros términos,
\sref[I]{I.10.8.8} identifica
\[
  \shExt^q_{\sh{O}_{\hat{X}}}
    (\hat{\sh{F}},\hat{\sh{G}})
\]
canónicamente con la completación de
$\shExt^q_{\sh{O}_X}(\sh{F},\sh{G})$, con la notación de
\sref{III.4.1.7}.

El teorema de comparación \sref{III.4.1.10} muestra ahora, para todo
$p\geq0$, que
\[
  \HH^p\bigl(\hat{X},
    \shExt^q_{\sh{O}_{\hat{X}}}
      (\hat{\sh{F}},\hat{\sh{G}})\bigr)
\]
se identifica canónicamente con la completación separada, para la
topología $\mathfrak J$-preádica, de
\[
  \HH^p\bigl(X,\shExt^q_{\sh{O}_X}(\sh{F},\sh{G})\bigr).
\]

\oldpage[III]{139}

Denotemos por $\mathrm{E}(\hat{\sh{F}},\hat{\sh{G}})$ la sucesión
espectral birregular de (T, 4.2) asociada a $\hat{\sh{F}}$ y
$\hat{\sh{G}}$. Si $\hat{A}$ es la completación separada de $A$ para
la topología $\mathfrak J$-preádica, entonces, salvo isomorfismo
canónico,
\[
  \mathrm{E}_2^{pq}(\hat{\sh{F}},\hat{\sh{G}})
  =
  \mathrm{E}_2^{pq}(\sh{F},\sh{G})\otimes_A\hat{A}
\]
\sref[0]{0.7.3.3}.

El morfismo plano $i$ define un homomorfismo canónico de sucesiones
espectrales
\[
  \varphi:\mathrm{E}(\sh{F},\sh{G})
  \longrightarrow
  \mathrm{E}(\hat{\sh{F}},\hat{\sh{G}})
  =
  \mathrm{E}(i^*(\sh{F}),i^*(\sh{G})).
\]
Sobre los términos $\mathrm{E}_2$ y sobre el término límite,
respectivamente, este es el homomorfismo
\[
  \varphi_2^{pq}:
  \HH^p\bigl(X,\shExt^q_{\sh{O}_X}(\sh{F},\sh{G})\bigr)
  \longrightarrow
  \HH^p\bigl(\hat{X},
    \shExt^q_{\sh{O}_{\hat{X}}}
      (\hat{\sh{F}},\hat{\sh{G}})\bigr)
\]
y el homomorfismo
\[
  \varphi^n:
  \Ext^n_{\sh{O}_X}(X;\sh{F},\sh{G})
  \longrightarrow
  \Ext^n_{\sh{O}_{\hat{X}}}
    (\hat{X};\hat{\sh{F}},\hat{\sh{G}}),
\]
deducidos de $u_q$ y $u_0$, respectivamente, por funtorialidad
\sref[0]{0.12.3.4}.

Tras tensorizar con $\hat{A}$, los $\varphi_r^{pq}$ y
$\varphi^n$ dan homomorfismos de $\hat{A}$-módulos
\[
  \psi_r^{pq}:
  \mathrm{E}_r^{pq}(\sh{F},\sh{G})\otimes_A\hat{A}
  \longrightarrow
  \mathrm{E}_r^{pq}(\hat{\sh{F}},\hat{\sh{G}})
\]
y
\[
  \psi^n:
  \Ext^n_{\sh{O}_X}(X;\sh{F},\sh{G})\otimes_A\hat{A}
  \longrightarrow
  \Ext^n_{\sh{O}_{\hat{X}}}
    (\hat{X};\hat{\sh{F}},\hat{\sh{G}}).
\]
Como $\hat{A}$ es un $A$-módulo plano \sref[0]{0.7.3.3}, los
$\hat{A}$-módulos
\[
  \mathrm{E}_r^{pq}(\sh{F},\sh{G})\otimes_A\hat{A}
\]
forman una sucesión espectral birregular cuyo término límite es
\[
  \Ext^n_{\sh{O}_X}(X;\sh{F},\sh{G})\otimes_A\hat{A},
\]
y los $\psi_r^{pq}$ y $\psi^n$ forman un morfismo de sucesiones
espectrales. Como los $\psi_2^{pq}$ son isomorfismos, lo mismo ocurre
con los $\psi^n$ \sref[0]{0.11.1.5}.
\end{proof}

\begin{corollary}[4.5.2]
\label{III.4.5.2}
Bajo las hipótesis de \sref{III.4.5.1}, supongamos además que $A$ es
un anillo noetheriano $\mathfrak J$-ádico. Entonces, para todo entero
$n\geq0$,
\[
  \Ext^n_{\sh{O}_X}(X;\sh{F},\sh{G})
  \isoto
  \Ext^n_{\sh{O}_{\hat{X}}}
    (\hat{X};\hat{\sh{F}},\hat{\sh{G}})
\]
canónicamente.
\end{corollary}

\begin{proof}
El $A$-módulo finito
$\Ext^n_{\sh{O}_X}(X;\sh{F},\sh{G})$ está separado y completo para la
topología $\mathfrak J$-preádica \sref[0]{0.7.3.6}.
\end{proof}

El caso particular $n=0$ de \sref{III.4.5.1} puede enunciarse como
sigue.

\begin{corollary}[4.5.3]
\label{III.4.5.3}
Bajo las hipótesis de \sref{III.4.5.1}, para todo homomorfismo
$u:\sh{F}\to\sh{G}$, denotemos por
$\hat{u}:\hat{\sh{F}}\to\hat{\sh{G}}$ el homomorfismo completado
\sref[I]{I.10.8.4}. Existe un isomorfismo canónico
\[
\label{III.4.5.3.1}
  \left(\Hom_{\sh{O}_X}(\sh{F},\sh{G})\right)^\wedge
  \isoto
  \Hom_{\sh{O}_{\hat{X}}}(\hat{\sh{F}},\hat{\sh{G}}),
  \tag{4.5.3.1}
\]
donde el miembro izquierdo es la completación separada, para la
topología $\mathfrak J$-preádica, del $A$-módulo
$\Hom_{\sh{O}_X}(\sh{F},\sh{G})$. Este isomorfismo se obtiene pasando
a las completaciones separadas en el homomorfismo
$u\mapsto\hat{u}$.
\end{corollary}


\subsection{Relaciones entre morfismos formales y morfismos ordinarios}
\label{subsection:III.4.6}

\begin{proposition}[4.6.1]
\label{III.4.6.1}
Sean $Y$ un esquema localmente noetheriano, $f:X\to Y$ un
morfismo propio, $\sh{F}$ un $\sh{O}_X$-módulo coherente que es
plano sobre $Y$, e $y\in Y$. Supongamos que, para algún entero
$n$,
\[
  \HH^n\bigl(f^{-1}(y),
    \sh{F}\otimes_{\sh{O}_Y}\kres(y)\bigr)=0.
\]
Entonces existe un entorno $U$ de $y$ en $Y$ tal que
\[
  \RR^n f_*(\sh{F})|U=0,
\]
y, para todo entero $p\geq0$, el homomorfismo canónico
\[
  \bigl(\RR^{n-1}f_*(\sh{F})\bigr)_y
  \longrightarrow
  \HH^{n-1}\bigl(f^{-1}(y),
    \sh{F}\otimes_{\sh{O}_Y}
      (\sh{O}_y/\mathfrak m_y^{p+1})\bigr)
\]
\eqref{III.4.2.1.1} es sobreyectivo.
\end{proposition}

\begin{proof}
El $\sh{O}_Y$-módulo $\RR^n f_*(\sh{F})$ es coherente
\sref{III.3.2.1}; por tanto, la primera afirmación se seguirá una vez que probemos
\[
  \bigl(\RR^n f_*(\sh{F})\bigr)_y=0
\]
\sref[0]{0.5.2.2}. Por \sref{III.4.2.1}, basta probar que
\[
  \HH^n\bigl(f^{-1}(y),
    \sh{F}\otimes_{\sh{O}_Y}
      (\sh{O}_y/\mathfrak m_y^{p+1})\bigr)=0
\]
para todo $p$. Esto es cierto por hipótesis para $p=0$; procedemos por
inducción sobre $p$.

Pongamos
\[
  X_p=X\times_Y
    \Spec(\sh{O}_y/\mathfrak m_y^{p+1}).
\]
Entonces $X_{p-1}$ es un subesquema cerrado de $X_p$ con el mismo
espacio subyacente \sref[I]{I.3.6.1}. La hipótesis de inducción
\[
  \HH^n\bigl(X_{p-1},
    \sh{F}\otimes_{\sh{O}_Y}
      (\sh{O}_y/\mathfrak m_y^p)\bigr)=0
\]
da, por tanto,
\[
  \HH^n\bigl(X_p,
    \sh{F}\otimes_{\sh{O}_Y}
      (\sh{O}_y/\mathfrak m_y^p)\bigr)=0.
\]
Por otra parte, la sucesión exacta de $\sh{O}_{X_p}$-módulos
\[
  0\longrightarrow
  \mathfrak m_y^p\sh{F}/\mathfrak m_y^{p+1}\sh{F}
  \longrightarrow
  \sh{F}/\mathfrak m_y^{p+1}\sh{F}
  \longrightarrow
  \sh{F}/\mathfrak m_y^p\sh{F}
  \longrightarrow0
\]
da la sucesión exacta de cohomología
\[
  \HH^n\bigl(X_p,
    \mathfrak m_y^p\sh{F}/\mathfrak m_y^{p+1}\sh{F}\bigr)
  \longrightarrow
  \HH^n\bigl(X_p,\sh{F}/\mathfrak m_y^{p+1}\sh{F}\bigr)
  \longrightarrow
  \HH^n\bigl(X_p,\sh{F}/\mathfrak m_y^p\sh{F}\bigr).
\]
Basta, pues, demostrar que
\[
\label{III.4.6.1.1}
  \HH^n\bigl(X_p,
    \mathfrak m_y^p\sh{F}/\mathfrak m_y^{p+1}\sh{F}\bigr)=0.
  \tag{4.6.1.1}
\]
En efecto, el término central será entonces un submódulo del último término,
que es cero por la hipótesis de inducción.

\oldpage[III]{140}

La fibra
\[
  Z=f^{-1}(y)=X\times_Y\Spec(\kres(y))
\]
es un subesquema cerrado de $X_p$, y
$\mathfrak m_y^p\sh{F}/\mathfrak m_y^{p+1}\sh{F}$ es anulado por
$\mathfrak m_y$. Por tanto, puede considerarse como un
$\sh{O}_Z$-módulo, y
\[
  \HH^n\bigl(Z,
    \mathfrak m_y^p\sh{F}/\mathfrak m_y^{p+1}\sh{F}\bigr)
  =
  \HH^n\bigl(X_p,
    \mathfrak m_y^p\sh{F}/\mathfrak m_y^{p+1}\sh{F}\bigr).
\]
Probaremos ahora que el $\sh{O}_Z$-homomorfismo canónico
\[
\label{III.4.6.1.2}
  (\sh{F}/\mathfrak m_y\sh{F})
    \otimes_{\kres(y)}
      (\mathfrak m_y^p/\mathfrak m_y^{p+1})
  \longrightarrow
  \mathfrak m_y^p\sh{F}/\mathfrak m_y^{p+1}\sh{F}
  \tag{4.6.1.2}
\]
es biyectivo. Una vez establecido esto, el hecho de que
$\mathfrak m_y^p/\mathfrak m_y^{p+1}$ sea un
$\kres(y)$-módulo libre da
\[
\begin{split}
  \HH^n\bigl(Z,
    \mathfrak m_y^p\sh{F}/\mathfrak m_y^{p+1}\sh{F}\bigr)
  &=
  \HH^n\bigl(Z,\sh{F}/\mathfrak m_y\sh{F}\bigr)
    \otimes_{\kres(y)}
      (\mathfrak m_y^p/\mathfrak m_y^{p+1})\\
  &=0
\end{split}
\]
\sref[0]{0.12.2.3}, puesto que
$\HH^n(Z,\sh{F}/\mathfrak m_y\sh{F})=0$ por hipótesis. Esto prueba
\eqref{III.4.6.1.1}.

Para terminar la demostración de la primera afirmación, solo queda demostrar
que \eqref{III.4.6.1.2} es biyectivo. La cuestión es puntual sobre
$X$, y $\sh{F}_x$ es un $\sh{O}_y$-módulo plano por hipótesis para
todo $x\in f^{-1}(y)$. Basta, pues, aplicar
\sref[0]{0.10.2.1}, c), puesto que
$\sh{F}_x/\mathfrak m_y\sh{F}_x$ es plano sobre el cuerpo
$\kres(y)=\sh{O}_y/\mathfrak m_y$.

Para la segunda afirmación de \sref{III.4.6.1}, nos reducimos de inmediato, como
en \sref{III.4.2.1}, al caso en que $Y$ es afín e $y$ es
cerrado. El mismo argumento demuestra, para todo $k>0$, que
\[
\label{III.4.6.1.3}
  \HH^n\bigl(X_{k+1},
    \mathfrak m_y^k\sh{F}/
      \mathfrak m_y^{k+p+1}\sh{F}\bigr)=0.
  \tag{4.6.1.3}
\]

\oldpage[III]{141}

De \sref{III.4.2.1} se sigue que
\[
\label{III.4.6.1.4}
  \bigl(\RR^n f_*(\mathfrak m_y^k\sh{F})\bigr)_y=0.
  \tag{4.6.1.4}
\]
La sucesión exacta de cohomología da ahora la exactitud de
\[
  \bigl(\RR^{n-1}f_*(\sh{F})\bigr)_y
  \longrightarrow
  \bigl(\RR^{n-1}f_*
    (\sh{F}/\mathfrak m_y^p\sh{F})\bigr)_y
  \longrightarrow
  \bigl(\RR^n f_*(\mathfrak m_y^p\sh{F})\bigr)_y=0.
\]
Como $y$ es cerrado e $Y$ es afín, tenemos
\sref{III.1.4.11}
\[
  \RR^{n-1}f_*(\sh{F}/\mathfrak m_y^p\sh{F})
  =
  \bigl(\HH^{n-1}
    (X,\sh{F}/\mathfrak m_y^p\sh{F})\bigr)^{\sim}
  =
  \bigl(\HH^{n-1}
    (f^{-1}(y),\sh{F}/\mathfrak m_y^p\sh{F})\bigr)^{\sim}
\]
(G, II, 4.9.1). El módulo
\[
  \HH^{n-1}
    (f^{-1}(y),\sh{F}/\mathfrak m_y^p\sh{F})
\]
es un $\sh{O}_y/\mathfrak m_y^p$-módulo, y por tanto
\[
  \bigl(\RR^{n-1}f_*
    (\sh{F}/\mathfrak m_y^p\sh{F})\bigr)_y
  =
  \HH^{n-1}
    (f^{-1}(y),\sh{F}/\mathfrak m_y^p\sh{F}).
\]
Esto termina la demostración.
\end{proof}

\begin{corollary}[4.6.2]
\label{III.4.6.2}
Sean $Y$ un esquema localmente noetheriano, $f:X\to Y$ un
morfismo propio y plano, $\sh{F}$ y $\sh{G}$ dos
$\sh{O}_X$-módulos localmente libres, e $y$ un punto de $Y$. Pongamos
\[
  X_y=f^{-1}(y)=X\times_Y\Spec(\kres(y)),\qquad
  \sh{F}_y=\sh{F}\otimes_{\sh{O}_Y}\kres(y),\qquad
  \sh{G}_y=\sh{G}\otimes_{\sh{O}_Y}\kres(y),
\]
y supongamos que
\[
\label{III.4.6.2.1}
  \HH^1\bigl(
    X_y,\shHom_{\sh{O}_{X_y}}(\sh{F}_y,\sh{G}_y)
  \bigr)=0.
  \tag{4.6.2.1}
\]
Entonces, para todo homomorfismo $u_0:\sh{F}_y\to\sh{G}_y$, existen
un entorno abierto $U$ de $y$ en $Y$ y un homomorfismo
\[
  u:\sh{F}|_{f^{-1}(U)}\longrightarrow
    \sh{G}|_{f^{-1}(U)}
\]
tal que $u_0=u\otimes1$.
\end{corollary}

\begin{proof}
La hipótesis nos permite aplicar \sref{III.4.6.1}, con $n=1$
y $p=0$, al $\sh{O}_X$-módulo coherente
\[
  \sh{H}=\shHom_{\sh{O}_X}(\sh{F},\sh{G}).
\]
En efecto, $\sh{H}$ es localmente libre y, por tanto, con mayor razón, plano sobre
$Y$, mientras que el $\sh{O}_{X_y}$-módulo
$\sh{H}\otimes_{\sh{O}_Y}\kres(y)$ se identifica con
$\shHom_{\sh{O}_{X_y}}(\sh{F}_y,\sh{G}_y)$
\sref[0]{0.6.2.2}. Podemos suponer que $Y=\Spec(A)$ es afín.
Entonces \sref{III.1.4.11} da
\[
  \RR^0f_*(\sh{H})
  =
  \bigl(\Hom_{\sh{O}_X}(\sh{F},\sh{G})\bigr)^{\sim},
\]
y, en consecuencia,
\[
  \bigl(\RR^0f_*(\sh{H})\bigr)_y
  =
  \Hom_{\sh{O}_X}(\sh{F},\sh{G})
    \otimes_A\sh{O}_y.
\]
Por \sref{III.4.6.1}, el homomorfismo canónico
\[
  \Hom_{\sh{O}_X}(\sh{F},\sh{G})
    \otimes_A\sh{O}_y
  \longrightarrow
  \Hom_{\sh{O}_{X_y}}(\sh{F}_y,\sh{G}_y)
\]
es sobreyectivo. Esto prueba el corolario, puesto que todo elemento del
módulo de la izquierda puede escribirse como $u\otimes(1/s)$, donde
$s\in A\setminus\mathfrak m_y$.
\end{proof}

\begin{corollary}[4.6.3]
\label{III.4.6.3}
Bajo las hipótesis de \sref{III.4.6.2}, si $u_0$ es inyectivo
(respectivamente sobreyectivo, biyectivo), se puede elegir $u$
inyectivo (respectivamente sobreyectivo, biyectivo).
\end{corollary}

\begin{proof}
Podemos restringirnos al caso $U=Y$. Basta probar
que, si $u_0$ es inyectivo (respectivamente sobreyectivo), entonces
\[
  \Ker u_x=0
  \quad\text{(respectivamente }\Coker u_x=0\text{)}
\]
para todo $x\in f^{-1}(y)$. En efecto, $\Ker u$ y $\Coker u$ son
$\sh{O}_X$-módulos coherentes \sref[0]{0.5.3.4}; por tanto, existe un
entorno $V$ de $f^{-1}(y)$ en $X$ sobre el cual la restricción de
$\Ker u$ (respectivamente $\Coker u$) es cero \sref[0]{0.5.2.2}.
Como $f$ es cerrado, existe un entorno abierto
$U'\subset U$ de $y$ tal que $f^{-1}(U')\subset V$, lo que prueba
la afirmación después de sustituir $U$ por $U'$.

\oldpage[III]{142}

Por hipótesis,
\[
  u_x\otimes1:
  \sh{F}_x\otimes_{\sh{O}_y}\kres(y)
  \longrightarrow
  \sh{G}_x\otimes_{\sh{O}_y}\kres(y)
\]
es inyectivo (respectivamente sobreyectivo), $\sh{F}_x$ y $\sh{G}_x$
son $\sh{O}_x$-módulos libres finitos, y $\sh{O}_x$ es plano sobre
$\sh{O}_y$. Si $u_x\otimes1$ es inyectivo, la inyectividad de $u_x$
se sigue de \sref[0]{0.10.2.4}. Si $u_x\otimes1$ es sobreyectivo,
entonces, con mayor razón, el homomorfismo inducido
\[
  \sh{F}_x/\mathfrak m_x\sh{F}_x
  \longrightarrow
  \sh{G}_x/\mathfrak m_x\sh{G}_x
\]
es sobreyectivo. Como $\sh{G}_x$ es un $\sh{O}_x$-módulo finito y
$\sh{O}_x$ es un anillo local de ideal maximal $\mathfrak m_x$, la
conclusión se sigue del lema de Nakayama (Bourbaki,
\emph{Alg\`ebre}, cap.~VIII, \S\,6, n.º~3, cor.~4 de la prop.~6).
\end{proof}

\begin{corollary}[4.6.4]
\label{III.4.6.4}
Sean $Y$ un esquema localmente noetheriano, $f:X\to Y$ un
morfismo propio y plano, e $y$ un punto de $Y$, y pongamos
$X_y=X\times_Y\Spec(\kres(y))$. Sea $\sh{E}_0$ un
$\sh{O}_{X_y}$-módulo localmente libre tal que
\[
\label{III.4.6.4.1}
  \HH^1\bigl(
    X_y,\shHom_{\sh{O}_{X_y}}(\sh{E}_0,\sh{E}_0)
  \bigr)=0.
  \tag{4.6.4.1}
\]
Sean $\sh{F}$ y $\sh{G}$ dos $\sh{O}_X$-módulos localmente libres
tales que $\sh{F}_y$ y $\sh{G}_y$, con la notación de
\sref{III.4.6.2}, son isomorfos a $\sh{E}_0$. Entonces existe un
entorno abierto $U$ de $y$ tal que
$\sh{F}|_{f^{-1}(U)}$ y $\sh{G}|_{f^{-1}(U)}$ son isomorfos.
\end{corollary}

\begin{corollary}[4.6.5]
\label{III.4.6.5}
Bajo las hipótesis de \sref{III.4.6.4} sobre $f$, $X$ e $Y$,
supongamos que
\[
  \HH^1(X_y,\sh{O}_{X_y})=0.
\]
Si $\sh{F}$ y $\sh{G}$ son dos $\sh{O}_X$-módulos invertibles tales
que $\sh{F}_y$ y $\sh{G}_y$ son isomorfos, entonces existe un entorno
abierto $U$ de $y$ tal que $\sh{F}|_{f^{-1}(U)}$ y
$\sh{G}|_{f^{-1}(U)}$ son isomorfos.
\end{corollary}

\begin{proof}
Basta aplicar \sref{III.4.6.4} a los módulos
$\sh{F}^{-1}\otimes\sh{G}$ y $\sh{O}_X$.
\end{proof}

\begin{remarks}[4.6.6]
\label{III.4.6.6}
\begin{enumerate}
\item[(i)] Utilizando \sref{III.4.6.5}, estableceremos en el capítulo~V
  la clasificación de los haces invertibles sobre un fibrado proyectivo
  anunciada en \sref[II]{II.4.2.7}.
\item[(ii)] El resultado de \sref{III.4.6.1} aparecerá en \S\,7 como
  consecuencia de proposiciones más generales.
\end{enumerate}
\end{remarks}

\begin{proposition}[4.6.7]
\label{III.4.6.7}
Sea $Z$ un esquema localmente noetheriano, y sean $X$ e $Y$
dos $Z$-esquemas cuyos morfismos estructurales
\[
  g:X\longrightarrow Z,\qquad h:Y\longrightarrow Z
\]
son propios. Sea $f:X\to Y$ un $Z$-morfismo, sea $z$ un punto
de $Z$, y pongamos
\[
  f_z=f\times_Z1:
  X\times_Z\Spec(\kres(z))
  \longrightarrow
  Y\times_Z\Spec(\kres(z)).
\]
\begin{enumerate}
\item[(i)] Si $f_z$ es finito (respectivamente una inmersión cerrada),
  existe un entorno abierto $U$ de $z$ tal que el morfismo
  \[
    g^{-1}(U)\longrightarrow h^{-1}(U)
  \]
  inducido por $f$ es finito (respectivamente una inmersión cerrada).
\item[(ii)] Supongamos además que $g$ es plano. Si $f_z$ es un
  isomorfismo, existe un entorno abierto $U$ de $z$ tal que
  el morfismo $g^{-1}(U)\to h^{-1}(U)$ inducido por $f$ es un
  isomorfismo.
\end{enumerate}
\end{proposition}

\begin{proof}
En ambos casos basta probar que, para todo
$y\in h^{-1}(z)$, existe un entorno $V_y$ de $y$ tal que
la restricción
\[
  f^{-1}(V_y)\longrightarrow V_y
\]
de $f$ es finita (respectivamente una inmersión cerrada, un isomorfismo).
En efecto, si $V$ es la unión de los $V_y$, entonces
$f^{-1}(V)\to V$ es finito (respectivamente una inmersión cerrada, un
isomorfismo) \sref[II]{II.6.1.1}, \sref[I]{I.4.2.4}. Como $h$ es
cerrado, existe un entorno abierto $U$ de $z$ tal que
$h^{-1}(U)\subset V$, y la proposición se sigue.

\noindent\textit{(i)} Obsérvese primero que $f$ es propio
\sref[II]{II.5.4.3}. Si
$f_z$ es finito, el entorno requerido $V_y$, para cada
$y\in h^{-1}(z)$, existe por \sref{III.4.4.11}.

\oldpage[III]{143}

Para tratar el caso en que $f_z$ es una inmersión cerrada, podemos
suponer ya, por tanto, que $f$ es finito. Así,
$X=\Spec(\sh{B})$, donde $\sh{B}$ es una
$\sh{O}_Y$-álgebra coherente, y $f$ corresponde \sref[II]{II.1.2.7} al
homomorfismo canónico
\[
  u:\sh{O}_Y\longrightarrow\sh{B}.
\]
Si, para todo $y\in h^{-1}(z)$, el homomorfismo
$u_y:\sh{O}_y\to\sh{B}_y$ es sobreyectivo, entonces $u$ es sobreyectivo en
algún entorno $V_y$ de $y$, puesto que $\Coker u$ es coherente
\sref[0]{0.5.3.4}, \sref[0]{0.5.2.2}. Ahora el morfismo finito
$f_z$ corresponde a
\[
  v=u\otimes1:
  \sh{O}_Y\otimes_{\sh{O}_Z}\kres(z)
  \longrightarrow
  \sh{B}\otimes_{\sh{O}_Z}\kres(z),
\]
y la hipótesis de que $f_z$ es una inmersión cerrada implica que
\[
  u_y\otimes1:
  \sh{O}_y\otimes_{\sh{O}_z}\kres(z)
  \longrightarrow
  \sh{B}_y\otimes_{\sh{O}_z}\kres(z)
\]
es sobreyectivo. Como $\sh{B}_y$ es un $\sh{O}_y$-módulo finito y
$\sh{O}_y$ es un anillo local noetheriano, la conclusión se sigue, como
en \sref{III.4.6.3}, del lema de Nakayama.

\noindent\textit{(ii)} El mismo razonamiento muestra que ahora
basta probar
que $u_y$ es biyectivo, sabiendo que $u_y\otimes1$ es biyectivo.
La inyectividad requerida se sigue del lema siguiente; la sobreyectividad
se sigue del argumento precedente.
\end{proof}

\begin{lemma}[4.6.7.1]
\label{III.4.6.7.1}
Sean $A$ y $B$ anillos locales noetherianos, sea
$\rho:A\to B$ un homomorfismo local, y sea $u:N\to M$ un
homomorfismo de $B$-módulos. Supongamos que $M$ es plano sobre $A$,
que $N$ es un $B$-módulo finito, y que
\[
  u\otimes1:
  N\otimes_A k\longrightarrow M\otimes_A k,
\]
donde $k$ es el cuerpo residual de $A$, es inyectivo. Entonces $N$ es
plano sobre $A$ y $u$ es inyectivo.
\end{lemma}

\begin{proof}
Para probar la primera afirmación, debemos demostrar que, para todo par de
$A$-módulos finitos $P,Q$ y todo $A$-homomorfismo inyectivo
$v:P\to Q$, el homomorfismo
\[
  1_N\otimes v:N\otimes_A P\longrightarrow N\otimes_A Q
\]
es inyectivo. Consideremos el diagrama conmutativo
\[
  \xymatrix{
    N\otimes_A P
      \ar[r]^-{1_N\otimes v}
      \ar[d]_{u\otimes1_P}
      &
    N\otimes_A Q
      \ar[d]^{u\otimes1_Q}
    \\
    M\otimes_A P
      \ar[r]_{1_M\otimes v}
      &
    M\otimes_A Q .
  }
\]
Como $1_M\otimes v$ es inyectivo por la planitud de $M$, basta
probar que $u\otimes1_P$ es inyectivo.

Sea $\mathfrak m$ el ideal maximal de $A$. La
filtración $\mathfrak m$-ádica sobre el $A$-módulo $N\otimes_A P$ es
también su filtración $\mathfrak mB$-ádica como $B$-módulo. La
topología definida por esta filtración es separada: $B$ es
noetheriano, $\mathfrak mB$ está contenido en el radical de $B$, y
$N\otimes_A P$ es un $B$-módulo finito, puesto que $N$ es finito sobre $B$
y $P$ es finito sobre $A$ \sref[0]{0.7.3.5}. Basta, por tanto,
probar que
\[
  \operatorname{gr}(u\otimes1_P):
  \operatorname{gr}_\bullet(N\otimes_A P)
  \longrightarrow
  \operatorname{gr}_\bullet(M\otimes_A P)
\]
es inyectivo, donde los módulos graduados asociados se toman con
respecto a las filtraciones $\mathfrak m$-ádicas (Bourbaki,
\emph{Alg\`ebre commutative}, cap.~III, \S\,2, n.º~8, cor.~1 del
teor.~1).

Como $M$ es plano sobre $A$, los homomorfismos
\[
  M\otimes_A(\mathfrak m^nP)
  \longrightarrow
  \mathfrak m^n(M\otimes_A P)
\]
son biyectivos. En consecuencia, el homomorfismo canónico
\[
  \vphi_M:
  \operatorname{gr}_0(M)\otimes_A\operatorname{gr}_\bullet(P)
  \longrightarrow
  \operatorname{gr}_\bullet(M\otimes_A P)
\]
es biyectivo.

\oldpage[III]{144}

Tenemos el diagrama conmutativo
\[
  \xymatrix{
    \operatorname{gr}_0(N)\otimes_A\operatorname{gr}_\bullet(P)
      \ar[r]^-{\operatorname{gr}_0(u)\otimes1}
      \ar[d]_{\vphi_N}
      &
    \operatorname{gr}_0(M)\otimes_A\operatorname{gr}_\bullet(P)
      \ar[d]^{\vphi_M}
    \\
    \operatorname{gr}_\bullet(N\otimes_A P)
      \ar[r]_{\operatorname{gr}(u\otimes1)}
      &
    \operatorname{gr}_\bullet(M\otimes_A P).
  }
\]
Aquí $\vphi_M$ es biyectivo y $\vphi_N$ es sobreyectivo.
Además, $\operatorname{gr}_0(u)$ es inyectivo por hipótesis, y
\[
\begin{split}
  \operatorname{gr}_0(N)\otimes_A\operatorname{gr}_\bullet(P)
    &=
  \operatorname{gr}_0(N)\otimes_k\operatorname{gr}_\bullet(P),\\
  \operatorname{gr}_0(M)\otimes_A\operatorname{gr}_\bullet(P)
    &=
  \operatorname{gr}_0(M)\otimes_k\operatorname{gr}_\bullet(P).
\end{split}
\]
Así, $\operatorname{gr}_0(u)\otimes1$ también es inyectivo. Del
diagrama se sigue que $\operatorname{gr}(u\otimes1_P)$ es
inyectivo, lo que prueba la primera afirmación. La segunda se sigue del
mismo argumento tomando $P=A$.
\end{proof}

\begin{proposition}[4.6.8]
\label{III.4.6.8}
Sea $Z$ un esquema localmente noetheriano, y sean $X$ e $Y$
dos $Z$-esquemas cuyos morfismos estructurales
$g:X\to Z$ y $h:Y\to Z$ son propios. Sea $Z'$ un subconjunto cerrado
de $Z$, pongamos
\[
  X'=g^{-1}(Z'),\qquad Y'=h^{-1}(Z'),
\]
y denotemos por
\[
  \widehat X=X_{/X'},\qquad
  \widehat Y=Y_{/Y'},\qquad
  \widehat Z=Z_{/Z'}
\]
las completaciones formales a lo largo de estos subconjuntos cerrados. Sea $f:X\to Y$
un $Z$-morfismo y sea
$\widehat f:\widehat X\to\widehat Y$ su extensión a las completaciones
formales. Para que $\widehat f$ sea un isomorfismo
(respectivamente una inmersión cerrada), es necesario y suficiente
que exista un entorno abierto $U$ de $Z'$ tal que el
morfismo
\[
  g^{-1}(U)\longrightarrow h^{-1}(U)
\]
inducido por $f$ sea un isomorfismo (respectivamente una inmersión cerrada).
\end{proposition}

\begin{proof}
La suficiencia de la condición es inmediata
\sref[I]{I.10.14.7}. Para probar su necesidad, basta, por el
mismo razonamiento que en \sref{III.4.6.7}, demostrar que para todo
$y\in Y'$ existe un entorno abierto $V_y$ de $y$ tal que
$f^{-1}(V_y)\to V_y$ sea un isomorfismo (respectivamente una
inmersión cerrada). Quedamos así reducidos al caso
\[
  Y=Z=\Spec(A),
\]
donde $A$ es noetheriano.

Por hipótesis \sref[I]{I.10.9.1}, \sref[I]{I.10.14.2}, la fibra
$f^{-1}(y)$ consta de un único punto para todo $y\in Y'$. Como
$f$ es propio \sref[II]{II.5.4.3}, \sref{III.4.4.11} da un entorno
abierto $U$ de $y$ para el cual $f^{-1}(U)\to U$ es finito. Podemos,
por tanto, suponer que $f$ es finito, de modo que
\[
  X=\Spec(B),
\]
donde $B$ es una $A$-álgebra finita.

Si $Y'=V(\mathfrak J)$, entonces
\[
  \widehat Y=\Spf(\widehat A),
  \qquad
  \widehat X=\Spf(\widehat B),
\]
donde $\widehat A$ es la completación separada de $A$ para la
topología $\mathfrak J$-ádica, y $\widehat B$ es la completación
separada de $B$ para la topología $\mathfrak JB$-ádica, o,
equivalentemente, la completación separada del $A$-módulo $B$ para
la topología $\mathfrak J$-ádica. Además, $\widehat f$
corresponde a la extensión continua
\[
  \widehat{\vphi}:\widehat A\longrightarrow\widehat B
\]
del homomorfismo canónico de anillos $\vphi:A\to B$, y la
hipótesis es que $\widehat{\vphi}$ es sobreyectivo (respectivamente
biyectivo) \sref[I]{I.10.14.2}. Pero $\widehat{\vphi}$ es también la
extensión continua de $\vphi$ considerado como un homomorfismo de
$A$-módulos. De \sref[I]{I.10.8.14} se sigue que existe un
entorno abierto $U$ de $Y'$ tal que la restricción a $U$ del
homomorfismo de $\sh{O}_Y$-módulos
\[
  \widetilde{\vphi}:\widetilde A\longrightarrow\widetilde B
\]
es sobreyectiva (respectivamente biyectiva), lo que termina la demostración.
\end{proof}


\subsection{Un criterio de amplitud}
\label{subsection:III.4.7}

\begin{theorem}[4.7.1]
\label{III.4.7.1}
Sean $Y$ un esquema localmente noetheriano, $f:X\to Y$ un
morfismo propio, $\sh{L}$ un $\sh{O}_X$-módulo invertible,
e $y$ un punto de $Y$. Pongamos
\[
  X_y=X\times_Y\Spec(\kres(y))=f^{-1}(y),
\]
y sea $g:X_y\to X$ el morfismo canónico. Si
\[
  \sh{L}_y=g^*(\sh{L})
  =\sh{L}\otimes_{\sh{O}_Y}\kres(y)
\]
es amplio sobre $X_y$, entonces existe un entorno abierto $U$ de $y$
en $Y$ tal que $\sh{L}|_{f^{-1}(U)}$ es amplio para la restricción
de $f$ a $f^{-1}(U)$.
\end{theorem}

\begin{proof}
\oldpage[III]{145}

\noindent\textit{I)} Pongamos
\[
  Y'=\Spec(\sh{O}_y),\qquad
  X'=X\times_Y Y',\qquad
  \sh{L}'=\sh{L}\otimes_{\sh{O}_Y}\sh{O}_y.
\]
Probamos primero que $\sh{L}'$ es amplio para $f'=f_{(Y')}$. Tenemos
el diagrama conmutativo
\[
  \xymatrix{
    X\ar[d]_f
      &
    X'\ar[l]\ar[d]_{f'}
      &
    X_y\ar[l]\ar[d]_{f_y}
    \\
    Y
      &
    Y'\ar[l]
      &
    \Spec(\kres(y))\ar[l].
  }
\]
Como $f'$ es propio \sref[II]{II.5.4.2}, (iii), y $\sh{O}_y$ es
noetheriano, \sref[II]{II.4.6.6} muestra que podemos reducirnos al caso
\[
  Y=Y'=\Spec(\sh{O}_y),\qquad X=X',
\]
suponer que $\sh{L}_y$ es amplio para $f_y$, y probar que $\sh{L}$
es amplio para $f$. Aplicamos el criterio
\sref{III.2.6.1}, c), y probamos que para todo
$\sh{O}_X$-módulo coherente $\sh{F}$ existe un entero $N$ tal que
\[
  \HH^1(X,\sh{F}(n))=0
  \quad\text{para todo }n\geq N,
  \qquad
  \sh{F}(n)=\sh{F}\otimes\sh{L}^{\otimes n}.
\]

El punto $y$ es cerrado en $Y$ y corresponde al ideal maximal
$\mathfrak m$ de $\sh{O}_y$. Por tanto, $X_y$ es el
subesquema cerrado de $X$ definido por el ideal coherente
\[
  \sh{J}
  =
  f^*(\widetilde{\mathfrak m})\sh{O}_X
  =
  \mathfrak m\sh{O}_X
\]
\sref[I]{I.4.4.5}, y $g$ es su inyección canónica. Consideremos
la $\kres(y)$-álgebra graduada
\[
  S=\operatorname{gr}(\sh{O}_y)
   =\bigoplus_{j\geq0}
      \mathfrak m^j/\mathfrak m^{j+1}.
\]
Es de tipo finito porque $\sh{O}_y$ es noetheriano. La
$\sh{O}_X$-álgebra
\[
  \sh{S}=f^*(\widetilde S)
\]
es cuasicoherente y de tipo finito, y es anulada por
$\sh{J}$. Así, si $\sh{S}_y=g^*(\sh{S})$, entonces $\sh{S}_y$ es una
$\sh{O}_{X_y}$-álgebra cuasicoherente de tipo finito y
\[
  \sh{S}=g_*(\sh{S}_y).
\]

Pongamos
\[
  \sh{M}_j
  =
  \mathfrak m^j\sh{F}/\mathfrak m^{j+1}\sh{F},
  \qquad
  \sh{M}
  =
  \bigoplus_{j\geq0}\sh{M}_j
  =
  \operatorname{gr}(\sh{F}).
\]
Como $\sh{F}$ es coherente, $\sh{M}$ es un
$\sh{S}$-módulo cuasicoherente de tipo finito \sref[0]{0.10.1.1}; también es
anulado por $\sh{J}$. En consecuencia, si
\[
  \sh{M}'_j=g^*(\sh{M}_j),
  \qquad
  \sh{M}'=g^*(\sh{M})
    =\bigoplus_{j\geq0}\sh{M}'_j,
\]
entonces $\sh{M}'$ es un $\sh{S}_y$-módulo graduado cuasicoherente de
tipo finito y $\sh{M}=g_*(\sh{M}')$. Además, poniendo
\[
  \sh{M}'_j(n)
  =
  \sh{M}'_j\otimes\sh{L}_y^{\otimes n},
\]
tenemos $\sh{M}'_j(n)=g^*(\sh{M}_j(n))$.

El morfismo $f_y$ es propio \sref[II]{II.5.4.2}, (iii), y
$\sh{L}_y$ es amplio; por tanto, $f_y$ es proyectivo
\sref[II]{II.5.5.4}, \sref[II]{II.4.6.11}. Podemos, pues, aplicar
\sref{III.2.4.1}, (ii), a
$\Spec(\kres(y))$, $f_y$, $\sh{S}_y$, $\sh{L}_y$ y $\sh{M}'$.
Existe un entero $N$ tal que, para $n\geq N$,
\[
  \HH^q(X_y,\sh{M}'_j(n))=0
\]
para todo $q>0$ y todo $j$. Se sigue que
\[
  \HH^q(X,\sh{M}_j(n))=0
\]
para todo $q>0$ y todo $j$ (G, II, 4.9.1).

Pongamos ahora
\[
  \sh{F}(n)_j
  =
  \sh{F}(n)/\mathfrak m^{j+1}\sh{F}(n).
\]
Para $j\geq1$ tenemos
\[
  \sh{F}(n)_{j-1}
  =
  \sh{F}(n)_j/\sh{M}_j(n),
  \qquad
  \sh{F}(n)_0=\sh{M}_0(n).
\]
Así, $\HH^1(X,\sh{F}(n)_0)=0$, y la sucesión exacta de cohomología
da, para todo $j\geq1$,
\[
  \HH^1(X,\sh{F}(n)_j)
  =
  \HH^1(X,\sh{F}(n)_{j-1}).
\]
Por tanto, $\HH^1(X,\sh{F}(n)_j)=0$ para todo $j\geq0$. Por
\sref{III.4.2.1}, se sigue que
\[
  \HH^1(X,\sh{F}(n))=0,
\]
lo que prueba la afirmación.

\oldpage[III]{146}

\noindent\textit{II)} Volvamos a la notación del comienzo de
la demostración. Podemos suponer que $Y=\Spec(A)$ es afín. Como
$f'$ es de tipo finito y $\sh{L}'$ es amplio para $f'$, existe un
entero $m>0$ tal que $\sh{L}'^{\otimes m}$ es muy amplio para
$f'$ \sref[II]{II.4.6.11}. Sustituyendo $\sh{L}$ por
$\sh{L}^{\otimes m}$, basta considerar el caso en que
$\sh{L}'$ es muy amplio para $f'$ y probar que
$\sh{L}|_{f^{-1}(U)}$ es muy amplio para la restricción de $f$.

Como $f'$ es propio, existe, para un entero adecuado $r>0$, una
$Y'$-inmersión cerrada
\[
  j:X'\longrightarrow P=\bb{P}_{Y'}^r
\]
tal que $\sh{L}'\simeq j^*(\sh{O}_P(1))$
\sref[II]{II.5.5.4}, (ii). Esta inmersión cerrada corresponde
canónicamente a un $\sh{O}_{X'}$-homomorfismo sobreyectivo
\[
  u:\sh{O}_{X'}^{r+1}\longrightarrow\sh{L}'
\]
\sref[II]{II.4.2.3}, o equivalentemente \sref[0]{0.5.1.1} a
$r+1$ secciones $s'_i$ $(0\leq i\leq r)$ de $\sh{L}'$ sobre $X'$
que generan este $\sh{O}_{X'}$-módulo.

Por definición, estas son también secciones de $f'_*(\sh{L}')$ sobre
$Y'$, y
\[
  f'_*(\sh{L}')
  =
  f_*(\sh{L})\otimes_{\sh{O}_Y}\sh{O}_{Y'}
\]
\sref[0]{0.5.4.10}. Como $Y$ es afín y $\sh{O}_y$ es el
anillo local de $A$ en el ideal primo $\mathfrak j_y$, cada $s'_i$
tiene la forma
\[
  s'_i=s''_i/t_i,
\]
donde $s''_i$ es una sección de $f_*(\sh{L})$ sobre $Y$ y
$t_i\in A\setminus\mathfrak j_y$. Se sigue que existen un
entorno abierto $V$ de $y$ en $Y$ y secciones $s_i$ de
$f_*(\sh{L})|_V$ tales que $s'_i=s_i/1$ (el espacio $Y'$ está
contenido en $V$; véase \sref[I]{I.2.4.2}).

Las secciones $s_i$ son secciones de $\sh{L}$ sobre $f^{-1}(V)$ y
definen, por tanto, un homomorfismo
\[
  v:
  \bigl(\sh{O}_X|_{f^{-1}(V)}\bigr)^{r+1}
  \longrightarrow
  \sh{L}|_{f^{-1}(V)}.
\]
Por construcción, es sobreyectivo en todo punto de $f^{-1}(y)$.
Como $\Coker(v)$ es coherente \sref[0]{0.5.3.4}, su soporte es
cerrado \sref[0]{0.5.2.2}; por tanto, existe un entorno abierto
$W\subset f^{-1}(V)$ de $f^{-1}(y)$ sobre el cual $v$ es sobreyectivo.
Como $f$ es cerrado, después de restringir alrededor de $y$ podemos elegir un
entorno abierto $U\subset V$ con $f^{-1}(U)\subset W$.
La conclusión se sigue ahora de \sref[II]{II.4.2.3}.
\end{proof}


\subsection{Morfismos finitos de esquemas formales}
\label{subsection:III.4.8}

\begin{proposition}[4.8.1]
\label{III.4.8.1}
Sea $\mathfrak{Y}$ un esquema formal localmente noetheriano,
$\sh{K}$ un ideal de definición de $\mathfrak{Y}$, y
$f:\mathfrak{X}\to\mathfrak{Y}$ un morfismo de esquemas formales.
Las condiciones siguientes son equivalentes:
\begin{enumerate}
  \item[\rm{a)}] $\mathfrak{X}$ es localmente noetheriano, $f$ es un
    morfismo ádico \sref[I]{I.10.12.1}, y, si
    $\sh{J}=f^*(\sh{K})\sh{O}_{\mathfrak{X}}$, el morfismo
    \[
      f_0:
      (\mathfrak{X},\sh{O}_{\mathfrak{X}}/\sh{J})
      \longrightarrow
      (\mathfrak{Y},\sh{O}_{\mathfrak{Y}}/\sh{K})
    \]
    inducido por $f$ es finito.
  \item[\rm{b)}] $\mathfrak{X}$ es localmente noetheriano y es el
    límite inductivo de un sistema inductivo ádico de $(Y_n)$, $(X_n)$,
    tal que el morfismo $X_0\to Y_0$ es finito.
  \item[\rm{c)}] Todo punto de $\mathfrak{Y}$ posee un entorno abierto
    formal afín noetheriano $V$ tal que $f^{-1}(V)$ es
    un abierto formal afín y
    $\Gamma(f^{-1}(V),\sh{O}_{\mathfrak{X}})$ es un
    $\Gamma(V,\sh{O}_{\mathfrak{Y}})$-módulo de tipo finito.
\end{enumerate}
\end{proposition}

\begin{proof}
De \sref[I]{I.10.12.3} se sigue inmediatamente que \emph{a)} implica
\emph{b)}. Para ver que \emph{b)} implica \emph{c)}, podemos suponer
que
\[
  \mathfrak{Y}=\Spf(B),\qquad \mathfrak{X}=\Spf(A),
\]
donde $B$ es un anillo noetheriano ádico y $\mathfrak{R}$ es un ideal
de definición de $B$. Por hipótesis, $X_0$ es un esquema afín
cuyo anillo $A_0$ es un $B/\mathfrak{R}$-módulo de tipo finito
\sref[II]{II.6.1.3}. Por \sref[I]{I.5.1.9}, cada $X_n$ es un esquema
afín; si $A_n$ denota su anillo, la condición \emph{b)} implica que,
para $m\leq n$,
\[
  A_m\simeq A_n/\mathfrak{R}^{m+1}A_n.
\]
Se sigue que $\mathfrak{X}\simeq\Spf(A)$, donde
$A=\varprojlim A_n$, y la conclusión se sigue de
\sref[0]{0.7.2.9}.

Finalmente, para probar que \emph{c)} implica \emph{a)}, podemos restringirnos
de nuevo al caso
\[
  \mathfrak{Y}=\Spf(B),\qquad \mathfrak{X}=\Spf(A),
\]
donde $A$ es una $B$-álgebra finita. Como
$A/\mathfrak{R}A$ es entonces una
$B/\mathfrak{R}$-álgebra finita, \sref[I]{I.10.10.9} muestra que se
satisfacen las condiciones de \emph{a)}.
\end{proof}

\begin{definition}[4.8.2]
\label{III.4.8.2}
Cuando se satisfacen las condiciones equivalentes \emph{a)}, \emph{b)} y \emph{c)}
de \sref{III.4.8.1}, se dice que el morfismo $f$ es
\emph{finito}; también se dice que $\mathfrak{X}$ es un
\emph{$\mathfrak{Y}$-esquema formal finito}, o un esquema
formal \emph{finito sobre $\mathfrak{Y}$}.
\end{definition}

\begin{proposition}[4.8.3]
\label{III.4.8.3}
\begin{enumerate}
  \item[\rm{(i)}] Una inmersión cerrada de esquemas formales localmente
    noetherianos es un morfismo finito.
  \item[\rm{(ii)}] El compuesto de dos morfismos finitos de esquemas formales
    localmente noetherianos es un morfismo finito.
  \item[\rm{(iii)}] Sean $\mathfrak{X}$, $\mathfrak{Y}$ y
    $\mathfrak{S}$ esquemas formales localmente noetherianos,
    $f:\mathfrak{X}\to\mathfrak{S}$ un morfismo finito, y
    $g:\mathfrak{Y}\to\mathfrak{S}$ un morfismo. Entonces la proyección
    \[
      \mathfrak{X}\times_{\mathfrak{S}}\mathfrak{Y}
      \longrightarrow\mathfrak{Y}
    \]
    es finita.
  \item[\rm{(iv)}] Sea $\mathfrak{S}$ un esquema formal localmente
    noetheriano, y sean $\mathfrak{X}'$ y $\mathfrak{Y}'$
    esquemas formales localmente noetherianos tales que
    $\mathfrak{X}'\times_{\mathfrak{S}}\mathfrak{Y}'$ es localmente
    noetheriano. Si $\mathfrak{X}$ y $\mathfrak{Y}$ son
    $\mathfrak{S}$-esquemas formales localmente noetherianos y
    \[
      f:\mathfrak{X}\to\mathfrak{X}',\qquad
      g:\mathfrak{Y}\to\mathfrak{Y}'
    \]
    son $\mathfrak{S}$-morfismos finitos, entonces
    $f\times_{\mathfrak{S}}g$ es finito.
  \item[\rm{(v)}] Sean
    \[
      f:\mathfrak{X}\to\mathfrak{Y},\qquad
      g:\mathfrak{Y}\to\mathfrak{Z}
    \]
    morfismos de esquemas formales localmente noetherianos tales que
    $g$ es de tipo finito y separado. Si $g\circ f$ es finito,
    entonces $f$ es finito.
\end{enumerate}
\end{proposition}

\begin{proof}
La afirmación \emph{(i)} es inmediata. Utilizando el criterio \emph{a)} de
\sref{III.4.8.1}, las demás afirmaciones se reducen de inmediato a las
proposiciones correspondientes para morfismos de esquemas ordinarios
\sref[II]{II.6.1.5}. Dejamos los detalles al lector, siguiendo
el modelo de \sref[I]{I.10.13.5}.
\end{proof}

\begin{corollary}[4.8.4]
\label{III.4.8.4}
Bajo las hipótesis de \sref[I]{I.10.9.9}, si $f$ es un morfismo
finito, también lo es su extensión $\widehat{f}$ a las
completaciones.
\end{corollary}

\begin{corollary}[4.8.5]
\label{III.4.8.5}
Sea $\mathfrak{X}$ un esquema formal finito sobre
$\mathfrak{Y}$, y sea
$f:\mathfrak{X}\to\mathfrak{Y}$ el morfismo estructural. Para
todo subconjunto abierto $U\subset\mathfrak{Y}$, el esquema formal
$f^{-1}(U)$ es finito sobre $U$.
\end{corollary}

\begin{proposition}[4.8.6]
\label{III.4.8.6}
Si $f:\mathfrak{X}\to\mathfrak{Y}$ es un morfismo finito de esquemas
formales localmente noetherianos, entonces
$f_*(\sh{O}_{\mathfrak{X}})$ es una
$\sh{O}_{\mathfrak{Y}}$-álgebra coherente.
\end{proposition}

\begin{proof}
Podemos considerar $f$ como el límite inductivo de un sistema inductivo
$(f_n)$ de morfismos $f_n:X_n\to Y_n$. Demostremos que los
$f_n$ son finitos y que $f_*(\sh{O}_{\mathfrak{X}})$ es
isomorfo al límite proyectivo de los
$(f_n)_*(\sh{O}_{X_n})$; esto probará la afirmación
\sref[I]{I.10.10.5}. Basta restringirse al caso
\[
  \mathfrak{Y}=\Spf(B),\qquad \mathfrak{X}=\Spf(A),
\]
y observar que, si $\mathfrak{R}$ es un ideal de definición
de $B$ y $A$ es un $B$-módulo de tipo finito, entonces
$A/\mathfrak{R}^{n+1}A$ es un módulo de tipo finito sobre
$B/\mathfrak{R}^{n+1}B$, y $A$ es el límite proyectivo de los
$A/\mathfrak{R}^{n+1}A$.
\end{proof}

Recíprocamente:

\begin{proposition}[4.8.7]
\label{III.4.8.7}
Sea $\mathfrak{Y}$ un esquema formal localmente noetheriano y
$\sh{A}$ una $\sh{O}_{\mathfrak{Y}}$-álgebra coherente. Existe
un esquema formal $\mathfrak{X}$ finito sobre $\mathfrak{Y}$,
único salvo $\mathfrak{Y}$-isomorfismo único, tal que
\[
  f_*(\sh{O}_{\mathfrak{X}})=\sh{A},
\]
donde $f:\mathfrak{X}\to\mathfrak{Y}$ es el morfismo estructural.
\end{proposition}

\begin{proof}
\oldpage[III]{148}
Sea $\sh{K}$ un ideal de definición de $\mathfrak{Y}$, y pongamos
\[
  Y_n=(\mathfrak{Y},
       \sh{O}_{\mathfrak{Y}}/\sh{K}^{n+1}),
  \qquad
  \sh{A}_n=\sh{A}/\sh{K}^{n+1}\sh{A}.
\]
Claramente, $\sh{A}_n$ es una $\sh{O}_{Y_n}$-álgebra finita y, por tanto,
define un $Y_n$-esquema finito
\[
  X_n=\Spec(\sh{A}_n)
\]
\sref[II]{II.6.1.3}. Para $m\leq n$, el homomorfismo canónico
sobreyectivo
\[
  h_{mn}:\sh{A}_n\longrightarrow\sh{A}_m
\]
define un morfismo $u_{mn}:X_m\to X_n$ para el cual el diagrama
\[
  \xymatrix{
    X_n\ar[d]_{f_n} &
    X_m\ar[l]_{u_{mn}}\ar[d]^{f_m}\\
    Y_n &
    Y_m\ar[l]
  }
\]
es conmutativo, donde $f_n$ denota el morfismo estructural; además,
identifica $X_m$ con
$X_n\times_{Y_n}Y_m$, como se sigue de inmediato de
\sref[II]{II.1.4.6}. El esquema formal $\mathfrak{X}$ obtenido
como límite inductivo del sistema inductivo $(X_n)$ es, por tanto,
localmente noetheriano, y su morfismo estructural
\[
  f:\mathfrak{X}\longrightarrow\mathfrak{Y},
\]
el límite inductivo del sistema $(f_n)$, es finito
\sref{III.4.8.1} \sref[I]{I.10.12.3.1}. Además, la demostración de
\sref{III.4.8.6} muestra que
$f_*(\sh{O}_{\mathfrak{X}})$ es el límite proyectivo de las
$\sh{A}_n$, y es, por tanto, igual a $\sh{A}$
\sref[I]{I.10.10.6}. La afirmación de unicidad se sigue del
resultado más general siguiente.
\end{proof}

\begin{proposition}[4.8.8]
\label{III.4.8.8}
Sea $\mathfrak{Y}$ un esquema formal localmente noetheriano, y
sean $\mathfrak{X}$ y $\mathfrak{X}'$
$\mathfrak{Y}$-esquemas formales finitos sobre $\mathfrak{Y}$, con
morfismos estructurales
\[
  f:\mathfrak{X}\to\mathfrak{Y},\qquad
  f':\mathfrak{X}'\to\mathfrak{Y}.
\]
Existe una biyección canónica
\[
  \Hom_{\mathfrak{Y}}(\mathfrak{X},\mathfrak{X}')
  \xrightarrow{\ \sim\ }
  \Hom_{\sh{O}_{\mathfrak{Y}}}
  \bigl(f'_*(\sh{O}_{\mathfrak{X}'}),
        f_*(\sh{O}_{\mathfrak{X}})\bigr)\footnotemark.
\]
\footnotetext{La última expresión denota el conjunto de homomorfismos de
$\sh{O}_{\mathfrak{Y}}$-álgebras
$f'_*(\sh{O}_{\mathfrak{X}'})\to f_*(\sh{O}_{\mathfrak{X}})$.}
\end{proposition}

\begin{proof}
La aplicación $h\mapsto\sh{A}(h)$ se define como en
\sref[II]{II.1.1.2}. Para probar que es biyectiva, nos reducimos de
inmediato al caso en que
$\mathfrak{Y}=\Spf(B)$ es un esquema formal afín noetheriano.
Entonces
\[
  \mathfrak{X}=\Spf(A),\qquad
  \mathfrak{X}'=\Spf(A'),
\]
donde $A$ y $A'$ son $B$-álgebras finitas, y
\[
  f_*(\sh{O}_{\mathfrak{X}})=A^\Delta,\qquad
  f'_*(\sh{O}_{\mathfrak{X}'})=A'^\Delta.
\]
La conclusión se sigue de la correspondencia biyectiva entre
$\mathfrak{Y}$-morfismos
$\mathfrak{X}\to\mathfrak{X}'$ y homomorfismos continuos de
$B$-álgebras $A'\to A$ \sref[I]{I.10.2.2}, junto con la
correspondencia biyectiva entre homomorfismos de $B$-módulos
$A'\to A$ y homomorfismos de
$\sh{O}_{\mathfrak{Y}}$-módulos
$A'^\Delta\to A^\Delta$ \sref[I]{I.10.10.2.3}.
\end{proof}

\begin{corollary}[4.8.9]
\label{III.4.8.9}
Bajo la biyección canónica de \sref{III.4.8.8}, las inmersiones
cerradas $\mathfrak{X}\to\mathfrak{X}'$ corresponden a
homomorfismos sobreyectivos de $\sh{O}_{\mathfrak{Y}}$-álgebras
\[
  f'_*(\sh{O}_{\mathfrak{X}'})
  \longrightarrow
  f_*(\sh{O}_{\mathfrak{X}}).
\]
\end{corollary}

\begin{proof}
La cuestión es local sobre $\mathfrak{Y}$ y, por tanto, se reduce a la
definición de las inmersiones cerradas de esquemas formales localmente
noetherianos \sref[I]{I.10.14.2}.
\end{proof}

\begin{corollary}[4.8.10]
\label{III.4.8.10}
Con la notación y las hipótesis de \sref{III.4.8.1}, un morfismo
ádico $f$ es una inmersión cerrada si y solo si $f_0$ es una
inmersión cerrada de esquemas ordinarios.
\end{corollary}

\begin{proof}
Esto se sigue de inmediato de \sref{III.4.8.9} y del criterio de
sobreyectividad para un homomorfismo de
$\sh{O}_{\mathfrak{Y}}$-módulos coherentes \sref[I]{I.10.11.5}.
\end{proof}

\begin{proposition}[4.8.11]
\label{III.4.8.11}
Un morfismo $f:\mathfrak{X}\to\mathfrak{Y}$ de esquemas formales
localmente noetherianos
\oldpage[III]{149}
es finito si y solo si es propio y todas sus fibras
$f^{-1}(y)$, para $y\in\mathfrak{Y}$, son finitas.
\end{proposition}

\begin{proof}
Por las definiciones \sref{III.3.4.1} y \sref{III.4.8.2}, la
afirmación se reduce inmediatamente a la misma afirmación para $f_0$
(con la notación de \sref{III.4.8.1}), que es precisamente
\sref{III.4.4.2}.
\end{proof}

\section{Un teorema de existencia para los haces algebraicos coherentes}
\label{section:III.5}


\subsection{Enunciado del teorema}
\label{subsection:III.5.1}

\begin{env}[5.1.1]
\label{III.5.1.1}
Sea $A$ un anillo noetheriano ádico y sea $\mathfrak{J}$ un ideal de
definición de $A$, de modo que $A$ sea separado y completo para la
topología $\mathfrak{J}$-ádica. Si $Y=\Spec(A)$, el preesquema formal
afín $\Spf(A)$ se identifica con la compleción
$\widehat{Y}$ de $Y$ a lo largo del subconjunto cerrado
$Y'=\mathrm{V}(\mathfrak{J})$ \sref[I]{I.10.10.1}. Sea $X$ un
$Y$-preesquema ordinario de tipo finito, de morfismo estructural
$f:X\to Y$. Denotemos por $\widehat{X}$ la compleción de $X$ a lo largo
del subconjunto cerrado $X'=f^{-1}(Y')$, o, equivalentemente, el
$\widehat{Y}$-preesquema formal $X\times_Y\widehat{Y}$, y por
$\widehat{f}:\widehat{X}\to\widehat{Y}$ la extensión de $f$ a las
compleciones. Finalmente, para todo $\sh{O}_X$-módulo coherente
$\sh{F}$, denotemos por $\widehat{\sh{F}}$ su compleción
$\sh{F}_{/X'}$, que es un
$\sh{O}_{\widehat{X}}$-módulo coherente.
\end{env}

\begin{proposition}[5.1.2]
\label{III.5.1.2}
Con las hipótesis y la notación de \sref{III.5.1.1}, sea $\sh{F}$
un $\sh{O}_X$-módulo coherente cuyo soporte es propio sobre $Y$
\sref[II]{II.5.4.10}. Los homomorfismos canónicos
\sref{III.4.1.4}
\[
  \rho_i:
  \HH^i(X,\sh{F})
  \longrightarrow
  \HH^i(\widehat{X},\widehat{\sh{F}})
\]
son isomorfismos.
\end{proposition}

\begin{proof}
El $A$-módulo $\HH^i(X,\sh{F})$ es de tipo finito
\sref{III.3.2.4} y, por tanto, se identifica con su compleción
separada para la topología $\mathfrak{J}$-preádica
\sref[0]{0.7.3.6}. La proposición es, por consiguiente, un caso particular
de \sref{III.4.1.10}.
\end{proof}

Recordemos que los isomorfismos canónicos $\rho_i$ conmutan con los
homomorfismos de conexión asociados a toda sucesión exacta
\[
  0\longrightarrow\sh{F}'
  \longrightarrow\sh{F}
  \longrightarrow\sh{F}''
  \longrightarrow0
\]
de $\sh{O}_X$-módulos coherentes
\sref[0]{0.12.1.6} \sref[I]{I.10.8.9}.

\begin{corollary}[5.1.3]
\label{III.5.1.3}
Sean $\sh{F}$ y $\sh{G}$ $\sh{O}_X$-módulos coherentes tales que
la intersección de sus soportes sea propia sobre $Y$. El
homomorfismo canónico
\[
  \Hom_{\sh{O}_X}(\sh{F},\sh{G})
  \longrightarrow
  \Hom_{\sh{O}_{\widehat{X}}}
  (\widehat{\sh{F}},\widehat{\sh{G}})
  \tag{5.1.3.1}
  \label{III.5.1.3.1.eq}
\]
que envía un homomorfismo $u:\sh{F}\to\sh{G}$ a su compleción
$\widehat{u}:\widehat{\sh{F}}\to\widehat{\sh{G}}$, es un
isomorfismo. Además, cuando $f$ es cerrado,
$\widehat{u}$ es inyectivo (respectivamente, sobreyectivo) si y solo
si $u$ es inyectivo (respectivamente, sobreyectivo).
\end{corollary}

\begin{proof}
La primera afirmación es un caso particular de \sref{III.4.5.3}; también
se deduce de que la fuente de \eqref{III.5.1.3.1.eq} es
un $A$-módulo de tipo finito y, por tanto, se identifica con su
compleción separada. Para la segunda afirmación, observemos que, por
\sref[I]{I.10.8.14}, $\widehat{u}$ es inyectivo (respectivamente,
sobreyectivo) si y solo si existe un entorno de $X'$ en
$X$ sobre el cual $u$ es inyectivo (respectivamente, sobreyectivo). La
conclusión se deduce entonces del lema siguiente.
\end{proof}

\oldpage[III]{150}

\begin{lemma}[5.1.3.1]
\label{III.5.1.3.1}
Bajo las hipótesis de \sref{III.5.1.1}, supongamos además que
$f$ es cerrado. Todo entorno de $X'$ en $X$ es entonces igual
a $X$.
\end{lemma}

\begin{proof}
Nos reducimos primero al caso $f(X)=Y$. En efecto, por hipótesis,
$f(X)$ es un subconjunto cerrado $Z$ de $Y$. Podemos sustituir $f$ por
$f_{\mathrm{red}}$ \sref[I]{I.6.3.4} y suponer así que $X$
e $Y$ son reducidos; podemos entonces sustituir $Y$ por el subpreesquema
cerrado reducido cuyo espacio subyacente es $Z$
\sref[I]{I.5.2.2}, pues todo ideal de $A$ es cerrado y todo
anillo cociente de $A$ es, por tanto, ádico y noetheriano.

Se tiene entonces $f(X')=Y'$. Si $V$ es un entorno abierto de $X'$
en $X$, el conjunto $f(X-V)$ es cerrado en $Y$ y no corta a $Y'$.
Esto es imposible salvo que $X-V$ sea vacío, pues $\mathfrak{J}$
está contenido en el radical de $A$
\sref[I]{I.1.1.15} \sref[0]{0.7.1.10}. Por tanto, $V=X$.
\end{proof}

Si nos restringimos a los $\sh{O}_X$-módulos coherentes cuyo
soporte es propio sobre $Y$, \sref{III.5.1.3} afirma, en lenguaje
categórico, que el funtor
\[
  \sh{F}\longmapsto\widehat{\sh{F}}
\]
es plenamente fiel de la categoría de tales $\sh{O}_X$-módulos en
la categoría de los $\sh{O}_{\widehat{X}}$-módulos coherentes. Proporciona así
una equivalencia de la primera categoría con una subcategoría plena
de la segunda \sref[0]{0.8.1.6}. El teorema de existencia
mostrará que, cuando $X$ es propio sobre $Y$, esta subcategoría es en
realidad la categoría de todos los
$\sh{O}_{\widehat{X}}$-módulos coherentes. Más precisamente:

\begin{theorem}[5.1.4]
\label{III.5.1.4}
Sea $A$ un anillo noetheriano ádico, sea $Y=\Spec(A)$, sea
$\mathfrak{J}$ un ideal de definición de $A$, y pongamos
$Y'=\mathrm{V}(\mathfrak{J})$. Sea $f:X\to Y$ un morfismo separado
de tipo finito y pongamos $X'=f^{-1}(Y')$. Sean
\[
  \widehat{Y}=Y_{/Y'}=\Spf(A),\qquad
  \widehat{X}=X_{/X'}
\]
las compleciones de $Y$ y $X$ a lo largo de $Y'$ y $X'$, y sea
$\widehat{f}:\widehat{X}\to\widehat{Y}$ la extensión de $f$ a
las compleciones. Entonces el funtor
\[
  \sh{F}\longmapsto\sh{F}_{/X'}=\widehat{\sh{F}}
\]
es una equivalencia de la categoría de los
$\sh{O}_X$-módulos coherentes cuyo soporte es propio sobre $\Spec(A)$ en la
categoría de los $\sh{O}_{\widehat{X}}$-módulos coherentes cuyo soporte
es propio sobre $\Spf(A)$.
\end{theorem}

En otras palabras, teniendo en cuenta \sref{III.5.1.3}:

\begin{corollary}[5.1.5]
\label{III.5.1.5}
Un $\sh{O}_{\widehat{X}}$-módulo es isomorfo a la compleción de
un $\sh{O}_X$-módulo coherente cuyo soporte es propio sobre
$\Spec(A)$ si y solo si es coherente y su soporte es propio
sobre $\Spf(A)$.
\end{corollary}

El caso más importante es el siguiente:

\begin{corollary}[5.1.6]
\label{III.5.1.6}
Supongamos que $X$ es propio sobre $Y=\Spec(A)$. Entonces el funtor
$\sh{F}\mapsto\widehat{\sh{F}}$ es una equivalencia entre la
categoría de los $\sh{O}_X$-módulos coherentes y la categoría de los
$\sh{O}_{\widehat{X}}$-módulos coherentes.
\end{corollary}

\begin{env}[5.1.7]
\label{III.5.1.7}
\emph{Escolio.}
Teniendo en cuenta la caracterización de los haces coherentes sobre
preesquemas formales \sref[I]{I.10.11.3}, vemos que, bajo las
condiciones de \sref{III.5.1.1}, dar un
$\sh{O}_X$-módulo coherente cuyo soporte es propio sobre $\Spec(A)$
equivale---si
\[
  Y_n=\Spec(A/\mathfrak{J}^{n+1}),\qquad
  X_n=X\times_Y Y_n
\]
---a dar un sistema proyectivo $(\sh{F}_n)$ de
$\sh{O}_{X_n}$-módulos coherentes tal que, para $m\leq n$,
\[
  \sh{F}_m
  =
  \sh{F}_n\otimes_{\sh{O}_{Y_n}}\sh{O}_{Y_m},
  \qquad\text{o equivalentemente}\qquad
  \sh{F}_m
  =
  \sh{F}_n/\mathfrak{J}^{m+1}\sh{F}_n,
\]
y tal que el soporte de $\sh{F}_0$ sea un subconjunto de $X_0$
propio sobre $Y_0$. Por \sref[I]{I.10.11.4}, los homomorfismos de
$\sh{O}_X$-módulos coherentes se interpretan de modo análogo como
homomorfismos de sistemas proyectivos de
$\sh{O}_{X_n}$-módulos coherentes.

\oldpage[III]{151}

En todas las aplicaciones conocidas, $A$ es de hecho un anillo
noetheriano ádico local. Así, los $Y_n$ son espectros de anillos
artinianos locales, y los resultados de esta subsección y de las
anteriores reducen, en gran medida, la geometría algebraica sobre un
anillo noetheriano ádico local a la geometría algebraica sobre anillos
artinianos locales.
\end{env}

\begin{corollary}[5.1.8]
\label{III.5.1.8}
Bajo las hipótesis de \sref{III.5.1.4}, la aplicación
\[
  Z\longmapsto\widehat{Z}=Z_{/(Z\cap X')}
\]
es una biyección del conjunto de los subpreesquemas cerrados $Z$ de $X$
que son propios sobre $Y$ sobre el conjunto de los subpreesquemas
formales cerrados de $\widehat{X}$ que son propios sobre
$\widehat{Y}$.
\end{corollary}

\begin{proof}
Un subpreesquema formal cerrado de $\widehat{X}$ es de la forma
\[
  \bigl(T,
    (\sh{O}_{\widehat{X}}/\sh{A})|_T\bigr),
\]
donde $\sh{A}$ es un ideal coherente de
$\sh{O}_{\widehat{X}}$ \sref[I]{I.10.14.2}. Si $T$ es propio sobre
$\widehat{Y}$, \sref{III.5.1.4} muestra que
$\sh{O}_{\widehat{X}}/\sh{A}$ es isomorfo a un
$\sh{O}_{\widehat{X}}$-módulo de la forma
$\widehat{\sh{F}}$, donde $\sh{F}$ es un
$\sh{O}_X$-módulo coherente cuyo soporte es propio sobre $Y$. Además,
\sref{III.5.1.3} muestra que el homomorfismo canónico
\[
  \sh{O}_{\widehat{X}}
  \longrightarrow
  \sh{O}_{\widehat{X}}/\sh{A}
\]
es de la forma $\widehat{u}$, donde
$u:\sh{O}_X\to\sh{F}$ es un homomorfismo sobreyectivo de
$\sh{O}_X$-módulos. Por consiguiente,
$\sh{F}=\sh{O}_X/\sh{N}$ para un ideal coherente $\sh{N}$ de
$\sh{O}_X$, y
$\sh{A}=\widehat{\sh{N}}$ \sref[I]{I.10.8.8}. La conclusión se
deduce ahora de \sref[I]{I.10.14.7}.
\end{proof}


\subsection{Demostración del teorema de existencia: caso proyectivo y cuasiproyectivo}
\label{subsection:III.5.2}

\begin{env}[5.2.1]
\label{III.5.2.1}
Bajo las condiciones de \sref{III.5.1.4}, diremos
provisionalmente que un $\sh{O}_{\widehat{X}}$-módulo coherente es
\emph{algebraizable} si es isomorfo a la compleción
$\widehat{\sh{F}}$ de un $\sh{O}_X$-módulo coherente $\sh{F}$ cuyo
soporte es propio sobre $Y$.
\end{env}

\begin{lemma}[5.2.2]
\label{III.5.2.2}
Sean $\sh{F}'$ y $\sh{G}'$ $\sh{O}_{\widehat{X}}$-módulos
algebraizables. Para todo homomorfismo
$u:\sh{F}'\to\sh{G}'$, los módulos
$\Ker(u)$, $\Im(u)$ y $\Coker(u)$ son algebraizables.
\end{lemma}

\begin{proof}
Escribamos
\[
  \sh{F}'=\widehat{\sh{F}},\qquad
  \sh{G}'=\widehat{\sh{G}},
\]
donde $\sh{F}$ y $\sh{G}$ son $\sh{O}_X$-módulos coherentes cuyos
soportes son propios sobre $Y$. Por \sref{III.5.1.3},
$u=\widehat{v}$ para un homomorfismo
$v:\sh{F}\to\sh{G}$. Como la compleción es exacta,
$\Ker(\widehat{v})\simeq\widehat{\Ker(v)}$; y como el soporte
de $\Ker(v)$ está contenido en el de $\sh{F}$, se deduce que
$\Ker(u)$ es algebraizable. Los argumentos para $\Im(u)$ y
$\Coker(u)$ son los mismos.
\end{proof}

\begin{proposition}[5.2.3]
\label{III.5.2.3}
Sea $A$ un anillo noetheriano ádico, sea $\mathfrak{J}$ un ideal de
definición de $A$, sea $\mathfrak{Y}=\Spf(A)$, y sea
$f:\mathfrak{X}\to\mathfrak{Y}$ un morfismo propio de preesquemas
formales. Pongamos
\[
  Y_k=\Spec(A/\mathfrak{J}^{k+1}),\qquad
  X_k=\mathfrak{X}\times_{\mathfrak{Y}}Y_k,
\]
y, para todo $\sh{O}_{\mathfrak{X}}$-módulo $\sh{F}$, pongamos
\[
  \sh{F}_k
  =
  \sh{F}\otimes_{\sh{O}_{\mathfrak{X}}}\sh{O}_{X_k}
  =
  \sh{F}/\mathfrak{J}^{k+1}\sh{F}.
\]
Sea $\sh{L}$ un $\sh{O}_{\mathfrak{X}}$-módulo invertible, y
supongamos que
$\sh{L}_0=\sh{L}/\mathfrak{J}\sh{L}$ es un
$\sh{O}_{X_0}$-módulo amplio. Para todo
$\sh{O}_{\mathfrak{X}}$-módulo $\sh{F}$ y todo entero $n$, pongamos
\[
  \sh{F}(n)=\sh{F}\otimes\sh{L}^{\otimes n}.
\]
Entonces, para todo $\sh{O}_{\mathfrak{X}}$-módulo coherente $\sh{F}$,
existe un entero $n_0$ tal que, para todo $n\geq n_0$:
\begin{enumerate}
  \item[\rm{(i)}] el homomorfismo canónico
    \[
      \HH^0(\mathfrak{X},\sh{F}(n))
      \longrightarrow
      \HH^0(\mathfrak{X},\sh{F}_k(n))
    \]
    es sobreyectivo para todo $k\geq0$;
  \item[\rm{(ii)}]
    $\HH^q(\mathfrak{X},\sh{F}(n))=0$ para todo $q>0$.
\end{enumerate}
\end{proposition}

\begin{proof}
Los espacios subyacentes a $\mathfrak{X}$ y $X_0$ son los mismos.
Los haces
\[
  \sh{M}_k
  =
  \mathfrak{J}^k\sh{F}/\mathfrak{J}^{k+1}\sh{F},
\]
al ser anulados por $\mathfrak{J}$, pueden considerarse
$\sh{O}_{X_0}$-módulos coherentes \sref[0]{0.5.3.10}. Además, si
\[
  \sh{M}_k(n)
  =
  \sh{M}_k\otimes_{\sh{O}_{X_0}}\sh{L}_0^{\otimes n},
\]
entonces
\[
  \sh{M}_k(n)
  =
  \mathfrak{J}^k\sh{F}(n)/
  \mathfrak{J}^{k+1}\sh{F}(n).
\]
Como $\sh{L}_0$ es amplio para $f_0:X_0\to Y_0$ y $f_0$ es
propio,
\oldpage[III]{152}
$f_0$ es proyectivo \sref[II]{II.5.5.4}. Sea
\[
  S=\bigoplus_{k\geq0}
  \mathfrak{J}^k/\mathfrak{J}^{k+1}
\]
la $A$-álgebra graduada asociada a la
filtración $\mathfrak{J}$-ádica de $A$. Es de tipo finito
porque $A$ es noetheriano. Pongamos
\[
  \sh{S}'=f_0^*(\widetilde{S}),\qquad
  \sh{M}=\bigoplus_{k\geq0}\sh{M}_k.
\]
Entonces $\sh{S}'$ es una $\sh{O}_{X_0}$-álgebra cuasi-coherente de
tipo finito, y $\sh{M}$ es un $\sh{S}'$-módulo graduado
cuasi-coherente de tipo finito, pues $\sh{F}_0$ es coherente y
genera $\sh{M}$ como $\sh{S}'$-módulo. El Teorema
\sref{III.2.4.1},~\emph{(ii)}, proporciona, por tanto, un entero $n_0$
tal que, para $n\geq n_0$ y todo $k$,
\[
  \HH^q(X_0,\sh{M}_k(n))=0
  \qquad(q>0).
  \tag{5.2.3.1}
  \label{III.5.2.3.1}
\]
Considerado ahora como $\sh{O}_{\mathfrak{X}}$-módulo,
$\sh{M}_k(n)$ satisface, por consiguiente, también
\[
  \HH^q(\mathfrak{X},\sh{M}_k(n))=0
  \qquad(q>0,\ n\geq n_0).
\]

Apliquemos la sucesión exacta de cohomología a
\[
  0\longrightarrow
  \frac{\mathfrak{J}^k\sh{F}(n)}
       {\mathfrak{J}^{k+1}\sh{F}(n)}
  \longrightarrow
  \frac{\mathfrak{J}^h\sh{F}(n)}
       {\mathfrak{J}^{k+1}\sh{F}(n)}
  \longrightarrow
  \frac{\mathfrak{J}^h\sh{F}(n)}
       {\mathfrak{J}^{k}\sh{F}(n)}
  \longrightarrow0.
\]
Para $0\leq h<k$, $n\geq n_0$ y $q>0$, una inducción sobre $k-h$
da
\[
  \HH^q\left(
    \mathfrak{X},
    \frac{\mathfrak{J}^h\sh{F}(n)}
         {\mathfrak{J}^{k}\sh{F}(n)}
  \right)=0,
  \tag{5.2.3.2}
  \label{III.5.2.3.2}
\]
y, en particular, para $h=0$,
\[
  \HH^q(\mathfrak{X},\sh{F}_k(n))=0
  \qquad(n\geq n_0,\ k\geq0,\ q>0).
  \tag{5.2.3.3}
  \label{III.5.2.3.3}
\]
Otra parte de la sucesión exacta de cohomología, con $h=0$, da
\[
  \HH^0(\mathfrak{X},\sh{F}_{k+1}(n))
  \longrightarrow
  \HH^0(\mathfrak{X},\sh{F}_k(n))
  \longrightarrow
  \HH^1\left(
    \mathfrak{X},
    \frac{\mathfrak{J}^k\sh{F}(n)}
         {\mathfrak{J}^{k+1}\sh{F}(n)}
  \right)
  =0.
  \tag{5.2.3.4}
  \label{III.5.2.3.4}
\]
Se deduce que, para $h\leq k$, la aplicación canónica
\[
  \HH^0(\mathfrak{X},\sh{F}_k(n))
  \longrightarrow
  \HH^0(\mathfrak{X},\sh{F}_h(n))
  \tag{5.2.3.5}
  \label{III.5.2.3.5}
\]
es sobreyectiva. Así, para todo $q$, el sistema proyectivo
\[
  \bigl(\HH^q(\mathfrak{X},\sh{F}_k(n))\bigr)_{k\geq0}
\]
satisface la condición (ML) cuando $n\geq n_0$.

Además, todo abierto formal afín $U$ de $\mathfrak{X}$ es un
abierto afín de cada $X_k$ \sref[I]{I.10.5.2}. Por tanto,
$\HH^q(U,\sh{F}_k(n))=0$ para $q>0$ \sref[I]{I.1.3.1}, y
\[
  \HH^0(U,\sh{F}_k(n))
  \longrightarrow
  \HH^0(U,\sh{F}_h(n))
\]
es sobreyectivo para $h\leq k$ \sref[I]{I.1.3.9}. Por tanto, se satisfacen
las hipótesis de \sref[0]{0.13.3.1}. Para $n\geq n_0$
obtenemos:
\begin{enumerate}
  \item[$1^\circ$] para todo $q>0$, la aplicación
    \[
      \HH^q(\mathfrak{X},\sh{F}(n))
      \longrightarrow
      \varprojlim_k
      \HH^q(\mathfrak{X},\sh{F}_k(n))
    \]
    es biyectiva; por tanto, por \eqref{III.5.2.3.3},
    $\HH^q(\mathfrak{X},\sh{F}(n))=0$;
  \item[$2^\circ$] el homomorfismo
    \[
      \HH^0(\mathfrak{X},\sh{F}(n))
      \longrightarrow
      \varprojlim_k
      \HH^0(\mathfrak{X},\sh{F}_k(n))
    \]
    es biyectivo. Como los homomorfismos
    \eqref{III.5.2.3.5} son sobreyectivos, también lo es cada homomorfismo
    \[
      \varprojlim_k
      \HH^0(\mathfrak{X},\sh{F}_k(n))
      \longrightarrow
      \HH^0(\mathfrak{X},\sh{F}_h(n)).
    \]
\end{enumerate}
Esto demuestra las dos afirmaciones.
\end{proof}

\begin{corollary}[5.2.4]
\label{III.5.2.4}
Bajo las hipótesis de \sref{III.5.2.3}, para todo
$\sh{O}_{\mathfrak{X}}$-módulo coherente $\sh{F}$ existe un entero $N$
tal que, para $n\geq N$, $\sh{F}(n)$ está generado por sus secciones
sobre
\oldpage[III]{153}
$\mathfrak{X}$. Equivalentemente, $\sh{F}$ es isomorfo a un cociente
de un $\sh{O}_{\mathfrak{X}}$-módulo de la forma
$(\sh{O}_{\mathfrak{X}}(-n))^k$.
\end{corollary}

\begin{proof}
Como $X_0$ es noetheriano, la hipótesis sobre $\sh{L}_0$ y
\sref[II]{II.4.5.5} proporcionan un entero $n_0$ tal que
$\sh{F}_0(n)$ está generado por sus secciones sobre $\mathfrak{X}$ para
$n\geq n_0$. Podemos tomar $n_0$ suficientemente grande para que
\[
  \Gamma(\mathfrak{X},\sh{F}(n))
  \longrightarrow
  \Gamma(\mathfrak{X},\sh{F}_0(n))
\]
sea también sobreyectivo para $n\geq n_0$ \sref{III.5.2.3}. Existen,
por tanto, un número finito de secciones
$s_i\in\Gamma(\mathfrak{X},\sh{F}(n))$ cuyas imágenes generan
$\sh{F}_0(n)$ \sref[0]{0.5.2.3}. Como $\mathfrak{J}$ está
contenido en el ideal maximal del anillo local en todo punto
de $\mathfrak{X}$, el lema de Nakayama aplicado a esos anillos locales
muestra que los $s_i$ generan $\sh{F}(n)$
\sref[0]{0.5.1.1}.
\end{proof}

\begin{env}[5.2.5]
\label{III.5.2.5}
\emph{Demostración del teorema de existencia: caso proyectivo.}
Con la notación de \sref{III.5.1.4}, supongamos que $f$ es
proyectivo, de modo que existe un $\sh{O}_X$-módulo amplio $\sh{L}$
\sref[I]{I.5.5.4}. Por definición,
\[
  X_n=\widehat{X}\times_{\widehat{Y}}Y_n
\]
es el subpreesquema cerrado
$X\times_Y Y_n=f^{-1}(Y_n)$ de $X$. Si
\[
  \sh{L}'=\sh{L}\otimes_{\sh{O}_X}\sh{O}_{\widehat{X}}
\]
es la compleción de $\sh{L}$, entonces
$\sh{L}'_0=\sh{L}/\mathfrak{J}\sh{L}$, considerado como
$\sh{O}_{X_0}$-módulo; este módulo es amplio
\sref[II]{II.4.6.13}[(i~\emph{bis})].

Podemos, por tanto, aplicar \sref{III.5.2.4} a $\sh{L}'$ y a todo
$\sh{O}_{\widehat{X}}$-módulo coherente $\sh{F}$. Se deduce que
$\sh{F}$ es isomorfo a un cociente de
\[
  \sh{G}=
  \bigl(\sh{L}'^{\otimes(-n)}\bigr)^k
\]
para enteros adecuados $n>0$ y $k>0$. El módulo $\sh{G}$ es la
compleción de $(\sh{L}^{\otimes(-n)})^k$
\sref[I]{I.10.8.10} y, por tanto, es algebraizable. Sea
$u:\sh{G}\to\sh{F}$ la sobreyección canónica, y pongamos
$\sh{H}=\Ker(u)$; este es un
$\sh{O}_{\widehat{X}}$-módulo coherente \sref[0]{0.5.3.4}. Aplicando el mismo
argumento a $\sh{H}$, obtenemos un
$\sh{O}_{\widehat{X}}$-módulo algebraizable $\sh{K}$ y un homomorfismo
$v:\sh{K}\to\sh{G}$ tal que $\sh{H}=\Im(v)$. Así,
$\sh{F}=\Coker(v)$, y $\sh{F}$ es algebraizable por
\sref{III.5.2.2}.
\end{env}

\begin{env}[5.2.6]
\label{III.5.2.6}
\emph{Demostración del teorema de existencia: caso cuasiproyectivo.}
Conservemos la notación de \sref{III.5.1.4} y supongamos ahora que $f$
es cuasiproyectivo. Existe un morfismo proyectivo
$g:Z\to Y$ tal que $X$ se identifica con el $Y$-preesquema
inducido sobre un subconjunto abierto de $Z$ \sref[II]{II.5.3.2}. Si
$Z'=g^{-1}(Y')$, entonces $X'=X\cap Z'$. Por consiguiente, la compleción
$\widehat{X}=X_{/X'}$ se identifica con el preesquema formal
inducido por $\widehat{Z}=Z_{/Z'}$ sobre el subconjunto abierto $X\cap Z'$ de
$\widehat{Z}$ \sref[I]{I.10.8.5}.

Sea $\sh{F}'$ un $\sh{O}_{\widehat{X}}$-módulo coherente cuyo
soporte $T'$ sea propio sobre $\widehat{Y}$. Por definición, existe
un subpreesquema cerrado de $X'$ de espacio subyacente
$T'\subset X'$ tal que la restricción $T'\to Y'$ de $f$ es
propia. Por tanto, $T'$ es propio sobre $Y$ y, en consecuencia, cerrado en
$Z'$ \sref[II]{II.5.4.10}. Se deduce que $\sh{F}'$ es inducido
sobre $\widehat{X}$ por el $\sh{O}_{\widehat{Z}}$-módulo $\sh{G}'$
obtenido pegando $\sh{F}'$, sobre el abierto $\widehat{X}$, con el
haz cero sobre el abierto $\widehat{Z}-T'$; ambos haces coinciden sobre
su intersección $\widehat{X}-T'$.

El módulo $\sh{G}'$ es coherente. Por \sref{III.5.2.5}, existe un
$\sh{O}_Z$-módulo coherente $\sh{G}$ tal que
$\sh{G}'=\widehat{\sh{G}}$. Sea $T$ el soporte de $\sh{G}$;
entonces $T'=T\cap Z'$ \sref[I]{I.10.8.12}. Si $h$ es la restricción
de $g$ al subpreesquema cerrado reducido de $Z$ cuyo espacio subyacente
es $T$, entonces
\[
  T'=h^{-1}(Y')=T\cap g^{-1}(Y').
\]
Así, $X\cap T$ es un subconjunto abierto de $T$ que contiene a $T'$.

\oldpage[III]{154}
Como $h$ es propio \sref[II]{II.5.4.2}, y por tanto cerrado,
\sref{III.5.1.3.1} da $X\cap T=T$. Así, $T\subset X$; y
como $T$ es cerrado en $Z$, es propio sobre $Y$. Si $\sh{F}$ es
el haz inducido sobre $X$ por $\sh{G}$, su compleción
$\widehat{\sh{F}}$ es inducida sobre $\widehat{X}$ por
$\widehat{\sh{G}}$ \sref[I]{I.10.8.4} y, por tanto, es igual a
$\sh{F}'$. Esto termina la demostración.
\end{env}


\subsection{Demostración del teorema de existencia: caso general}
\label{subsection:III.5.3}

\begin{lemma}[5.3.1]
\label{III.5.3.1}
Bajo las condiciones de \sref{III.5.1.4}, sea
\[
  0\longrightarrow\sh{H}\longrightarrow\sh{F}
  \longrightarrow\sh{G}\longrightarrow0
\]
una sucesión exacta de $\sh{O}_{\widehat{X}}$-módulos coherentes.
Si $\sh{G}$ y $\sh{H}$ son algebraizables, entonces $\sh{F}$ es
algebraizable.
\end{lemma}

\begin{proof}
Supongamos que
\[
  \sh{G}=\widehat{\sh{B}},
  \qquad
  \sh{H}=\widehat{\sh{C}},
\]
donde $\sh{B}$ y $\sh{C}$ son $\sh{O}_X$-módulos coherentes cuyos
soportes son propios sobre $Y$. El módulo $\sh{F}$ define canónicamente
un elemento del $A$-módulo
\[
  \Ext^1_{\sh{O}_{\widehat{X}}}
  (\widehat{X};\widehat{\sh{B}},\widehat{\sh{C}})
\]
\sref[0]{0.12.3.2}. Las hipótesis implican que este $A$-módulo es
canónicamente isomorfo a
\[
  \Ext^1_{\sh{O}_X}(X;\sh{B},\sh{C})
\]
\sref{III.4.5.2}. Existe, por tanto, una sucesión exacta
\[
  0\longrightarrow\sh{C}\longrightarrow\sh{A}
  \longrightarrow\sh{B}\longrightarrow0
\]
de $\sh{O}_X$-módulos coherentes cuya clase canónica se envía en la
clase correspondiente a $\sh{F}$. Por definición, teniendo en cuenta
\sref[I]{I.10.8.8}, (ii), esto significa que $\sh{F}$ es isomorfo a
$\widehat{\sh{A}}$. El lema se deduce, pues
$\Supp(\sh{A})$ está contenido en
$\Supp(\sh{B})\cup\Supp(\sh{C})$ y es, por tanto, propio sobre $Y$.
\end{proof}

\begin{corollary}[5.3.2]
\label{III.5.3.2}
Bajo las condiciones de \sref{III.5.1.1}, sea
$u:\sh{F}\to\sh{G}$ un homomorfismo de
$\sh{O}_{\widehat{X}}$-módulos coherentes. Si $\sh{G}$, $\Ker(u)$ y
$\Coker(u)$ son algebraizables, entonces $\sh{F}$ es algebraizable.
\end{corollary}

\begin{proof}
El Lema \sref{III.5.2.2}, aplicado a
$\sh{G}\to\Coker(u)$, muestra que $\Im(u)$ es algebraizable. Basta entonces
aplicar el Lema \sref{III.5.3.1} a la sucesión exacta
\[
  0\longrightarrow\Ker(u)\longrightarrow\sh{F}
  \longrightarrow\Im(u)\longrightarrow0.
\]
\end{proof}

\begin{lemma}[5.3.3]
\label{III.5.3.3}
Bajo las condiciones de \sref{III.5.1.1}, sea $h:Z\to Y$ un
morfismo de tipo finito, sea $\widehat{Z}$ la compleción de $Z$
a lo largo de $Z'=h^{-1}(Y')$, sea $g:Z\to X$ un $Y$-morfismo propio,
y sea $\widehat{g}:\widehat{Z}\to\widehat{X}$ su extensión a
las compleciones. Para todo $\sh{O}_{\widehat{Z}}$-módulo
algebraizable $\sh{F}'$, el módulo
$\widehat{g}_*(\sh{F}')$ es un
$\sh{O}_{\widehat{X}}$-módulo algebraizable.
\end{lemma}

\begin{proof}
Si $\sh{F}$ es un $\sh{O}_Z$-módulo coherente tal que
$\sh{F}'=\widehat{\sh{F}}$, el primer teorema de comparación
\sref{III.4.1.5} identifica
$\widehat{g}_*(\widehat{\sh{F}})$ con la compleción de
$g_*(\sh{F})$.
\end{proof}

\begin{lemma}[5.3.4]
\label{III.5.3.4}
Sea $X$ un esquema ordinario noetheriano, sea $X'$ un subconjunto
cerrado de $X$, sea $f:Z\to X$ un morfismo propio, y pongamos
\[
  Z'=f^{-1}(X'),\qquad
  \widehat{X}=X_{/X'},\qquad
  \widehat{Z}=Z_{/Z'}.
\]
Sea $\widehat{f}:\widehat{Z}\to\widehat{X}$ la extensión de $f$
a las compleciones. Sea $\sh{M}$ un ideal coherente de
$\sh{O}_X$ tal que, si
\[
  U=X-\Supp(\sh{O}_X/\sh{M}),
\]
la restricción $f^{-1}(U)\to U$ de $f$ es un isomorfismo. Entonces,
para todo $\sh{O}_{\widehat{X}}$-módulo coherente $\sh{F}$, existe
un entero $n>0$ tal que el núcleo y el conúcleo del homomorfismo
canónico
\[
  \rho_{\sh{F}}:
  \sh{F}\longrightarrow
  \widehat{f}_*\bigl(\widehat{f}^{\,*}(\sh{F})\bigr)
\]
\sref[0]{0.4.4.3} son anulados por
$\widehat{\sh{M}}^{\,n}$.
\end{lemma}

\begin{proof}
Podemos restringirnos al caso $X=\Spec(B)$, donde $B$ es un
anillo noetheriano; entonces $X'=\mathrm{V}(\mathfrak{K})$ para un ideal
$\mathfrak{K}$ de $B$. Nos reducimos primero al caso en que $B$ es
noetheriano ádico y $\mathfrak{K}$ es un ideal de definición.

Sea $B_1$ la compleción separada de $B$ para la
topología $\mathfrak{K}$-preádica.
\oldpage[III]{155}
Si $\mathfrak{K}_1=\mathfrak{K}B_1$, entonces $B_1$ es un anillo
noetheriano ádico y $\mathfrak{K}_1$ es un ideal de definición.
Pongamos $X_1=\Spec(B_1)$, y sea $h:X_1\to X$ el morfismo
correspondiente al homomorfismo canónico $B\to B_1$. Si
$X'_1=h^{-1}(X')$, entonces $X'_1=\mathrm{V}(\mathfrak{K}_1)$. Finalmente,
pongamos
\[
  Z_1=Z\times_X X_1,\qquad f_1:Z_1\longrightarrow X_1.
\]
El morfismo $f_1$ es propio \sref[II]{II.5.4.2}. Denotemos por
$\widehat{X}_1$ la compleción de $X_1$ a lo largo de $X'_1$, por
\[
  \widehat{Z}_1=
  Z_1\times_{X_1}\widehat{X}_1
\]
la compleción de $Z_1$ a lo largo de
$Z'_1=f_1^{-1}(X'_1)$, y por
$\widehat{f}_1$ la extensión de $f_1$ a las compleciones.

La extensión
$\widehat{h}:\widehat{X}_1\to\widehat{X}$ corresponde a la
aplicación identidad de $B_1$ y es, por tanto, un isomorfismo
\sref[I]{I.10.9.1}. El morfismo correspondiente
$\widehat{Z}_1\to\widehat{Z}$ es también, por consiguiente, un isomorfismo,
y estos isomorfismos identifican $\widehat{f}_1$ con
$\widehat{f}$. Además,
$\sh{M}_1=h^*(\sh{M})$ es un ideal coherente de $\sh{O}_{X_1}$ y
\[
  \Supp(\sh{O}_{X_1}/\sh{M}_1)
  =h^{-1}\bigl(\Supp(\sh{O}_X/\sh{M})\bigr)
\]
\sref[I]{I.9.1.13}. Así, si
$U_1=X_1-\Supp(\sh{O}_{X_1}/\sh{M}_1)$, entonces
$U_1=h^{-1}(U)$, y la restricción
$f_1^{-1}(U_1)\to U_1$ es un isomorfismo
\sref[I]{I.3.2.7}. Finalmente, $\widehat{\sh{M}}$ y
$\widehat{\sh{M}}_1$ son identificados por $\widehat{h}$
\sref[I]{I.10.9.5}. Todas las hipótesis del lema son, pues,
satisfechas por $X_1$, $X'_1$, $f_1$ y $\sh{M}_1$; en lo sucesivo
podemos suponer que $B$ es noetheriano ádico y que $\mathfrak{K}$ es
un ideal de definición.

Se tiene entonces $\widehat{X}=\Spf(B)$ y
$\sh{F}=N^\Delta$, donde $N$ es un $B$-módulo de tipo finito. Así,
$\sh{F}=\widehat{\sh{G}}$, donde $\sh{G}$ es el
$\sh{O}_X$-módulo coherente $\widetilde{N}$ \sref[I]{I.10.10.5}, y
\[
  \widehat{f}^{\,*}(\sh{F})
  =\widehat{f^*(\sh{G})}
\]
\sref[I]{I.10.9.5}. Por el primer teorema de comparación
\sref{III.4.1.5},
\[
  \widehat{f}_*
  \bigl(\widehat{f^*(\sh{G})}\bigr)
  \simeq
  \widehat{f_*\bigl(f^*(\sh{G})\bigr)},
\]
y, por \sref{III.5.1.3}, el homomorfismo canónico
$\rho_{\sh{F}}$ es precisamente $\widehat{\rho}_{\sh{G}}$.

El núcleo $\sh{P}$ y el conúcleo $\sh{R}$ de
\[
  \rho_{\sh{G}}:
  \sh{G}\longrightarrow f_*\bigl(f^*(\sh{G})\bigr)
\]
son $\sh{O}_X$-módulos coherentes, y sus restricciones a $U$ son
nulas. Existe, por tanto, un entero $n>0$ tal que
\[
  \sh{M}^n\sh{P}=\sh{M}^n\sh{R}=0
\]
\sref[I]{I.9.3.4}. Se deduce que
\[
  \widehat{\sh{M}}^{\,n}\widehat{\sh{P}}
  =
  \widehat{\sh{M}}^{\,n}\widehat{\sh{R}}
  =0
\]
\sref[I]{I.10.8.8} \sref[I]{I.10.8.10}, lo que demuestra el lema.
\end{proof}

\begin{env}[5.3.5]
\label{III.5.3.5}
\emph{Fin de la demostración del teorema de existencia.}
Conservemos las hipótesis de \sref{III.5.1.4}. Utilicemos la inducción
noetheriana \sref[0]{0.2.2.2}, suponiendo conocido el teorema para todo
subpreesquema cerrado $T$ de $X$ cuyo espacio subyacente es distinto
de $X$; aquí $\widehat{T}$ denota la compleción de $T$ a lo largo de
$T\cap X'$. Podemos suponer que $X$ no es vacío.

Como $f$ es separado y de tipo finito, el lema de Chow
\sref[II]{II.5.6.1} proporciona un $Y$-esquema $Z$ y un $Y$-morfismo
$g:Z\to X$ tales que el morfismo estructural $h:Z\to Y$ es
cuasiproyectivo, $g$ es proyectivo y sobreyectivo, y existe un
subconjunto abierto no vacío $U$ de $X$ para el cual
$g^{-1}(U)\to U$ es un isomorfismo. Sea $\sh{M}$ un ideal coherente
de $\sh{O}_X$ que define un subpreesquema cerrado de espacio subyacente
$X-U$ \sref[I]{I.5.2.2}, y sea $\sh{F}$ un
$\sh{O}_{\widehat{X}}$-módulo coherente cuyo soporte $E$ es propio sobre $Y$.
Denotemos por $\widehat{Z}$ la compleción de $Z$ a lo largo de
$h^{-1}(Y')$, y por
$\widehat{g}:\widehat{Z}\to\widehat{X}$ la extensión de $g$ a
las compleciones.

El módulo $\widehat{g}^{\,*}(\sh{F})$ es un
$\sh{O}_{\widehat{Z}}$-módulo coherente cuyo soporte está contenido en
$g^{-1}(E)$ y es, por tanto, propio sobre $Y$, pues $g$ es
proyectivo y, por tanto, propio \sref[II]{II.5.4.6}. Como $h$ es
cuasiproyectivo, $\widehat{g}^{\,*}(\sh{F})$ es algebraizable por
\sref{III.5.2.6}.
\oldpage[III]{156}
Se deduce de \sref{III.5.3.3}, puesto que $g$ es propio, que
\[
  \widehat{g}_*
  \bigl(\widehat{g}^{\,*}(\sh{F})\bigr)
\]
es un $\sh{O}_{\widehat{X}}$-módulo algebraizable.

Apliquemos \sref{III.5.3.4} a $\sh{F}$ y $g$. El núcleo $\sh{P}$
y el conúcleo $\sh{R}$ de
\[
  \rho_{\sh{F}}:
  \sh{F}\longrightarrow
  \widehat{g}_*\bigl(\widehat{g}^{\,*}(\sh{F})\bigr)
\]
son anulados por una potencia $\widehat{\sh{M}}^{\,n}$. Sea $T$
el subpreesquema cerrado de $X$ definido por $\sh{M}^n$, cuyo
espacio subyacente es $X-U$, y sea $j:T\to X$ la inmersión
canónica. Su extensión
$\widehat{j}:\widehat{T}\to\widehat{X}$ es la inmersión canónica
\sref[I]{I.10.14.7}, y podemos escribir
\[
  \sh{P}=
  \widehat{j}_*\bigl(\widehat{j}^{\,*}(\sh{P})\bigr),
  \qquad
  \sh{R}=
  \widehat{j}_*\bigl(\widehat{j}^{\,*}(\sh{R})\bigr).
\]
Como $U$ no es vacío, la hipótesis de inducción muestra que
$\widehat{j}^{\,*}(\sh{P})$ y
$\widehat{j}^{\,*}(\sh{R})$ son $\sh{O}_{\widehat{T}}$-módulos
algebraizables. El Lema \sref{III.5.3.3} muestra entonces
que $\sh{P}$ y $\sh{R}$ son algebraizables, y
\sref{III.5.3.2} demuestra finalmente que $\sh{F}$ es algebraizable.
\end{env}


\subsection{Aplicación: comparación entre morfismos de preesquemas ordinarios y morfismos de preesquemas formales. Preesquemas formales algebraizables}
\label{subsection:III.5.4}

\begin{theorem}[5.4.1]
\label{III.5.4.1}
Sea $A$ un anillo noetheriano ádico, sea $\mathfrak{J}$ un ideal
de definición de $A$, y pongamos
\[
  S=\Spec(A),\qquad S'=\mathrm{V}(\mathfrak{J}).
\]
Sea $u:X\to S$ un morfismo propio y sea $v:Y\to S$ un
morfismo separado de tipo finito. Sean
$\widehat{S}$, $\widehat{X}$ y $\widehat{Y}$ las compleciones
de $S$, $X$ e $Y$ a lo largo de $S'$, $u^{-1}(S')$ y $v^{-1}(S')$,
respectivamente. Para todo $S$-morfismo $f:X\to Y$, sea
$\widehat{f}:\widehat{X}\to\widehat{Y}$ su extensión a las
compleciones. Entonces la aplicación $f\mapsto\widehat{f}$ es una biyección
\[
  \Hom_S(X,Y)
  \xrightarrow{\ \sim\ }
  \Hom_{\widehat{S}}(\widehat{X},\widehat{Y}).
\]
\end{theorem}

\begin{proof}
Demostremos primero la inyectividad. Supongamos que $f,g:X\to Y$ son
$S$-morfismos tales que $\widehat{f}=\widehat{g}$. Existe entonces
un entorno abierto $V$ de $X'=u^{-1}(S')$ sobre el cual $f$ y $g$
coinciden \sref[I]{I.10.9.4}. Como $u$ es cerrado,
\sref{III.5.1.3.1} da $V=X$, y por tanto $f=g$.

Demostremos ahora la sobreyectividad. Sea
$h:\widehat{X}\to\widehat{Y}$ un
$\widehat{S}$-morfismo. Pongamos
\[
  Z=X\times_S Y,
\]
y denotemos las proyecciones canónicas por $p:Z\to X$ y
$q:Z\to Y$. El preesquema $Z$ es de tipo finito sobre $S$
\sref[I]{I.6.3.4} y, por tanto, noetheriano. Sea
$\widehat{Z}$ su compleción a lo largo de
\[
  Z'=p^{-1}\bigl(u^{-1}(S')\bigr).
\]
El preesquema formal $\widehat{Z}$ se identifica canónicamente con
$\widehat{X}\times_{\widehat{S}}\widehat{Y}$, y sus proyecciones
sobre $\widehat{X}$ y $\widehat{Y}$ se identifican con
$\widehat{p}$ y $\widehat{q}$ \sref[I]{I.10.9.7}.

Como $Y$ es separado sobre $S$, $\widehat{Y}$ es separado sobre
$\widehat{S}$ \sref[I]{I.10.15.7}. El morfismo grafo
\[
  \Gamma_h=(1_{\widehat{X}},h):
  \widehat{X}\longrightarrow\widehat{Z}
\]
es, por consiguiente, una inmersión cerrada \sref[I]{I.10.15.4}. Sea
$\mathfrak{T}$ el subpreesquema formal cerrado de
$\widehat{Z}$ asociado a esta inmersión, y sea
$j:\mathfrak{T}\to\widehat{Z}$ la inmersión canónica. Así,
\[
  \Gamma_h=j\circ w,
\]
donde $w:\widehat{X}\to\mathfrak{T}$ es un isomorfismo
\sref[I]{I.10.14.3} cuyo inverso es
$\widehat{p}\circ j$. Además, $\mathfrak{T}$ es propio sobre
$\widehat{S}$ porque $\widehat{X}$ lo es. Se deduce de
\sref{III.5.1.8} que existe un subpreesquema cerrado $T$ de $Z$
tal que
\[
  \mathfrak{T}
  =\widehat{T}
  =T_{/(T\cap Z')},
\]
y que $j=\widehat{i}$, donde $i:T\to Z$ es la inmersión
canónica \sref[I]{I.10.14.7}.

El morfismo $p\circ i:T\to X$ es un isomorfismo: su extensión a
las compleciones es el isomorfismo
$\widehat{p}\circ\widehat{i}$, y basta aplicar
\sref{III.4.6.8}, observando como antes que $S$ es el único
entorno de $S'$ en $S$.
\oldpage[III]{157}
Sea $g:X\to T$ el inverso de $p\circ i$, y pongamos
\[
  f=q\circ i\circ g.
\]
Entonces $f:X\to Y$ tiene por grafo $\Gamma_f=i\circ g$. Como
$\widehat{g}$ es el inverso de
$\widehat{p\circ i}$, se tiene
\[
  \widehat{\Gamma_f}
  =\widehat{i}\circ\widehat{g}
  =j\circ w
  =\Gamma_h.
\]
Pero $\widehat{\Gamma_f}=\Gamma_{\widehat{f}}$
\sref[I]{I.10.9.8}; por tanto, $h=\widehat{f}$, lo que termina la
demostración.
\end{proof}

En lenguaje categórico, el funtor
$X\mapsto\widehat{X}$ es, por tanto, plenamente fiel
\sref[0]{0.8.1.6} de la categoría de los preesquemas propios sobre
$\Spec(A)$ en la categoría de los preesquemas formales propios sobre
$\Spf(A)$, para todo anillo noetheriano ádico $A$. Proporciona, por
consiguiente, una equivalencia entre la primera categoría y una
subcategoría plena de la segunda. Los objetos de esta subcategoría se
denominarán en lo sucesivo \emph{preesquemas formales algebraizables}.
Para tal preesquema formal $\mathfrak{X}$, existe un preesquema
ordinario $X$, propio sobre $\Spec(A)$ y determinado salvo isomorfismo
único, tal que $\mathfrak{X}\simeq\widehat{X}$.

\begin{env}[5.4.2]
\label{III.5.4.2}
\emph{Escolio.}
Con la notación de \sref{III.5.4.1}, pongamos
\[
  S_n=\Spec(A/\mathfrak{J}^{n+1}),\qquad
  X_n=X\times_S S_n,\qquad
  Y_n=Y\times_S S_n.
\]
Se deduce de \sref{III.5.4.1} y \sref[I]{I.10.12.3} que
dar un $S$-morfismo $f:X\to Y$ equivale a dar un sistema inductivo
ádico $(S_n)$ \sref[I]{I.10.12.2} de
$S_n$-morfismos $f_n:X_n\to Y_n$.
\end{env}

\begin{remark}[5.4.3]
\label{III.5.4.3}
Contrariamente a lo que podría sugerir el teorema de existencia
\sref{III.5.1.6}, existen preesquemas formales propios sobre $\Spf(A)$
que \emph{no son algebraizables}, así como existen espacios analíticos
compactos que no proceden de variedades algebraicas complejas. Más
adelante encontraremos tales preesquemas en la «teoría de moduli»,
que, cuando el cuerpo base es $\mathbb{C}$, trata precisamente de las
variaciones infinitesimales de la estructura compleja de una variedad
algebraica completa; es bien sabido que tales variaciones pueden dar
lugar a variedades analíticas que no son algebraicas.
\end{remark}

\begin{proposition}[5.4.4]
\label{III.5.4.4}
Sea $A$ un anillo noetheriano ádico, pongamos
$\mathfrak{S}=\Spf(A)$, y sean
\[
  g:\mathfrak{X}\to\mathfrak{S},
  \qquad
  h:\mathfrak{Y}\to\mathfrak{S}
\]
dos morfismos propios de preesquemas formales. Sea
$f:\mathfrak{X}\to\mathfrak{Y}$ un
$\mathfrak{S}$-morfismo. Si $f$ es finito y
$\mathfrak{Y}$ es algebraizable, entonces $\mathfrak{X}$ es
algebraizable.
\end{proposition}

\begin{proof}
Las hipótesis sobre $g$ y $h$ implican ya que $f$ es propio
\sref{III.3.4.1}; para que $f$ sea finito, basta que la fibra
$f^{-1}(y)$ sea finita para todo $y\in\mathfrak{Y}$
\sref[0]{0.4.8.11}. Por hipótesis,
\[
  \sh{B}=f_*(\sh{O}_{\mathfrak{X}})
\]
es una $\sh{O}_{\mathfrak{Y}}$-álgebra coherente
\sref[0]{0.4.8.6}. Escribamos
$\mathfrak{Y}=\widehat{Y}$ y $h=\widehat{w}$, donde
$w:Y\to\Spec(A)$ es un morfismo propio de preesquemas ordinarios. El
teorema de existencia proporciona una $\sh{O}_Y$-álgebra coherente $\sh{C}$
tal que $\sh{B}=\widehat{\sh{C}}$. Pongamos
\[
  X=\Spec(\sh{C}),
\]
y sea $u:X\to Y$ el morfismo estructural. La definición de
$\mathfrak{X}$ a partir de $\sh{B}$ \sref[0]{0.4.8.7} muestra entonces que
$\mathfrak{X}$ es canónicamente isomorfo a $\widehat{X}$ y
que $f$ se identifica con $\widehat{u}$; basta comprobarlo
cuando $Y$ es afín.
\end{proof}

La Proposición \sref{III.5.1.8} es un caso particular de
\sref{III.5.4.4}.

\begin{theorem}[5.4.5]
\label{III.5.4.5}
Sea $A$ un anillo noetheriano ádico, sea $\mathfrak{J}$ un ideal
de definición de $A$, pongamos
\[
  S=\Spec(A),\qquad
  \mathfrak{S}=\widehat{S}=\Spf(A),
\]
y sea $f:\mathfrak{X}\to\mathfrak{S}$ un morfismo propio de
preesquemas formales. Pongamos
\[
  S_k=\Spec(A/\mathfrak{J}^{k+1}),\qquad
  X_k=\mathfrak{X}\times_{\mathfrak{S}}S_k,
\]
y, para todo $\sh{O}_{\mathfrak{X}}$-módulo $\sh{F}$, pongamos
\[
  \sh{F}_k
  =\sh{F}\otimes_{\sh{O}_{\mathfrak{X}}}\sh{O}_{X_k}
  =\sh{F}/\mathfrak{J}^{k+1}\sh{F}.
\]
\oldpage[III]{158}
Sea $\sh{L}$ un
$\sh{O}_{\mathfrak{X}}$-módulo invertible, y supongamos que
\[
  \sh{L}_0=\sh{L}/\mathfrak{J}\sh{L}
\]
es un $\sh{O}_{X_0}$-módulo amplio. Entonces $\mathfrak{X}$ es
algebraizable. Además, si $X$ es un $S$-preesquema propio tal
que $\mathfrak{X}\simeq\widehat{X}$, existe un
$\sh{O}_X$-módulo amplio $\sh{M}$ tal que
$\sh{L}\simeq\widehat{\sh{M}}$; en particular, $X$ es proyectivo
sobre $S$.
\end{theorem}

\begin{proof}
Apliquemos \sref{III.5.2.3} con
$\sh{F}=\sh{O}_{\mathfrak{X}}$. Existe un entero $n_0$ tal
que, para $n\geq n_0$, el homomorfismo canónico
\[
  \Gamma(\mathfrak{X},\sh{L}^{\otimes n})
  \longrightarrow
  \Gamma(X_0,\sh{L}_0^{\otimes n})
\]
es sobreyectivo. Elijamos $n\geq n_0$ suficientemente grande para que
$\sh{L}_0^{\otimes n}$ sea muy amplio para $S_0$
\sref[II]{II.4.5.10}. Como
$f_0:X_0\to S_0$ es propio,
$\Gamma(X_0,\sh{L}_0^{\otimes n})$ es un $A$-módulo de tipo finito
\sref{III.3.2.1}. Existe, por consiguiente, un
$A$-submódulo de tipo finito
\[
  E\subset\Gamma(\mathfrak{X},\sh{L}^{\otimes n})
\]
cuya imagen en $\Gamma(X_0,\sh{L}_0^{\otimes n})$ es el módulo
entero.

Para todo $k\geq0$, consideremos el homomorfismo
\[
  u_k:
  E/\mathfrak{J}^{k+1}E
  \longrightarrow
  \Gamma(X_k,\sh{L}_k^{\otimes n})
\]
inducido por la inyección canónica
$E\to\Gamma(\mathfrak{X},\sh{L}^{\otimes n})$. El módulo
$(f_k)_*(\sh{L}_k^{\otimes n})$ es cuasi-coherente, y
\[
  \Gamma\bigl(S_k,(f_k)_*(\sh{L}_k^{\otimes n})\bigr)
  =
  \Gamma(X_k,\sh{L}_k^{\otimes n}).
\]
Así, $u_k$ define un homomorfismo
\[
  \widetilde{u}_k:
  \widetilde{E/\mathfrak{J}^{k+1}E}
  \longrightarrow
  (f_k)_*(\sh{L}_k^{\otimes n}),
\]
y, por tanto, por adjunción, un homomorfismo
\[
  \widetilde{u}_k^{\,\sharp}:
  f_k^*\bigl(\widetilde{E/\mathfrak{J}^{k+1}E}\bigr)
  \longrightarrow
  \sh{L}_k^{\otimes n}.
\]
Pongamos
\[
  \sh{G}_k=
  f_k^*\bigl(\widetilde{E/\mathfrak{J}^{k+1}E}\bigr).
\]
Se tiene
$\sh{G}_0=\sh{G}_k/\mathfrak{J}\sh{G}_k$
\sref[I]{I.9.1.5}, y
\[
  \widetilde{u}_0^{\,\sharp}:
  \sh{G}_0/\mathfrak{J}\sh{G}_0
  \longrightarrow
  \sh{L}_0^{\otimes n}/
  \mathfrak{J}\sh{L}_0^{\otimes n}
\]
se obtiene de $\widetilde{u}_k^{\,\sharp}$ pasando a los
cocientes.

Por la definición de $E$, $\widetilde{u}_0^{\,\sharp}$ es el
homomorfismo canónico
\[
  \sigma:
  f_0^*\bigl((f_0)_*(\sh{L}_0^{\otimes n})\bigr)
  \longrightarrow
  \sh{L}_0^{\otimes n}.
\]
Como $\sh{L}_0^{\otimes n}$ es muy amplio,
$\widetilde{u}_0^{\,\sharp}$ es sobreyectivo
\sref[II]{II.4.4.3}. El lema de Nakayama muestra entonces que todo
$\widetilde{u}_k^{\,\sharp}$ es sobreyectivo.

Cada $\widetilde{u}_k^{\,\sharp}$ define, por tanto,
\sref[II]{II.4.2.2} un $S_k$-morfismo
\[
  g_k:X_k\longrightarrow
  P_k=\mathbf{P}(E/\mathfrak{J}^{k+1}E).
\]
Para $h\leq k$,
\[
  P_h=P_k\times_{S_k}S_h
\]
por \sref[II]{II.4.1.3}. Las relaciones entre los
$\widetilde{u}_k^{\,\sharp}$ y \sref[II]{II.4.2.10} muestran que
$(g_k)$ es un sistema inductivo ádico $(S_k)$
\sref[I]{I.10.12.2}. Los $g_k$ definen, por consiguiente, un
$\mathfrak{S}$-morfismo de preesquemas formales
\[
  g:\mathfrak{X}\longrightarrow\mathfrak{P},
\]
donde $\mathfrak{P}$ es el límite inductivo del sistema $(P_k)$,
o, equivalentemente, la compleción $\widehat{P}$ de
$P=\mathbf{P}(E)$. El hecho de que
$\sh{L}_0^{\otimes n}$ sea muy amplio implica que $g_0$ es una inmersión cerrada
\sref[II]{II.4.4.3}. Por tanto, $g$ es una inmersión cerrada de
preesquemas formales \sref[0]{0.4.8.10}, y $\mathfrak{X}$ es algebraizable
por \sref{III.5.1.8}.

El teorema de existencia \sref{III.5.1.6} muestra ahora que $\sh{L}$ es
la compleción $\widehat{\sh{M}}$ de un
$\sh{O}_X$-módulo invertible $\sh{M}$. Además,
$\sh{L}^{\otimes n}$ es la compleción de
$\sh{M}^{\otimes n}$ \sref[I]{I.10.8.10}, y los homomorfismos
$\widetilde{u}_k^{\,\sharp}$ definen un único homomorfismo
\[
  v:f^*(\widetilde{E})
  \longrightarrow
  \sh{M}^{\otimes n}
\]
\sref{III.5.1.7}. Como los
$\widetilde{u}_k^{\,\sharp}$ son sobreyectivos,
$\widehat{v}$ es sobreyectivo \sref[I]{I.10.11.5}, y por tanto también
lo es $v$ \sref{III.5.1.3}. Sea $r:X\to P$ el morfismo definido
por $v$ \sref[II]{II.4.2.2}. Su extensión a las compleciones es
$g$. Como $g$ es una inmersión cerrada, también lo es $r$, por
\sref{III.5.1.8} y \sref{III.5.4.1}. Se deduce que
$\sh{M}^{\otimes n}$ es muy amplio \sref[II]{II.4.4.6} y, por tanto,
que $\sh{M}$ es amplio \sref[II]{II.4.5.10}.
\end{proof}

\begin{remark}[5.4.6]
\label{III.5.4.6}
Sea $A$ un anillo noetheriano ádico, pongamos
$\mathfrak{S}=\Spf(A)$, y sea
$f:\mathfrak{X}\to\mathfrak{S}$ un morfismo propio de preesquemas
formales. Sea $\sh{N}$ el ideal coherente de
$\sh{O}_{\mathfrak{X}}$ tal que, para todo subconjunto abierto
formal afín $U$ de $\mathfrak{X}$,
\[
  \Gamma(U,\sh{N})
\]
es el nilradical de
$\Gamma(U,\sh{O}_{\mathfrak{X}})$. La existencia de este ideal
se deduce fácilmente de \sref[I]{I.10.10.2} y del hecho de que
todo homomorfismo de anillos $B\to C$ envía el nilradical de $B$ en
el de $C$. Sea $\mathfrak{X}'$ el subpreesquema formal cerrado de
$\mathfrak{X}$ definido por $\sh{N}$
\sref[I]{I.10.14.2}. Sería interesante saber si la
algebraizabilidad de $\mathfrak{X}'$ implica la algebraizabilidad del
propio $\mathfrak{X}$.

Sin duda se llegaría a una solución de este problema si se supiera
clasificar, por ejemplo mediante invariantes
\oldpage[III]{159}
de naturaleza cohomológica, las \emph{extensiones} de un haz estructural
$\sh{O}_{\mathfrak{X}}$---para un preesquema ordinario o
formal---por un ideal de cuadrado nulo. Equivalentemente, estas son las
$\sh{O}_{\mathfrak{X}}$-álgebras $\sh{A}$ tales que
$\sh{O}_{\mathfrak{X}}\simeq\sh{A}/\sh{J}$, donde $\sh{J}$ es un
ideal de cuadrado nulo de $\sh{A}$.
\end{remark}


\subsection{Una descomposición de ciertos preesquemas}
\label{subsection:III.5.5}

\begin{proposition}[5.5.1]
\label{III.5.5.1}
Sea $A$ un anillo noetheriano ádico, sea $\mathfrak{J}$ un ideal
de definición de $A$, y pongamos $Y=\Spec(A)$. Sea $f:X\to Y$ un
morfismo separado de tipo finito, y pongamos
\[
  Y_0=\Spec(A/\mathfrak{J}),\qquad
  X_0=X\times_Y Y_0=f^{-1}(Y_0).
\]
Sea $Z_0$ un subconjunto abierto de $X_0$ que es propio sobre $Y_0$.
Entonces existe un subconjunto abierto y cerrado $Z$ de $X$, propio sobre
$Y$, tal que
\[
  Z\cap X_0=Z_0.
\]
\end{proposition}

\begin{proof}
Por hipótesis, existe un subconjunto abierto $T$ de $X$ tal que
$T\cap X_0=Z_0$. Sea $\widehat{T}$ la compleción a lo largo de $Z_0$
del preesquema inducido por $X$ sobre $T$. El soporte de
$\sh{O}_{\widehat{T}}$ es $Z_0$, que es propio sobre $Y_0$;
por tanto, $\widehat{T}$ es propio sobre
$\widehat{Y}=\Spf(A)$ \sref{III.3.4.1}.

Se deduce de \sref{III.5.1.8} que existe un subpreesquema cerrado
$Z$ de $T$, propio sobre $Y$, tal que, si
$i:Z\to T$ es la inmersión canónica, entonces
\[
  \widehat{i}:\widehat{Z}\longrightarrow\widehat{T}
\]
es un isomorfismo; aquí $\widehat{Z}$ es la compleción de $Z$
a lo largo de $Z_0$. Por \sref{III.4.6.8}, existe un entorno abierto
$V$ de $Z_0$ en $T$ tal que
$i^{-1}(V)\to V$ es un isomorfismo. Pero $i^{-1}(V)$ es un
entorno de $Z_0$ en $Z$ y, por tanto, es igual a $Z$ por
\sref{III.5.1.3.1}. Así, $Z$ es abierto en $T$, y por tanto en $X$.
También es cerrado en $X$, pues $Z$ es propio sobre $Y$ y $X$ es
separado sobre $Y$. Esto demuestra la proposición.
\end{proof}

\begin{corollary}[5.5.2]
\label{III.5.5.2}
Si $X_0$ es propio sobre $Y_0$, entonces $X$ es la unión de dos
subconjuntos abiertos disjuntos $Z$ y $Z'$ tales que $Z$ es propio sobre
$Y$ y contiene a $X_0$. Además, todo subconjunto cerrado $P$ de $X$
que es propio sobre $Y$ está contenido en $Z$.
\end{corollary}

\begin{proof}
Apliquemos \sref{III.5.5.1} con $Z_0=X_0$, y pongamos $Z'=X-Z$. Para la
última afirmación, $P\cap Z'$ es cerrado en $P$ y, por tanto, propio
sobre $Y$. Si no fuera vacío, $f(P\cap Z')$ sería un subconjunto cerrado
no vacío de $Y$ y, por tanto, cortaría a $Y_0$
\sref{III.5.1.3.1}. Esto contradice
$Z'\cap X_0=\varnothing$ y demuestra que $P\subset Z$.
\end{proof}

\bigskip
\noindent\emph{(Continuará.)}

% La página impresa 504 está en blanco.
\clearpage
\thispagestyle{empty}
\null
\clearpage
\thispagestyle{plain}

\oldpage[III]{161}
\section*{BIBLIOGRAFÍA (continuación)}

\begin{flushleft}
[27] P. Gabriel, \emph{Des catégories abéliennes}, tesis, París (1961).

\medskip
[28] J. E. Roos, \emph{Sur les foncteurs dérivés de $\varprojlim$.
Applications}, C. R. Acad. Sci. Paris, vol. CCLII, 1961, pp. 3702--3704.

\medskip
[29] H. Cartan, \emph{Séminaire de l'École Normale Supérieure},
13.er año (1960--1961), exposición n.º 11.
\end{flushleft}

\clearpage
\oldpage[III]{162}
\section*{ÍNDICE DE NOTACIONES}

\begin{flushleft}
$\Set$: 0, 8.1.1.

$h_X(Y)$, $h_X(u)$ ($X$, $Y$ objetos de una categoría $\C$, $u$ un morfismo
de $\C$): 0, 8.1.1.

$h_w(Y)$ ($Y$ un objeto y $w$ un morfismo de $\C$): 0, 8.1.2.

$Z_k(E_r^{pq})$, $B_k(E_r^{pq})$ ($k\geq r$): 0, 11.1.1.

$Z_\infty(E_2^{pq})$, $B_\infty(E_2^{pq})$: 0, 11.1.1.

$K^\bullet$, $K_\bullet$ (complejos): 0, 11.2.1.

$K^{\bullet\bullet}$, $K_{\bullet\bullet}$ (bicomplejos): 0, 11.2.1.

$B^i(K^\bullet)$, $Z^i(K^\bullet)$, $H^i(K^\bullet)$,
$H^\bullet(K^\bullet)$: 0, 11.2.1.

$B_i(K_\bullet)$, $Z_i(K_\bullet)$, $H_i(K_\bullet)$,
$H_\bullet(K_\bullet)$: 0, 11.2.1.

$T(K^\bullet)$, $T(K_\bullet)$ ($T$ un funtor): 0, 11.2.1.

$E(K^\bullet)$ ($K^\bullet$ un complejo filtrado): 0, 11.2.2.

$K^{i,\bullet}$, $K^{\bullet,j}$, $Z_{\mathrm{II}}^p(K^{i,\bullet})$,
$B_{\mathrm{II}}^p(K^{i,\bullet})$, $H_{\mathrm{II}}^p(K^{i,\bullet})$,
$Z_{\mathrm I}^p(K^{\bullet,j})$, $B_{\mathrm I}^p(K^{\bullet,j})$,
$H_{\mathrm I}^p(K^{\bullet,j})$: 0, 11.3.1.

$H_{\mathrm{II}}^p(K^{\bullet\bullet})$,
$Z_{\mathrm I}^q(H_{\mathrm{II}}^p(K^{\bullet\bullet}))$,
$B_{\mathrm I}^q(H_{\mathrm{II}}^p(K^{\bullet\bullet}))$,
$H_{\mathrm I}^q(H_{\mathrm{II}}^p(K^{\bullet\bullet}))$,
$H^n(K^{\bullet\bullet})$: 0, 11.3.1.

$F_{\mathrm I}^p(K^{\bullet\bullet})$,
$F_{\mathrm{II}}^p(K^{\bullet\bullet})$, ${}'E(K^{\bullet\bullet})$,
${}''E(K^{\bullet\bullet})$: 0, 11.3.2.

$K_{i,\bullet}$, $K_{\bullet,j}$, $H_p^{\mathrm{II}}(K_{i,\bullet})$,
$H_p^{\mathrm I}(K_{\bullet,j})$, $H_p^{\mathrm{II}}(K_{\bullet\bullet})$,
$H_p^{\mathrm I}(K_{\bullet\bullet})$,
$H_q^{\mathrm I}(H_p^{\mathrm{II}}(K_{\bullet\bullet}))$,
$H_q^{\mathrm{II}}(H_p^{\mathrm I}(K_{\bullet\bullet}))$,
$H_n(K_{\bullet\bullet})$: 0, 11.3.5.

$\RR^\bullet T(K^\bullet)$ ($T$ un funtor covariante aditivo): 0, 11.4.3.

$\RR^\bullet T(K^\bullet,K'^{\bullet})$ ($T$ un bifuntor covariante
aditivo): 0, 11.4.6.

$\LL_\bullet T(K_\bullet)$ ($T$ un funtor covariante aditivo): 0, 11.6.2.

$\LL_\bullet T(K_\bullet,K'_\bullet)$ ($T$ un bifuntor covariante
aditivo): 0, 11.6.6.

$\LL_\bullet T(K_{\bullet\bullet})$ ($T$ un funtor covariante aditivo):
0, 11.7.2.

$|\sigma|$, $\Sigma(A)$, $\mathrm C_\bullet(A)$, $\mathrm L_\bullet(A)$
($A$ un conjunto finito, $\sigma$ un símplex de $A$): 0, 11.8.1.

$\mathrm C^\bullet(A,B;\sh{S})$, $\mathrm L^\bullet(A,B;\sh{S})$
($A$, $B$ conjuntos finitos, $\sh{S}$ un sistema de coeficientes): 0, 11.8.4.

$\mathrm P^\bullet(A,B;\sh{S})$ ($A$, $B$ conjuntos finitos, $\sh{S}$ un
sistema de coeficientes): 0, 11.8.6.

$\mathrm C^\bullet(A,B;\sh{S}^\bullet)$,
$\mathrm L^\bullet(A,B;\sh{S}^\bullet)$ ($A$, $B$ conjuntos finitos,
$\sh{S}^\bullet$ un complejo de sistemas de coeficientes): 0, 11.8.8.

$\mathrm P^\bullet(A,B;\sh{S}^\bullet)$ ($A$, $B$ conjuntos finitos,
$\sh{S}^\bullet$ un complejo de sistemas de coeficientes): 0, 11.8.10.

$\chi(M^\bullet)$ ($M^\bullet$ un complejo finito de módulos de longitud finita):
0, 11.10.1.

$\check C^\bullet(\mathfrak U,\sh{F})$ ($\sh{F}$ un haz de grupos abelianos
sobre $X$, $\mathfrak U$ un recubrimiento abierto de $X$): 0, 12.1.2.

$\RR^p f_*(\sh{F})$ ($f:X\to Y$ un morfismo de espacios anillados,
$\sh{F}$ un haz de $\sh{O}_X$-módulos): 0, 12.2.1.

$\sh{H}^p(f,\sh{K}^\bullet)$,
$\sh{H}^p_f(\sh{K}^\bullet)$, ${}'E(f,\sh{K}^\bullet)$,
${}''E(f,\sh{K}^\bullet)$ ($f:X\to Y$ un morfismo de espacios anillados,
$\sh{K}^\bullet$ un complejo de $\sh{O}_X$-módulos): 0, 12.4.1.

$\mathbf H^\bullet(X,\sh{K}^\bullet)$,
$\mathbf H^\bullet(V,\sh{K}^\bullet)$ ($X$ un espacio anillado,
$\sh{K}^\bullet$ un complejo de $\sh{O}_X$-módulos, $V$ abierto en $X$):
0, 12.4.2.

$\operatorname{gr}^\bullet(A)$ ($A$ un sistema proyectivo que satisface (ML)):
0, 13.4.2.

$E(A)$ ($A$ un objeto filtrado): 0, 13.6.4.

$K_\bullet(\mathbf f)$, $i_{\mathbf f}$ ($\mathbf f$ una familia finita de
elementos de un anillo): III, 1.1.1.

$K^\bullet(\mathbf f,M)$, $K_\bullet(\mathbf f,M)$ ($\mathbf f$ una familia
finita de elementos de un anillo $A$, $M$ un $A$-módulo): III, 1.1.2.

$H^\bullet(\mathbf f,M)$, $H_\bullet(\mathbf f,M)$ ($\mathbf f$ una familia
finita de elementos de un anillo $A$, $M$ un $A$-módulo): III, 1.1.3.

$C^\bullet((\mathbf f),M)$, $H^\bullet((\mathbf f),M)$ ($\mathbf f$ una
familia finita de elementos de un anillo $A$, $M$ un $A$-módulo): III, 1.1.6.

$H^\bullet(U,\sh{F}^{(*)})$,
$H^\bullet(\mathfrak U,\sh{F}^{(*)})$ ($\sh{F}$ un
$\sh{O}_X$-módulo cuasi-coherente, $U$ abierto en $X$, $\mathfrak U$ un recubrimiento abierto de $X$):
III, 2.1.1.
\end{flushleft}

\clearpage
\oldpage[III]{163}
\section*{ÍNDICE TERMINOLÓGICO}

\begin{flushleft}
Término de convergencia de una sucesión espectral: 0, 11.1.1.

Algebraizable ($\sh{O}_X$-módulo): III, 5.2.1.

Algebraizable (esquema formal): III, 5.4.2.

Analíticamente íntegro (anillo): III, 4.3.6.

Aplicación cuasicompacta: 0, 9.1.1.

Aumentación de una resolución: 0, 11.4.1.

Característica de Euler--Poincaré de un complejo: 0, 11.10.1.

Característica de Euler--Poincaré de un $\sh{O}_X$-módulo coherente
($X$ proyectivo sobre $\Spec(A)$, $A$ un anillo artiniano): III, 2.5.1.

Cocadena bialternada: 0, 11.8.4.

Complejo definido por un bicomplejo: 0, 11.3.1.

Complejo de álgebra exterior: III, 1.1.1.

Condición (ML): 0, 13.1.2.

Constructible (subconjunto, conjunto): 0, 9.1.2.

Constructible (función): 0, 9.3.1.

Producto cup: 0, 12.1.2.

Dihomomorfismo para estructuras algebraicas sobre categorías: 0, 8.2.1.

Exacto (subconjunto) en una categoría abeliana: III, 3.1.1.

Esencialmente constante (sistema proyectivo): 0, 13.4.2.

Factorización de Stein: III, 4.3.3.

Codiscreta, coseparada, discreta, exhaustiva, finita, separada
(filtración): 0, 11.1.3.

Final (objeto) de una categoría: 0, 8.1.10.

Finito (morfismo) de preesquemas formales: III, 4.8.2.

Finito (preesquema formal) sobre $\mathfrak Y$,
$\mathfrak Y$-finito (preesquema formal): III, 4.8.2.

Funtor covariante canónico $\C\to\CHom(\C\op,\Set)$: 0, 8.1.2.

Género aritmético: III, 2.5.1.

Geométricamente conexa (fibra): III, 4.3.4.

Homomorfismo de estructuras algebraicas sobre categorías: 0, 8.2.1.

Hipercohomología de un funtor respecto de un complejo $K^\bullet$:
0, 11.4.3.

Hipercohomología de un bifuntor respecto de dos complejos
$K^\bullet$, $K'^{\bullet}$: 0, 11.4.6.

Hiperhomología de un funtor respecto de un complejo $K_\bullet$:
0, 11.6.2.

Hiperhomología de un bifuntor respecto de dos complejos $K_\bullet$,
$K'_\bullet$: 0, 11.6.6.

Hiperhomología de un funtor respecto de un bicomplejo: 0, 11.7.2.

Localmente constructible (conjunto): 0, 9.1.11.

Ley de composición externa (interna) sobre los objetos de una categoría: 0, 8.2.1.

\oldpage[III]{164}

$S$-$\C$-módulo filtrado, $\operatorname{gr}^\bullet(S)$-$\C$-módulo graduado,
$\operatorname{gr}^{\bullet\bullet}(S)$-$\C$-módulo bigraduado: 0, 13.6.6.

Morfismo de sucesiones espectrales: 0, 11.1.2.

Número geométrico de componentes conexas de una fibra: III, 4.3.4.

Objeto grupo en $\C$, $\C$-grupo, $\C$-anillo, $\C$-módulo: 0, 8.2.3.

Objeto de bordes, objeto de cobordes, objeto de cociclos, objeto de
ciclos: 0, 11.2.1.

Objeto de imágenes universales: 0, 13.1.1.

Objeto graduado asociado a un sistema proyectivo que satisface (ML): 0, 13.4.2.

Objeto graduado acotado inferiormente (superiormente): 0, 11.2.1.

Subconjunto propio (sobre $\mathfrak Y$) de un preesquema formal: III, 3.4.1.

Plena (subcategoría): 0, 8.1.5.

Plenamente fiel (funtor): 0, 8.1.5.

Polinomio de Hilbert: III, 2.5.3.

Propio (morfismo) de preesquemas formales: III, 3.4.1.

Representable (funtor): 0, 8.1.8.

Cohomológica, derecha, izquierda, homológica, inyectiva, libre, plana, proyectiva
(resolución): 0, 11.4.1.

Resolución de Cartan--Eilenberg derecha, inyectiva: 0, 11.4.2.

Resolución de Cartan--Eilenberg izquierda, proyectiva: 0, 11.6.1.

Retrocompacto (conjunto): 0, 9.1.1.

Estricto (sistema proyectivo): 0, 13.4.2.

Sucesión espectral en una categoría abeliana: 0, 11.1.1.

Débilmente convergente, regular, coregular, birregular (sucesión espectral):
0, 11.1.3.

Sucesión espectral degenerada: 0, 11.1.6.

Sucesión espectral de un funtor relativa a un objeto filtrado: 0, 13.6.4.

Sucesiones espectrales de un bicomplejo: 0, 11.3.2 y 11.3.5.

Sucesiones espectrales de hipercohomología de un funtor respecto de un complejo:
0, 11.4.3.

Sucesiones espectrales de hiperhomología de un funtor respecto de un complejo:
0, 11.6.2.

Sucesiones espectrales de hiperhomología de un funtor respecto de un bicomplejo:
0, 11.7.2.

Sistema de coeficientes: 0, 11.8.4.

Unirrama (anillo, punto): III, 4.3.6.

Universalmente abierto (morfismo): III, 4.3.9.
\end{flushleft}

\clearpage
\oldpage[III]{165}
\section*{ÍNDICE GENERAL}

\begin{flushleft}
\textsc{Capítulo 0.} --- \textbf{Preliminares (continuación)} \dotfill 5

\medskip
\S~8. Funtores representables \dotfill 5

8.1. Funtores representables \dotfill 5

8.2. Estructuras algebraicas en categorías \dotfill 9

\medskip
\S~9. Conjuntos constructibles \dotfill 12

9.1. Conjuntos constructibles \dotfill 12

9.2. Subconjuntos constructibles de espacios noetherianos \dotfill 14

9.3. Funciones constructibles \dotfill 16

\medskip
\S~10. Suplemento sobre los módulos planos \dotfill 17

10.1. Relaciones entre módulos planos y módulos libres \dotfill 17

10.2. Criterios locales de planitud \dotfill 18

10.3. Existencia de extensiones planas de anillos locales \dotfill 20

\medskip
\S~11. Suplemento sobre álgebra homológica \dotfill 23

11.1. Recordatorio sobre las sucesiones espectrales \dotfill 23

11.2. La sucesión espectral de un complejo filtrado \dotfill 27

11.3. Las sucesiones espectrales de un bicomplejo \dotfill 29

11.4. Hipercohomología de un funtor respecto de un complejo $K^\bullet$
\dotfill 32

11.5. Paso al límite inductivo en hipercohomología \dotfill 35

11.6. Hiperhomología de un funtor respecto de un complejo $K_\bullet$
\dotfill 39

11.7. Hiperhomología de un funtor respecto de un bicomplejo
$K_{\bullet\bullet}$ \dotfill 41

11.8. Suplemento sobre la cohomología de los complejos simpliciales \dotfill 43

11.9. Un lema sobre complejos de tipo finito \dotfill 46

11.10. Característica de Euler--Poincaré de un complejo de módulos de longitud finita
\dotfill 48

\medskip
\S~12. Suplemento sobre la cohomología de haces \dotfill 49

12.1. Cohomología de haces de módulos sobre espacios anillados \dotfill 49

12.2. Imágenes directas superiores \dotfill 57

12.3. Suplementos sobre los funtores Ext de haces \dotfill 60

12.4. Hipercohomología del funtor de imagen directa \dotfill 62

\oldpage[III]{166}

\S~13. Límites proyectivos en álgebra homológica \dotfill 64

13.1. La condición de Mittag--Leffler \dotfill 64

13.2. La condición de Mittag--Leffler para grupos abelianos \dotfill 65

13.3. Aplicación: cohomología de un límite proyectivo de haces \dotfill 68

13.4. Condición de Mittag--Leffler y objetos graduados asociados a sistemas
proyectivos \dotfill 69

13.5. Límites proyectivos de sucesiones espectrales de complejos filtrados
\dotfill 71

13.6. Sucesión espectral de un funtor relativa a un objeto dotado de una
filtración finita \dotfill 73

13.7. Funtores derivados de un límite proyectivo de argumentos \dotfill 75

\medskip
\textsc{Capítulo III.} --- \textbf{Estudio cohomológico de los haces coherentes}
\dotfill 81

\medskip
\S~1. Cohomología de los esquemas afines \dotfill 82

1.1. Recordatorio sobre el complejo de álgebra exterior \dotfill 82

1.2. Cohomología de \v Cech de un recubrimiento abierto \dotfill 85

1.3. Cohomología de un esquema afín \dotfill 88

1.4. Aplicación a la cohomología de preesquemas cualesquiera \dotfill 89

\medskip
\S~2. Estudio cohomológico de los morfismos proyectivos \dotfill 95

2.1. Cálculos explícitos de ciertos grupos de cohomología \dotfill 95

2.2. El teorema fundamental de los morfismos proyectivos \dotfill 100

2.3. Aplicación a haces graduados de álgebras y módulos \dotfill 102

2.4. Una generalización del teorema fundamental \dotfill 107

2.5. Característica de Euler--Poincaré y polinomio de Hilbert \dotfill 109

2.6. Aplicación: criterios de amplitud \dotfill 111

\medskip
\S~3. Teorema de finitud para los morfismos propios \dotfill 115

3.1. El lema de dévissage \dotfill 115

3.2. El teorema de finitud: caso de los esquemas ordinarios \dotfill 116

3.3. Generalización del teorema de finitud (esquemas ordinarios) \dotfill 118

3.4. Teorema de finitud: caso de los esquemas formales \dotfill 119

\medskip
\S~4. El teorema fundamental de los morfismos propios. Aplicaciones \dotfill 122

4.1. El teorema fundamental \dotfill 122

4.2. Casos particulares y variantes \dotfill 129

4.3. Teorema de conexión de Zariski \dotfill 130

4.4. «Teorema principal» de Zariski \dotfill 135

4.5. Compleciones de módulos de homomorfismos \dotfill 138

4.6. Relaciones entre morfismos formales y morfismos ordinarios \dotfill 139

4.7. Un criterio de amplitud \dotfill 145

4.8. Morfismos finitos de preesquemas formales \dotfill 146

\oldpage[III]{167}

\S~5. Un teorema de existencia para los haces algebraicos coherentes \dotfill 149

5.1. Enunciado del teorema \dotfill 149

5.2. Demostración del teorema de existencia: caso proyectivo y cuasiproyectivo
\dotfill 151

5.3. Demostración del teorema de existencia: caso general \dotfill 154

5.4. Aplicación: comparación entre morfismos de preesquemas ordinarios y
morfismos de preesquemas formales. Preesquemas formales algebraizables \dotfill 156

5.5. Una descomposición de ciertos preesquemas \dotfill 159

\medskip
\textsc{Bibliografía (continuación)} \dotfill 161

\textsc{Índice de notaciones} \dotfill 162

\textsc{Índice terminológico} \dotfill 163

\begin{flushright}
\emph{Recibido el 28 de junio de 1961.}
\end{flushright}
\end{flushleft}


\clearpage
\phantomsection
\addcontentsline{toc}{section}{Referencias consolidadas de la edición}
\section*{Referencias consolidadas de la edición}
\bibliography{the}
\bibliographystyle{amsalpha}

\end{document}
